% SOURCE main.tex
\documentclass[a4paper,10pt]{article}
\usepackage[top=24mm,bottom=24mm,left=26mm,right=26mm,headheight=14pt]{geometry}
\usepackage{fontspec,amsmath,amsthm,unicode-math}
\usepackage[spanish,es-nodecimaldot,es-noquoting]{babel}
\usepackage{microtype,xcolor,graphicx,tikz,booktabs,array,needspace,fancyhdr}
\usepackage{longtable,tabularx,enumitem,float,import,xurl,multicol}
\usetikzlibrary{arrows.meta,calc,positioning,decorations.pathreplacing,matrix}
\usepackage[unicode,hidelinks]{hyperref}
\usepackage[nameinlink,noabbrev]{cleveref}
\crefname{theorem}{teorema}{teoremas}\Crefname{theorem}{Teorema}{Teoremas}
\crefname{proposition}{proposición}{proposiciones}\Crefname{proposition}{Proposición}{Proposiciones}
\crefname{lemma}{lema}{lemas}\Crefname{lemma}{Lema}{Lemas}
\crefname{section}{sección}{secciones}\Crefname{section}{Sección}{Secciones}
\crefname{equation}{ecuación}{ecuaciones}\Crefname{equation}{Ecuación}{Ecuaciones}
\setmainfont{STIXTwoText-Regular.otf}[BoldFont=STIXTwoText-Bold.otf,ItalicFont=STIXTwoText-Italic.otf,BoldItalicFont=STIXTwoText-BoldItalic.otf]
\setmathfont{STIXTwoMath-Regular.otf}
\newtheorem{theorem}{Teorema}[section]
\newtheorem{proposition}[theorem]{Proposición}
\newtheorem{lemma}[theorem]{Lema}
\newtheorem{corollary}[theorem]{Corolario}
\newtheorem{falsifier}[theorem]{Criterio de refutación}
\newtheorem{ablation}[theorem]{Prueba de necesidad estructural}
\newtheorem{criterion}[theorem]{Criterio}
\theoremstyle{definition}\newtheorem{definition}[theorem]{Definición}
\newtheorem{axiom}[theorem]{Axioma}
\newtheorem{principle}[theorem]{Principio}
\theoremstyle{remark}\newtheorem{remark}[theorem]{Observación}
\newcommand{\RR}{\mathbb R}\newcommand{\QQ}{\mathbb Q}
\newcommand{\ZZ}{\mathbb Z}\newcommand{\CC}{\mathbb C}
\newcommand{\Gr}{\operatorname{Gr}}\newcommand{\Res}{\operatorname{Res}}
\newcommand{\APP}{\mathrm{APP}}\newcommand{\TPK}{\mathrm{TPK}}
\newcommand{\TRIT}{\mathrm{TRIT}}\newcommand{\id}{\operatorname{id}}
\newcommand{\diag}{\operatorname{diag}}\newcommand{\conv}{\operatorname{conv}}
\newcommand{\FF}{\mathbb F}\newcommand{\Id}{\operatorname{id}}
\newcommand{\tr}{\operatorname{tr}}\newcommand{\dd}{\mathrm d}
\newcommand{\Span}{\operatorname{span}}\newcommand{\Cyl}{\operatorname{Cyl}}
\providecommand{\llbracket}{\lBrack}\providecommand{\rrbracket}{\rBrack}
\providecommand{\square}{\mdlgwhtsquare}
\providecommand{\dashrightarrow}{\mathrel{\rightdasharrow}}
\providecommand{\index}[1]{}
\newcommand{\HMT}{\mathrm{HMT}}\newcommand{\Tr}{\operatorname{Tr}}\newcommand{\Spec}{\operatorname{Spec}}
\newcommand{\Cstar}{C^{*}}\newcommand{\R}{\mathbb R}\newcommand{\Z}{\mathbb Z}\newcommand{\I}{\mathrm I}
\setcounter{tocdepth}{1}
\setlist{nosep,topsep=4pt}
\numberwithin{equation}{section}
\setlength{\parskip}{4pt plus 1pt}\setlength{\parindent}{1em}
\setlength{\emergencystretch}{2em}
\renewcommand{\normalsize}{\fontsize{10.5}{13.5}\selectfont}\normalsize
\clubpenalty=10000\widowpenalty=10000\displaywidowpenalty=10000
\AddToHook{cmd/section/before}{\Needspace{10\baselineskip}}
\AddToHook{cmd/subsection/before}{\Needspace{6\baselineskip}}
\AddToHook{env/theorem/before}{\Needspace{6\baselineskip}}
\AddToHook{env/proposition/before}{\Needspace{6\baselineskip}}
\pagestyle{fancy}\fancyhf{}
\fancyhead[L]{\footnotesize Grassmannianos positivos y amplituedros: formas canónicas y refinamiento}
\fancyfoot[C]{\thepage}\renewcommand{\headrulewidth}{.2pt}
\hypersetup{pdftitle={Grassmannianos positivos y amplituedros: formas canónicas y refinamiento},pdfsubject={Dualidad reticular y brecha de masa de Yang–Mills en Holografía Modular Triádica},pdfauthor={Oumar Haidara Fall; Rubén Ramos Balsa}}
\makeatletter
\newcommand{\compactcontents}{%
  \begingroup
  \fontsize{9.2}{11.1}\selectfont
  \setlength{\parskip}{0pt}\setlength{\columnsep}{7mm}
  \renewcommand{\l@section}{\@dottedtocline{1}{0pt}{1.8em}}
  \renewcommand{\l@part}[2]{\addpenalty{-\@highpenalty}\addvspace{6pt}%
    {\bfseries\noindent ##1\nobreak\hfill\nobreak ##2\par}\nobreak}
  \section*{Índice}
  \begin{multicols}{2}\@starttoc{toc}\end{multicols}
  \endgroup}
\makeatother
\begin{document}
\begin{center}
{\small Geometría positiva y física matemática}\par\vspace{6mm}
{\LARGE\bfseries Grassmannianos positivos y amplituedros:\\formas canónicas y refinamiento}\par\vspace{3mm}
{\large Dualidad reticular y brecha de masa de Yang--Mills\\en Holografía Modular Triádica}\par\vspace{5mm}
Oumar Haidara Fall \textperiodcentered\ Rubén Ramos Balsa\par
{\fontsize{8}{10}\selectfont Coautoría sin prelación de contribución\par
Integración y recopilación documental generadas por ChatGPT - Astra.}
\end{center}
\begin{abstract}
La Holografía Modular Triádica construye la estructura discreta conjunta
del continuo mediante APP, TRIT y el transporte con memoria del TPK.
La reducción compatible de cilindros produce la realización aritmética,
mientras la proyección conjugada conserva la incidencia excepcional.
Sobre esta genealogía se demuestra una familia de realizaciones positivas
recuperables. La carga uniforme y dos memorias nonádicas determinan el
registro de doce coordenadas; sus separaciones ordenadas producen una
medición planar de frontera y una representación grassmanniana con
inversa marcada. Los lectores que factorizan por el registro se transportan
a esa representación y se conservan bajo cambios de coordenadas del mismo
punto. El refinamiento nonádico construye banderas positivas y puntos
amplituhedrales en todo rango y dimensión externa finitos admisibles.
En rango uno se obtiene cobertura simplicial completa y cancelación de
residuos interiores en toda dimensión finita. Las formas con valores en
la recta de acción satisfacen una condición exacta de pegado en las
fronteras comunes. Los promedios celulares conservan una medida
cronológica y descomponen el incremento de los momentos en detalles
ortogonales; su límite determina un operador de Jacobi cuyas compresiones
finitas tienen espectro simple, residuos positivos y resolventes en
fracción continua.
La reducción de Schur conserva la energía efectiva de la partición,
con reconstrucción del detalle y realización límite de Sobolev.
El flujo del registro induce además una deformación reticular de
veinticuatro dimensiones que conserva volumen y simetría ternaria,
y transforma la inversión del parámetro en dualidad reticular y
reciprocidad theta. Se reúnen los Hamiltonianos elíptico, hiperbólico
y parabólico, sus cartas simplécticas, términos de conexión y
transformaciones de Legendre, junto con los lectores materiales
correspondientes. El desarrollo de Yang--Mills incorpora la medida
común, la escala areal--volumétrica, la coercividad uniforme, el
testigo espectral persistente y la reconstrucción del Hamiltoniano.
La articulación de estos resultados distingue y relaciona las
positividades de Plücker, de los momentos, de las constitutivas
y del espectro mediante operaciones explícitas sobre el mismo
estado enriquecido.
\end{abstract}
\clearpage
\compactcontents
\clearpage

\section{Introducción: genealogía de la forma y conservación del refinamiento}
\label{sec:genealogy}
\input{sections/introduction_addition.tex}
\subsection{Genealogía aritmética y sistema de realizaciones geométricas}
La construcción empieza por las dos hojas aritméticas de APP. Si
\(p,q\in\{1,\ldots,9\}\), cada operación conserva una división completa,
\[
 p+q=R_\Sigma+9q_\Sigma,\qquad
 pq=R_\Pi+9q_\Pi.
\]
TRIT añade el régimen local y la orientación. El selector trítico de fase
operativa elige la operación compatible con ese estado, mientras ambas hojas
permanecen en el soporte. Sobre el espacio de estados enriquecidos, el
Topological Phase Kernel compone selección, transporte y actualización:
\[
 U_t=\operatorname{Upd}_t\circ\operatorname{Tra}_t\circ\operatorname{Sel}_t.
\]
El registro conserva los residuos, cocientes, acarreos, hojas, rutas,
fronteras, memoria y supervivencia que utiliza cada prolongación. La
holonomía nonádica \(\Gamma_9=U_{t+8}\circ\cdots\circ U_t\) pertenece
a esta dinámica enriquecida. Su publicación reducida de memoria se efectúa
después. El retorno de fase conserva el avance del registro y por ello
distingue estados situados a profundidades diferentes.

La prolongación
\[
 w_6\longrightarrow w_{12}\longrightarrow w_{18}\longrightarrow
 w_{24}\longrightarrow w_{30}\longrightarrow R_{36}\longrightarrow G_9
\]
organiza los refinamientos compatibles del mismo estado. Las cinco
construcciones del continuo emergen conjuntamente, con sus aspectos
espectral, solenoidal, gauge, cohomológico y determinantal. Los lectores
de ese estado publican correlativamente \(\pi,\varphi,e,\alpha\).
Su origen común admite un orden interno de evaluación: las coordenadas
regionales ya generadas actúan junto con el registro dodecafásico en los
lectores de clausura de \(\alpha\). En este artículo, Grassmannianos,
polígonos y formas diferenciales son realizaciones posteriores de datos
generados por esta cadena. Sus parámetros quedan fijados antes de toda
comparación geométrica o metrológica.

El problema concreto consiste en describir lo que conserva una realización
positiva y lo que puede variar al cambiar su forma. Una misma lista positiva
admite varios usos matemáticos: pesos de aristas, longitudes de intervalos,
coeficientes de una forma cuadrática o coordenadas de un lector material.
Cada uso posee un dominio, una operación y unos invariantes propios. La
expresión \emph{multimorfología} designa aquí el sistema de estas realizaciones
con sus mapas explícitos de compatibilidad. Una identidad de sus números
requiere una comprobación menor que una identidad de sus operadores; el
desarrollo siguiente formula y demuestra cada nivel por separado.

La base de procedencia es el artículo de generación coinductiva de las
constantes y estructura algebraica del electrón, junto con el desarrollo
de extrema y media razón holográfica. Sus construcciones anteriores fijan
el registro utilizado abajo. Los resultados geométricos de este texto comienzan
con esa salida finita y prueban sus consecuencias mediante datos explícitos. Las
proposiciones de transferencia material declaran adicionalmente los
lectores que emplean. Esta organización preserva el alcance específico de
cada prueba.

\clearpage
\part{Construcción aritmética, estado enriquecido y continuo conjunto}
\begingroup
\sloppy
\renewcommand{\TRIT}{\{-1,0,+1\}}
\import{foundation/}{sections/nucleo.tex}
\import{foundation/}{sections/extension.tex}
\import{foundation/}{sections/continuo_conjunto.tex}
\import{foundation/}{sections/generacion.tex}
\import{foundation/}{sections/registro_k.tex}
\import{foundation/}{sections/alpha.tex}
\import{foundation/}{sections/excepcional.tex}
\import{foundation/}{sections/01_hilbert_nonadico.tex}
\endgroup
\clearpage
\part{Geometría positiva, refinamiento y dualidad reticular}
\input{sections/trit_projection.tex}
\input{sections/incidence_transport.tex}

\section{Extracción dodecafásica y proyección isotípica de la dirección}
\label{sec:register}
Antes de evaluar números, los lectores distinguen regiones dinámicas.
La simetría especular del corpus actúa sobre el bloque
\(B=(u^\pi,u^e,u^\varphi)\in(\mathbb F_3^6)^3\) por
\[
 \mathfrak cB=(u^e,u^\pi,u^\varphi).
\]
Con \(q=u^\pi+u^e+u^\varphi\),
\(a=u^\pi+u^e-u^\varphi\) y \(d=u^\pi-u^e\), las fórmulas
\[
 u^\varphi=2(q-a),\quad s=q-u^\varphi,\quad
 u^\pi=2(s+d),\quad u^e=2(s-d)
\]
son identidades componente a componente en \(\mathbb F_3\). El espejo
conserva el eje y el agregado en una fibra e invierte el reparto
antisimétrico. El transporte nonádico hace evolucionar la fibra misma.
La quinta fase tiene dato de frontera \(111111\), firma residual
\(000\), una solución local y una extensión: la marca de
\(\varphi\) en cinco identifica el eje regional de esa operación.
En las fases seis, siete y ocho los agregados regionales son
\(012100\), \(101001\) y \(112000\), respectivamente. Estos datos
pertenecen al dominio ternario; los lectores de historias compatibles
producen después sus evaluaciones reales.

El registro de doce ventanas y los bloques regionales conservan así
sus dominios y su enlace en el mismo estado. La realización positiva
que sigue utiliza el registro publicado, con su orden y su escala.
Para transportar el espejo regional a una deformación geométrica debe
componerse primero su acción en ese registro; la igualdad de genealogía
por sí sola expresa menos que la conmutatividad de ese cuadrado. Esta
precisión permite conservar la dinámica regional completa al estudiar
una de sus realizaciones finitas.

El calendario divide las ciento ocho posiciones en doce ventanas nonádicas.
La extracción terminal tiene la composición
\[
 x_{\rm term}\longmapsto\mathcal R_{12}(x_{\rm term})
 \longmapsto(b^{90},b^{120},Q)\longmapsto U_{\rm sgn}\longmapsto K.
\]
Las dos primeras flechas conservan la historia orientada en los lectores
del corpus. Para mostrar íntegramente la reconstrucción finita que utiliza
este artículo, sus salidas enteras son
\begin{align*}
 b^{90}&=(-495,-116,-684,108,601,-90,-173,-88,313,560,-397,461),\\
 b^{120}&=(-425,-281,-481,671,-255,30,475,-543,680,251,6,-128),
 \qquad Q=6263.
\end{align*}
Se fija el orden cíclico de doce posiciones. Sea
\(B_r=\sum_{j=0}^3 b^{120}_{r+3j}\) y
\(d_{r,j}=b^{90}_{r+3j}\), para \(r=1,2,3\). Entonces
\[
 s_1=\frac{Q+2B_1+B_2}{3},\quad s_2=s_1-B_1,\quad s_3=s_2-B_2,
 \qquad
 U^{(r)}=(s_r,d_{r,1}-d_{r,3},d_{r,0}+d_{r,2},d_{r,0}-d_{r,2}).
\]
La evaluación da
\[
 U^{(1)}=(2378,-452,-668,-322),\quad
 U^{(2)}=(1406,998,-204,-28),\quad
 U^{(3)}=(2479,-551,-371,-997).
\]
La matriz
\[
 H_4=\begin{pmatrix}1&1&1&1\\1&1&-1&-1\\1&-1&1&-1\\1&-1&-1&1\end{pmatrix},
 \qquad H_4^2=4I_4,
\]
reconstruye cada órbita mediante \(H_4U^{(r)}/4\). Restituir el orden
de las posiciones produce
\begin{equation}
 K=(234,543,140,729,659,824,621,58,914,794,146,601),
 \qquad Q=\sum_{j=1}^{12}K_j=6263.
 \label{eq:K-main}
\end{equation}
Cada componente es positiva. La división por tres procede de las tres
órbitas y la división por cuatro de la inversa de Hadamard. El control
directo adicional es \((I-S_3)K=b^{90}\),
\((I-S_4)K=b^{120}\), donde \((S_aK)_j=K_{j+a}\) con índices cíclicos.
Estas identidades hacen reproducible la reconstrucción sin seleccionar
coeficientes mediante una geometría objetivo.

La carta incidencial del corpus realiza las doce posiciones como cosets
de \(A_5/C_5\). Si \(\rho\) es esa representación y \(\chi_3\) su
carácter tridimensional real, su proyector es
\[
 P_3=\frac3{60}\sum_{a\in A_5}\chi_3(a^{-1})\rho(a),\qquad
 P_{11}=I-\frac1{12}\mathbf1\mathbf1^{\mathsf T}.
\]
La enumeración y la orientación de cosets se conservan en la fuente de
esta carta. Para que la extensión geométrica sea evaluable directamente,
se escribe su salida exacta \(u=P_3P_{11}K\), con \(r=\sqrt5\):
\begin{equation}
\begin{split}
u=\bigl(&-495/4-233r/10,\quad403/4+169r/10,\\
&-403/4-169r/10,\quad495/4+233r/10,\\
&601/4+1031r/10,\quad203/4+211r/5,\\
&-203/4-211r/5,\quad-601/4-1031r/10,\\
&313/4+569r/10,\quad162+219r/20,\\
&-162-219r/20,\quad-313/4-569r/10\bigr).
\end{split}
\label{eq:u-main}
\end{equation}
Se comprueba por suma y producto en \(\QQ(\sqrt5)\) que
\[
 \mathbf1^{\mathsf T}u=0,\qquad
 \|u\|^2=\frac{6638585+2275584\sqrt5}{20}>0.
\]
En particular, sus coordenadas son no constantes. La dirección \(u\),
el registro íntegro \(K\) y el proyector
\(P_{10}=P_{11}-uu^{\mathsf T}/\|u\|^2\) conservan diferentes cantidades
de información. Esta distinción permite utilizar la dirección para un
flujo sin sustituir por ella el registro que emplean los lectores de acción.

\input{sections/memory_geometry.tex}
\input{appendices/network_cluster.tex}
\input{appendices/polygon_amplituhedron.tex}
\input{sections/higher_rank.tex}
\input{sections/refinamiento_dimensional_momentos.tex}
\input{sections/jacobi_spectral.tex}
\input{appendices/figures.tex}
\input{appendices/convex_refinement.tex}
\input{sections/energy_refinement.tex}
\input{sections/energy_limit.tex}
\input{sections/nonadic_transport.tex}
\input{sections/leech_foundations.tex}
\input{sections/leech_dualidad.tex}
\input{appendices/mass_bridge.tex}
\input{sections/lectores_formas_compatibles.tex}
\input{appendices/multimorphology_figure.tex}

\clearpage
\part{Realización gauge, refinamiento y brecha espectral}
\input{sections/ym_bridge.tex}
\input{sections/ym_scale.tex}
\input{sections/ym_construction.tex}
\input{sections/ym_gauge_object.tex}
\input{sections/ym_complete.tex}
\input{sections/hamiltonianos_curvatura.tex}

\section{Discusión: reconstrucción, invariantes y dualidad}
\label{sec:discussion}
El resultado estructural prioritario es la recuperabilidad de la forma.
La carga uniforme y las dos memorias nonádicas reconstruyen el registro
mediante una inversa explícita; el cono positivo describe exactamente
el dominio donde esa reconstrucción induce menores estrictamente positivos.
La cota de estabilidad cuantifica la persistencia de ese dominio. Por
tanto, los pesos geométricos admiten una procedencia y una recuperación
operatorias, además de su evaluación numérica.

La composición desarrollada tiene una secuencia explícita. El registro
terminal genera doce pesos positivos; sus sumas parciales dan trece puntos
ordenados; una red planar realiza sus menores; la curva de momentos
construye polígonos y politopos; sus triangulaciones proporcionan cobertura
y formas con residuos orientados. El refinamiento de una celda fija conserva
la suma de sus formas. La introducción de nuevos vértices en la curva de
momentos produce, por su parte, otra geometría y exige comparar dominios.

La dirección algebraica del registro distingue dos operaciones. El
reescalado positivo de columnas conserva las razones proyectivas pertinentes.
Su acción normalizada sobre los intervalos genera en cambio una deformación
efectiva, certificada por una derivada exacta no nula. Esta deformación admite
un potencial convexo cuya transformada de Legendre se conserva al subdividir
con dirección heredada. La invariancia tiene así un mecanismo comprobable:
cada hijo conserva el generador de su padre y la suma de sus pesos recupera
el peso anterior.

El refinamiento energético proporciona una segunda conservación exacta.
La reducción de Schur recupera el laplaciano efectivo y la prolongación
minimizante separa cada configuración en su parte observada y un detalle
reconstructible. Las energías de los detalles se suman ortogonalmente.
Con mallas decrecientes y energía uniformemente acotada, esa sucesión
converge en la realización de Sobolev declarada. La invariancia de la
forma canónica, la identidad del potencial convexo y la conservación de
energía son tres resultados distintos de compatibilidad, demostrados
para los refinamientos especificados.

La prolongación a los doce planos reticulares conserva una simetría
ternaria y el volumen. El resultado decisivo es la reciprocidad
\(\Lambda_s^*=\Lambda_{-s}\), que induce la inversión
\(\Theta_s(t)=t^{-12}\Theta_{-s}(t^{-1})\). La realización partida de
signatura \((24,24)\) conserva integralidad, paridad y autodualidad para
todo el flujo. En la proyección euclídea, la variación explícita de una
norma distingue deformación y automorfismo de la red inicial. Esta
distinción determina el contenido exacto de la conexión con la dualidad
de momentos y devanados.

La realización material conserva sus propios lectores. La publicación de
acción y el carácter de masas reciben el registro por composiciones ya
construidas; la congruencia preserva positividad e inercia, mientras la
métrica determina la interpretación de sus autovalores. Estas diferencias
precisan el significado de multimorfología: varias formas matemáticas
proceden de una genealogía común y se relacionan por operaciones verificables.
La identificación completa de amplituedros de rango superior requerirá
las pruebas de cobertura y residuos para esos mapas concretos. El resultado
de rango uno demostrado aquí fija un caso completo, con coordenadas
racionales, triangulación y controles reproducibles.

La orientación temporal del TRIT integra el retorno con memoria, la
actualización de umbral y las prolongaciones conjugadas. Sus generadores
elíptico, nilpotente e hiperbólico proporcionan leyes distintas de
evolución; la composición del TPK conserva la genealogía que las
concatena. En particular, la actividad del régimen cero queda expresada
por un transporte unipotente, y la doble proyección conserva simultáneamente
la convergencia de las coordenadas y la incidencia excepcional. La tesis
geométrica del artículo reside en esta articulación: la estructura de
partida genera registros recuperables y éstos sostienen realizaciones
positivas cuyas transformaciones, fronteras y memorias se pueden calcular
y comparar con precisión.

El entrelazamiento \(B_FM=MB_V\) añade un transporte exacto desde el
bloque modular hasta el cociente incidencial. En él, las multiplicidades
uno y tres se deducen de la matriz de Gram, y la forma firmada declarada
se expresa como \(J_L=(7I_Q-\overline B_F)/2\). Las matrices de
Vandermonde y de Hankel aportan, por otra parte, realizaciones explícitas
\(Y^{(k)}\) para \(2\le k\le9\), con positividad y rango probados
en la partición inicial. La prolongación nonádica extiende después
esa construcción a todo par finito admisible de rango y dimensión
externa, mediante los sistemas de nodos y momentos desarrollados aquí.
Estos resultados amplían el sistema de realizaciones sin identificar
la dimensión de un cociente, el rango de una matriz y el número de
generadores de una superálgebra.

En la realización gauge, la escala areal--volumétrica precede al
Hamiltoniano. El generador producto conserva todos los sectores de la
descomposición ortogonal, mientras su restricción física contiene un
testigo que alcanza el primer nivel positivo. Los entrelazadores de
refinamiento conservan la coercividad y ese testigo. El desarrollo
continuo reúne además la medida común, la curvatura, la continuidad
euclídea, las formas de reflexión y el Hamiltoniano reconstruido. Así,
el argumento de brecha queda situado junto a la construcción que le da
dominio y significado, y la positividad de Plücker mantiene su realización
geométrica propia. Su articulación exige esos transportes explícitos;
la palabra positividad designa propiedades tipadas que el estado común
permite comparar sin sustituir unas por otras.

\clearpage
\part{Comprobación exacta y enumeración de las emisiones}
\input{sections/exact_reproduction.tex}
\import{foundation/}{libro/censo_firmas_emision.tex}
\input{sections/conclusions.tex}

\begin{thebibliography}{99}
\bibitem{hmtI} O. Haidara Fall y R. Ramos Balsa,
\emph{Generación coinductiva de \(\pi,\varphi,e\) a profundidad arbitraria,
derivación de la constante de estructura fina y estructura algebraica del
electrón}. Serie HMT--MD, artículo I, edición consultada de septiembre de 2026.
\bibitem{hmtX} O. Haidara Fall y R. Ramos Balsa,
\emph{Extrema y media razón holográfica}. Desarrollo del acoplamiento
aritmético--geométrico, registro dodecafásico y dirección dimensional,
edición consultada de septiembre de 2026.
\bibitem{hmtVI} O. Haidara Fall y R. Ramos Balsa,
\emph{Partículas, persistencia y espectro de masas}. Serie HMT--MD, artículo VI,
edición consultada de septiembre de 2026.
\bibitem{HMT-IV} O. Haidara Fall y R. Ramos Balsa,
\emph{Moonshine, dualidad T y teoría M desde la estructura discreta del
continuo}. Serie HMT--MD, artículo IV, edición consultada de septiembre de 2026.
\bibitem{HMT-VIII} O. Haidara Fall y R. Ramos Balsa,
\emph{Aritmética genealógica de los números primos y estructura espectral
de la función zeta}. Serie HMT--MD, artículo VIII, edición consultada de septiembre de 2026.
\bibitem{HMT-IX} O. Haidara Fall y R. Ramos Balsa,
\emph{Resolución de la hipótesis del continuo en la clausura genealógica
holográfica}. Serie HMT--MD, artículo IX, edición consultada de septiembre de 2026.
\bibitem{Elkies} N. D. Elkies,
\emph{Rational lattices and their theta functions}. Arizona Winter School,
2009, teorema 2.
\href{https://people.math.harvard.edu/~elkies/aws09.pdf}{Notas del autor}.
\bibitem{ConwaySloane} J. H. Conway y N. J. A. Sloane,
\emph{Sphere Packings, Lattices and Groups}, tercera edición, Springer, 1999.
\href{https://doi.org/10.1007/978-1-4757-6568-7}{DOI: 10.1007/978-1-4757-6568-7}.
\bibitem{postnikov} A. Postnikov,
\emph{Total positivity, Grassmannians, and networks}, 2006,
\href{https://arxiv.org/abs/math/0609764}{arXiv:math/0609764}.
\bibitem{scott} J. S. Scott,
\emph{Grassmannians and Cluster Algebras}, 2003,
\href{https://arxiv.org/abs/math/0311148}{arXiv:math/0311148}.
\bibitem{amplituhedron} N. Arkani-Hamed y J. Trnka,
\emph{The Amplituhedron}, 2013,
\href{https://arxiv.org/abs/1312.2007}{arXiv:1312.2007}.
\bibitem{positive} N. Arkani-Hamed, Y. Bai y T. Lam,
\emph{Positive Geometries and Canonical Forms}, 2017,
\href{https://arxiv.org/abs/1703.04541}{arXiv:1703.04541}.
\bibitem{HMT-YM} O. Haidara Fall y R. Ramos Balsa,
\emph{El cierre holográfico del infinito}, edición integral de 2.249 páginas,
capítulo 99, «Imagen gauge y composición hacia Yang--Mills», pp. 1742--1762;
y \emph{Yang--Mills: medida común, reconstrucción y brecha espectral},
monografía especializada, desarrollo de la escala areal--volumétrica.
\bibitem{hmtIntegral} O. Haidara Fall y R. Ramos Balsa,
\emph{El cierre holográfico del infinito}, edición integral de 2.249 páginas,
septiembre de 2026. Desarrollos de los regímenes cuadráticos, el vacío
electromagnético, la realización electrodébil y las correspondencias
espirales de Mecánica Dimensional.
\input{sections/bibliography_additions.tex}
\end{thebibliography}
\end{document}


% SOURCE sections/introduction_addition.tex
La cuestión rectora de este trabajo es la conservación de una estructura
generada cuando cambia su realización matemática. La Holografía Modular
Triádica produce, mediante operaciones aritméticas levantadas, estados que
retienen simultáneamente orientación, incidencia, frontera y memoria. Su
estructura discreta del continuo reúne las prolongaciones compatibles de
esos estados. La realización arquimediana publica coordenadas mediante el
refinamiento de cilindros; la realización conjugada conserva relaciones
excepcionales de incidencia. La genealogía precede, por tanto, a los valores
y a las formas en que se los representa. Ésta es la tesis de partida que el
desarrollo formal siguiente expone con sus operaciones y dependencias.

El problema geométrico propio del artículo comienza cuando una salida
positiva de aquella construcción se convierte en separación entre nodos,
peso de una red, coordenada de Plücker, momento de una medida, coeficiente
de energía o dirección de deformación reticular. Para justificar ese paso
se necesita determinar qué información recupera cada representación, cómo
se compara con las restantes y qué conserva su refinamiento. El interés
de una realización positiva aumenta de forma decisiva cuando admite una
inversa marcada: la geometría deja entonces de ofrecer sólo una evaluación
del estado y pasa a conservar la información que utilizan sus lectores.
La marca comprende el orden de las posiciones, la escala total y los datos
enriquecidos que el lector particular requiere.

La aportación central se organiza alrededor de tres operaciones. La primera
es la reconstrucción: la carga y la memoria recuperan las coordenadas
del registro, y los cocientes de menores recuperan sus separaciones.
La segunda es el refinamiento compatible: una partición se subdivide
conservando sus extremos, sus marcas y un archivo exacto del detalle.
La tercera es el transporte de estructuras: una forma, un operador o un
lector se lleva a otra representación mediante un mapa explícito. Estas
operaciones permiten demostrar conservación de residuos, energía efectiva,
momentos y espectros en los dominios concretos donde cada identidad tiene
sentido. El calificativo positivo adquiere así varias realizaciones
matemáticas precisas, articuladas por la reconstrucción del mismo dato.

La geometría positiva posee antecedentes propios que permiten identificar
con precisión el nivel de esta contribución. Las mediciones de frontera
de redes dirigidas parametrizan celdas del Grassmanniano totalmente no
negativo en el trabajo de Postnikov~\cite{postnikov}; Scott describe la
estructura de álgebra de conglomerados del anillo de coordenadas
grassmanniano~\cite{scott}. Arkani-Hamed y Trnka introducen el amplituedro
en la descripción geométrica de amplitudes de la teoría supersimétrica
\(\mathcal N=4\) de Yang--Mills en el límite planar~\cite{amplituhedron}.
La formulación de geometrías positivas caracteriza después sus formas
canónicas por polos logarítmicos y residuos recursivos de
frontera~\cite{positive}. El presente desarrollo compone esos lenguajes
de realización con datos previamente generados por HMT. Su resultado
diferencial es una familia marcada, recuperable y refinable, junto con
los mapas que enlazan su forma positiva y sus lecturas materiales.

El refinamiento amplía además el alcance dimensional. El número de nodos
crece mediante subdivisiones nonádicas y permite construir matrices
positivas de todo rango y dimensión externa finitos compatibles con ese
número. La prueba se apoya en determinantes de Vandermonde y formas de
momentos de rango completo. En rango uno se obtiene la geometría convexa
completa, con cobertura por símplices y residuos en toda dimensión finita.
Para rangos superiores se construye una familia explícita de imágenes
amplituhedrales y se preservan sus banderas marcadas. La cobertura global
de un amplituedro de rango superior constituye un enunciado de alcance
distinto del teorema de existencia y compatibilidad de esta familia; las
pruebas posteriores conservan esa distinción en su formulación.

La dimensión física se introduce mediante los operadores correspondientes.
El calendario trítico distingue rotación elíptica, apertura hiperbólica y
transporte parabólico de memoria. Sus realizaciones hamiltonianas conservan
la forma simpléctica y permiten calcular la transformación de Legendre y
los términos de conexión originados por una carta variable. En el sector
gauge, la escala areal--volumétrica, la medida común y las formas de
reflexión anteceden a la reconstrucción del Hamiltoniano de Yang--Mills.
El argumento espectral se presenta junto a la coercividad uniforme, el
testigo que alcanza el primer nivel positivo y los entrelazadores que lo
preservan. La relación con las formas canónicas se establece mediante sus
lectores y sus residuos, mientras la brecha pertenece al Hamiltoniano y a
su espacio físico. Esta jerarquía sitúa cada resultado en el lugar donde
su demostración produce efectivamente la conclusión.

La exposición incorpora primero la construcción formal necesaria para
obtener el estado y su registro. Desarrolla después las realizaciones
positivas, la extensión dimensional, la energía y la dualidad reticular;
reúne finalmente las realizaciones hamiltonianas y el desarrollo gauge.
Los programas de comprobación acompañan a identidades determinadas y a
controles negativos explícitos. La demostración simbólica conserva el
alcance general; la ejecución exacta comprueba las especializaciones
declaradas y su correspondencia con las fórmulas impresas.


% SOURCE foundation/sections/nucleo.tex
\section{Núcleo formal de la Holografía Modular Triádica}
\label{sec:nucleo}

La Holografía Modular Triádica, abreviada HMT, comienza con una estructura finita de posiciones y operaciones. El orden de construcción importa: las posiciones reciben evaluaciones aritméticas, las transiciones adquieren una firma ternaria y la dinámica conserva la información necesaria para prolongarlas. Sólo después se forman las realizaciones numéricas. Una constante reconocida al final de este proceso no sustituye a la regla que produjo su región ni a la historia que determina su evaluación.

La denominación \emph{All Possible Paths}, APP, expresa el papel de los caminos sobre el soporte. Se adopta como denominación de la construcción, con una referencia conceptual a la formulación por caminos de la mecánica cuántica; no supone que su ley aritmética sea la integral de caminos de Feynman. El \emph{Topological Phase Kernel}, TPK,\footnote{Los autores dedican asimismo las iniciales TPK a Tan Pengkiang. Esta dedicación es independiente de la definición matemática.} designa la composición de selección, transporte, actualización de memoria y prolongación. Las siglas se mantienen únicamente después de haber definido los objetos que representan.

\subsection{APP: soporte posicional y evaluaciones aritméticas}

Consideremos las nueve marcas interiores de una recta graduada entre cero y diez,
\[
 D_9=\{1,2,3,4,5,6,7,8,9\}.
\]
Su disposición horizontal y vertical produce las ochenta y una posiciones de \(D_9^2\). Las marcas interiores no deben confundirse con los diez intervalos unitarios de la recta de referencia. Esta disposición fija un alfabeto y dos direcciones; todavía no utiliza una expansión real para decidir una trayectoria.

La identificación cíclica de las posiciones conduce al soporte
\[
 X_9=(\ZZ/9\ZZ)^2.
\]
La carta de representantes positivos conserva el símbolo nueve para la clase nula. Para todo entero positivo se definen
\begin{equation}
 \rho_9(n)=1+((n-1)\bmod9),\qquad
 q_9(n)=\frac{n-\rho_9(n)}9,
 \qquad n=\rho_9(n)+9q_9(n).
 \label{eq:residuo-cociente}
\end{equation}
La raíz digital positiva es, por tanto, una elección de representante de la clase residual. No conserva por sí sola el entero anterior a la reducción; la coordenada \(q_9\) proporciona precisamente la información que falta.

En cada posición \((i,j)\), la acumulación de \(i\) repetida \(j\) veces produce \(ij\). Las dos evaluaciones enteras y sus proyecciones son
\[
 S(i,j)=i+j,\qquad P(i,j)=ij,\qquad
 \Sigma(i,j)=\rho_9(i+j),\quad \Pi(i,j)=\rho_9(ij).
\]
La simetría \((i,j)\mapsto(j,i)\) intercambia las direcciones de acumulación y conserva ambas evaluaciones. Es una propiedad de un único soporte, no una identificación posterior entre dos matrices independientes.

\begin{definition}[Evaluación aritmética levantada]
La evaluación de APP que conserva las dos hojas es
\begin{equation}
 \mathcal A(i,j)=
 \bigl((R^+_{ij},q^+_{ij}),(R^\times_{ij},q^\times_{ij})\bigr)
 =\bigl((\rho_9(i+j),q_9(i+j)),(\rho_9(ij),q_9(ij))\bigr).
 \label{eq:app-levantada}
\end{equation}
El término \emph{hoja} indica aquí una evaluación sobre el mismo soporte: aditiva o multiplicativa. Las dos coordenadas de cociente pertenecen al estado aritmético y no se descartan al reducir los residuos.
\end{definition}

\input{figures/app.tex}

\begin{proposition}[Reconstrucción de las evaluaciones]
Los dos pares de \eqref{eq:app-levantada} determinan exactamente \(i+j\) e \(ij\). Los residuos aislados no determinan esas evaluaciones enteras ni sus cocientes.
\end{proposition}
\begin{proof}
La primera afirmación es \eqref{eq:residuo-cociente} aplicada a cada hoja. Para la segunda, \(10=1+9\) y \(19=1+2\cdot9\) tienen el mismo residuo y distinto cociente. En el soporte bidimensional, las posiciones \((5,5)\) y \((8,2)\) dan la misma suma y el mismo residuo producto, pero productos \(25=7+18\) y \(16=7+9\). La proyección residual identifica datos que la evaluación levantada distingue.
\end{proof}

\subsubsection{Grafo de Cayley, topología toroidal y grado de enrollamiento}

Las traslaciones cardinales están dadas, en coordenadas de fila y columna, por
\[
 \nu_N=(-1,0),\quad\nu_E=(0,1),\quad
 \nu_S=(1,0),\quad\nu_O=(0,-1),\qquad T_d(x)=x+\nu_d.
\]
La primera coordenada aumenta hacia el sur y la segunda hacia el este. El grafo
\[
 \mathscr G_9=\operatorname{Cay}(X_9,\{\nu_N,\nu_E,\nu_S,\nu_O\})
\]
es la retícula cíclica bidimensional: tiene ochenta y un vértices y cuatro aristas orientadas salientes por vértice. Es un grafo de Cayley para la estructura \emph{aditiva} del soporte. La multiplicación de residuos tiene divisores de cero y no convierte todas las filas de APP en un grupo multiplicativo.

Una palabra cardinal \(w=d_1\cdots d_r\) y un origen \(x_0\) determinan
\[
 x_k=x_0+\sum_{j=1}^k\nu_{d_j}\pmod9.
\]
Existen \(81\cdot4^r\) caminos orientados de longitud \(r\), incluidos retornos y retrocesos. Cuando el camino se cierra, su desplazamiento entero previo a la reducción define
\begin{equation}
 \operatorname{wind}(w)=\frac19
 \bigl(\#S(w)-\#N(w),\#E(w)-\#O(w)\bigr)\in\ZZ^2.
\end{equation}
La inversión del camino cambia el signo de este vector. El retorno a la misma posición residual no obliga, por tanto, a que el enrollamiento sea nulo. Éste es el primer ejemplo de una diferencia que reaparecerá en el retorno de fase: una coordenada puede cerrarse mientras otra registra una evolución no trivial.

\subsubsection{Cociente de transporte e identidad de cociclo}

Sea \(s_9:\ZZ/9\ZZ\to D_9\) la sección de representantes positivos. El cociente de transporte asociado a la composición aditiva, denominado también \emph{acarreo} en aritmética posicional, es
\[
 c_9(a,b)=\frac{s_9(a)+s_9(b)-s_9(a+b)}9.
\]
El carácter composicional de esta memoria queda expresado por una identidad exacta.

\begin{lemma}[Identidad de cociclo]
Para \(a,b,c\in\ZZ/9\ZZ\),
\[
 c_9(a,b)+c_9(a+b,c)=c_9(b,c)+c_9(a,b+c).
\]
\end{lemma}
\begin{proof}
Ambas sumas son iguales a
\[
 \frac{s_9(a)+s_9(b)+s_9(c)-s_9(a+b+c)}9.
\]
La asociatividad de la suma de representantes levantados requiere precisamente este término adicional al componer en el cociente.
\end{proof}

Con representantes positivos, este cociclo no está normalizado en la identidad: \(c_9(0,a)=1\). Si \(s_0\) toma representantes en \(\{0,\ldots,8\}\), su acarreo normalizado \(c_0\) satisface
\[
 c_9(a,b)=c_0(a,b)+\mathbf1_{a=0}+\mathbf1_{b=0}-\mathbf1_{a+b=0}.
\]
La diferencia es el cambio de sección; no altera la identidad de cociclo. La descomposición posterior del calendario en fase y ventana utilizará los representantes de cero a ocho.

Los totales de ambas evaluaciones distinguen la contribución visible y la conservada en el cociente:
\[
 \sum S=810,\quad\sum P=2025,\quad
 \sum\Sigma=405,\quad\sum\Pi=459,
\]
\[
 \sum q^+=45,\qquad\sum q^\times=174.
\]
De aquí se obtiene
\begin{equation}
 2025-810=(459-405)+9(174-45)=54+9\cdot129.
 \label{eq:defecto-integrado}
\end{equation}
La diferencia suma--producto no se agota en el defecto visible \(54\). La identidad conserva también el defecto de cociente \(129\). Suprimirlo cambiaría la aritmética que se transporta, aunque la tabla residual permaneciera idéntica.

\subsubsection{Reducción ternaria y radical multiplicativo}

El anillo \(R=\ZZ/9\ZZ\) tiene grupo de unidades \(\{1,2,4,5,7,8\}\) e ideal máximo
\[
 \mathfrak m=(3)=\{0,3,6\},\qquad\mathfrak m^2=0.
\]
Toda fila aditiva es una permutación de los nueve representantes. Una fila multiplicativa lo es exactamente cuando su índice es una unidad. Las filas tres y seis contienen tres valores repetidos; la fila nueve contiene únicamente el representante de la clase nula. Por ello el plegamiento multiplicativo ya distingue, antes de cualquier geometría continua, las unidades del sector nilpotente.

La aplicación \(x\mapsto3x\) tiene núcleo e imagen iguales a \(\mathfrak m\). Induce un isomorfismo de espacios vectoriales sobre \(\FF_3\),
\[
 R/\mathfrak m\longrightarrow\mathfrak m,\qquad [x]\longmapsto3x,
\]
pero no un isomorfismo de anillos. En particular, la secuencia visible \(3,6,9\) admite dos lecturas diferentes: su reducción directa módulo tres es \(0,0,0\), mientras que la extracción del factor tres seguida de la reducción de su coordenada interna da \(1,2,0\). Esta segunda operación conserva una coordenada del ideal; no es la misma proyección que reducir directamente el número original.

\subsubsection{Realización latina de la evaluación aditiva}

La diferencia entre las dos evaluaciones puede expresarse también mediante
operadores. Sea \((e_r)_{r\in\ZZ/9\ZZ}\) la base ortonormal de \(\CC^9\).
La asignación \(v_{ab}=e_{a+b}\) produce una base ortonormal en cada fila
y columna. Los proyectores \(P_{ab}=v_{ab}v_{ab}^{*}\) satisfacen
\[
 \sum_bP_{ab}=I_9=\sum_aP_{ab},\qquad P_{ab}^2=P_{ab}.
\]
En efecto, cada traslación \(b\mapsto a+b\) permuta los nueve residuos.
Se obtiene una realización por proyectores de rango uno de un cuadrado
latino, en la terminología de los cuadrados latinos cuánticos
\cite{mustovicary,delascuevas2023}. Esta construcción parte de la hoja
aditiva ya definida y explicita el mapa hacia su representación operatoria.

La evaluación multiplicativa conserva otro contenido: para \(a=3\), los
vectores \(e_{3b}\) repiten los residuos \(0,3,6\), y
\[
 \sum_{b\in\ZZ/9\ZZ}|e_{3b}\rangle\langle e_{3b}|
 =3\bigl(|e_0\rangle\langle e_0|+|e_3\rangle\langle e_3|
                    +|e_6\rangle\langle e_6|\bigr).
\]
La multiplicidad queda registrada en el operador y caracteriza el radical
multiplicativo. Por ello el soporte toroidal, la propiedad latina y la
resolución operatoria de la identidad son propiedades que requieren
definiciones diferentes. Los cuadrados mágicos cuánticos estudiados por
De las Cuevas, Drescher y Netzer exigen entradas positivas y sumas de filas
y columnas iguales a la identidad; su teoría distingue además las
dilataciones con entradas conmutativas\cite{delascuevas2020}. Esta referencia
proporciona un lenguaje preciso para comparar realizaciones de APP.

\subsection{TRIT: orientación y regímenes cuadráticos}

\begin{definition}[Firma ternaria y levantamiento local]
La firma TRIT utiliza la identificación balanceada
\[
 \FF_3\longrightarrow\TRIT,\qquad 0\mapsto0,\quad1\mapsto+1,\quad2\mapsto-1.
\]
Para \(n\in\ZZ\), su levantamiento conserva \(n=r+3c\), donde \(r\in\TRIT\) y \(c\in\ZZ\). En una transición de amplitud local acotada, los tres estados \((0,-1),(0,0),(0,+1)\) pueden presentar el mismo residuo nulo y distintos acarreos.
\end{definition}

La firma no reemplaza a la trayectoria. Un cero residual no significa ausencia de información de orientación, fase o cociente. Tampoco introduce por sí solo una magnitud física. Su función es distinguir algebraicamente las clases de transporte que se realizarán a continuación.

\subsubsection{Tres regímenes cuadráticos y una exponencial común}

Una vez obtenida la firma ternaria, su realización cuadrática se define por
\begin{equation}
 \mathbb A_\tau=\RR[X]/(X^2+\tau),\qquad \tau\in\TRIT.
 \label{eq:algebras-trit}
\end{equation}
Si \(J_\tau\) es la clase de \(X\), entonces \(J_\tau^2=-\tau I\). La estructura distingue el tipo elíptico \(\tau=+1\), el parabólico \(\tau=0\) y el hiperbólico \(\tau=-1\). La realización real no genera la firma: expresa los tipos que la reducción orientada ya ha producido.

\begin{proposition}[Exponenciación de los tres tipos]
En cada álgebra \eqref{eq:algebras-trit},
\begin{equation}
 \exp(tJ_\tau)=C_\tau(t)I+S_\tau(t)J_\tau,
\end{equation}
donde
\[
 C_\tau(t)=\sum_{n\ge0}\frac{(-\tau)^nt^{2n}}{(2n)!},\qquad
 S_\tau(t)=\sum_{n\ge0}\frac{(-\tau)^nt^{2n+1}}{(2n+1)!}.
\]
Las tres realizaciones son
\[
 \begin{array}{c|c|c}
 \tau&J_\tau^2&\exp(tJ_\tau)\\\hline
 +1&-I&\cos t\,I+\sin t\,J_\tau\\
 0&0&I+tJ_\tau\\
 -1&I&\cosh t\,I+\sinh t\,J_\tau.
 \end{array}
\]
\end{proposition}
\begin{proof}
Separar los términos pares e impares de la serie exponencial y sustituir \(J_\tau^{2n}=(-\tau)^nI\) y \(J_\tau^{2n+1}=(-\tau)^nJ_\tau\) da las fórmulas. Para \(\tau=0\), todos los términos de grado al menos dos se anulan. En los otros dos casos se recuperan las series trigonométricas o hiperbólicas.
\end{proof}

La exponencial es común a los tres regímenes: no se asigna exclusivamente a la rama parabólica. Los invariantes
\[
 C_\tau(t)^2+\tau S_\tau(t)^2=1
\]
se deducen de \(\exp(tJ_\tau)\exp(-tJ_\tau)=I\). El mismo principio algebraico conserva la unidad a través de realizaciones que difieren por su tipo cuadrático.

En el régimen hiperbólico, \(P_\pm=(I\pm J_{-1})/2\) son proyectores complementarios y
\[
 \exp(tJ_{-1})=e^tP_++e^{-t}P_-.
\]
El crecimiento y la contracción aparecen conjuntamente en dos direcciones propias. En el elíptico, el parámetro describe una rotación; en el parabólico, una traslación nilpotente. La interpretación temporal adoptada por HMT asocia \(+1\) al retorno con memoria, \(0\) al umbral y \(-1\) a la apertura bidireccional. Esta interpretación exige un orden de parametrización; no añade un cuarto tipo algebraico.

\input{sections/trit_desarrollo.tex}

\input{figures/regimenes_trit.tex}

\newpage
\subsection{TPK: núcleo topológico de fase}

La dinámica utiliza el selector de fase
\[
 M_{\rm ph}(r)=
 \begin{cases}
 +1,&r\in\{1,4,7\},\\
 -1,&r\in\{2,5,8\},\\
 0,&r\in\{3,6,9\}.
 \end{cases}
\]
Su valor decide qué evaluación opera. El transporte cardinal desplaza el cursor correspondiente; el registro conserva el valor anterior, el nuevo cociente, la orientación y la transición. Por tanto, selección y transporte son funciones distintas: un signo ternario no es un desplazamiento del toro.

Denotaremos por \(\widehat X\) el espacio de estados aritméticos con memoria de trayectoria. Un elemento contiene las posiciones y direcciones de los dos cursores, el modo activo, la firma TRIT, la fase, los residuos y cocientes de ambas evaluaciones, los acarreos, el prefijo construido y los datos de frontera. Cuando se introduce un cilindro de refinamiento, éste forma parte del estado. La notación evita sustituir el objeto por una versión reducida que conserve únicamente el símbolo visible.

\begin{definition}[Actualización del TPK]
Sean \(\mathsf{Sel}_t\) la selección de modo mediante \(M_{\rm ph}\), \(\mathsf{Tra}_t\) el transporte cardinal levantado y \(\mathsf{Mem}_t\) la actualización de los registros. La transición es
\begin{equation}
 \mathcal U_t=\mathsf{Mem}_t\circ\mathsf{Tra}_t\circ\mathsf{Sel}_t.
 \label{eq:actualizacion}
\end{equation}
Cada aplicación conserva los campos que no modifica y actualiza los restantes mediante la división euclídea y la concatenación de la trayectoria.
\end{definition}

La notación \eqref{eq:actualizacion} expresa el orden de composición, no reemplaza las reglas efectivas. La siguiente sección especifica completamente su realización finita de dos cursores. Después, la prolongación no se obtendrá borrando su memoria y repitiendo la primera emisión, sino transportando los prefijos y sus fronteras compatibles. De este modo, APP proporciona las evaluaciones, TRIT distingue sus regímenes y TPK compone su evolución.

\input{sections/tpk_desarrollo_integrado.tex}

\input{sections/memoria_resolvente.tex}


% SOURCE foundation/figures/app.tex
\begin{figure}[tbp]
\centering
\begin{tikzpicture}[x=.53cm,y=.53cm,font=\small]
\fill[black!5] (1.5,-.5) rectangle (2.5,8.5);
\fill[black!5] (4.5,-.5) rectangle (5.5,8.5);
\fill[black!5] (-.5,2.5) rectangle (8.5,3.5);
\fill[black!5] (-.5,5.5) rectangle (8.5,6.5);
\fill[black!12] (7.5,-.5) rectangle (8.5,8.5);
\fill[black!12] (-.5,-.5) rectangle (8.5,.5);
\foreach \i in {1,...,9}{
 \node at (\i-1,9.05) {\(\i\)};
 \node at (-1.05,9-\i) {\(\i\)};
 \foreach \j in {1,...,9}{
  \pgfmathtruncatemacro{\v}{mod(\i*\j-1,9)+1}
  \node at (\j-1,9-\i) {\(\v\)};
 }
}
\foreach \k in {0,...,9}{
 \draw[black!25] (\k-.5,-.5)--(\k-.5,8.5);
 \draw[black!25] (-.5,\k-.5)--(8.5,\k-.5);
}
\draw[line width=.65pt] (2.5,3.5) rectangle (4.5,5.5);
\node[anchor=west,align=left,text width=6.5cm] at (10,7.2) {Un soporte: \(81\) posiciones.\\[5pt]Dos evaluaciones:\\ \(S(i,j)=i+j\), \(P(i,j)=ij\).\\[5pt]Cada fila se genera por acumulación:\\ \(P(i,j+1)-P(i,j)=i\).};
\node[anchor=west,align=left,text width=6.5cm] at (10,1.8) {La carta representa \(0\pmod9\) por \(9\).\\[5pt]El bloque señalado es\\ \(B_c=\left(\begin{smallmatrix}7&2\\2&7\end{smallmatrix}\right)\),\\ con autovalores \(9\) y \(5\).};
\end{tikzpicture}
\caption{Tabla multiplicativa de APP en representantes positivos. Las filas y columnas no unitarias se distinguen por trama gris; el marco identifica el bloque diagonal adyacente de diagonales iguales. La tabla dibujada es un observable del toro aritmético; sus bordes de representación no constituyen una frontera topológica intrínseca del toro.}
\label{fig:app}
\end{figure}


% SOURCE foundation/sections/trit_desarrollo.tex
% Ampliación acumulativa del TRIT; se inserta antes de la figura de regímenes.
% No sustituye ni elimina el desarrollo precedente.
\subsubsection{División balanceada y composición de los levantamientos}
\label{trit-ampl-div-balanceada}

La reducción ternaria permite expresar una evaluación entera de APP en una
coordenada visible y una coordenada de prolongación. Su función no consiste
únicamente en abreviar la escritura de un entero: proporciona una regla de
composición que determina cómo se modifica el cociente cuando se suman o se
multiplican evaluaciones. La elección balanceada hace explícita, además, la
compatibilidad de esa regla con la inversión de signo.

\begin{definition}[Coordenadas ternarias balanceadas]
\label{trit-ampl-def-coordenadas}
Para \(n\in\mathbb Z\), sean
\[
 q_3(n)=\left\lfloor\frac{n+1}{3}\right\rfloor,\qquad
 r_3(n)=n-3q_3(n).
\]
La aplicación
\[
 \mathcal B:\mathbb Z\longrightarrow\TRIT\times\mathbb Z,\qquad
 n\longmapsto(r_3(n),q_3(n))
\]
se denomina descomposición ternaria balanceada. Su inversa es
\(\mathcal L(r,c)=r+3c\).
\end{definition}

En efecto, \(r_3(n)\in\{-1,0,+1\}\). Si dos pares representan el mismo entero,
\(r+3c=s+3d\), entonces \(r-s\) es múltiplo de \(3\) y pertenece al intervalo
\([-2,2]\); necesariamente \(r=s\) y \(c=d\). La representación es, por tanto,
única. La división puede repetirse sobre \(q_3(n)\), de manera que cada nivel
conserva el dato necesario para reconstruir el anterior. El cociente no es una
perturbación del residuo: constituye su coordenada complementaria.

La firma obtenida de una evaluación \(F\), con \(F=S\) o \(F=P\) en APP,
es \(r_3\circ F\). Cuando la reducción actúa sobre el cursor de fase, la misma
sección balanceada produce las clases \(\{1,4,7\}\), \(\{2,5,8\}\) y
\(\{3,6,9\}\), con valores \(+1,-1,0\). El selector del TPK que se define
a continuación utiliza esas clases para determinar el modo operativo.
La lectura de una evaluación y la selección de un modo tienen así la misma
reducción ternaria, pero dominios y funciones diferentes.

\Needspace{16\baselineskip}
\begin{proposition}[Aritmética de los pares balanceados]
\label{trit-ampl-prop-aritmetica}
Sean \(n=r+3c\) y \(m=s+3d\), con \(r,s\in\TRIT\). Defínase
\[
 \chi(r,s)=\frac{r+s-r_3(r+s)}3.
\]
Las operaciones enteras se transportan a los pares mediante
\begin{align}
 (r,c)\boxplus(s,d)
 &=\bigl(r_3(r+s),c+d+\chi(r,s)\bigr),
 \label{trit-ampl-eq-suma}\\
 (r,c)\boxtimes(s,d)
 &=\bigl(rs,rd+sc+3cd\bigr).
 \label{trit-ampl-eq-producto}
\end{align}
La aplicación \(\mathcal L\) conserva ambas operaciones.
\end{proposition}
\begin{proof}
La igualdad \(r+s=r_3(r+s)+3\chi(r,s)\) da
\[
 n+m=r_3(r+s)+3\bigl(c+d+\chi(r,s)\bigr).
\]
Por otra parte,
\[
 nm=rs+3(rd+sc+3cd).
\]
Como el producto de dos representantes balanceados pertenece de nuevo a
\(\{-1,0,+1\}\), no necesita una reducción residual adicional. La unicidad
de la descomposición demuestra las dos fórmulas.
\end{proof}

Por ejemplo, \(2=(-1)+3\) y \(4=1+3\). Su suma y su producto se reconstruyen
como
\[
 (-1,1)\boxplus(1,1)=(0,2),\qquad
 (-1,1)\boxtimes(1,1)=(-1,3).
\]
Las evaluaciones resultantes son \(6\) y \(8\). La suma de cocientes basta
en este ejemplo aditivo; en el producto intervienen los términos cruzados
\(rd\), \(sc\) y \(3cd\). La discrepancia entre ambas operaciones reside
también en la coordenada de prolongación, no solamente en el residuo.

El término \(\chi\) es un cociclo de la sección balanceada. Si
\(r\oplus s=r_3(r+s)\), satisface
\[
 \chi(r,s)+\chi(r\oplus s,t)
 =\chi(s,t)+\chi(r,s\oplus t).
\]
Ambos miembros representan el cociente de la misma suma de tres enteros.
Esta identidad garantiza que reagrupar las operaciones no cambia la
evaluación reconstruida. Al mismo tiempo, la proyección sobre el primer
componente suprime esa contabilidad: \(1+1=-1+3\), mientras que la lectura
residual conserva únicamente \(-1\). La información eliminada está
determinada, en este caso, por \(\chi(1,1)=1\).

La relación con el módulo \(9\) conserva una segunda escala discreta.
La aplicación
\[
 \beta:\TRIT\times\mathbb F_3\longrightarrow\mathbb Z/9\mathbb Z,
 \qquad (r,\bar c)\longmapsto r+3c\pmod9
\]
está bien definida y es biyectiva. Cambiar el representante de \(\bar c\)
añade un múltiplo de \(9\); recíprocamente, una igualdad de imágenes fija
primero \(r\) por reducción módulo \(3\) y después \(\bar c\). La suma
transportada por esta biyección contiene el término
\(\chi(r,s)\) de \eqref{trit-ampl-eq-suma}, reducido ahora módulo \(3\).
Así, los dos niveles ternarios no se componen mediante una suma
independiente de coordenadas: la normalización del primer nivel modifica
el segundo.

En particular, las clases \(3\), \(6\) y \(0\) de
\(\mathbb Z/9\mathbb Z\) tienen primer componente \(0\), pero segundo
componente \(1\), \(2\) y \(0\), respectivamente. La distinción ya observada
entre reducción directa y extracción de la coordenada del ideal aparece
como una propiedad de \(\beta\). Conservar esa coordenada impide que el
sector residual nulo se convierta artificialmente en una sola posición.
El módulo \(1000\), utilizado después para organizar bloques decimales,
cumple otra función de representación posicional; su cociente se conserva
mediante la división correspondiente. Un cambio de base conserva el entero
cuando retiene residuo y cociente, mientras que la coincidencia de residuos
en bases distintas no identifica por sí sola sus estados de transporte.

\Needspace{11\baselineskip}
\subsubsection{Inversión, clase residual y coordenadas no visibles}
\label{trit-ampl-inversion}

La inversión aditiva se levanta sin ambigüedad:
\[
 \mathcal B(-n)=(-r_3(n),-q_3(n)).
\]
En particular, \(3\), \(0\) y \(-3\) tienen coordenadas
\((0,1)\), \((0,0)\) y \((0,-1)\). Comparten firma nula, pero representan
evaluaciones distintas. En una transición restringida a un cociente de
amplitud máxima \(1\), son las tres posibilidades de la fibra visible sobre
cero; en un levantamiento entero general, esa fibra contiene todos los pares
\((0,c)\), con \(c\in\mathbb Z\).

La coordenada de cociente recupera el entero evaluado, pero no recupera por
sí sola la trayectoria que lo produjo. Dos caminos de APP pueden compartir
evaluación, residuo y cociente, y diferir en orientación, sucesión de modos o
enrollamiento. Por ello la descomposición \(\mathcal B\) se integra en el
estado de transporte, en vez de sustituirlo. La memoria aritmética de una
división y la memoria de una trayectoria son coordenadas relacionadas, no
sinónimos.

Para una sucesión de firmas de longitud \(n\), la inversión orientada es
\[
 C_{\rm rev}(\tau_1,\ldots,\tau_n)
   =(-\tau_n,\ldots,-\tau_1).
\]
Aplicarla dos veces restituye la sucesión. La operación invierte el orden de
las posiciones y el signo de cada firma; conserva su longitud y las
posiciones nulas con el orden invertido. En la parametrización adoptada,
intercambia el retorno con memoria, asociado a \(+1\), y la apertura
bidireccional, asociada a \(-1\), manteniendo el umbral \(0\). El transporte
de los restantes datos de una trayectoria se especifica mediante la
involución correspondiente del TPK.

Esta inversión de firmas debe distinguirse de la conjugación dentro de un
álgebra cuadrática. La primera puede cambiar el tipo \(\tau\); la segunda,
que se construye a continuación, actúa dentro del tipo ya fijado. La
distinción impide atribuir a una inversión de signo una equivalencia
algebraica entre regímenes diferentes.

\subsubsection{Producto cuadrático, conjugación y norma algebraica}
\label{trit-ampl-norma}

Un elemento de \(\mathbb A_\tau\) tiene expresión única \(aI+bJ_\tau\).
La relación \(J_\tau^2=-\tau I\) determina su producto:
\begin{equation}
 (aI+bJ_\tau)(cI+dJ_\tau)
   =(ac-\tau bd)I+(ad+bc)J_\tau.
 \label{trit-ampl-producto-cuadratico}
\end{equation}
Se trata de una realización real del tipo ya seleccionado. Las coordenadas
\(a,b\) no son residuos de APP, y el producto
\eqref{trit-ampl-producto-cuadratico} no sustituye a
\eqref{trit-ampl-eq-producto}; expresa otra operación, en el codominio
cuadrático declarado.

\begin{definition}[Conjugación y norma algebraica]
\label{trit-ampl-def-norma}
La conjugación del tipo \(\tau\) y su norma algebraica son
\[
 \overline{aI+bJ_\tau}=aI-bJ_\tau,\qquad
 N_\tau(aI+bJ_\tau)=a^2+\tau b^2.
\]
\end{definition}

\begin{proposition}[Multiplicatividad e invertibilidad]
\label{trit-ampl-prop-norma}
Para \(z,w\in\mathbb A_\tau\),
\[
 z\bar z=N_\tau(z)I,\qquad
 N_\tau(zw)=N_\tau(z)N_\tau(w).
\]
Además, \(z\) es invertible si y sólo si \(N_\tau(z)\ne0\), y en ese caso
\[
 z^{-1}=\frac{\bar z}{N_\tau(z)}.
\]
\end{proposition}
\begin{proof}
El producto de \(aI+bJ_\tau\) y \(aI-bJ_\tau\) es
\((a^2+\tau b^2)I\). La conjugación es multiplicativa y el álgebra es
conmutativa, por lo que
\((zw)\overline{zw}=(z\bar z)(w\bar w)\). Si la norma no se anula, la
fórmula indicada proporciona la inversa. Recíprocamente, si \(zw=I\), la
multiplicatividad implica \(N_\tau(z)N_\tau(w)=1\).
\end{proof}

La forma cuadrática \(N_\tau\) es positiva definida únicamente en el tipo elíptico.
Para \(\tau=0\), vale \(a^2\), y el elemento \(J_0\) es no nulo aunque su
cuadrado y su norma sean cero. Para \(\tau=-1\), vale \(a^2-b^2\); los
elementos no nulos \(I+J_{-1}\) e \(I-J_{-1}\) tienen norma cero y su
producto se anula. La norma algebraica es aquí una forma cuadrática
multiplicativa: no presupone una norma positiva de espacio vectorial.

Los tres tipos pueden realizarse fielmente mediante
\[
 aI+bJ_\tau\longmapsto
 \begin{pmatrix}a&-\tau b\\b&a\end{pmatrix}.
\]
El determinante de esta matriz es \(N_\tau(aI+bJ_\tau)\). Quedan
representadas así, con una misma expresión matricial, la estructura
compleja, la estructura de números duales y la estructura real escindida.
Cada una conserva su ley y sus divisores de cero. La forma de firma
\((1,1)\) del último caso tiene interpretación lorentziana en dimensión
\(1+1\); la identificación de una métrica física requiere además un espacio,
un campo y una regla de transporte específicos.

\subsubsection{Conservación multiplicativa y composición de evoluciones}
\label{trit-ampl-evoluciones}

La exponencial ya obtenida pertenece al conjunto de elementos de norma uno:
\[
 N_\tau\bigl(\exp(tJ_\tau)\bigr)=1.
\]
Esta conservación se mantiene al componer parámetros. En efecto, todos los
factores son funciones del mismo generador, y
\[
 \exp(tJ_\tau)\exp(sJ_\tau)=\exp((t+s)J_\tau).
\]
Al comparar las dos componentes se deducen las leyes de adición
\begin{align}
 C_\tau(t+s)&=C_\tau(t)C_\tau(s)-\tau S_\tau(t)S_\tau(s),\\
 S_\tau(t+s)&=S_\tau(t)C_\tau(s)+C_\tau(t)S_\tau(s).
\end{align}
La conjugación cambia \(t\) por \(-t\) y coincide con la inversión dentro
de este subgrupo. No cambia \(\tau\): invierte una evolución del régimen
elegido.

En el tipo elíptico, la forma conservada es \(a^2+b^2\), y la trayectoria
exponencial recorre el círculo de norma uno. En el hiperbólico, las
coordenadas propias son \(u=a+b\), \(v=a-b\), con \(uv=N_{-1}(z)\).
La exponencial da
\[
 (u,v)=(e^t,e^{-t}),\qquad uv=1.
\]
El crecimiento de una coordenada se acompaña de la contracción de la otra.
La conservación expresa una relación entre ambas coordenadas; no afirma que
cada una permanezca constante.

En el tipo parabólico,
\[
 (I+tJ_0)(I+sJ_0)=I+(t+s)J_0,\qquad
 (I+tJ_0)^{-1}=I-tJ_0.
\]
La evolución no se detiene porque \(J_0^2=0\). La nilpotencia elimina los
términos de orden al menos dos, pero conserva el término lineal que
transporta el parámetro. Son distintos el símbolo visible \(0\), el
elemento nulo del álgebra y el generador no nulo \(J_0\). Esta distinción
es necesaria para describir el umbral sin identificarlo con ausencia de
estado o de dinámica.

\begin{proposition}[Composición de coordenadas afines]
\label{trit-ampl-prop-pendientes}
En la carta donde la componente escalar no se anula, represéntese una
dirección mediante \(I+uJ_\tau\). Si \(1-\tau uv\ne0\), la dirección de su
producto con \(I+vJ_\tau\) tiene coordenada
\[
 u\star_\tau v=\frac{u+v}{1-\tau uv}.
\]
\end{proposition}
\begin{proof}
Se tiene
\[
 (I+uJ_\tau)(I+vJ_\tau)
   =(1-\tau uv)I+(u+v)J_\tau.
\]
Dividir por la componente escalar da la expresión anunciada.
\end{proof}

La carta afín conserva la razón entre componentes, no el factor escalar
dividido. Cuando \(1-\tau uv=0\), debe examinarse el producto antes de
cambiar de carta: puede tener dirección definida con componente escalar
nula, o ser el elemento cero en presencia de divisores de cero. La ley de
composición sigue perteneciendo al álgebra completa; la singularidad de
una coordenada no equivale automáticamente a una singularidad del objeto.
Esta fórmula reaparecerá en la lectura regional elíptica. Su uso requiere
mantener el denominador y, cuando corresponda, el factor de normalización.

\subsubsection{Transporte del tipo y realizaciones posteriores}
\label{trit-ampl-puentes}

Las álgebras \(\mathbb A_{+1}\), \(\mathbb A_0\) y
\(\mathbb A_{-1}\) no son isomorfas como álgebras reales unitarias.
La primera es un cuerpo; la segunda contiene un nilpotente no nulo; la
tercera posee idempotentes no triviales y es isomorfa a
\(\mathbb R\times\mathbb R\). Por ello, el cambio de firma que organiza una
ruta no puede interpretarse como un isomorfismo indiferenciado entre las
tres realizaciones. El TPK conserva el índice de régimen, los dominios y el
operador efectivo que actúa en cada transición.

Dentro del capítulo electrónico, el par de proyecciones de las hojas de APP
produce un conmutador normalizado cuyo cuadrado es \(-I\). El mismo bloque
contiene una involución de cuadrado \(I\) y dos direcciones nilpotentes. La
realización conjunta se demostrará sobre ese par concreto de proyecciones;
la clasificación precedente explica por qué las tres relaciones cuadráticas
pueden aparecer en una misma álgebra matricial sin identificarse entre sí.

En la prolongación excepcional, la matriz de Paley reducida módulo \(3\)
satisface \(A_W^2=-I_6\). La relación dota al espacio ternario de una
estructura de dimensión \(3\) sobre \(\mathbb F_9\). Su vínculo con una raíz
operatorial real se expresa mediante la presentación integral
\[
 \mathbb Z[u]/(u^2+1).
\]
Los cambios de base a \(\mathbb F_3\) y a \(\mathbb R\) producen
representaciones de esa misma relación. Cada una conserva su característica,
su dimensión y su información de incidencia. El desarrollo de los
respectivos operadores establece el enlace, no la igualdad de sus
cardinales.

La conservación de norma uno describe, por consiguiente, las realizaciones
cuadráticas de las evoluciones. La conservación de cocientes, orientación y
trayectoria pertenece al estado aritmético transportado. Ambas propiedades
intervienen en la construcción, con funciones diferentes: una determina las
identidades de las representaciones; la otra permite componer, reconstruir y
prolongar las operaciones que las originan.


% SOURCE foundation/figures/regimenes_trit.tex
\begin{figure}[htbp]
\centering
\begin{tikzpicture}[scale=.88,>=Stealth,font=\small]
\foreach \c/\titulo in {0/{Elíptico},4.6/{Parabólico},9.2/{Hiperbólico}}{
\begin{scope}[xshift=\c cm]
\draw[->,gray] (-1.65,0)--(1.7,0) node[right] {\(u\)};
\draw[->,gray] (0,-1.35)--(0,1.5) node[above] {\(v\)};
\node at (0,2.05) {\titulo};
\end{scope}}
\draw[thick] (0,0) circle[radius=1];
\draw[->,thick] (1,0) arc[start angle=0,end angle=60,radius=1];
\node at (0,-1.85) {\(u^2+v^2=1\)};
\node at (0,-2.35) {\(J^2=-I\)};
\draw[thick] (3.6,-1.2)--(3.6,1.2);
\draw[->,thick] (5.6,-1.2)--(5.6,1.2);
\node at (4.6,-1.85) {\(u^2=1\)};
\node at (4.6,-2.35) {\(J^2=0\)};
\draw[thick,domain=-.9:.9,samples=50] plot ({9.2+cosh(\x)},{sinh(\x)});
\draw[thick,domain=-.9:.9,samples=50] plot ({9.2-cosh(\x)},{sinh(\x)});
\node at (9.2,-1.85) {\(u^2-v^2=1\)};
\node at (9.2,-2.35) {\(J^2=I\)};
\end{tikzpicture}
\caption{Realizaciones cuadráticas de TRIT. Para
\(J_\tau=\left(\begin{smallmatrix}0&-\tau\\1&0\end{smallmatrix}\right)\),
la exponencial conserva \(u^2+\tau v^2\). La degeneración parabólica
conserva \(u\) y transporta \(v\) linealmente. Los tres diagramas representan
órbitas de formas cuadráticas; su tipo geométrico corresponde a la relación
algebraica del generador.}
\label{fig:trit-regimenes}
\end{figure}


% SOURCE foundation/sections/tpk_desarrollo_integrado.tex
% Recuperación expositiva desde los propietarios del núcleo TPK.
% La procedencia y el control de conservación se documentan en technical/TPK_INTEGRACION.md.
% Este archivo amplía la sección TPK; no sustituye extension.tex ni generacion.tex.

\subsubsection{Estado dinámico, observables y dependencia de las operaciones}
\label{tpk-int:estado}

El núcleo topológico de fase organiza la evolución de las dos hojas de APP
con la orientación determinada por TRIT. Su definición comprende el estado
sobre el que actúa, las reglas de transición y las relaciones que hacen
compatibles las prolongaciones. La ecuación \eqref{eq:actualizacion} adquiere
contenido operativo al especificar qué coordenadas recibe y modifica cada
uno de sus factores. La posición pertenece al soporte toroidal; la fase
pertenece al calendario; la orientación determina el transporte; los
cocientes registran la evaluación aritmética anterior a su reducción.
Estas coordenadas se relacionan mediante la transición y conservan funciones
matemáticas diferentes.

La proyección observable de un estado es
\begin{equation}
 q=(x_+,d_+,x_\times,d_\times,\theta)
 \in\mathcal Q_{\rm obs}
 :=(X_9\times\mathcal D)^2\times\Theta_{108},
 \qquad \mathcal D=\{N,E,S,O\},\quad
 \Theta_{108}=\ZZ/108\ZZ.
 \label{tpk-int:qobs}
\end{equation}
Los subíndices distinguen los cursores aditivo y multiplicativo. En las
fibras de estados completos sobre \(q\) se conservan la trayectoria ordenada, las evaluaciones
enteras de ambas hojas, sus residuos y cocientes, y los registros de
orientación y frontera. En una prolongación de precisión se añaden el
prefijo y el cilindro compatible. Por tanto, dos estados con la misma
proyección \(q\) pueden conservar registros distintos.

La dependencia de las operaciones tiene tres niveles. En el nivel local,
el selector activa una hoja, el transporte modifica su posición y la
actualización calcula las nuevas coordenadas aritméticas. En el nivel de
trayectorias, estas transiciones se componen y conservan los registros de
los pasos anteriores. En el nivel de refinamiento, las relaciones de
extensión determinan qué continuación conserva el prefijo, la firma y la
frontera del antecedente. El calendario finito describe la fase de esas
operaciones; la holonomía describe su composición sobre la historia.

\subsubsection{Selección trítica y transporte cardinal levantado}
\label{tpk-int:seleccion}

Sea \(\bar\theta\in\{0,\ldots,107\}\) el representante del calendario.
La fase operativa y el sector dodecafásico son
\begin{equation}
 r(\theta)=1+(\bar\theta\bmod9),\qquad
 \beta(\theta)=\left[\left\lfloor\frac{\bar\theta}{9}\right\rfloor\right]_{12}.
 \label{tpk-int:fase}
\end{equation}
El selector ya definido puede escribirse \(\varepsilon_9=M_{\rm ph}\).
Su vector de valores, en la carta ordenada de \(D_9\), es
\[
 m_{\rm ph}=(1,-1,0,1,-1,0,1,-1,0).
\]
La función \(M_{\rm ph}:D_9\to\TRIT\) y este vector son dos
presentaciones del mismo selector, con sus dominios respectivos.

Definamos los indicadores de activación
\[
 a_+(\theta)=\mathbf1_{\{1,4,7\}}(r(\theta)),\qquad
 a_\times(\theta)=\mathbf1_{\{2,5,8\}}(r(\theta)).
\]
El transporte de los cursores es entonces
\begin{equation}
 x_+'=x_++a_+(\theta)\nu_{d_+},\qquad
 x_\times'=x_\times+a_\times(\theta)\nu_{d_\times}
 \quad\text{en }X_9.
 \label{tpk-int:cursores}
\end{equation}
En las fases \(3,6,9\) permanecen ambas posiciones; el evento neutral
forma parte de la cronología y de su registro. Las direcciones cardinales
determinan los vectores \(\nu_d\), mientras la firma de fase determina
cuál de esos vectores se aplica.

El levantamiento entero conserva el cruce de frontera. Si
\(s:X_9\to D_9^2\) elige los representantes positivos y
\(x'=x+a\nu_d\), con \(a\in\{0,1\}\), se define
\begin{equation}
 b_d(x,a)=\frac{s(x)+a\nu_d-s(x')}{9}\in\ZZ^2.
 \label{tpk-int:borde}
\end{equation}
La divisibilidad por \(9\) se sigue de la igualdad en \(X_9\).
Si \(z\in\ZZ^2\) registra los cruces anteriores, la actualización es
\(z'=z+b_d(x,a)\). De este modo,
\[
 s(x')+9z'=s(x)+9z+a\nu_d.
\]
El dato entero de transporte permanece reconstruible después de reducir
la posición sobre el toro.

\begin{proposition}[Conservación del desplazamiento levantado]
\label{tpk-int:prop-desplazamiento}
Para una trayectoria cardinal de \(n\) pasos, cuyos indicadores de
activación son \(a_t\), se cumple
\[
 s(x_n)+9z_n=s(x_0)+9z_0+
 \sum_{t=0}^{n-1}a_t\nu_{d_t}.
\]
La misma igualdad se conserva al concatenar dos trayectorias compatibles.
\end{proposition}
\begin{proof}
La identidad de un paso es \eqref{tpk-int:borde}. Al sumarla para los
\(n\) pasos, las coordenadas intermedias se cancelan. Dividir el intervalo
de suma en dos tramos produce exactamente la ley de concatenación. En una
trayectoria cerrada, la diferencia \(z_n-z_0\) es el grado de
enrollamiento definido previamente.
\end{proof}

El mismo principio rige las evaluaciones de APP. Si una hoja toma el
valor entero \(N_t\), su registro conserva simultáneamente
\(N_t=R_t+9q_t\). La diferencia entre dos pasos verifica
\[
 N_{t+1}-N_t=(R_{t+1}-R_t)+9(q_{t+1}-q_t).
\]
La memoria de cruce espacial \(z_t\) y el cociente aritmético \(q_t\)
tienen definiciones diferentes; su conservación conjunta mantiene la
procedencia geométrica y aritmética de cada transición.

\subsubsection{Cronología de dos cursores y retornos compuestos}
\label{tpk-int:cronologia}

La involución cardinal \(\jmath\) intercambia \(N\) con \(S\) y
\(E\) con \(O\). Tras aplicar \eqref{tpk-int:cursores}, ambas
direcciones se invierten cuando finaliza un bloque de \(27\) pasos.
Escribiendo
\[
 h(\theta)=\mathbf1_{\{0,27,54,81\}}
             ((\bar\theta+1)\bmod108),
\]
la regla completa sobre la proyección observable es
\begin{equation}
 F_{\rm obs}(q)=
 \bigl(x_+',\jmath^{h(\theta)}d_+,
       x_\times',\jmath^{h(\theta)}d_\times,
       \theta+1\bigr).
 \label{tpk-int:fobs}
\end{equation}
La lectura que forma la emisión utiliza las posiciones anteriores al
transporte, según la recurrencia explícita de la
sección~\ref{sec:extension}. La actualización observable y la emisión
se componen, por tanto, con un orden fijado.

\begin{proposition}[Reversibilidad y periodo del calendario observable]
\label{tpk-int:periodo}
La aplicación \(F_{\rm obs}\) es una permutación de
\(\mathcal Q_{\rm obs}\) y satisface \(F_{\rm obs}^{108}=I\).
En un bloque de \(54\) pasos se restauran posiciones y direcciones;
la coordenada de \(\Theta_{108}\) avanza \(54\).
\end{proposition}
\begin{proof}
Dado el estado final, se recupera primero la fase anterior restando uno
en \(\Theta_{108}\). Esa fase determina si se invirtieron las
direcciones y qué cursor se desplazó. Aplicar la involución correspondiente
y restar el vector cardinal recupera el estado inicial de manera única.

Cada intervalo de \(27\) pasos contiene nueve activaciones de cada
cursor. El bloque siguiente conserva las activaciones e invierte las
direcciones, de modo que los desplazamientos enteros se cancelan en
\(54\) pasos. La regla tiene periodo \(54\) en sus incrementos de
posición y orientación; la misma cancelación vale con cualquier origen
del calendario. Tras \(108\) pasos se restaura también \(\theta\).
\end{proof}

Se reserva
\[
 \mathcal K_{27}=F_{\rm obs}^{27}:
 \mathcal Q_{\rm obs}\longrightarrow\mathcal Q_{\rm obs}
\]
para el iterado de \(27\) transiciones. Su cuarto iterado es la
identidad sobre el estado observable definido en \eqref{tpk-int:qobs}.
La proyección de la historia al calendario permite efectuar esta
comprobación finita. El registro de la trayectoria, sus evaluaciones y
los nuevos prefijos se conservan en el estado dinámico sobre el que
actúa \(\mathcal U_t\); la igualdad observable no los reinicializa.

\subsubsection{Transporte afín de la memoria y composición de estados}
\label{tpk-int:memoria}

La proyección observable admite fibras de datos aritméticos y geométricos.
Sobre una configuración \(q\), escribimos un estado en la forma
\[
 x=(q,o,h,c,\sigma,C,b,\lambda).
\]
Aquí \(o\) conserva la orientación y la hoja; \(h\), los residuos
y cocientes; \(c\), la memoria aditiva de transporte; \(\sigma\),
la firma de frontera; \(C\), el cilindro de precisión; \(b\), su
etiqueta de frontera; y \(\lambda\), la trayectoria registrada. La
firma está determinada por una aplicación
\[
 \operatorname{Sig}_q:
 \mathsf H_q\times\mathsf C_q\times\mathsf{Hist}_q
 \longrightarrow\mathsf S_q,
 \qquad \sigma=\operatorname{Sig}_q(h,c,\lambda).
\]
Los símbolos \(\mathsf H_q,\mathsf C_q,\mathsf{Hist}_q\) y
\(\mathsf S_q\) denotan, respectivamente, los conjuntos de datos
residuo--cociente, el grupo abeliano de memoria de transporte, las
trayectorias registradas y las firmas de frontera.
Las realizaciones de esta firma se especifican con sus componentes en
las construcciones regional y terminal. Esta dependencia impide tratar
la firma como un parámetro libre una vez fijados los datos que la producen.

La carta elemental de residuo y cociente es explícita. Para
\(n\in\ZZ\), sean
\[
 r=1+((n-1)\bmod9),\qquad u=\frac{n-r}{9}.
\]
Entonces \(n=r+9u\), con \(r\in D_9\) y \(u\in\ZZ\).
La unicidad se sigue de que dos representantes de \(D_9\) congruentes
módulo \(9\) son iguales. Esta carta proporciona la reconstrucción
entera local. Las torres de precisión incorporan, además, las aplicaciones
de truncamiento y los cilindros desarrollados en
\ref{sec:extension} y \ref{sec:generacion}.

Sea \(\gamma:q\to q'\) una trayectoria observable. El transporte
de las fibras se expresa mediante aplicaciones
\[
 \mathsf H(\gamma):\mathsf H_q\to\mathsf H_{q'},\qquad
 \mathsf C(\gamma):\mathsf C_q\to\mathsf C_{q'},\qquad
 \mathsf{Hist}(\gamma):\mathsf{Hist}_q\to\mathsf{Hist}_{q'}.
\]
Las aplicaciones respetan identidades y composición, y
\(\mathsf C(\gamma)\) es un homomorfismo de grupos abelianos.
El incremento \(k_{\rm tr}(\gamma)\in\mathsf C_{q'}\) cumple
\begin{equation}
 \begin{split}
 k_{\rm tr}(\operatorname{id}_q)&=0,\\
 k_{\rm tr}(\gamma_2\circ\gamma_1)
 &=k_{\rm tr}(\gamma_2)+
   \mathsf C(\gamma_2)k_{\rm tr}(\gamma_1).
 \end{split}
 \label{tpk-int:cociclo}
\end{equation}
El símbolo \(k_{\rm tr}\) designa aquí el cociclo de transporte;
la coordenada racional \(\kappa\) del registro \(K\) conserva
su notación en el capítulo de la clausura dodecafásica.

Las coordenadas transportables se actualizan por
\begin{equation}
 \begin{aligned}
 h'&=\mathsf H(\gamma)h,&
 c'&=\mathsf C(\gamma)c+k_{\rm tr}(\gamma),\\
 \lambda'&=\gamma\circ\lambda,&
 \sigma'&=\operatorname{Sig}_{q'}(h',c',\lambda').
 \end{aligned}
 \label{tpk-int:accion-memoria}
\end{equation}
La segunda ecuación es una acción afín: el término lineal transporta
la memoria anterior y el cociclo añade el incremento producido por la
trayectoria. En la realización de cruce espacial de cada cursor por separado,
\(\mathsf C(\gamma)=I\) y \(k_{\rm tr}(\gamma)\) es la suma
de los vectores de \eqref{tpk-int:borde}. Esta instancia concreta
vincula la ley abstracta con las operaciones cardinales ya definidas.

\begin{proposition}[Compatibilidad de la actualización con la concatenación]
\label{tpk-int:composicion-memoria}
La actualización \eqref{tpk-int:accion-memoria} por
\(\gamma_1\), seguida de la actualización por \(\gamma_2\),
coincide con la actualización por \(\gamma_2\circ\gamma_1\).
\end{proposition}
\begin{proof}
Para la coordenada afín, dos actualizaciones dan
\[
 c''=\mathsf C(\gamma_2)\mathsf C(\gamma_1)c
       +\mathsf C(\gamma_2)k_{\rm tr}(\gamma_1)
       +k_{\rm tr}(\gamma_2).
\]
La functorialidad de \(\mathsf C\) y
\eqref{tpk-int:cociclo} convierten esta expresión en
\(\mathsf C(\gamma_2\circ\gamma_1)c+
k_{\rm tr}(\gamma_2\circ\gamma_1)\).
La afirmación para \(h\) se sigue de la composición de
\(\mathsf H\), y para \(\lambda\), de la asociatividad de las
trayectorias. La firma final coincide porque se evalúa sobre las mismas
tres coordenadas reconstruidas.
\end{proof}

El cilindro y su frontera completan esta ley mediante las relaciones
\(\operatorname{Ext}_\gamma\), definidas en
\ref{tpk-int:bidireccionalidad}. Sus elementos determinan cuáles de
las actualizaciones anteriores conservan una prolongación admisible.
Los estados, junto con las ternas \((\gamma,x,x')\) admitidas por
esas relaciones, forman una categoría sobre las trayectorias observables:
\[
 (\gamma_2,x',x'')\circ(\gamma_1,x,x')
     =(\gamma_2\circ\gamma_1,x,x'').
\]
La identidad de \(x\) es \((\operatorname{id}_q,x,x)\).
La composición conserva a la vez el transporte aritmético y la
prolongabilidad de frontera. Esta es la estructura en la que se efectúan
los retornos con memoria y los refinamientos sucesivos.

\subsubsection{Registro dodecafásico, fase y memoria acumulada}
\label{tpk-int:dodecafase}

Las \(108\) transiciones se distribuyen en las ventanas ordenadas
\[
 E_m=\{9m-8,\ldots,9m\},\qquad 1\leq m\leq12.
\]
El registro conserva el origen de fase, la orientación y los datos
producidos durante cada ventana. Su aplicación es
\[
 \mathcal R_{12}(x)=
 \bigl((E_m,\tau_m,\sigma_m,\kappa_m,\Lambda_m)\bigr)_{m=1}^{12},
\]
donde \(\tau_m\) es la subruta de la ventana, \(\sigma_m\) su
orientación, \(\kappa_m\) el incremento entero de frontera y
\(\Lambda_m\) el registro local de incidencias. Se adopta la carta
que se desarrollará en \ref{subsec:k-registro-integral}; sus índices
de ventana corresponden a \(m=\beta+1\).
El dominio comprende las historias compatibles
del TPK y el codominio, sus registros temporales orientados. La
composición de segmentos y sus incrementos se rige por
\eqref{tpk-int:accion-memoria}. El grupo \(C_{12}\) describe, en
cambio, la traslación del índice de ventana; \(\Theta_{108}\) conserva
el calendario de pasos. Esta distinción permite precisar qué se transforma
y qué se observa al completar un ciclo.

En las coordenadas \(\theta=9b+r\), con \(0\leq r<9\), el avance
de \(9\) pasos conserva \(r\) y aumenta \(b\) en una unidad.
La adición de fases introduce el cociente
\[
 c_0(r,r')=\left\lfloor\frac{r+r'}9\right\rfloor,
\]
cuya identidad de cociclo asegura la asociatividad de la composición.
En el calendario finito se conserva \([b]_{12}\); en un registro de
vueltas se conserva \(b\in\ZZ\), o \(b\in\mathbb N_0\) para
una historia de origen fijado. La sección~\ref{sec:extension} demuestra
la sucesión exacta no escindida \eqref{eq:extension108}. Su función es
describir la relación entre estas coordenadas, antes de emplearlas en la
construcción del estado firmado.

El registro firmado recibe los datos de fase y memoria mediante la
composición que se desarrolla en la sección~\ref{sec:k-registro}:
\begin{equation}
 R_{\rm sgn}=\Pi_H\circ\mathcal E_{108}^{90,120}
                         \circ\mathcal R_{12}.
 \label{tpk-int:registro-causal}
\end{equation}
Su dominio es el estado terminal con memoria; su salida pertenece a
\(\ZZ^{12}\). La transformación reversible entre ese estado firmado
y \(K\), así como la lectura racional de \(K\), se demuestran en
el capítulo correspondiente. Las operaciones de registro anteceden a
esas transformaciones. La reconstrucción de un registro desde sus
diferencias cíclicas verifica su recuperabilidad; la procedencia dinámica
debe conservarse mediante los eventos que lo produjeron.

Para la traza codificada \((d_t)\) conservada por el registro, el
lector de eventos cuenta los códigos \(\{2,4,6,8\}\) de cada
ventana con la orientación
\((+,+,+,-,-,-,+,+,+,-,-,-)\). Estos códigos y los símbolos
cardinales de \(\mathcal D\) son coordenadas distintas del
registro. El capítulo~\ref{sec:k-registro} desarrolla este lector,
la descomposición integral \(1+5+6\), las congruencias que
caracterizan su imagen y su inversa exacta. La prueba de inversión
establece qué datos conserva esa carta; la definición de
\(\mathcal R_{12}\) especifica la historia orientada sobre la que
actúa.

\subsubsection{Incidencia de fases y construcción interna de los canales}
\label{tpk-int:incidencia}

El selector determina tres subconjuntos de \(D_9\):
\[
 H_+=\{1,4,7\},\qquad H_-=\{2,5,8\},\qquad H_0=\{3,6,9\}.
\]
Sea \(H_{\rm act}=H_+\sqcup H_-\) el conjunto de fases activas y
sea \(O_{\rm ph}=D_9\setminus\{9\}\) la carta anterior al retorno
marcado. Denotemos por \(P_2(H_{\rm act})\) el conjunto de pares
no ordenados de fases activas. La incidencia interna está dada por
\begin{equation}
 \mathcal F^{\rm ph}_6=H_{\rm act}\times P_2(H_{\rm act}),\qquad
 \mathcal F^{\rm ph}_8=O_{\rm ph}\times P_2(H_{\rm act}),\qquad
 \Theta_{\rm ph}=D_9\times P_2(H_{\rm act}).
 \label{tpk-int:fases90120}
\end{equation}
Los índices describen la cantidad de posiciones de la primera componente.
La construcción utiliza las fases y el origen del calendario ya fijados;
las realizaciones excepcionales posteriores transportarán estas
incidencias a sus propios dominios.

\begin{proposition}[Cardinales e incidencia orientada de fase]
\label{tpk-int:cardinales}
Los conjuntos de \eqref{tpk-int:fases90120} tienen cardinales
\(90,120,135\), respectivamente. Si
\[
 \partial_{\rm ph}=\Theta_{\rm ph}\times\{-1,+1\},\qquad
 E_{\rm ph}=\partial_{\rm ph}\sqcup\{\infty\},\qquad
 L_{\rm ph}=H_0\times P_2(H_{\rm act}),
\]
entonces
\[
 |\partial_{\rm ph}|=270,\quad |E_{\rm ph}|=271,\quad
 |L_{\rm ph}|=45,\quad
 |\partial_{\rm ph}\sqcup L_{\rm ph}|=315.
\]
Trasladar simultáneamente el origen de fase y el retorno marcado conserva
todos estos cardinales y sus inclusiones.
\end{proposition}
\begin{proof}
Hay seis fases activas, por lo que
\(|P_2(H_{\rm act})|=\binom62=15\). Los productos por
\(6,8,9\) dan \(90,120,135\). La orientación duplica \(135\);
la marca de retorno añade un elemento; las tres fases neutrales dan
\(3\cdot15=45\). Una traslación del calendario transporta cada
factor por una biyección y lleva el punto marcado a su imagen. Los
productos y las uniones disjuntas conservan así la incidencia.
\end{proof}

El cardinal \(271\) tiene aquí una realización por incidencia de
fase. La completación cúbica de la sección~\ref{sec:extension} tiene
su propia construcción de la corona \(1000-729\). La correspondencia
entre ambas realizaciones debe conservar los mapas y las marcas que las
identifican, además del cardinal. Análogamente, el selector
\(\varepsilon_9\), la incidencia \((90,120)\) y el extractor
\(\mathcal E_{108}^{90,120}\) conservan los dominios y operaciones
que se acaban de distinguir.

\subsubsection{Inversión de trayectorias y prolongaciones conjugadas}
\label{tpk-int:bidireccionalidad}

La bidireccionalidad comienza en una operación sobre trayectorias, antes
de asignarles valores escalares. Para
\(\gamma=(x_0;d_1,\ldots,d_n)\), con extremo \(x_n\), definimos
\[
 \operatorname{rev}\gamma
   =(x_n;\jmath d_n,\ldots,\jmath d_1).
\]
El origen de la trayectoria inversa es el extremo de la original.
La composición satisface
\[
 \operatorname{rev}^2\gamma=\gamma,\qquad
 \operatorname{rev}(\gamma_2\circ\gamma_1)
   =\operatorname{rev}\gamma_1\circ\operatorname{rev}\gamma_2,
\]
y el desplazamiento entero cambia de signo. Estas identidades se obtienen
invirtiendo el orden de los pasos y sustituyendo cada vector cardinal por
su opuesto. En un lazo se sigue
\(\operatorname{wind}(\operatorname{rev}\gamma)
=-\operatorname{wind}(\gamma)\).

Sobre los registros, la inversión requiere transportar asimismo el orden
de sus componentes y la orientación de sus eventos. La inversión de
ruta y una conjugación de firma son operaciones distinguibles: su
composición se especifica cuando se aplica a un registro concreto. En
particular, la involución regional que intercambia las componentes de
clausura y propagación conserva el eje de autoescala; su fórmula se
presenta en la sección~\ref{mono-ampl:regiones}. Esta operación regional
actúa sobre otro dominio que la inversión de una trayectoria cardinal.

Las extensiones sobre una ruta \(r:q\to q'\) se expresan mediante
relaciones entre fibras de estados compatibles. Si \(\mathcal F_q\)
denota la fibra de estados completos sobre \(q\), se exige
\[
 \operatorname{Ext}_r\subseteq\mathcal F_q\times\mathcal F_{q'},
 \qquad \Delta_{\mathcal F_q}\subseteq
        \operatorname{Ext}_{\operatorname{id}_q},
\]
\[
 \operatorname{Ext}_{r_2}\circ\operatorname{Ext}_{r_1}
       \subseteq\operatorname{Ext}_{r_2\circ r_1}.
\]
Su composición conserva las evaluaciones de ambas hojas, la orientación,
los cocientes, el prefijo y la frontera. La inversión de la ruta no
implica por sí sola que toda relación de extensión sea una biyección.
La prolongación conjugada se define en el dominio donde el transporte de
esos datos satisface las compatibilidades correspondientes.

\subsubsection{Holonomía nonádica y continuidad del refinamiento}
\label{tpk-int:holonomia}

Una vuelta de fase compone las transiciones sobre el estado con memoria:
\begin{equation}
 \Gamma_{9,k}=\mathcal U_{k+8}\circ\cdots\circ\mathcal U_k.
 \label{tpk-int:gamma}
\end{equation}
Sobre una historia individualizada, esta composición funcional realiza
la relación de holonomía; en el dominio completo se conserva
\(\Gamma_{9,k}\subseteq\widehat X_k\times\widehat X_{k+9}\),
con las prolongaciones admisibles desarrolladas más adelante.
La base cíclica tiene nueve fases, mientras
la fibra conserva la trayectoria, los cocientes y la frontera. La
monodromía es la acción de una vuelta de esa base sobre las fibras; la
holonomía \(\Gamma_{9,k}\) expresa su realización por las transiciones
efectivas del TPK.

Si \(g\) observa la fase y \(M\) el número de vueltas conservadas,
\[
 g(\Gamma_{9,k}x)=g(x),\qquad M(\Gamma_{9,k}x)=M(x)+1.
\]
La primera identidad expresa clausura de fase y la segunda identifica
la nueva profundidad temporal. La prolongación coinductiva conserva
además sus cilindros y prefijos. Las restricciones entre profundidades
componen un sistema inverso; una sección compatible reúne sus
antecedentes sin reemplazarlos por la última observación. El desarrollo
de esta aritmética de historias y de su unidad ocupa
\ref{hist:seccion}--\ref{hist:doble-varianza}; la construcción de las
regiones y sus lectores se realiza en la sección~\ref{sec:generacion}.

La transferencia de fase sobre emisiones ya construidas tiene otro
dominio. Sean \(\mathcal A_+\) y \(\mathcal A_\times\) los
alfabetos de firmas aditivas y multiplicativas, de cardinales \(18\)
y \(26\). Su producto identifica el catálogo \(\mathcal U_6\).
Para \(f:\mathcal A_+\times\mathcal A_\times\to\RR\),
las medias de hoja son
\[
 (E_+f)(s,e)=\frac1{18}\sum_{s'\in\mathcal A_+}f(s',e),\qquad
 (E_\times f)(s,e)=\frac1{26}\sum_{e'\in\mathcal A_\times}f(s,e').
\]
La construcción de estas emisiones y sus cardinales se efectúa mediante
la recurrencia de la sección~\ref{sec:extension}. Una vez obtenidas,
el operador de transferencia se define sobre
\(\mathcal U_6\times C_9\) por
\[
 P_+=E_+,\quad P_-=E_\times,\quad P_0=I,\qquad
 \mathcal K_{\rm ph}((u,r),(v,r+1))=P_{M_{\rm ph}(r)}(u,v),
\]
con peso cero para las restantes transiciones de fase. El índice
\(r\) utiliza el representante positivo de la clase cíclica.
El orden de fases \((+1,-1,0)^3\) produce entonces la renovación
\[
 \mathcal K_{\rm ph}^{9}\big|_r=E_+E_\times.
\]
En efecto, ambas medias son idempotentes, conmutan y promedian factores
diferentes; su composición asigna peso \(1/(18\cdot26)=1/468\)
a cada emisión. Esta igualdad caracteriza un observable de transferencia
posterior. El transporte de las historias sigue dado por
\eqref{tpk-int:gamma}, con su memoria y su frontera.

La jerarquía resultante establece la función de los capítulos siguientes.
La emisión finita construye el dominio y las regiones; las relaciones de
extensión conservan sus prefijos; la conexión nonádica los prolonga; el
registro dodecafásico determina los datos que intervienen en \(K\) y
en la clausura de \(\alpha\). La realización numérica y la realización
incidencial reciben esa estructura anterior. Este orden permite exponer
cada consecuencia en su capítulo conservando, desde el núcleo formal,
los objetos que la hacen posible.


% SOURCE foundation/sections/memoria_resolvente.tex
\subsection{Memoria acumulada y resolventes del transporte}
\label{subsec:memoria-resolvente}

El retorno del calendario no elimina los incrementos producidos durante
el recorrido. Para expresar esta diferencia operatoriamente, se construye
primero un registro acumulativo de eventos de la trayectoria y después su
serie generadora. La reducción de variables de memoria se realizará sobre
ese transporte, conservando también la fuente y el lector de salida.

\subsubsection{Eventos orientados y actualización del registro}

Sea \((d_t)\) la secuencia de códigos enteros de dirección conservada
por una trayectoria del TPK. El código \(d_t\) y la dirección cardinal
del cursor son coordenadas distintas; no se supone una identificación
entre sus alfabetos. Numeramos desde cero las ventanas
\[
 I_m=\{9m+1,\ldots,9m+9\},\qquad m\geq0.
\]
En el primer ciclo, \(I_m\) es la ventana \(E_{m+1}\) del registro
dodecafásico. Su orientación es
\(\sigma=(1,1,1,-1,-1,-1,1,1,1,-1,-1,-1)\).
Para una historia prolongada, \(\sigma_m\in\{-1,1\}\) es el signo
conservado en la ventana correspondiente. El evento firmado es
\begin{equation}
 a_m=\sigma_m\sum_{t\in I_m}
       \mathbf1_{\{2,4,6,8\}}(d_t),\qquad |a_m|\leq9.
 \label{mem:evento}
\end{equation}
Cada sumando se incorpora cuando ocurre su transición. Esta definición
cuenta los eventos de la traza; no escoge direcciones mediante un valor
que deba obtenerse después.

Sean \(\mathbf e_0,\ldots,\mathbf e_{11}\) la base de \(\ZZ^{12}\)
y \(S\mathbf e_j=\mathbf e_{j+1\bmod12}\). Escribimos
\(R=S^{-1}\). El registro acumulado en coordenadas fijas satisface
\[
z_{m+1}=z_m+a_m\mathbf e_{m\bmod12}.
\]
Un prefijo sin historia acumulada empieza en \(z_0=0\); en otro
origen, \(z_0\) conserva el registro previo.
El marco que acompaña al calendario es \(y_m=S^{-m}z_m\).

\begin{proposition}[Actualización y retorno con memoria]
\label{mem:prop-actualizacion}
El registro en el marco móvil cumple
\begin{equation}
 y_{m+1}=Ry_m+\eta_m,\qquad
 \eta_m=a_mR\mathbf e_0.
 \label{mem:actualizacion}
\end{equation}
Para todo \(n\geq1\),
\begin{equation}
 y_n=R^ny_0+\sum_{j=0}^{n-1}R^{n-1-j}\eta_j.
 \label{mem:iteracion}
\end{equation}
En particular,
\begin{equation}
 y_{12}=y_0+\sum_{j=0}^{11}a_j\mathbf e_j.
 \label{mem:retorno12}
\end{equation}
\end{proposition}
\begin{proof}
Multiplicar la actualización de \(z_m\) por \(S^{-(m+1)}\) da
\(y_{m+1}=S^{-1}y_m+a_mS^{-1}\mathbf e_0\), porque
\(S^{-m}\mathbf e_{m\bmod12}=\mathbf e_0\).
La fórmula \eqref{mem:iteracion} se obtiene por inducción.
Para \(n=12\), \(R^{12}=I\) y
\(R^{11-j}\eta_j=a_jR^{12-j}\mathbf e_0=a_j\mathbf e_j\),
lo que prueba \eqref{mem:retorno12}.
\end{proof}

Así, doce ventanas restauran el marco y conservan el vector de eventos.
El recuento no permite recuperar por sí solo el orden de todos los pasos
dentro de cada ventana; ese orden permanece en la trayectoria registrada.

Para precisar el balance, definimos
\[
 D_3=I-S^3,\quad D_4=I-S^4,\quad
 B=D_3^{\mathsf T}D_3+D_4^{\mathsf T}D_4,
\]
\[
 \mathcal M(y)=\tfrac12y^{\mathsf T}By,
 \qquad q(y)=\sum_{j=0}^{11}y_j.
\]
La primera es una forma cuadrática del registro; la segunda registra
la suma de sus coordenadas. No se identifican aquí con energía ni carga físicas.

\begin{proposition}[Balance producido por el evento]
\label{mem:prop-balance}
La actualización \eqref{mem:actualizacion} satisface
\begin{equation}
 \begin{aligned}
 \mathcal M(y_{m+1})-\mathcal M(y_m)
   &=a_m(By_m)_0+2a_m^2,\\
 q(y_{m+1})-q(y_m)&=a_m.
 \end{aligned}
 \label{mem:balance}
\end{equation}
\end{proposition}
\begin{proof}
Al desarrollar los productos,
\(B=4I-S^3-S^{-3}-S^4-S^{-4}\); por tanto
\(R^{\mathsf T}BR=B\) y \(B_{00}=4\).
Como \(y_{m+1}=R(y_m+a_m\mathbf e_0)\), la diferencia cuadrática
es \(a_m\mathbf e_0^{\mathsf T}By_m+\tfrac12a_m^2B_{00}\).
La suma de coordenadas es invariante bajo la permutación \(R\),
lo que da la segunda igualdad.
\end{proof}

\Needspace{20\baselineskip}
\subsubsection{Serie generadora y control uniforme de la memoria}

La serie conserva todos los incrementos, no únicamente su suma al final
de un ciclo. Con una variable formal \(u\), sean
\[
 Y(u)=\sum_{m\geq0}y_mu^m,\qquad
 H_\eta(u)=\sum_{m\geq0}\eta_mu^m.
\]

\begin{proposition}[Resolvente y cota de cola]
\label{mem:prop-serie}
Se cumple la identidad formal
\begin{equation}
 Y(u)=(I-uR)^{-1}\bigl(y_0+uH_\eta(u)\bigr).
 \label{mem:serie}
\end{equation}
Ambas series convergen para \(|u|<1\). Si \(\ell:\RR^{12}\to\RR\)
es lineal, \(L=\|\ell\|_{1\to\RR}\), \(M_0=\|y_0\|_1\) y
\(0<r<1\), entonces, para todo entero \(N\geq0\) y \(|u|\leq r\),
\begin{equation}
 \left|\sum_{m>N}\ell(y_m)u^m\right|
 \leq Lr^{N+1}\left(
 \frac{M_0+9(N+1)}{1-r}+\frac{9r}{(1-r)^2}\right).
 \label{mem:cota}
\end{equation}
La serie de derivadas converge uniformemente en cada disco
\(|u|\leq r<1\).
\end{proposition}
\begin{proof}
Multiplicar \eqref{mem:actualizacion} por \(u^{m+1}\) y sumar da
\(Y-y_0=uRY+uH_\eta\). El término constante de \(I-uR\) es
invertible, de modo que su inversa formal es \(\sum_{k\geq0}u^kR^k\).
Además, \(R\) preserva la norma \(\ell^1\) y
\(\|\eta_m\|_1\leq9\), de donde
\(\|y_m\|_1\leq M_0+9m\).
Para la cola se suman las mayorantes
\(L(M_0+9m)r^m\), usando
\[
 \sum_{m>N}r^m=\frac{r^{N+1}}{1-r},\qquad
 \sum_{m>N}mr^m=r^{N+1}
 \left(\frac{N+1}{1-r}+\frac r{(1-r)^2}\right).
\]
Por último, \(\sum_{m\geq1}Lm(M_0+9m)r^{m-1}<\infty\)
mayora la serie derivada y prueba su convergencia uniforme.
\end{proof}

La variable \(u\) cuenta ventanas de nueve transiciones. Su identificación
con la variable de otro lector requiere conservar ese reloj y el mapa
que relaciona ambas lecturas.

\subsubsection{Eliminación de memoria: operador, fuente y lector}

En un dominio donde la actualización está fijada como \(U:X\to X\),
el espacio vectorial libre \(V=\QQ^{(X)}\) sobre los estados admite el operador
\(T\delta_x=\delta_{U(x)}\). El reloj se incluye en \(x\) si la
transición depende de él. Esta extensión lineal conserva estados
distintos aunque compartan una fase observable.

Consideremos una descomposición \(V=V_0\oplus V_1\), donde \(V_1\)
es el sector auxiliar que se desea eliminar, y escribamos
\[
 T=\begin{pmatrix}A&B_1\\C&D\end{pmatrix},\qquad
 v=\binom{v_0}{v_1},\qquad \ell=(\ell_0,\ell_1).
\]
Aquí \(A:V_0\to V_0\), \(B_1:V_1\to V_0\),
\(C:V_0\to V_1\) y \(D:V_1\to V_1\).
Las identidades siguientes son formales; también valen donde los
resolventes convergen.

\begin{proposition}[Reducción de Schur con fuente y lector]
\label{mem:prop-schur}
Sean
\[
 D_u=(I-uD)^{-1},\qquad
 \mathscr F(u)=I-uA-u^2B_1D_uC.
\]
La componente retenida de \((I-uT)^{-1}v\) es
\begin{equation}
 w_0=\mathscr F(u)^{-1}\bigl(v_0+uB_1D_uv_1\bigr),
 \qquad w_1=D_u(v_1+uCw_0).
 \label{mem:schur-fuente}
\end{equation}
Su lectura completa satisface
\begin{equation}
 \begin{split}
 \ell(I-uT)^{-1}v
 ={}&(\ell_0+u\ell_1D_uC)
       \mathscr F(u)^{-1}(v_0+uB_1D_uv_1)\\
    &+\ell_1D_uv_1.
 \end{split}
 \label{mem:schur-lector}
\end{equation}
\end{proposition}
\begin{proof}
Las dos ecuaciones de \((I-uT)w=v\) son
\(w_0=v_0+uAw_0+uB_1w_1\) y
\(w_1=v_1+uCw_0+uDw_1\).
Resolver la segunda y sustituirla en la primera da
\eqref{mem:schur-fuente}; aplicar \(\ell_0\) y \(\ell_1\)
a ambas componentes da \eqref{mem:schur-lector}.
Todas las inversas formales existen porque sus términos constantes
son identidades.
\end{proof}

El término \(u^2B_1D_uC=\sum_{k\geq0}u^{k+2}B_1D^kC\)
describe una salida hacia la memoria, \(k\) pasos dentro de ella y
un regreso. En cambio, \(u\ell_1D_uC\) lee la memoria antes de ese
regreso. Por ello, reducir sólo el operador no conserva necesariamente
la observación: también deben transformarse fuente y lector.

La diferencia es visible en el ciclo ya construido. Si \(P\) retiene
\(\mathbf e_0\), entonces \(PRP=0\), pero
\begin{equation}
 \left\langle\mathbf e_0,(I-uR)^{-1}\mathbf e_0\right\rangle
 =\frac1{1-u^{12}},\qquad
 \left\langle\mathbf e_{11},(I-uR)^{-1}\mathbf e_0\right\rangle
 =\frac u{1-u^{12}}.
 \label{mem:ejemplo-ciclo}
\end{equation}
En efecto, \(R^n\mathbf e_0=\mathbf e_{-n\bmod12}\): las
dos series recogen, respectivamente, \(n\equiv0\) y
\(n\equiv1\pmod{12}\). El resolvente de la compresión es \(I\),
mientras la compresión del resolvente conserva los retornos omitidos.
Este ejemplo corresponde al registro de doce ventanas, no a la totalidad
del estado del TPK.

\subsubsection{Composición de segmentos y coherencia de la reducción}

Sean tres transportes afines composables
\(F_i(y)=R_i y+b_i\), con dominios y codominios compatibles.
La memoria producida por el recorrido completo es
\begin{equation}
 F_3F_2F_1(y)
 =R_3R_2R_1y+b_3+R_3b_2+R_3R_2b_1.
 \label{mem:cociclo-tres}
\end{equation}
La prueba es la sustitución sucesiva de las tres fórmulas afines.
Agrupar primero los dos primeros segmentos o los dos últimos produce
la misma expresión; los factores transportan cada incremento desde
el punto donde se originó. En \eqref{mem:actualizacion},
\(R_i=R\) y \(b_i\) es el evento de su ventana.

La eliminación de varias capas conserva una coherencia semejante.
Para hacerla explícita, descomponemos ahora
\(V=V_0\oplus V_1\oplus V_2\) y escribimos
\[
 T=\begin{pmatrix}
 A&B_1&B_2\\C_1&D_{11}&D_{12}\\C_2&D_{21}&D_{22}
 \end{pmatrix},\qquad Q_2=(I-uD_{22})^{-1}.
\]
Al eliminar la tercera componente quedan los cuatro bloques
\begin{equation}
 \begin{aligned}
 A'&=A+uB_2Q_2C_2,&
 B_1'&=B_1+uB_2Q_2D_{21},\\
 C_1'&=C_1+uD_{12}Q_2C_2,&
 D_{11}'&=D_{11}+uD_{12}Q_2D_{21}.
 \end{aligned}
 \label{mem:schur-intermedio}
\end{equation}

\begin{proposition}[Transitividad de la eliminación]
Si \(D_*=(D_{ij})_{i,j=1}^2\), el complemento obtenido en dos etapas
coincide con el obtenido al eliminar juntas ambas memorias:
\begin{equation}
 \begin{split}
 I-uA'-u^2B_1'(I-uD_{11}')^{-1}C_1'
 ={}&I-uA\\
 &-u^2(B_1\ B_2)(I-uD_*)^{-1}\binom{C_1}{C_2}.
 \end{split}
 \label{mem:schur-transitivo}
\end{equation}
\end{proposition}
\begin{proof}
Resolver la tercera ecuación de \((I-uT)w=v\) y sustituirla en las
dos primeras produce \eqref{mem:schur-intermedio}. Eliminar después la
segunda componente da el miembro izquierdo de
\eqref{mem:schur-transitivo}. Resolver a la vez las dos ecuaciones de
memoria da el derecho, por la proposición~\ref{mem:prop-schur}.
La solución formal es única porque \(I-uT\) tiene término constante
\(I\). En ambos procedimientos se transforman también la fuente y
el lector según \eqref{mem:schur-fuente}--\eqref{mem:schur-lector}.
\end{proof}

El balance \eqref{mem:balance} describe los incrementos del registro.
Las identidades anteriores permiten transportarlos al reutilizar la memoria
en la clausura dodecafásica; no determinan por sí solas coeficientes
analíticos adicionales.


% SOURCE foundation/sections/extension.tex
\section{Desarrollo del núcleo formal: condiciones iniciales, geometría y refinamiento}
\label{sec:extension}

El soporte de ochenta y una posiciones no coincide con el dominio de condiciones iniciales de la dinámica. Una condición inicial debe indicar además una orientación; y las hojas aditiva y multiplicativa disponen de cursores independientes. Esta distinción permite reconstruir exhaustivamente la emisión sin seleccionar de antemano una región numérica.

\subsection{Los estados iniciales y la recurrencia de emisión}

Para la enumeración utilizamos índices \(I_9=\{0,\ldots,8\}\), relacionados con las etiquetas positivas por \(i\mapsto i+1\). Sea
\[
 \mathcal S=I_9^2\times\{N,E,S,O\},\qquad
 \Omega_{\rm seed}=\mathcal S_+\times\mathcal S_\times.
\]
Cada condición inicial \(\omega=(s_+,s_\times)\) determina dos posiciones y dos direcciones. El dominio comprende todas las parejas ordenadas, por lo que
\begin{equation}
 |\mathcal S|=324,\qquad |\Omega_{\rm seed}|=324^2=104\,976.
 \label{eq:semillas}
\end{equation}
El censo de firmas del emisor residual, sección~\ref{app:firmas-emision}, enumera estas condiciones mediante identificadores y presenta sus fibras; aquí se define la ley que las transforma. Las filas de una tabla de emisiones son fibras de este dominio de condiciones iniciales.

La realización finita utiliza seis ventanas de nueve transiciones. En cada instante \(t\), la firma de fase es \(M_{\rm ph}(1+((t-1)\bmod9))\). Si vale \(+1\), se lee el residuo positivo \(\Sigma(i+1,j+1)=\rho_9(i+j+2)\) del primer cursor y después se desplaza ese cursor; si vale \(-1\), se lee el residuo positivo \(\Pi(i+1,j+1)=\rho_9((i+1)(j+1))\) del segundo cursor y después se desplaza el segundo. El valor cero incrementa el contador neutral. Al finalizar las transiciones veintisiete y cincuenta y cuatro, se invierten las orientaciones cardinales. Las evaluaciones enteras y sus cocientes se conservan en el registro, pero el emisor utiliza las lecturas residuales aquí especificadas.

En el sector de unidades, el lector multiplicativo utiliza el logaritmo discreto de base dos en \((\ZZ/9\ZZ)^\times\):
\[
 \ell_9(1)=0,\quad\ell_9(2)=1,\quad\ell_9(4)=2,\quad
 \ell_9(8)=3,\quad\ell_9(7)=4,\quad\ell_9(5)=5.
\]
Para la emisión finita se extiende con valor cero al sector no unitario. Esta extensión es un observable y no una identificación de los elementos \(3,6,9\) con la unidad multiplicativa. La hoja y su pertenencia al sector radical permanecen en el registro de trayectoria.

Sean \(S_m\) la suma de los tres residuos positivos aditivos \(\Sigma\) de la ventana \(m\), \(E_m\) la suma de los tres exponentes discretos de los residuos multiplicativos \(\Pi\) y \(Z_m=3\) el número de eventos neutrales. Así, \(S_m\) no suma las evaluaciones enteras \(S=i+j\), cuyos cocientes permanecen en el registro. Con \([a]_{10}\in\{0,\ldots,9\}\), la regla de emisión decimal es
\begin{align}
 d_{m,1}&=[S_m]_{10},\notag\\
 d_{m,2}&=[S_m+E_m]_{10},\notag\\
 d_{m,3}&=[3S_m+5E_m+7Z_m]_{10},\notag\\
 U_m&=100d_{m,1}+10d_{m,2}+d_{m,3}.
 \label{eq:emisor-seis}
\end{align}
Los coeficientes enteros y el calendario de esta recurrencia son parte explícita de la ley emisora declarada. No se reciben valores de \(\pi,\varphi,e\) ni de \(\alpha\). La distinción entre una regla de emisión y una constante reconocida a partir de su salida evita que la genealogía quede escondida en una tabla de decimales.

Si \(s_m=[S_m]_{10}\) y \(e_m=[E_m]_{10}\), la misma regla se escribe
\begin{equation}
 U_m=100s_m+10[s_m+e_m]_{10}+[3s_m+5e_m+1]_{10}.
 \label{eq:acoplamiento}
\end{equation}
El uno final es el residuo de \(7Z_m=21\) módulo diez. La reducción ternaria orientada del catálogo adopta la convención
\begin{equation}
 \Psi(U_1,\ldots,U_6)=([-U_1]_3,\ldots,[-U_6]_3).
 \label{eq:psi-catalogo}
\end{equation}
El signo de esta carta es explícito: cambiarlo modifica las palabras y exige transportar las orientaciones posteriores. La identificación balanceada de cada componente de \(\FF_3\) con el TRIT se realiza después de \eqref{eq:psi-catalogo}.

\begin{proposition}[Factorización inyectiva de la emisión]
La palabra \(U=(U_1,\ldots,U_6)\) determina las dos firmas \(s=(s_1,\ldots,s_6)\) y \(e=(e_1,\ldots,e_6)\). La imagen de la emisión tiene \(18\cdot26=468\) elementos. La proyección ternaria tiene \(243\) valores distintos en el ambiente de \(729\) palabras.
\end{proposition}
\begin{proof}
Las centenas de \(U_m\) recuperan \(s_m\). Si \(b_m\) es su cifra de decenas, entonces \(e_m=[b_m-s_m]_{10}\). La cifra de unidades comprueba la tercera congruencia de \eqref{eq:acoplamiento}. Por ello el acoplamiento es inyectivo en el producto de firmas.

La enumeración de las \(324\) condiciones de un cursor mediante la recurrencia anterior produce dieciocho firmas aditivas, cada una con dieciocho preimágenes; y veintiséis multiplicativas, de las cuales veinticuatro tienen ocho preimágenes, una tiene veinticuatro y una tiene ciento ocho. La sección~\ref{app:firmas-emision} reúne las firmas, sus multiplicidades y la regla exacta de recuperación de sus condiciones iniciales. La inyectividad del acoplamiento da \(18\cdot26\) emisiones y multiplicidades
\[
 \underbrace{432}_{18\cdot24}\text{ fibras de }144,
 \qquad18\text{ fibras de }432,
 \qquad18\text{ fibras de }1944.
\]
La suma es \(432\cdot144+18\cdot432+18\cdot1944=104\,976\). Aplicar \eqref{eq:psi-catalogo} a las \(468\) emisiones produce exactamente \(243\) palabras. Esta última cifra es el cardinal de la imagen, no el del espacio \(\FF_3^6\).
\end{proof}

Esta factorización explica la diferencia entre exhaustividad finita y profundidad ilimitada. El dominio \eqref{eq:semillas} es completo para el soporte y el emisor declarados; no contiene anticipadamente cada prefijo de cada prolongación. La profundidad pertenece al sistema de extensiones compatibles que se construye sobre esas condiciones.

\subsection{Clausura del calendario y avance de memoria}

El calendario de ciento ocho transiciones consta de doce ventanas nonádicas. Su estructura de grupo es más precisa que un producto informal de dos períodos:
\begin{equation}
 0\longrightarrow C_{12}\xrightarrow{\iota}C_{108}
 \xrightarrow{p}C_9\longrightarrow0,
 \quad\iota([b])=[9b],\quad p([t])=[t]_9.
 \label{eq:extension108}
\end{equation}
Si \(t=9b+r\) con \(0\le r<9\), la suma de dos fases produce el acarreo
\[
 (b,r)+(b',r')=
 \left(b+b'+\left\lfloor\frac{r+r'}9\right\rfloor,\ [r+r']_9\right),
\]
donde la primera coordenada se reduce módulo doce. La composición conserva el término que enlaza las dos coordenadas.

\begin{proposition}[No escisión del calendario]
La sucesión \eqref{eq:extension108} es exacta y no admite una sección que sea homomorfismo de grupos.
\end{proposition}
\begin{proof}
El núcleo de la reducción módulo nueve es el subgrupo de múltiplos de nueve, imagen de \(\iota\). Si la extensión se escindiera, por conmutatividad se tendría \(C_{108}\cong C_{12}\times C_9\). El primer grupo tiene exponente \(108\), mientras el segundo tiene exponente \(\operatorname{mcm}(12,9)=36\), contradicción.
\end{proof}

Nueve transiciones restauran la fase nonádica y avanzan una ventana; ciento ocho restauran el calendario finito. Ninguno de estos hechos obliga a borrar la trayectoria o a restaurar el prefijo de una prolongación. La holonomía nonádica \(\Gamma_9\) designa el transporte de una vuelta sobre estados con memoria, de modo que
\begin{equation}
 g(\Gamma_9x)=g(x),\qquad M(\Gamma_9x)=M(x)+1.
 \label{eq:holonomia-memoria}
\end{equation}
La igualdad de fase y el incremento de memoria son compatibles y constituyen dos observaciones del mismo transporte.

\input{sections/temporalidad_trit.tex}

La representación monodrómica de este retorno se desarrolla en la
sección~\ref{mono-ampl:conexion}. Allí se distingue la holonomía del
estado de su representación sobre fase y contador. La suspensión
simbólica de la sección~\ref{mono-ampl:cristal} define el modelo de
cristal temporal aperiódico, y la sección~\ref{mono-ampl:euler}
compone esa temporalidad con las tres realizaciones de Euler.

\subsection{Compresión observable y conservación ortogonal}

Una realización isométrica permite expresar de manera cuantitativa qué pierde una observación reducida. Sea \(U:\mathcal H\to\mathcal H\) una realización isométrica del transporte de una vuelta, \(U^*U=I\), y \(J:\mathcal H_{\rm vis}\to\mathcal H\) una inclusión isométrica. Definimos
\[
 T=J^*UJ,\qquad D=(I-JJ^*)UJ.
\]
Entonces
\begin{equation}
 UJ=JT+D,\qquad J^*D=0,\qquad I-T^*T=D^*D.
 \label{eq:conservacion-ortogonal}
\end{equation}
La última igualdad se deduce tomando el producto adjunto de la primera y usando la ortogonalidad. La contracción de la componente visible no constituye por sí sola destrucción de información del estado total.

\begin{proposition}[Conservación a horizonte finito]
Si \(I-T^*T=D^*D\), entonces, para todo \(n\ge1\),
\begin{equation}
 I=\sum_{r=0}^{n-1}T^{*r}D^*DT^r+T^{*n}T^n.
 \label{eq:telescopia}
\end{equation}
\end{proposition}
\begin{proof}
Multiplicar \(D^*D=I-T^*T\) por \(T^{*r}\) y \(T^r\) y sumar produce una suma telescópica. Se cancelan los términos interiores y queda \(I-T^{*n}T^n\).
\end{proof}

El último sumando conserva el estado terminal. Suprimirlo antes de justificar un límite alteraría la igualdad. Esta distinción es el contenido operativo de la conservación de la unidad en la compresión: la norma se distribuye entre la componente visible, la memoria y el terminal, sin ser reemplazada por un único residuo.

\subsection{Desarrollo local: conmutación, firma lorentziana y curvatura discreta}
\label{vii:sec:bloque-compresion}
\label{iv:extension-curvatura-mixta}

El bloque diagonal adyacente con diagonales iguales módulo nueve es único. En efecto,
\[
 k^2\equiv(k+1)^2\pmod9\quad\Longleftrightarrow\quad2k+1\equiv0\pmod9
\]
da \(k=4\). Su matriz visible es
\[
 B_c=\begin{pmatrix}7&2\\2&7\end{pmatrix}=7I+2R,
 \quad R=\begin{pmatrix}0&1\\1&0\end{pmatrix},\quad R^2=I.
\]
Con \(P_\pm=(I\pm R)/2\), se obtiene \(B_c=9P_++5P_-\). La unidad de la normalización estocástica y la conservación de una forma indefinida son dos realizaciones de este bloque:
\[
 \frac{B_c}{9}=P_++\frac59P_-,\qquad
 \frac{B_c}{\sqrt{45}}=\exp(\eta R),\quad\eta=\frac12\log\frac95.
\]
Para \(Q=\operatorname{diag}(1,-1)\), la segunda matriz satisface \(B_c^{\mathsf T}QB_c=45Q\). De aquí procede una realización local en \(SO^+(1,1)\); no se requiere introducir de antemano un espacio-tiempo de cuatro dimensiones para demostrarla.

El promediado uniforme sobre las clases \(C_1=\{1,4,7\}\), \(C_2=\{2,5,8\}\), \(C_0=\{3,6,9\}\) da
\[
 M_\Pi=\begin{pmatrix}4&5&6\\5&4&6\\6&6&9\end{pmatrix},\quad
 M_\Sigma=\begin{pmatrix}5&6&4\\6&4&5\\4&5&6\end{pmatrix}.
\]
El polinomio característico de \(M_\Pi\) es
\[
 (\lambda+1)((\lambda-9)^2-72),
\]
de donde su forma cuadrática tiene firma \((2,1)\). El salto espectral local \(9-5=4\), el coeficiente \(2\), la diferencia agregada \(51-45=6\) y el defecto integrado \(9\cdot6=54\) describen niveles distintos de la misma comparación aritmética. No son índices intercambiables de un mismo operador.

La curvatura discreta puede describirse directamente antes de esta realización espectral. En residuos, sean \(S(x,y)=x+y\), \(P(x,y)=xy\) y \(\Delta_x,\Delta_y\) las diferencias progresivas. Las variables de enlace son diferencias de potencial, \(A_x^T=\Delta_xT\), \(A_y^T=\Delta_yT\), y se fija
\[
 F(T_x,T_y)=\Delta_x\Delta_yT_y-\Delta_y\Delta_xT_x.
\]
Como \(\Delta_xS=\Delta_yS=1\), \(\Delta_xP=y\) y \(\Delta_yP=x\),
\begin{equation}
 F(S,S)=F(P,P)=0,\qquad F(S,P)=+1,\quad F(P,S)=-1\pmod9.
\end{equation}
La suma y el producto, tomados separadamente, inducen conexiones planas; la mezcla ordenada cambia el signo de la curvatura de plaqueta. Al sumar sobre un rectángulo de lados \(a,b\), las aristas interiores se cancelan y la suma orientada de borde vale \(\pm ab\) módulo nueve. Así aparece una relación exacta entre orden de composición y orientación de frontera.

\subsection{Completación cúbica y dos nociones de frontera}

La ampliación del soporte cúbico de lado nueve al de lado diez admite una estratificación posicional exacta:
\begin{equation}
 10^3=9^3+3\cdot9^2+3\cdot9+1=729+243+27+1.
 \label{eq:cubo}
\end{equation}
El primer término corresponde a las tres coordenadas interiores; el segundo, a una coordenada nueva; el tercero, a dos; y el último, al vértice de tres coordenadas nuevas. La corona de ampliación contiene \(271\) posiciones. En cambio, la frontera interna del cubo discreto de lado nueve contiene \(9^3-7^3=386\). La diferencia entre estos recuentos es una diferencia de dominio, no un conflicto de valores.

Dividir \eqref{eq:cubo} por \(1000\) da una partición normalizada de la unidad. La incidencia del cubo distingue seis caras, doce aristas y ocho vértices. El enlace ulterior con hexadas y octadas de Steiner requiere los mapas combinatorios correspondientes; los cardinales del cubo no los sustituyen. Esta completación aporta aquí la relación entre las bases \(729\) y \(1000\), mientras la prolongación excepcional se desarrollará después de la generación numérica, del registro dodecafásico y de las construcciones de \(\alpha\).

\input{sections/revision_emrh.tex}

\input{sections/geometria_modular.tex}
\input{sections/aritmetica_historias.tex}


% SOURCE foundation/sections/temporalidad_trit.tex
% Recuperación de 02_trit_geometrias y de la holonomía con memoria.
\subsubsection{Parametrización temporal de la compatibilidad y orientación trítica}
\label{trit-temporal-compatibilidad}

La interpretación temporal del TRIT se establece sobre las transiciones del
estado completo. Sea \(x\xrightarrow{\rm adm}y\) la relación de admisibilidad determinada
por las operaciones del TPK en \(\widehat X\). Una trayectoria parametrizada
por un conjunto discretamente ordenado \(\mathcal T\) es una aplicación
\[
 \gamma:\mathcal T\longrightarrow\widehat X,\qquad
 \gamma(t)\xrightarrow{\rm adm}\gamma(t^+),
\]
donde \(t^+\) es el sucesor de \(t\). La relación de compatibilidad precede a
la elección del parámetro: el orden permite expresar sus transiciones como
una evolución. Cuando se utiliza el término sección, se declara además la
proyección \(p:E\to\mathcal T\) y la condición
\(p\circ\gamma=\operatorname{id}_{\mathcal T}\). Así quedan especificados
el espacio de estados, el orden temporal y la información que se transporta.

La firma \(\tau\) selecciona un régimen cuadrático y admite una lectura
temporal relativa a esa parametrización. El tipo \(+1\), con
\(J_{+1}^{2}=-I\), representa retorno y memoria; el tipo \(0\), con
\(J_0^{2}=0\), representa el umbral de transición; el tipo \(-1\), con
\(J_{-1}^{2}=I\), representa apertura en dos direcciones propias. En la
terminología temporal del modelo, estas funciones corresponden,
respectivamente, a pasado, presente y futuro. La correspondencia se refiere
a la posición operacional de una transición dentro de una historia:
antecedente conservado, actualización y prolongación admisible.

Esta interpretación adquiere contenido mediante tres operaciones distintas.
La restricción de una trayectoria recupera el prefijo ya construido. La
actualización \(\mathcal U_t\) compone selección, transporte y registro en
la transición observada. La prolongación considera las continuaciones que
conservan las condiciones de frontera de ese prefijo. Si
\(\mathcal P_n(x)\) denota el conjunto de continuaciones admisibles de
longitud \(n\) desde \(x\), la supresión del último paso define
\[
 r_n:\mathcal P_{n+1}(x)\longrightarrow\mathcal P_n(x).
\]
Una continuación ilimitada es una familia \((\gamma_n)_{n\geq1}\) que
satisface \(r_n(\gamma_{n+1})=\gamma_n\). El sistema inverso conserva de
este modo las restricciones que cada profundidad impone sobre las
anteriores. La existencia de una familia compatible se establece en el
dominio de prolongación correspondiente; el mero establecimiento de los
conjuntos \(\mathcal P_n(x)\) no sustituye esa demostración.

La bidireccionalidad combina, por tanto, la reconstrucción de antecedentes
con la compatibilidad de continuaciones. Sobre la sucesión de firmas, la
involución \(C_{\rm rev}\) invierte el orden y cambia el signo de cada
componente. Sobre la trayectoria actúan, además, las transformaciones de
posición, hoja y frontera especificadas por el transporte. La firma
invertida es una observación de la trayectoria conjugada, no un reemplazo
de sus restantes coordenadas. El umbral permanece distinguido: el valor
visible \(0\) conserva su función de transición incluso cuando su cociente
o su registro histórico son distintos de cero.

La relación con las geometrías exponenciales es igualmente precisa. En
cada régimen fijo, \(\exp(tJ_\tau)\) compone parámetros y conserva la
norma algebraica. El cierre elíptico pertenece a esa representación local.
El retorno de fase de \(\Gamma_9\), en cambio, actúa sobre la trayectoria
con memoria, conforme a \eqref{eq:holonomia-memoria}. Una trayectoria puede
recuperar su fase y conservar un nuevo antecedente: el dato histórico ha
cambiado aunque la observación de fase sea la misma. El umbral parabólico
retiene un desplazamiento lineal; la apertura hiperbólica separa dos
direcciones propias de expansión y contracción. Estas tres realizaciones
proporcionan leyes de composición, mientras que el TPK determina su
selección y su concatenación efectiva.

La suspensión aperiódica de la sección \ref{sec:generacion} hará explícita
esta distinción mediante un factor de fase finito, una rotación irracional
y un grado de memoria. El vínculo entre TRIT y temporalidad queda así
expresado por operaciones sobre trayectorias y por sus representaciones
cuadráticas. Su contenido excede una clasificación de signos: conserva el
régimen, el sentido del transporte y la relación entre antecedente,
transición y continuación.


% SOURCE foundation/sections/revision_emrh.tex
\subsubsection{Estratificación de la completación y restitución del ápice}
\label{sec:revision-emrh}

La completación admite una descripción conjuntista que hace explícita su
operación. Sea \(E_9\) el conjunto de las nueve clases residuales, representadas
positivamente por \(1,\ldots,9\), y sea \(\ast\) un elemento adicional.
El símbolo \(\ast\) es distinto de la clase nula, ya representada por \(9\).
En consecuencia, la ampliación de \(E_9\) a
\(\widehat E=E_9\sqcup\{\ast\}\) aumenta el cardinal de nueve a diez
sin modificar la identificación modular interna. La completación cúbica es
\(\widehat C=\widehat E^3\), y el cubo inicial es \(C=E_9^3\).

\begin{definition}[Estratos de codimensión combinatoria]
Para \(0\leq r\leq3\), se define
\[
 S_r=\{(x_1,x_2,x_3)\in\widehat E^3:
             \#\{j:x_j=\ast\}=r\}.
\]
El ápice de la completación es \(a_\ast=(\ast,\ast,\ast)\), y la
completación perforada es \(\widehat C^{\circ}=\widehat C\setminus\{a_\ast\}\).
\end{definition}

\begin{proposition}[Descomposición exacta de la completación]
Los estratos son disjuntos y satisfacen
\begin{align}
 \widehat C&=S_0\sqcup S_1\sqcup S_2\sqcup S_3,
 &|S_r|&=\binom3r9^{3-r},\label{eq:revision-estratos}\\
 |\widehat C^{\circ}|&=729+243+27=999,
 &|\widehat C|&=999+1=1000.\label{eq:revision-999}
\end{align}
La ampliación perforada contiene \(270\) elementos adicionales al cubo
inicial; la ampliación completa contiene \(271\).
\end{proposition}
\begin{proof}
Un elemento de \(S_r\) se obtiene eligiendo las \(r\) coordenadas que
contienen \(\ast\), de \(\binom3r\) maneras, y asignando a las restantes
una de las nueve clases de \(E_9\). La clasificación por el número de
coordenadas adicionales es exhaustiva y sus clases son disjuntas. El único
elemento de \(S_3\) es \(a_\ast\). Retirarlo resta una unidad al cardinal
total; restituirlo recupera el producto cartesiano completo.
\end{proof}

La expresión \emph{extrema y media razón holográfica}, abreviada EMRH en
el corpus, designa aquí esta articulación entre el cuerpo inicial, los
estratos de ampliación y la restitución unitaria. Su contenido inmediato
es la partición exacta de un producto finito y su normalización. El número
\(999\) representa un objeto determinado, la completación perforada;
\(1000\) representa su clausura por restitución del ápice. La unidad
añadida tiene, por tanto, una residencia de incidencia explícita.

\begin{figure}[htbp]
\centering
\begin{tikzpicture}[x=1cm,y=1cm,>=Stealth,font=\small]
\draw[->] (0,0)--(12.4,0);
\foreach \x/\a/\b in {0/729/{S_0},4.0/243/{S_1},7.5/27/{S_2},10.6/1/{S_3}}{
  \draw (\x,0.12)--(\x,-0.12);
  \node[above=5pt] at (\x,0) {\(\a\)};
  \node[below=5pt] at (\x,0) {\(\b\)};
}
\draw[decorate,decoration={brace,amplitude=5pt}] (0,1)--(7.7,1);
\node[above=8pt] at (3.85,1) {Completación perforada: \(999\)};
\draw[decorate,decoration={brace,mirror,amplitude=5pt}] (0,-1)--(10.8,-1);
\node[below=8pt] at (5.4,-1) {Completación íntegra: \(1000=999+1\)};
\node[align=center] at (10.6,1.5) {Ápice\\\((\ast,\ast,\ast)\)};
\end{tikzpicture}
\caption{Estratos de la completación cúbica. La disposición representa
la descomposición conjuntista, no una escala métrica. Los \(243\) elementos
de \(S_1\) son tres estratos coordenados de \(81\) elementos; los \(27\)
de \(S_2\) son tres estratos de \(9\). El ápice completa la partición.}
\label{fig:revision-emrh}
\end{figure}

\subsubsection{Conservación de la unidad y compatibilidad dimensional}

La normalización no se limita a dividir una identidad aislada. En dimensión
\(d\), el producto \(\widehat E^d\) lleva la medida uniforme
\(\mu_d(\{x\})=10^{-d}\). La proyección coordenada
\(p_d:\widehat E^{d+1}\to\widehat E^d\) suprime la última coordenada.

\begin{proposition}[Compatibilidad proyectiva de las medidas normalizadas]
Para todo \(d\geq1\),
\[
 (p_d)_*\mu_{d+1}=\mu_d,
 \qquad \mu_d(\widehat E^d)=1.
\]
La masa del estrato con \(r\) coordenadas adicionales es
\(\binom dr9^{d-r}/10^d\), y la suma de las masas de todos los estratos
es uno.
\end{proposition}
\begin{proof}
Cada punto de \(\widehat E^d\) tiene diez preimágenes. Su masa total es
\(10\,10^{-(d+1)}=10^{-d}\). La igualdad para cualquier subconjunto
se obtiene sumando sobre sus puntos. La fórmula estratificada procede
del mismo recuento que \eqref{eq:revision-estratos}; el binomio
\((9+1)^d\) da la normalización total.
\end{proof}

En dimensión tres, retirar el ápice deja masa \(999/1000\), y
restituirlo añade \(1/1000\). Renormalizar antes de la restitución
definiría otra medida: la distinción es operacionalmente relevante.
En particular, conservación de masa probabilística, conservación de los
registros de transporte y disipación termodinámica son afirmaciones con
objetos diferentes. La proposición demuestra la primera y proporciona el
soporte compatible para estudiar la segunda; la tercera requiere una
realización dinámica y energética especificada.

También permanece diferenciada la frontera geométrica interna.
Sobre el representante cartesiano \(\{1,\ldots,9\}^3\), retirar
\(\{2,\ldots,8\}^3\) deja \(386\) puntos. Esta frontera pertenece al
cubo inicial, mientras \(\widehat C\setminus C\) pertenece a su
ampliación. La completación cúbica, la expansión posicional de base mil y
la lectura de incidencia excepcional comparten datos del soporte;
sus mapas respectivos determinan qué información de esos datos utiliza
cada representación.


% SOURCE foundation/sections/geometria_modular.tex
\subsection{Incidencia cúbica y realización lorentziana}

La completación cúbica aporta una relación entre interior y frontera.
Su geometría puede desarrollarse mediante la incidencia de las seis caras
con los ocho vértices. Etiquetemos las caras por \((i,\varepsilon)\),
\(i\in\{1,2,3\}\), \(\varepsilon\in\{+1,-1\}\), y los vértices por
\(\nu\in\{+1,-1\}^3\). La matriz de incidencia es
\[
 M_{(i,\varepsilon),\nu}=\mathbf1_{\{\nu_i=\varepsilon\}}.
\]
Sean \(u=(1,\ldots,1)^{\mathsf T}\in\RR^6\), \(R_F\) la oposición de
caras y \(P_u=uu^{\mathsf T}/6\), \(P_-=(I-R_F)/2\).

\begin{proposition}[Cociente de incidencia de dimensión cuatro]
El Gramiano de incidencia satisface
\[
 MM^{\mathsf T}=12P_u+4P_-.
\]
Sus autovalores \(12,4,0\) tienen multiplicidades \(1,3,2\).
En consecuencia,
\[
 Q_\square=\RR^6/\ker(MM^{\mathsf T})
 \simeq\RR u\oplus E_-,\qquad E_-=\operatorname{im}P_-.
\]
\end{proposition}
\begin{proof}
Una cara contiene cuatro vértices; dos caras opuestas tienen intersección
vacía, y dos caras distintas no opuestas comparten dos vértices.
La matriz \(12P_u+4P_-\) tiene precisamente esas entradas.
El espacio uniforme tiene dimensión uno; las diferencias entre caras
opuestas generan un espacio de dimensión tres, ortogonal al uniforme.
Su complemento tiene dimensión dos y pertenece al núcleo.
\end{proof}

La descomposición \(1+3\) procede de la incidencia. La elección temporal
se declara mediante la forma bilineal
\[
 B_\square([x],[y])=\langle x,(P_--P_u)y\rangle.
\]
Está bien definida en el cociente, porque ambos proyectores anulan el
núcleo. En la base
\[
 t=[u]/\sqrt6,\qquad
 d_i=[e_{i,+}-e_{i,-}]/\sqrt2,
\]
su matriz es \(\operatorname{diag}(-1,1,1,1)\). Quedan explicitados
el dominio, la descomposición y la convención de signo que producen
la realización de Minkowski. El Gramiano positivo determina el cociente;
la asignación temporal determina la forma indefinida sobre él.

El operador
\[
 W_\square=\sqrt6P_u+\sqrt2P_-,\qquad
 W_\square e_{i,\varepsilon}=t+\varepsilon d_i
\]
envía las seis etiquetas de cara a direcciones nulas, puesto que
\(B_\square(t+\varepsilon d_i,t+\varepsilon d_i)=-1+1=0\).
La inversión de una dirección y la oposición de caras disponen así
de una realización geométrica concreta.

\subsubsection{Soldadura, inversión y subdivisión}

La realización celular utiliza un transporte invertible \(U_m(a)\),
que conserva la forma lorentziana, entre las fibras de los extremos de
una arista dual orientada \(a\), una etiqueta
de cara \(\lambda_m(a)\) y una orientación \(\sigma_m(a)\).
Con \(\ell_m=\ell_0/9^m\), la aplicación de soldadura es
\[
 \beta_m(a)=\ell_m\sigma_m(a)
 W_{\square,s(a)}e_{\lambda_m(a)}.
\]
La inversión compatible satisface
\[
 \beta_m(\bar a)=-U_m(a)\beta_m(a).
\]
Identifiquemos las fibras del origen en dos niveles consecutivos mediante
la isometría de refinamiento declarada. Si los nueve subsegmentos
\(a_1,\ldots,a_9\), de nivel \(m+1\), conservan la etiqueta
y la orientación al ser transportados a esa fibra común, entonces
\begin{equation}
 \beta_m(a)=\sum_{j=1}^9
 U_{m+1}(a_1\cdots a_{j-1})^{-1}\beta_{m+1}(a_j).
 \label{eq:soldadura-refinamiento}
\end{equation}
Cada término transportado equivale a \(\ell_m/9\) por la misma
dirección nula, lo que demuestra la identidad. La hipótesis de
compatibilidad determina qué subdivisiones realizan este refinamiento.
La ecuación reúne dirección, escala y transporte, en lugar de limitar
la correspondencia geométrica a un recuento de dimensiones.

Fijadas además la orientación y la forma lorentziana, la dualidad de
Hodge en \(\Lambda^2Q_\square^*\) satisface \(\star^2=-I\).
En dimensión cuatro, el signo es \((-1)^{2(4-2)+1}=-1\).
Se obtiene una estructura compleja real sobre las dos-formas y,
tras complejificación, los subespacios propios de valores \(+i\) y
\(-i\). La realización electromagnética utilizará este dominio
después de construir la sección electrónica y declarar su acoplamiento.

\subsection{Compatibilidad entre las reducciones modular y posicional}

La base decimal por bloques interviene mediante una división euclídea
que conserva simultáneamente residuo y cociente:
\[
 n=1000q+r,\qquad 0\le r<1000.
\]
Puesto que \(1000\equiv1\pmod9\) y \(1000\equiv1\pmod3\),
\[
 n\equiv q+r\pmod9,\qquad n\equiv q+r\pmod3.
\]
Por tanto, la reducción de un bloque requiere la contribución del
cociente para recuperar la clase del entero. Los números \(0\) y
\(1000\) tienen el mismo resto por \(1000\), pero clases diferentes
por \(9\) y por \(3\). Ésta es una razón operatoria para conservar
la memoria durante el cambio de resolución.

El refinamiento ternario y la determinación de un prefijo decimal
pueden coordinarse mediante un único residuo entero. Si \(N/3^L\)
es un extremo de cilindro y \(D/1000^r\) un extremo decimal, definimos
\[
 A=1000^rN-D3^L.
\]
Al añadir seis símbolos ternarios, \(N'=729N+\nu\),
\(L'=L+6\), con \(0\le\nu<729\). Si el prefijo decimal permanece
fijo, se obtiene
\[
 A'=729A+\nu1000^r.
\]
Si se incorpora además un bloque decimal \(\delta\), entonces
\(D'=1000D+\delta\), \(r'=r+1\), y
\begin{equation}
 A'=729000A+\nu1000^{r+1}-729\delta3^L.
 \label{eq:residuo-dos-bases}
\end{equation}
Ambas identidades se demuestran sustituyendo las definiciones. La
comparación de extremos se reduce así a operaciones enteras que
retienen la contribución de cada prolongación. Esta compatibilidad
es la que permite refinar una región antes de evaluar su coordenada real.


% SOURCE foundation/sections/aritmetica_historias.tex
% Incorporación al final de extension.tex; no contiene una sección principal.
% Procedencia: RESULTADO_RECUPERADO; explicitación sectorial de las hipótesis.
% Fuentes leídas: monografía Doble proyección holográfica (cantera editorial),
% manuscrito/local/01_tesis_recuperada.tex y
% manuscrito/generated/algebra_holonomica.tex, especialmente §§1--7 y 21.
% Definición rectora de número: manuscrito/flattened/main_autosuficiente.tex,
% secciones que comienzan en líneas 22137 y 23000 de la cantera de 775 páginas.
% No se adopta su presentación global como prueba de todas las flechas TPK.

\subsection{El número como sección compatible y el valor como lectura}
\label{hist:seccion}

La distinción entre estado y observación permite precisar qué objeto se
prolonga cuando aumenta la resolución. A profundidad \(n\), sea \(X_n\)
el conjunto de estados admisibles producidos por APP--TRIT--TPK. Sus datos
incluyen posición, hojas suma y producto, residuos y cocientes, régimen,
orientación, acarreo, fase, ruta, frontera, memoria y supervivencia. No son
coordenadas independientes: las evaluaciones aritméticas y los transportes
anteriores imponen sus relaciones. Una prolongación conserva esas relaciones
en cada antecedente, aunque añada nuevos pasos y nuevos datos de memoria.
Éste es el sentido de la aritmética de historias desarrollada en
(secciones~\ref{sec:extension} y~\ref{mono-ampl:conexion}).

Sean \(r_{m,n}:X_m\to X_n\), para \(m\geq n\), las restricciones de
profundidad, con \(r_{n,n}=\operatorname{id}\) y
\(r_{\ell,n}=r_{m,n}\circ r_{\ell,m}\). Cada restricción conserva el
prefijo y la información ya determinada en él.

\begin{definition}[Sección numérica compatible]
Un número HMT es una historia compatible del sistema de estados:
\begin{equation}
 x=(x_n)_{n\geq0}\in X_\infty:=\varprojlim_n X_n,
 \qquad r_{m,n}(x_m)=x_n.
 \label{hist:limite-inverso}
\end{equation}
Un lector es una operación adicional sobre un dominio especificado de estas
historias. Una cifra es una observación finita; un valor escalar es una
evaluación del lector. El par formado por la sección y su lector constituye
una presentación leída, no una redefinición de la identidad de la sección.
\end{definition}

En el dominio arquimediano \(X_\infty^{\mathrm{Arch}}\), el lector asocia
a cada prefijo un cilindro cerrado, conserva su encajamiento y hace tender
su diámetro a cero. Así determina
\begin{equation}
 \operatorname{ev}_{\mathbb R}:X_\infty^{\mathrm{Arch}}\longrightarrow
 \mathbb R,\qquad
 \{\operatorname{ev}_{\mathbb R}(x)\}=\bigcap_n I_n(x).
 \label{hist:evaluacion}
\end{equation}
La construcción concreta de estos cilindros y la decisión de sus fronteras
se desarrollan en la sección~\ref{sec:generacion}. La definición
\eqref{hist:limite-inverso} no afirma por sí sola que la imagen de
\eqref{hist:evaluacion} sea toda la recta: una afirmación de cobertura
requiere su propio teorema. Tampoco una igualdad de valores establece, sin
un resultado de inyectividad, la igualdad de sus historias. Esta separación
permite hablar de precisión como profundidad conservando la diferencia
entre el antecedente aritmético y su realización escalar.

\subsection{Composición de historias y multiplicador de memoria}
\label{hist:algebra}

Fijemos un sector de estados a profundidad \(n\) y una categoría
\(\mathcal C\) cuyas flechas sean sus historias admisibles. La composición
se escribe \(vu=v\circ u\): primero se recorre \(u\) y después \(v\).
Los pasos elementales proceden de la selección, el transporte y la
actualización del TPK. Una fibra trítica local admite la escritura
\(\mathbb R[J_x]/(J_x^2+\tau(x)1_x)\), donde
\(\tau(x)\in\{+1,0,-1\}\). El régimen y la orientación pertenecen al
estado que se transporta; una transición entre regímenes no equivale a
identificar sus álgebras cuadráticas.

La siguiente formulación precisa un sector algebraico de esa composición.
Para cada objeto \(x\), se da un álgebra real asociativa unitaria \(A_x\),
que puede ser no conmutativa. Para cada \(u:x\to y\), se da un
homomorfismo unital \(\rho_u:A_x\to A_y\). Suponemos una acción estricta:
\begin{equation}
 \rho_{\operatorname{id}_x}=\operatorname{id}_{A_x},\qquad
 \rho_v\rho_u=\rho_{vu}.
 \label{hist:accion-estricta}
\end{equation}
Para cada pareja componible \(u:x\to y\), \(v:y\to z\), fijamos un
multiplicador central invertible
\(\Omega(v,u)\in Z(A_z)^\times\). Se exige la normalización
\(\Omega(\operatorname{id}_y,u)=\Omega(u,\operatorname{id}_x)=1_y\)
y, para cada triple componible, la identidad
\begin{equation}
 \rho_w\bigl(\Omega(v,u)\bigr)\Omega(w,vu)
   =\Omega(w,v)\Omega(wv,u).
 \label{hist:cociclo}
\end{equation}
Son hipótesis explícitas sobre los transportes y sus coeficientes, no una
consecuencia de llamar holonómica a la dinámica. Al aplicar esta presentación
a un sector TPK deben identificarse sus \(\rho_u\) y \(\Omega(v,u)\).
Las transiciones expresadas mediante bimódulos requieren conservar esa
estructura y no quedan reemplazadas por \eqref{hist:accion-estricta}.

El espacio de sumas finitas
\[
 \mathcal H(\mathcal C,A,\rho,\Omega)
   =\bigoplus_{u\in\operatorname{Mor}(\mathcal C)} A_{t(u)}T_u
\]
recibe el producto bilineal
\begin{equation}
 (aT_v)(bT_u)=a\rho_v(b)\Omega(v,u)T_{vu},
 \label{hist:producto}
\end{equation}
si las flechas son componibles, y cero en otro caso. Aquí \(t(u)\)
designa el estado terminal. La suma reúne historias con coeficientes; el
producto concatena historias conservando el multiplicador del registro.

\begin{proposition}[Asociatividad sectorial con multiplicador central]
Bajo \eqref{hist:accion-estricta} y \eqref{hist:cociclo}, el producto
\eqref{hist:producto} es asociativo. Si el conjunto de objetos es finito,
su unidad es \(\sum_x1_xT_{\operatorname{id}_x}\).
\end{proposition}
\begin{proof}
Para \(u:x\to y\), \(v:y\to z\), \(w:z\to t\), sean
\(a\in A_y\), \(b\in A_z\), \(c\in A_t\). Los dos coeficientes
de \(T_{wvu}\) son
\begin{align*}
 ((cT_w)(bT_v))(aT_u):\quad&
 c\rho_w(b)\Omega(w,v)\rho_{wv}(a)\Omega(wv,u),\\
 (cT_w)((bT_v)(aT_u)):\quad&
 c\rho_w(b)\rho_w\rho_v(a)\rho_w(\Omega(v,u))\Omega(w,vu).
\end{align*}
La acción estricta iguala \(\rho_w\rho_v(a)\) con \(\rho_{wv}(a)\).
La centralidad de \(\Omega(w,v)\) permite desplazarlo a la derecha de
ese factor; \eqref{hist:cociclo} iguala entonces los productos restantes.
No se permutan los coeficientes \(a,b,c\). Si el triple no es componible,
ambas parentizaciones son cero por sus estados iniciales y terminales.
La bilinealidad concluye la prueba. Las identidades locales y la
normalización de \(\Omega\) verifican la unidad indicada.
\end{proof}

El acarreo del calendario proporciona un ejemplo efectivo. En la categoría
de un objeto y flechas \(C_9\), tomemos coeficientes
\(A=\mathbb R[z,z^{-1}]\), acción identidad y
\[
 \Omega(b,a)=z^{c_0(a,b)},\qquad
 c_0(a,b)=\left\lfloor\frac{a+b}{9}\right\rfloor,
 \quad 0\leq a,b<9.
\]
Las dos sumas de acarreos de un triple son
\(\lfloor(a+b+c)/9\rfloor\); por ello se cumple
\eqref{hist:cociclo}. La variable formal \(z\) registra la vuelta sin
asignarle un valor numérico. Imponer después \(z^{12}=1\) recupera el
registro dodecafásico del calendario \eqref{eq:extension108}; conservar
\(z\) sin esa relación distingue vueltas sucesivas. El ejemplo realiza
la memoria de este calendario, sin identificarla con la totalidad de la
memoria TPK.

\subsection{Doble varianza del refinamiento y conservación de la unidad}
\label{hist:doble-varianza}

Al restringir estados se prolongan observables en sentido contrario.
Consideremos las álgebras de funciones, con producto puntual,
\(D_n=\operatorname{Fun}(X_n,\mathbb R)\). La precomposición define
\begin{equation}
 r_{m,n}^*:D_n\longrightarrow D_m,
 \qquad r_{m,n}^*f=f\circ r_{m,n}.
 \label{hist:precomposicion}
\end{equation}
Son morfismos unitales; son además inyectivos cuando las restricciones
correspondientes son sobreyectivas. El límite directo
\(D_{\mathrm{cyl}}=\varinjlim_n(D_n,r_{n+1,n}^*)\) reúne las
observaciones de profundidad finita. No se identifica aquí este álgebra
con toda el álgebra de rutas: prolongar los operadores de transporte exige
también sus compatibilidades con el refinamiento.

\begin{lemma}[Unidad y naturalidad de la evaluación]
Si \(e_x\) es el indicador de \(x\in X_n\), se cumple
\begin{equation}
 r_{m,n}^*e_x=\sum_{y:\,r_{m,n}(y)=x}e_y,
 \qquad r_{m,n}^*1_n=1_m,
 \label{hist:unidad}
\end{equation}
y, para todo \(f\in D_n\), \(y\in X_m\),
\[
 \langle r_{m,n}^*f,y\rangle=\langle f,r_{m,n}y\rangle.
\]
Cada sección \(x\in X_\infty\) determina una evaluación bien definida
\(D_{\mathrm{cyl}}\to\mathbb R\), dada por \([f]\mapsto f(x_n)\)
para \(f\in D_n\).
\end{lemma}
\begin{proof}
Las fibras de \(r_{m,n}\) forman una partición de \(X_m\). Sus
indicadores son ortogonales y suman \(1_m\), lo que prueba
\eqref{hist:unidad}. La naturalidad es la definición de precomposición.
Finalmente, \((r_{m,n}^*f)(x_m)=f(x_n)\), de modo que dos representantes
compatibles dan la misma evaluación en el límite directo.
\end{proof}

La unidad global es la función constante uno sobre todas las posibilidades;
la sección individual selecciona una historia compatible. Conservar la
primera no significa inmovilizar la segunda. El refinamiento puede aumentar
la cantidad de extensiones y la profundidad de cada historia manteniendo
la unidad y todas las evaluaciones de sus antecedentes.

\subsection{Retorno de fase, representación de Weyl y dos lecturas}
\label{hist:puente-lecturas}

La ley \eqref{eq:holonomia-memoria} adquiere una realización operatoria
en la suspensión bilateral que se construirá en la
sección~\ref{sec:generacion}. Allí \(T\) es el avance de un paso,
\(V_9\) observa la fase nonádica, \(W_\delta\) la fase de rotación y
\(D_M\) el grado de memoria. Escribiendo
\(\widehat\Gamma_9=T^9\) para la vuelta en esta representación,
las relaciones de Weyl \eqref{eq:weyl-relojes} implican, sobre vectores
de soporte finito,
\begin{equation}
 \begin{aligned}
 V_9\widehat\Gamma_9&=\widehat\Gamma_9V_9,\qquad
 W_\delta\widehat\Gamma_9=e^{2\pi i9\delta}
     \widehat\Gamma_9W_\delta,\\
 [D_M,\widehat\Gamma_9]&=\widehat\Gamma_9.
 \end{aligned}
 \label{hist:tres-retornos}
\end{equation}
La primera igualdad sigue de \(\zeta_9^9=1\); la segunda, de nueve
aplicaciones de la covariancia de \(W_\delta\); la tercera expresa
el incremento unitario de memoria. La fase finita retorna, mientras la
rotación irracional y la memoria distinguen el estado siguiente. Los
caracteres y el parámetro \(\delta\) se realizan después desde la
construcción numérica allí especificada: no seleccionan retrospectivamente
las semillas. Esta representación describe observables y traslaciones de
la suspensión; no reemplaza todas las coordenadas ni todos los operadores
del TPK.

La doble lectura prolonga la misma distinción entre historia y observable.
En su dominio arquimediano se aplica \eqref{hist:evaluacion}; en el
dominio incidencial calibrado, los mapas \(Q_n\) de la
sección~\ref{exc:doble-lectura} conservan las palabras y los marcos de
frontera. La igualdad allí demostrada
\(r_{n,m}Q_n=Q_mp_{n,m}\) determina su prolongación \(Q_\infty\).
Sobre el dominio común de ambas lecturas pueden considerarse conjuntamente
\(\operatorname{ev}_{\mathbb R}(x)\) y \(Q_\infty(x)\): la
primera retiene el valor; la segunda, la incidencia transportada que su
codominio conserva. La sección antecede a ambas. Esta articulación permite
pasar de la aritmética de historias a las constantes y a la lectura
excepcional sin convertir sus realizaciones en nuevos datos iniciales.


% SOURCE foundation/sections/continuo_conjunto.tex
\section{Construcción correlativa del continuo: historias, transporte, incidencia y terminal}
\label{iv:continuo-conjunto}
La aritmética publicada por una forma normal conserva una lectura precisa del estado que la produce. Para situar esa lectura dentro de HMT es necesario mantener también las historias, las hojas y las operaciones que actúan sobre ellas. Este capítulo reúne ese sustrato común. Sus distintas construcciones se obtienen del mismo árbol de ejecuciones APP–TRIT–TPK; los cortes de exposición permiten estudiar sus propiedades sin convertirlas en dominios independientes elegidos por una aplicación posterior.

El punto de partida es el núcleo formal expuesto anteriormente. Una emisión admisible contiene la operación de APP, la orientación del TRIT, el cambio de hoja, el residuo y el cociente, junto con la actualización efectiva de memoria del TPK. Cuando se publica solamente un residuo, la historia de ejecución sigue perteneciendo al objeto de partida. Dos ejecuciones con idéntica lectura final no se identifican por ese solo motivo.

\subsection{Historias admisibles y memoria por composición}
\label{iv:cont:historias}
Sea \(X_k\) el conjunto de historias admisibles de longitud \(k\). Sus elementos se construyen por emisiones sucesivas desde el dominio inicial del núcleo. Escribimos \(h\varepsilon\) para la prolongación de \(h\) por una emisión admisible \(\varepsilon\), y \(\rho_k(h\varepsilon)=h\) para la supresión de la última emisión. Se conserva la etiqueta de la emisión, incluso cuando diferentes etiquetas producen un mismo estado publicado.

En las dos hojas aritméticas, los datos locales satisfacen las divisiones exactas

\[
\delta_\diamond=r_\diamond+9q_\diamond,
\qquad 0\leq r_\diamond<9,
\qquad \diamond\in\{\Sigma,\Pi\}.
\]
La coordenada orientada del TRIT admite, análogamente, una escritura \(\Delta=o_r+3\kappa\), con el representante orientado declarado por el núcleo. La conservación del cociente permite componer esas publicaciones sin sustituir una emisión por su residuo aislado.

En una carta de memoria, la actualización de una emisión se escribe

\[
U_\varepsilon(m)=C_\varepsilon m+\kappa_\varepsilon.
\]
Para una trayectoria \(\alpha=(\varepsilon_1,\ldots,\varepsilon_k)\), la operación efectiva es la composición ordenada de esas actualizaciones. Su parte lineal y su traslación son

\[
C_\alpha=C_{\varepsilon_k}\cdots C_{\varepsilon_1},
\qquad
\kappa_\alpha=
\sum_{j=1}^{k}C_{\varepsilon_k}\cdots
C_{\varepsilon_{j+1}}\kappa_{\varepsilon_j}.
\]
El producto vacío de la última suma es la identidad.

\begin{proposition}[Compatibilidad afín de la concatenación]
\label{iv:cont:concatenacion}
Si \(\alpha\) precede a \(\beta\), entonces

\[
C_{\beta\alpha}=C_\beta C_\alpha,
\qquad
\kappa_{\beta\alpha}=C_\beta\kappa_\alpha+\kappa_\beta,
\qquad
U_{\beta\alpha}=U_\beta U_\alpha.
\]
\end{proposition}
\begin{proof}
Sustituir \(U_\alpha(m)\) en \(U_\beta\) produce \(C_\beta C_\alpha m+C_\beta\kappa_\alpha+\kappa_\beta\). La descomposición en parte lineal y traslación da las tres identidades. La asociatividad de la composición demuestra que el resultado no depende de una subdivisión auxiliar de la trayectoria.
\end{proof}

Esta identidad de cociclo es el contenido preciso de la memoria acumulada. La cantidad transportada por una segunda trayectoria depende de la actualización de carta que ésta realiza; no es, en general, la suma sin transporte de dos registros.

El retorno nonádico conserva una segunda distinción. En las coordenadas de fase y vuelta, con \(\bar r\in\{0,\ldots,8\}\), su publicación es

\[
\sigma_9(w,r)=
\left(w+\left\lfloor\frac{\bar r+1}{9}\right\rfloor,
r+1\pmod9\right),
\qquad
\sigma_9^9(w,r)=(w+1,r).
\]
Durante nueve pasos se produce exactamente un cruce de \(8\) a \(0\). La fase retorna y la coordenada de vuelta aumenta. La holonomía enriquecida retiene además las actualizaciones de memoria que acompañan ese recorrido.

\subsection{Refinamiento del árbol y medida de historias}
\label{iv:cont:medida}
Se fija el conjunto inicial finito no vacío de semillas admisibles, con una masa inicial estrictamente positiva y de suma \(1\). En este apartado cada historia tiene un número finito y positivo de prolongaciones. La supresión de emisiones define las aplicaciones de refinamiento. La admisibilidad del núcleo proporciona las prolongaciones; cuando una residencia utiliza una emisión neutra, ésta da explícitamente una sección de \(\rho_k\). La existencia de esa sección se refiere al árbol que incluye dicha emisión, no a cualquier restricción arbitraria de rutas.

Sea \(d(h)\geq1\) el número de emisiones admisibles en \(h\). La elección equiprobable entre sus hijos define

\[
\mu(h\varepsilon)=\frac{\mu(h)}{d(h)}.
\]
Más generalmente, una distribución local positiva \(p(\varepsilon\mid h)\), cuya suma es \(1\), define \(\mu(h\varepsilon)=\mu(h)p(\varepsilon\mid h)\). Esta segunda formulación permite conservar las ponderaciones del lector cuando los números de prolongaciones no son uniformes.

\begin{proposition}[Compatibilidad cilíndrica y esperanza condicional]
\label{iv:cont:esperanza}
En cada nivel se cumple

\[
\sum_{\rho_k(y)=x}\mu(y)=\mu(x).
\]
En los espacios \(L^2(X_k,\mu_k)\), las aplicaciones

\[
(J_kf)(y)=f(\rho_k(y)),
\qquad
(E_kg)(x)=\sum_{\rho_k(y)=x}\frac{\mu(y)}{\mu(x)}g(y)
\]
satisfacen \(E_k=J_k^*\), \(E_kJ_k=I\) y \(\|J_kf\|=\|f\|\).

\end{proposition}
\begin{proof}
La primera identidad es la normalización de las probabilidades locales. Aplicándola a \(|f(x)|^2\) se obtiene la isometría. Al agrupar por padres la suma que define \(\langle J_kf,g\rangle\), se obtiene la fórmula de \(E_k\); sobre una función constante en cada fibra, esa fórmula devuelve su valor.
\end{proof}

Las masas compatibles definen una medida de probabilidad sobre el espacio de historias infinitas: el cilindro de una historia finita recibe su masa \(\mu(h)\). La medida se extiende desde la álgebra de cilindros, y su unicidad procede de que éstos generan la sigma-álgebra. Si el árbol es finitamente ramificado y cada historia tiene una prolongación, el límite inverso de sus niveles por supresión de emisiones es no vacío. En la topología cilíndrica, cada cilindro no vacío tiene medida positiva.

Esta medida distingue historias que convergen a una misma publicación. La medida uniforme sobre valores terminales sería una lectura distinta, pues eliminaría las multiplicidades que la construcción acaba de conservar.

\subsection{Transporte por hojas y contabilidad exacta de las salidas}
\label{iv:cont:transporte}
Consideremos primero el espacio con base ortonormal formada por las historias del nivel \(k\). El transporte que conserva la etiqueta de cada prolongación es

\[
\mathcal U_ke_h=
\sum_{\varepsilon\ \mathrm{admisible\ en}\ h}
\sqrt{p(\varepsilon\mid h)}\,e_{h\varepsilon}.
\]
Los conjuntos de hijos de padres distintos son disjuntos. La normalización de \(p\) demuestra, por tanto, \(\mathcal U_k^*\mathcal U_k=I\). La representación mediante masas de historias se relaciona con ésta por el cambio de base diagonal que incorpora \(\sqrt{\mu(h)}\).

La etiqueta de hoja forma parte del registro completo de una historia. Se escribe
\[
\widehat h=((h_0,s_0),\ldots,(h_k,s_k)),\qquad
s_j\in\{\Sigma,\Pi\},
\]
donde \(s_j\) es la hoja activa en el paso \(j\). La prolongación añade la emisión y la hoja \(s_{k+1}\) que determina el selector operativo del TPK, conservando íntegro el prefijo, incluida la hoja inicial. En la base de esas historias levantadas,
\[
\mathcal U_ke_{\widehat h}
=\sum_{\varepsilon\ \mathrm{admisible\ en}\ \widehat h}
\sqrt{p(\varepsilon\mid\widehat h)}\,e_{\widehat h\varepsilon},
\qquad
P_ke_{\widehat h}=\mathbf1_{s_k=\Sigma}e_{\widehat h},\quad
Q_k=I-P_k.
\]
Así, \(P_k\) lee la hoja actual, mientras la distinción entre historias se mantiene por su prefijo completo. Dos historias iniciales que lleguen a la misma hoja activa siguen teniendo hijos distintos. La prueba de isometría por familias disjuntas de prolongaciones se aplica literalmente a esta base.
La orientación trítica es otra coordenada del registro: el valor \(0\) conserva la lectura de umbral en cada hoja y no añade una tercera hoja a esta descomposición. Separamos ahora la conservación de hoja y el intercambio de hoja:

\[
A_k=P_{k+1}\mathcal U_kP_k+Q_{k+1}\mathcal U_kQ_k,
\qquad
\eta_k=Q_{k+1}\mathcal U_kP_k+P_{k+1}\mathcal U_kQ_k.
\]
\begin{proposition}[Identidad de Gram por hojas]
\label{iv:cont:gram}
Para un transporte lineal \(U\), sea o no isométrico, la descomposición anterior satisface

\[
A^*A+\eta^*\eta=PU^*UP+QU^*UQ.
\]
En particular, para \(U=\mathcal U_k\),

\[
A_k^*A_k+\eta_k^*\eta_k=I.
\]
\end{proposition}
\begin{proof}
En \(A^*A\) y \(\eta^*\eta\), los términos mixtos que contienen \(P_{k+1}Q_{k+1}\) se anulan. Al sumar los términos restantes aparece \(P_kU^*(P_{k+1}+Q_{k+1})UP_k\), y su análogo con \(Q_k\). La primera identidad resulta de la complementariedad de los proyectores; la segunda utiliza además \(U^*U=I\).
\end{proof}

La primera identidad es la fórmula general. Sustituir su lado derecho por \(U^*U\) requiere que los bloques cruzados de ese Gram se anulen. La conservación isométrica del transporte de historias proporciona aquí esa condición; no debe atribuirse automáticamente a una compresión posterior.

Escribamos \(F_{0,0}=I\), \(F_{0,n}=A_{n-1}\cdots A_0\) y \(T_N=F_{0,N}^*F_{0,N}\).

\begin{theorem}[Descomposición finita con residuo terminal]
\label{iv:cont:terminal}
Para cada horizonte \(N\),

\[
I=\sum_{n=0}^{N-1}
F_{0,n}^*\eta_n^*\eta_nF_{0,n}+T_N.
\]
La sucesión \(T_N\) es positiva y decreciente, y posee un límite fuerte positivo \(T_\infty\). La suma de salidas converge fuertemente a \(I-T_\infty\).

\end{theorem}
\begin{proof}
La proposición~\ref{iv:cont:gram} da \(T_n-T_{n+1}=F_{0,n}^*\eta_n^*\eta_nF_{0,n}\). La suma telescópica prueba la igualdad finita y la monotonía. Para cada vector, las formas cuadráticas decrecen a un límite; la polarización define una forma acotada positiva y, por el teorema de representación de formas acotadas, un operador \(T_\infty\). La convergencia es fuerte porque, para un operador positivo \(S\leq I\), \(\|Sv\|^2\leq\langle v,Sv\rangle\); se aplica a \(S=T_N-T_\infty\).
\end{proof}

El residuo terminal se conserva hasta demostrar su extinción en el dominio considerado. Esta distinción impide que un retorno de fase, por sí solo, se convierta en una hipótesis de agotamiento de todas las historias.

\subsection{Balance orientado y cancelación en el refinamiento}
\label{iv:cont:balance}
Una historia lleva su orientación acumulada \(o(h)\) y los recuentos de sus visitas a las dos hojas. El balance local correspondiente es

\[
b(h)=o(h)\bigl(N_\Sigma(h)-N_\Pi(h)\bigr),
\qquad
b^+=\frac{|b|+b}{2},\quad b^-=\frac{|b|-b}{2}.
\]
Cuando un lector publica la media sobre las prolongaciones, se define \(b_c=Eb_f\). La orientación sigue estando en el argumento de \(b_f\); no se elimina antes de promediar.

\begin{proposition}[Defecto exacto de cancelación]
\label{iv:cont:cancelacion}
La cantidad

\[
\kappa_{\mathrm{can}}=
\frac{E|b_f|-|Eb_f|}{2}
\]
es no negativa y satisface

\[
E(b_f^+)=(b_c)^++\kappa_{\mathrm{can}},
\qquad
E(b_f^-)=(b_c)^-+\kappa_{\mathrm{can}}.
\]
\end{proposition}
\begin{proof}
La desigualdad triangular ponderada da \(|Eb_f|\leq E|b_f|\). Al escribir las partes positiva y negativa mediante \((|b|\pm b)/2\), se obtienen ambas igualdades.
\end{proof}

La información omitida al publicar solamente el balance es la cancelación común de las contribuciones opuestas. Esta cantidad se determina desde las mismas prolongaciones utilizadas por el transporte anterior. El factor \(1/9\) aparece sólo cuando la fibra contiene nueve hijos con pesos iguales; la esperanza condicional es la regla que sigue siendo válida para ramificación y ponderaciones generales.

\subsection{Incidencia ternaria de las hojas y complejo rectangular}
\label{iv:cont:incidencia-rectangular}
Las dos hojas determinan el módulo \(V=\mathbb F_3e_\Sigma\oplus\mathbb F_3e_\Pi\). Sus cuatro direcciones proyectivas son las rectas generadas por \(e_\Sigma,e_\Pi,e_\Sigma+e_\Pi,e_\Sigma-e_\Pi\). De ellas proceden seis pares no ordenados. Los ocho vectores no nulos son sus representantes orientados. Así, los recuentos \(4,6,8\) pertenecen a operaciones distintas sobre un mismo módulo de hojas.

Sobre las seis posiciones de pares, definimos

\[
(Ba)_{\{L,L'\}}=a_L+a_{L'},
\qquad
s(q)=\sum_{\{L,L'\}}q_{\{L,L'\}}.
\]
Cada dirección interviene en tres pares; por ello \(sB=0\) en \(\mathbb F_3\). El lector \(W(q)=\mathbf1_{s(q)=0}-1/3\) es invariante bajo \(q\mapsto q+Ba\). Sobre el ambiente uniforme \(\mathbb F_3^6\), sus momentos son

\[
\mathbb EW=0,\qquad
\mathbb EW^2=\frac29,\qquad
\mathbb EW^3=\frac2{27}.
\]
En efecto, \(s\) es sobreyectiva y cada fibra contiene \(3^5\) elementos. El lector vale \(2/3\) con probabilidad \(1/3\) y \(-1/3\) con probabilidad \(2/3\). Estas medias describen el ambiente uniforme; la distribución inducida por un subconjunto de historias debe calcularse con sus propias multiplicidades, ya conservadas en \(X_k\).

Para fijar los puertos sin perder multiplicidades, a profundidad \(k\) se conservan las lecturas etiquetadas \(\lambda_\Sigma(h)=(h,\Sigma)\) y \(\lambda_\Pi(h)=(h,\Pi)\). Sus conjuntos son
\[
I_k=\{\lambda_\Sigma(h):h\in X_k\},\qquad
J_k=\{\lambda_\Pi(h):h\in X_k\}.
\]
Son finitos y no vacíos porque lo es \(X_k\). Los valores, cocientes y residuos se publican sobre esos puertos; una igualdad de valores no identifica sus etiquetas. El producto \(I_k\times J_k\) es el dominio de la lectura rectangular, dentro del cual se registra el soporte de las incidencias admisibles.

Fijado este nivel, escribimos \(I=I_k\), \(J=J_k\) y \(A_N=\mathbb Z/3^N\mathbb Z\). Elegimos puertos de referencia \(i_0,j_0\). Esta elección sólo proporciona coordenadas para las fórmulas de reconstrucción. Las aplicaciones

\[
\kappa(u,v)_{ij}=u_i-v_j,
\qquad
(\delta w)_{ij}=w_{ij}-w_{ij_0}-w_{i_0j}+w_{i_0j_0}
\]
definen la siguiente secuencia, válida sobre el anillo finito \(A_N\):

\[
0\longrightarrow A_N\longrightarrow
A_N^I\oplus A_N^J\mathrel{\mathop{\longrightarrow}^{\kappa}}
A_N^{I\times J}\mathrel{\mathop{\longrightarrow}^{\delta}}
A_N^{(I\setminus\{i_0\})\times(J\setminus\{j_0\})}
\longrightarrow0.
\]
\begin{proposition}[Exactitud de la incidencia rectangular]
\label{iv:cont:rectangular-exacta}
La secuencia anterior es exacta, con primera flecha dada por la diagonal constante.

\end{proposition}
\begin{proof}
La igualdad \(u_i-v_j=0\) para cada par obliga a que todas las coordenadas de \(u\) y \(v\) sean una misma constante. La sustitución directa da \(\delta\kappa=0\). Si \(\delta w=0\), tómense \(u_i=w_{ij_0}\) y \(v_j=w_{i_0j_0}-w_{i_0j}\). Entonces \(u_i-v_j=w_{ij}\). Finalmente, cualquier matriz del último término se prolonga con fila y columna de referencia nulas; su imagen por \(\delta\) es la matriz inicial. La prueba sólo utiliza sumas y restas, por lo que no exige dividir por una unidad del anillo.
\end{proof}

Reducir de \(A_{N+1}\) a \(A_N\), manteniendo fijos \(k,I_k,J_k,i_0,j_0\), conmuta con las dos aplicaciones. Las fórmulas de reconstrucción son las mismas en cada profundidad modular. Ésta es la compatibilidad aritmética demostrada en este apartado. La prolongación de historias \(k\to k+1\) conserva las etiquetas de prefijo; sus aplicaciones entre conjuntos de puertos son un dato distinto del cambio de anillo \(N\to N+1\).

Esas aplicaciones se obtienen de la supresión de emisiones:
\[
\rho_k^\Sigma(h\varepsilon,\Sigma)=(h,\Sigma),\qquad
\rho_k^\Pi(h\varepsilon,\Pi)=(h,\Pi).
\]
Se fija una historia de referencia con prolongaciones compatibles, cuya existencia se demostró para el árbol sin hojas terminales. Sus dos etiquetas definen puertos de referencia compatibles en todos los niveles. El transporte de coordenadas es la precomposición:
\[
(L_ku)_{i'}=u_{\rho_k^\Sigma(i')},\quad
(L_kv)_{j'}=v_{\rho_k^\Pi(j')},\quad
(L_kw)_{i'j'}=w_{\rho_k^\Sigma(i'),\rho_k^\Pi(j')}.
\]
Para el último módulo de la secuencia, se prolonga primero la matriz por cero en la fila y la columna de referencia, se aplica esta misma fórmula y se restringe al nuevo complemento de los puertos de referencia. Esta regla define \(L_k\) también cuando un puerto nuevo tiene por padre el puerto marcado.

\begin{proposition}[Naturalidad del complejo rectangular bajo prolongación]
\label{iv:rectangular-natural}
Las aplicaciones anteriores verifican
\(L_k\kappa_k=\kappa_{k+1}L_k\) y
\(L_k\delta_k=\delta_{k+1}L_k\).
Además conmutan con la reducción de coeficientes
\(A_{N+1}\to A_N\) y se componen conforme a la supresión de prefijos.
\end{proposition}
\begin{proof}
La primera identidad se obtiene de
\(u_{\rho_k^\Sigma(i')}-v_{\rho_k^\Pi(j')}\).
En la segunda, cada uno de los cuatro términos que define \(\delta\)
se transporta por la misma precomposición. La compatibilidad de los
puertos marcados identifica sus filas y columnas; cuando un padre es
el puerto marcado, los cuatro términos se cancelan y coinciden con
la prolongación por cero. Todas las fórmulas son sumas y restas con
coeficientes enteros, de modo que conmutan con la reducción modular.
La composición de precomposiciones es la precomposición con el
prefijo compuesto, lo que demuestra la última afirmación.
\end{proof}

\subsection{Complejo de memoria y compatibilidad de las incidencias}
\label{iv:cont:memoria}
La misma coordenada acumulada de memoria \(c_y\in\mathbb Z\), en una residencia \(y\), define un complejo libre entero. Su reducción a \(A_N\) utiliza el mismo registro entero antes de reducirlo; por tanto, las profundidades sucesivas son compatibles. Sus generadores son \(f_y\) en grado \(2\), \(e_y,b_y,r_y\) en grado \(1\), y \(g_y\) en grado \(0\):

\[
\partial f_y=3c_y e_y+b_y,\qquad
\partial e_y=g_y,\qquad
\partial b_y=-3c_y g_y,\qquad
\partial r_y=0.
\]
La identidad \(\partial^2=0\) se verifica sobre \(f_y\) mediante \(3c_yg_y-3c_yg_y=0\), y es inmediata en los otros generadores.

Si una trayectoria \(\alpha:y\to y'\) actualiza la memoria, su transporte en el complejo envía los generadores homónimos a sus nuevos representantes, salvo

\[
\Psi_\alpha(b_y)=b_{y'}+3(c_{y'}-c_y)e_{y'}.
\]
\begin{proposition}[Naturalidad del complejo de memoria]
\label{iv:cont:memoria-natural}
Se cumplen \(\partial\Psi_\alpha=\Psi_\alpha\partial\) y \(\Psi_{\beta\alpha}=\Psi_\beta\Psi_\alpha\). Las aplicaciones se prolongan linealmente sobre el anillo de coeficientes elegido.

\end{proposition}
\begin{proof}
Sobre \(f_y\), la imagen de su borde es \(3c_ye_{y'}+b_{y'}+3(c_{y'}-c_y)e_{y'} =3c_{y'}e_{y'}+b_{y'}\). Sobre \(b_y\), el borde de su imagen es \(-3c_{y'}g_{y'}+3(c_{y'}-c_y)g_{y'}=-3c_yg_{y'}\). Los otros dos casos son inmediatos. En una concatenación, las diferencias de memoria se suman telescópicamente, lo que prueba la segunda identidad.
\end{proof}

La relación entre dos representantes de una misma residencia también puede escribirse como un borde explícito. Para una coordenada entera \(v\), reducida en el mismo anillo cuando corresponda, sean

\[
z^-=r_y,\qquad z^+=r_y+3c_yv e_y+v b_y,
\qquad h_*=-vf_y.
\]
Entonces \(\partial h_*=z^--z^+\). Bajo un transporte que actualiza también \(v\), la corrección \(K_\alpha=(v_{y'}-v_y)f_{y'}\) satisface \(K_{\beta\alpha}=K_\beta+\Psi_\beta K_\alpha\). Los símbolos \(K_\alpha\) de esta homotopía local no designan el sello dodecafásico \(K\): pertenecen a otro dominio y su función está especificada por la identidad anterior. Su función puede comprobarse también en la igualdad

\[
z^+_{y'}-\Psi_\alpha z^+_y=\partial K_\alpha.
\]
Ambos miembros valen \((v_{y'}-v_y)(3c_{y'}e_{y'}+b_{y'})\).

\begin{proposition}[Homología y memoria del complejo local]
\label{iv:cont:memoria-homologia}
Sobre \(\mathbb Z\), o después de reducir a \(A_N\), se tiene

\[
H_0(\Gamma)=H_2(\Gamma)=0,
\qquad H_1(\Gamma)\cong A_N[r_y]
\]
en la realización modular, y \(H_1(\Gamma)\cong\mathbb Z[r_y]\) en la realización entera.

\end{proposition}
\begin{proof}
Defínase la homotopía de grado \(1\) mediante \(h(g_y)=e_y\), \(h(b_y)=f_y\), y cero en los otros generadores. Sean \(p\) la proyección al sumando generado por \(r_y\) e \(\iota\) su inclusión. La comprobación sobre cada generador da

\[
\partial h+h\partial=I-\iota p.
\]
En particular, sobre \(b_y\) los términos \(3c_ye_y\) y \(-3c_ye_y\) se cancelan. El complejo se retrae por homotopía al módulo generado por \(r_y\), concentrado en grado \(1\), lo que demuestra la afirmación.
\end{proof}

El transporte \(\Psi_\alpha\) es una cizalla de determinante \(1\). Así, este complejo conserva explícitamente cómo cambia \(c_y\) a nivel de cadenas, pero su homología local no distingue los valores de esa coordenada. La memoria completa permanece en las etiquetas y en los transportes del estado fuente; no se recupera a partir de \(H_1\) aisladamente. Una torsión numérica dependiente de memoria requiere el ensamblaje y la lectura determinantal especificados, además de este complejo local.

\subsection{Ensamblaje con caminos terminales}
\label{iv:cont:ensamblaje}
Para un horizonte finito, las residencias y sus prolongaciones forman una categoría de caminos. A cada residencia \(\nu\) se asigna el módulo libre \(W(\nu)=A_N[\operatorname{Path}(\nu,\mathrm{Term})]\), que conserva cada continuación terminal como generador. Una trayectoria inicial actúa en \(W\) por precomposición. El complejo \(\Gamma(\nu)\) del apartado anterior actúa covariantemente mediante \(\Psi\).

El ensamblaje utiliza ambos transportes en el mismo módulo bigraduado:

\[
\mathcal B_{q,r}=
\bigoplus_{\nu_0\to\cdots\to\nu_q}
W(\nu_q)\otimes_{A_N}\Gamma_r(\nu_0).
\]
Se utiliza el complejo de caminos normalizado: se omiten las cadenas degeneradas que contienen identidades. El horizonte finito acota entonces el número de prolongaciones estrictas y el grado horizontal. El borde horizontal suprime una residencia intermedia componiendo las flechas adyacentes. En el primer extremo transporta \(\Gamma\); en el último precompone \(W\). Con los signos alternados se obtiene el borde de caminos \(d_{\mathrm{bar}}\). El diferencial total es

\[
D=d_{\mathrm{bar}}+(-1)^q\partial_\Gamma.
\]
Las identidades de caras cancelan por pares los términos de \(d_{\mathrm{bar}}^2\). La naturalidad de la proposición~\ref{iv:cont:memoria-natural} implica que \(d_{\mathrm{bar}}\) conmuta con \(\partial_\Gamma\); el factor \((-1)^q\) cancela los términos cruzados del cuadrado. Por tanto, \(D^2=0\). La compatibilidad se demuestra sobre generadores de caminos y no se introduce por una etiqueta de pertenencia a un complejo.

La evaluación terminal tampoco requiere elegir una historia particular. Para cada continuación \(h\), se aplica la composición efectiva de sus actualizaciones. El resultado conserva el par formado por \(h\) y su lectura, con su multiplicidad. Si \(k<\ell<N\), las continuaciones se descomponen como una unión disjunta de prefijos hasta \(\ell\) y sus continuaciones hasta \(N\). La proposición~\ref{iv:cont:concatenacion} demuestra que evaluar directamente o evaluar esas dos partes produce la misma lectura. Así se obtiene el cuadrado de naturalidad del terminal.

Los módulos de caminos, las incidencias, el transporte de memoria y sus diferenciales proceden de las mismas prolongaciones. Su ensamblaje no consiste en añadir a un estado cinco coordenadas que se conservarían por definición: sus relaciones se construyen mediante sumas, bordes y composiciones efectivas, y se verifican sobre sus generadores.

\subsection{Determinantes, límite y lectura aritmética}
\label{iv:cont:determinantes}
En cada residencia, el complejo local libre \(\Gamma\) admite su línea determinante graduada

\[
\operatorname{Det}\Gamma=
\bigotimes_r\left(\bigwedge^{\mathrm{rk}\Gamma_r}\Gamma_r\right)^{(-1)^r}.
\]
El ensamblaje total del apartado anterior tiene, por su parte, los módulos \(\operatorname{Tot}_d\mathcal B=\bigoplus_{q+r=d}\mathcal B_{q,r}\). A horizonte finito y sobre el mismo anillo, son libres de rango finito y sólo un número finito de ellos es no nulo. Su línea es
\[
\operatorname{Det}(\operatorname{Tot}\mathcal B)=
\bigotimes_{q,r}
\left(\bigwedge^{\operatorname{rk}\mathcal B_{q,r}}
\mathcal B_{q,r}\right)^{(-1)^{q+r}}.
\]
Esta fórmula resulta de la descomposición por suma directa en cada grado. Distingue la línea del complejo local y la del ensamblaje con todos los caminos; sus bases se ordenan mediante las etiquetas conservadas. La reducción modular conmuta con estas potencias exteriores porque los módulos y sus bases proceden de los mismos módulos libres enteros a horizonte fijado.
La línea está definida por los módulos y sus grados, tanto sobre \(\mathbb Z\) como después de reducir a \(A_N\); el exponente negativo designa el módulo dual. Una torsión numérica no nula exige, adicionalmente, especificar el complejo sobre un cuerpo, las bases y las identificaciones de homología pertinentes. Para calcular factores de Smith se usa la matriz entera producida antes de la reducción, no un levantamiento arbitrario de una matriz modular. Un bloque cuadrado entero de determinante no nulo es invertible sobre \(\mathbb Q\); su reducción es invertible sobre \(A_N\) solamente cuando ese determinante es una unidad, es decir, cuando no es divisible por \(3\). El refinamiento de memorias e incidencias permite formular y comprobar la estabilidad de los factores invariantes pertinentes; estas condiciones no se deducen de la existencia de la línea determinante.

El paso al límite conserva los proyectivos compatibles de residuos, las historias y sus multiplicidades. Los cuadrados probados en los apartados anteriores permanecen conmutativos porque sus aplicaciones están dadas por las mismas fórmulas en cada nivel. Para una coordenada arquimediana, el paso adicional consiste en publicar intervalos cilíndricos no vacíos, compatibles por inclusión y de diámetro tendente a cero. Sus clausuras compactas anidadas tienen una intersección de un solo punto: la existencia se obtiene por compacidad y la unicidad por la cota del diámetro. Ése es el valor de la lectura límite. Si los intervalos son semiabiertos, el punto puede pertenecer sólo a sus clausuras; la regla de representación de frontera determina entonces su escritura posicional. El valor no sustituye la historia que lo ha producido. La otra publicación conserva las incidencias del mismo estado. Los dos destinos comparten su origen sin recibir por ello la misma información.

En los capítulos siguientes, las regiones coinductivas, el registro dodecafásico y la incidencia excepcional se publican desde ese sustrato con sus respectivos mapas. La realización hilbertiana conserva después las etiquetas de fase y especifica su propia inclusión unital; la dualidad y las pantallas actúan sobre los codominios que allí se construyen. Las operaciones conjuntas aquí reunidas transmiten historias, hojas, orientación y memoria a esas publicaciones. La identificación de un funcional concreto con una forma analítica externa requiere, a su vez, escribir su aplicación sobre los generadores y verificar sus productos escalares: las identidades de conservación del presente capítulo no reemplazan esa composición.

\subsection{Procedencia, comprobación y alcance}
\label{iv:cont:mapa-correlativo}
% RESULTADO_RECUPERADO: 04b2_operaciones_intrinsecas_correlativas.tex,
% eq:pdfv2-cor-accion-sp/sol/gau/coh/det y eq:pdfv2-cor-memorias.
% Las marcas internas sp/sol/gau/coh/det corresponden, en ese orden,
% a las cinco filas siguientes; no son cinco dominios independientes.
Las cinco operaciones se reconocen por su acción sobre las mismas
historias y por las relaciones demostradas en los apartados anteriores.
El transporte afín \(U_\alpha\) de la memoria y el transporte
lineal \(\mathcal U_k\) de las historias tienen los dominios
respectivos de las secciones~\ref{iv:cont:historias}
y~\ref{iv:cont:transporte}. La graduación de hojas es
\(\sigma_k=P_k-Q_k\). Para una historia
\(h_j=(\varepsilon_1,\ldots,\varepsilon_j)\), la acumulación del
balance queda explícita mediante
\[
 M_0^\pm=0,\qquad
 M_k^\pm(h_k)=\sum_{j=1}^k b^\pm(h_j),\qquad
 M_{k+1}^\pm(h_k\varepsilon)
 =M_k^\pm(h_k)+b^\pm(h_k\varepsilon).
\]
Cada sumando es la lectura firmada del prefijo correspondiente,
definida en la sección~\ref{iv:cont:balance}. Separar el último
sumando demuestra la actualización; las identidades
\(M_k^+-M_k^-=\sum_{j=1}^kb(h_j)\) y
\(M_k^++M_k^-=\sum_{j=1}^k|b(h_j)|\) se deducen de las
definiciones de las partes positiva y negativa. La acumulación y la
cancelación condicional son, así, operaciones distintas sobre las
mismas contribuciones.

\begin{center}
\small
\begin{tabular}{@{}p{.21\linewidth}p{.48\linewidth}p{.23\linewidth}@{}}
\toprule
Operación conjunta & Objetos y acción efectiva & Desarrollo y prueba\\
\midrule
Transporte y graduación de hojas
& \(\mathcal U_k:\ell^2(\widehat X_k)\to
\ell^2(\widehat X_{k+1})\), con \(\widehat X_k\) formado por
las historias levantadas; \(\sigma_k=P_k-Q_k\), \(A_k\) y
\(\eta_k\) conservan o intercambian la hoja actual.
& Sección~\ref{iv:cont:transporte}; proposición~\ref{iv:cont:gram}
y teorema~\ref{iv:cont:terminal}.\\[5pt]
Balance, cancelación y memoria
& \(b:X_k\to\ZZ\), \(b^\pm:X_k\to\mathbb N_0\);
acumulación \(M_k^\pm\) sobre prefijos y cancelación
\(\kappa_{\mathrm{can}}=(E|b_f|-|Eb_f|)/2\) al publicar
la media condicional.
& Sección~\ref{iv:cont:balance};
proposición~\ref{iv:cont:cancelacion}.\\[5pt]
Incidencia ternaria de las hojas
& \(B:\FF_3^4\to\FF_3^6\),
\((Ba)_{\{L,L'\}}=a_L+a_{L'}\); se conserva el ambiente
completo de \(q\), con \(sB=0\) y \(W(q+Ba)=W(q)\).
& Primera parte de la
sección~\ref{iv:cont:incidencia-rectangular}.\\[5pt]
Puertos y complejo rectangular
& \(\kappa:A_N^{I_k}\oplus A_N^{J_k}\to
A_N^{I_k\times J_k}\), \(\kappa(u,v)_{ij}=u_i-v_j\);
la diferencia \(\delta\) y la precomposición \(L_k\)
conservan las etiquetas y sus multiplicidades.
& Proposiciones~\ref{iv:cont:rectangular-exacta}
y~\ref{iv:rectangular-natural}.\\[5pt]
Memoria celular, relleno y línea determinante
& \(\Psi_\alpha:\Gamma_y\to\Gamma_{y'}\),
\(K_\alpha\in\Gamma_{y',2}\),
\(\partial h_*=z^--z^+\); las líneas
\(\operatorname{Det}\Gamma\) y
\(\operatorname{Det}(\operatorname{Tot}\mathcal B)\)
se forman sobre los complejos basados correspondientes.
& Secciones~\ref{iv:cont:memoria}--\ref{iv:cont:determinantes};
proposiciones~\ref{iv:cont:memoria-natural}
y~\ref{iv:cont:memoria-homologia}.\\
\bottomrule
\end{tabular}
\end{center}

El cuadro identifica cinco acciones de una misma construcción: cada
prefijo conserva sus hojas, orientación, residuos, cocientes, memoria
y prolongaciones terminales. La precomposición de puertos lleva
coordenadas de un nivel al refinado; la reducción de coeficientes
lleva \(A_{N+1}\) a \(A_N\). Sus sentidos y dominios son
distintos y sus cuadrados conmutativos están explicitados en la
proposición~\ref{iv:rectangular-natural}. El ensamblaje de la
sección~\ref{iv:cont:ensamblaje} utiliza simultáneamente las
continuaciones terminales y el transporte del complejo de memoria;
sus fórmulas de composición reúnen las acciones del cuadro sobre el
mismo sistema de historias. La lectura de una línea determinante
como escalar conserva las condiciones de bases, homología y
no singularidad de la sección~\ref{iv:cont:determinantes}.

Este desarrollo reúne las construcciones de estado, árbol y medida común; las operaciones intrínsecas correlativas; y el ensamblaje terminal con naturalidad y límite demostrados en esta sección. Se incorporan sus identidades de balance, incidencia y memoria, con demostraciones dentro de este capítulo. Las cinco construcciones del continuo se mantienen como operaciones sobre un mismo árbol. La exposición amplía sus pruebas locales y explicita las condiciones de cada especialización: transporte isométrico, pesos uniformes, extinción del terminal y no singularidad de un bloque determinantal.

El comprobador adjunto contrasta instancias exactas de los cociclos, la cancelación ponderada, la secuencia rectangular y el complejo de memoria. Esos controles acompañan las demostraciones generales del texto y verifican los operadores aquí definidos en sus dominios respectivos. El paso posterior hacia las representaciones excepcionales y la realización hilbertiana conserva sus aplicaciones sobre generadores: la sección~\ref{iv:hilbert} especifica la fibra de fase, los operadores de Weyl y las inclusiones que determinan sus clausuras. La conservación documental y los controles finitos mantienen su función de trazabilidad y contraste de esas pruebas.


% SOURCE foundation/sections/generacion.tex
\section{Generación coinductiva de \texorpdfstring{\(\pi,\varphi,e\)}{pi, phi, e} a profundidad arbitraria}
\label{sec:generacion}

La emisión de seis bloques es una etapa finita de la construcción, no una expansión decimal completa. Su función consiste en producir regiones con multiplicidad, orientación y reglas de transporte. La prolongación debe conservar esos datos y determinar una familia compatible de prefijos de longitud arbitraria. La palabra \emph{generación} designa esta operación; \emph{genealogía} designa la reconstrucción de sus dependencias. Ambas nociones intervienen, pero cumplen funciones distintas.

\subsection{Órbitas, regiones y caracteres de clausura, propagación y autoescala}

Sobre una palabra \(w=(w_1,\ldots,w_6)\), definimos
\[
 r(w)=(w_3,w_1,w_2,w_5,w_6,w_4),\qquad
 s(w)=(w_4,w_5,w_6,w_1,w_2,w_3).
\]
Se cumple \(r^3=s^2=I\) y \(srs=r^{-1}\); por tanto estas permutaciones realizan el grupo diedral \(D_3\). La acción conserva el catálogo y las multiplicidades y conmuta con \(\Psi\). El cociente de las \(243\) palabras visibles tiene \(43\) órbitas: treinta y ocho de cardinal seis y cinco de cardinal tres.

La clase de clausura se reconoce por un predicado de fibra: sus seis orientaciones tienen cinco emisiones por palabra, multiplicidad total \(1008\) y distribución \(432+4\cdot144\). Existe una sola órbita con esas propiedades en el catálogo. Para precisar las otras dos selecciones, escribimos \(f(w)=\#\Psi^{-1}(w)\), \(m(w)=\sum_{U\in\Psi^{-1}(w)}\operatorname{mult}(U)\) y \(\nu(w)=(n_0,n_1,n_2)\). La clase diagonal aditiva de una emisión es \(d(U)=[i+j]_9\) en los índices de cero a ocho; su firma aditiva conserva esta clase y la orientación \(ES\) o \(NO\). Sea \(\mathcal D(\mathcal O)\) el conjunto de esas clases para todas las emisiones sobre una órbita.

Entre las órbitas de tamaño seis, el predicado conjunto
\[
 f(w)=1,\qquad m(w)=144\quad(w\in\mathcal O),\qquad
 \mathcal D(\mathcal O)=\{0,3,6\}
\]
deja exactamente seis órbitas. En este sector, el censo \(\nu=(2,3,1)\), igual al de clausura, selecciona únicamente la propagación; el censo \(\nu=(2,2,2)\) selecciona únicamente la autoescala. El soporte diagonal no puede omitirse: sin él hay cuatro órbitas unipuntuales equilibradas de multiplicidad \(144\), no una. Finalmente, la existencia de una emisión con \(d=6\) y orientación \(ES\) selecciona un representante de cada órbita y conduce a
\begin{equation}
 w^{\rm cl}=010211,\qquad w^{\rm pr}=201101,\qquad w^{\rm au}=121200.
 \label{eq:palabras-regionales}
\end{equation}
Estos símbolos designan primero regiones del catálogo. La identificación escalar se demostrará mediante sus lectores; no se eligen porque coincidan con cifras de una constante conocida.

La palabra de clausura \(010211\) tiene las cinco preimágenes siguientes:
\begin{center}\small
\begin{tabular}{@{}lll r@{}}\toprule
Emisión de seis bloques&Firma multiplicativa&Papel&Multiplicidad\\\midrule
\(501,614,498,169,272,272\)&\(555555\)&transversal&432\\
\(810,923,870,169,272,272\)&\(339555\)&orientación positiva&144\\
\(810,983,810,169,272,272\)&\(393555\)&orientación positiva&144\\
\(870,923,810,169,272,272\)&\(933555\)&orientación positiva&144\\
\(870,923,810,418,674,521\)&\(933717\)&retorno orientado&144\\\bottomrule
\end{tabular}
\end{center}
La multiplicidad \(432\) y la firma constante \(555555\) caracterizan la región transversal. No identifican la posición central \((5,5)\) de APP. Confundir una firma de seis ventanas con una posición del soporte eliminaría la distinción entre región numérica y estructura electrónica.

La microfibra posee además una morfología bipartita precisa. Cada emisión es un producto de una clase de cursor aditivo y una clase multiplicativa. Las tres regiones positivas comparten la misma clase aditiva de dieciocho condiciones; juntas contienen tres clases multiplicativas de ocho condiciones cada una. La región transversal posee un producto \(18\times24\); las tres positivas reunidas, otro \(18\times24\); y el retorno, un producto \(18\times8\). La descomposición por emisiones tiene cinco partes, mientras la descomposición por componentes conexas del grafo bipartito tiene tres. Ambas conservan las mismas \(1008\) semillas sin reducir la forma a un cardinal.

\subsection{Transporte finito y frontera de prolongación}

Los endomorfismos \(L_0,L_1\) del espacio de palabras ternarias actúan antes de la realización arquimediana. Su determinación comienza por las transiciones del estado aritmético, no por la prescripción de dos matrices. Si \(U_t\) es el paso compuesto de selección, transporte y actualización, el bloque observado satisface
\[
 b_j^\chi=\Pi_6(x_j^\chi),\qquad
 x_{j+1}^\chi=U_{9j+9}\cdots U_{9j+1}(x_j^\chi),
 \qquad \chi\in\{\mathrm{cl},\mathrm{pr},\mathrm{au}\}.
\]
Se forman seis transiciones por cada régimen: dos consecutivas en cada una de las tres regiones. Los pares de matrices de orígenes y destinos obtenidos en el registro son
\[
\begin{array}{c|c|c|c}
X_0&Y_0&X_1&Y_1\\\hline
010211&012222&010211&002111\\
012222&010211&002111&110221\\
201101&121221&102011&012222\\
121221&102011&012222&102011\\
121200&112202&121020&010210\\
112202&121020&010210&010200
\end{array}
\]
Cada cadena de seis dígitos representa una fila en \(\FF_3^6\). Puesto que \(\det X_0=1\) y \(\det X_1=-1\) en \(\FF_3\), las ecuaciones \(X_aL_a=Y_a\) determinan unívocamente \(L_a=X_a^{-1}Y_a\). En esa carta de vectores fila resultan
\[
 L_0=\begin{pmatrix}
 2&2&2&1&2&1\\2&1&2&2&1&1\\1&0&0&1&0&2\\
 0&0&1&1&1&0\\2&2&2&0&2&1\\2&1&2&1&0&0
 \end{pmatrix},\quad
 L_1=\begin{pmatrix}
 2&1&2&0&2&2\\1&2&1&1&2&0\\1&2&1&1&1&2\\
 0&1&0&0&2&1\\2&0&2&1&0&1\\0&2&2&2&1&1
 \end{pmatrix}\quad(\bmod3).
\]
El calendario \((L_0,L_0,L_1,L_1)\) produce cuatro nuevos bloques de seis componentes. Sus concatenaciones forman
\[
 w_6\longrightarrow w_{12}\longrightarrow w_{18}
 \longrightarrow w_{24}\longrightarrow w_{30}.
\]
Los subíndices de \(w\) indican longitud acumulada. La frontera \(R_{36}\) conserva la estructura terminal y abre la relación de prolongación; no se renombra como un último bloque periódico que pudiera repetirse indefinidamente.

La identidad \(L=X^{-1}Y\) demuestra unicidad una vez fijadas las transiciones; no sustituye la producción de esas transiciones por \(U_t\). La procedencia de su registro y el examen de los dieciséis calendarios binarios se documentan en el suplemento técnico. El calendario indicado es el único de esos dieciséis que reproduce conjuntamente las tres sucesiones de cinco bloques.

\subsection{Cinco regiones de clausura y compensación angular}

La región de clausura contiene una componente transversal y cuatro acompañantes. El lector simétrico actúa en el espacio de sus cinco posiciones mediante
\[
 P_5=\frac15\mathbf1\mathbf1^{\mathsf T},\qquad
 \mathbf1=(1,1,1,1,1)^{\mathsf T}.
\]
La invariancia \(P_5\lambda=\lambda\) y la conservación \(\mathbf1^{\mathsf T}\lambda=1\) fuerzan \(\lambda=\mathbf1/5\). Esta normalización distribuye la unidad entre posiciones del lector; no sustituye las multiplicidades de semillas \(432,144,144,144,144\) por frecuencias uniformes. La carta elíptica del lector realiza cada coordenada mediante \(1+\lambda_jJ\); su pendiente es \(\lambda_j\), por lo que la pendiente primitiva resulta \(1/5\). Se hace explícita así la regla que pasa de una coordenada normalizada a una pendiente, en lugar de identificar ambas nociones sin mapa. El lector compone las cuatro posiciones acompañantes por la ley
\begin{equation}
 a\oplus b=\frac{a+b}{1-ab}.
 \label{eq:suma-pendientes}
\end{equation}
Esta ley se obtiene al multiplicar \((1+aJ)(1+bJ)\) en el régimen \(J^2=-I\) y dividir por su componente escalar. Por ello la composición de pendientes es una realización de la estructura cuadrática de TRIT, no una suma ordinaria de cuatro fracciones.

\begin{proposition}[Compensador de la clausura]
La composición de cuatro pendientes \(1/5\) es \(120/119\). La condición de retorno a pendiente uno determina un único compensador positivo de la forma \(1/q\), con \(q>1\), y da \(q=239\).
\end{proposition}
\begin{proof}
La primera duplicación da \((1/5)\oplus(1/5)=5/12\); la segunda, \((5/12)\oplus(5/12)=120/119\). La condición
\[
 \frac{120/119-1/q}{1+(120/119)/q}=1
\]
equivale a \((120-119)q=120+119\), por lo que \(q=239\).
\end{proof}

El carácter de clausura resultante tiene realización
\begin{equation}
 \Pi_{\rm cl}=4\left(4\arctan\frac15-\arctan\frac1{239}\right).
 \label{eq:lector-pi}
\end{equation}
La identidad angular de Machin reconoce \(\Pi_{\rm cl}=\pi\). El sentido de la construcción es preciso: la pendiente y el número de composiciones están declarados por el lector de la pentafibra; la ecuación de compensación fuerza \(239\). Una vez producido este carácter, su evaluación por series alternadas no selecciona de nuevo la órbita de semillas.

Para \(0<x<1\), las sumas alternadas
\[
 A_n(x)=\sum_{j=0}^n\frac{(-1)^jx^{2j+1}}{2j+1}
\]
encierran \(\arctan x\) entre dos truncamientos consecutivos, con diferencia \(x^{2n+3}/(2n+3)\). Aplicadas a las dos pendientes de \eqref{eq:lector-pi}, proporcionan extremos racionales de error explícito. El valor numérico se calcula así como consecuencia del carácter exacto, sin introducir un prefijo objetivo.

\subsection{Lector de propagación y normalización de la unidad}

La región de propagación se realiza mediante una familia formal \(P_t\in\mathcal A[[t]]\), donde \(\mathcal A\) es un álgebra conmutativa unitaria sobre \(\QQ\). Su composición satisface \(P_{t+s}=P_tP_s\), con unidad \(P_0=1\), y el transporte elemental orientado fija el generador escalar \(P'_0=+1\). La condición fija su signo, no sólo su magnitud. Escribiendo
\[
 P_t=\sum_{n\ge0}a_nt^n,
\]
la normalización determina \(a_0=a_1=1\). La comparación del coeficiente de \(s^nt\) en \(P_{s+t}=P_sP_t\) da \((n+1)a_{n+1}=a_na_1\); equivalentemente, \(P'_t=P_t\). Así,
\begin{equation}
 (n+1)a_{n+1}=a_n,\qquad a_n=\frac1{n!}.
 \label{eq:propagacion}
\end{equation}
El carácter en una unidad de parámetro es \(\Pi_{\rm pr}=P_1\), reconocido posteriormente como \(e\).

La normalización del generador importa. La ley semigrupal sin ese dato admitiría \(P_t=e^{ct}\) para distintos \(c\); por ello no se omite la condición que fija la unidad del parámetro en el lector. Esta condición es anterior a la comparación decimal y pertenece a la exposición de la regla operacional, no a la elección de un valor experimental.

\begin{proposition}[Cotas racionales de propagación]
Para \(n\ge1\), sea \(E_n=\sum_{j=0}^n1/j!\). Entonces
\[
 E_n<\Pi_{\rm pr}<E_n+\frac{n+2}{(n+1)(n+1)!}.
\]
\end{proposition}
\begin{proof}
El primer término omitido es \(1/(n+1)!\). Cada cociente posterior es a lo sumo \(1/(n+2)\). La suma geométrica mayorante de la cola es
\[
 \frac1{(n+1)!}\frac1{1-1/(n+2)}
 =\frac{n+2}{(n+1)(n+1)!}.
\]
La desigualdad estricta se sigue de la disminución de los cocientes sucesivos.
\end{proof}

\subsection{Incidencia modal y generación de la autoescala}

El selector trítico produce el ciclo \((+1,-1,0)\). En los modos \(\Sigma,\Pi\), la regla es \(+1:m\mapsto\Sigma\), \(-1:m\mapsto\Pi\), \(0:m\mapsto m\). Con cierre cíclico, el modo observado es \((\Sigma,\Pi,\Pi)\). Sus aristas son \(\Sigma\to\Pi\), \(\Pi\to\Pi\) y \(\Pi\to\Sigma\). Por tanto su incidencia, en ese orden de modos, es
\begin{equation}
 F_{\rm av}=\begin{pmatrix}0&1\\1&1\end{pmatrix}.
 \label{eq:fav}
\end{equation}
Las cuatro entradas se obtienen del calendario, no de una ecuación que haya recibido previamente el valor de \(\varphi\). La repetición del calendario da conteos \(3F_{\rm av},9F_{\rm av},36F_{\rm av}\) en longitudes nueve, veintisiete y ciento ocho. Estos conteos no son las potencias de la matriz.

\begin{theorem}[Carácter positivo de autoescala]
La incidencia \eqref{eq:fav} posee un único rayo propio estrictamente positivo. Al normalizarlo como \((1,x)^{\mathsf T}\), su coordenada \(x\) satisface \(x^2-x-1=0\), \(x>0\), y es \(\varphi=(1+\sqrt5)/2\).
\end{theorem}
\begin{proof}
La ecuación propia da \(x=\lambda\) y \(1+x=\lambda x\). Eliminar \(\lambda\) produce el polinomio. De sus dos raíces sólo una es positiva. Como \(F_{\rm av}^2\) tiene todas sus entradas positivas, el rayo positivo es único y simple. La expresión radical reconoce la coordenada obtenida; no participa en la producción de \(F_{\rm av}\).
\end{proof}

Con \(F_0=0,F_1=1,F_{n+1}=F_n+F_{n-1}\), la incidencia induce, para \(n\geq1\),
\[
 F_{\rm av}^n=\begin{pmatrix}F_{n-1}&F_n\\F_n&F_{n+1}\end{pmatrix}.
\]
Los cocientes consecutivos \(F_{n+1}/F_n\) alternan alrededor de \(\varphi\). La identidad de determinante
\[
 F_{n+1}^2-F_nF_{n+2}=(-1)^n
\]
da una anchura exacta \(1/(F_nF_{n+1})\) entre dos cocientes consecutivos. Su crecimiento hace arbitrariamente pequeña esta horquilla racional.

La medida sobre todas las historias compatibles no coincide con la frecuencia de modos del calendario de tres eventos. La frecuencia cronológica es \((1/3,2/3)\). La medida de máxima entropía del envolvente simbólico definido por \(F_{\rm av}\) se construye con sus vectores propios y tiene razón de pesos \(\varphi^{-2}\). Esta última interviene en la lectura areal--volumétrica; no se obtiene sustituyendo sin más una frecuencia de conteo por otra.

\input{sections/retorno_areal_volumetrico.tex}

Después de generados los caracteres, las dos orientaciones del lector
\[
 R(x,y)=\frac{x}{y^2}
\]
proporcionan
\begin{equation}
 \rho_A=\frac\pi{\varphi^2},\qquad
 \rho_E=\frac\varphi{\pi^2},\qquad
 \varphi^3\rho_A^2\rho_E=1.
 \label{eq:areal-electron}
\end{equation}
La segunda expresión interviene en el exponente de la construcción electrónica desarrollada en \cite{hmtX}, en la sección sobre centro electrónico y registros. La primera se obtiene también como límite de la lectura marcada \(\pi F_{n-1}/F_{n+1}\). El intercambio de argumentos relaciona ambas, pero no convierte una reparametrización de la ecuación electrónica en una segunda derivación independiente de todos sus coeficientes.

\input{sections/euler_trit_articulacion.tex}

\subsection{Cilindros compatibles y profundidad arbitraria}

La realización de un bloque ternario \(b=(b_1,\ldots,b_6)\) es el entero
\[
 \nu(b)=\sum_{j=1}^6b_j3^{6-j}\in\{0,\ldots,728\}.
\]
Una sucesión de bloques define prefijos enteros por \(N_{n+1}=729N_n+\nu(b_{n+1})\), comenzando en \(N_0=m\), donde \(m\) es la parte entera determinada por el lector. Las horquillas anteriores sitúan los caracteres de clausura, propagación y autoescala en \((3,4)\), \((2,3)\) y \((1,2)\), respectivamente; por tanto, sus valores iniciales son \(m=3,2,1\). Los bloques siguientes refinan la parte fraccionaria. El cilindro correspondiente es
\begin{equation}
 I_n=\left[\frac{N_n}{729^n},\frac{N_n+1}{729^n}\right],
 \qquad |I_n|=729^{-n}.
 \label{eq:cilindro}
\end{equation}
La convención semiabierta se utiliza para seleccionar un representante cuando un valor cae en una frontera. Para el paso al límite se utilizan las clausuras de estos intervalos y se identifican las dos expansiones fronterizas equivalentes. Esta precisión impide suponer que toda intersección de intervalos semiabiertos encajados sea automáticamente no vacía.

\begin{theorem}[Determinación de la realización escalar]
Una historia de bloques compatible determina un único límite de los extremos izquierdos de \eqref{eq:cilindro}. Si un carácter exacto dispone de horquillas racionales de anchura arbitrariamente pequeña y se separa de la frontera de cada cilindro decimal examinado, cada prefijo decimal de longitud finita se determina después de un cálculo finito. Los prefijos así obtenidos son compatibles entre sí.
\end{theorem}
\begin{proof}
La extensión de un bloque conserva el prefijo anterior por división euclídea. Los cilindros cerrados están encajados y su diámetro tiende a cero; sus extremos izquierdos son de Cauchy y tienen un único límite. Para una longitud decimal fijada, un valor separado de las fronteras de la partición posee distancia positiva a ellas. Una horquilla suficientemente estrecha queda contenida en una única celda y determina su prefijo. La unicidad de la celda y el encajamiento de las particiones prueban la compatibilidad por truncamiento. En los puntos con expansión terminal debe añadirse la decisión exacta de frontera y su convención de representación; no se deduce esa decisión de una anchura numérica solamente.
\end{proof}

\begin{corollary}[Obtención finita y compatibilidad de los prefijos regionales]
\label{gen:prefijos-regionales}
Para cada una de las realizaciones \(\pi,\varphi,e\) y para toda longitud
decimal finita, las horquillas construidas determinan el prefijo correspondiente
mediante un cálculo finito. Los prefijos obtenidos a distintas profundidades
son compatibles por truncamiento.
\end{corollary}
\begin{proof}
Los caracteres de clausura, propagación y autoescala poseen las horquillas
anteriores. Tras su reconocimiento, la irracionalidad de \(\pi,e,\varphi\)
garantiza su separación de las fronteras racionales de cada partición decimal
finita. Se aplica el teorema precedente a cada carácter.
\end{proof}
La profundidad arbitraria expresa el cuantificador del corolario: para todo
horizonte finito existe un cálculo que determina su prefijo. El objeto
matemático es la familia compatible completa; cada ejecución obtiene una
proyección finita de ella.

\subsection{Involución especular y sincronización de las regiones}

La conexión nonádica transporta conjuntamente los tres caracteres. En la carta de la novena fase, el eje de autoescala y la suma ternaria de clausura y propagación permanecen fijados bajo el intercambio especular. La inversión helicoidal cambia el sentido del transporte; el espejo de los dos canales cambia su orientación relativa. Son involuciones diferentes, aunque ambas actúen sobre el mismo sistema de trayectorias.

Las coincidencias de censo ilustran esa coordinación:
\begin{center}
\begin{tabular}{@{}c rrr l@{}}\toprule
Fase&\(N_\pi\)&\(N_e\)&\(N_\varphi\)&Igualdad de conteos\\\midrule
4&207&207&122&\(N_\pi=N_e\)\\
5&39&151&151&\(N_e=N_\varphi\)\\
6&110&110&69&\(N_\pi=N_e\)\\
7&81&29&81&\(N_\varphi=N_\pi\)\\
8&59&59&59&sincronización triple\\
9&43&43&43&sincronización triple\\\bottomrule
\end{tabular}
\end{center}
Estos enteros cuentan coordenadas compatibles de las regiones. No afirman una igualdad entre \(\pi\), \(e\) y \(\varphi\) como números reales. Si
\[
 \partial(a,b,c)=(a-b,b-c,c-a),
\]
su imagen pertenece al retículo \(A_2=\{(u,v,w)\in\ZZ^3:u+v+w=0\}\). En las fases ocho y nueve la componente relativa se anula, pero la componente diagonal cambia de \((59,59,59)\) a \((43,43,43)\). La sincronización relativa coexiste con un desplazamiento de \(-16\) por canal. El cierre que conducirá a \(\alpha\) debe conservar ambas informaciones: coincidencia relativa y evolución de la memoria.

\subsection{Suspensión helicoidal esturmiana}

La relación entre las bases nueve y diez de la completación define el parámetro adimensional
\begin{equation}
 \delta=\frac{\log(10/9)}{\log10}.
 \label{eq:delta-sturm}
\end{equation}
No se obtiene de una constante física objetivo. Es la escala logarítmica relativa de dos bases ya presentes en la construcción. Su empleo como lectura simbólica del refinamiento tiene una consecuencia exacta.

\begin{lemma}[Irracionalidad de la pendiente]
El número \(\delta\) es irracional.
\end{lemma}
\begin{proof}
Si \(\delta=p/q\), \(q>0\), entonces \((10/9)^q=10^p\) y \(9^q=10^{q-p}\). La factorización del primer miembro contiene el primo tres y la del segundo sólo los primos dos y cinco. La igualdad es imposible.
\end{proof}

La palabra mecánica inferior, con intercepto \(\theta\), es
\[
 z_n(\theta)=\lfloor(n+1)\delta+\theta\rfloor-\lfloor n\delta+\theta\rfloor\in\{0,1\}.
\]
Se obtiene al observar la rotación \(\theta\mapsto\theta+\delta\pmod1\) mediante el intervalo \([1-\delta,1)\). La sucesión de eventos tiene densidad \(\delta\), y la suma sobre cualquier ventana de longitud \(L\) vale \(\lfloor L\delta\rfloor\) o \(\lceil L\delta\rceil\). Como \(21<1/\delta<22\), las separaciones entre eventos son veintiuno o veintidós. La palabra que indica cuáles son los intervalos cortos es otra palabra mecánica de pendiente \(\sigma=22-1/\delta\); sus eventos están separados por seis o siete intervalos primarios. La segunda escala es inducida por la primera y no un segundo parámetro libre.

\input{sections/vacancias_capacidad.tex}

\begin{theorem}[Complejidad de las lecturas de fase]
Para \(m\ge1\), la palabra mecánica tiene complejidad factorial \(p(m)=m+1\). Al conservar la fase \(n\bmod9\), la complejidad es \(p_9(m)=9(m+1)\); al conservar sólo su clase ternaria, es \(p_3(m)=3(m+1)\). Las tres entropías topológicas son nulas.
\end{theorem}
\begin{proof}
Los predecesores de los dos extremos del intervalo de codificación forman, para factores de longitud \(m\), \(m+1\) puntos distintos de la circunferencia: la coincidencia entre extremos desplazados elimina las duplicaciones y la irracionalidad impide coincidencias adicionales. Esos puntos determinan \(m+1\) intervalos de itinerario constante y distintos. Toda órbita irracional visita cada intervalo, de donde \(p(m)=m+1\).

Fijada una fase módulo nueve, las posiciones de esa fase recorren una rotación de paso \(9\delta\), también irracional. Por densidad aparecen todos los factores en cada fase, y la primera coordenada distingue sus nueve copias. El mismo argumento con paso \(3\delta\) da tres copias tras la reducción ternaria. Finalmente, \(m^{-1}\log(c(m+1))\to0\) para \(c=1,3,9\).
\end{proof}

El producto de fases \(C_9\times\mathbb T\), trasladado por \((1,\delta)\), tiene caracteres de frecuencia
\[
 \frac a9+k\delta\pmod1,\qquad a\in\ZZ/9\ZZ,\ k\in\ZZ.
\]
El módulo espectral no es \(\log10\): es el módulo aditivo generado por \(1/9\) y \(\delta\), módulo los enteros. El factor \(\log10\) pertenece a la normalización de \eqref{eq:delta-sturm}.

Una presentación operatoria hace explícita la memoria. En \(\ell^2(\ZZ)\), sea \(T|n\rangle=|n+1\rangle\), \(V_9|n\rangle=\zeta_9^n|n\rangle\), \(W_\delta|n\rangle=e^{2\pi i n\delta}|n\rangle\) y \(D_M|n\rangle=\lfloor n/9\rfloor|n\rangle\). Sobre vectores de soporte finito,
\begin{equation}
 V_9T=\zeta_9TV_9,\qquad
 W_\delta T=e^{2\pi i\delta}TW_\delta,\qquad
 [D_M,T^9]=T^9.
 \label{eq:weyl-relojes}
\end{equation}
Las dos primeras relaciones son covariancias de Weyl; la tercera expresa que una vuelta nonádica aumenta el grado de memoria. La rotación de base es abeliana, mientras el álgebra de observables y traslaciones no lo es. La inversión \(|n\rangle\mapsto|-n\rangle\) conjuga el avance con el retroceso. Esta realización proporciona una suspensión helicoidal aperiódica con retorno finito de fase, no una palabra periódica de longitud nueve.

\input{figures/holonomia.tex}

Para este recuento fijamos el intercepto \(\theta=0\) y los extremos acumulados \((N_0,\ldots,N_4)=(0,729,972,999,1000)\), correspondientes al orden cúbico \((729,243,27,1)\). Los conteos inferiores son \(\lfloor N_j\delta\rfloor-\lfloor N_{j-1}\delta\rfloor\), y dan \((33,11,1,0)\). Los superiores son \(\lceil N_j\delta\rceil-\lceil N_{j-1}\delta\rceil\), con la misma frontera inicial cero, y dan \((34,11,1,0)\). Estos valores dependen del intercepto y del orden de los estratos; no se afirman para una fase inicial arbitraria. Sus totales son cuarenta y cinco y cuarenta y seis. El incremento corresponde a una única unidad de frontera, situada en el primer estrato de esta partición ordenada. Las cifras treinta y cuatro y once aparecen así en un mismo recuento previo a toda carta de unidades; identificarlas con exponentes físicos exige, además, la derivación dimensional correspondiente.

La denominación \emph{cristal aperiódico} se emplea aquí para este orden simbólico, con complejidad, espectro y retorno especificados. La entropía topológica nula no equivale por sí sola a disipación física nula: una realización termodinámica necesita una dinámica energética y una operación de borrado definidas. Del mismo modo, esta lectura esturmiana no sustituye a todos los operadores del TPK ni demuestra por sí sola la periodicidad o aperiodicidad de la salida completa de \(\alpha\). Su función es hacer explícita una estructura de orden que el simple retorno del calendario no describe.

\input{sections/revision_monodromia.tex}

\input{sections/primos_retornos.tex}


% SOURCE foundation/sections/retorno_areal_volumetrico.tex
\subsubsection{Espacio de trayectorias y medida invariante entre hojas}
\label{sec:medida-retorno-hojas}

La topología toroidal de APP fija el soporte de las traslaciones cardinales
y sus clases de enrollamiento. La correspondencia entre hojas procede de
otra operación sobre ese mismo soporte: la selección temporal de la
evaluación aditiva o multiplicativa. Su tipado geométrico es
\[
 \eta_{\rm av}(\Sigma)=\mathscr H_{\rm ar},\qquad
 \eta_{\rm av}(\Pi)=\mathscr H_{\rm vol}.
\]
La primera hoja corresponde a la lectura aditiva de frontera; la segunda,
a la lectura multiplicativa de interior. Este mapa transporta las etiquetas
de modo y conserva la incidencia \eqref{eq:fav}. El cuadrado de la razón
de autoescala se determina mediante las trayectorias de esa incidencia.

\Needspace{8\baselineskip}
\begin{definition}[Envolvente simbólico de memoria uno]
Sea
\[
 X_{\rm av}=\{(s_k)_{k\in\mathbb Z}\in
 \{\mathscr H_{\rm ar},\mathscr H_{\rm vol}\}^{\mathbb Z}:
 (F_{\rm av})_{s_k,s_{k+1}}=1\text{ para todo }k\}.
\]
El desplazamiento de índices actúa sobre este espacio. Es el menor sistema
simbólico de memoria uno que contiene las sucesiones modales y conserva
sus pares consecutivos: admite las dos transiciones cruzadas y la
autorretención volumétrica. La fase, la orientación cardinal y los registros
de la trayectoria permanecen en el estado TPK del que se obtiene este
envolvente.
\end{definition}

\begin{proposition}[Ponderación de las trayectorias entre hojas]
En el orden areal--volumétrico, la medida invariante de entropía máxima
de \(X_{\rm av}\) tiene matriz de transición y distribución estacionaria
\begin{equation}
 P_{\rm av}=\begin{pmatrix}0&1\\\varphi^{-2}&\varphi^{-1}\end{pmatrix},
 \qquad
 (\mu_{\rm ar},\mu_{\rm vol})
 =\frac{(1,\varphi^2)}{1+\varphi^2}.
 \label{eq:medida-hojas-exacta}
\end{equation}
En particular, \(\mu_{\rm ar}/\mu_{\rm vol}=\varphi^{-2}\).
\end{proposition}
\begin{proof}
La incidencia ya construida tiene vector propio positivo
\(r=(1,\varphi)^{\mathsf T}\) y autovalor \(\varphi\). La
normalización
\[
 (P_{\rm av})_{ij}
 =\frac{(F_{\rm av})_{ij}r_j}{\varphi r_i}
\]
da la matriz indicada. Sus filas suman uno por
\(\varphi^{-2}+\varphi^{-1}=1\). La simetría de \(F_{\rm av}\)
implica que los pesos proporcionales a \(r_i^2\) satisfacen balance
detallado; por tanto, \(\mu P_{\rm av}=\mu\). La irreducibilidad y la
autorretención volumétrica determinan una única distribución estacionaria.

Para verificar la maximalidad, en toda transición permitida se cumple
\(\log(P_{\rm av})_{ij}=\log r_j-\log r_i-\log\varphi\).
Bajo cualquier medida invariante, los términos de extremo se cancelan
en promedio. Si \(\nu_i\) son sus pesos de vértice y
\(q_i(j)=\nu(s_1=j\mid s_0=i)\), para \(\nu_i>0\), entonces
\[
 h_\nu\leq H_\nu(s_1\mid s_0)
 =\log\varphi-\sum_{i:\nu_i>0}\nu_iD(q_i\Vert P_i)
 \leq\log\varphi,
\]
donde \(D(q\Vert p)=\sum_jq_j\log(q_j/p_j)\) es la entropía
relativa en los pares permitidos y se adopta \(0\log0=0\).
La medida markoviana de \eqref{eq:medida-hojas-exacta} alcanza la cota.
La igualdad en la primera desigualdad exige la propiedad de Markov;
la igualdad en la segunda exige \(q_i=P_i\). La estacionariedad y
la irreducibilidad fijan entonces \(\mu\), lo que establece la
unicidad. Esta es la medida de Parry del envolvente definido.
\end{proof}

La frecuencia \((1/3,2/3)\) cuenta los instantes del calendario primitivo;
\(\mu\) pondera las trayectorias del envolvente. Su razón cuadrática
expresa el producto de las amplitudes de incidencia de entrada y salida.
La memoria y las restricciones de fase del TPK conservan su dominio
propio; la entropía del envolvente corresponde al sistema simbólico
especificado en esta definición.

\subsubsection{Operador de clausura y límite de los retornos marcados}
\label{sec:operador-retorno-marcado}

La coordenada de clausura \(\pi\), generada por el lector regional,
interviene después de la construcción de la incidencia modal. Sean
\(e_{\rm ar}=(1,0)^{\mathsf T}\),
\(e_{\rm vol}=(0,1)^{\mathsf T}\) y
\(P_i=e_ie_i^{\mathsf T}\) los proyectores de hoja. Las condiciones
del lector son conservación de ambas hojas, normalización volumétrica
unitaria y marca areal igual a la coordenada de clausura. Estas condiciones
determinan el operador positivo
\begin{equation}
 D_\pi=\pi P_{\rm ar}+P_{\rm vol}
       =\begin{pmatrix}\pi&0\\0&1\end{pmatrix}.
 \label{eq:operador-marca-pi}
\end{equation}
La diagonalidad expresa la conservación de hojas y las dos entradas
expresan sus normalizaciones; de este modo queda explícita la unicidad
del operador para el lector fijado.

\Needspace{8\baselineskip}
\begin{proposition}[Razón de retornos de frontera]
Para \(n\geq1\), la evaluación marcada
\begin{equation}
 \Theta_n=
 \frac{\langle e_{\rm ar},D_\pi F_{\rm av}^{\,n}e_{\rm ar}\rangle}
 {\langle e_{\rm vol},F_{\rm av}^{\,n}e_{\rm vol}\rangle}
 =\pi\frac{F_{n-1}}{F_{n+1}}
 \label{eq:retorno-marcado-finito}
\end{equation}
satisface
\begin{equation}
 \lim_{n\to\infty}\Theta_n
 =\pi\frac{\mu_{\rm ar}}{\mu_{\rm vol}}
 =\frac{\pi}{\varphi^2}.
 \label{eq:retorno-marcado-limite}
\end{equation}
\end{proposition}
\begin{proof}
Las entradas diagonales de \(F_{\rm av}^{\,n}\) cuentan los caminos
de longitud \(n\) que regresan a su hoja inicial. La fórmula de potencias
ya demostrada proporciona \(F_{n-1}\) y \(F_{n+1}\). El operador
\(D_\pi\) multiplica por \(\pi\) el primer recuento. Finalmente,
\(F_{n-1}/F_{n+1}\) es el producto de dos cocientes consecutivos
que convergen a \(\varphi^{-1}\); su límite es
\(\varphi^{-2}\). La identidad con la razón de pesos se sigue de
\eqref{eq:medida-hojas-exacta}.
\end{proof}

El operador de clausura aparece una sola vez, en la evaluación terminal
del retorno. La transferencia \(F_{\rm av}\) permanece determinada
por el calendario. Así se compone la información de dos lectores del
mismo estado: la incidencia modal determina la ponderación cuadrática
y la región de clausura aporta su coordenada a la frontera.

La composición conserva también el refinamiento. Para cilindros positivos
\(I=[a,b]\) y \(J=[c,d]\), la imagen del lector
\(\mathcal R(x,y)=x/y^2\) es el intervalo
\[
 \mathcal Q(I,J)=\left[\frac{a}{d^2},\frac{b}{c^2}\right].
\]
La monotonía en cada argumento prueba que los cilindros encajados de
clausura y autoescala producen intervalos encajados de la razón; la
positividad de sus límites y la contracción de sus diámetros determinan
\(\pi/\varphi^2\). La expresión escalar retiene así una cadena
de operaciones compatible con la profundidad, que se aplica a
continuación en sus dos orientaciones.


% SOURCE foundation/sections/euler_trit_articulacion.tex
% Articulación posterior a la generación de pi, phi y e.
% Contenido acumulativo; no reemplaza el desarrollo previo del TRIT.
\subsection{Realizaciones exponenciales de las coordenadas generadas}
\label{euler-trit-articulacion}

Las coordenadas regionales ya determinadas permiten reconocer las
realizaciones concretas de la familia cuadrática del TRIT. El orden es
sustantivo: la construcción de las regiones precede a la evaluación de una
exponencial en los parámetros \(\pi\) o \(2\log\varphi\). Estas expresiones
identifican la función de las coordenadas en operadores construidos desde
APP y TPK; no intervienen en la selección de los prefijos que las producen.

\subsubsection{Ciclo ternario, transporte nilpotente e incidencia de hojas}

La multiplicación \(a(r)=4r\) en \(\mathbb Z/9\mathbb Z\) tiene orden
\(3\). Sobre las unidades determina las órbitas \(\{1,4,7\}\) y
\(\{2,8,5\}\). En el módulo de aumento de cualquiera de ellas, formado por
las combinaciones enteras cuyos coeficientes suman cero, la base
\((e_r-e_{4r},e_{4r}-e_{7r})\) da
\[
 C=\begin{pmatrix}0&-1\\1&-1\end{pmatrix},
 \qquad C^2+C+I=0.
\]
La forma restringida tiene matriz de Gram
\(\bigl(\begin{smallmatrix}2&-1\\-1&2\end{smallmatrix}\bigr)\).
En las evaluaciones residuales módulo \(9\), la acción distingue las hojas:
\(S(ax,ay)=aS(x,y)\), mientras \(P(ax,ay)=a^{-1}P(x,y)\).
La realización cíclica conserva, por tanto, una orientación contragrediente
entre suma y producto. Los cocientes del levantamiento entero se conservan
en el registro aritmético correspondiente.

El transporte parabólico elemental se representa mediante
\[
 N=\begin{pmatrix}0&1\\0&0\end{pmatrix}.
\]
La incidencia areal--volumétrica ya construida en \eqref{eq:fav} proporciona
\[
 F_{\rm av}=\begin{pmatrix}0&1\\1&1\end{pmatrix},
 \qquad A=F_{\rm av}^2=\begin{pmatrix}1&1\\1&2\end{pmatrix}.
\]
Los tres operadores actúan en espacios reales declarados: el módulo de
aumento tensorizado con \(\mathbb R\), el plano del transporte nilpotente y
el plano de incidencia modal. Compartir dimensión matricial no identifica
estos espacios ni sus coordenadas.

\begin{proposition}[Generadores concretos de los tres tipos]
\label{euler-trit-generadores}
Los operadores
\[
 J_{\mathrm e}=\frac{2C+I}{\sqrt3},\qquad
 J_{\mathrm p}=N,\qquad
 J_{\mathrm h}=\frac{2A-3I}{\sqrt5}
\]
satisfacen, respectivamente,
\[
 J_{\mathrm e}^{\,2}=-I,\qquad
 J_{\mathrm p}^{\,2}=0,\qquad
 J_{\mathrm h}^{\,2}=I.
\]
\end{proposition}
\begin{proof}
La relación de \(C\) implica \((2C+I)^2=-3I\).
La multiplicación directa da \(N^2=0\).
Finalmente, \(A\) tiene traza \(3\) y determinante \(1\), por lo que
\(A^2-3A+I=0\) y \((2A-3I)^2=5I\).
\end{proof}

\begin{proposition}[Tres evaluaciones de la identidad exponencial]
\label{euler-trit-evaluaciones}
En las realizaciones anteriores,
\begin{equation}
 \boxed{
 \exp(\pi J_{\mathrm e})+I=0,\qquad
 \exp(J_{\mathrm p})=I+J_{\mathrm p},\qquad
 \exp(2\log\varphi\,J_{\mathrm h})=F_{\rm av}^{\,2}.}
 \label{euler-trit-tres-evaluaciones}
\end{equation}
\end{proposition}
\begin{proof}
La especialización elíptica de la exponencial trítica da
\(\cos\pi\,I+\sin\pi\,J_{\mathrm e}=-I\).
La nilpotencia elimina exactamente los términos de grado al menos \(2\)
en la segunda expresión. Para la tercera, \(\varphi^2=\varphi+1\) implica
\[
 \cosh(2\log\varphi)=\frac32,\qquad
 \sinh(2\log\varphi)=\frac{\sqrt5}{2}.
\]
Por tanto, la exponencial vale
\(\frac32I+\frac{\sqrt5}{2}J_{\mathrm h}=A\).
\end{proof}

La primera igualdad representa la media vuelta elíptica; la vuelta completa
satisface \(\exp(2\pi J_{\mathrm e})=I\). La segunda mantiene un término
lineal no nulo: el régimen parabólico expresa transporte, no ausencia de
evolución. La tercera realiza un avance hiperbólico unimodular. Sus
autovalores son \(\varphi^2\) y \(\varphi^{-2}\); una dirección se expande y
la otra se contrae conservando el determinante. Las tres evaluaciones
emplean parámetros distintos de una familia exponencial común.

En esa familia, \(e\) interviene mediante la evaluación escalar
\(\exp(I)=eI\) y la ley \(\exp(sI)\exp(tI)=\exp((s+t)I)\).
Su papel no se restringe al tipo parabólico. La coordenada \(e\), obtenida
previamente de su región, es reconocida aquí como la evaluación unitaria
de la exponencial; esta identidad no sustituye su genealogía coinductiva.

\subsubsection{Resolución polinómica de los regímenes}

Las tres realizaciones admiten una presentación compacta sobre la suma
directa de sus espacios. Defínanse
\[
 \mathbb J=J_{\mathrm e}\oplus J_{\mathrm p}\oplus J_{\mathrm h},
 \qquad Q=\mathbb J^2.
\]
Las relaciones anteriores implican
\[
 \mathbb J^6=\mathbb J^2,\qquad Q^3=Q.
\]
El uso de \(Q\) mantiene separado este cuadrado operatorial del registro
dodecafásico \(K\).

\begin{proposition}[Proyectores polinómicos de régimen]
\label{euler-trit-proyectores}
Las aplicaciones
\[
 p_{\mathrm e}=\frac{Q^2-Q}{2},\qquad
 p_{\mathrm p}=I-Q^2,\qquad
 p_{\mathrm h}=\frac{Q^2+Q}{2}
\]
son idempotentes mutuamente ortogonales, suman \(I\) y recuperan los tres
sumandos de la representación.
\end{proposition}
\begin{proof}
Los tres polinomios toman los valores \((1,0,0)\), \((0,1,0)\) y
\((0,0,1)\) en los puntos \(-1,0,+1\), respectivamente. Las diferencias
entre sus cuadrados y ellos mismos, así como sus productos cruzados, son
múltiplos de \(X^3-X\). Sustituir \(Q\), que anula ese polinomio, demuestra
las identidades. En los tres bloques, \(Q\) actúa como \(-I,0,I\).
\end{proof}

Se obtiene así
\begin{align}
 \exp(t\mathbb J)
 ={}&p_{\mathrm e}(\cos t\,I+\sin t\,\mathbb J)
   +p_{\mathrm p}(I+t\mathbb J)\notag\\
   &+p_{\mathrm h}(\cosh t\,I+\sinh t\,\mathbb J).
 \label{euler-trit-exponencial-total}
\end{align}
Cada sumando conserva su relación cuadrática y la conjugación invierte
su evolución. La representación total reúne los tres regímenes, pero no
define por sí sola un operador de transición entre ellos. Tal transición
requiere los mapas de transporte del TPK con sus dominios, orientación y
memoria. La articulación exponencial permite reconocer sus tipos sin
reemplazar esa dinámica por una identificación de las fibras.

% PROCEDENCIA FOCAL (no impresa)
% Resultado recuperado: tres generadores y tres evaluaciones exponenciales.
% Propietario completo:
% BASE/manuscrito/sections/hmt/09b_algebras_cuadraticas.tex:289-516.
% Antecedente causal de C y acción contragrediente:
% BASE/manuscrito/sections/hmt/03d_esqueleto_ciclotomico_app.tex:13-47,145-246.
% Los tres tipos y la orientación están en:
% BASE/manuscrito/sections/hmt/02_trit_geometrias.tex:83-196.
% Resolución polinómica recuperada:
% output/INVESTIGACION_HOLONOMIA_NONADICA_CRISTAL_APERIODICO_20260903/
% ALGEBRA_HOLONOMICA_UNIVERSAL_Y_EULER_TRITICO.md, secciones 8-11 y 22.
% Los cuatro propietarios se leyeron completos. BASE designa:
% /Users/ruben/Documents/HMT2/
% EL_CIERRE_HOLOGRAFICO_DEL_INFINITO_SUCESOR_102_CAPITULOS_2026-08-28_EN_TRABAJO/
% No se incorpora la identidad P9(alpha)=delta del documento especializado:
% esa identificación no pertenece a este bloque.
% Tampoco se identifica una dirección nilpotente con la vacancia electrónica,
% ni se presume una representación global del transporte por la suma directa.
% COMPROBACIONES: Python estándar, enteros y Fraction.
% C^2+C+I=0; (2C+I)^2=-3I; N^2=0; (2Fav^2-3I)^2=5I.
% 81 pares de residuos verifican las dos acciones contragredientes.
% Los 3 proyectores se verificaron en Q=-1,0,+1; prueba general en texto.
% 6 labels únicos y entornos equilibrados. Sin compilación.


% SOURCE foundation/sections/vacancias_capacidad.tex
% Desarrollo reunido de DESARROLLO_MATEMATICO, apartados 38--40 y 50.
\subsubsection{Índice de capacidad, vacancias e inducción del régimen trítico}
\label{vac-capacidad-trit}

La sucesión de vacancias está vinculada al refinamiento mediante una
identidad de incrementos. Para mostrarla se fija el origen de fase
\(\theta=0\). Sean
\[
 \rho=\log_{10}9=1-\delta,\qquad
 \mathcal C_t=\lfloor(t+1)\rho\rfloor\quad(t\geq0).
\]
El índice \(\mathcal C_t\) expresa la capacidad decimal del refinamiento
en la normalización que compara un bloque ternario de \(6\) posiciones,
con \(729\) posibilidades, y un bloque decimal de \(3\) posiciones, con
\(1000\) posibilidades. En efecto,
\(\log_{1000}729=\log_{10}9=\rho\). Se reserva la letra \(K\) para
el registro dodecafásico que se construirá posteriormente.

\begin{proposition}[Complementariedad de capacidad y vacancia]
\label{prop:capacidad-vacancia}
Para \(t\geq1\),
\begin{equation}
 \mathcal C_t-\mathcal C_{t-1}=1-z_t(0).
 \label{eq:capacidad-vacancia}
\end{equation}
En consecuencia, si \(0\leq a<b\),
\[
 \mathcal C_b-\mathcal C_a+
 \sum_{t=a+1}^{b}z_t(0)=b-a.
\]
\end{proposition}
\begin{proof}
La irracionalidad de \(\delta\) da
\[
 \mathcal C_t
 =t+1-\lceil(t+1)\delta\rceil
 =t-\lfloor(t+1)\delta\rfloor.
\]
La resta de dos términos consecutivos produce
\eqref{eq:capacidad-vacancia}. La suma sobre el intervalo es telescópica.
\end{proof}

Cada transición aumenta en una unidad el contador de bloques. Si
\(z_t=0\), aumenta también la capacidad; si \(z_t=1\), la capacidad se
mantiene mientras se incorpora una transición a la historia. La vacancia
es aquí un evento de retención de capacidad, con una posición y un
antecedente determinados. El balance anterior conserva exactamente la
longitud entre capacidad adquirida y vacancias registradas. La fase
observada \(g_t^{\rm cap}=[\mathcal C_t]_9\) satisface
\[
 g_t^{\rm cap}-g_{t-1}^{\rm cap}\equiv1-z_t\pmod9.
\]
Esta fase de capacidad y la fase cronológica \([t]_9\) son dos
observaciones diferentes. Una retención local de capacidad y una vuelta
completa de la holonomía tienen mecanismos distintos, aunque ambas
permitan conservar la observación de fase con modificación del estado.

La relación con TRIT puede seguirse en las propias posiciones de vacancia.
Para \(j\geq1\), defínanse
\[
 p_j=\left\lfloor\frac j\delta\right\rfloor,\qquad
 v_j=\mathcal C_{p_j},\qquad h_j=p_{j+1}-p_j.
\]
La desigualdad \(p_j\delta<j<(p_j+1)\delta\) identifica \(p_j\)
como la posición del evento \(j\). Como \(0<\delta<1\), da también
\(\lfloor(p_j+1)\delta\rfloor=j\), y por tanto
\[
 v_j=p_j-j,\qquad
 v_{j+1}-v_j=h_j-1.
\]
Si \(s_j=22-h_j\), entonces \(s_j\in\{0,1\}\) y
\begin{equation}
 [v_{j+1}]_3=[v_j-s_j]_3.
 \label{eq:vacancia-regimen}
\end{equation}
Un intervalo de longitud \(22\) conserva la clase; uno de longitud
\(21\) la desplaza una unidad en sentido negativo. El selector canónico
se aplica a la fase positiva \(r_j=1+[v_j]_9\):
\[
 \tau_j=M_{\rm ph}(r_j).
\]
Queda determinado así un régimen trítico por el índice de capacidad
retenido en cada evento, además de la codificación binaria de la vacancia.

Las posiciones de los intervalos cortos forman la codificación mecánica
inducida de pendiente \(22-1/\delta\). Sus separaciones \(6\) y \(7\)
determinan las longitudes de las mesetas de la clase \([v_j]_3\), una
vez fijados el origen y la convención de extremos. Las primeras clases
son \(2\) durante \(6\) eventos, \(1\) durante \(7\) y \(0\) durante
\(7\); el primer tramo de \(6\) es la restricción inicial de una meseta.
Sus firmas recorren \(0,-1,+1\). Las ecuaciones
\eqref{eq:capacidad-vacancia} y \eqref{eq:vacancia-regimen} explicitan el
enlace entre refinamiento, vacancias y selección del régimen. La lectura
cuadrática del TRIT realiza después esos tipos mediante sus respectivos
generadores; el acontecimiento temporal y su representación geométrica
conservan así una dependencia identificable.

La inversión del índice monótono proporciona otra observación útil:
\[
 \mathcal T(k)=\min\{t:\mathcal C_t=k\}
   =\left\lceil\frac k\rho\right\rceil-1\qquad(k\geq1).
\]
Si \(\sigma=\rho^{-1}-1\), el número de bloques necesarios para ganar
\(m\) unidades de capacidad es
\[
 \mathcal T(k+m)-\mathcal T(k)
 =m+\lceil(k+m)\sigma\rceil-\lceil k\sigma\rceil
 \in\{m+\lfloor m\sigma\rfloor,m+\lceil m\sigma\rceil\}.
\]
En particular, \(9\) unidades de capacidad requieren \(9\) o \(10\)
bloques. El período de una fase finita permanece exacto, mientras el
tiempo de retorno medido por el otro índice admite dos valores. Esta
distinción es parte de la dinámica del refinamiento y no una corrección
de sus valores numéricos a posteriori.


% SOURCE foundation/figures/holonomia.tex
\begin{figure}[htbp]
\centering
\begin{tikzpicture}[x=1cm,y=1cm,>=Latex,font=\small]
\begin{scope}[xshift=1.8cm,yshift=.3cm]
\draw[->,gray] (0,-.5)--(0,3.65) node[above] {\(n/9\)};
\foreach \k in {0,1,2,3}{
 \draw[gray!60,dashed] (-1.1,\k)--(1.1,\k);
 \node[left,font=\scriptsize] at (-1.1,\k){\(\k\)};
}
\draw[black,line width=.65pt,domain=0:27,samples=260,smooth,variable=\n]
 plot ({cos(40*\n)},{\n/9+.14*sin(40*\n)});
\foreach \n in {0,...,27}{
 \fill ({cos(40*\n)},{\n/9+.14*sin(40*\n)}) circle (.026);
}
\foreach \k in {0,1,2,3}{\fill (1,\k) circle (.045);}
\node[align=center,text width=3.3cm,font=\scriptsize] at (0,-1.02)
 {Misma fase en \(n=0,9,18,27\);\\grado de memoria distinto.};
\end{scope}
\begin{scope}[xshift=4.2cm,yshift=.5cm]
\node[anchor=west] at (0,3.1){\(z_n(0)=\lfloor(n+1)\delta\rfloor-\lfloor n\delta\rfloor\)};
\draw[->] (0,1.7)--(9,1.7) node[right]{\(n\)};
\foreach \n in {21,43,65,87,109,131,152,174,196,218}{
 \draw[line width=.6pt] ({\n*.039},1.7)--({\n*.039},2.25);
 \fill ({\n*.039},2.25) circle (.033);
 \node[below,font=\tiny] at ({\n*.039},1.7){\n};
}
\foreach \a/\b/\g in {21/43/22,43/65/22,65/87/22,87/109/22,109/131/22,131/152/21,152/174/22,174/196/22,196/218/22}{
 \node[font=\scriptsize] at ({(\a+\b)*.0195},2.55){\g};
}
\node[anchor=west,align=left,text width=8.3cm] at (0,.55)
 {La codificación irracional determina los intervalos entre eventos. El retorno de la fase nonádica no convierte esta secuencia en una palabra periódica.};
\end{scope}
\end{tikzpicture}
\caption{Representaciones complementarias del transporte. A la izquierda, proyección de la hélice \((\cos(2\pi n/9),\sin(2\pi n/9),n/9)\): el avance de nueve pasos recupera la fase e incrementa la altura. A la derecha, los diez primeros eventos de la palabra inferior para \(\delta=\log_{10}(10/9)\), con separaciones exactas \(21\) o \(22\). La representación helicoidal no introduce una métrica física adicional.}
\label{fig:holonomia-esturmiana}
\end{figure}


% SOURCE foundation/sections/revision_monodromia.tex
% Adición acumulativa. Insertar después de la suspensión esturmiana existente.
% Conservar íntegro el texto y la figura fig:holonomia-esturmiana.
\subsection{Conexión nonádica, holonomía y representación monodrómica}
\label{mono-ampl:conexion}

La conexión regional ya construida se estudia ahora mediante tres objetos:
el transporte del estado completo, la representación de sus vueltas y la
suspensión de sus itinerarios. Las operaciones de APP--TRIT--TPK y el
catálogo inicial conservan las definiciones anteriores. La distinción
necesaria concierne aquí a lo que cada representación retiene de una
misma historia (secciones~\ref{sec:extension} y~\ref{sec:generacion}).

En el índice de prolongación \(k\), una carta del estado conserva
\[
 \widehat x_k=(B_k,C_k,p_k,c_k,\Sigma_k,W_k,g_k,m_k).
\]
Aquí \(B_k\) reúne los tres bloques regionales, \(C_k\) es el cilindro,
\(p_k\) el prefijo, \(c_k\) el registro de acarreo, \(\Sigma_k\) la firma
de frontera y \(W_k\) la información incidencial. La fase es
\(g_k=1+[(k-1)]_9\); \(m_k\) cuenta las vueltas completas. La reducción
\(\beta_k=[m_k]_{12}\) conserva la posición dodecafásica, mientras el
entero \(m_k\) distingue sus sucesivas vueltas.

\begin{definition}[Transporte de una vuelta nonádica]
\label{mono-ampl:def-vuelta}
Sean \(\mathcal R_k\subseteq\widehat X_k\times\widehat X_{k+1}\)
las relaciones de extensión admisible del TPK. La realización de la
holonomía en el nivel \(k\) es la composición relacional
\[
 \Gamma_{9,k}=\mathcal R_{k+8}\circ\cdots\circ\mathcal R_k.
\]
Para \(s\geq1\), su prolongación a \(s\) vueltas es
\[
 \Gamma_{9s,k}=
 \Gamma_{9,k+9(s-1)}\circ\cdots\circ\Gamma_{9,k}.
\]
\end{definition}

La formulación relacional conserva las continuaciones admisibles. Sobre
una historia individualizada, cada flecha lleva al estado efectivamente
prolongado; el bloque visible aislado puede seguir admitiendo más de una
continuación. Esta diferencia es precisamente la función de la memoria y
de las condiciones de frontera en la determinación de la trayectoria.

\begin{proposition}[Iteración del retorno y memoria acumulada]
\label{mono-ampl:vueltas}
En toda composición admisible de \(s\) vueltas,
\[
 g_{k+9s}=g_k,\qquad m_{k+9s}=m_k+s,\qquad
 \beta_{k+9s}=\beta_k+s\pmod{12}.
\]
\end{proposition}
\begin{proof}
La regla de una vuelta conserva la fase e incrementa el contador en una
unidad. Su aplicación al estado ya transportado repite exactamente esas
dos operaciones. La inducción da las dos primeras igualdades; reducir la
segunda módulo \(12\) da la tercera.
\end{proof}

El contador proporciona así un testigo entero de la diferencia entre las
vueltas. Tras \(12\) vueltas la pareja de fases finitas se restaura, pero
el contador aumenta \(12\). La extensión no escindida
\eqref{eq:extension108} expresa la composición del calendario mediante su
cociclo; la trayectoria conserva además la elevación entera de ese dato.
La conservación significa aquí transporte de las contribuciones previas,
no inmovilidad de todas las coordenadas.

\subsubsection{Representación bilateral del índice de retorno}

El desplazamiento \(T\) de \eqref{eq:weyl-relojes} admite una
expresión en coordenadas de fase y vuelta. Numeramos las fases de \(0\) a
\(8\), mediante un desplazamiento del origen respecto de \(g_k\), y
escribimos \(k=9m+r\), con \(0\leq r<9\), en la recta entera
que recubre el ciclo de fases. Sobre
\(\mathcal H_{\rm ind}=\ell^2(\mathbb Z)\), la base
\(|k\rangle=|r,m\rangle\) expresa ese mismo desplazamiento unitario
\[
 T|r,m\rangle=
 \begin{cases}
 |r+1,m\rangle,&r<8,\\
 |0,m+1\rangle,&r=8.
 \end{cases}
\]
La inversa reduce \(r\) en una unidad cuando \(r>0\), y lleva
\(|0,m\rangle\) a \(|8,m-1\rangle\). El paso por el origen de fase
incorpora, por tanto, la unidad exacta de cociente.

\Needspace{8\baselineskip}
\begin{proposition}[Representación de la monodromía del ciclo]
\label{mono-ampl:representacion}
Sea \(S=T^9\). Entonces
\[
 S|r,m\rangle=|r,m+1\rangle,\qquad
 \mathbb Z\longrightarrow\mathcal U(\mathcal H_{\rm ind}),
 \quad s\longmapsto S^s
\]
es una representación unitaria fiel del grupo de vueltas del ciclo.
\end{proposition}
\begin{proof}
Nueve desplazamientos conservan el resto de la división por \(9\) e
incrementan el cociente. Las potencias satisfacen
\(S^{s+t}=S^sS^t\), y \(S^s|r,m\rangle=|r,m+s\rangle\) distingue
todo entero \(s\). La permutación de una base ortonormal es unitaria.
\end{proof}

Esta representación corresponde al recubrimiento del ciclo, cuyo grupo
de vueltas es \(\mathbb Z\), y representa la fase con su contador. El
estado del TPK contiene además las coordenadas de trayectoria ya descritas.
La extensión bilateral admite índices negativos como coordenadas del
recubrimiento; una trayectoria iniciada en una condición inicial utiliza
su semieje admisible. De este modo se dispone de una inversa operatoria
del índice sin convertir cada relación de extensión en una biyección del
estado completo.

Con \(D_M|k\rangle=\lfloor k/9\rfloor|k\rangle\), la identidad
\([D_M,S]=S\) expresa el incremento de memoria en el dominio denso de
vectores de soporte finito. Si \(\mathcal I|k\rangle=|-k\rangle\) y
\(P_0\) proyecta sobre los índices divisibles por \(9\), se obtiene
además
\begin{equation}
 \mathcal I T\mathcal I=T^{-1},\qquad
 \mathcal I D_M\mathcal I=-D_M-(I-P_0).
 \label{mono-ampl:inversion-contador}
\end{equation}
En efecto, para \(k=9m+r\),
\(\lfloor-k/9\rfloor=-m\) si \(r=0\) y vale \(-m-1\) si \(r>0\).
La inversión conserva así el término de frontera de la división euclídea.

\subsection{Reloj de bloques y reloj de capacidad}
\label{mono-ampl:relojes}

El índice de una aplicación de refinamiento y el número de posiciones
decimales disponibles tienen leyes relacionadas, pero diferentes. Para
expresarlas reservamos \(n\) al índice de bloques ternarios y definimos
\[
 \rho=\log_{1000}729=\log_{10}9,\qquad
 \delta=1-\rho,\qquad
 \mathcal C_n=\lfloor(n+1)\rho\rfloor,\quad n\geq0.
\]
La letra \(\mathcal C_n\) designa aquí la capacidad, no el registro
dodecafásico \(K\). En el origen de fase canónico, la palabra mecánica
anterior proporciona
\[
 Z_n=\sum_{j=0}^{n}z_j(0)=\lfloor(n+1)\delta\rfloor.
\]

\begin{proposition}[Balance exacto entre capacidad y vacancias]
\label{mono-ampl:balance-capacidad}
Para \(n\geq0\) y, en la segunda identidad, \(n\geq1\),
\[
 \mathcal C_n=n-Z_n,\qquad
 \mathcal C_n-\mathcal C_{n-1}=1-z_n(0).
\]
\end{proposition}
\begin{proof}
La irracionalidad de \(\delta\) implica
\(\lceil(n+1)\delta\rceil=\lfloor(n+1)\delta\rfloor+1\).
Entonces
\[
 \lfloor(n+1)(1-\delta)\rfloor
 =n+1-\lceil(n+1)\delta\rceil=n-Z_n.
\]
Restar dos identidades consecutivas concluye la prueba.
\end{proof}

Un evento \(z_n=1\) añade un bloque a la historia manteniendo la capacidad
\(\mathcal C_n\); un evento \(z_n=0\) incrementa ambos contadores. En
cualquier intervalo, los incrementos de capacidad y las vacancias suman
exactamente su longitud. Reduciendo módulo \(q\), para \(q=9\) o
\(q=108\), resulta
\[
 [\mathcal C_n]_q=[n]_q-[Z_n]_q.
\]
La fase uniforme de bloques y la fase de capacidad se relacionan mediante
la memoria integrada de vacancias. En particular, una retención de
capacidad y una vuelta completa de la conexión tienen mecanismos
distintos, aunque ambas puedan conservar una observación de fase.

La inversión monótona del reloj da un tiempo de primera llegada:
\[
 T_{\rm cap}(a)=\min\{n:\mathcal C_n=a\}
 =\left\lceil\frac a\rho\right\rceil-1,
 \qquad a\geq1.
\]
Con \(\sigma=\rho^{-1}-1\), el tiempo empleado en ganar \(b\) unidades
de capacidad es
\[
 T_{\rm cap}(a+b)-T_{\rm cap}(a)
 =b+\lceil(a+b)\sigma\rceil-\lceil a\sigma\rceil.
\]
La diferencia de techos toma uno de los dos enteros adyacentes a
\(b\sigma\). Así se obtienen \(9\) o \(10\) bloques para ganar
\(9\) unidades de capacidad, y \(113\) o \(114\) para ganar \(108\).
Son tiempos del reloj inverso, medidos desde una primera llegada; el
contador de vueltas se interpreta siempre en el índice al que pertenece.

\subsection{Suspensión simbólica y cristal temporal aperiódico}
\label{mono-ampl:cristal}

La codificación por rotación reúne periodicidad de fase, recurrencia local
y aperiodicidad global. Sea \(\mathcal X_\delta\subseteq\{0,1\}^{\mathbb Z}\)
la clausura, en la topología producto, de las traslaciones de la palabra
bilateral \(z_n(\theta)\). Denotemos por \(\sigma_{\rm sh}\) el
desplazamiento de sus sucesiones. Conservando una fase nonádica, el paso es
\[
 F(r,z)=([r+1]_9,\sigma_{\rm sh}z),\qquad
 (r,z)\in C_9\times\mathcal X_\delta.
\]

\begin{definition}[Cristal temporal aperiódico simbólico]
\label{mono-ampl:def-cristal}
La suspensión de techo unitario del sistema anterior es el espacio
\[
 \mathcal S_\delta=
 \bigl((C_9\times\mathcal X_\delta)\times\mathbb R\bigr)/\sim,
 \qquad (y,t+1)\sim(Fy,t),
\]
con flujo \(\Phi_s[y,t]=[y,t+s]\). Se denomina \emph{cristal temporal
aperiódico simbólico} a este modelo de orden temporal, junto con su
codificación de vacancias y la representación del contador de vueltas.
\end{definition}

El adjetivo temporal designa la acción del flujo y de su sección de retorno;
la aperiodicidad caracteriza sus itinerarios. La frecuencia de unos de cada
sucesión es \(\delta\): la cota uniforme de discrepancia menor que uno,
obtenida por telescopía en cualquier ventana, persiste al tomar la clausura.
Una sucesión periódica tendría frecuencia racional. Por tanto,
\(\mathcal X_\delta\) no contiene órbitas periódicas. La suspensión
tampoco tiene órbitas cerradas: un retorno del flujo exigiría un tiempo
entero y una periodicidad del itinerario.

Las palabras finitas sí reaparecen. Los itinerarios de longitud \(m\)
proceden de una partición de la circunferencia por \(m+1\) puntos; cada
intervalo es visitado recurrentemente por la rotación irracional. Incluso
al restringir los tiempos a una fase fija, el paso \(9\delta\) sigue
siendo irracional. Ésta es la razón de que cada palabra aparezca en las
\(9\) fases y de las complejidades ya establecidas
\[
 p(m)=m+1,\qquad p_9(m)=9(m+1),\qquad p_3(m)=3(m+1).
\]
La recurrencia de configuraciones locales coexiste con la ausencia de un
período global. Su entropía topológica es nula porque estas complejidades
crecen linealmente.

La realización por los observables \(V_9,W_\delta\) y el desplazamiento
\(T\) expresa el mismo orden mediante las covariancias
\eqref{eq:weyl-relojes}. Las frecuencias del reloj uniforme pertenecen a
\(\frac19\mathbb Z+\delta\mathbb Z\), módulo \(1\). El reloj de
capacidad conserva adicionalmente la relación
\([\mathcal C_n]_9=[n]_9-[Z_n]_9\); no se sustituye por la fase uniforme
al interpretar los retornos. El término \emph{cristal temporal} queda así
definido mediante un sistema dinámico matemático. Su realización física
requiere especificar una dinámica energética, sus observables temporales
y su estabilidad, datos distintos de la construcción simbólica presente.

\subsection{Simetría especular de las regiones y doble realización}
\label{mono-ampl:regiones}

El intercambio especular actúa en el bloque regional
\(B=(u^\pi,u^e,u^\varphi)\in(\mathbb F_3^6)^3\) por
\[
 \mathfrak cB=(u^e,u^\pi,u^\varphi).
\]
Su eje fijo y su agregado son, respectivamente, \(u^\varphi\) y
\(s=u^\pi+u^e\). Esta fijación se refiere a la involución en una fibra;
ambas coordenadas pueden evolucionar durante el transporte nonádico.
Con \(q=u^\pi+u^e+u^\varphi\),
\(a=u^\pi+u^e-u^\varphi\) y \(d=u^\pi-u^e\), se recupera
\[
 u^\varphi=2(q-a),\qquad s=q-u^\varphi,\qquad
 u^\pi=2(s+d),\qquad u^e=2(s-d).
\]
Las igualdades son componente a componente en \(\mathbb F_3\), donde
\(2^{-1}=2\). El espejo conserva \(q,a,s\) e invierte \(d\); esta
coordenada antisimétrica individualiza el reparto que el agregado conserva
sin distinguir. La inversión \(\mathcal I\) del índice temporal y la
involución \(\mathfrak c\) de canales actúan, por consiguiente, sobre
datos diferentes.

El retorno de fase posee un testigo explícito:
\[
 \begin{aligned}
 B_1&=(012222\mid121221\mid112202),\\
 B_{10}&=(101222\mid112221\mid112002).
 \end{aligned}
\]
Las dos posiciones tienen fase \(1\); el bloque, el eje y la memoria han
cambiado. Asimismo, la sincronización de los censos en las fases \(8\) y
\(9\) anula sus diferencias relativas, mientras la componente diagonal
pasa de \((59,59,59)\) a \((43,43,43)\). La igualdad de un observable
relativo no agota el estado que lo determina.

La doble realización conserva esa arquitectura. En la sección
\ref{exc:doble-lectura}, el mismo prefijo determina por una parte cilindros
encajados y por otra una sucesión de palabras codificadas con sus marcos.
El primer mapa pasa al límite por contracción de diámetros; el segundo
satisface \(r_{n,m}Q_n=Q_mp_{n,m}\) y define \(Q_\infty\).
La figura \ref{mono-ampl:fig-regional} anticipa esas dos aplicaciones.
El recorrido incidencial posterior construye códigos, retículos y la
realización de Moonshine en sus dominios propios; el grupo de automorfismos
actúa sobre esa realización, no sobre las cifras por una identificación
implícita.

\input{figures/monodromia_regional_ampliacion.tex}

\input{sections/recta_nonadica_recuperada.tex}

\subsection{Transporte temporal y composición exponencial trítica}
\label{mono-ampl:euler}

Las relaciones de Euler de la sección \ref{euler-trit-articulacion}
describen las realizaciones locales de los tres regímenes. Para una firma
fija \(\tau\), el generador satisface \(J_\tau^2=-\tau I\); la
conjugación algebraica lleva \(J_\tau\) a \(-J_\tau\). En consecuencia,
\[
 \exp(sJ_\tau)\exp(tJ_\tau)=\exp((s+t)J_\tau),\qquad
 \overline{\exp(tJ_\tau)}=\exp(-tJ_\tau),
\]
y el producto de una evolución con su conjugada es la unidad. En el tipo
elíptico se realiza mediante \(\cos^2t+\sin^2t=1\); en el hiperbólico,
mediante \(\cosh^2t-\sinh^2t=1\); en el parabólico,
\((I+tN)(I-tN)=I\), pues \(N^2=0\). La misma ley de conservación
admite tres expresiones cuadráticas.

La realización elíptica cierra una vuelta; la nilpotente conserva un
desplazamiento lineal; la hiperbólica separa expansión y contracción con
determinante unitario. En los parámetros regionales ya generados,
\[
 \exp(\pi J_{\mathrm e})=-I,\qquad
 \exp(J_{\mathrm p})=I+J_{\mathrm p},\qquad
 \exp(2\log\varphi\,J_{\mathrm h})=F_{\rm av}^{\,2}.
\]
La evaluación \(\exp(I)=eI\) reconoce la unidad de la ley exponencial
común. El número \(e\) no se asigna exclusivamente al régimen parabólico.

El transporte de una realización local por un isomorfismo \(U\) conserva
estas relaciones: si \(J'=UJU^{-1}\), la serie exponencial da
\(U\exp(tJ)=\exp(tJ')U\). La identidad transporta el régimen y su
norma. Una transición entre tipos cuadráticos diferentes pertenece, en
cambio, a la selección de hojas y regímenes del TPK; no es una conjugación
que cambie el valor de \(J^2\). El producto temporal conserva por ello el
orden de las transiciones y la información de sus extremos.

Esta articulación distingue la clausura de una representación local y la
holonomía de una historia. Puede cumplirse \(\exp(2\pi J_{\mathrm e})=I\)
al mismo tiempo que \(S|r,m\rangle=|r,m+1\rangle\). La primera igualdad
restaura la orientación del régimen; la segunda registra otra vuelta. La
estructura discreta conserva ambas: el retorno local y la procedencia
acumulada del estado que lo realiza.


% SOURCE foundation/figures/monodromia_regional_ampliacion.tex
% Figura acumulativa; destino figures/monodromia_regional_ampliacion.tex.
\begin{figure}[H]
\centering
\begin{tikzpicture}[x=1cm,y=1cm,>=Latex,font=\small,
  flow/.style={->,line width=.55pt},
  ann/.style={align=center,font=\scriptsize}]
\node[font=\bfseries] at (6.9,9.2)
  {Conexión nonádica y coordenadas regionales};
\begin{scope}[shift={(3.4,6.8)}]
 \foreach \k in {1,...,9}{
  \pgfmathsetmacro{\ang}{90-40*(\k-1)}
  \coordinate (p\k) at (\ang:1.52);
  \fill (p\k) circle (1.4pt);
  \node[fill=white,inner sep=1.1pt,font=\scriptsize]
    at (\ang:1.83) {\(\k\)};
 }
 \foreach \k in {1,...,8}{
  \pgfmathtruncatemacro{\next}{\k+1}
  \draw[flow,shorten >=3pt,shorten <=3pt] (p\k)--(p\next);
 }
 \draw[flow,shorten >=3pt,shorten <=3pt] (p9)--(p1);
 \node[ann] at (0,.1) {\(g_k\in C_9\)\\\(\Gamma_{9,k}\)};
 \node[ann] at (0,-2.23) {Una vuelta: \(g\mapsto g\), \(m\mapsto m+1\)};
\end{scope}
\node[ann,text width=5.25cm] at (10.1,7.75)
 {\(B_k=(u_k^\pi,u_k^e,u_k^\varphi)\)\\
 bloque regional en una fibra};
\node at (8.55,6.5) {\(u^\pi\)};
\node at (11.65,6.5) {\(u^e\)};
\draw[<->,line width=.65pt] (8.95,6.5)--node[above]{\(\mathfrak c\)}(11.25,6.5);
\draw[densely dashed,gray] (10.1,5.72)--(10.1,7.23);
\node[fill=white,inner sep=2pt] at (10.1,5.6){\(u^\varphi\)};
\node[ann] at (10.1,4.6)
 {Eje \(u^\varphi\) y agregado \(u^\pi+u^e\)\\
 invariantes bajo \(\mathfrak c\)};
\draw[flow] (5.45,6.8)--(7.25,6.8);
\draw[gray!60,line width=.3pt] (.55,4.02)--(13.25,4.02);
\node[ann,text width=12cm] (hist) at (6.9,3.47)
 {Historia compatible: prefijos, cilindros, orientación, cocientes,
 marcos de incidencia y memoria};
\coordinate (fork) at (6.9,2.85);
\draw[flow] (hist.south)--(fork);
\draw[flow] (fork)--(3.6,2.2);
\draw[flow] (fork)--(10.3,2.2);
\node[ann,text width=5.6cm] at (3.6,1.65)
 {Realización arquimediana\\
 \(I_{n+1}\subseteq I_n,\quad |I_n|=729^{-n}\)\\
 límite de los prefijos regionales};
\node[ann,text width=5.6cm] at (10.3,1.65)
 {Realización incidencial\\
 \(r_{n,m}Q_n=Q_mp_{n,m}\)\\
 palabras codificadas y marcos compatibles};
\node[ann,text width=5.6cm] at (3.6,.42)
 {\(\pi,\varphi,e\); determinación dodecafásica de \(\alpha\)};
\node[ann,text width=5.6cm] at (10.3,.42)
 {Witt--Golay--Leech--\(V^\natural\)\\
 automorfismos de la realización excepcional};
\draw[flow] (3.6,1.03)--(3.6,.78);
\draw[flow] (10.3,1.03)--(10.3,.78);
\end{tikzpicture}
\caption{Nueve fases de una conexión y dos realizaciones de la historia
compatible. El ciclo representa la fase; cada vuelta incrementa el contador.
El diagrama regional representa la involución \(\mathfrak c\), que
intercambia clausura y propagación y fija el eje de autoescala en una fibra.
Esta invariancia no impone que el eje sea constante durante el transporte.
Las dos aplicaciones inferiores utilizan el mismo prefijo: una determina
su límite numérico y la otra conserva palabras y marcos de incidencia. Los
mapas de esta segunda realización se especifican en la sección
\ref{exc:doble-lectura}; el recorrido excepcional se desarrolla después.}
\label{mono-ampl:fig-regional}
\end{figure}


% SOURCE foundation/sections/recta_nonadica_recuperada.tex
\subsubsection{Representación rectilínea de las fases y del retorno}
\label{mono-recta:desarrollo}

La representación sobre el segmento \([0,10]\) conserva el origen
posicional de APP y ordena sus nueve marcas interiores. En este diagrama,
los enteros \(1,\ldots,9\) son representantes de fase. La posición
horizontal indica el orden de las cartas; las etiquetas superiores
especifican operaciones regionales. La figura
\ref{mono-recta:figura} reúne así la recta de partida, la fijación del
eje de autoescala, la distribución especular de los dos canales y el
retorno con memoria (sección~\ref{mono-ampl:conexion}).

\input{figures/recta_nonadica_0_10.tex}

La fijación del eje se formula en la fibra de una frontera dada. Si
\(q=u^\pi+u^e+u^\varphi\) y
\(a=u^\pi+u^e-u^\varphi\), entonces
\[
 u^\varphi=2(q-a),\qquad s=u^\pi+u^e=q-u^\varphi
 \quad\text{en }\mathbb F_3^6.
\]
Los datos \((q,a)\) determinan el eje y el agregado; el reparto ordenado
de \(s\) conserva la libertad especular resuelta por la prolongación.
La quinta fase posee una solución local y una extensión, con dato
de frontera \(c=111111\) y firma residual \(\rho_W=000\). Éste es el
sentido de su primera saturación fuerte. El símbolo \(\varphi\)
situado sobre la marca \(5\) identifica el eje regional que interviene
en esa operación.

Las fases intermedias conservan esta descomposición dentro de cada fibra,
mientras sus datos de frontera evolucionan. El censo completo de
multiplicidades locales y extensiones de la rama distinguida es
\[
\begin{array}{c|rrrrrrrrr}
\text{fase}&1&2&3&4&5&6&7&8&9\\\hline
N_{\rm loc}&3&2&5&4&1&1&3&3&2\\
N_{\rm ext}&2&5&4&1&1&3&3&2&1
\end{array}
\]
La fase \(6\) combina unicidad local y tres prolongaciones; la fase
\(7\) resuelve una fibra especular triple; la fase \(8\) sincroniza
los tres censos regionales. La fase \(9\) conserva dos distribuciones
locales con el mismo eje y el mismo agregado, de las cuales una se
prolonga. Estos recuentos expresan funciones diferentes dentro de una
misma conexión, y justifican las marcas de la recta.

En particular, para las fases \(6,7,8\) las sumas regionales son
respectivamente \(012100\), \(101001\) y \(112000\). La conservación
representada por la llave es una invariancia respecto del reparto
especular en cada fibra. La evaluación real posterior aplica los
lectores de los cilindros compatibles a las historias completas; la
suma interna representada en la figura pertenece al dominio ternario.

Finalmente, el arco \(9\to1\) representa el paso de una vuelta a la
siguiente. Las fórmulas de la sección
\ref{mono-ampl:conexion} conservan el resto de fase e incrementan el
contador de vueltas. El diagrama rectilíneo y el diagrama cíclico de la
figura \ref{mono-ampl:fig-regional} expresan, por tanto, dos aspectos
complementarios de la misma holonomía: orden de las posiciones y
composición del retorno sobre el estado con memoria.


% SOURCE foundation/figures/recta_nonadica_0_10.tex
% Recuperación de fig_nonadic_route_mini.tex de la síntesis académica.
% Los rótulos designan operaciones regionales; la suma se efectúa en F_3^6.
\begin{figure}[H]
\centering
\begin{tikzpicture}[x=1.12cm,y=1cm,font=\small,>=Latex,
  guide/.style={line width=.45pt,black!55}]
 \draw[line width=.8pt] (0,0)--(10,0);
 \foreach \k in {0,...,10}{
  \draw[guide] (\k,-.09)--(\k,.09);
  \node[below=2.5mm] at (\k,0) {$\k$};
 }
 \foreach \k in {1,...,9}{\fill (\k,0) circle (1.7pt);}
 \node[anchor=west,font=\footnotesize] at (0,3.05)
   {Segmento generador $[0,10]$ y nueve representantes de fase};
 \draw[guide] (5,.12)--(5,1.7);
 \node[above,align=center,font=\footnotesize] at (5,1.7)
   {fijación del eje\\$u^\varphi$};
 \draw[decorate,decoration={brace,amplitude=4pt},guide]
   (6,.65)--(8,.65);
 \node[above=5pt,align=center,font=\footnotesize] at (7,.65)
   {agregado regional\\$u^e+u^\pi$};
 \draw[guide] (9,.12)--(9,1.7);
 \node[above,align=center,font=\footnotesize] at (9,1.7)
   {involución especular\\$\pi\leftrightarrow e$};
 \draw[->,line width=.7pt] (9,-.7)
   .. controls (9,-1.55) and (1,-1.55) .. (1,-.7);
 \node[fill=white,inner sep=3pt,font=\footnotesize] at (5,-1.23)
   {retorno $9\to1$;\quad $m\mapsto m+1$};
 \node[font=\footnotesize] at (5,-1.95)
   {$w_6\to w_{12}\to w_{18}\to w_{24}\to w_{30}\to R_{36}\to G_9$};
\end{tikzpicture}
\caption{Representación rectilínea de la conexión nonádica. La marca
\(5\) sitúa la primera saturación fuerte y el eje regional de
autoescala; la llave \(6\)--\(8\) identifica el agregado ternario
fijado en cada fibra; la fase \(9\) selecciona la orientación de
clausura. El retorno conserva la fase e incrementa la memoria. Las
abscisas numeran posiciones operativas, mientras las etiquetas
\(u^\pi,u^e,u^\varphi\) designan bloques de \(\mathbb F_3^6\).}
\label{mono-recta:figura}
\end{figure}


% SOURCE foundation/sections/primos_retornos.tex
% Recuperación del producto nativo y su selector, no incorporación de RH.
\subsection{Irreducibilidad multiplicativa y períodos de retorno}
\label{primos-retornos}

La distinción entre estructura, evaluación y fase aparece también en el
sector multiplicativo. El producto local de APP contiene tanto el residuo
como el cociente; su iteración permite construir formas normales y
factorizaciones antes de examinar sus períodos decimales. Este desarrollo
sitúa la expresión \emph{modo primo} en un dominio preciso y relaciona la
irreducibilidad con las operaciones ya definidas, no con una coincidencia
aislada de períodos (sección~\ref{primos-retornos}).

Sea
\[
 \mathcal W_9=\{\varnothing\}\sqcup
 \bigsqcup_{\ell\geq1}\{1,\ldots,9\}^{\ell}.
\]
Las secuencias se ordenan desde el dígito menos significativo. La evaluación
\(\operatorname{ev}_9(u)=\sum_j u_j9^j\), con
\(\operatorname{ev}_9(\varnothing)=0\), es biyectiva sobre los enteros no
negativos: para un entero positivo, la división de \(n-1\) por \(9\)
determina de manera única el primer dígito y la continuación. Esta es la
numeración nonádica biyectiva; conserva el representante \(9\) de la
frontera.

La suma \(\boxplus_\#\) se obtiene aplicando la hoja aditiva de APP a los
dígitos de cada posición y propagando sus cocientes. Si sólo hay un dígito,
éste constituye el residuo provisional; un cociente entrante no nulo se
normaliza mediante una segunda operación aditiva. La suma de los
cocientes se transporta a la posición siguiente. Para multiplicar una
secuencia por un dígito \(d\), la celda multiplicativa produce
\(u_jd=9q_j+r_j\), y el cociente entrante \(c_j\) se incorpora mediante
\[
 (w_j,c_{j+1})=
 \begin{cases}
 (r_j,q_j),&c_j=0,\\
 \bigl(R^+(r_j,c_j),q_j+q^+(r_j,c_j)\bigr),&c_j>0.
 \end{cases}
\]
Se comienza con \(c_0=0\) y se incorpora el cociente final cuando es no
nulo. Durante la suma, \(c_j\leq2\); durante el producto por un dígito,
\(c_j\leq9\). Las operaciones locales reciben así argumentos dentro
de su dominio. Denotemos por \(d\odot_\#u\) este segundo procedimiento,
con \(d\odot_\#\varnothing=\varnothing\).

Para \(v=(v_0,\ldots,v_{m-1})\), el producto completo se construye por
la recurrencia
\[
 a_m=\varnothing,\qquad
 a_j=(9\odot_\#a_{j+1})\boxplus_\#(v_j\odot_\#u),
 \qquad u\boxtimes_\#v=a_0.
\]
La regla compone las dos hojas locales. Su evaluación verifica
\begin{equation}
 \operatorname{ev}_9(u\boxtimes_\#v)
   =\operatorname{ev}_9(u)\operatorname{ev}_9(v).
 \label{eq:primos-entrelazamiento}
\end{equation}
En efecto, cada transición conserva la suma parcial más el cociente
pendiente ponderado por la potencia de \(9\). La recurrencia da entonces
\(\operatorname{ev}_9(a_j)=9\operatorname{ev}_9(a_{j+1})+
v_j\operatorname{ev}_9(u)\), cuya iteración prueba
\eqref{eq:primos-entrelazamiento}. La unicidad de la forma normal transporta
asociatividad y conmutatividad al producto construido; su unidad es \((1)\).

\begin{definition}[Coproducto multiplicativo reducido]
Sobre el espacio vectorial de soporte finito generado por las formas
normales no vacías se define
\[
 \widetilde\Delta_\# e_w
 =\sum_{\substack{u\boxtimes_\#v=w\\u,v\ne(1)}}e_u\otimes e_v.
\]
Una forma normal no unitaria es irreducible cuando su imagen por
\(\widetilde\Delta_\#\) es nula.
\end{definition}

Las fibras de factorización son finitas por
\eqref{eq:primos-entrelazamiento}. En la completación
\(\mathcal H_\#=\ell^2(\mathcal W_9\setminus\{\varnothing\})\),
las imágenes de formas distintas tienen soportes disjuntos. Si
\(d_\#(w)\) cuenta las factorizaciones ordenadas no unitarias, el dominio
de la clausura es
\[
 \operatorname{Dom}(\overline{\widetilde\Delta_\#})
 =\left\{x\in\mathcal H_\#:
       \sum_w d_\#(w)|x_w|^2<\infty\right\}.
\]
En efecto, la norma cuadrada de la imagen es esa suma ponderada. Por ello,
el operador
\[
 A_\#=(\overline{\widetilde\Delta_\#})^*
                  \overline{\widetilde\Delta_\#}
\]
es diagonal, positivo y autoadjunto: su valor propio en \(e_w\) es el
número de factorizaciones ordenadas no unitarias de \(w\). La proyección
\[
 P_{\rm irr}=(I-|e_{(1)}\rangle\langle e_{(1)}|)
                    \mathbf1_{\{0\}}(A_\#)
\]
selecciona exactamente las formas irreducibles. Mediante la evaluación
multiplicativa, éstas corresponden a los primos ordinarios. La selección
depende de la estructura de factorización, y mantiene separado el estado
unitario.

Para un valor \(n\geq2\) coprimo con \(10\), la división decimal induce la
permutación de residuos \(r\mapsto10r\pmod n\). El retorno del residuo
\(1\) tiene período
\[
 \tau_n=\operatorname{ord}_n(10).
\]
Si \(n=\prod_p p^{a_p}\), el teorema chino de los restos da
\[
 \tau_n=\operatorname{mcm}_{p^{a_p}\parallel n}
                    \operatorname{ord}_{p^{a_p}}(10).
\]
El retorno compuesto sincroniza los retornos de las potencias primas.
Para un primo \(p\notin\{2,3,5\}\), se tiene
\[
 \tau_p\mid p-1,\qquad
 M_p=\frac{10^{\tau_p}-1}{p}\in\mathbb N,
 \qquad M_p\equiv0\pmod9.
\]
La última congruencia expresa el cierre nonádico de la repetición decimal:
\(10^{\tau_p}-1\) es múltiplo de \(9\) y \(p\) es invertible módulo
\(9\). La misma congruencia vale para cualquier \(n\) coprimo con
\(90\). En particular, \(7,13,77,91,143\) tienen período decimal \(6\),
aunque sólo los dos primeros son irreducibles. El coproducto y el período
registran, por tanto, propiedades diferentes de un mismo valor construido.

Esta diferenciación precisa la relación con el refinamiento aperiódico.
La secuencia de vacancias determina los incrementos del índice de
capacidad; la factorización determina los modos multiplicativos; la
división decimal determina sus tiempos de retorno. El representante
nonádico conserva el cierre de fase, mientras que el cociente y la
factorización conservan información que ese cierre no distingue. El
enlace con las construcciones espectrales del tratado reside en esos
operadores y sus entrelazamientos. La demostración de sus resultados
analíticos posteriores mantiene su desarrollo propio; no se sustituye por
la periodicidad de una observación residual.

\subsubsection{Discrepancia de vacancias y lecturas de Dirichlet}

El vínculo admite también una expresión analítica exacta. Se conserva la
misma codificación \(z_n=z_n(0)\), para \(n\geq1\), y se consideran
dos observables posteriores: la serie sobre todos los índices y el
producto sobre los índices seleccionados por irreducibilidad. Para
\(\operatorname{Re}s>1\),
\[
 Z_{\rm vac}^{+}(s)=\sum_{n\geq1}\frac{z_n}{n^s},\qquad
 Z_{\rm vac}^{\times}(s)=\prod_p(1-p^{-s})^{-z_p}.
\]
Ambos convergen absolutamente en ese semiplano. Su factor común es la
densidad \(\delta\), pero sus discrepancias son diferentes.

\Needspace{21\baselineskip}
\begin{proposition}[Separación de la densidad de vacancias]
Se tiene
\[
 Z_{\rm vac}^{+}(s)=\delta\zeta(s)+H_{\rm vac}(s),
 \qquad H_{\rm vac}\in\mathcal O(\operatorname{Re}s>0).
\]
En \(\operatorname{Re}s>1\), con logaritmos definidos por las series
absolutamente convergentes de Euler,
\[
 Z_{\rm vac}^{\times}(s)=\zeta(s)^\delta G_{\rm vac}(s),\qquad
 \log G_{\rm vac}(s)=
 \sum_p(z_p-\delta)\log(1-p^{-s})^{-1}.
\]
\end{proposition}
\begin{proof}
Las sumas parciales centradas satisfacen
\[
 A_N:=\sum_{n=1}^N(z_n-\delta)
 =\lfloor(N+1)\delta\rfloor-N\delta,\qquad |A_N|<1.
\]
La sumación de Abel representa el término centrado por
\(s\int_1^\infty A_{\lfloor x\rfloor}x^{-s-1}\,dx\). La integral
converge localmente de manera uniforme en \(\operatorname{Re}s>0\),
lo que prueba la holomorfía. La segunda identidad resulta de separar
\(z_p=\delta+(z_p-\delta)\) en el logaritmo del producto, dentro de
su dominio de convergencia absoluta.
\end{proof}

La descomposición hace visible qué añade la lectura sobre primos: examina
la misma sucesión de vacancias sobre un subconjunto determinado por la
factorización. El control de la discrepancia sobre todos los enteros da
la primera prolongación holomorfa; el segundo observable conserva una
discrepancia propia sobre primos. Son dos aplicaciones de la misma
estructura discreta, con dominios analíticos expresos. Esta articulación
explica su pertinencia en el núcleo aritmético sin trasladar a este
artículo la demostración de los resultados terminales del tratado.


% SOURCE foundation/sections/registro_k.tex
\section{Extracción dodecafásica de la constante de acoplamiento aritmético-geométrico}
\label{sec:k-registro}

\begin{definition}[Constante holográfica de acoplamiento aritmético-geométrico]
\label{lib01:definicion}
Se denomina así al registro finito ordenado producido por la extracción
orientada del estado APP--TRIT--TPK, con su origen de fase, normalización y
lectores declarados, cuya representación posicional interviene en la
clausura aritmética y cuyas realizaciones incidencial y geométrica actúan
sobre el mismo registro. El símbolo \(K\) designa su representación
integral canónica; \(\kappa_{\rm ag}\) designa su lectura racional
periódica en la carta fijada.
\end{definition}

La generación de las coordenadas regionales no agota la información del
estado APP--TRIT--TPK. Una lectura arquimediana determina un valor mediante
cilindros compatibles; otra lectura conserva la oposición de las hojas, la
orientación de la ruta y el registro de las incidencias atravesadas. El
registro dodecafásico pertenece a esta segunda lectura antes de intervenir en
la clausura de alfa. Por ello no se define restando una constante física a
una combinación de valores previamente conocidos. Se reconstruye desde una
coordenada firmada del mismo estado aritmético con memoria de trayectoria que produce las coordenadas
regionales.

La distinción es necesaria incluso cuando dos objetos tienen doce
coordenadas. Una palabra de doce bloques del catálogo, una lista de doce
eventos orientados, un vector entero firmado y un registro dodecafásico en base mil son
objetos distintos. Su comparación requiere las aplicaciones que los
relacionan. La igualdad de longitud no proporciona esas aplicaciones, y la
pérdida de signos no es una mera simplificación de notación.

Este capítulo mantiene el corte causal fijado en los capítulos precedentes:
las nueve posiciones APP de los ejes horizontal y vertical producen las
hojas aritméticas; TRIT tipa régimen y orientación; TPK selecciona,
transporta, actualiza y prolonga, conservando residuo, cociente, acarreo,
hoja, ruta, frontera y memoria. Desde la estructura discreta conjunta del
continuo se obtienen las coordenadas que se utilizan aquí. La cadena
\(w_6\to w_{12}\to w_{18}\to w_{24}\to w_{30}\to R_{36}\to G_9\)
no se sustituye por un calendario finito. Tampoco se invierte el orden
\(U_t\to\Gamma_9\to K_{\mathrm{ph}}\): la última aplicación es una
lectura reducida de memoria, no el generador del registro firmado.

\input{sections/revision_k.tex}

\subsection{Extensión cíclica y graduación de memoria}
\label{subsec:k-calendario}

El calendario observable contiene ciento ocho posiciones, organizadas en
doce ventanas nonádicas. Con índices positivos, sean
\[
 E_m=\{9(m-1)+1,\ldots,9m\},\qquad 1\leq m\leq12.
\]
El soporte ordenado \(E_1\sqcup\cdots\sqcup E_{12}\) conserva posición y
pertenencia a ventana. El grupo \(C_{108}=\mathbb Z/108\mathbb Z\)
conserva, en cambio, la ley de traslación cíclica. Para el representante
\(0\leq\widetilde\theta<108\), las coordenadas
\[
 \beta(\theta)=\left[\left\lfloor\frac{\widetilde\theta}{9}\right\rfloor\right]_{12},
 \qquad r(\theta)=1+(\widetilde\theta\bmod9)
\]
separan el sector dodecafásico y la posición nonádica interior. Ninguna de
ellas conserva por sí sola el número total de vueltas.

\begin{proposition}[El acarreo del calendario]
\label{prop:k-extension}
Las aplicaciones
\[
 \iota:C_{12}\longrightarrow C_{108},\quad[b]\longmapsto[9b],
 \qquad p:C_{108}\longrightarrow C_9,\quad[\theta]\longmapsto[\theta]_9
\]
forman una sucesión exacta no escindida
\[
 0\longrightarrow C_{12}\xrightarrow{\iota}C_{108}
 \xrightarrow{p}C_9\longrightarrow0.
\]
Para la sección conjuntista que elige los representantes
\(0,\ldots,8\), el defecto de aditividad es el acarreo
\[
 c(r,r')=\left[\left\lfloor\frac{\widetilde r+\widetilde r'}9\right\rfloor\right]_{12}.
\]
\end{proposition}

\begin{proof}
El núcleo de \(p\) está formado por los múltiplos de nueve y coincide con
la imagen de \(\iota\). La reducción es sobreyectiva y la multiplicación
por nueve es inyectiva en el dominio indicado. Si la sucesión se
escindiera, el grupo central sería isomorfo a \(C_{12}\times C_9\), cuyo
exponente es \(\operatorname{mcm}(12,9)=36\); el exponente de
\(C_{108}\) es 108. La contradicción prueba la no escisión. Finalmente,
la división euclídea de \(\widetilde r+\widetilde r'\) por nueve da el
acarreo escrito y verifica la identidad de la sección.
\end{proof}

El resultado explica una diferencia que reaparecerá en alfa: regresar a la
fase no equivale a regresar al estado. Si \(q_{\mathrm{mem}}\) cuenta las
vueltas nonádicas, el calendario conserva solamente
\(\beta=[q_{\mathrm{mem}}]_{12}\). Después de ciento ocho pasos,
\(\beta\) retorna, pero \(q_{\mathrm{mem}}\) aumenta en doce. El
registro conserva además el cilindro, el registro de incidencias y la memoria
vertical. El cierre de una carta finita no convierte esa historia en un
ciclo sin memoria.

\subsection{El registro orientado y la carta integral \texorpdfstring{\(1+5+6\)}{1+5+6}}
\label{subsec:k-registro-integral}

Para una historia producida por TPK, sea \(\tau_m\) la subruta sobre
\(E_m\), \(\sigma_m\) su orientación, \(\kappa_m\) el acarreo de
frontera y \(\Lambda_m\) el artículo local de incidencias. El registro es
\[
 \mathcal R_{12}(x)=
 \bigl((E_m,\tau_m,\sigma_m,\kappa_m,\Lambda_m)\bigr)_{m=1}^{12}.
\]
Su dominio es el espacio de historias aritméticas con memoria de trayectoria, no el conjunto de
etiquetas de fase. La proyección al calendario olvida precisamente algunos
de los datos que el lector firmado necesita.

Una carta inicial del registro puede escribirse mediante eventos
orientados. La extracción orientada fija el conjunto de códigos de evento
\(\{2,4,6,8\}\) y la orientación sectorial
\[
 \Sigma_{12}=(+,+,+,-,-,-,+,+,+,-,-,-).
\]
Si \(d_t\) es el código de dirección en el paso \(t\), se obtiene
\[
 e_i=\Sigma_{12}(i+1)
 \#\{t\in E_{i+1}:d_t\in\{2,4,6,8\}\},
 \qquad0\leq i<12.
\]
La fórmula mantiene separados el censo no negativo de eventos y la
orientación. No identifica una dirección cardinal con un signo TRIT. El
vector \(e\), cuyos componentes son censos locales firmados, tampoco es
todavía el vector \(U\) que aparecerá más abajo.

Las posiciones opuestas determinan, para \(0\leq i<6\),
\[
 p_i=e_i+e_{i+6},\qquad a_i=e_i-e_{i+6}.
\]
Con
\[
 s_C=\sum_{i=0}^{5}p_i,\qquad M_i=p_i-p_5\quad(0\leq i<5),
\]
se define la carta integral
\[
 \mathcal L_{12}(e)
 =(s_C,M_0,\ldots,M_4,a_0,\ldots,a_5).
\]
La descomposición \(1+5+6\) no es un reparto de coordenadas arbitrario:
una coordenada conserva la carga, cinco comparan los pares respecto de un
par de referencia y seis conservan la oposición de las semicoronas.

\begin{lemma}[Inversión de la carta integral]
\label{lem:k-carta-integral}
Una tupla entera \((s_C,M_0,\ldots,M_4,a_0,\ldots,a_5)\) pertenece a la
imagen de \(\mathcal L_{12}\) si y sólo si
\[
 s_C-\sum_{i=0}^{4}M_i\equiv0\pmod6,
 \qquad p_i\equiv a_i\pmod2\quad(0\leq i<6),
\]
donde
\[
 p_5=\frac{s_C-\sum_{i=0}^{4}M_i}{6},\qquad p_i=p_5+M_i\quad(i<5).
\]
En ese caso su preimagen es única y viene dada por
\[
 e_i=\frac{p_i+a_i}{2},\qquad e_{i+6}=\frac{p_i-a_i}{2}.
\]
\end{lemma}

\begin{proof}
La suma de los cinco márgenes es \(s_C-6p_5\), de donde se obtiene la
división por seis. Reconstruidos los \(p_i\), cada par de ecuaciones
\(p_i=e_i+e_{i+6}\), \(a_i=e_i-e_{i+6}\) tiene la solución indicada.
Esta solución es entera exactamente bajo la congruencia de paridad. Todas
las operaciones se invierten, por lo que no queda libertad residual.
\end{proof}

Esta reversibilidad muestra qué información se conserva; no autoriza a
suprimir el resto del artículo. El registro completo tiene un dominio mayor
que esta carta de censos. La extracción de los canales que alimentan a
\(U\) utiliza ese registro completo y no se deduce de la lista
\((e_i)\) por una identificación de nombres.

\input{sections/registro_incidencias.tex}
\input{sections/censo_terminal_completo.tex}

\subsection{Dos lecturas del estado terminal}
\label{subsec:k-terminal}

Denotemos por \(x_{\mathrm{term}}\) el estado producido por la composición
de selección, transporte y actualización hasta el corte terminal. La
escritura operatoria es
\[
 x_{\mathrm{term}}=
 (\mathcal U_{108}\circ\cdots\circ\mathcal U_1)(x_{\mathrm{seed}}),
 \qquad \mathcal U_t=\mathsf{Upd}_t\circ\mathsf{Tra}_t\circ\mathsf{Sel}_t.
\]
Sus dos lecturas pertinentes son
\[
 x_{\mathrm{term}}
 \xrightarrow{(\pi_{\mathrm{term}},R_{\mathrm{sgn}})}
 (B_{\mathrm{term}},U).
\]
La matriz \(B_{\mathrm{term}}\) retiene la lectura ternaria visible; el
vector \(U\) conserva otra coordenada de la misma historia. No se inserta
entre ellos una flecha \(B_{\mathrm{term}}\to U\).

El selector terminal opera sobre los tres catálogos regionales
de tamaños 196, 197 y 196, construidos en la
sección~\ref{censo:terminal-completo}. Las firmas locales reducen sus
\(196\cdot197\cdot196=7\,567\,952\) ternas a 331, y el margen de
columnas \(q_{\partial}=(1,2,2,1,2,0)\) deja dos matrices:
\[
 B_0=\begin{pmatrix}0&2&0&2&2&0\\0&0&1&2&2&0\\1&0&1&0&1&0\end{pmatrix},
 \qquad
 B_1=\begin{pmatrix}0&2&1&0&2&0\\0&0&1&2&2&0\\1&0&0&2&1&0\end{pmatrix}.
\]
La orientación de filas es \(\sigma=(1,1,-1)=(1,1,2)\) en
\(\mathbb F_3^3\). Las cargas de filas de las candidatas son
\((0,2,0)\) y \((2,2,1)\). Sólo la segunda pertenece a
\(\langle\sigma\rangle=\{(0,0,0),(1,1,2),(2,2,1)\}\); por tanto
\(B_{\mathrm{term}}=B_1\). Esta comprobación prueba el último paso del
selector a partir del censo declarado. No convierte el margen de columnas
en una consecuencia de la carga de filas, ni reemplaza la enumeración
anterior por la inspección de esas dos matrices.

El lector firmado tiene la composición tipada
\begin{equation}
 R_{\mathrm{sgn}}
 =\Pi_H\circ\mathcal E_{108}^{90,120}\circ\mathcal R_{12}:
 \mathcal X_{\mathrm{TPK}}^{\mathrm{term,mem}}
 \longrightarrow\mathbb Z^{12}.
 \label{eq:k-lector-firmado}
\end{equation}
El registro de canales determina
\[
 \mathcal E_{108}^{90,120}(\mathcal R_{12}(x_{\mathrm{term}}))
 =(b^{90},b^{120},Q_{\mathrm{TPK}}),
\]
con las coordenadas
\begin{align*}
 b^{90}={}&(-495,-116,-684,108,601,-90,\\
            &\hspace{19mm}-173,-88,313,560,-397,461),\\
 b^{120}={}&(-425,-281,-481,671,-255,30,\\
             &\hspace{19mm}475,-543,680,251,6,-128),\\
 Q_{\mathrm{TPK}}={}&6263.
\end{align*}
Estas cantidades se utilizan como coordenadas del registro previo, no
como parámetros que se ajusten mediante alfa. La reconstrucción demostrada
en los apartados siguientes comienza exactamente en esta salida tipada.
El rastro material del extractor anterior y el alcance de su transcripción
autónoma se conservan en la nota técnica de procedencia; la aritmética
posterior no se presenta como sustituto de ese rastro.

\subsection{Del registro de canales al vector firmado}
\label{subsec:k-pi-h}

El lector armónico \(\Pi_H\) puede escribirse sin utilizar ni \(K\) ni
alfa. Para evitar confundir sus sumas orbitales con los bloques de alfa,
las denotaremos por \(s_1,s_2,s_3\). Sean
\[
 B_r=\sum_{j=0}^{3}b^{120}_{r+3j},\qquad
 d_{r,j}=b^{90}_{r+3j},\qquad1\leq r\leq3,
\]
y
\begin{equation}
 s_1=\frac{Q_{\mathrm{TPK}}+2B_1+B_2}{3},\qquad
 s_2=s_1-B_1,\qquad s_3=s_2-B_2.
 \label{eq:k-sumas-armonicas}
\end{equation}
La división por tres pertenece al lector de las tres órbitas. Su
integralidad se comprueba sobre las salidas compatibles del registro; no
se supone para un elemento arbitrario de \(\mathbb Z^{25}\).

Los tres bloques firmados son
\begin{equation}
 U^{(r)}=(s_r,d_{r,1}-d_{r,3},d_{r,0}+d_{r,2},d_{r,0}-d_{r,2}).
 \label{eq:k-bloques-firmados}
\end{equation}
La sustitución entera da
\[
 (B_1,B_2,B_3)=(972,-1073,101),\qquad
 (s_1,s_2,s_3)=(2378,1406,2479),
\]
y, por tanto,
\begin{align*}
 U^{(1)}&=(2378,-452,-668,-322),\\
 U^{(2)}&=(1406,998,-204,-28),\\
 U^{(3)}&=(2479,-551,-371,-997).
\end{align*}
Intercalando por columnas se obtiene
\begin{equation}
 U=(2378,1406,2479,-452,998,-551,
       -668,-204,-371,-322,-28,-997).
 \label{eq:k-u-firmado}
\end{equation}
La presencia de valores negativos y de componentes superiores a 999
muestra que \(U\) no es una palabra de dígitos en base mil. Tampoco es
\(U_6\Vert U_6\), concatenación catalogal de período seis. La diferencia
es de dominio y de información conservada, no solamente de valores.

El orden constructivo hasta aquí es
\[
 x_{\mathrm{term}}\longrightarrow\mathcal R_{12}
 \longrightarrow(b^{90},b^{120},Q_{\mathrm{TPK}})
 \longrightarrow U.
\]
La inversión que sigue produce el registro dodecafásico. Que después podamos recuperar
los canales desde el registro dodecafásico constituye un control de reversibilidad; no
invierte este orden causal.

\subsection{Hadamard por órbitas y reconstrucción entera}
\label{subsec:k-hadamard}

En el ciclo de doce posiciones, el paso tres determina tres órbitas:
\[
 (1,4,7,10),\qquad(2,5,8,11),\qquad(3,6,9,12).
\]
La separación de cuatro sectores de cada órbita es la que, una vez fijados
origen y orientación, admite la lectura angular de noventa grados. La
transformación no actúa sobre tres cuaternas contiguas del orden natural.

Sea
\[
 H_4=\begin{pmatrix}
 1&1&1&1\\1&1&-1&-1\\1&-1&1&-1\\1&-1&-1&1
 \end{pmatrix}.
\]
Si \(P_{\mathrm{orb}}\) reordena por las tres órbitas indicadas,
definimos
\[
 \mathcal H_{12}=P_{\mathrm{orb}}^{-1}
 (H_4\oplus H_4\oplus H_4)P_{\mathrm{orb}}.
\]

\begin{lemma}[Inversión de Hadamard e integralidad]
\label{lem:k-hadamard}
Se tiene
\[
 H_4^2=4I_4,\qquad H_4^{-1}=\frac14H_4,\qquad\det H_4=-16.
\]
Para \(u\in\mathbb Z^4\), la preimagen \(H_4^{-1}u\) es entera si y
sólo si \(H_4u\equiv0\pmod4\).
\end{lemma}

\begin{proof}
Las filas son ortogonales y tienen norma cuadrada cuatro; la matriz es
simétrica. Esto da \(H_4^2=4I_4\) y la inversa. Una eliminación entera,
o la expansión del determinante, da el signo negativo y el valor \(-16\).
La condición de integralidad es exactamente la divisibilidad por cuatro
de las cuatro coordenadas de \(H_4u\).
\end{proof}

\begin{definition}[Transformación integral dodecafásica]
La subred de registros firmados y su coordenatización integral son
\[
 \mathcal L_H=\{u\in\ZZ^{12}:\mathcal H_{12}u\equiv0\pmod4\},
 \qquad \mathscr K(u)=\tfrac14\mathcal H_{12}u.
\]
La aplicación \(\mathscr K:\mathcal L_H\to\ZZ^{12}\) es un isomorfismo
de grupos abelianos, con inversa \(k\mapsto\mathcal H_{12}k\).
El registro canónico \(K\) es la imagen del registro firmado generado
por las operaciones anteriores, con origen de fase y orientación fijados.
\end{definition}

\begin{proof}
La congruencia define exactamente la integralidad de \(\mathscr K(u)\).
La identidad \(\mathcal H_{12}^2=4I\) demuestra las dos composiciones
inversas. En particular, todo \(k\in\ZZ^{12}\) tiene preimagen
\(\mathcal H_{12}k\in\mathcal L_H\).
\end{proof}

\begin{proposition}[Registro integral canónico]
\label{prop:k-sello}
La inversión \(K^{(r)}=H_4U^{(r)}/4\) de los bloques de
\eqref{eq:k-bloques-firmados} produce
\begin{align*}
 K^{(1)}&=(234,729,621,794),\\
 K^{(2)}&=(543,659,58,146),\\
 K^{(3)}&=(140,824,914,601).
\end{align*}
En el orden natural, el registro dodecafásico es
\begin{equation}
 K=(234,543,140,729,659,824,621,058,914,794,146,601).
 \label{eq:k-vector}
\end{equation}
Es la única preimagen racional de \(U\), y pertenece a
\(\{0,\ldots,999\}^{12}\).
\end{proposition}

\begin{proof}
Los tres productos antes de dividir son
\begin{align*}
 H_4U^{(1)}&=(936,2916,2484,3176),\\
 H_4U^{(2)}&=(2172,2636,232,584),\\
 H_4U^{(3)}&=(560,3296,3656,2404).
\end{align*}
Todas las coordenadas son divisibles por cuatro. Dividir y deshacer la
permutación de órbitas da \eqref{eq:k-vector}. La unicidad procede de la
invertibilidad racional, no de una búsqueda entre candidatos próximos a
un valor físico.
\end{proof}

La integralidad también puede formularse en términos de redes. Los
divisores determinantes de \(H_4\) son \(1,2,4,16\): las entradas
tienen máximo común divisor uno; los menores de orden dos son pares y
alguno tiene módulo dos; los de orden tres son múltiplos de cuatro y
alguno tiene módulo cuatro. De ahí se deduce
\[
 \operatorname{SNF}(H_4)=\operatorname{diag}(1,2,2,4),
\]
y para la transformación global,
\[
 \operatorname{coker}\mathcal H_{12}
 \cong(\mathbb Z/2\mathbb Z)^6\oplus(\mathbb Z/4\mathbb Z)^3.
\]
El índice de la imagen es \(16^3=4096\). Así, no todo vector firmado
entero admite una preimagen entera; el vector obtenido satisface las
congruencias precisas que la inversión exige. El factor \(1/4\) está
forzado por la matriz, y no tiene el estatuto de un coeficiente libre.

\subsection{Diferencias cíclicas y componente uniforme}
\label{subsec:k-transversal}

Con índices cíclicos módulo doce, sean
\[
 (D_3z)_m=z_m-z_{m+3},\qquad(D_4z)_m=z_m-z_{m+4},
 \qquad Q(z)=\sum_{m=1}^{12}z_m.
\]
El primer operador compara sectores separados por tres posiciones; el
segundo, por cuatro. Los nombres de canal 90 y 120 se refieren aquí a
esas separaciones en el ciclo orientado, no a una identificación con un
operador físico de otro dominio.

\begin{proposition}[Completitud del sistema transversal]
\label{prop:k-transversal}
El operador
\[
 \mathcal D=(D_3,D_4,Q):\mathbb Z^{12}\longrightarrow\mathbb Z^{25}
\]
tiene rango doce y factores invariantes no nulos
\[
 \operatorname{diag}(1,\ldots,1,12),
\]
con once unidades. En particular, sus salidas determinan un único registro dodecafásico.
\end{proposition}

\begin{proof}
Si \(D_3z=D_4z=0\), el vector es constante a lo largo de ambos pasos.
Como \(3\) y \(4\) generan el ciclo módulo doce,
\(\ker D_3\cap\ker D_4=\langle\mathbf1\rangle\). La carga total
elimina esta última libertad porque \(Q(\mathbf1)=12\).

Para la afirmación integral, las filas de \((D_3,D_4)\) son incidencias
orientadas de un grafo conexo. Un árbol generador proporciona once
factores invariantes unitarios. Todo menor de orden doce debe contener la
fila \(Q\). Sumando las primeras once columnas a la última, esta última
se anula en las filas de incidencia y vale doce en la fila \(Q\). Todos
los menores máximos son divisibles por doce y el árbol produce uno de
módulo exactamente doce. El último factor es, por tanto, doce.
\end{proof}

La comprobación sobre \eqref{eq:k-vector} da
\[
 D_3K=b^{90},\qquad D_4K=b^{120},\qquad Q(K)=6263.
\]
Las fórmulas \eqref{eq:k-sumas-armonicas} y
\eqref{eq:k-bloques-firmados} reconstruyen a su vez \(U\) desde esas
diferencias. Se cierra así un triángulo de representaciones exactas del
mismo objeto: registro transversal, vector firmado y registro dodecafásico. La igualdad
de resultados no convierte sus codominios en uno solo; demuestra que
las aplicaciones declaradas recuperan la información pertinente.

\input{sections/registro_imagen_integral.tex}

\subsection{Dos lecturas racionales: \texorpdfstring{\(\kappa_{36}\)}{kappa36} y \texorpdfstring{\(\kappa_{\mathrm{per}}\)}{kappa periódica}}
\label{subsec:k-kappas}

Una vez construido \(K\), la carta en base mil permite dos operaciones
distintas. La primera lee una sola vuelta; la segunda prolonga el registro dodecafásico
periódicamente. Sea \(B=1000\) y definamos su entero de concatenación
\begin{equation}
 N_K=\sum_{m=1}^{12}K_mB^{12-m}
 =234543140729659824621058914794146601.
 \label{eq:k-numerador}
\end{equation}
Los ceros iniciales de un bloque, como el de 058, son parte de su escritura
de tres cifras. No cambian su valor entero, pero sí permiten concatenar
sin ambigüedad los doce bloques.

\begin{definition}[Lecturas finita y periódica del registro dodecafásico]
\label{def:k-kappas}
La lectura de una vuelta es
\[
 \kappa_{36}=\sum_{m=1}^{12}K_mB^{-m}=\frac{N_K}{B^{12}}.
\]
La prolongación periódica es
\[
 K_j^{\mathrm{per}}=K_{1+((j-1)\bmod12)},\qquad
 \kappa_{\mathrm{per}}=\sum_{j\geq1}K_j^{\mathrm{per}}B^{-j}.
\]
\end{definition}

\begin{proposition}[Valor y período exactos del registro dodecafásico prolongado]
\label{prop:k-periodo}
Se cumple
\begin{equation}
 \kappa_{\mathrm{per}}=\frac{N_K}{B^{12}-1},\qquad
 \kappa_{\mathrm{per}}-\kappa_{36}
 =\frac{N_K}{B^{12}(B^{12}-1)}
 =B^{-12}\kappa_{\mathrm{per}}.
 \label{eq:k-diferencia-kappas}
\end{equation}
La palabra infinita del registro dodecafásico tiene período mínimo doce en base mil y
período mínimo treinta y seis en base diez.
\end{proposition}

\begin{proof}
La suma de la primera vuelta es \(N_KB^{-12}\); cada nueva vuelta
multiplica su contribución por \(B^{-12}\). La serie geométrica produce
\(N_K/(B^{12}-1)\), y la resta da la segunda identidad.

Los doce componentes de \(K\) son distintos. Un período propio del ciclo
identificaría dos de ellos; por tanto el período mínimo de bloques es
doce. La expansión decimal es canónica, pues el registro dodecafásico no tiene una cola
de bloques 999. Si su período mínimo decimal fuera \(d\), dividiría
36. Agrupar de tres en tres produciría un período de bloques
\(d/\gcd(d,3)\). Ningún divisor propio de 36 da doce mediante esa
operación. Luego el período decimal mínimo es 36.
\end{proof}

Ambas lecturas son racionales, pero no son el mismo racional. La primera
termina después de doce bloques; la segunda repite la vuelta. Además,
\(0<\kappa_{\mathrm{per}}-\kappa_{36}<B^{-12}\). Su cercanía no
permite sustituir una por otra en una identidad exacta. En particular,
el cierre finito de alfa utiliza \(\kappa_{36}\), mientras que la
suma de la sección prolongada utiliza \(\kappa_{\mathrm{per}}\).

\input{sections/k_reversibilidad.tex}

\subsection{Periodicidad y conúcleo del operador de transporte}
\label{subsec:k-clase-acarreo}

La periodicidad que acabamos de demostrar pertenece al registro dodecafásico. No implica
periodicidad de las coordenadas regionales, de los acarreos de la clausura
ni de alfa. Es útil precisar esta separación mediante un segundo
cociente, puramente integral.

Sea \(S\) el desplazamiento cíclico de doce coordenadas, con orientación
fijada, y sea \(b\geq2\) una base entera. El operador de acarreo
periódico es \(\partial_b=S-bI\). Si una identidad periódica se escribe
\[
 A=P+E-\Phi-K+\partial_bC,
\]
entonces las transformaciones
\[
 K\longmapsto K+\partial_bh,\qquad C\longmapsto C+h
\]
conservan su lado derecho. Esta invariancia compara representantes de una
clase de acarreo; no es una autorización para modificar el registro dodecafásico ya
seleccionado por el registro.

\begin{proposition}[Cociente del acarreo periódico]
\label{prop:k-cociente-acarreo}
Con la orientación anterior,
\[
 \operatorname{coker}(S-bI)\cong\mathbb Z/(b^{12}-1)\mathbb Z.
\]
\end{proposition}

\begin{proof}
El módulo cíclico se presenta como
\(\mathbb Z[t]/(t^{12}-1)\). Imponer \(S=bI\) equivale a imponer
\(t=b\), y deja la relación \(b^{12}-1=0\). La evaluación de los
coeficientes en \(b\), con el orden de concatenación elegido, representa
la clase resultante.
\end{proof}

El denominador \(B^{12}-1\) de \(\kappa_{\mathrm{per}}\) y este
cociente proceden de la misma longitud cíclica, pero sus objetos siguen
siendo distintos: uno es una evaluación racional y el otro un cociente
de un módulo de acarreos. La clausura finita de alfa será una cadena
orientada con frontera derecha. Su prolongación transportará una
frontera remota compatible, no un acarreo que se identifique por decreto
entre la primera y la decimotercera posición.

Queda así preparado el capítulo siguiente. El registro dodecafásico se ha fijado antes
de resolver la recurrencia de alfa; su reconstrucción por Hadamard es
entera y única; sus diferencias transversales recuperan el registro; y
sus dos evaluaciones racionales tienen dominios y usos separados. Ninguno
de estos hechos utiliza alfa como entrada, y ninguno decide todavía el
tipo aritmético de la salida de la clausura completa.


% SOURCE foundation/sections/revision_k.tex
% Adición propuesta al inicio de registro_k.tex, después de su introducción.
% Conserva todas las definiciones, fórmulas y demostraciones actuales.
% Procedencia: registro_k.tex, §§k-hadamard/k-transversal/k-kappas;
% k_reversibilidad.tex; incidencia_articulacion.tex, inc-ampl-cuaterna.
\paragraph{Coordenatización reversible y determinación incidencial.}
\label{ampl-alfa-k-intro-reversible}
El resultado que organiza este capítulo es la equivalencia entre tres
coordenatizaciones de un registro previamente producido por
APP--TRIT--TPK. La transformación de Hadamard relaciona el registro
firmado con el registro entero; el extractor de diferencias relaciona
este último con sus variaciones transversales y su componente uniforme:
\[
 U\underset{\mathcal H_{12}}{\overset{\mathcal H_{12}/4}{\longleftrightarrow}}K
 \overset{(D_3,D_4,Q)}{\longleftrightarrow}
 (D_3K,D_4K,Q(K)).
\]
La primera equivalencia se establece sobre la subred integral
$\mathcal L_H$; la segunda, sobre la imagen del extractor. Fijadas la
base $1000$, las $12$ posiciones y el origen de fase, la evaluación
$K\mapsto\kappa_{\mathrm{per}}$ añade una cuarta coordenatización
reversible: multiplicar por $1000^{12}-1$ y efectuar las divisiones
euclídeas recupera el registro completo. Las demostraciones de
\S\ref{subsec:k-hadamard}, \S\ref{subsec:k-transversal} y
\S\ref{subsec:k-kappas} establecen estas equivalencias en sus respectivos
dominios.

Sobre ese registro actúan después dos composiciones. La composición
posicional combina sus componentes con las coordenadas de clausura,
propagación y autoescala, y su normalización determina $\alpha_A$. La
composición incidencial aplica el umbral cúbico y obtiene
\[
 K\longmapsto B_K=\{j:K_j\geq729\}=\{4,6,9,10\}.
\]
El registro íntegro permite reconstruir ambas composiciones. El soporte
$B_K$ conserva, específicamente, la pertenencia a la clase seleccionada
por el umbral. La identidad $B_K=H_e\cap P_\pi$ de
\eqref{inc-ampl-cuaterna} enlaza esta extracción entera con la incidencia
de los códigos regionales. El soporte negativo de su balance con $K$
completa la bandera $B_K\subset H^-_\alpha$ que se desarrolla en la
prolongación excepcional. Así, la función de $K$ queda definida mediante
una transformación invertible y mediante aplicaciones posteriores
precisas; la información adicional de profundidad y transporte permanece
en el estado del que procede el registro.


% SOURCE foundation/sections/registro_incidencias.tex
% Suplemento de trabajo. No sustituye registro_k.tex ni cierra por decreto E108.
% Procedencia y alcance: INFORME_RECUPERACION.md, misma carpeta.
\subsection{Evaluación incidencial anterior al registro dodecafásico}
\label{subsec:registro-evaluacion-incidencial}

La recuperación de un registro a partir de sus diferencias cíclicas y su
carga total es una cuestión distinta de la evaluación de esas coordenadas
sobre el registro de incidencias. El primer problema se resuelve mediante el
extractor transversal de la proposición~\ref{prop:k-transversal}.
Para el segundo, una presentación histórica del
registro proporciona tres recuentos enteros en cada ventana. Es útil
explicitar su composición y determinar exactamente la información que
requiere su evaluación.

Sea \(\mathscr L\) un libro de visitas y trazas orientadas previamente
generado. Cada visita conserva la ruta, la posición APP, la hoja, el
instante, la multiplicidad y la incidencia de frontera. La evaluación
aritmética produce el entero y su descomposición
\[
 n_t=r_t+9q_t,\qquad 1\leq r_t\leq9.
\]
El entero \(n_t\) se calcula a partir de la hoja de suma o producto; su
residuo no se recibe como un valor independiente. La distinción entre
incidencia de corona e incidencia interior permanece explícita para las
visitas de residuo \(9\).

En la ventana \(E_m\), la presentación por recuentos considera
\begin{align}
 A_m&=\sum_{v\in E_m}w(v)
             \mathbf1_{\{1,3,5,7\}}(r(v)),\\
 C_m&=\sum_{j\geq0}\bigl(\nu^-_{120,m}(j)-\nu^+_{120,m}(j)\bigr),\\
 V_m&=\sum_{v\in E_m}w(v)\mathbf1_{\{9\}}(r(v))\epsilon_\partial(v).
\end{align}
Aquí \(w(v)\) es la multiplicidad incidencial y
\(\epsilon_\partial(v)=1\) en la corona, \(-1\) en el interior.
Las trazas contadas por \(\nu^\pm\) deben enumerarse mediante la regla
correspondiente del artículo. La suma entera requiere soporte finito o una
construcción alternativa que determine su significado. El número de
incidencias cronológicas y la longitud reducida \(120(1+3j)\) tienen
tipos distintos; su relación exige la unidad de longitud del lector.

Se define \([x]_{10}\in\{0,\ldots,9\}\) como el representante
euclídeo y se conservan los cocientes
\[
 x=[x]_{10}+10\lfloor x/10\rfloor.
\]
El bloque incidencial es
\begin{equation}
 z_m=100[A_m]_{10}+10[C_m]_{10}+[V_m]_{10}.
 \label{eq:registro-bloque-incidencial}
\end{equation}
Por composición se obtiene la aplicación
\begin{equation}
 \mathcal E_{\rm cnt}(\mathscr L)=
 \left((z_m-z_{m+3})_{m=1}^{12},
       (z_m-z_{m+4})_{m=1}^{12},\sum_{m=1}^{12}z_m\right).
\label{eq:registro-composicion-incidencial}
\end{equation}
Los índices se leen cíclicamente módulo \(12\). En este apartado,
\((D_jz)_m=z_m-z_{m+j}\), para \(j=3,4\), y
\(Q(z)=\sum_{m=1}^{12}z_m\); por tanto, la composición anterior es
\((D_3z,D_4z,Q(z))\). Esta convención se conserva en
\S\ref{subsec:k-transversal}, donde se demuestra la reconstrucción
integral completa.
Esta fórmula produce sus coordenadas a partir de los recuentos. No utiliza
\(K\), \(U\), \(\alpha\) ni las coordenadas transversales esperadas
como argumentos. Su identificación con
\(\mathcal E_{108}^{90,120}\) requiere componerla con la familia y
las incidencias efectivas del estado terminal; no se deduce únicamente
de la invertibilidad del sistema transversal.

\begin{proposition}[Información aritmética suficiente y necesaria]
\label{prop:registro-36-residuos}
Sean \(N,N'\in\mathbb Z^{12\times3}\), ordenados por los recuentos
\((A,C,V)\). Para la aplicación definida por
\eqref{eq:registro-bloque-incidencial}--\eqref{eq:registro-composicion-incidencial},
\[
 \mathcal E_{\rm cnt}(N)=\mathcal E_{\rm cnt}(N')
 \quad\Longleftrightarrow\quad
 N-N'\in10\mathbb Z^{36}.
\]
En particular, los \(36\) residuos sectoriales determinan exactamente
la salida de \(25\) coordenadas. La conservación de los cocientes y
de la procedencia constituye información adicional del registro.
\end{proposition}

\begin{proof}
La igualdad módulo \(10\) produce los mismos bloques y, por tanto,
las mismas coordenadas transversales. Recíprocamente, si las coordenadas
transversales coinciden, \(D_3(z-z')=D_4(z-z')=0\). Los pasos
\(3\) y \(4\) generan el ciclo de longitud \(12\), de modo que
\(z-z'\) es constante. La igualdad de cargas totales anula esa
constante. Finalmente, la representación decimal de un entero entre
\(0\) y \(999\) tiene tres dígitos únicos. Se recuperan así los
tres residuos de cada ventana. La afirmación es una identidad algebraica
para todo \(N,N'\), no una inferencia obtenida de pruebas finitas.
\end{proof}

\subsection{Conjugación de los recuentos y términos de frontera}
\label{subsec:registro-conjugacion-incidencial}

Sea \(\rho(m)=13-m\). Considérese una conjugación de incidencias
que invierte el orden de las ventanas, conserva la clase aritmética y la
clase de frontera de las visitas e intercambia los canales de las trazas.
Entonces
\[
 (A_m^\dagger,C_m^\dagger,V_m^\dagger)
 =(A_{\rho(m)},-C_{\rho(m)},V_{\rho(m)}).
\]
Si \(c_m=[C_m]_{10}\), la reducción decimal produce la identidad
\begin{equation}
 z_m^\dagger=z_{\rho(m)}
       +100\mathbf1_{\{c_{\rho(m)}\ne0\}}-20c_{\rho(m)}.
 \label{eq:registro-conjugacion-decimal}
\end{equation}
En efecto, \([-C]_{10}=0\) cuando \([C]_{10}=0\), y
\([-C]_{10}=10-[C]_{10}\) en otro caso. La fórmula conserva
exactamente esa diferencia. La negación del canal no equivale a negar
el bloque decimal completo.

La aplicación de esta identidad a una involución concreta del TPK exige
verificar su acción sobre las incidencias. Para un lector de ocupaciones
evaluado antes del avance, la inversión de una ruta
\(\gamma=(x_0,\ldots,x_n)\) satisface
\[
 \sum_{k=0}^{n-1}f(x_{n-k})
 =\sum_{k=0}^{n-1}f(x_k)+f(x_n)-f(x_0).
\]
Los términos de los extremos se incorporan antes de reducir módulo
\(10\). En rutas cerradas se anulan; en ventanas abiertas pueden
alterar sus residuos. Esta corrección, ya presente en el desarrollo
bidireccional del transporte, especifica qué debe conservar el empalme
entre la ruta y los recuentos.

\begin{proposition}[Necesidad de la información no periódica del registro]
Supóngase que el observable de una familia fija de rutas satisface
\(o(t+54)=o(t)\), que la multiplicidad se conserva bajo ese retorno
y que el mismo lector local actúa en las ventanas de \(9\) instantes.
Entonces sus recuentos y bloques satisfacen
\[
 (A,C,V)_{m+6}=(A,C,V)_m,\qquad z_{m+6}=z_m.
\]
\end{proposition}

\begin{proof}
La traslación \(t\mapsto t+54\) lleva \(E_m\) biyectivamente
a \(E_{m+6}\), conserva cada incidencia contada y su multiplicidad.
La igualdad se transmite a la suma y a la reducción decimal.
\end{proof}

Por tanto, una presentación con \(12\) bloques distintos requiere que
el lector conserve alguna información que no factorice por ese observable
de período \(54\): selección de incidencias, orientación efectiva,
extremos, multiplicidad o memoria. La proposición delimita una reducción
inadmisible del estado; no afirma que la dinámica completa tenga
período \(54\) ni determina por sí misma cuál es su lector terminal.


% SOURCE foundation/sections/censo_terminal_completo.tex
\subsection{Censo completo de cilindros y selección de la frontera terminal}
\label{censo:terminal-completo}

La lectura terminal utiliza conjuntamente dos informaciones: la compatibilidad
de los cilindros regionales y un perfil de incidencia con canales y posiciones
etiquetados. En este apartado se construyen los catálogos completos, se
enumeran sus ternas y se demuestra el efecto de cada condición del lector.
El punto de partida son las salidas regionales ya producidas en la
sección~\ref{sec:generacion}, no valores introducidos para definir el
transporte APP--TRIT--TPK. La reducción a una matriz pertenece a una lectura
posterior de esas salidas y conserva el orden de los tres canales.

El resultado es un teorema censal con estructura de partida explícita.
El perfil de Gram, pesos, distancias y márgenes se declara como dato del
lector terminal. La exhaustividad de la enumeración demuestra qué selecciona
ese perfil sobre los catálogos fijados; no demuestra, por sí sola, que el
perfil sea una ley universal ni que pueda reconstruirse desde la única terna
que finalmente lo satisface. Esta separación permite incorporar la prueba
finita completa sin invertir su genealogía.

\subsubsection{De las salidas regionales a los catálogos ternarios}

Sean \(x_\pi,x_e,x_\varphi\in[0,1)\) las partes fraccionarias de las
tres coordenadas regionales ya construidas. Se utilizan sus cilindros
publicados de mil posiciones decimales,
\[
 J_\chi^{(1000)}=
 \left[\frac{d_\chi}{10^{1000}},
       \frac{d_\chi+1}{10^{1000}}\right),
 \qquad \chi\in\{\pi,e,\varphi\}.
\]
Los enteros \(d_\chi\) son lecturas finitas de salida. La longitud mil
es suficiente para los enteros que siguen, como se comprueba mediante las
dos cotas de cada intervalo; no es una profundidad supuesta infinita.
Definimos
\begin{equation}
 p=30,\quad r=1050,\quad L=p+r=1080,\quad k=171,
 \qquad D=3^L,\quad Q=1000^k,
 \label{censo:parametros}
\end{equation}
y
\[
 P_\chi=\lfloor3^{30}x_\chi\rfloor,\qquad
 T_\chi=\lfloor Qx_\chi\rfloor.
\]
Aquí \(k\) cuenta tríadas decimales y no designa el registro dodecafásico
\(K\). Se tiene \(1000^{171}\leq3^{1080}<1000^{172}\).

\begin{lemma}[Extracción estable de un entero]
\label{censo:estabilidad}
Para \(d,m,h\in\mathbb Z\), con \(m,h>0\), el valor de
\(\lfloor mx\rfloor\) es constante en \([d/h,(d+1)/h)\) si y sólo si
\begin{equation}
 \left\lfloor\frac{md}{h}\right\rfloor
 =\left\lfloor\frac{m(d+1)-1}{h}\right\rfloor.
 \label{censo:prueba-estabilidad}
\end{equation}
Cuando se cumple, cualquiera de los dos miembros da ese valor.
\end{lemma}
\begin{proof}
El mínimo de \(mx\) se alcanza en el extremo izquierdo. Su supremo
es \(m(d+1)/h\) y no se alcanza. El mayor entero que puede tomar
\(\lfloor mx\rfloor\) es
\(\lceil m(d+1)/h\rceil-1\), igual al segundo miembro porque
el numerador y el denominador son enteros. La igualdad de los extremos
del intervalo de valores prueba la afirmación.
\end{proof}

La prueba de estabilidad se cumple en los tres canales para
\(m=3^{30},Q,3^{1080}\). La última elección se utilizará únicamente
para una comparación posterior, no para seleccionar el perfil. Las
primeras treinta posiciones recuperadas son precisamente las cinco
etapas de seis trits del transporte regional:
\begin{center}\small
\begin{tabular}{@{}c l r@{}}\toprule
Canal&Prefijo de treinta trits&\(P_\chi\)\\\midrule
\(\pi\)&\texttt{010211012222010211002111110221}&29152671743887\\
\(e\)&\texttt{201101121221102011012222102011}&147887858824447\\
\(\varphi\)&\texttt{121200112202121020010210010200}&127247717616687\\\bottomrule
\end{tabular}
\end{center}
La coincidencia conserva la prolongación ya construida; la lectura decimal
no se usa para sustituir retroactivamente sus matrices de transporte.

Fijado un canal, una cola de \(r\) trits se representa por
\(0\leq R<3^r\). Su cilindro y el cilindro de las primeras \(k\)
tríadas son
\begin{equation}
 I_{\chi,R}=\left[\frac{P_\chi3^r+R}{D},
                       \frac{P_\chi3^r+R+1}{D}\right),\qquad
 J_{\chi,T}=\left[\frac{T_\chi}{Q},\frac{T_\chi+1}{Q}\right).
 \label{censo:dos-cilindros}
\end{equation}
El catálogo incluye todas las colas para las que estos dos cilindros
tienen intersección no vacía.

\begin{proposition}[Catálogo por intersección exacta]
\label{censo:catalogos}
Las colas compatibles son los enteros del intervalo \([a_\chi,b_\chi]\),
con
\begin{align}
 a_\chi&=\max\left(0,
       \left\lfloor\frac{T_\chi D}{Q}\right\rfloor-P_\chi3^r\right),
       \label{censo:extremo-a}\\
 b_\chi&=\min\left(3^r-1,
       \left\lfloor\frac{(T_\chi+1)D-1}{Q}\right\rfloor-P_\chi3^r\right).
       \label{censo:extremo-b}
\end{align}
Para los cilindros regionales fijados, los recortes no actúan y se obtiene
\begin{equation}
\begin{array}{c|r|r|r}
 \chi&a_\chi\bmod729&b_\chi\bmod729&b_\chi-a_\chi+1\\\hline
 \pi&67&262&196\\
 e&584&51&197\\
 \varphi&117&312&196
\end{array}
\label{censo:tabla-catalogos}
\end{equation}
\end{proposition}
\begin{proof}
Escribamos \(N=P_\chi3^r+R\). Dos intervalos semiabiertos de
\eqref{censo:dos-cilindros} se cortan exactamente cuando
\[
 NQ<(T_\chi+1)D,\qquad (N+1)Q>T_\chi D.
\]
La segunda desigualdad equivale a
\(N\geq\lfloor T_\chi D/Q\rfloor\); la primera, a
\(N\leq\lfloor((T_\chi+1)D-1)/Q\rfloor\). Restar
\(P_\chi3^r\) y conservar \(0\leq R<3^r\) prueba las fórmulas.
La evaluación por divisiones enteras de los \(T_\chi\) extraídos por
\eqref{censo:prueba-estabilidad} da la tabla. Como \(r\geq6\),
\(P_\chi3^r\equiv0\pmod{729}\), de modo que el residuo del primer
extremo también se obtiene directamente de
\(\lfloor T_\chi D/Q\rfloor\bmod729\).
\end{proof}

Los últimos seis trits de una cola forman la palabra
\(\nu_6(R\bmod729)\), donde \(\nu_6\) es la escritura ternaria
de longitud seis, incluidos sus ceros iniciales. Por la tabla, cada
catálogo tiene menos de \(729\) elementos consecutivos, de manera que
esta reducción es inyectiva sobre él. Así, los tres catálogos de frontera
quedan descritos completamente, sin elegir las ternas finales:
\begin{equation}
 \begin{split}
 \mathcal B_\pi&=\{\nu_6(67+i):0\leq i<196\},\\
 \mathcal B_e&=\{\nu_6((584+j)\bmod729):0\leq j<197\},\\
 \mathcal B_\varphi&=\{\nu_6(117+\ell):0\leq\ell<196\}.
 \end{split}
 \label{censo:fronteras-explicitas}
\end{equation}
Exigir que el extremo izquierdo de \(I_{\chi,R}\) pertenezca a
\(J_{\chi,T}\) sería otra condición: perdería el primer cilindro
que corta la frontera y daría \(195,196,195\). El censo utiliza la
intersección de los cilindros completos, de acuerdo con su definición.

\subsubsection{Perfil declarado e indicadores de compatibilidad}

Para \(u,v\in\mathbb F_3^6\), sean
\(g(u,v)=\sum u_iv_i\in\mathbb F_3\),
\(n(u)=g(u,u)\), \(w(u)=\#\{i:u_i\ne0\}\) y
\(h(u,v)=\#\{i:u_i\ne v_i\}\). Las dos últimas funciones toman
valores enteros; las dos primeras, valores en \(\mathbb F_3\).
El perfil local del lector, en orden \((\pi,e,\varphi)\), es
\begin{equation}
 \begin{split}
 (g_{\pi e},g_{\pi\varphi},g_{e\varphi})&=(2,2,0),\qquad
 (n_\pi,n_e,n_\varphi)=(0,0,0),\\
 (w_\pi,w_e,w_\varphi)&=(3,3,3),\qquad
 (h_{\pi e},h_{\pi\varphi},h_{e\varphi})=(2,5,3).
 \end{split}
 \label{censo:perfil-local}
\end{equation}
Las palabras se ordenan como en \eqref{censo:fronteras-explicitas}.
Usaremos \([E]\in\{0,1\}\) para el indicador de una condición \(E\).
Para cada par de canales \(\chi,\psi\), la matriz rectangular
\(A_{\chi\psi}\) tiene entrada \([g(u,v)=g_{\chi\psi}]\).
La matriz \(A^{h}_{\chi\psi}\) multiplica esa entrada por
\([h(u,v)=h_{\chi\psi}]\). Sean \(D_\chi^n\) la matriz diagonal
de \([n(u)=0]\), y \(D_\chi^{nw}\) la de
\([n(u)=0][w(u)=3]\).

\begin{theorem}[Enumeración exhaustiva por sumas finitas]
\label{censo:teorema-conteo}
El producto de los tres catálogos contiene \(7\,567\,952\) ternas.
Las condiciones sucesivas de \eqref{censo:perfil-local} dejan
\begin{equation}
 7\,567\,952\longrightarrow275\,794\longrightarrow6235
 \longrightarrow3127\longrightarrow331.
 \label{censo:cadena-local}
\end{equation}
Estos cuatro conteos son, respectivamente,
\begin{align}
 N_g&=\operatorname{tr}(A_{\pi e}A_{e\varphi}A_{\varphi\pi}),
 \nonumber\\
 N_{gn}&=\operatorname{tr}(D_\pi^n A_{\pi e}
             D_e^n A_{e\varphi}D_\varphi^n A_{\varphi\pi}),
 \nonumber\\
 N_{gnw}&=\operatorname{tr}(D_\pi^{nw} A_{\pi e}
             D_e^{nw} A_{e\varphi}D_\varphi^{nw} A_{\varphi\pi}),
 \nonumber\\
 N_{gnwh}&=\operatorname{tr}(D_\pi^{nw} A^h_{\pi e}
             D_e^{nw} A^h_{e\varphi}D_\varphi^{nw} A^h_{\varphi\pi}).
 \label{censo:trazas}
\end{align}
Todos los productos y trazas de esta fórmula se calculan en los enteros,
no en el cuerpo de tres elementos.
\end{theorem}
\begin{proof}
El tamaño del producto es \(196\cdot197\cdot196\). Al expandir
una traza de \eqref{censo:trazas}, sus índices recorren exactamente una
terna de palabras. El sumando vale uno si satisface todas las condiciones
representadas y cero en otro caso. Cada terna se cuenta una vez.
Las matrices diagonales añaden las condiciones individuales; los factores
rectangulares añaden las condiciones por pares. Esto demuestra que las
trazas son precisamente los cardinales anunciados.

Para explicitar su evaluación, fijados los índices \(i,j\) de los dos
primeros canales, se forman los conjuntos de índices \(\ell\) del tercero
que cumplen cada condición por pares. Se intersectan sus indicadores,
se añaden los indicadores individuales pertinentes y se suma su cardinal.
La suma final recorre \(0\leq i<196\), \(0\leq j<197\).
La sustitución de las palabras de \eqref{censo:fronteras-explicitas} y
del perfil \eqref{censo:perfil-local} da, en ese orden,
\(275794,6235,3127,331\). El suplemento registra las cuatro contribuciones
de cada uno de los \(196\) índices \(i\) y las \(331\) ternas completas;
su programa evalúa estas sumas con indicadores binarios exactos y contrasta
escalarmente cada superviviente. No interviene una muestra ni una tolerancia.
\end{proof}

\subsubsection{Incidencia etiquetada y carga de orientación}

Para una terna superviviente, sea \(B\in\mathbb F_3^{3\times6}\) su
matriz de filas ordenadas. El segundo perfil exige
\begin{equation}
 \operatorname{rank}_{\mathbb F_3}B=3,\qquad
 w_{\mathrm{col}}(B)=(1,1,2,2,3,0),\qquad
 q_{\partial}(B)=(1,2,2,1,2,0),
 \label{censo:perfil-columnas}
\end{equation}
donde \(w_{\mathrm{col}}\) cuenta los elementos no nulos de cada
columna y \(q_{\partial}\) suma sus entradas módulo tres. Es un perfil
de columnas etiquetadas; permutarlas sin transportar el perfil define otro
lector. Denotemos por \(\mathcal F\) el conjunto de las \(331\) ternas
del teorema anterior.

\begin{proposition}[Reducción del censo al par terminal]
\label{censo:dos-matrices}
Las sumas de indicadores sobre \(\mathcal F\), al exigir sucesivamente
rango, pesos de columna y sumas de columna, son
\begin{equation}
 \sum_{B\in\mathcal F}[\operatorname{rank}B=3]=331,
 \qquad N_{\mathrm{rango,pesos}}=18,
 \qquad N_{\mathrm{rango,pesos,sumas}}=2.
 \label{censo:conteo-columnas}
\end{equation}
Con índices de catálogo comenzando en uno, las dos ternas finales son
\((120,197,157)\) y \((129,197,148)\), y sus matrices son
\[
 B_0=\begin{pmatrix}0&2&0&2&2&0\\0&0&1&2&2&0\\1&0&1&0&1&0\end{pmatrix},
 \qquad
 B_1=\begin{pmatrix}0&2&1&0&2&0\\0&0&1&2&2&0\\1&0&0&2&1&0\end{pmatrix}.
\]
\end{proposition}
\begin{proof}
En cada terna del conjunto ya enumerado se calcula el rango por eliminación
en \(\mathbb F_3\), se cuentan los elementos no nulos de las seis
columnas y se suman éstas módulo tres. Los productos de los indicadores
de \eqref{censo:perfil-columnas} dan \eqref{censo:conteo-columnas}.
La sustitución de los dos índices finales en
\eqref{censo:fronteras-explicitas} da las matrices expuestas. La lista
completa de supervivientes y los filtros exactos del suplemento permiten
reproducir la eliminación sin presuponer esas dos matrices como dominio.
\end{proof}

La orientación \(\sigma=(1,1,-1)=(1,1,2)\in\mathbb F_3^3\)
selecciona la recta
\(\langle\sigma\rangle=\{(0,0,0),(1,1,2),(2,2,1)\}\).
Las cargas de filas de \(B_0,B_1\) son \((0,2,0)\) y \((2,2,1)\);
sólo la segunda pertenece a esa recta. Se recupera así el último paso
del selector de la sección~\ref{subsec:k-terminal}, ahora precedido por
la construcción y enumeración de su dominio completo. La lectura estable
de \(\lfloor3^{1080}x_\chi\rfloor\), efectuada después de los filtros,
da las fronteras \texttt{021020}, \texttt{001220}, \texttt{100210},
con los mismos rangos \((129,197,148)\). Este cotejo posterior no define
el perfil ni participa en las sumas de indicadores.

\subsubsection{Información conservada y alcance del resultado}

El procedimiento conserva el prefijo de treinta trits, la cola compatible,
el orden regional, la posición de cada trit y la orientación. La reducción
módulo \(729\) conserva aquí la identidad de cada candidato porque el
intervalo de colas tiene longitud menor que \(729\); esa inyectividad
no se afirma para una historia arbitraria. El censo de \(331\) demuestra
también que el perfil local aislado no selecciona una sola matriz: los
márgenes etiquetados y la carga aportan condiciones efectivas adicionales.

Se ha recuperado, por tanto, una enumeración finita íntegra y su selector
sobre cilindros de salida, con perfil de lectura explícito. La construcción
del registro firmado utiliza además su lector de memoria y sus canales,
como se especifica a continuación; la matriz visible no se identifica con
ese registro por igualdad de número de componentes. Tampoco este conteo
demuestra la universalidad del perfil, genera de nuevo las coordenadas
regionales ni decide una afirmación sobre alfa o sobre el continuo completo.


% SOURCE foundation/sections/registro_imagen_integral.tex
% Formalización nueva integrada en la revisión III; originales preservados.
% Propietarios: U016, ecuaciones extractor-diferencias-k y lector-armonico-*;
% registro_k.tex, subsecciones k-pi-h, k-hadamard y k-transversal.
\subsection{Compatibilidad integral y determinación del registro transversal}
\label{img09:seccion}

El dominio de esta construcción es la publicación del registro de una
historia APP--TRIT--TPK. Se distinguen la evaluación de los eventos del
registro, la compatibilidad de sus canales y la inversión de sus
coordenadas. El resultado siguiente da un criterio integral completo para
el segundo nivel y una factorización explícita para conectar el primero con él.
La demostración utiliza los operadores del registro y establece sus
relaciones para todo elemento de los dominios enteros indicados.

\subsubsection{Caracterización por incrementos adyacentes}

Los índices pertenecen a $\mathbb Z/12\mathbb Z$. Sean
\[
 (Sz)_i=z_{i+1},\qquad D_3=I-S^3,\qquad D_4=I-S^4,
 \qquad Q(z)=\sum_{i=0}^{11}z_i.
\]
Se conserva el extractor transversal
\[
 \mathcal D:\mathbb Z^{12}\longrightarrow\mathbb Z^{25},
 \qquad k\longmapsto(D_3k,D_4k,Q(k)).
\]
Para una terna de coordenadas enteras $(b,c,q)$, escribimos
\begin{equation}
 w_i=c_i-b_{i+1},\qquad
 P_0=0,\quad P_i=\sum_{j=0}^{i-1}w_j\ (1\leq i\leq12),
 \qquad T(w)=\sum_{i=0}^{11}P_i.
 \label{img09:incremento}
\end{equation}

\begin{theorem}[Imagen integral del extractor transversal]
\label{img09:imagen}
Una terna $(b,c,q)\in\mathbb Z^{25}$ pertenece a
$\operatorname{im}\mathcal D$ si y sólo si satisface
\begin{align}
 b_i&=w_i+w_{i+1}+w_{i+2}&& (0\leq i<12),
 \label{img09:relaciones}\\
 \sum_{i=0}^{11}w_i&=0,
 \label{img09:cierre}\\
 q+T(w)&\equiv0\pmod {12}.
 \label{img09:congruencia}
\end{align}
Su única preimagen entera es
\begin{equation}
 k_i=\frac{q+T(w)}{12}-P_i,\qquad0\leq i<12.
 \label{img09:inversa}
\end{equation}
Las doce ecuaciones de \eqref{img09:relaciones} y la ecuación
\eqref{img09:cierre} son linealmente independientes sobre $\mathbb Q$.
\end{theorem}
\begin{proof}
Para una salida de $\mathcal D$ se tiene
\[
 c_i-b_{i+1}=(k_i-k_{i+4})-(k_{i+1}-k_{i+4})=k_i-k_{i+1}.
\]
La suma de tres incrementos da $b_i$; la suma cíclica de todos ellos
es cero. Además, $P_i=k_0-k_i$, por lo que
$q=12k_0-T(w)$. Se obtienen todas las condiciones y la fórmula inversa.

Recíprocamente, \eqref{img09:congruencia} hace entera la fórmula
\eqref{img09:inversa}; \eqref{img09:cierre} permite prolongarla
cíclicamente. Sus diferencias adyacentes son $w_i$. Por tanto,
$D_3k=b$ y
\[
 (D_4k)_i=w_i+w_{i+1}+w_{i+2}+w_{i+3}
 =w_i+b_{i+1}=c_i.
\]
La suma de la inversa da $Q(k)=q$. La misma fórmula prueba la unicidad.
Para la independencia, el cambio entero invertible
$(b,c)\mapsto(b,w=c-Sb)$ convierte las doce primeras relaciones en
$b=(I+S+S^2)w$. Cada una despeja una coordenada diferente de $b$.
La condición $\sum w_i=0$ es una relación adicional independiente.
\end{proof}

La imagen racional tiene dimensión doce. Su estructura integral añade
una congruencia de orden doce: ésta es la manifestación constructiva del
último factor de Smith de $\mathcal D$. Los once incrementos
$w_0,\ldots,w_{10}$ y la carga $q$ son coordenadas suficientes, con
$w_{11}=-\sum_{i=0}^{10}w_i$ y la congruencia
\eqref{img09:congruencia}. Las veinticinco coordenadas originales
conservan esos mismos doce grados de libertad mediante trece relaciones.

\subsubsection{Conmutatividad con la lectura armónica}

Sea $\mathcal H_{12}$ la transformación de Hadamard del registro,
con órbitas $(r,r+3,r+6,r+9)$ para $r=0,1,2$. Su bloque es
\[
 H_4=\begin{pmatrix}
 1&1&1&1\\1&1&-1&-1\\1&-1&1&-1\\1&-1&-1&1
 \end{pmatrix},\qquad H_4^2=4I_4.
\]
Para $(b,c,q)$ compatible, se definen
\[
 B_r=\sum_{j=0}^3c_{r+3j},\qquad
 A_0=\frac{q+2B_0+B_1}{3},\quad A_1=A_0-B_0,
 \quad A_2=A_1-B_1.
\]
El lector $\Pi_H$ de \eqref{eq:k-bloques-firmados} tiene los bloques
\[
 (\Pi_H(b,c,q))^{(r)}=
 (A_r,b_{r+3}-b_{r+9},b_r+b_{r+6},b_r-b_{r+6}).
\]

\begin{proposition}[Concordancia de las dos reconstrucciones]
\label{img09:triangulo}
Sobre $\operatorname{im}\mathcal D$ se cumple
\[
 \Pi_H\mathcal D=\mathcal H_{12},\qquad
 \mathcal D^{-1}=\tfrac14\mathcal H_{12}\Pi_H.
\]
La imagen de $\Pi_H$ restringida a ese dominio es exactamente
\[
 \mathcal L_H=
 \{u\in\mathbb Z^{12}:\mathcal H_{12}u\equiv0\pmod4\}.
\]
\end{proposition}
\begin{proof}
Para $k=\mathcal D^{-1}(b,c,q)$, sea
$a_r=\sum_{j=0}^3k_{r+3j}$. El desplazamiento de cuatro posiciones
lleva cada órbita a la siguiente, y por ello $B_r=a_r-a_{r+1}$.
Junto con $q=a_0+a_1+a_2$, estas identidades dan $A_r=a_r$.
En una órbita, escribamos $(x_0,x_1,x_2,x_3)$ para las coordenadas
de $k$ y $d_j=x_j-x_{j+1}$. Entonces
\[
 d_1-d_3=x_0+x_1-x_2-x_3,\quad
 d_0+d_2=x_0-x_1+x_2-x_3,\quad
 d_0-d_2=x_0-x_1-x_2+x_3.
\]
Son las tres filas restantes de $H_4k^{(r)}$. La inversión y la
caracterización integral se siguen de $\mathcal H_{12}^2=4I$.
\end{proof}

La condición $\Pi_H(b,c,q)\in\mathcal L_H$ sólo controla la
salida armónica. La pertenencia de las veinticinco coordenadas a
$\operatorname{im}\mathcal D$ exige además
\eqref{img09:relaciones}--\eqref{img09:congruencia}.
Por ejemplo, añadir a $c$ el vector $e_0-e_3$ conserva las tres sumas
orbitales y deja $\Pi_H$ inalterado. Partiendo de $(0,0,0)$ se
obtiene así una terna con $\Pi_H=0$ que incumple
\eqref{img09:relaciones}. Éste es un control negativo exacto de la
compatibilidad del canal, independiente de toda evaluación terminal.

\subsubsection{Determinación y transporte sobre las órbitas del registro}

El Teorema~\ref{img09:imagen} determina $k$ cuando se han producido
sus canales y su carga. Como condición sobre salidas aún no
individualizadas, la misma imagen es un grupo isomorfo a
$\mathbb Z^{12}$. Fijada únicamente una carga $q$, las soluciones
forman una clase afín de
\[
 L_0=\{v\in\mathbb Z^{12}:\sum v_i=0\}
 =\bigoplus_{i=0}^{10}\mathbb Z(e_i-e_{11}).
\]
Por ejemplo, $100\mathbf1+e_0$ y $100\mathbf1+e_1$ son dos
registros distintos, de carga $1201$, con componentes en
$\{0,\ldots,999\}$. Cada uno satisface todas las condiciones de
imagen y todas las congruencias de Hadamard mediante sus propios
canales. Estos testigos comparan las ecuaciones de compatibilidad;
no se les atribuye la condición de historias terminales HMT.

\begin{proposition}[Covariancia cíclica y estabilizador del registro]
\label{img09:covariancia}
Sean $\rho$ la traslación de ventanas en un dominio de registros
$\mathscr R$ y $S$ la traslación anterior en $\mathbb Z^{12}$.
Sobre la órbita de $R\in\mathscr R$, un lector covariante queda
determinado por un vector
\[
 k\in\operatorname{Fix}(S^p),
\]
donde $p\mid12$ es el período mínimo de $R$ bajo $\rho$.
La condición de carga $Q(k)=q$ admite soluciones enteras si y sólo si
$12/p$ divide a $q$; cuando existen, forman una clase afín de rango
$p-1$. En particular, sobre una órbita libre, covariancia y carga
dejan once grados de libertad.
\end{proposition}
\begin{proof}
La regla $F(\rho^jR)=S^jk$ está bien definida exactamente cuando
$S^pk=k$. Esos vectores repiten $12/p$ veces una palabra de
longitud $p$. Su carga es $\frac{12}{p}\sum_{i=0}^{p-1}k_i$.
La divisibilidad y el rango se siguen de esta fórmula. La
covariancia de los canales resulta de $D_aS=SD_a$ y $QS=Q$.
Sobre $\mathcal L_H$, la acción correspondiente es
$\widehat S_H=\mathcal H_{12}S\mathcal H_{12}/4$;
por la Proposición~\ref{img09:triangulo}, ambas reconstrucciones
entrelazan estas acciones.
\end{proof}

Si el dominio lleva un origen marcado, debe utilizarse la acción que
transporta también esa marca. La proposición se refiere a esa acción
cíclica declarada; no presupone la periodicidad de la historia
enriquecida. La naturalidad frente a truncamientos posteriores al corte
$108$ se obtiene, para cualquier lector $F_{108}$ ya definido, mediante
$F_N=F_{108}\circ\tau_{N,108}$. Esta naturalidad conserva la regla
de evaluación inicial; por sí sola no elige una de las reglas posibles
sobre los primeros ciento ocho eventos.

\subsubsection{Evaluación por eventos y composición de contribuciones}

La caracterización anterior proporciona una factorización de la extracción mediante dos
lecturas efectivas del registro:
\begin{equation}
 \omega:\mathscr R\longrightarrow
 \{w\in\mathbb Z^{12}:\sum w_i=0\},\qquad
 q_{\rm reg}:\mathscr R\longrightarrow\mathbb Z,
 \quad q_{\rm reg}(R)+T(\omega(R))\equiv0\pmod {12}.
 \label{img09:contrato-productor}
\end{equation}
Estas aplicaciones deben evaluarse sobre los eventos, la orientación,
la incidencia y la memoria que produjo APP--TRIT--TPK. Una vez dadas sus
acciones sobre esos datos, la composición
\[
 R\longmapsto(w,q)\longmapsto
 ((I+S+S^2)w,(I+S+S^2+S^3)w,q)
 \xrightarrow{\Pi_H}u\xrightarrow{\mathcal H_{12}/4}k
\]
está determinada y satisface todos los controles inversos. La fórmula
$w=c-Sb$ también permite conectar un productor que ya entregue ambos
canales: es una comprobación sobre su salida, anterior a toda comparación
con el testigo terminal.

\begin{proposition}[Acumulación de contribuciones y unicidad prospectiva]
\label{img09:acumulacion}
Supóngase que el estado inicial y cada evento $a_t$ del registro
producen, mediante reglas fijadas previamente, contribuciones
$(\delta w_t,\delta q_t)$ con
\[
 \sum_i\delta w_{t,i}=0,\qquad
 \delta q_t+T(\delta w_t)\equiv0\pmod {12}.
\]
Con condiciones iniciales compatibles $(w^{(0)},q^{(0)})$, la
actualización
\[
 w^{(t+1)}=w^{(t)}+\delta w_t,\qquad
 q^{(t+1)}=q^{(t)}+\delta q_t
\]
produce un único registro en cada etapa. Su incremento es
\[
 \delta k_{t,i}=
 \frac{\delta q_t+T(\delta w_t)}{12}-P_i(\delta w_t).
\]
La construcción respeta la concatenación de registros, y la lectura
de un prefijo permanece igual bajo prolongaciones posteriores.
\end{proposition}
\begin{proof}
Las condiciones de compatibilidad se conservan por suma. La fórmula
\eqref{img09:inversa} es aditiva en $(w,q)$ y produce el
incremento indicado. La inducción fija sucesivamente cada estado.
La suma sobre dos segmentos es la suma concatenada, y un prefijo recibe
exactamente los mismos sumandos antes y después de prolongar la historia.
\end{proof}

La proposición proporciona un criterio suficiente, no una exigencia de
linealidad para todo lector TPK. Si los eventos se transportan por acciones
no triviales, sus contribuciones pueden expresarse en la carta terminal
mediante el cociclo afín correspondiente. La ecuación de cociclo conserva
entonces la misma independencia de la partición de la trayectoria.
En ambos casos, las reglas de evento y el estado inicial constituyen los
datos constructivos que han de provenir del artículo: escogerlos para
reproducir un vector esperado sustituiría esa procedencia por otra tarea.

El resultado de este desarrollo es la condición exacta de determinación
\eqref{img09:contrato-productor}, su inversión demostrada y sus identidades
de transporte. Su alcance es el registro transversal; las restantes
construcciones de HMT conservan sus dominios y sus pruebas propias.


% SOURCE foundation/sections/k_reversibilidad.tex
\subsection{Reversibilidad de la constante racional y cambio de fase}

La constante racional \(\kappa_{\mathrm{per}}\) determina íntegramente
el registro \(K\) cuando se conservan la base, la longitud y el origen
de lectura. Su función numérica y su función estructural se relacionan
mediante una inversión exacta.

\begin{proposition}[Recuperación integral y transformación cíclica]
Sea \(I(K)=\sum_{j=1}^{12}K_j1000^{12-j}\). Entonces
\[
 I(K)=(1000^{12}-1)\kappa_{\mathrm{per}}.
\]
La expansión entera de \(I(K)\) en doce bloques de base \(1000\)
recupera todos los \(K_j\), incluidos los ceros iniciales de cada bloque.
Si \(\rho(K)=(K_2,\ldots,K_{12},K_1)\), se cumple
\begin{equation}
 \kappa(\rho K)=1000\kappa(K)-K_1.
 \label{eq:k-rotacion-racional}
\end{equation}
\end{proposition}
\begin{proof}
La primera identidad invierte la definición de la evaluación periódica.
La división euclídea iterada por \(1000\), con doce posiciones fijadas,
recupera el vector. Por otra parte,
\[
 I(\rho K)=1000I(K)-K_1(1000^{12}-1).
\]
Dividir por \(1000^{12}-1\) da \eqref{eq:k-rotacion-racional}.
\end{proof}

La recuperación compone, por tanto, \(\kappa\mapsto K\mapsto U\)
y las transformaciones \(D_3K,D_4K,Q\) ya construidas. La evaluación
racional conserva el registro firmado y las diferencias cíclicas que
dependen de él. La profundidad total y los operadores de transporte de
una historia permanecen en sus coordenadas propias: la inversión anterior
tiene por dominio el registro integral dodecafásico.

Escribamos \(\kappa_j=\kappa(\rho^jK)\), con índices módulo \(12\).
La misma identidad da
\[
 K_{j+1}=1000\kappa_j-\kappa_{j+1},\qquad
 \sum_{j=0}^{11}\kappa_j=\frac{\sum_{j=1}^{12}K_j}{999}
 =\frac{6263}{999}.
\]
La fase actúa de manera exacta sobre la constante: su valor es fijo
en la lectura canónica, mientras sus rotaciones siguen una ley afín.
La suma de la órbita es invariante bajo la elección del origen.

También la incidencia posterior dispone de una expresión en estas
coordenadas. El umbral cúbico determina
\[
 B_K=\{j+1:1000\kappa_j-\kappa_{j+1}\ge729\}
 =\{4,6,9,10\}.
\]
Se recupera la tétrada a partir de la órbita racional porque ésta
recupera primero las componentes enteras. El soporte firmado asociado
a \(\alpha\) utiliza además sus registros regionales y la normalización
de frontera. La relación entre coordenada e incidencia queda así
expresada mediante operaciones identificables.

Esta es una definición operacional de la función de \(K\): proporciona
una coordenatización integral reversible del registro orientado, cuya
evaluación periódica y cuyas funciones de incidencia participan en
construcciones posteriores. El valor, la transformación y el dominio
se conservan conjuntamente.


% SOURCE foundation/sections/alpha.tex
\section{Derivación de la constante de estructura fina}
\label{sec:alpha}

La constante de estructura fina recibe en este desarrollo una determinación
aritmética que reúne las tres coordenadas regionales y el registro integral
dodecafásico. La operación combina sus bloques con orientación fijada y
transporta los cocientes de la normalización posicional. El parámetro
resultante expresa la compatibilidad de esa diferencia orientada con
la prolongación de los cilindros. Su valor y los datos de incidencia
utilizados en su construcción permanecerán relacionados durante el
desarrollo excepcional.

El orden interno es importante. Las coordenadas \(\pi\),
\(e\) y \(\varphi\), ya construidas en la sección precedente, son coordenadas
previas del mismo estado aritmético con memoria de trayectoria, y pueden actuar legítimamente como
entradas de un lector posterior. No son constantes convencionales
inyectadas para seleccionar alfa. La generación correlacionada de las cuatro coordenadas comparte origen y
compatibilidad coinductiva, sin que ello imponga evaluar sus cuatro
componentes en un solo paso.

Las dos vías expresan funciones matemáticas diferentes del mismo parámetro.
La primera realiza una clausura entera por división euclídea en base
\(1000\). La segunda estudia el equilibrio analítico de vacancias mediante
el resolvente \(E_{21/22}\) y sus prolongaciones graduadas. Una vía
determina bloques y cocientes; la otra determina una raíz mediante una
ecuación y un dominio de unicidad. Su identificación se formula después
en los cilindros comunes a toda profundidad. La exposición distingue
estos tres resultados: construcción posicional, determinación analítica
y compatibilidad de ambas.

\subsection{Primera vía: determinación posicional por clausura entera}
\subsubsection{Dominio y orientación de la clausura entera}
\label{subsec:alpha-dominio}

Sean \(P_j,E_j,\Phi_j\in\{0,\ldots,999\}\) los bloques canónicos de
las tres coordenadas regionales, y prolonguemos el registro dodecafásico por
\[
 K_j=K_{1+((j-1)\bmod12)}.
\]
La orientación del canal es \((1,1,-1)\). La contribución del registro dodecafásico se
sustrae según la clausura dodecafásica ya tipada. Para \(B=1000\), la
operación local distingue el balance firmado
\[
 Z_j:=P_j+E_j-\Phi_j-K_j
\]
de la suma que incorpora el cociente entrante. La normalización es
\begin{equation}
 s_j=Z_j+c_{j+1},\qquad
 A_j=s_j-B\left\lfloor\frac{s_j}{B}\right\rfloor,
 \qquad c_j=\left\lfloor\frac{s_j}{B}\right\rfloor.
 \label{eq:alpha-recurrencia}
\end{equation}
Así, \(A_j\in\{0,\ldots,B-1\}\), incluso cuando \(s_j\) es
negativo. El acarreo que entra desde la derecha pertenece al dominio de
la operación; no puede omitirse porque la tabla se imprima de izquierda
a derecha.

\begin{definition}[Clausura de profundidad finita]
\label{def:alpha-finita}
Para una profundidad \(N\) y una frontera entera \(t\), la aplicación
\(\mathcal C_N\) recibe
\((P_{1:N},E_{1:N},\Phi_{1:N},K_{1:N},c_{N+1}=t)\) y aplica
\eqref{eq:alpha-recurrencia} en el orden \(N,N-1,\ldots,1\). Su salida
es la pareja de palabras
\[
 (A_{1:N},c_{1:N+1}).
\]
\end{definition}

La división euclídea no posee una elección de signo residual: el resto se
fija en \(\{0,\ldots,999\}\) y el cociente queda determinado. Por
ejemplo, \(-431=569-1000\) y \(-1200=800-2\cdot1000\). Los
acarreos negativos de la clausura no son anomalías; son la información
entera necesaria para mantener la igualdad orientada.

\subsubsection{Ventana dodecafásica y condiciones de acarreo}
\label{subsec:alpha-ventana}

Las primeras doce tríadas obtenidas son
\begin{align*}
 P={}&(141,592,653,589,793,238,462,643,383,279,502,884),\\
 E={}&(718,281,828,459,045,235,360,287,471,352,662,497),\\
 \Phi={}&(618,033,988,749,894,848,204,586,834,365,638,117).
\end{align*}
En el canal de clausura, las cinco regiones que determinan la coordenada
\(P\) conservan sus antecedentes distintos. Compartir esta lectura no
las identifica como regiones o como historias.

Con frontera \(c_{13}=0\), la recurrencia produce la tabla siguiente.
Se muestra cada entrada, el acarreo entrante, el total antes de normalizar,
el bloque obtenido y el acarreo saliente. La tabla es un cálculo entero
reproducible, no una lista de bloques seleccionados por su parecido con
un objetivo.

\begin{table}[htbp]
\centering
\small
\setlength{\tabcolsep}{3.5pt}
\begin{tabular}{r|rrrr|rr|rr}
\hline
\(j\)&\(P_j\)&\(E_j\)&\(\Phi_j\)&\(K_j\)&\(c_{j+1}\)&\(s_j\)&\(A_j\)&\(c_j\)\\
\hline
1&141&718&618&234&0&7&007&0\\
2&592&281&033&543&0&297&297&0\\
3&653&828&988&140&\(-1\)&352&352&0\\
4&589&459&749&729&\(-1\)&\(-431\)&569&\(-1\)\\
5&793&045&894&659&\(-2\)&\(-717\)&283&\(-1\)\\
6&238&235&848&824&\(-1\)&\(-1200\)&800&\(-2\)\\
7&462&360&204&621&0&\(-3\)&997&\(-1\)\\
8&643&287&586&058&\(-1\)&285&285&0\\
9&383&471&834&914&\(-1\)&\(-895\)&105&\(-1\)\\
10&279&352&365&794&0&\(-528\)&472&\(-1\)\\
11&502&662&638&146&0&380&380&0\\
12&884&497&117&601&0&663&663&0\\
\hline
\end{tabular}
\caption{Clausura inicial con frontera derecha nula. La dependencia se propaga de la última fila a la primera.}
\label{tab:alpha-ventana}
\end{table}

La palabra de acarreos es
\[
 (c_1,\ldots,c_{13})=(0,0,0,-1,-1,-2,-1,0,-1,-1,0,0,0),
\]
y la lectura de la ventana es el racional
\begin{equation}
 \alpha_{12}^{(0)}
 =\frac{7297352569283800997285105472380663}{10^{36}}
 =0.007297352569283800997285105472380663.
 \label{eq:alpha-ventana-cero}
\end{equation}
El superíndice recuerda la frontera elegida. Esta ventana es exacta en su
dominio finito; no se la identifica sin más con el truncamiento de la
sección infinita. La diferencia entre ambas lecturas se hará visible al
transportar el acarreo de la cola remota.

\subsubsection{Telescopado y reversibilidad en una fibra de frontera}
\label{subsec:alpha-telescopia}

La forma local de la igualdad es
\begin{equation}
 P_j+E_j-\Phi_j-K_j-A_j=Bc_j-c_{j+1}.
 \label{eq:alpha-identidad-local}
\end{equation}
El término de acarreo es un enlace entre posiciones contiguas. Si se lo
elimina coordenada a coordenada, se obtiene una igualdad falsa. Sólo una
evaluación que conserve sus extremos permite telescoparlo.

\begin{proposition}[Identidad entera de profundidad arbitraria]
\label{prop:alpha-telescopia}
Para \(\iota_N(z)=\sum_{j=1}^{N}z_jB^{N-j}\), toda ejecución de
\(\mathcal C_N\) satisface
\begin{equation}
 \iota_N(P)+\iota_N(E)-\iota_N(\Phi)-\iota_N(K)-\iota_N(A)
 =B^Nc_1-c_{N+1}.
 \label{eq:alpha-telescopia}
\end{equation}
\end{proposition}

\begin{proof}
Multiplicar \eqref{eq:alpha-identidad-local} por \(B^{N-j}\) y sumar
produce en el lado derecho
\[
 \sum_{j=1}^{N}(Bc_j-c_{j+1})B^{N-j}=B^Nc_1-c_{N+1}.
\]
Cada acarreo interior aparece dos veces con signo opuesto; sólo sobreviven
los dos extremos.
\end{proof}

Para la tabla~\ref{tab:alpha-ventana}, ambos extremos son cero. Si
\(p_{12},e_{12},f_{12}\) son las tres sumas fraccionarias truncadas,
\[
 \alpha_{12}^{(0)}=p_{12}+e_{12}-f_{12}-\kappa_{36}.
\]
La constante del registro dodecafásico que aparece es \(\kappa_{36}\), no
\(\kappa_{\mathrm{per}}\). El dominio de esta identidad contiene
truncamientos y una frontera finita; su forma no autoriza a borrar ambas
condiciones al pasar a la sección completa.

La inversión local es también exacta:
\[
 K_j=P_j+E_j-\Phi_j+c_{j+1}-A_j-Bc_j.
\]
Dados \(P,E,\Phi\), los bloques \(A\) y los acarreos completos,
recupera el registro dodecafásico. La concatenación recupera además
\[
 \iota_N(K)=\iota_N(P)+\iota_N(E)-\iota_N(\Phi)-\iota_N(A)
             -B^Nc_1+c_{N+1}.
\]
Esta reversibilidad es una propiedad de la fibra con datos y fronteras
fijados. No debe convertirse en una nueva receta causal que empiece por
alfa y fabrique después el registro. La dirección de construcción sigue
siendo registro firmado, registro dodecafásico, clausura y salida.

\subsubsection{Prolongación y normalización canónica de la sección completa}
\label{subsec:alpha-completa}

La prolongación conserva el registro dodecafásico periódico y recibe las secuencias
regionales completas. Sean
\[
 p=\sum_{j\geq1}P_jB^{-j},\qquad
 e_0=\sum_{j\geq1}E_jB^{-j},\qquad
 f=\sum_{j\geq1}\Phi_jB^{-j}.
\]
Cada serie es una realización del carácter interno correspondiente.
Las partes enteras ya obtenidas son 3, 2 y 1, de modo que
\(p=\pi-3\), \(e_0=e-2\) y
\(f=\varphi-1\).

\begin{proposition}[Existencia y unicidad de la normalización completa]
\label{prop:alpha-normalizacion}
La serie
\begin{equation}
 y=p+e_0-f-\kappa_{\mathrm{per}}
   =\sum_{j\geq1}(P_j+E_j-\Phi_j-K_j)B^{-j}
 \label{eq:alpha-serie}
\end{equation}
converge absolutamente. En el dominio de los bloques obtenidos,
\(0<y<10^{-2}\). Existe una única expansión canónica \((A_j)\) de
\(y\), sin cola finalmente constante 999, y una única sucesión entera
acotada \((c_j)\) que satisface \eqref{eq:alpha-identidad-local} con
\(c_1=0\). Estos acarreos pertenecen al conjunto invariante
\(\mathcal C=\{-2,-1,0,1,2\}\).
\end{proposition}

\begin{proof}
Los coeficientes de la serie están acotados por \(2(B-1)\) en módulo,
de modo que domina una serie geométrica. Si
\(Z_j=P_j+E_j-\Phi_j-K_j\) y
\(R_j(Z)=\sum_{r\geq0}Z_{j+r}B^{-r-1}\), cada cola es la diferencia
de dos sumas de dos colas canónicas y pertenece a \((-2,2)\).
El primer coeficiente es
\(Z_1=141+718-618-234=7\); por tanto
\(y=(7+R_2(Z))/1000\in(0.005,0.009)\).

Sea \((A_j)\) la expansión canónica de \(y\). Definamos
\[
 c_j=R_j(Z)-R_j(A).
\]
Al multiplicar la igualdad de las dos series por \(B^{j-1}\) y separar
sus primeros \(j-1\) términos se obtiene
\[
 c_j=\sum_{r=1}^{j-1}(A_r-Z_r)B^{j-1-r}\in\mathbb Z,
\]
y \(c_1=0\). Además, \(R_j(A)\in[0,1)\), luego
\(-3<c_j<2\), lo que incluye todos los acarreos en \(\mathcal C\).
Las identidades de desplazamiento de las colas dan
\(Bc_j=Z_j-A_j+c_{j+1}\), que es la recurrencia requerida.

Recíprocamente, cualquier solución acotada con salida canónica y
\(c_1=0\) telescopa en el límite a la misma serie \(y\). La expansión
canónica es única y, una vez fijada, la fórmula de las colas determina
todos los acarreos. Finalmente, si se aplica una ejecución finita con
\(c_{j+1}\in\mathcal C\), se tiene \(-2000\leq s_j\leq2000\),
por lo que su cociente euclídeo vuelve a pertenecer a \(\mathcal C\).
\end{proof}

La proposición explicita la normalización del lector posterior: no
reemplaza la generación de las tres secuencias por una serie convencional
introducida como dato. Una vez producidas esas secuencias por el estado
coinductivo, su combinación y la frontera de la normalización quedan
determinadas sin consultar alfa.

El procedimiento finito conservado en el corpus propaga las cinco
fronteras de \(\mathcal C\) y conserva solamente el prefijo común a
todas ellas. Si dos profundidades han coalescido antes de una frontera,
la recurrencia determinista de derecha a izquierda da los mismos bloques
a su izquierda. Ésa es la compatibilidad con el truncamiento. No se
escoge una frontera porque dé cifras favorables, ni se deduce la
compatibilidad a toda profundidad de haber observado una corrida larga.

La comparación inicial muestra por qué esta precaución importa. En la
prolongación estabilizada se obtiene \(c_{13}^{\infty}=-1\), mientras
que la ventana aislada imponía \(c_{13}=0\). Así,
\[
 884+497-117-601-1=662
\]
en el último bloque de la primera vuelta. Las primeras once tríadas no
cambian, pero la prolongación comienza
\[
 (007,297,352,569,283,800,997,285,105,472,380,662,999,\ldots).
\]
Una aparición del bloque 999 no viola la convención canónica: lo
excluido es una cola que sea finalmente constante 999. La ventana con
frontera cero termina en 663; el truncamiento de la sección completa
termina en 662. Ambas afirmaciones describen operaciones diferentes y
exactas. No deben fusionarse en una sola notación de proyección.

\subsubsection{Definición operativa y ecuación aditiva completa}
\label{subsec:alpha-definicion}

\begin{definition}[Coordenada alfa de la clausura]
\label{def:alpha-operativa}
La vía entera determina la coordenada
\begin{equation}
 \alpha_A=\sum_{j\geq1}A_jB^{-j}
 =p+e_0-f-\kappa_{\mathrm{per}},
 \label{eq:alpha-via-a}
\end{equation}
donde \((A_j,c_j)\) es la normalización canónica compatible de
\eqref{eq:alpha-recurrencia}. En la realización escalar de las
coordenadas HMT,
\begin{equation}
 \boxed{\alpha_A=
 \pi+e-\varphi
 -4-\kappa_{\mathrm{per}}.}
 \label{eq:alpha-identidad-completa}
\end{equation}
\end{definition}

El entero cuatro no se calibra: es
\(3+2-1\), combinación de las partes enteras de las tres coordenadas.
La base mil procede de la carta cúbica y el período doce del registro dodecafásico. Los
signos son los del canal orientado y del cierre, y la frontera remota es
la que impone la sección compatible. Éstos son los orígenes efectivos de
los coeficientes de la ecuación.

En este nivel puede describirse alfa como la coordenada escalar del
defecto de clausura entre las tres realizaciones regionales y el retorno
racional dodecafásico, calculado con acarreo compatible. El término
\emph{defecto} designa una diferencia algebraica determinada, no una
pérdida arbitraria ni una imperfección del sistema. Esta descripción no
suprime la genealogía: sin las regiones, el registro firmado, la
orientación y la ley de acarreo, la igualdad escalar no sería una
exposición suficiente de la construcción.

También deben mantenerse separados tres objetos que la notación de alfa
puede acompañar en lectores posteriores: \(A\) es una palabra de
bloques; el soporte negativo de preacarreo se obtiene del signo de los
totales antes de normalizar; y un vector radial excepcional pertenece a
un codominio de incidencia. Ninguno de ellos es intercambiable con el
valor real \(\alpha_A\).

\input{sections/revision_alpha_notacion.tex}
\input{sections/revision_alpha_frontera.tex}

\input{sections/alpha_dependencias.tex}

\subsection{Segunda vía: determinación analítica mediante vacancias}
\subsubsection{Invariantes del transporte y resolvente graduado}
\label{subsec:alpha-via-b}

La vía analítica recibe invariantes del estado orientado, no la salida
\(\alpha_A\). Su truncamiento de orden seis es
\begin{equation}
 P_6(x)=2\pi x-\frac74x^2
 +\frac{x^3}{2\pi}
 +\frac{x^4}{20}-\frac{2x^5}{21}-\frac{x^6}{46},
 \label{eq:alpha-p6}
\end{equation}
y la condición de vacancia utiliza
\[
 \delta=\log_{10}\frac{10}{9}.
\]
La razón \(10/9\) compara la carta decimal con las nueve posiciones
APP de la construcción nonádica. Su logaritmo pertenece a esta lectura
posterior de compactificación; no es un valor objetivo de alfa.

La genealogía de los coeficientes debe conservarse completa, pero con su
fuerza lógica exacta. Los desarrollos recientes reúnen una firma marcada
del estado. La matriz de sincronización y el censo asociado son
\[
 D_{\mathrm{syn}}=\begin{pmatrix}85&112\\41&52\end{pmatrix},
 \qquad N_9=43.
\]
De su determinante se extrae
\[
 a=-\frac{\det D_{\mathrm{syn}}}{N_9}
   =\frac{172}{43}=4.
\]
El reloj inducido sobre los huecos cortos tiene espectro secundario
\(\{6,7\}\); su extremo superior da \(b=7\). La cardinalidad
trítica \(m=3\) determina
\[
 t_*=mb=21,\qquad k_*=mb-1=20.
\]
La lectura de vacancia en la carta mil da
\[
 v_-=\lfloor1000\delta\rfloor=45,\qquad
 v_+=\lceil1000\delta\rceil=46.
\]
Estos enteros son antecedentes tipados de la evaluación del truncamiento. No se
obtienen de su raíz ni de una aproximación metrológica a ella.

La rama inferior coincide además con el determinante del bloque APP
\[
 B_c=\begin{pmatrix}7&2\\2&7\end{pmatrix},\qquad\det B_c=45.
\]
El rango transversal del registro del capítulo anterior da
\(r_{\mathrm{dod}}=11\), y la semicorona de fase proporciona
\(f_{\mathrm{ph}}=8\binom62=120\). En particular,
\(45r_{\mathrm{dod}}=495\).

Con \(\omega=2\pi\), la firma
\[
 \mathcal I_9=(\omega,m,a,b,k_*,t_*,v_-,v_+,f_{\mathrm{ph}},r_{\mathrm{dod}})
\]
alimenta el lector graduado
\begin{align}
 \mathfrak E_9(\mathcal I_9)(x)
 ={}&\omega x-\frac ba x^2+\omega^{-1}x^3
 +\frac{x^4}{k_*}-\frac{2x^5}{t_*}-\frac{x^6}{v_+}\notag\\
 &-\frac{x^7}{f_{\mathrm{ph}}}+\frac{x^8}{v_-}
 +\frac{2x^9}{v_-r_{\mathrm{dod}}}.
 \label{eq:alpha-firma-coeficientes}
\end{align}
La sustitución recupera los coeficientes de \(P_6\) y de su
prolongación de orden nueve. El par \(45/46\) atraviesa la separación
entre ambos: la rama superior ocupa el grado seis y la inferior reaparece
en los grados ocho y nueve.

La factorización de \eqref{eq:alpha-firma-coeficientes} está formulada
para un estado que conserva origen, orden de canales y orientación. La
naturalidad respecto de morfismos que preservan esos datos no equivale a
la unicidad de todo lector después de olvidarlos. El dominio de esta
factorización está determinado por la firma
\eqref{eq:alpha-firma-coeficientes}; la prolongación en todos los grados
se construye en la sección~\ref{subsec:alpha-publicaciones-completas},
con su clase de recurrencia y su dependencia de la precoordenada del estado
expresamente definidas.

\subsubsection{Resolvente nilpotente en los grados 7, 8 y 9}
\label{subsec:alpha-resolvente}

La primera prolongación admite una construcción operatoria finita
especialmente explícita. Sobre el espacio con base \(e_7,e_8,e_9\),
sea
\[
 Ne_7=-\frac83e_8,\qquad Ne_8=\frac2{11}e_9,\qquad Ne_9=0,
 \qquad v_7=-\frac1{120}e_7.
\]
El símbolo \(N\) designa aquí un operador nilpotente, no una profundidad
de truncamiento. Su dominio y graduación están fijados antes de evaluar
el polinomio en un número.

\begin{lemma}[Prolongación nilpotente exacta]
\label{lem:alpha-resolvente}
Se cumple \(N^3=0\) y
\[
 (I-N)^{-1}v_7
 =-\frac{e_7}{120}+\frac{e_8}{45}+\frac{2e_9}{495}.
\]
La evaluación \(e_n\mapsto x^n\) produce
\begin{equation}
 T_{7:9}(x)=-\frac{x^7}{120}+\frac{x^8}{45}+\frac{2x^9}{495}.
 \label{eq:alpha-t79}
\end{equation}
\end{lemma}

\begin{proof}
El operador aumenta el grado y se anula en el último vector de la base,
por lo que su cubo es cero. La inversa es el resolvente finito
\(I+N+N^2\). Además,
\[
 Nv_7=\frac1{45}e_8,\qquad N^2v_7=\frac2{495}e_9.
\]
Sumar los tres términos y evaluar da la fórmula.
\end{proof}

Aquí los denominadores 120, 45 y 495 no son ajustes decimales: proceden
de la semicorona, del determinante APP y de su producto por el rango
dodecafásico. Las transferencias \(-8/3\) y \(2/11\), junto con el
signo y el grado de la semilla, pertenecen al lector tipado. Conservar
sólo la simetría del ciclo y sustituir este operador por otro operador
simétrico no conserva necesariamente los coeficientes.

\subsubsection{Raíz de multiplicidad uno del truncamiento de orden nueve}
\label{subsec:alpha-jet}

Definamos
\[
 P_9=P_6+T_{7:9},\qquad E_{21/22}^{(9)}(x)=P_9(x)-\delta,
 \qquad I_\alpha=(0,10^{-2}).
\]
Este operador finito tiene un resultado propio de existencia y unicidad.

\begin{proposition}[Raíz de multiplicidad uno del truncamiento]
\label{prop:alpha-raiz-jet}
El operador \(E_{21/22}^{(9)}\) tiene una única raíz
\(\xi_9\) en \(I_\alpha\), y esa raíz es simple. Se obtiene el aislamiento racional
\begin{equation}
 \frac{729735256928380099}{10^{20}}
 <\xi_9<\frac{729735256928380100}{10^{20}}.
 \label{eq:alpha-intervalo-jet}
\end{equation}
\end{proposition}

\begin{proof}
Las cotas de las coordenadas internas
\(3.14<\pi<3.15\) y
\(0.045<\delta<0.046\) son suficientes para la existencia. En cero,
\(E_{21/22}^{(9)}(0)=-\delta<0\). En \(10^{-2}\), la contribución
lineal menos los términos negativos es mayor que \(0.0626\), luego el
valor del operador es positivo.

Para \(0\leq x\leq10^{-2}\), al descartar los términos positivos de
la derivada se obtiene
\[
 (E_{21/22}^{(9)})'(x)
 \geq2\pi-\frac72\cdot10^{-2}
 -\frac{10}{21}\cdot10^{-8}-\frac6{46}\cdot10^{-10}
 -\frac7{120}\cdot10^{-12}>6.24.
\]
La continuidad da existencia; la cota derivativa da unicidad y
simplicidad. El intervalo más estrecho
\eqref{eq:alpha-intervalo-jet} resulta de una evaluación racional dirigida.
Sean \(x_-\) y \(x_+\) sus extremos. Para calcular
\(\delta=\log(10/9)/\log 10\), se emplea, para \(z>1\),
\[
 L_N(z)=2\sum_{j=0}^{N-1}\frac{y^{2j+1}}{2j+1},\qquad
 y=\frac{z-1}{z+1},\qquad
 0<\log z-L_N(z)<\frac{2y^{2N+1}}{(2N+1)(1-y^2)}.
\]
El resto se acota mediante la serie geométrica de los términos positivos.
Se aplica esta fórmula a \(z=10/9\) y \(z=3\), con \(N=520\),
usando \(\log 10=\log(10/9)+2\log3\). La evaluación de
\(\pi\) utiliza el cilindro regional generado de mil cifras, incluyendo
su anchura \(10^{-1000}\). La aritmética de intervalos en
\eqref{eq:alpha-firma-coeficientes} da
\[
 -\frac{4559}{10^{23}}<E_{21/22}^{(9)}(x_-)< -\frac{4558}{10^{23}},
 \qquad
 \frac{1698}{10^{23}}<E_{21/22}^{(9)}(x_+)<\frac{1699}{10^{23}}.
\]
Estos signos y la monotonía probada establecen el aislamiento. Los extremos
son resultados de la evaluación del operador; sus coeficientes proceden de
la firma previamente construida.
\end{proof}

El resultado concierne a \(\xi_9\). No autoriza a escribir
\(\xi_9=\alpha_A\) como una igualdad exacta de números completos.
La raíz de un truncamiento puede compartir un prefijo con la raíz del
operador prolongado y diferir en una posición posterior. La simplicidad
permite controlar el movimiento de la raíz bajo perturbaciones; no
anula esas perturbaciones.

\subsection{Compatibilidad de las dos determinaciones a toda profundidad}
\label{subsec:alpha-limite-dos-vias}

\begin{remark}[Truncamiento, prolongación y compatibilidad]
Las operaciones de los grados \(6\) a \(9\) determinan un jet finito.
La normalización posicional determina la coordenada completa mediante
su sucesión compatible de acarreos. La proposición siguiente expresa
el criterio general de identificación por cilindros. La
sección~\ref{subsec:alpha-publicaciones-completas} construye después
una carta analítica correlacionada de todos los grados, con recurrencia,
convergencia uniforme, unicidad de la raíz y compatibilidad proyectiva.
La construcción recibe la coordenada compartida del estado y la clase
de prolongaciones declarada; conserva así la diferencia entre una
representación analítica de esa coordenada y una selección independiente.
\end{remark}

La vía analítica completa se organiza en una familia de lectores
graduados compatibles
\[
 E_{21/22}^{(6)}\longleftarrow E_{21/22}^{(9)}
 \longleftarrow E_{21/22}^{(12)}\longleftarrow\cdots.
\]
El símbolo
\[
 E_{21/22}=\varprojlim_m E_{21/22}^{(6+3m)}
\]
designa ese lector completo, con sus transportes y truncamientos, no una
serie formal elegida libremente para hacerla pasar por \(\alpha_A\).
La segunda vía determina su raíz positiva única en \(I_\alpha\),
denotada por \(\alpha_B\), con la unicidad y la compatibilidad del
lector completo demostradas en la sección~\ref{subsec:alpha-publicaciones-completas}.

Las demostraciones de la sección~\ref{subsec:alpha-publicaciones-completas}
establecen que ambos sistemas determinan, a cada profundidad, el mismo
cilindro superviviente \(C_N^\alpha\), y que sus diámetros tienden a
cero. Conviene exponer por separado la consecuencia matemática de esta
compatibilidad.

\begin{proposition}[Identificación de los límites compatibles]
\label{prop:alpha-dos-limites}
Si los valores de las vías entera y analítica pertenecen a los
mismos cilindros \(C_N^\alpha\) para todo \(N\), y
\(\operatorname{diam}C_N^\alpha\to0\), entonces
\[
 \alpha_A=\alpha_B=\alpha.
\]
\end{proposition}

\begin{proof}
Para todo \(N\), la distancia entre ambos valores está acotada
por \(\operatorname{diam}C_N^\alpha\). Hacer tender \(N\) a infinito
da distancia cero. La existencia de las secciones compatibles asegura
que no se trata de la intersección vacía de un sistema de aproximaciones.
\end{proof}

Esta prueba identifica la consecuencia de los cuadrados de compatibilidad.
Su construcción a profundidad arbitraria aparece en la
sección~\ref{subsec:alpha-publicaciones-completas}. Las evaluaciones
finitas contrastan realizaciones de esa construcción; la recurrencia
y sus cotas uniformes establecen el paso al límite.

La diferencia puede expresarse algebraicamente. Sea
\(F=\mathbb Q(\pi,\delta)\) y sea
\(j_9:F[[x]]\to F[x]/(x^{10})\) el truncamiento. Entonces
\[
 \ker j_9=x^{10}F[[x]].
\]
Todos los operadores
\[
 E_h(x)=E_{21/22}^{(9)}(x)+x^{10}h(x)
\]
tienen el mismo truncamiento de orden nueve. Sin embargo, tomando
\(h=0\) y \(h=1/729\),
\[
 E_{1/729}(\xi_9)=\frac{\xi_9^{10}}{729}>0,
\]
de modo que no comparten esa raíz. La observación no permite elegir una
cola ajena al transporte HMT; precisamente demuestra por qué la cola
generada y su compatibilidad deben conservarse. Tampoco basta una serie
formal arbitraria para hablar de raíces analíticas sin su dominio de
convergencia o sin la realización del sistema de cilindros.
\input{sections/alpha_publicacion_completa.tex}


% SOURCE foundation/sections/revision_alpha_notacion.tex
\paragraph{Balance firmado y suma con cociente entrante.}
\label{ampl-alfa-z-s-aclaracion}
La división posicional consta de dos etapas consecutivas. Con
\(K_j^{\mathrm{per}}=K_{1+((j-1)\bmod12)}\), que designa la
prolongación periódica denotada también por \(K_j\) en
\eqref{eq:alpha-recurrencia}, el balance anterior al transporte del
cociente es
\[
 Z_j:=P_j+E_j-\Phi_j-K_j^{\mathrm{per}}.
\]
La variable $s_j$ de \eqref{eq:alpha-recurrencia} designa la suma que ya
incluye la contribución entrante, por lo que
\[
 s_j=Z_j+c_{j+1},\qquad
 A_j=Z_j+c_{j+1}-1000c_j.
\]
La configuración incidencial
$H^-_\alpha$ utiliza los signos del balance $Z_j$ de la ventana inicial;
la determinación de las cifras utiliza además los cocientes de enlace
$c_{j+1}$ y $c_j$. Esta distinción conserva conjuntamente el soporte
firmado y la normalización de la expansión.


% SOURCE foundation/sections/revision_alpha_frontera.tex
% Borrador focal de adición, NO demostración de la igualdad de dos vías.
% Puede incorporarse su proposición positiva sin alterar fórmulas actuales.
% La nota PROCEDENCIA_Y_ALCANCE.md delimita la compatibilidad no reconstruida.
\paragraph{Dependencia exacta respecto de la frontera de truncamiento.}
\label{ampl-alfa-frontera-exacta}
Las coordenadas regionales y el registro dodecafásico proceden de la
generación común APP--TRIT--TPK. Una vez producidas, su composición
posicional posee una dependencia explícita respecto de la profundidad
de lectura. Sean $B=1000$,
\[
 Z_j=P_j+E_j-\Phi_j-K_j^{\mathrm{per}},\qquad
 y_N=\sum_{j=1}^{N}Z_jB^{-j},\qquad
 R_{N+1}(Z)=\sum_{r\geq0}Z_{N+1+r}B^{-r-1}.
\]
La serie completa de \eqref{eq:alpha-serie} satisface exactamente
\[
 \alpha_A=y_N+B^{-N}R_{N+1}(Z),\qquad
 -2<R_{N+1}(Z)<2.
\]
Esta identidad reúne la contribución de la ventana y la de su
prolongación, ambas expresadas en las mismas coordenadas.

\begin{proposition}[Transporte de la frontera y cilindro canónico]
\label{ampl-alfa-prop-frontera}
Sea $A_j^{(N,t)}$ la salida de la división posicional finita con
$c_{N+1}=t\in\mathcal C=\{-2,-1,0,1,2\}$, y sea
$a_N^{(t)}=\sum_{j=1}^{N}A_j^{(N,t)}B^{-j}$. Para los registros del
dominio considerado, se cumple
\begin{equation}
 \alpha_A-a_N^{(t)}
 =B^{-N}\bigl(R_{N+1}(Z)-t\bigr),
 \qquad
 |\alpha_A-a_N^{(t)}|<4B^{-N}.
 \label{ampl-alfa-eq-frontera}
\end{equation}
El cociente de la prolongación completa es
$c_{N+1}^{\infty}=\lfloor R_{N+1}(Z)\rfloor$. Con esta frontera,
\[
 0\leq\alpha_A-a_N^{(c_{N+1}^{\infty})}<B^{-N},
\]
de modo que la ventana obtenida es el prefijo canónico de orden $N$.
\end{proposition}
\begin{proof}
El telescopado de \eqref{eq:alpha-identidad-local} da
$a_N^{(t)}=y_N+tB^{-N}-c_1$. Como $Z_1=7$, cada estado entrante
pertenece a $\mathcal C$ y el primer dividendo queda entre $5$ y $9$,
se tiene $c_1=0$. Sustituir esta igualdad en la descomposición de la
serie prueba \eqref{ampl-alfa-eq-frontera}. La cota resulta de
$R_{N+1}(Z)\in(-2,2)$ y $t\in\mathcal C$.

Para la normalización canónica completa,
$c_{N+1}^{\infty}=R_{N+1}(Z)-R_{N+1}(A)$ y
$R_{N+1}(A)\in[0,1)$. La integralidad del cociente implica
$c_{N+1}^{\infty}=\lfloor R_{N+1}(Z)\rfloor$. La fórmula de error
se convierte entonces en
$\alpha_A-a_N^{(c_{N+1}^{\infty})}=B^{-N}R_{N+1}(A)$, que prueba
la pertenencia al cilindro canónico.
\end{proof}

La proposición identifica el significado de la frontera derecha: es la
parte entera de la cola firmada expresada en la escala de la posición
siguiente. El cálculo sobre las cinco fronteras conserva esta información
como una familia finita de posibilidades hasta que los bloques a la
izquierda se estabilizan. La contracción cuantitativa procede de la
profundidad y del transporte de la cola; su objeto es la normalización de
las coordenadas ya generadas.

\paragraph{Dominio graduado del resolvente analítico.}
\label{ampl-alfa-dominio-resolvente}
La composición analítica exhibida en
\S\ref{subsec:alpha-resolvente} actúa sobre
$\operatorname{span}\{e_7,e_8,e_9\}$ y satisface $Ne_9=0$. Su
evaluación produce exactamente $T_{7:9}$; el grado $9$ es la última
componente de ese dominio. La identidad $N^3=0$ explica por qué el
resolvente se expresa mediante tres sumandos. Para describir una
prolongación en grados superiores hay que especificar, además, el
transporte entre dominios graduados sucesivos y su acción sobre la
componente inicial de cada dominio. La nilpotencia local y la
compatibilidad global desempeñan así funciones matemáticas distintas.


% SOURCE foundation/sections/alpha_dependencias.tex
\subsubsection{Determinación constructiva y función del registro dodecafásico}
\label{alpha-ampl-dependencias}

La definición anterior permite precisar qué significa determinar alfa desde
la estructura. El dominio de la operación contiene cuatro registros de
procedencia conocida: las tres secuencias regionales y el registro integral
de fase. La normalización incorpora, además, una condición de frontera
compatible con la prolongación. Una vez fijados esos datos, la división
euclídea determina cada resto y su cociente. La proposición
\ref{prop:alpha-normalizacion} extiende esa determinación al sistema
completo. La independencia respecto de un valor físico objetivo es, por
tanto, una propiedad del orden de construcción: el valor que se compara al
final no interviene en la selección de ninguno de los argumentos de esta
aplicación.

Esta independencia causal debe distinguirse de la independencia algebraica
entre números. La primera se establece identificando el dominio, las
operaciones y la procedencia de los coeficientes; la segunda sería una
afirmación sobre relaciones polinómicas. El resultado utilizado aquí es el
primero. La construcción determina una coordenada mediante registros
producidos por APP--TRIT--TPK y permite comprobar posteriormente las
relaciones que satisface. Ese orden hace posible utilizar una coordenada ya
generada en una operación ulterior sin convertirla en un parámetro externo.

El registro $K$ tiene una función más precisa que la de una corrección
numérica añadida a una fórmula. Sus doce posiciones proceden de la
transformación reversible del registro firmado. La ecuación
$U=\mathcal H_{12}K$ conserva la posibilidad de volver a ese registro,
mientras las diferencias transversales y la suma total permiten otra
reconstrucción. En la evaluación periódica, el orden de las posiciones
determina el peso posicional de cada componente. El mismo conjunto de
componentes, dispuesto en otro orden, produce en general otra coordenada
racional. Por ello el origen de fase y la orientación forman parte de la
definición del número que interviene en \eqref{eq:alpha-identidad-completa}.

La normalización de alfa compone esa información con las regiones. Su
identidad local distingue dos operaciones sucesivas. Primero se forma el
balance orientado de bloques; después se representa ese balance mediante
un resto canónico y un cociente que enlaza posiciones contiguas. La primera
operación conserva el signo del balance y permite la extracción posterior
de un soporte. La segunda permite formar los prefijos de la coordenada
real. Ambas lecturas permanecen relacionadas porque el estado conserva los
datos anteriores y posteriores a la normalización. El soporte firmado y la
expansión posicional son, así, realizaciones diferentes de una composición
aritmética explícita.

La identidad telescópica expresa el contenido global de esta composición.
Los cocientes interiores se cancelan por pares al evaluar un prefijo,
mientras sus extremos permanecen en la igualdad. Esa cancelación es una
propiedad de la suma ponderada; no elimina los cocientes del estado ni
permite suprimirlos antes de efectuar la suma. La frontera derecha de una
ventana recibe la contribución de la continuación. Por esa razón, la
evaluación de una ventana aislada y la restricción de una sección
prolongada son operaciones distintas. El cambio del bloque final $663$ a
$662$, demostrado anteriormente, ofrece una realización explícita de esa
diferencia sin alterar los bloques ya estabilizados a su izquierda.

En este sentido, la bidireccionalidad tiene una formulación aritmética
concreta. La generación proporciona prolongaciones compatibles; la
normalización transmite la información entera de su frontera desde las
posiciones de menor peso hacia las de mayor peso. La ecuación conserva la
contribución de ambas direcciones mediante sus índices y sus condiciones
de contorno. La reversibilidad en la fibra fijada se comprueba reconstruyendo
$K$ a partir de los demás registros y de los cocientes completos. Esa
inversión certifica la consistencia de la operación original; la dirección
genealógica continúa siendo la que produce primero el registro y después
la coordenada de alfa.

La segunda determinación tiene un contenido complementario. En ella, el
objeto inmediato es un balance de vacancias y el parámetro aparece como
su raíz en un dominio aislante. La vía posicional proporciona una
normalización canónica de registros; la vía analítica estudia la anulación
de un operador. La exposición de ambas requiere conservar precisamente
esa diferencia de funciones. Sus relaciones con el estado común han de
expresarse mediante las aplicaciones de refinamiento y la comparación de
sus realizaciones completas, tal como se desarrolla en la subsección
correspondiente. El acuerdo de una aproximación finita constituye un
control localizado; la identificación de las construcciones completas
tiene el dominio y el alcance de su propio enunciado.

Finalmente, esta organización explica la posición de alfa en el artículo.
Su ecuación aparece después de la generación regional y de la construcción
de $K$, porque utiliza ambos resultados. La sección electrónica utiliza
después esa coordenada como un factor de su composición. La prolongación
excepcional recibe los registros y sus soportes, donde se conserva la
información de incidencia. Estas dos continuaciones retoman una misma
procedencia y desarrollan funciones matemáticas distintas: evaluación de
una sección central y construcción de una configuración de frontera.
El desarrollo de cada una permite seguir la estructura a través de la
ecuación, en lugar de limitar la comparación a su valor escalar.


% SOURCE foundation/sections/alpha_publicacion_completa.tex
\subsection{Compatibilidad entre las publicaciones dodecafásica y analítica de alfa}
\label{subsec:alpha-publicaciones-completas}

La compatibilidad se establece a toda profundidad. La publicación
dodecafásica reversible y la representación analítica de la
coordenada de alfa
\(E_{21/22}\) son dos publicaciones internas con codominios distintos,
producidas a partir del mismo estado enriquecido APP--TRIT--TPK. El jet
\(E_{21/22}^{(9)}\) es un testigo finito exacto, pero no determina por sí solo
la cola. La completación que se construye después no se deduce del jet
desnudo: recibe la coordenada compartida del estado y publica su imagen
analítica mediante una recurrencia explícita. Por ello no constituye una
selección causal independiente; constituye una representación analítica
tipada del mismo punto generado.

Esta separación conserva el orden causal
\[
 \mathrm{APP}\longrightarrow\mathrm{TRIT}\longrightarrow\mathrm{TPK}
 \longrightarrow\mathcal X_{\alpha}^{\mathrm{enr}}
 \longrightarrow (P,E,\Phi,U_{12}^{\mathrm{sgn}},K)
 \longrightarrow\alpha_{\mathrm{HMT}}.
\]
Las coordenadas \(P,E,\Phi\) y el registro \(K\) son salidas anteriores del
mismo estado enriquecido. Su uso en la clausura de alfa es, por tanto, una
composición interna posterior y no una inyección de constantes
convencionales.

Para \(r\in6\mathbb N\), \(r\ge36\), sea
\(\widehat{\mathcal X}^{\mathrm{enr},\alpha}_r\) el sector de estados
truncados que prolongan la semilla interna de alfa fijada en \(R_{36}\). Se
escribe
\[
 x_r^\alpha=(\widetilde x_r,\chi_\alpha,
              \varepsilon_r^\alpha,N_r^\alpha),
\]
donde \(\chi_\alpha\) es el carácter interno del sector, no el valor escalar
publicado. Cada estado conserva residuo, cociente, acarreo, hoja, orientación,
fase, ruta, cilindro, frontera, firma, registro, memoria, supervivencia,
procedencia y profundidad. Si
\[
 d_r^\alpha=\operatorname{Dig}_{729}(\varepsilon_r^\alpha),\qquad
 \varepsilon_{r+6}^\alpha=729\varepsilon_r^\alpha-d_r^\alpha,
 \qquad N_{r+6}^\alpha=729N_r^\alpha+d_r^\alpha,
\]
la transición efectiva es
\[
 \mathcal U_r^\alpha=
 \operatorname{Upd}_r^\alpha\circ
 \operatorname{Tra}_r^\alpha\circ
 \operatorname{Sel}_r^\alpha,
 \qquad
 \operatorname{Ext}^{\alpha,\to}_{r+6,r}
 =\operatorname{graph}(\mathcal U_r^\alpha).
\]
La funcionalidad de estas extensiones, con la semilla \(R_{36}\) fijada,
produce una única sección compatible: la fase retorna mientras memoria y
profundidad avanzan. El dominio global de ambas determinaciones es
\begin{equation}
 \mathcal X_\alpha^{\mathrm{enr}}:=\widehat\Omega_\alpha
 :=\varprojlim_{r\ge36}
 \bigl(\widehat{\mathcal X}^{\mathrm{enr},\alpha}_r,
       \operatorname{Ext}^{\alpha,\to}_{r+6,r}\bigr).
 \label{eq:alpha-dominio-enriquecido-completo}
\end{equation}

En esta sección, \(\kappa_{\mathrm{per}}\) es la lectura periódica \(\kappa_{\mathrm{ag}}\) del registro ya construido; ambas notaciones designan la misma fracción en la carta de base mil fijada.

\subsubsection{La clausura dodecafásica como sección completa}

Para \(N\geq1\), pongamos
\[
 \mathcal W_N:=\{0,\ldots,999\}^{N},
 \qquad
 \mathsf T_{\alpha,N}(\mathsf s):=(A_1,\ldots,A_N),
 \qquad \mathsf s\in\mathcal X_{\alpha}^{\mathrm{enr}},
\]
donde los bloques \(A_j\) se obtienen mediante
\eqref{eq:alpha-recurrencia} con la frontera compatible
\(c_{N+1}^{\infty}=\lfloor R_{N+1}(Z)\rfloor\). Si \(N'\geq N\), sea
\(\rho_{N',N}:\mathcal W_{N'}\to\mathcal W_N\) la restricción a los
primeros \(N\) bloques.

\begin{theorem}[Salida dodecafásica completa y compatibilidad proyectiva]
\label{thm:alpha-salida-dodecafasica-completa}
La familia \((\mathsf T_{\alpha,N})_{N\geq1}\) es compatible:
\begin{equation}
 \rho_{N',N}\circ\mathsf T_{\alpha,N'}
 =\mathsf T_{\alpha,N}
 \qquad(N'\geq N).
 \label{eq:alpha-via-a-proyectiva-exacta}
\end{equation}
Sus cilindros canónicos
\begin{equation}
 C_N^A(\mathsf s):=
 \left[
  \sum_{j=1}^{N}A_j1000^{-j},
  \sum_{j=1}^{N}A_j1000^{-j}+1000^{-N}
 \right)\cap I_\alpha
 \label{eq:alpha-cilindros-via-a-exactos}
\end{equation}
son anidados, tienen diámetro a lo sumo \(1000^{-N}\) y satisfacen
\begin{equation}
 \bigcap_{N\geq1}C_N^A(\mathsf s)
 =\{\alpha_A(\mathsf s)\},
 \qquad
 \alpha_A
 =\pi_{\mathrm{HMT}}+e_{\mathrm{HMT}}-\varphi_{\mathrm{HMT}}
  -4-\kappa_{\mathrm{per}}.
 \label{eq:alpha-salida-a-completa}
\end{equation}
En consecuencia, la primera vía publica una única sección infinita, que
denotamos \(\alpha_{\mathrm{HMT}}:=\alpha_A\).
\end{theorem}

\begin{proof}
La proposición~\ref{prop:alpha-normalizacion} proporciona la expansión
canónica única y la sucesión acotada única de acarreos. La identidad local
\eqref{eq:alpha-identidad-local} y su telescopado
\eqref{eq:alpha-telescopia} conservan la frontera remota; la proposición
\ref{ampl-alfa-prop-frontera} prueba que restringir una ejecución prolongada
recupera exactamente el prefijo anterior. Esto da
\eqref{eq:alpha-via-a-proyectiva-exacta}. La pertenencia de la suma completa
a cada cilindro se sigue de la identidad de colas, y
\(\operatorname{diam}C_N^A\leq1000^{-N}\to0\); la intersección contiene un
solo punto. Finalmente, \eqref{eq:alpha-identidad-completa} identifica ese
punto con la expresión de \eqref{eq:alpha-salida-a-completa}. Todos sus
coeficientes y signos se fijan antes de la publicación de alfa.
\end{proof}

La ventana \(\alpha_{12}^{(0)}\), cerrada con \(c_{13}=0\), sigue siendo un
cálculo finito exacto; no debe confundirse con la restricción de la sección
completa, cuya frontera compatible satisface \(c_{13}^{\infty}=-1\). En
particular,
\begin{equation}
 \rho_{36,12}(\alpha_{A,36})=\alpha_{A,12}
 \label{eq:alpha-rho-36-12-proyectiva}
\end{equation}
es una igualdad de palabras compatibles, no una igualdad entre palabras de
longitudes distintas ni una identificación de la ventana aislada con el
límite.


\input{sections/alpha_carta_analitica_completa.tex}


% SOURCE foundation/sections/alpha_carta_analitica_completa.tex
% Representación analítica correlacionada de alfa: acción, convergencia e
% intervalos.
\subsubsection{Representación analítica de la coordenada de alfa}
\label{subsec:alpha-lector-completo-rev08}

La prolongación analítica se especifica mediante todos sus coeficientes.
Sea \(\mathsf s\in\mathcal X_\alpha^{\mathrm{enr}}\). Antes
de la normalización decimal, el mismo estado ya publica la precoordenada
\begin{equation}
 a(\mathsf s):=p(\mathsf s)+e_0(\mathsf s)-f(\mathsf s)
                 -\kappa_{\mathrm{per}}(\mathsf s)\in I_\alpha,
 \label{eq:alpha-precoordenada-comun}
\end{equation}
con \(p=\pi_{\mathrm{HMT}}-3\),
\(e_0=e_{\mathrm{HMT}}-2\) y
\(f=\varphi_{\mathrm{HMT}}-1\). Equivalentemente,
\(a=\pi_{\mathrm{HMT}}+e_{\mathrm{HMT}}-
\varphi_{\mathrm{HMT}}-4-\kappa_{\mathrm{per}}\). No se introduce aquí
\(\alpha\) convencional: \(p,e_0,f\) son los tres caracteres regionales ya
generados y \(\kappa_{\mathrm{per}}\) es la lectura racional del registro
terminal. La clausura entera aplica a
\eqref{eq:alpha-precoordenada-comun} la normalización de
\eqref{eq:alpha-recurrencia}; la carta que sigue publica sobre la misma
precoordenada una representación analítica compatible. Su grafo de
dependencia recibe \((p,e_0,f,\kappa_{\mathrm{per}})\) directamente: no
contiene una arista \(\alpha_A\to\lambda\), aunque la igualdad demostrada más
abajo permita identificar después ambos nombres de publicación.

Escribamos \(q=1/729\), sea \(F_{\mathsf s}\) el cuerpo ordenado generado por
las coordenadas de \(\mathsf s\), y pongamos
\(E_{9,\mathsf s}(X)=E_{21/22}^{(9)}(\mathsf s;X)\). Definimos el coeficiente
de enlace
\begin{equation}
 \lambda(\mathsf s):=
 -\frac{E_{9,\mathsf s}(a(\mathsf s))
         \bigl(1-q\,a(\mathsf s)^3\bigr)}{a(\mathsf s)^{10}}.
 \label{eq:alpha-lambda-estado}
\end{equation}
Todos sus argumentos pertenecen al estado antes del reconocimiento
metrológico. En particular, \(\lambda\) no se obtiene ajustando una cadena
decimal objetivo. La ecuación~\eqref{eq:alpha-lambda-estado} debe leerse con
su tipo exacto: una vez que la primera publicación del estado ha producido
\(a(\mathsf s)\), fija la única prolongación de la clase geométrica que se
define a continuación. No constituye una segunda selección independiente de
\(a(\mathsf s)\), sino una representación analítica correlacionada del mismo
estado enriquecido.

En efecto, para \(\mu\in F_{\mathsf s}\) pongamos
\[
 E_{\mathsf s,\mu}(X)
 :=E_{9,\mathsf s}(X)+\mu\frac{X^{10}}{1-qX^3}.
\]
Como \(0<a(\mathsf s)<10^{-2}\) y \(qa(\mathsf s)^3<1\), se tiene
\begin{equation}
 E_{\mathsf s,\mu}(a(\mathsf s))=0
 \quad\Longleftrightarrow\quad
 \mu=\lambda(\mathsf s).
 \label{eq:alpha-unicidad-lambda-clase-geometrica}
\end{equation}
Así, la semilla de coeficientes no contiene libertad una vez fijados el estado
y la clase de prolongaciones \(q\)-geométricas. Esta unicidad es relativa a la
clase declarada; no afirma que el jet de orden nueve seleccione por sí solo una
cola entre todas las series analíticas posibles.

La acción local sobre bloques de tres coeficientes es
\begin{equation}
 \mathcal C_\alpha:F_{\mathsf s}^{3}\longrightarrow F_{\mathsf s}^{3},
 \qquad \mathcal C_\alpha(u,v,w)=q(u,v,w),
 \qquad c^{(0)}=(\lambda(\mathsf s),0,0).
 \label{eq:alpha-accion-coeficientes}
\end{equation}
Por tanto, para todo \(n\geq0\),
\begin{equation}
 b_{10+3n}=q^n\lambda(\mathsf s),\qquad
 b_{11+3n}=b_{12+3n}=0.
 \label{eq:alpha-recurrencia-coeficientes}
\end{equation}
Ésta es la regla que calcula cada coeficiente posterior al grado nueve. No
queda un parámetro libre en cada prolongación.

Para \(M\geq0\), sea
\begin{equation}
 E_{\mathsf s,M}(X):=E_{9,\mathsf s}(X)
 +\lambda(\mathsf s)\sum_{n=0}^{M}q^nX^{10+3n}.
 \label{eq:alpha-truncamientos-completos}
\end{equation}
La función completa es
\begin{equation}
 E_{\mathsf s,\infty}(X)
 :=E_{9,\mathsf s}(X)
 +\lambda(\mathsf s)\frac{X^{10}}{1-X^3/729}.
 \label{eq:alpha-funcion-completa-explicita}
\end{equation}

Para formular el sistema proyectivo sin tipos implícitos, fijamos
\(d_M=12+3M\) y definimos la fibración de jets
\begin{equation}
 \mathfrak J_M:=
 \coprod_{\mathsf s\in\mathcal X_\alpha^{\mathrm{enr}}}
 \{\mathsf s\}\times F_{\mathsf s}[X]/(X^{d_M+1}).
 \label{eq:alpha-espacio-jets-rev08}
\end{equation}
Si \(M'\ge M\), la restricción es
\begin{equation}
 \tau_{M',M}(\mathsf s,[P]_{d_{M'}})
 :=(\mathsf s,[P]_{d_M}).
 \label{eq:alpha-truncamiento-jets-rev08}
\end{equation}

\begin{proposition}[Naturalidad de la recurrencia analítica]
\label{prop:alpha-naturalidad-recurrencia-explicita}
Sea
\[
 \Gamma_M(\mathsf s)
 =\bigl(\mathsf s,[E_{\mathsf s,M}]_{d_M}\bigr).
\]
Si \(M'\geq M\), el truncamiento \(\tau_{M',M}\) satisface
\begin{equation}
 \tau_{M',M}\circ\Gamma_{M'}=\Gamma_M,
 \qquad
 \Gamma_{E21}(\mathsf s)=\varprojlim_M\Gamma_M(\mathsf s).
 \label{eq:alpha-naturalidad-explicita}
\end{equation}
La proyección sobre los tres grados nuevos es precisamente
\(\mathcal C_\alpha^{M+1}(c^{(0)})\).
\end{proposition}

\begin{proof}
La ecuación~\eqref{eq:alpha-recurrencia-coeficientes} fija el bloque de
grados \(10+3n,11+3n,12+3n\). Al pasar de \(M\) a \(M+1\) sólo se añade
el bloque siguiente; truncarlo recupera literalmente el polinomio anterior.
La compatibilidad para \(M'\geq M\) se obtiene por composición. Así, el
límite inverso no es una sección escogida del jet desnudo, sino la órbita
única de \(c^{(0)}\) bajo \(\mathcal C_\alpha\).
\end{proof}

\subsubsection{Convergencia uniforme, raíz simple e intervalos racionales}
\label{subsec:alpha-cotas-completas-rev08}

Los cilindros regionales y el registro periódico proporcionan cotas mediante
aritmética de intervalos racionales. Para cada coordenada
\[
 c\in\{\pi_{\mathrm{HMT}},e_{\mathrm{HMT}},\varphi_{\mathrm{HMT}}\},
\]
sea \(P_c\) su prefijo nonádico certificado de mil cifras. El dato utilizado
no es el racional \(P_c\) como si fuese el valor completo, sino el cilindro
semiabierto y su cierre exterior
\[
 \mathcal C_c^{(1000)}=[P_c,P_c+\varepsilon),
 \qquad
 \overline{\mathcal C}_c^{(1000)}=[P_c,P_c+\varepsilon],
 \qquad \varepsilon=10^{-1000}.
\]
La aritmética dirigida opera conservadoramente con esos cierres. Si
\(P_\pi,P_e,P_\varphi\) denotan los tres prefijos, se obtiene
\begin{equation}
\begin{aligned}
 a_-&:=P_\pi+P_e-(P_\varphi+\varepsilon)-4-
       \kappa_{\mathrm{per}},\\
 a_+&:=(P_\pi+\varepsilon)+(P_e+\varepsilon)-P_\varphi-4-
       \kappa_{\mathrm{per}},\\
 A_{1000}&:=[a_-,a_+],
 \qquad a(\mathsf s)\in[a_-,a_+]=A_{1000},
 \qquad \operatorname{diam}A_{1000}=3\varepsilon.
\end{aligned}
\label{eq:alpha-propagacion-colas-rev08}
\end{equation}
La extensión intervalar natural de
\eqref{eq:alpha-lambda-estado} transporta después \(A_{1000}\), el cilindro
de \(\pi_{\mathrm{HMT}}\) y la cota racional dirigida de \(\delta\) a un
intervalo cerrado \(\Lambda_{1000}\ni\lambda(\mathsf s)\); no se congela
ninguna de las tres colas. En particular,
\begin{equation}
\begin{aligned}
 \frac{7297352569283800997285105472380662}{10^{36}}
 &<a(\mathsf s)<
 \frac{7297352569283800997285105472380663}{10^{36}},\\
 -8.711\mathbin{\times}10^{-6}
 &<\lambda(\mathsf s)<-8.709\mathbin{\times}10^{-6}.
\end{aligned}
 \label{eq:alpha-cotas-precoordenada-lambda}
\end{equation}
Estas desigualdades usan prefijos certificados con su cota racional de
cola; no reciben una tabla de alfa. El verificador del paquete las reproduce
desde las tres salidas nonádicas y desde la carta terminal.

Sea
\[
 r:=\frac{10^{-6}}{729}<1.372\mathbin{\times}10^{-9}.
\]
Para \(0\leq x\leq10^{-2}\), la cola posterior a \(M\) satisface
\begin{equation}
 \left|E_{\mathsf s,\infty}(x)-E_{\mathsf s,M}(x)\right|
 \leq
 8.711\mathbin{\times}10^{-26}\,\frac{r^{M+1}}{1-r}.
 \label{eq:alpha-cota-cola-uniforme}
\end{equation}
Además,
\begin{equation}
 \left|\frac{\mathrm d}{\mathrm dx}
  \left(\lambda\frac{x^{10}}{1-x^3/729}\right)\right|
 \leq 8.711\mathbin{\times}10^{-6}
 \left(\frac{10\,10^{-18}}{1-r}
 +\frac{3\,10^{-24}/729}{(1-r)^2}\right)
 <8.712\mathbin{\times}10^{-23}.
 \label{eq:alpha-cota-derivada-cola}
\end{equation}
La cola de las derivadas posterior al bloque \(M\) admite, además, la
estimación dependiente de la profundidad
\begin{equation}
 \sup_{0\le x\le10^{-2}}
 \left|E'_{\mathsf s,\infty}(x)-E'_{\mathsf s,M}(x)\right|
 \le 8.711\mathbin{\times}10^{-24}\,r^{M+1}
 \left(\frac{13+3M}{1-r}+\frac{3r}{(1-r)^2}\right)
 \xrightarrow[M\to\infty]{}0.
 \label{eq:alpha-cota-resto-derivadas-rev08}
\end{equation}

\begin{theorem}[Función límite y raíz simple única de la carta correlacionada]
\label{thm:alpha-raiz-completa-explicita}
La sucesión \(E_{\mathsf s,M}\) converge uniformemente, junto con sus
derivadas, a \(E_{\mathsf s,\infty}\) en
\(\overline I_\alpha=[0,10^{-2}]\). Se cumple
\begin{equation}
 E_{\mathsf s,\infty}(a(\mathsf s))=0,
 \qquad
 \inf_{x\in\overline I_\alpha}
 E'_{\mathsf s,\infty}(x)>6.239.
 \label{eq:alpha-raiz-y-derivada-completas}
\end{equation}
Por consiguiente, \(a(\mathsf s)\) es la raíz simple única de la función
completa en \(I_\alpha\).
\end{theorem}

\begin{proof}
La razón de dos términos consecutivos de la cola está acotada por \(r<1\),
lo que da \eqref{eq:alpha-cota-cola-uniforme}; la derivación término a
término y la suma de las dos series geométricas dan
\eqref{eq:alpha-cota-derivada-cola} y
\eqref{eq:alpha-cota-resto-derivadas-rev08}. La prueba del jet establece
\(E'_{9,\mathsf s}>6.24\) en \(\overline I_\alpha\); por tanto, la
perturbación completa deja la derivada por encima de \(6.239\).
Finalmente, al sustituir \eqref{eq:alpha-lambda-estado} en
\eqref{eq:alpha-funcion-completa-explicita}, los dos sumandos evaluados en
\(a(\mathsf s)\) se cancelan exactamente. La función es estrictamente
creciente; su única raíz es, pues, simple y coincide con esa precoordenada.
\end{proof}

Para hacer efectiva la identificación a toda profundidad, las evaluaciones
dirigidas de los extremos certifican primero
\[
 E_{\mathsf s,\infty}(0)<0<
 E_{\mathsf s,\infty}(10^{-2}),
 \qquad
 E'_{\mathsf s,\infty}\geq\mu:=6239/1000>c:=6
 \quad\hbox{en }[0,10^{-2}].
\]
Por continuidad existe una raíz en el intervalo y, por la segunda desigualdad,
es única. Sólo después de fijar esta existencia, la monotonía y la invariante
de que la horquilla contiene la raíz se aplica el rescate por el teorema del
valor medio. Partimos de \(D_0=[0,10^{-2}]\). Si
\(D_j=[r_j,s_j]\), sea \(m_j=(r_j+s_j)/2\), y sea
\([F_j^-,F_j^+]\) una inclusión racional dirigida de
\(E_{\mathsf s,\infty}(m_j)\), obtenida propagando los cilindros de
\(\pi_{\mathrm{HMT}},e_{\mathrm{HMT}},\varphi_{\mathrm{HMT}}\),
\(\delta\) y \(\lambda\), con
\begin{equation}
 0\leq F_j^+-F_j^-\leq \frac{c}{4}\,(s_j-r_j).
 \label{eq:alpha-anchura-evaluacion-rescate-rev08}
\end{equation}
Cuando \(F_j^+<0\) se conserva la mitad derecha;
cuando \(F_j^->0\), la mitad izquierda. Si
\(F_j^-\leq0\leq F_j^+\), no se intenta decidir si el punto medio es
exactamente la raíz. Sea \(w_j=s_j-r_j\). La cota uniforme
\(E'_{\mathsf s,\infty}\geq c=6\) y el teorema del valor medio implican que
cualquier raíz contenida en \(D_j\) pertenece a
\begin{equation}
 D_{j+1}=D_j\cap
 \left[m_j-\frac{w_j}{4},
             m_j+\frac{w_j}{4}\right].
 \label{eq:alpha-rescate-derivada-rev08}
\end{equation}
En efecto, sea \(\xi_j\) la raíz y sea
\(y_j=E_{\mathsf s,\infty}(m_j)\in[F_j^-,F_j^+]\). Como el recinto contiene
cero y satisface \eqref{eq:alpha-anchura-evaluacion-rescate-rev08}, se tiene
\(|y_j|\leq cw_j/4\). El teorema del valor medio proporciona
\(|m_j-\xi_j|\leq |y_j|/c\leq w_j/4\), lo que prueba
\eqref{eq:alpha-rescate-derivada-rev08}. El argumento incluye el caso
\(y_j=0\): la raíz está entonces en el punto medio, pero el algoritmo no
necesita reconocer esa igualdad para conservarla.

Si el recinto disponible no satisface
\eqref{eq:alpha-anchura-evaluacion-rescate-rev08}, se refinan primero los
cilindros fuente y se repite la evaluación; no se declara un signo ni se
acepta una contracción insuficiente. En cualquiera de las tres ramas el nuevo
intervalo conserva la raíz y tiene diámetro a lo sumo \(w_j/2\). Para cada
profundidad objetivo finita, el refinamiento dirigido permite satisfacer la
cota de anchura sin comparar de forma no certificada un número real con cero.
La monotonía demuestra inductivamente
\begin{equation}
 a(\mathsf s)\in D_{j+1}\subseteq D_j,\qquad
 \operatorname{diam}D_{j+1}\leq
 \tfrac12\operatorname{diam}D_j,
 \qquad \operatorname{diam}D_j\longrightarrow0,
 \label{eq:alpha-biseccion-racional}
\end{equation}
donde la desigualdad es una cota de contracción y no una igualdad impuesta.
El certificado ejecutable aplica
esta regla a veinticuatro niveles en base mil y comprueba que el cierre
exterior completo \(A_{1000}\), no sólo su extremo inferior, queda dentro de
cada cilindro publicado. Sus pruebas negativas rechazan un recinto demasiado
ancho, un recinto desordenado, un punto que no sea el medio, un dominio sin
cambio de signo y una horquilla degenerada. La prolongación coinductiva permite
reemplazar mil por cualquier profundidad requerida.

Concretamente, para \(S_N=1000^N\) el verificador calcula
\begin{equation}
 j_N^-:=\lfloor S_Na_-\rfloor,
 \qquad j_N^+:=\lfloor S_Na_+\rfloor
 \label{eq:alpha-indices-cilindro-dirigidos-rev08}
\end{equation}
y sólo publica el nivel si \(j_N^-=j_N^+\). Esta igualdad, que incluye el
extremo superior del cierre exterior, prueba
\(A_{1000}\subset[j_N^-/S_N,(j_N^-+1)/S_N)\) y evita atribuir a la celda
izquierda un extremo situado exactamente sobre su frontera derecha.

Sea \(C_N^A(\mathsf s)\) el cilindro entero de
\eqref{eq:alpha-cilindros-via-a-exactos}. Elegimos recursivamente
\(m(N)>m(N-1)\) de modo que
\(\operatorname{diam}D_{m(N)}<1000^{-N}/4\), y definimos
\begin{equation}
I_{B,N}:=D_{m(N)}\cap C_N^A(\mathsf s).
 \label{eq:alpha-intervalos-b-completos}
\end{equation}
Definimos
\(\ell_N:=\inf I_{B,N}\) y \(u_N:=\sup I_{B,N}\). Ambos son racionales;
el intervalo es no vacío porque los dos factores contienen
\(a(\mathsf s)\). Puesto que el cilindro es semiabierto, \(u_N\) puede no
pertenecer a \(I_{B,N}\). La evaluación en ese extremo se entiende mediante
la extensión continua de \(E_{\mathsf s,\infty}\) al cierre exterior. Así,
\begin{equation}
 I_{B,N+1}\subseteq I_{B,N}\subseteq C_N^A(\mathsf s),
 \qquad \operatorname{diam}I_{B,N}\longrightarrow0,
 \qquad E_{\mathsf s,\infty}(\ell_N)\leq0\leq
 E_{\mathsf s,\infty}(u_N).
 \label{eq:alpha-inclusion-signos-completa}
\end{equation}
La cota~\eqref{eq:alpha-cota-cola-uniforme} permite escoger un
truncamiento que conserve el signo en cada extremo donde la función completa
no se anula. Si un extremo coincide con la raíz, se utiliza la identidad de
la función completa y la pertenencia ya demostrada; no se atribuye ese mismo
signo a sus truncamientos. En efecto, el resto exacto es
\[
 E_{\mathsf s,\infty}(x)-E_{\mathsf s,M}(x)
 =\lambda(\mathsf s)x^{10}
   \frac{(qx^3)^{M+1}}{1-qx^3}.
\]
En el dominio positivo considerado, donde \(\lambda(\mathsf s)<0\), se
obtiene en la raíz
\[
 E_{\mathsf s,M}(a(\mathsf s))
 =-\lambda(\mathsf s)a(\mathsf s)^{10}
   \frac{(qa(\mathsf s)^3)^{M+1}}{1-qa(\mathsf s)^3}>0
 \qquad(M<\infty).
\]
Por ello, la certificación de los signos débiles en
\eqref{eq:alpha-inclusion-signos-completa} se refiere a
\(E_{\mathsf s,\infty}\). Su evaluación utiliza la función completa o el
truncamiento acompañado del resto con signo. Esto cubre tanto el ínfimo
perteneciente como el supremo eventualmente excluido del cilindro.

\begin{proposition}[Cuadrados de compatibilidad y cilindros comunes]
\label{prop:alpha-cuadrados-compatibilidad}
Las restricciones de palabras y de publicaciones analíticas conmutan:
\[
\begin{array}{ccc}
 \mathcal X_\alpha^{\mathrm{enr}}&
 \xrightarrow{\ \mathsf T_{\alpha,N'}\ }&\mathcal W_{N'}\\[1mm]
 \Big\Vert&&\Big\downarrow\rho_{N',N}\\[1mm]
 \mathcal X_\alpha^{\mathrm{enr}}&
 \xrightarrow{\ \mathsf T_{\alpha,N}\ }&\mathcal W_N
\end{array}
\qquad
\begin{array}{ccc}
 \mathcal X_\alpha^{\mathrm{enr}}&
 \xrightarrow{\ \Gamma_{M'}\ }&\mathfrak J_{M'}\\[1mm]
 \Big\Vert&&\Big\downarrow\tau_{M',M}\\[1mm]
 \mathcal X_\alpha^{\mathrm{enr}}&
 \xrightarrow{\ \Gamma_M\ }&\mathfrak J_M .
\end{array}
\]
La familia \(I_{B,N}\) proporciona la conexión efectiva entre ambos
cuadrados y satisface materialmente
\(I_{B,N}\subseteq C_N^A\) para todo \(N\).
\end{proposition}

\begin{proof}
El primer cuadrado es la compatibilidad de la normalización entera; el segundo
es la proposición~\ref{prop:alpha-naturalidad-recurrencia-explicita}. La
inclusión se sigue de la definición~\eqref{eq:alpha-intervalos-b-completos},
y la no vacuidad y los signos se siguen del
teorema~\ref{thm:alpha-raiz-completa-explicita}. Ningún paso identifica la
raíz de un truncamiento con la de una cola arbitraria.
\end{proof}

Definimos ahora, sin dejar implícito el lector de prefijos,
\begin{equation}
 \alpha_B(\mathsf s):=
 \operatorname*{arg\,zero}_{x\in I_\alpha}E_{\mathsf s,\infty}(x),
 \qquad
 \operatorname{pref}_N(x):=
 \frac{\lfloor1000^Nx\rfloor}{1000^N},
 \qquad
 \alpha_{B,N}:=\operatorname{pref}_N(\alpha_B).
 \label{eq:alpha-definicion-b-y-prefijos-rev08}
\end{equation}

\begin{theorem}[Igualdad proyectiva de las dos publicaciones correlacionadas]
\label{thm:alpha-dos-publicaciones-completas}
Para todo \(N\), la vía entera y la vía analítica publican el mismo prefijo:
\begin{equation}
 \boxed{\alpha_{A,N}=\operatorname{pref}_N(\alpha_A)
 =\operatorname{pref}_N(\alpha_B)=\alpha_{B,N}.}
 \label{eq:alpha-dos-vias-cada-nivel}
\end{equation}
En el límite,
\begin{equation}
 \boxed{\alpha_A
 =\operatorname*{arg\,zero}_{x\in I_\alpha}
       E_{\mathsf s,\infty}(x)
 =\alpha_B=\alpha_{\mathrm{HMT}}.}
 \label{eq:alpha-dos-vias-limite-coinductivo}
\end{equation}
\end{theorem}

\begin{proof}
La normalización entera publica exactamente \(a(\mathsf s)\). El
teorema~\ref{thm:alpha-raiz-completa-explicita} demuestra que la carta
analítica representa ese mismo punto como raíz simple única. Los
cilindros semiabiertos \(C_N^A\) fijan un prefijo canónico y
\(I_{B,N}\subseteq C_N^A\) contiene \(\alpha_B=\alpha_A\). Por tanto, sus
prefijos coinciden para cada \(N\); el anidamiento y los diámetros tendentes a
cero dan la igualdad de los límites. El reconocimiento
metrológico ocurre después.
\end{proof}

La función analítica constituye una publicación correlacionada del mismo
estado, no una segunda derivación causalmente independiente de la clausura
entera. Esta precisión no reduce la igualdad demostrada: identifica su
mecanismo exacto, impide atribuir al jet una cola que no determina y evita la
promoción de un valor convencional a generador.


% SOURCE foundation/sections/excepcional.tex
\section{Prolongación excepcional: de la incidencia dodecafásica a Moonshine}
\label{exc:seccion}

La expresión numérica no agota el estado que la ha generado. En la
construcción precedente, APP conserva las dos hojas y la división con
residuo y cociente; el TRIT determina el régimen local y su orientación;
el TPK compone selección, transporte, acarreo, memoria y retorno. La
lectura arquimediana contrae cilindros compatibles. La lectura que ahora
estudiamos conserva, en cambio, la incidencia de las posiciones y su
transporte entre marcos. Ambas parten del mismo estado aritmético con
memoria de trayectoria. No se
reconstruye una geometría excepcional a partir de los decimales finales
de $\pi$, $\varphi$, $e$ o $\alpha$.

Este punto permite precisar el sentido de la prolongación
$\alpha/K\longrightarrow\text{incidencia excepcional}$. La flecha no
afirma que un número real aislado determine un código, un retículo o un
grupo. Transporta el registro dodecafásico y los datos firmados que han
intervenido en la determinación de $\alpha$, junto con las palabras y los
marcos procedentes de APP--TRIT--TPK. En particular, el registro ordenado
$K$ puede recuperarse desde su constante racional $\kappa$ con las
convenciones de base, longitud y origen ya fijadas. La proyección a un
soporte conserva, en cambio, sólo las incidencias seleccionadas.
Por su parte, $K_{\rm ph}$
designa otra lectura, de memoria reducida, posterior a la holonomía
$\Gamma_9$.

La exposición separa tres operaciones matemáticas: construir los objetos
finitos y sus mapas; prolongarlos a un retículo y a un álgebra de operadores
de vértice; reconocer después los objetos obtenidos mediante los teoremas
clásicos pertinentes. Los nombres Paley, Witt, Mathieu, Golay, Leech y
Monstruo designan esos reconocimientos y las realizaciones explícitas
que los hacen verificables. No son selectores retrospectivos de las
semillas ni de los coeficientes de la generación.

\input{sections/revision_incidencia_argumento.tex}



\subsection{Compatibilidad de las realizaciones arquimediana e incidencial}
\label{exc:doble-lectura}

Sea $X_n^{\rm cal}$ el espacio de prefijos admisibles de longitud $n$ del
estado aritmético con memoria de trayectoria, con truncaciones
$p_{n,m}:X_n^{\rm cal}\to X_m^{\rm cal}$ para $m\leq n$. El superíndice
recuerda que se conserva el marco inicial declarado y su transporte, no
que se ajuste un parámetro contra un valor experimental. Escribimos
$(w_j,f_j)$ para la palabra visible de seis trits y el marco en el paso
$j$. Los marcos satisfacen una ley causal
\[
 f_{j+1}=g_j(x)f_j,
\]
donde $g_j(x)$ depende del prefijo ya construido. En consecuencia, una
prolongación no altera los marcos anteriores.

La cadena
\[
 w_6\longrightarrow w_{12}\longrightarrow w_{18}
 \longrightarrow w_{24}\longrightarrow w_{30}
 \longrightarrow R_{36}\longrightarrow G_9
\]
se conserva en esta lectura. El retorno de la fase no reinicia ni el
marco ni la memoria del estado. Asimismo, se mantienen distintos el
censo de $468$ emisiones decimales, las $243$ palabras visibles de seis
trits y el ambiente $\mathbb F_3^6$, de $729$ palabras. El ambiente se
utilizará para construir un código lineal; no sustituye al dominio
admisible de emisiones.

La lectura arquimediana toma, para un bloque $b$ de seis trits, su
dirección entera $\nu(b)\in\{0,\ldots,728\}$. Un prefijo determina
\[
 q_n=m+\sum_{j=0}^{n-1}\frac{\nu(b_j)}{729^{j+1}},
 \qquad I_n=[q_n,q_n+729^{-n}].
\]
La compatibilidad de prefijos da cilindros encajados cuyo diámetro tiende
a cero. Su clase de Cauchy pertenece primero a la completación interna
$\mathbb R_{\rm HMT}$; su identificación posterior con la recta real
usual no interviene en la elección de los bloques.

Para la lectura incidencial introducimos, sobre el mismo prefijo, el
mapa de gráfico
\begin{equation}
 \gamma_A:\mathbb F_3^6\longrightarrow\mathbb F_3^{12},
 \qquad w\longmapsto(w,wA).
 \label{exc:gamma}
\end{equation}
Las palabras se escriben como vectores fila. Si $A^f=fAf^{-1}$, el
cambio de marco satisface la identidad concreta
\begin{equation}
 \gamma_{A^f}(wf^{-1})
   =\gamma_A(w)(f^{-1}\oplus f^{-1}).
 \label{exc:covariancia}
\end{equation}
El codominio calibrado conserva $f$ junto a la palabra codificada.
Definimos entonces
\begin{equation}
 Q_n(x)=\bigl((f_j,\gamma_{f_jA_Wf_j^{-1}}(w_j))\bigr)_{0\leq j<n}.
 \label{exc:qn}
\end{equation}

\begin{proposition}[Compatibilidad y prolongación de la lectura incidencial]
\label{exc:limite-inverso}
Si $r_{n,m}$ trunca las sucesiones codificadas, se cumple
\[
 r_{n,m}Q_n=Q_m p_{n,m}.
\]
Por tanto, los $Q_n$ determinan un único mapa $Q_\infty$ entre los
respectivos límites inversos. En cada marco, el mapa conserva exactamente
la palabra visible, pues $\operatorname{pr}_1\gamma_A=\operatorname{id}$.
\end{proposition}

\begin{proof}
Dos prefijos compatibles tienen las mismas palabras y los mismos
operadores de transporte hasta el índice $m-1$. Tienen, por inducción,
los mismos marcos en esos índices. Las primeras $m$ componentes de
\eqref{exc:qn} coinciden, lo que prueba el cuadrado. La propiedad
universal del límite inverso da $Q_\infty$. La afirmación de recuperación
es la proyección sobre las seis primeras coordenadas de cada gráfico.
\end{proof}

Esta inyectividad local tiene un dominio preciso. No demuestra que
$Q_\infty$ recupere toda la memoria profunda del estado TPK si sólo se
ha retenido la sucesión de palabras visibles y marcos. Mucho menos
demuestra inyectividad después de pasar de palabras a soportes.
Además, un cambio general de base en $\operatorname{GL}_6(\mathbb F_3)$
no conserva por sí mismo el peso de Hamming ordinario. En una carta
general se transportan también la incidencia y la métrica de la carta
inicial. Los cálculos de soportes que siguen se realizan en la carta
Paley declarada; los relabelados monomiales conservan directamente su
interpretación coordenada.

\subsection{La realización finita y el código ternario}
\label{exc:paley-codigo}

El operador elíptico del TRIT admite la siguiente realización finita en
la carta de seis posiciones. Éstas coordinatizan las seis unidades
$\{1,2,4,5,7,8\}$ módulo nueve, conservadas por el transporte modular.
Sobre $\mathbb P^1(\mathbb F_5)$, una vez
declarado el orden $\infty,0,1,2,3,4$, sea $C$ la matriz con diagonal
nula, $C_{\infty,a}=C_{a,\infty}=1$ y
$C_{a,b}=\chi(b-a)$ en las cinco posiciones finitas. Aquí
$\chi(1)=\chi(4)=1$, $\chi(2)=\chi(3)=-1$ y $\chi(0)=0$.
Es una realización en coordenadas de la estructura cuadrática ya
transportada; la identificación de las seis posiciones con la recta
proyectiva es un dato de carta, no una nueva entrada numérica en el
generador.

Las identidades de la suma cuadrática dan $C^2=5I_6$. Al reducir módulo
tres obtenemos
\begin{equation}
 A_W=
 \begin{pmatrix}
 0&1&1&1&1&1\\
 1&0&1&2&2&1\\
 1&1&0&1&2&2\\
 1&2&1&0&1&2\\
 1&2&2&1&0&1\\
 1&1&2&2&1&0
 \end{pmatrix},
 \qquad A_W^2=2I_6=-I_6.
 \label{exc:aw}
\end{equation}
La igualdad cuadrática realiza el régimen elíptico; no identifica el
cuerpo $\mathbb F_3$ con una realización arquimediana de la unidad
imaginaria. Son realizaciones distintas de una relación operatoria
común, con sus correspondientes dominios.

\Needspace{8\baselineskip}
\begin{definition}[Código de la carta dodecafásica]
\label{exc:codigo}
El código ternario asociado a \eqref{exc:aw} es
\[
 C_W=\{(w,wA_W):w\in\mathbb F_3^6\}
       \subset\mathbb F_3^{12}.
\]
Su conjunto de hexadas es
\[
 \mathcal H=
 \{\operatorname{supp}(c):c\in C_W,
                           \operatorname{wt}(c)=6\}.
\]
\end{definition}

\begin{proposition}[Autodualidad, distancia y diseño de Witt]
\label{exc:codigo-witt}
El código $C_W$ es autodual de parámetros $[12,6,6]_3$. Su enumerador
de pesos es
\begin{equation}
 \sum_{c\in C_W}z^{\operatorname{wt}(c)}
   =1+264z^6+440z^9+24z^{12}.
 \label{exc:enumerador}
\end{equation}
El conjunto $\mathcal H$ tiene $132$ elementos y constituye un sistema
$S(5,6,12)$: cada conjunto de cinco posiciones está contenido en una
única hexada.
\end{proposition}

\begin{proof}
La matriz generadora $[I_6\mid A_W]$ tiene rango seis y, por simetría
de $A_W$, satisface
\[
 [I_6\mid A_W][I_6\mid A_W]^{\mathsf T}
   =I_6+A_W^2=0.
\]
El código es autoortogonal de dimensión mitad de la longitud y, por
tanto, autodual. La evaluación finita de las $3^6$ palabras da
\eqref{exc:enumerador}; es una enumeración sobre la matriz explícita,
sin valores reales ni objetivos externos. En particular, la distancia
mínima es seis.

Si dos palabras de peso seis tienen el mismo soporte, entre sus seis
cocientes de coordenadas no nulas uno de los signos $1,-1$ aparece al
menos tres veces. Restando el múltiplo correspondiente se obtendría
una palabra no nula de peso a lo sumo tres, imposible. Luego un soporte
corresponde exactamente a la pareja $\{c,-c\}$ y hay $264/2=132$
soportes. Si dos soportes distintos compartieran cinco posiciones,
el mismo argumento cancelaría al menos tres coordenadas de una unión
de tamaño siete, produciendo peso a lo sumo cuatro. Por ello una
quintupla pertenece a lo sumo a una hexada. Finalmente,
\[
 132\binom65=792=\binom{12}5,
\]
de modo que pertenece a una y sólo una.
\end{proof}

El reconocimiento posterior es el código de Golay ternario extendido y
el pequeño diseño de Witt \cite{conwaysloane}. La prueba anterior explica qué objeto se
reconoce. No reemplaza la genealogía APP--TRIT--TPK por una búsqueda
entre códigos conocidos.

\subsection{Bandera dodecafásica asociada a \texorpdfstring{$K$ y $\alpha$}{K y alfa}}
\label{exc:bandera}

\input{sections/incidencia_registro.tex}

En el marco heredado del registro dodecafásico, los datos son
\begin{align}
 K={}&(234,543,140,729,659,824,621,\mathtt{058},914,794,146,601),
 \label{exc:k}\\
 w_\pi={}&010211, &w_e={}&201101,\nonumber\\
 w_\varphi={}&121200, &w_{\rm APP}={}&022020.\nonumber
\end{align}
Al aplicar $\gamma_{A_W}$ se obtienen las siguientes incidencias:
\begin{align*}
 P_\pi&=\{2,4,5,6,7,8,9,10,11\},\\
 H_e&=\{1,3,4,6,9,10\},\\
 H_\varphi&=\{1,2,3,4,7,9\},\\
 H_{\rm APP}&=\{2,3,5,10,11,12\}.
\end{align*}
La primera es el soporte de una palabra de peso nueve; las otras tres
son hexadas. El registro aporta, por su lectura de bloque alto,
\begin{equation}
 B_K=\{i:K_i\geq729\}=\{4,6,9,10\}=H_e\cap P_\pi.
 \label{exc:bk}
\end{equation}
El entero $729$ señala aquí una frontera del formato de bloques
declarado. No se identifica por el mero cardinal con el ambiente
ternario usado en \eqref{exc:codigo}. En este registro concreto, todos
los umbrales enteros entre $660$ y $729$ producen el mismo soporte:
el soporte no recupera el umbral ni los doce bloques.

La vía aritmética de $\alpha$ conserva, antes del acarreo, el vector
firmado
\begin{equation}
 Z_0=(7,297,353,-430,-715,-1199,-3,286,-894,-528,380,663).
 \label{exc:precarry}
\end{equation}
Su soporte negativo es
\begin{equation}
 H^-_\alpha=\{4,5,6,7,9,10\}.
 \label{exc:halpha}
\end{equation}
No se extraen estos signos de la expansión decimal del escalar
$\alpha$; se conservan en el estado que precede a su expresión numérica.

\input{sections/incidencia_normalizacion.tex}

\begin{proposition}[Coincidencia de los dos selectores de la bandera]
\label{exc:selector-bandera}
Dentro de $\mathcal H$, el soporte \eqref{exc:halpha} es la única hexada
$H$ que satisface
\[
 B_K\subset H\subset P_\pi,
 \qquad
 (|H\cap H_e|,|H\cap H_\varphi|,|H\cap H_{\rm APP}|)=(4,3,2).
\]
Así, la lectura firmada y la lectura incidencial seleccionan la misma
bandera $B_K\subset H^-_\alpha$.
\end{proposition}

\begin{proof}
Las cuatro hexadas que contienen $B_K$ se obtienen añadiendo,
respectivamente, los pares
\[
 \{1,3\},\qquad\{2,11\},\qquad\{5,7\},\qquad\{8,12\}.
\]
Sólo los pares segundo y tercero están contenidos en $P_\pi$. Para el
segundo, la intersección con $H_\varphi$ tiene tamaño tres y con
$H_{\rm APP}$ tamaño tres; no cumple la última condición. Para el
tercero, las tres intersecciones tienen tamaños $4,3,2$. Este último
produce precisamente el soporte negativo de \eqref{exc:precarry}.
\end{proof}

Conviene conservar tanto la bandera como su marco marcado:
\[
 \mathcal I_*=(B_K\subset H^-_\alpha),\qquad
 \widetilde{\mathcal I}_*
   =(\mathcal I_*;P_\pi,H_e,H_\varphi,H_{\rm APP}).
\]
Las operaciones de unión, intersección y diferencia son equivariantes
bajo un relabelado simultáneo del diseño. La clase marcada, y no los
números de las posiciones tomados aisladamente, es el objeto transportado.

\begin{lemma}[Origen seleccionado por la incidencia marcada]
\label{exc:origen}
Para el marco anterior,
\[
 o_*=(H_e\cap H_\varphi)
           \setminus(H^-_\alpha\cup H_{\rm APP})=\{1\}.
\]
La regla es equivariante bajo todo automorfismo aplicado simultáneamente
a los datos de $\widetilde{\mathcal I}_*$.
\end{lemma}

\begin{proof}
La primera intersección es $\{1,3,4,9\}$. La unión que se sustrae
contiene $3,4,9$ y no contiene $1$. La equivariancia se sigue de que una
biyección conmuta con las operaciones de conjuntos.
\end{proof}

Este origen seleccionará una dirección del vecino reticular cuando se
declaren la métrica $A_2$ y la cámara positiva. La incidencia por sí sola
no contiene esa métrica. El paso reticular conservará explícitamente
ambos datos.

\input{sections/incidencia_bandera.tex}

\subsection{Entrelazamiento isométrico de las representaciones de incidencia}
\label{exc:entrelazador}

La incidencia define un mapa lineal cuya acción puede comprobarse
posición por posición:
\begin{equation}
 I:\mathbb R^{12}\longrightarrow\mathbb R^{\mathcal H},
 \qquad (Iv)(H)=\sum_{p\in H}v_p.
 \label{exc:incidencia}
\end{equation}
Sea $V_{11}=\mathbf1^\perp\subset\mathbb R^{12}$. Las dos realizaciones
euclídeas son aquí lenguajes posteriores para el mismo diseño finito.

\begin{proposition}[Isometría de la pantalla dodecafásica]
\label{exc:isometria}
Se cumple
\[
 I^{\mathsf T}I=36I_{12}+30J_{12},
\]
donde $J_{12}$ es la matriz cuyas entradas son uno. En consecuencia,
\[
 U_W=\frac16 I\big|_{V_{11}}:
 V_{11}\xrightarrow{\ \sim\ }E_{\rm inc}:=I(V_{11})
\]
es una isometría equivariante para las acciones por permutación de los
automorfismos del diseño.
\end{proposition}

\begin{proof}
Cada punto pertenece a
$r=\binom{11}{4}/\binom54=66$ hexadas, y cada pareja a
$\lambda_2=\binom{10}{3}/\binom43=30$. Son, respectivamente, las
entradas diagonales y no diagonales de $I^{\mathsf T}I$. Sobre
$\mathbf1^\perp$, el término $J_{12}$ se anula y queda $36I$.
Por último, si $g$ es un automorfismo, $p\in H$ equivale a
$g(p)\in g(H)$; por ello $I\rho_{12}(g)=\rho_{\mathcal H}(g)I$.
\end{proof}

La información de incidencia también tiene una lectura espectral
concreta. Definamos
\[
 G_{H,H'}=\binom{|H\cap H'|}{4}.
\]

\Needspace{8\baselineskip}
\begin{lemma}[Acción del núcleo de intersección]
\label{exc:espectro}
Si $J_{132,12}$ es la matriz rectangular de unos,
\[
 GI=30I+15J_{132,12}.
\]
Por tanto, $G$ actúa por el escalar $30$ en $E_{\rm inc}$.
\end{lemma}

\begin{proof}
La entrada $(H,p)$ de $GI$ cuenta parejas $(T,H')$ con
$T\subset H\cap H'$, $|T|=4$ y $p\in H'$. Si $p\in H$, hay diez
cuaternas $T$ que contienen $p$, cada una contenida en cuatro hexadas,
y cinco que no lo contienen, cada una con una única extensión de la
quintupla $T\cup\{p\}$. El total es $10\cdot4+5=45$.
Si $p\notin H$, las quince cuaternas de $H$ dan una extensión cada
una: el total es $15$. De ahí la identidad. El término rectangular
se anula en $V_{11}$.
\end{proof}

Aquí sí existe un entrelazador material: tiene dominio, codominio,
entradas y normalización determinados. Por ejemplo, cualquier
proyector ortogonal $P$ ya construido sobre $V_{11}$ se transporta a
$Q=U_WPU_W^*$ sobre $E_{\rm inc}$, con
$U_WP=QU_W$. Esta afirmación no identifica por su sola traza a todos
los operadores que actúan en esos espacios.

\input{sections/k_direccion_dimensional.tex}

\subsection{La estrella de Witt y la generación de \texorpdfstring{$M_{12}$}{M12}}
\label{exc:estrella}

\begin{definition}[Estrella sobre una cuaterna]
\label{exc:def-estrella}
Para $B\in\binom{[12]}4$, la estrella de Witt es la familia de las
cuatro hexadas que contienen $B$. Sus pétalos son los pares $H\setminus B$.
La permutación $\sigma_B$ fija $B$ e intercambia los dos puntos de cada
pétalo.
\end{definition}

La existencia de cuatro hexadas se deduce de
\[
 \lambda_4=\frac{\binom81}{\binom21}=4.
\]
Dos pétalos no pueden compartir un punto, pues entonces dos hexadas
contendrían la misma quintupla. Son, por tanto, cuatro pares disjuntos
que particionan $[12]\setminus B$. Esta descripción no es una estrella
de cinco puntas ni una metáfora geométrica: es una incidencia
cuatro-a-uno del diseño.

\begin{figure}[htbp]
\centering
\begin{tikzpicture}[x=1cm,y=1cm, every node/.style={font=\small}]
 \node[draw,rounded corners,minimum width=3.8cm,minimum height=.7cm]
  (b) at (0,0) {$B_K=\{4,6,9,10\}$};
 \node[draw,minimum width=2.7cm,minimum height=.65cm]
  (p1) at (-4.6,-1.7) {$\{1,3\}$};
 \node[draw,minimum width=2.7cm,minimum height=.65cm]
  (p2) at (-1.55,-1.7) {$\{2,11\}$};
 \node[draw,very thick,minimum width=2.7cm,minimum height=.65cm]
  (p3) at (1.55,-1.7) {$\{5,7\}$};
 \node[draw,minimum width=2.7cm,minimum height=.65cm]
  (p4) at (4.6,-1.7) {$\{8,12\}$};
 \draw (b) -- (p1); \draw (b) -- (p2);
 \draw[very thick] (b) -- (p3); \draw (b) -- (p4);
 \node at (-4.6,-2.35) {$H_e$};
 \node at (-1.55,-2.35) {$B_K\cup\{2,11\}$};
 \node at (1.55,-2.35) {$H^-_\alpha$};
 \node at (4.6,-2.35) {$B_K\cup\{8,12\}$};
 \node at (0,-3.05)
  {$\sigma_{B_K}=(1\ 3)(2\ 11)(5\ 7)(8\ 12)$};
\end{tikzpicture}
\caption{La estrella concreta de la bandera dodecafásica. Cada línea
indica la unión de la cuaterna central con el par de su extremo; la
línea gruesa señala la hexada seleccionada también por los signos
anteriores al acarreo de $\alpha$. Los cuatro pares agotan las ocho
posiciones exteriores.}
\label{exc:fig-estrella}
\end{figure}

\begin{proposition}[La familia de estrellas]
\label{exc:m12}
Las $\binom{12}4=495$ permutaciones $\sigma_B$ preservan $\mathcal H$.
El grupo que generan tiene orden $95040$ y es el grupo de Mathieu
$M_{12}$ en su acción sobre el pequeño diseño de Witt.
\end{proposition}

\begin{proof}
La afirmación de preservación es una comprobación finita sobre las
$132$ hexadas de \eqref{exc:codigo}: para cada cuaterna se construyen
los cuatro pétalos y se verifica
$\{\sigma_B(H):H\in\mathcal H\}=\mathcal H$.
La clausura por composición de estas permutaciones produce $95040$
elementos. En la enumeración lexicográfica basta añadir sucesivamente
seis estrellas no contenidas aún en el grupo; los órdenes intermedios
son $2,6,18,360,7920,95040$, y las $495$ estrellas pertenecen al grupo
final. El procedimiento y los seis generadores están explicitados en
la nota técnica de procedencia. Finalmente se aplica el reconocimiento
del grupo de automorfismos de $S(5,6,12)$ como $M_{12}$, cuyo orden es
$12\cdot11\cdot10\cdot9\cdot8=95040$.
\end{proof}

Una estrella individual genera un grupo cíclico de orden dos. No
genera por sí sola $M_{12}$, y menos aún el Monstruo. Tampoco la
existencia de $495$ generadores de un grupo de permutaciones prueba que
cada uno sea el transporte de un lazo específico del TPK. Esa afirmación
requiere el mapa de lazos correspondiente; el resultado aquí establecido
es la acción incidencial finita.

\begin{remark}[La igualdad de trazas no identifica operadores]
Sobre $V_{11}$, una estrella tiene espectro $+1$ con multiplicidad siete
y $-1$ con multiplicidad cuatro. Su traza normalizada es $3/11$.
Un proyector de rango tres en el mismo espacio tiene también traza
normalizada $3/11$, pero espectro $1^3,0^8$. No son conjugados: sus
polinomios mínimos y sus espectros son diferentes. El entrelazador
\eqref{exc:incidencia} transporta operadores definidos; no convierte
esta coincidencia escalar en una equivalencia de operadores.
\end{remark}

\subsection{El residuo binario y las cinco regiones de \texorpdfstring{$\pi$}{pi}}
\label{exc:cinco-regiones}

El siguiente reconocimiento utiliza el código de Golay binario
extendido $G_{24}$, de parámetros $[24,12,8]_2$. Sus pesos son
$0,8,12,16,24$, y los soportes de peso ocho forman el diseño
$S(5,8,24)$. Fijamos una dodecada $D$, soporte de una palabra de peso
doce. Este dato y el isomorfismo de incidencia que se declarará a
continuación pertenecen a la carta de realización; no se usan para
generar ni seleccionar los decimales de $\pi$.

\begin{proposition}[Residuo de una dodecada y elevación de hexadas]
\label{exc:residuo-binario}
El conjunto
\[
 \mathcal H(D)=
 \{O\cap D: O\text{ es una octada},\ |O\cap D|=6\}
\]
es un sistema $S(5,6,12)$ sobre $D$. Cada una de sus hexadas tiene
una única elevación a una octada. Para cualquier cuaterna $B\subset D$,
las cinco octadas que la contienen se distribuyen en cuatro elevaciones
residuales y una única octada $O_B^\perp$ con $O_B^\perp\cap D=B$.
\end{proposition}

\begin{proof}
Una quintupla $A\subset D$ pertenece a una única octada $O$. Si
$j=|O\cap D|$, la palabra binaria suma de $O$ y $D$ tiene peso
$20-2j$. Entre los pesos permitidos, y con $0\leq j\leq8$, sólo
son posibles $j=2,4,6$. Como $A\subset O\cap D$, se tiene $j\geq5$
y por tanto $j=6$. Esto da existencia y unicidad de la hexada
residual, así como unicidad de su elevación al escoger cinco de sus
puntos. Por otra parte,
\[
 \lambda_4(S(5,8,24))=\frac{\binom{20}1}{\binom41}=5,
 \qquad
 \lambda_4(S(5,6,12))=4.
\]
Las cuatro hexadas residuales sobre $B$ dan cuatro octadas distintas.
La quinta no puede cortar $D$ en seis puntos, pues sería otra hexada
residual; como contiene $B$, su intersección tiene tamaño cuatro y
coincide con $B$.
\end{proof}

Una carta $\ell:[12]\to D$ que preserve las hexadas identifica
$\mathcal H$ con $\mathcal H(D)$. Su existencia es el reconocimiento
del diseño ya construido como el pequeño diseño de Witt; una elección
concreta de $\ell$ fija el relabelado. La elevación
\[
 \iota_D:\mathcal H\longrightarrow
       \{\text{octadas de }G_{24}\},
 \qquad H\longmapsto O\text{ tal que }O\cap D=\ell(H),
\]
es inyectiva porque la intersección con $D$ recupera $\ell(H)$.
Esta inyectividad tiene por dominio las hexadas, no el registro $K$.

La relación con las cinco regiones de $\pi$ conserva datos más ricos
que el número cinco. Las emisiones y sus firmas son:
\begin{center}
\small
\begin{tabular}{@{}crll@{}}
\toprule
Región & Multiplicidad & Emisión $U_6$ & Firma \\
\midrule
$\perp$ & $432$ & $(501,614,498,169,272,272)$ & $555555$ \\
$++_1$ & $144$ & $(810,923,870,169,272,272)$ & $339555$ \\
$++_2$ & $144$ & $(810,983,810,169,272,272)$ & $393555$ \\
$++_3$ & $144$ & $(870,923,810,169,272,272)$ & $933555$ \\
$--$ & $144$ & $(870,923,810,418,674,521)$ & $933717$ \\
\bottomrule
\end{tabular}
\end{center}
En cada bloque decimal de tres cifras, la firma usa la diferencia
modular de las dos primeras cifras. Las cinco emisiones tienen la
misma palabra trítica visible $w_\pi=010211$, pero no la misma
genealogía regional. Las tres regiones positivas conservan la cola
$169|272|272$; la región negativa la cambia por $418|674|521$;
la región transversal tiene firma constante y multiplicidad mayor.

La extensión marcada de la carta regional envía la región transversal
a $O_{\ell(B_K)}^\perp$, la negativa a $\iota_D(H^-_\alpha)$ y las
tres positivas a las otras tres elevaciones de la estrella. Así se
realiza la partición $4+1$ de las cinco octadas sobre $\ell(B_K)$.
La correspondencia individual entre los tres nombres positivos y
los tres pétalos restantes exige un orden de carta; no la determina
la incidencia sin ese dato. Por ello la afirmación precisa es una
correspondencia de regiones marcada, compatible con sus roles y con
la bandera, no una biyección canónica deducida del cardinal cinco.

La identidad $432+4\cdot144=1008$ conserva las multiplicidades
regionales. También $1008=36\cdot28$ y una octada tiene
$\binom82=28$ parejas. Estas igualdades son controles aritméticos de
catálogo, no sustitutos del mapa $\iota_D$ ni de la carta regional.

\input{sections/estrella_pentadica_recuperada.tex}

\subsection{Del código ternario al retículo de Niemeier}
\label{exc:niemeier}

El paso a dimensión veinticuatro no necesita recorrer primero el
código binario. Se realiza directamente desde $C_W$ por pegado de
discriminantes. Así quedan distinguidas dos prolongaciones compatibles:
una conserva el residuo $W_{12}\subset W_{24}$; la otra construye el
retículo sobre el que se realizará el álgebra de operadores de vértice.

Sea $A_2=\mathbb Z a_1\oplus\mathbb Z a_2$ con matriz de Gram
\[
 \begin{pmatrix}2&-1\\-1&2\end{pmatrix},
 \qquad \rho=a_1+a_2,\qquad \rho^2=2.
\]
Representamos las tres clases de $A_2^*/A_2$ mediante
\[
 \varpi_0=0,\qquad
 \varpi_1=\frac{2a_1+a_2}{3},\qquad
 \varpi_2=\frac{a_1+2a_2}{3}.
\]
La aplicación $j\mapsto\varpi_j+A_2$ identifica aditivamente
$\mathbb F_3$ con ese discriminante. Para $c\in C_W$, escribimos
$\phi(c)=(\varpi_{c_1},\ldots,\varpi_{c_{12}})$ y definimos
\begin{equation}
 N(C_W)=\bigcup_{c\in C_W}\bigl(A_2^{12}+\phi(c)\bigr).
 \label{exc:niemeier-def}
\end{equation}

\begin{theorem}[Pegado ternario]
\label{exc:niemeier-teorema}
El objeto \eqref{exc:niemeier-def} es un retículo positivo, par y
unimodular de rango $24$. Su sistema de raíces es exactamente
$A_2^{12}$; por tanto, se reconoce como el retículo de Niemeier con
ese sistema de raíces.
\end{theorem}

\begin{proof}
La linealidad del código hace que la unión de clases sea estable por
suma. La autoortogonalidad anula el emparejamiento discriminante:
éste es un múltiplo de $c\cdot d/3$ módulo $\mathbb Z$. Así, el
retículo es integral. La norma de una clase de pegado es
$2\operatorname{wt}(c)/3$ módulo $2\mathbb Z$. Los pesos de
\eqref{exc:enumerador} son múltiplos de tres, de modo que es par.

El índice de $A_2^{12}$ en $N(C_W)$ es $|C_W|=3^6$. En consecuencia,
\[
 \det N(C_W)=\frac{3^{12}}{(3^6)^2}=1.
\]
La positividad procede de la suma ortogonal de las doce métricas $A_2$.
En una clase no nula de $A_2^*/A_2$ la norma mínima es $2/3$;
toda palabra no nula tiene al menos seis posiciones no nulas.
Por tanto, un vector de una clase de pegado no trivial tiene norma
al menos $6\cdot2/3=4$. No aparecen raíces de norma dos fuera de
$A_2^{12}$. Dentro de éste, las raíces son precisamente las de sus
doce componentes.
\end{proof}

\subsection{Selección incidencial del vecino unimodular sin raíces}
\label{exc:leech}

\input{sections/incidencia_reticulo.tex}

La etiqueta $o_*=1$ de la bandera se realiza ahora como una de las
doce componentes $A_2$. Esta operación declara la métrica anterior y
la cámara radial positiva
\[
 v(a)=\sum_{i=1}^{12}a_i\rho_i,
 \qquad a_i=1+3b_i,\quad b_i\in\mathbb Z_{\geq0}.
\]
El requisito de paridad del vecino de índice tres es
$v(a)^2\equiv0\pmod{18}$. Puesto que
\[
 v(a)^2=2\sum_i(1+3b_i)^2,
\]
esa congruencia equivale a $\sum_i b_i\equiv1\pmod3$. El menor
valor de la norma en la cámara se obtiene con un único $b_i=1$ y
todos los demás nulos; es $54$. El origen marcado selecciona
\begin{equation}
 v_\alpha=4\rho_1+\rho_2+\cdots+\rho_{12},
 \qquad v_\alpha^2=54.
 \label{exc:vector}
\end{equation}
La palabra ``menor'' se refiere a esta cámara positiva. No afirma
minimalidad en toda la clase residual: el vector
$(a_1-2a_2,\rho,\ldots,\rho)$ representa la misma clase módulo tres
y tiene norma $14+11\cdot2=36$. La restricción de cámara es parte
de los datos de la construcción, no una conclusión omitible.

Sea $N=N(C_W)$. Definimos
\begin{equation}
 N_{\alpha,0}=\{x\in N:\langle x,v_\alpha\rangle\equiv0\pmod3\},
 \qquad
 N_\alpha=N_{\alpha,0}+\mathbb Z\frac{v_\alpha}{3}.
 \label{exc:vecino}
\end{equation}

\begin{theorem}[Vecino de Leech determinado por la bandera marcada]
\label{exc:teorema-leech}
Con la métrica y la cámara declaradas, $N_\alpha$ es un retículo
positivo, par y unimodular de rango $24$, sin vectores de norma dos.
Por el teorema de clasificación correspondiente, es isométrico al
retículo de Leech $\Lambda_{24}$.
\end{theorem}

\begin{proof}
El carácter $x\mapsto\langle x,v_\alpha\rangle\bmod3$ no es nulo:
una raíz simple de una componente se empareja con $a_i\rho_i$ en
$a_i\equiv1\pmod3$. Por ello $[N:N_{\alpha,0}]=3$ y
$\det N_{\alpha,0}=9$. Además,
$v_\alpha/3\in N_{\alpha,0}^*$ y $v_\alpha\in N_{\alpha,0}$.
Si $v_\alpha/3$ perteneciera a $N$, la integridad de $N$ haría que
todos los emparejamientos con $v_\alpha$ fueran divisibles por tres,
contradiciendo el carácter no nulo. La clase de $v_\alpha/3$ tiene
por tanto orden tres. Se deduce
\[
 [N_\alpha:N_{\alpha,0}]=3,
 \qquad \det N_\alpha=9/3^2=1.
\]
La extensión es integral, y para $x\in N_{\alpha,0}$,
\[
 \left(x+m\frac{v_\alpha}{3}\right)^2
   =x^2+2m\frac{\langle x,v_\alpha\rangle}{3}+6m^2
   \in2\mathbb Z.
\]
Luego es par.

Queda excluir las raíces, que es el paso decisivo. Las raíces de $N$
pertenecen a $A_2^{12}$. Sus emparejamientos con $v_\alpha$ son
congruentes con $\pm1$ o $\pm2$ módulo tres; ninguna pertenece a
$N_{\alpha,0}$. Para las otras dos clases del vecino, cada componente
pertenece a alguna de las clases locales
\[
 A_2+\varpi_j\pm\rho/3,\qquad j=0,1,2.
\]
Como $4\rho/3\equiv\rho/3\pmod{A_2}$, esta descripción vale
también en la componente distinguida. Escribiendo un vector local
como $(u a_1+v a_2)/3$, su norma es
$2(u^2-uv+v^2)/9$. Las seis clases anteriores tienen residuos no
nulos y mínimo $2/9$: los representantes de mínima forma cuadrática
son $(\pm1,0),(0,\pm1),(1,1),(-1,-1)$ en los residuos apropiados.
Un vector de una clase nueva tiene, por tanto, norma al menos
$12\cdot2/9=8/3>2$. No hay raíces en ninguna de las tres clases.
La clasificación de los retículos pares unimodulares positivos de
rango veinticuatro identifica el único caso sin raíces con
$\Lambda_{24}$.
\end{proof}

El argumento no reconoce Leech por una coincidencia de dimensión.
Construye el pegado, prueba integridad, paridad y determinante, y
elimina las raíces mediante una cota local. Sólo entonces aplica el
teorema de clasificación. La dependencia de $\alpha$ se conserva en
la bandera que selecciona $o_*$ y en el vector \eqref{exc:vector};
no se convierte el valor escalar de la constante en una longitud del
retículo.

\subsection{El álgebra de operadores de vértice y el Monstruo}
\label{exc:voa}

Sea $\Lambda=N_\alpha$, ya reconocido como Leech. El paso siguiente
utiliza su retículo integral, no una expansión decimal. La construcción
de álgebra de operadores de vértice de retículo se escribe
\[
 V_\Lambda=M(1)\otimes\mathbb C_\varepsilon[\Lambda].
\]
Aquí $M(1)$ es el espacio de osciladores de Heisenberg en rango
veinticuatro; $\mathbb C_\varepsilon[\Lambda]$ es el álgebra de grupo
torcida por el cociclo de la extensión central del retículo. El
cociclo y el levantamiento de las isometrías forman parte de la
construcción: una suma vectorial desnuda no determina los productos
de vértice.

La negación $\lambda\mapsto-\lambda$ admite el levantamiento estándar
$\theta$. Con su módulo torcido y la elección de paridad compatible,
la construcción de Frenkel--Lepowsky--Meurman produce
\begin{equation}
 V^\natural_{\rm HMT}
   =V_\Lambda^+\oplus(V_\Lambda^T)^+.
 \label{exc:flm}
\end{equation}
El signo $+$ denota la parte par de la acción levantada. La expresión
incluye el producto de vértice del orbifold; no se limita a una suma
de espacios graduados.

\begin{theorem}[Prolongación hasta el módulo Moonshine]
\label{exc:moonshine}
El álgebra \eqref{exc:flm}, construida sobre el vecino
\eqref{exc:vecino}, es una realización del álgebra Moonshine
$V^\natural$. Tiene carga central $24$, componente de peso uno nula
y grupo de automorfismos isomorfo al Monstruo $\mathbb M$. Su carácter
graduado es
\begin{equation}
 \operatorname{Tr}_{V^\natural_{\rm HMT}}q^{L_0-1}
   =J(\tau)=j(\tau)-744
   =q^{-1}+196884q+21493760q^2+\cdots,
 \qquad q=e^{2\pi i\tau}.
 \label{exc:j}
\end{equation}
\end{theorem}

\begin{proof}
El teorema \ref{exc:teorema-leech} verifica las hipótesis reticulares:
positividad, paridad, unimodularidad, rango veinticuatro y ausencia de
raíces. Se aplica entonces la construcción y el teorema de
Frenkel--Lepowsky--Meurman para el retículo de Leech. Éstos identifican
el orbifold con $V^\natural$, establecen su carácter y su grupo de
automorfismos. La composición de la bandera, el vecino y el functor
de construcción de VOA da la realización aquí indicada. No se deduce
el grupo a partir del coeficiente $196884$.
\end{proof}

El teorema utilizado es un reconocimiento clásico posterior con
hipótesis verificadas, no una nueva demostración en este artículo de
la construcción FLM \cite{flm}.
Además, la notación $q=e^{2\pi i\tau}$ es la variable convencional
de Fourier empleada para reconocer el carácter. No introduce $\pi$
ni $e$ como entradas del generador coinductivo que ya los ha determinado.

Puede verse una parte del mecanismo de traza antes del orbifold.
Sólo el vector reticular cero queda fijo por la negación; cada
oscilador contribuye con signo alterno. De ahí
\[
 \operatorname{Tr}_{V_\Lambda}\theta q^{L_0-1}
   =t_2(\tau)
   :=q^{-1}\prod_{n\geq1}(1+q^n)^{-24}.
\]
Esta traza en el álgebra de retículo no debe confundirse con el
carácter sin inserción de $V^\natural$. Tampoco identifica por sí
sola una clase de conjugación del Monstruo: el espacio, la inserción
y el levantamiento deben especificarse antes de comparar funciones.

\input{sections/incidencia_automorfismos.tex}

\input{sections/fibras_app_witt.tex}
\input{sections/moonshine_comparacion.tex}

\input{sections/pantallas_dualidad.tex}

\subsection{Trazas graduadas de Moonshine y representación del Monstruo}
\label{exc:alcance}

Para $g\in\operatorname{Aut}(V^\natural_{\rm HMT})$, la traza
con inserción es
\[
 T_g(\tau)=
 \operatorname{Tr}_{V^\natural_{\rm HMT}}gq^{L_0-1}.
\]
El teorema de Moonshine monstruoso identifica estas series como
Hauptmoduln de los grupos de género cero correspondientes
\cite[teorema~1.1]{borcherds}.
El carácter $J$, el grupo $\mathbb M$ y la familia $g\mapsto T_g$
son realizaciones del álgebra ya construida. No forman una sucesión
causal en la que primero se adivina $J$ y después se deduce el Monstruo.

La unidad tiene aquí una localización algebraica precisa. El subespacio
de peso cero es la línea del vacío, \(V^\natural_0=\CC\mathbf1\), y
el coeficiente de \(q^{-1}\) de \(J\) es por ello \(1\). En peso uno
la dimensión es cero, de donde procede la ausencia del término constante
en \(J=j-744\). Un morfismo de álgebras de vértices conserva el vector
vacío. Esta propiedad proporciona el sentido técnico de conservación
de la unidad en este nivel; la unidad de los observables cilíndricos
se conserva por la precomposición demostrada anteriormente. La relación
entre ambos niveles debe seguir los mapas de construcción con sus unidades.

En peso dos aparece la igualdad
\[
 196884=1+196883=196560+324=54(5\cdot729+1).
\]
La primera separación corresponde a la línea conforme trivial y al
componente irreducible de dimensión $196883$ de la representación
del Monstruo. Las otras expresiones tienen lecturas reticulares o
aritméticas. Su igualdad no identifica automáticamente sus sumandos
con regiones, palabras o rutas del TPK. Para hacerlo habría que
presentar los mapas entre esos espacios y probar que respetan las
acciones relevantes.

En particular, una isometría $N_\alpha\simeq\Lambda_{24}$ induce
una realización de \eqref{exc:flm}, y una elección de isomorfismo
con $V^\natural$ transporta su grupo de automorfismos. No se ha
construido con ello una acción del Monstruo sobre los dígitos de
$\pi$, sobre los bloques de $K$ o sobre todas las historias TPK.
Tal afirmación adicional requeriría una representación
$\rho_{\rm TPK}$ y un mapa $F$ con un cuadrado explícito
\[
 F\rho_{\rm TPK}(g)=\rho_{V^\natural}(g)F,
\]
incluyendo dominio, codominio e información conservada. La ausencia
de ese paso en la afirmación presente no retira el entrelazador
incidencial $U_W$, la bandera ni el vecino ya construidos; sólo
delimita una promoción distinta, de acción sobre el estado fuente.

La cadena establecida en esta sección puede resumirse, manteniendo
visibles sus datos de transporte, como
\begin{align*}
 &\text{APP--TRIT--TPK y estado aritmético con memoria de trayectoria}
   \longrightarrow (w_j,f_j;K,Z_0)\\
 &\longrightarrow (C_W,\mathcal H;B_K\subset H^-_\alpha,
                   P_\pi,H_e,H_\varphi,H_{\rm APP})\\
 &\longrightarrow \bigl(N(C_W),o_*,\text{métrica y cámara}\bigr)
   \longrightarrow N_\alpha
   \longrightarrow V^\natural_{\rm HMT}.
\end{align*}
En el nivel de incidencia, las estrellas generan $M_{12}$ y la
carta residual realiza la estructura regional $4+1$ en $W_{24}$.
En el nivel reticular y de VOA, la construcción da Leech y el módulo
Moonshine. Son prolongaciones de una misma fuente marcada, con
objetos y mapas diferentes. La continuidad causal no obliga a
identificar sus codominios ni a borrar las elecciones explícitas
de marco, cámara y levantamiento.

El resultado de la sección es, por tanto, una derivación reunida:
las pruebas finitas y reticulares hacen efectivo el tránsito desde
la incidencia dodecafásica hasta las hipótesis de los teoremas de
reconocimiento. La conclusión comprende las estructuras y los
transportes especificados en las proposiciones precedentes, con sus
dominios y elecciones explícitos. La comparación de constantes y la
validación experimental corresponden a afirmaciones distintas de esta
construcción algebraica.


% SOURCE foundation/sections/revision_incidencia_argumento.tex
El argumento comienza con la compatibilidad de prefijos y marcos,
construye el código ternario y determina la bandera mediante \(K\) y el
registro firmado de \(\alpha\). La incidencia marcada selecciona después
el origen reticular y el vecino de Leech. El último paso realiza el álgebra
de vértices y sus automorfismos. En cada etapa se precisará qué estructura
se conserva: combinaciones lineales, soportes, clases discriminantes, forma
bilineal o producto graduado. Las explicaciones de la función de \(K\),
de la normalización y del origen de norma \(54\) acompañan sus respectivas
construcciones en las secciones \ref{exc:bandera} y \ref{exc:leech}.


% SOURCE foundation/sections/incidencia_registro.tex
\paragraph{Articulación de las coordenadas regionales y la incidencia de clausura.}
\label{inc-ampl-articulacion}
La función de esta prolongación se comprende a partir de una distinción
entre tres operaciones: conservar un registro, evaluar una coordenada y
extraer una relación de incidencia. Las tres actúan sobre información
producida por APP--TRIT--TPK. Su articulación hace explícita la tesis
estructural del artículo: una realización numérica posee una procedencia
operatoria, y determinados datos de esa procedencia admiten, además, una
realización combinatoria y reticular. Las ecuaciones permiten seguir ambos
recorridos con sus dominios y con la información que permanece en cada uno.

El punto de partida es el estado terminal con sus prolongaciones
compatibles. En él se conservan los residuos y cocientes de las dos hojas
APP, el régimen y la orientación TRIT, y el transporte TPK con su memoria.
Las regiones de clausura, propagación y autoescala se determinan en ese
dominio antes de que sus lectores establezcan las coordenadas
$\pi,e,\varphi$. Al mismo tiempo, el registro orientado de las doce
ventanas determina el vector firmado $U$. La aparición posterior de $K$
procede de la transformación integral demostrada en
\S\ref{subsec:k-hadamard}:
\begin{equation}
 U=\mathcal H_{12}K,\qquad
 K=\frac14\mathcal H_{12}U,\qquad
 \mathcal H_{12}^{2}=4I_{12}.
 \label{inc-ampl-registro}
\end{equation}
La congruencia que define $\mathcal L_H$ garantiza la integralidad de
la inversa. De este modo, $K$ constituye una coordenatización reversible
del registro firmado en la subred pertinente. Sus posiciones conservan el
orden de fase; sus diferencias $D_3K,D_4K$ conservan las relaciones
transversales; la suma $Q(K)$ determina la componente uniforme. Esa
estructura explica por qué el registro interviene tanto en la composición
numérica como en la selección de una configuración de frontera.

La evaluación periódica
\[
 K\longmapsto\kappa_{\mathrm{per}}
 =\frac{\sum_{m=1}^{12}K_m1000^{12-m}}{1000^{12}-1}
\]
mantiene recuperables los doce bloques cuando se fijan base, longitud y
origen de fase. La extracción de un soporte realiza una operación
diferente: conserva la pertenencia de cada posición a una clase definida
por un umbral y agrupa los registros que tienen el mismo soporte. La
distinción entre estas dos aplicaciones resulta decisiva. La precisión
escalar de $\kappa_{\mathrm{per}}$ y la incidencia de $B_K$ son dos
contenidos matemáticos específicos, relacionados mediante el registro
completo del que proceden.



% SOURCE foundation/sections/incidencia_normalizacion.tex
\paragraph{La normalización de alfa y su configuración firmada.}
\label{inc-ampl-clausura}
La coordenada de $\alpha$ se determina en la vía posicional mediante la
composición de los bloques regionales y del registro dodecafásico. Antes
de la división posicional, la combinación orientada tiene componentes
$Z_j=P_j+E_j-\Phi_j-K_j^{\mathrm{per}}$ y el balance con cociente entrante
es $s_j=Z_j+c_{j+1}$. La normalización satisface
\begin{equation}
 A_j=s_j-1000c_j=Z_j+c_{j+1}-1000c_j,\qquad 0\leq A_j<1000.
 \label{inc-ampl-normalizacion}
\end{equation}
Esta identidad conserva simultáneamente el bloque normalizado y la
contribución del cociente transportado. El límite de los prefijos
compatibles determina la coordenada escalar de la clausura. La
configuración firmada conserva, por su parte, qué posiciones de la
combinación anterior a la normalización presentan defecto negativo.
En la ventana inicial, esa configuración es el vector $Z_0$ de
\eqref{exc:precarry}, y su soporte negativo es $H^-_\alpha$.

Así quedan definidos dos contenidos de una misma operación. La
normalización posicional determina una coordenada real mediante sus
restos y cocientes compatibles; el soporte negativo determina una
configuración de posiciones en el ciclo de doce componentes. El signo se
refiere al balance de registros anterior a la normalización. Su lectura
incidencial permanece disponible porque el estado conserva ese balance
junto con el resultado normalizado. Esta articulación corresponde a la
vía posicional demostrada; la comparación con los operadores analíticos
de vacancias conserva el alcance preciso establecido en el capítulo de
$\alpha$.


% SOURCE foundation/sections/incidencia_bandera.tex
\paragraph{De la codificación ternaria a la bandera marcada.}
\label{inc-ampl-bandera}
Las regiones aportan también los códigos $w_\pi,w_e,w_\varphi$ y
$w_{\rm APP}$ en el marco heredado. La aplicación
\[
 \gamma_{A_W}:w\longmapsto(w,wA_W)
\]
conserva la palabra de seis componentes como primera coordenada y añade
su transformación por $A_W$. La relación $A_W^2=-I_6$ fija la
compatibilidad cuadrática de esta realización. El paso siguiente,
$c\mapsto\operatorname{supp}(c)$, conserva las posiciones no nulas.
En las palabras de peso seis, el soporte identifica exactamente la pareja
$\{c,-c\}$, como prueba la distancia mínima del código. De esta forma,
el cociente por signo produce las hexadas, mientras el código completo
sigue disponible para las operaciones que requieren sus símbolos.

En esta carta, las dos construcciones de la cuaterna coinciden:
\begin{equation}
 \{i:K_i\geq729\}
 =H_e\cap P_\pi=B_K=\{4,6,9,10\}.
 \label{inc-ampl-cuaterna}
\end{equation}
El lado izquierdo procede de una comparación entera de los componentes
de $K$; el derecho, de la intersección de los soportes codificados de
propagación y clausura. El umbral pertenece a la carta de bloques
declarada y permanece explicitado en la definición del lector. El
registro íntegro conserva los valores sobre los que actúa, incluso cuando
distintos umbrales producen el mismo subconjunto. La igualdad expresa,
por tanto, una compatibilidad entre aplicaciones especificadas.

El soporte $H^-_\alpha$ completa esa cuaterna hasta una hexada. La
condición de pertenencia al soporte $P_\pi$, junto con los tamaños de
intersección $(4,3,2)$ respecto de $H_e,H_\varphi,H_{\rm APP}$,
selecciona una única hexada. Las pruebas siguientes reúnen las dos
determinaciones: la obtenida mediante el signo de $Z_0$ y la obtenida
mediante incidencia. El resultado se organiza como
\begin{equation}
 \widetilde{\mathcal I}_*
 =\bigl(B_K\subset H^-_\alpha;
         P_\pi,H_e,H_\varphi,H_{\rm APP}\bigr).
 \label{inc-ampl-marco}
\end{equation}
La inclusión es la bandera; los cuatro soportes adicionales constituyen
su marco marcado. Las permutaciones simultáneas transportan el marco y
conservan sus inclusiones e intersecciones. La información relevante reside
en esta configuración transportada y en su clase de isomorfía.

Las cinco regiones de $\pi$ recuperan aquí su diferenciación interna.
Comparten $w_\pi$ como código ternario visible, pero conservan emisiones,
firmas y multiplicidades distintas. En la realización binaria sobre una
dodecada marcada, la región transversal corresponde a la octada cuya
intersección con la dodecada es la cuaterna; las otras cuatro regiones
corresponden a las elevaciones de las cuatro hexadas de su estrella. La
región negativa queda distinguida por $H^-_\alpha$ y el orden de carta
individualiza las tres regiones positivas. Esta correspondencia conserva
los papeles regionales y la distribución $432+4\cdot144=1008$, además
de realizar la partición incidencial $1+4$.


% SOURCE foundation/sections/k_direccion_dimensional.tex
\subsection{La dirección seleccionada por el registro dodecafásico}
\label{exc:k-direccion}

El registro \(K\) tiene una función geométrica que complementa su
evaluación racional y su lectura de soportes. Una vez construida la
representación de las doce posiciones, sus componentes determinan una
dirección en el espacio de variaciones de media nula. La reducción
\(12\to11\) elimina el modo uniforme; la reducción siguiente
\(11\to10\) elimina la dirección que selecciona el propio registro.
Las dos operaciones tienen, por tanto, antecedentes diferentes.

Para declarar completamente la carta utilizada, consideremos el grupo
\(A_5\) de permutaciones pares de cinco objetos y el subgrupo
\(C_5=\langle(1\,2\,3\,4\,5)\rangle\). Las doce clases izquierdas
\(r_pC_5\) se ordenan mediante estos representantes, escritos como
listas de imágenes:
\[
\begin{array}{c|c@{\qquad}c|c}
p&r_p&p&r_p\\ \hline
1&(4,3,5,2,1)&7&(5,4,3,2,1)\\
2&(4,2,1,3,5)&8&(1,3,4,2,5)\\
3&(1,2,4,5,3)&9&(4,2,3,5,1)\\
4&(2,5,3,4,1)&10&(3,2,5,4,1)\\
5&(4,3,1,5,2)&11&(4,5,2,3,1)\\
6&(5,1,2,3,4)&12&(5,3,2,4,1).
\end{array}
\]
La acción por multiplicación izquierda define una representación
ortogonal \(\rho:A_5\to O(\RR^{12})\):
\(\rho(a)e_p=e_q\) cuando \(ar_pC_5=r_qC_5\).
Esta carta de representación actúa sobre el registro ya obtenido;
los caracteres del grupo no se utilizan para elegir sus doce valores.

Sean \(\chi_3\) los valores del carácter tridimensional
\[
\begin{array}{c|ccccc}
\text{clase}&1&2^2&3&5_a&5_b\\ \hline
\chi_3&3&-1&0&(1+\sqrt5)/2&(1-\sqrt5)/2.
\end{array}
\]
La clase \(5_a\) contiene el generador de \(C_5\) mostrado y
\(5_b\) contiene su cuadrado. Esta convención distingue los dos
sectores tridimensionales conjugados de la representación. Definimos
\begin{equation}
 P_3=\frac3{60}\sum_{a\in A_5}\chi_3(a^{-1})\rho(a),
 \qquad P_{11}=I_{12}-\frac1{12}J_{12}.
 \label{eq:k-proyectores-incidencia}
\end{equation}
Ambas matrices quedan determinadas por sumas finitas en
\(\QQ(\sqrt5)\). La aparición de este cuerpo cuadrático pertenece
a la realización de la representación, posterior a la generación
regional de \(\varphi\).

\begin{lemma}[Proyección del registro sobre el sector tridimensional]
\label{lem:k-sector-no-nulo}
Se tienen
\[
 P_3=P_3^{\mathsf T}=P_3^2,\quad
 P_3\mathbf1=0,\quad \operatorname{tr}P_3=3,
 \qquad P_3P_{11}=P_{11}P_3=P_3.
\]
Para el registro \eqref{exc:k}, el vector
\begin{equation}
 u_K=P_3P_{11}K
 \label{eq:k-vector-direccional}
\end{equation}
es no nulo. Su norma exacta es
\begin{equation}
 \lVert u_K\rVert^2
 =\frac{6638585+2275584\sqrt5}{20}>0.
 \label{eq:k-norma-direccion}
\end{equation}
\end{lemma}
\begin{proof}
La suma de carácter en \eqref{eq:k-proyectores-incidencia} es el
proyector isótipo. El carácter real y la sustitución \(a\mapsto a^{-1}\)
dan la simetría. La multiplicidad del sector en la acción sobre
\(A_5/C_5\) es la dimensión de sus vectores fijos por \(C_5\),
\[
 \frac15\left(3+2\frac{1+\sqrt5}{2}
                 +2\frac{1-\sqrt5}{2}\right)=1.
\]
Por ello el rango es tres. Su ortogonalidad al carácter trivial
da \(P_3\mathbf1=0\) y las dos composiciones con \(P_{11}\).
También pueden verificarse estas identidades multiplicando las
matrices de la suma finita.

Finalmente, \(\sum K_p=6263\) y la sustitución de las doce
componentes en esa misma suma da
\[
 \lVert u_K\rVert^2=K^{\mathsf T}P_3K
 =\frac1{20}\sum_{a\in A_5}
                \chi_3(a^{-1})K^{\mathsf T}\rho(a)K
 =\frac{6638585+2275584\sqrt5}{20}.
\]
Esta es una evaluación exacta de sesenta términos especificados,
reproducida por el certificado adjunto con aritmética racional
cuadrática. Los dos coeficientes positivos prueban la no nulidad.
\end{proof}

\begin{definition}[Dirección dodecafásica y complemento transversal]
La dirección seleccionada por \(K\) en la carta declarada es
\(\ell_K=\RR u_K\). Su complemento transversal dentro de
\(V_{11}=\mathbf1^\perp\) es
\[
 V_{10}(K)=V_{11}\cap u_K^\perp.
\]
\end{definition}

\begin{proposition}[Reducción ortogonal determinada por \(K\)]
\label{prop:k-reduccion-diez}
El operador
\begin{equation}
 P_{10}(K)=P_{11}
       -\frac{u_Ku_K^{\mathsf T}}{\lVert u_K\rVert^2}
 \label{eq:k-proyector-diez}
\end{equation}
es el proyector ortogonal de rango diez sobre \(V_{10}(K)\). Se cumple
\[
 V_{11}=\ell_K\mathbin{\oplus^{\perp}}V_{10}(K),
 \qquad P_{10}P_{11}=P_{10},\quad P_{10}u_K=0.
\]
\end{proposition}
\begin{proof}
Sea \(E_K=u_Ku_K^{\mathsf T}/\lVert u_K\rVert^2\). La
no nulidad demostrada garantiza que está definido. Es simétrico,
\(E_K^2=E_K\), tiene rango uno y, como \(u_K\in V_{11}\),
\(P_{11}E_K=E_KP_{11}=E_K\). De aquí
\[
 (P_{11}-E_K)^2=P_{11}-2E_K+E_K=P_{11}-E_K.
\]
Su imagen es \(V_{11}\cap u_K^\perp\) y su traza es
\(11-1=10\). La descomposición y las identidades restantes se
obtienen por proyección sobre los dos sumandos ortogonales.
\end{proof}

Este resultado precisa una función de \(K\) que su denominación como
registro no agotaba: además de representar reversiblemente datos
orientados y de seleccionar la tétrada \(B_K\), determina una
polarización lineal. La dirección depende del registro completo a
través de \(P_3K\), no únicamente de los cuatro puntos de su soporte.
Se ha mantenido así el orden registro, representación, dirección y
reducción dimensional (proposición~\ref{prop:k-reduccion-diez}).

La transformación tiene una ley de covariancia precisa. Si se cambia
simultáneamente la carta por una matriz ortogonal de permutación \(S\),
\(K'=SK\) y \(P'_3=SP_3S^{-1}\), entonces
\[
 u_{K'}'=Su_K,\qquad P'_{10}(K')=SP_{10}(K)S^{-1}.
\]
Por otra parte, sustituir \(u_K\) por \(a u_K\), con \(a\ne0\),
no cambia \(E_K\). La dirección retiene una polarización y no
permite recuperar por sí sola ni la amplitud de \(u_K\) ni el
registro íntegro. La reversibilidad de \(\kappa\leftrightarrow K\)
permanece en su dominio propio.

La realización geométrica aquí establecida es una reducción de espacios
con producto interno. Su empleo posterior como dirección compacta
requiere transportar esta recta a la realización física correspondiente.
El apartado \ref{exc:pantallas-dualidad} expone el mapa algebraico
que mantiene coordinadas esta reducción y la pantalla reticular.


% SOURCE foundation/sections/estrella_pentadica_recuperada.tex
\subsubsection{Representación pentádica de la microfibra y de su incidencia}
\label{pi-estrella:desarrollo}

La disposición de las cinco regiones permite visualizar simultáneamente
la unidad del cociente ternario y la diferenciación que conserva el
catálogo. La aplicación relevante es la restricción
\[
 \Psi\big|_{\mathscr R_\pi^{\rm micro}}:
 \mathscr R_\pi^{\rm micro}\longrightarrow\{010211\},
 \qquad \Psi(U)=(-U)\bmod3.
\]
Sus cinco antecedentes son los vectores \(U_6\) de la tabla anterior.
La figura \ref{pi-estrella:figura} recupera su organización radial,
conservando el orden de carta, las firmas y las multiplicidades del
catálogo (ecuación~\eqref{eq:palabras-regionales}).

\input{figures/estrella_cinco_regiones.tex}

La proyección al centro identifica las cinco imágenes ternarias; el
registro de las posiciones exteriores retiene la información regional
que hace posible la realización incidencial. La distribución
\(1_\perp+3_{++}+1_{--}\) especifica una región transversal, tres
regiones positivas y una negativa. Sus pesos normalizados son
\(3/7\) y \(1/7\) para cada una de las cuatro restantes: la
multiplicidad transversal es triple, aun cuando las cinco imágenes bajo
\(\Psi\) coincidan.

\Needspace{11\baselineskip}
El enlace con la estrella de Witt se formula mediante la bandera
\(B_K\subset H^-_\alpha\), el alineamiento \(\ell\) y la
elevación \(\iota_D\) ya construidos. Las cinco posiciones regionales
se realizan en las cinco octadas sobre \(\ell(B_K)\):
\[
\begin{aligned}
 U_\pi^\perp&\longmapsto O^\perp_{\ell(B_K)},\\
 U_\pi^{--}&\longmapsto\iota_D(H^-_\alpha),\\
 \{U_{\pi,1}^{++},U_{\pi,2}^{++},U_{\pi,3}^{++}\}
 &\longmapsto
 \{\iota_D(B_K\cup P):P\in\mathcal P_+\},\\
 \mathcal P_+&=\bigl\{\{1,3\},\{2,11\},\{8,12\}\bigr\}.
\end{aligned}
\]
La asignación de las tres posiciones positivas a los tres pares se
efectúa en una carta ordenada. La fila transversal y la fila negativa
quedan distinguidas por sus predicados regionales y por la bandera.

Así se articulan las dos figuras: la figura
\ref{exc:fig-estrella} muestra las cuatro hexadas del pequeño diseño
de Witt sobre \(B_K\); la figura \ref{pi-estrella:figura} muestra
las cinco regiones que, bajo el alineamiento, ocupan las cuatro octadas
residuales y la octada transversal del diseño binario. La ley
\(4+1\), demostrada en la proposición
\ref{exc:residuo-binario}, especifica el cambio de incidencia. En esta
composición, \(K\) determina la cuaterna marcada y la firma de
\(\alpha\) distingue una de sus hexadas. Las regiones conservan,
además de su valor ternario común, la estructura necesaria para ese
transporte hacia la realización excepcional.


% SOURCE foundation/figures/estrella_cinco_regiones.tex
% Orden de carta y datos recuperados de fig_cinco_regiones_pi.tex.
\begin{figure}[H]
\centering
\begin{tikzpicture}[x=1cm,y=1cm,font=\small,>=Latex]
 \coordinate (a) at (90:2.2);
 \coordinate (b) at (18:2.2);
 \coordinate (c) at (-54:2.2);
 \coordinate (d) at (-126:2.2);
 \coordinate (e) at (162:2.2);
 \draw[black!35,line width=.45pt]
    (a)--(c)--(e)--(b)--(d)--cycle;
 \node[draw,fill=white,circle,minimum size=1.6cm,align=center,
    inner sep=3pt] (z) at (0,0) {$w_\pi$\\$010211$};
 \foreach \v in {a,b,c,d,e}{
    \fill (\v) circle (2pt);
    \draw[->,line width=.6pt] (\v)--(z);
 }
 \node[above=3mm,align=center] at (a)
   {$U_\pi^\perp$\\$555555$\\multiplicidad $432$};
 \node[right=3mm,align=left] at (b)
   {$U_{\pi,1}^{++}$\\$339555$\\multiplicidad $144$};
 \node[below right=2mm and 1mm,align=left] at (c)
   {$U_{\pi,2}^{++}$\\$393555$\\multiplicidad $144$};
 \node[below left=2mm and 1mm,align=right] at (d)
   {$U_{\pi,3}^{++}$\\$933555$\\multiplicidad $144$};
 \node[left=3mm,align=right] at (e)
   {$U_\pi^{--}$\\$933717$\\multiplicidad $144$};
 \node[align=center,font=\footnotesize] at (0,-3.55)
   {$1_\perp+3_{++}+1_{--}$\qquad $432+4\cdot144=1008$};
\end{tikzpicture}
\caption{Las cinco regiones de \(\pi\) en disposición pentádica.
Cada posición conserva su firma multiplicativa y su multiplicidad;
las flechas radiales representan la proyección común
\(\Psi(U)=010211\). El contorno estrellado organiza las cinco
posiciones de la carta. Su realización en las cinco octadas de Witt
se efectúa mediante el alineamiento regional y la descomposición
\(4+1\) del texto.}
\label{pi-estrella:figura}
\end{figure}


% SOURCE foundation/sections/incidencia_reticulo.tex
\paragraph{Del origen incidencial a la modificación reticular.}
\label{inc-ampl-reticulo}
La bandera marcada determina una posición mediante una operación
conjuntista equivariante:
\[
 o_*=(H_e\cap H_\varphi)
       \setminus(H^-_\alpha\cup H_{\rm APP})=\{1\}.
\]
Su función consiste en distinguir una componente entre las doce que
intervienen en la realización reticular. El código ternario completo
actúa como código de pegado en el módulo discriminante
$(A_2^*/A_2)^{12}$. Su preimagen define $N(C_W)$, y la
autodualidad, los pesos y la distancia del código se traducen,
respectivamente, en propiedades de integralidad y determinante, paridad
y cotas de norma de las clases de pegado. Aquí la información simbólica
del código cumple una función que el soporte aislado ya no desempeña.

La métrica $A_2$ y la cámara radial positiva se fijan en el desarrollo
reticular. Dentro de esa cámara, la congruencia necesaria para un vecino
par de índice tres tiene doce realizaciones de norma mínima $54$,
permutadas por el cambio de componente distinguida. El origen $o_*$
selecciona
\[
 v_\alpha=4\rho_1+\rho_2+\cdots+\rho_{12},\qquad
 v_\alpha^2=2(4^2+11)=54.
\]
El coeficiente $4$ pertenece a esa solución de la congruencia en la
cámara declarada. El subíndice $\alpha$ registra la procedencia
incidencial de la selección: su contenido es el marco que intervino en
la clausura y que ahora individualiza una dirección reticular.

El emparejamiento con $v_\alpha$ determina el subretículo
$N_{\alpha,0}$ de índice tres, y la adición de la clase
$v_\alpha/3$ produce el vecino $N_\alpha$ de
\eqref{exc:vecino}. Esta modificación conserva el rango y recupera
unimodularidad y paridad; la selección por congruencia elimina las raíces
anteriores, y las cotas locales excluyen raíces en las nuevas clases.
La identificación con Leech aparece después de esas verificaciones. La
realización binaria de las cinco regiones y esta construcción reticular
son dos prolongaciones del marco común: la segunda recibe directamente
el código ternario como dato de pegado.



% SOURCE foundation/sections/incidencia_automorfismos.tex
\paragraph{Graduación, unidad y automorfismos.}
\label{inc-ampl-automorfismos}
La construcción de Frenkel--Lepowsky--Meurman se aplica al retículo
$N_\alpha$ una vez verificadas sus propiedades. Incorpora el espacio
de osciladores, la extensión central del retículo, el producto de vértice
y el sector torcido por el levantamiento de la negación. El resultado es
el álgebra $V^\natural_{\rm HMT}$ de \eqref{exc:flm}, cuyo grupo
de automorfismos es isomorfo al Monstruo. Su dominio de acción queda así
fijado por una estructura algebraica con graduación y producto. El teorema
de Moonshine de Borcherds determina posteriormente las trazas graduadas
con inserción de esos automorfismos \cite{flm,borcherds}.

La unidad adquiere en este nivel una realización precisa: la componente
de peso cero es la línea del vacío y aporta el coeficiente $1$ de
$q^{-1}$ en el carácter graduado. En el nivel cilíndrico, las
aplicaciones de refinamiento conservan la función constante mediante
precomposición. Cada afirmación de conservación tiene, por tanto, su
objeto y su aplicación. El desarrollo de las pruebas hace visible esa
continuidad de construcción, desde las operaciones aritméticas y la
memoria hasta la incidencia, la métrica y los productos de vértice.

Esta articulación sitúa a $K$ y a la clausura de $\alpha$ en el argumento
principal. $K$ conserva reversiblemente una coordenada integral de la
historia; la normalización de $\alpha$ compone los registros regionales
con esa coordenada; la configuración firmada determina una bandera; y
la bandera marcada selecciona la modificación reticular sobre la que se
realiza Moonshine. Las demostraciones siguientes desarrollan estas
aplicaciones, distinguiendo las equivalencias reversibles, los cocientes
de información y las realizaciones que añaden métrica o estructura
algebraica. Así se expresa una relación demostrable entre estructura,
ecuación y valor, conservando la atribución de cada construcción y el
alcance de sus teoremas.


% SOURCE foundation/sections/fibras_app_witt.tex
\subsection{Construcción estatal de las doce fibras ciclotómicas}
\label{nuclear:fibras-app-witt}
Sobre las marcas de APP, la multiplicación \(a:r\mapsto4r\) en
\(\ZZ/9\ZZ\) satisface \(a^3=1\). Las hojas verifican
\[
 S(ax,ay)=aS(x,y),\qquad P(ax,ay)=a^{-1}P(x,y),
\]
pues \(16\equiv7\equiv4^{-1}\pmod9\). Las órbitas unitarias
ordenadas son \(\mathcal O_0=(1,4,7)\), \(\mathcal O_1=(2,8,5)\).
En cada una, el módulo de aumento
\(A(\mathcal O)=\ker(\ZZ[\mathcal O]\xrightarrow{\sum}\ZZ)\)
hereda el producto que hace ortonormales los \(e_r\). En la base
\(b_1=e_r-e_{4r}\), \(b_2=e_{4r}-e_{7r}\), sus matrices son
\[
 G_{A_2}=\begin{pmatrix}2&-1\\-1&2\end{pmatrix},\qquad
 C=\begin{pmatrix}0&-1\\1&-1\end{pmatrix},\qquad
 C^{\mathsf T}G_{A_2}C=G_{A_2}.
\]
La forma reticular procede, por tanto, de las órbitas de las marcas.
Para los pares \((1,2),(2,3),(3,1)\) y las hojas
\(\mu\in\{\Sigma,\Pi\}\), se obtiene
\[
 W_{\APP}=\bigoplus_{(i,j)}\bigoplus_\mu\bigoplus_{\varepsilon=0}^1
 A(\mathcal O_\varepsilon),\qquad
 \ell^\Sigma_{ij}(x)=x_i+x_j,\quad
 \ell^\Pi_{ij}(x)=x_ix_j\pmod9.
\]
Cada sumando conserva par, hoja y órbita. La aplicación estatal es
\[
 \Psi:(\ZZ/9\ZZ)^3\longrightarrow W_{\APP},\qquad
 \Psi(x)_{ij,\mu,\varepsilon}=\beta_\varepsilon(\ell^\mu_{ij}(x)),
 \quad
 \beta_\varepsilon(r)=
 \begin{cases}e_r-e_{4r},&r\in\mathcal O_\varepsilon,\\0,&r\notin\mathcal O_\varepsilon.
 \end{cases}
\]
\begin{proposition}[Equivarianza y rango del módulo estatal]
\label{nuclear:psi-rango}
La acción de \(C\) en las seis fibras aditivas y de \(C^{-1}\) en
las seis multiplicativas cumple \(\Psi(ax)=\rho(a)\Psi(x)\).
Las imágenes generan \(W_{\APP}\otimes\QQ\), de dimensión veinticuatro.
\end{proposition}
\begin{proof}
Las identidades de las hojas y \(\beta(4r)=C\beta(r)\) prueban la
equivarianza. Para el rango, se ordenan las coordenadas por los tres
pares indicados, por \(\Sigma,\Pi\), por \(\mathcal O_0,\mathcal O_1\)
y por \(b_1,b_2\). En cada órbita, \(\beta\) toma sucesivamente
\((1,0),(0,1),(-1,-1)\). La matriz de las imágenes de estos estados,
leídos por filas, tiene determinante \(9\):
\[
\begin{array}{rrrr}
(0,0,1)&(0,0,2)&(0,0,4)&(0,0,5)\\
(0,1,0)&(0,1,1)&(0,1,2)&(0,1,3)\\
(0,1,4)&(0,1,5)&(0,1,6)&(0,2,0)\\
(0,2,2)&(0,2,3)&(0,4,0)&(0,5,0)\\
(1,0,1)&(1,0,2)&(1,0,4)&(1,0,5)\\
(1,1,0)&(1,2,0)&(1,4,0)&(1,5,0).
\end{array}
\]
La eliminación sucesiva, reduciendo cada fila por las anteriores y
normalizando su primer coeficiente no anulado, da columnas pivote
numeradas desde cero
\[
\begin{split}
&(8,10,9,11,0,12,14,16,13,15,17,2,\\
&\hspace{2em}18,19,1,3,20,22,21,23,4,6,5,7)
\end{split}
\]
y pivotes
\[
(1,1,1,-1,1,1,1,1,1,-1,3,1,1,3,1,-1,1,1,1,-1,1,1,1,-1).
\]
Su producto, multiplicado por el signo de la permutación de columnas,
es \(9\). El menor invertible acredita el rango total.
\end{proof}

\subsection{Alineamiento dodecafásico con las posiciones de Witt}
\label{nuclear:alineamiento-witt}
El índice del par es \(i\in\ZZ/3\ZZ\), el de la hoja
\(\mu\in\{0,1\}\) y el de la órbita \(\varepsilon\in\{0,1\}\).
Fijado el origen dodecafásico, \(\lambda(i,\mu,\varepsilon)
=4i+2\mu+\varepsilon\pmod{12}\) es biyectiva: los cuatro valores
de cada \(i\) forman un bloque y los tres bloques son disjuntos.
Sea \(R=\left(\begin{smallmatrix}0&1\\1&0\end{smallmatrix}\right)\).
La multiplicación prueba \(RCR=C^{-1}\) y
\(R^{\mathsf T}G_{A_2}R=G_{A_2}\). Para las trivializaciones
\(\iota_b\) y los marcos \(P_b\) transportados conjuntamente con
la incidencia de Witt, ponemos
\[
g_b=\iota_b^{-1}P_bCP_b^{-1}\iota_b,\qquad
T_{i\mu\varepsilon}=\iota_{\lambda(i,\mu,\varepsilon)}^{-1}
P_{\lambda(i,\mu,\varepsilon)}R^\mu.
\]
Entonces
\[
g_{\lambda(i,\mu,\varepsilon)}T_{i\mu\varepsilon}
=T_{i\mu\varepsilon}C^{(-1)^\mu}.
\]
La suma directa es un isomorfismo \(C_3\)-equivariante entre las doce
fibras estatales y las doce posiciones reticulares. Conserva par,
hoja, órbita y forma bajo el transporte simultáneo del origen,
orden y orientación. Estos datos permanecen en el pegado reticular
y el paso al vecino que se construyen a continuación.


% SOURCE foundation/sections/moonshine_comparacion.tex
% Incorporación aditiva para el Artículo I; no modifica la fuente anterior.
% RESULTADO_RECUPERADO: integral, capítulo 22, propietario extenso
% colaboracion/partes_i_ii/source/base_residencias/
% capitulo_22_c20_lineas_2613_3270.tex, líneas 176--576.
% Equivalencia 2/3: public/residencias/capitulo_22_voa_leech_moonshine.tex,
% líneas 55--127. Orden dos y Fricke: hmt/15d_voa_moonshine_orden_dos_rev11.tex.
% Prerrequisitos residentes en I: código C_W, matriz A_W, N(A_2^12),
% origen o_*, vector v_alpha de norma54, N_alpha reconocido como Leech,
% VOA reticular y presentación FLM de las secciones precedentes.
% Bibliografía nueva: chenlamshimakura2016, abelamyamada2017,
% kmconwaynorton1979 (bibitems listos para incorporar al final de este archivo).

\subsection{Orientación ternaria del vecino de Leech}
\label{km:orden-tres-reticular}

La presentación de orden dos construida anteriormente utiliza la negación
del retículo de Leech. El mismo retículo, obtenido de la incidencia marcada
por \(K\) y por el soporte dodecafásico de \(\alpha\), admite una segunda
presentación compatible con el carácter ternario de la construcción.
Para hacer explícita esta continuidad se conserva el pegado, se construye
su isometría de orden tres y se comprueba que ésta preserva el vecino ya
seleccionado. El paso no requiere alterar \(K\), su lectura racional ni
la definición previa de \(\alpha\). Se desarrolla aquí la prolongación de orden tres, incluyendo las operaciones
reticulares y las condiciones precisas de los teoremas clásicos utilizados.

Trabajamos en la base de raíces simples de cada factor \(A_2\), con
\begin{equation}
 B_2=\begin{pmatrix}2&-1\\-1&2\end{pmatrix},\qquad
 \mathsf C=\begin{pmatrix}0&-1\\1&-1\end{pmatrix},\qquad
 \mathsf R=\begin{pmatrix}0&1\\1&0\end{pmatrix},\qquad
 \rho=\begin{pmatrix}1\\1\end{pmatrix}.
 \label{km:coxeter-reflexion}
\end{equation}
La multiplicación da
\[
 \mathsf C^3=I_2,\quad
 \mathsf C^2+\mathsf C+I_2=0,\quad
 \mathsf C^{\mathsf T}B_2\mathsf C=B_2,\quad
 \mathsf R\mathsf C\mathsf R=\mathsf C^{-1},\quad
 \mathsf R^{\mathsf T}B_2\mathsf R=B_2.
\]
Estas identidades fijan tanto la métrica como la orientación de la acción.
En el discriminante \(A_2^*/A_2\cong\mathbb F_3\) tomamos como
generador la clase de \(\varpi=(2,1)^{\mathsf T}/3\). Se cumple
\(\mathsf C\varpi-\varpi=(-1,0)^{\mathsf T}\), por lo que
\(\mathsf C\) actúa trivialmente sobre el discriminante; la reflexión
\(\mathsf R\) cambia su signo.

La matriz de Paley--Witt ya construida produce un elemento de soporte
completo mediante la siguiente operación:
\begin{equation}
 \begin{gathered}
 q_*=(1,1,1,1,1,1),\qquad q_*A_W=(2,1,1,1,1,1),\\
 c_*=(q_*,q_*A_W)
 =(1,1,1,1,1,1,2,1,1,1,1,1)\in C_W.
 \end{gathered}
 \label{km:palabra-total}
\end{equation}
Cada coordenada de \(c_*\) es no nula. Para las doce posiciones
definimos los marcos y la isometría
\begin{equation}
 P_b(c_*)=
 \begin{cases}\mathsf R,&(c_*)_b=1,\\I_2,&(c_*)_b=2,\end{cases}
 \qquad
 g_c=\bigoplus_{b=1}^{12}P_b(c_*)\mathsf C P_b(c_*)^{-1}.
 \label{km:isometria-ternaria}
\end{equation}
El símbolo \(P_b(c_*)\) designa aquí un marco invertible de un factor
\(A_2\), no el proyector regional \(P_3\).

\begin{proposition}[Isometría de orden tres sobre el vecino marcado]
\label{km:estabilidad-vecino}
El operador \(g_c\) es una isometría de orden tres sin vectores fijos
no nulos. Preserva el retículo de Niemeier \(N=N(C_W)\) y el vecino
\(N_\alpha\) construido en \eqref{exc:vecino}.
\end{proposition}
\begin{proof}
Cada bloque de \(g_c\) es \(\mathsf C\) o \(\mathsf C^{-1}\);
su polinomio mínimo es \(X^2+X+1\). Esto prueba el orden y la
ausencia de vectores fijos. Ambos bloques actúan trivialmente sobre el
discriminante, de modo que conservan cada clase de pegado y, por tanto,
el retículo \(N\).

Escribamos \(v=v_\alpha=4\rho_{o_*}+\sum_{b\ne o_*}\rho_b\),
con los marcos y el origen incidencial ya fijados. Sus coeficientes
radiales son \(a_{o_*}=4\) y \(a_b=1\) para \(b\ne o_*\).
Todos satisfacen \(a_b\equiv1\pmod3\). En un factor,
\[
 (\mathsf C-I_2)\rho=(-2,-1)^{\mathsf T},\qquad
 (\mathsf C^{-1}-I_2)\rho=(-1,-2)^{\mathsf T}.
\]
Después de dividir por tres, sus clases discriminantes son,
respectivamente, \(2\) y \(1\). En consecuencia, los vectores
\begin{equation}
 y=\frac{g_cv-v}{3},\qquad z=\frac{g_c^{-1}v-v}{3}
 \label{km:desplazamientos}
\end{equation}
tienen palabras discriminantes \(c_*\) y \(-c_*\), ambas en
\(C_W\). Por ello \(y,z\in N\). El producto local es
\[
 \left\langle
 \frac{(\mathsf C^{\pm1}-I_2)a_b\rho}{3},a_b\rho
 \right\rangle=-a_b^2,
\]
y su suma da
\(\langle y,v\rangle=\langle z,v\rangle=-(16+11)=-27\).
Así, \(y,z\in N_0=\{x\in N:\langle x,v\rangle\equiv0\pmod3\}\).
Para \(x\in N_0\),
\[
 \langle g_cx,v\rangle
 =\langle x,g_c^{-1}v\rangle
 =\langle x,v\rangle+3\langle x,z\rangle\equiv0\pmod3.
\]
El argumento aplicado a \(g_c^{-1}\) prueba \(g_cN_0=N_0\).
Finalmente, \(g_c(v/3)=v/3+y\) conserva la clase que genera
\(N_\alpha=N_0+\mathbb Z(v/3)\). Se obtiene
\(g_cN_\alpha=N_\alpha\).
\end{proof}

La prueba conserva más que la dimensión del retículo: utiliza el código
de pegado, su orientación, el origen de la bandera y la congruencia del
vector radial. Estos datos explican por qué la segunda presentación
permanece vinculada al registro dodecafásico del que partió la construcción.

\subsection{Traza ternaria, peso torcido y extensión de orden tres}
\label{km:orbifold-tres}

Sea \(\Lambda=N_\alpha\) y sea \(\widehat g_c\) el levantamiento
estándar de \(g_c\) al álgebra reticular
\(V_\Lambda=M(1)\otimes\mathbb C_\varepsilon[\Lambda]\), con la
extensión central y el cociclo ya declarados. El orden impar permite
elegir este levantamiento de orden tres. Cada factor complejo \(A_2\)
aporta los autovalores \(\zeta_3\) y \(\zeta_3^2\), donde
\(\zeta_3^2+\zeta_3+1=0\).

\begin{proposition}[Traza graduada y tipo del levantamiento]
\label{km:traza-tipo}
Para \(j=1,2\),
\begin{equation}
 \operatorname{Tr}_{V_\Lambda}
       \bigl(\widehat g_c^{\,j}q^{L_0-1}\bigr)
 =t_3(\tau)
 :=q^{-1}\prod_{n\ge1}(1+q^n+q^{2n})^{-12}.
 \label{km:tres-traza}
\end{equation}
El peso mínimo de los módulos torcidos no triviales es \(4/3\), y
el levantamiento es de tipo cero.
\end{proposition}
\begin{proof}
En la parte reticular de la traza sólo contribuye el vector cero,
porque \(g_c\) y \(g_c^2\) carecen de vectores fijos no nulos.
En cada grado oscilador, la traza de potencias simétricas es el inverso
del determinante \(I-g_c^jq^n\). Cada factor \(A_2\) contribuye
\(1+q^n+q^{2n}\), lo que demuestra \eqref{km:tres-traza}.

La fórmula del peso del vacío torcido para un automorfismo reticular de
orden tres, aplicada a los dos espacios propios de dimensión doce, da
\[
 \rho_{\widehat g_c}
 =\frac{1}{4\cdot3^2}
   \bigl(1\cdot2\cdot12+2\cdot1\cdot12\bigr)=\frac43.
\]
El tipo es la clase de \(3^2\rho_{\widehat g_c}=12\) módulo
tres, que es cero. La fórmula de peso es un resultado de la teoría de
módulos torcidos; las multiplicidades y la isometría que la alimentan
se han construido aquí.
\end{proof}

\begin{theorem}[Segunda presentación del álgebra Moonshine]
\label{km:moonshine-tres}
La extensión cíclica de tipo cero
\begin{equation}
 V_{(3)}=\operatorname{Orb}_{\mathbb Z/3}
                 (V_\Lambda,\widehat g_c)
 \label{km:extension-orden3}
\end{equation}
es isomorfa al álgebra de operadores de vértice \(V^\natural\).
\end{theorem}
\begin{proof}
El teorema \ref{exc:teorema-leech} proporciona un retículo positivo,
par, unimodular, de rango veinticuatro y sin raíces. La
Proposición~\ref{km:estabilidad-vecino} construye sobre ese mismo
retículo una isometría de orden tres sin vectores fijos; la
Proposición~\ref{km:traza-tipo} verifica el levantamiento y su tipo.
Se aplica entonces el teorema de Chen--Lam--Shimakura sobre el
orbifold de orden tres del álgebra reticular de Leech
\cite[Teorema~1.1]{chenlamshimakura2016}. La identificación también queda comprendida
en el tratamiento uniforme de Abe--Lam--Yamada
\cite[Teorema~4.4 y secciones~2--4]{abelamyamada2017}. Esta aplicación es posterior a la construcción
de los datos reticulares; la igualdad de caracteres no sustituye las
hipótesis del teorema utilizado.
\end{proof}

Para precisar la simetría dual, escribamos
\(V_{(3)}=W_0\oplus W_1\oplus W_2\), donde \(W_i\) es el sector
integral seleccionado en el módulo \(\widehat g_c^{\,i}\)-torcido.
El producto satisface
\(Y(W_i,z)W_j\subset W_{i+j\bmod3}((z))\). Por ello
\begin{equation}
 g^\vee|_{W_i}=\zeta_3^i I
 \label{km:dual-ciclico}
\end{equation}
define un automorfismo de orden tres. En efecto, el factor sobre un
producto es \(\zeta_3^{i+j}=\zeta_3^i\zeta_3^j\), y el vacío y el
vector conforme pertenecen a \(W_0\). La identificación clásica de
este automorfismo en el Monstruo es la clase \(3B\), con serie
\(T_{3B}=t_3+12\) \cite{kmconwaynorton1979,borcherds}.
Esta normalización puede comprobarse en la construcción. El carácter
reticular es \(\operatorname{ch}V_\Lambda=J+24\): su polo inicial
es \(q^{-1}\), su espacio de peso uno tiene dimensión 24 y el
carácter modular de peso cero queda determinado por esos datos.
Si \(F\) es el carácter del subespacio fijo y \(H\) el de cada
sector torcido integral, la proyección sobre invariantes y la conjugación
de los dos sectores dan
\[
 F=\frac{J+24+2t_3}{3},\qquad F+2H=J,
 \qquad \operatorname{Tr}g^\vee q^{L_0-1}=F-H=t_3+12.
\]
Su dominio es el álgebra extendida, mientras
\(\widehat g_c\) actúa sobre el álgebra reticular anterior.

\subsection{Equivalencia de las presentaciones de órdenes dos y tres}
\label{km:equivalencia-dos-tres}

Denotemos por
\[
 V_{(2)}=(V_\Lambda)^+\oplus(V_\Lambda^T)^+
\]
la presentación FLM de \eqref{exc:flm}. Ambas presentaciones conservan
el mismo vecino de Leech como antecedente. Sus extensiones emplean
automorfismos diferentes y productos de vértice determinados mediante
los respectivos datos de levantamiento.

\begin{theorem}[Equivalencia graduada y transporte de automorfismos]
\label{km:equivalencia-voa}
Fijados isomorfismos de álgebras de operadores de vértice
\[
 \phi_2:V_{(2)}\xrightarrow{\ \cong\ }V^\natural,
 \qquad
 \phi_3:V_{(3)}\xrightarrow{\ \cong\ }V^\natural,
\]
la composición
\begin{equation}
 \Phi_{23}=\phi_2^{-1}\phi_3:
 V_{(3)}\xrightarrow{\ \cong\ }V_{(2)}
 \label{km:Phi23}
\end{equation}
conserva el vacío, el vector conforme, la gradación y el producto de
vértice. La conjugación por \(\Phi_{23}\) identifica los grupos de
automorfismos y conserva sus trazas graduadas.
\end{theorem}
\begin{proof}
La composición de los dos isomorfismos preserva las estructuras
declaradas. En particular,
\[
 \Phi_{23}Y_{(3)}(a,z)b
 =Y_{(2)}(\Phi_{23}a,z)\Phi_{23}b,
 \qquad
 \Phi_{23}L_0^{(3)}=L_0^{(2)}\Phi_{23}.
\]
Si \(h\) es un automorfismo de \(V_{(3)}\), entonces
\(h' =\Phi_{23}h\Phi_{23}^{-1}\) es un automorfismo de
\(V_{(2)}\). La igualdad de trazas en cada componente homogénea
finito-dimensional da
\[
 \operatorname{Tr}_{V_{(3)}}h q^{L_0-1}
 =\operatorname{Tr}_{V_{(2)}}h' q^{L_0-1}.
\]
\end{proof}

La elección de \(\phi_2\) y \(\phi_3\) forma parte del enunciado;
\(\Phi_{23}\) no se presenta como una identificación canónica
independiente de esas elecciones. Tampoco transforma un automorfismo
de orden tres en la involución de orden dos: la conjugación conserva
el orden. Su contenido es que dos extensiones distintas realizan la
misma estructura graduada, con transporte explícito de sus productos,
representaciones y trazas.

\subsection{Funciones eta, involución de Fricke y normalización del vacío}
\label{km:fricke}

Para \(\operatorname{Im}\tau>0\), las trazas de los dos automorfismos
reticulares admiten las expresiones
\begin{equation}
 \eta(\tau)=q^{1/24}\prod_{n\ge1}(1-q^n),\qquad
 t_2=\left(\frac{\eta(\tau)}{\eta(2\tau)}\right)^{24},\qquad
 t_3=\left(\frac{\eta(\tau)}{\eta(3\tau)}\right)^{12}.
 \label{km:eta-trazas}
\end{equation}
Se deducen directamente de las factorizaciones
\(1-q^{2n}=(1-q^n)(1+q^n)\) y
\(1-q^{3n}=(1-q^n)(1+q^n+q^{2n})\). El orden en \(q\) del
cociente de nivel tres, antes de elevarlo a la potencia doce, es
\[
 \frac1{24}-\frac3{24}=-\frac1{12};
 \qquad 12\left(-\frac1{12}\right)=-1.
\]
Aquí \(-1/12\) designa la diferencia exacta de los exponentes
iniciales de dos funciones eta. La operación está especificada y no
requiere interpretar una suma divergente como una suma convergente.
La coincidencia con otros defectos normalizados del corpus conserva
las realizaciones respectivas; no identifica sus operadores por
igualdad del escalar.

\begin{proposition}[Involución de Fricke sobre la traza de orden dos]
\label{km:fricke-prop}
Sea \(w_2\tau=-1/(2\tau)\). Entonces
\begin{equation}
 t_2(w_2\tau)=\frac{2^{12}}{t_2(\tau)},\qquad
 T_{2A}=t_2+24+\frac{2^{12}}{t_2}
 \label{km:fricke-t2}
\end{equation}
y \(T_{2A}(w_2\tau)=T_{2A}(\tau)\).
\end{proposition}
\begin{proof}
La ley de transformación
\(\eta(-1/\tau)=(-i\tau)^{1/2}\eta(\tau)\) da
\[
 \frac{\eta(-1/(2\tau))}{\eta(-1/\tau)}
 =\sqrt2\,\frac{\eta(2\tau)}{\eta(\tau)}.
\]
La potencia veinticuatro produce la primera identidad. La sustitución
intercambia los dos términos variables de \(T_{2A}\) y deja fijo
el término constante. El reconocimiento de esta expresión como serie
de McKay--Thompson de la clase \(2A\) corresponde a la
simetrización de nivel dos de Conway--Norton
\cite[\S\S4 y 6, tabla~3]{kmconwaynorton1979} y al teorema
de Moonshine \cite{borcherds}.
\end{proof}

Las uniformizaciones modulares de niveles dos y tres, en las
normalizaciones anteriores, son
\begin{align}
 j&=\frac{(t_2+256)^3}{t_2^2},
 \label{km:j-nivel2}\\
 J=j-744
 &=t_3+12+\frac{10\cdot3^9}{t_3}
       +\frac{4\cdot3^{14}}{t_3^2}
       +\frac{3^{18}}{t_3^3}.
 \label{km:j-nivel3}
\end{align}
Se utilizan como identidades modulares clásicas sobre las trazas ya
construidas, no como reglas para seleccionar el código o el vecino.
La segunda también se escribe
\(j=(t_3+27)(t_3+243)^3/t_3^3\).
El desarrollo directo del producto de \eqref{km:tres-traza} comienza por
\[
 t_3=q^{-1}-12+54q-76q^2-243q^3+1188q^4+O(q^5).
\]
En particular, \([q]J=54+10\cdot3^9=196884\); ésta es una
consecuencia de las trazas y la uniformización. La construcción de
\(V^\natural\) conserva su prueba reticular y de orbifold,
independientemente de esta comprobación del primer coeficiente.

Las tres operaciones que comparecen en esta prolongación tienen
dominios distintos. La simetría dual \(g^\vee\) actúa sobre sectores
del álgebra de vértices; Fricke actúa sobre el parámetro modular y
su función; la dualidad de radio intercambia momento y enrollamiento
en el sector reticular compacto. Sus fórmulas de reciprocidad pueden
coordinarse mediante los mapas correspondientes, pero el nombre
«dualidad» no sustituye esos mapas. Esta distinción permite continuar
la construcción hacia las pantallas dimensionales manteniendo la
incidencia, la gradación y el origen de cada operación.

% BIBITEMS NUEVOS: copiar al entorno thebibliography del artículo.
% Verificación primaria: arXiv 1606.05961v2 (teorema 1.1);
% arXiv 1705.09022v4 (secciones 2--4);
% Conway--Norton, pp. 311 y 314--315, tabla 3.
% \bibitem{chenlamshimakura2016}
% H.-Y. Chen, C. H. Lam y H. Shimakura,
% ``$\mathbb Z_3$-orbifold construction of the Moonshine vertex operator
% algebra and some maximal $3$-local subgroups of the Monster'',
% prepublicación, arXiv:1606.05961 (2016), versión 2 (2017).
% \href{https://arxiv.org/abs/1606.05961}{arXiv:1606.05961}.
% \bibitem{abelamyamada2017}
% T. Abe, C. H. Lam y H. Yamada,
% ``A remark on $\mathbb Z_p$-orbifold constructions of the Moonshine
% vertex operator algebra'', prepublicación, arXiv:1705.09022 (2017).
% \href{https://arxiv.org/abs/1705.09022}{arXiv:1705.09022}.
% \bibitem{kmconwaynorton1979}
% J. H. Conway y S. P. Norton, ``Monstrous Moonshine'',
% \emph{Bulletin of the London Mathematical Society} \textbf{11},
% 308--339 (1979).
% \href{https://doi.org/10.1112/blms/11.3.308}{doi:10.1112/blms/11.3.308}.


% SOURCE foundation/sections/pantallas_dualidad.tex
% REV02: pruebas reunidas en 02_dualidad_t.tex y 03_pantallas_coordinadas.tex.
\subsection{Incidencia excepcional y coordinación de las pantallas}
\label{exc:pantallas-dualidad}

La dirección seleccionada por \(K\) y el retículo de Leech proceden
de representaciones distintas de la incidencia construida. La primera
determina la reducción ortogonal \(12\to11\to10\); la segunda
interviene en la reducción isotrópica \(26\to24\). La coordinación
conserva el espacio previo, sus relaciones y la información transmitida
por cada proyección.

Las secciones de dualidad T y pantallas coordinadas de \cite{hmtX} desarrollan
conjuntamente esa coordinación. La primera reúne la involución de
coordenadas enteras, la conjugación unitaria con sus dominios y el
cambio de sección residual con memoria del cociente. La segunda
construye el morfismo \(\mathcal R_M^{\mathrm{alg}}\) de
allí definido y demuestra el isomorfismo de las presentaciones
cocientadas en el teorema correspondiente de pantallas isomorfas.
El grafo conserva la coordenada en \(V_{11}\) junto con su proyección
en \(V_{10}\), mientras el cociente isotrópico recupera el retículo
de Leech. Así se preservan los respectivos núcleos y codominios.

La elipse de acción se construye antes de aplicar la dualidad y fija
su radio adimensional junto con el radio recíproco
(sección de elipse de acción de \cite{hmtX}). El sector compacto enlaza entonces
la incidencia excepcional con una forma cuadrática determinada. Los
campos, la acción y los operadores de la realización física posterior
se presentan en su dominio propio, conservando este antecedente
algebraico y sus demostraciones internas.


% SOURCE foundation/sections/01_hilbert_nonadico.tex
% Procedencia: integral 2249, cap. 23, pp. 601--608; propietario base_c21_sin_encabezado.tex.
% Recuperación y desarrollo de pruebas. Requiere el núcleo común APP--TRIT--TPK.
\section{Realización hilbertiana de la fase nonádica y conservación de la unidad}
\label{iv:hilbert}

La construcción de Holografía Modular Triádica proporciona primero los
estados, sus transiciones y los mapas de restricción. La realización
hilbertiana añade a sus coordenadas discretas una estructura lineal y un
producto interno. Este orden determina qué representa cada operador:
las etiquetas de la base proceden de la construcción nonádica; las
operaciones lineales actúan después sobre esas etiquetas. La memoria
genealógica permanece en el estado del TPK y acompaña a las proyecciones
que se utilicen para representarlo.

\subsection{Fibra de fase, base etiquetada y producto interno}

Para una profundidad entera \(d\geq1\), sea
\begin{equation}
 \mathcal C_{9,d}=(\ZZ/9\ZZ)^d,
 \qquad \mathcal H_d=\ell^2(\mathcal C_{9,d}),
 \qquad \mathcal A_d=B(\mathcal H_d).
 \label{iv:eq:hilbert-finito}
\end{equation}
El producto interno es
\(\langle f,g\rangle=\sum_x\overline{f(x)}g(x)\). Las funciones
\(\delta_x\) forman una base ortonormal etiquetada por las coordenadas
de fase. Así, \(\dim\mathcal H_d=9^d\) y
\(\mathcal A_d\cong M_{9^d}(\CC)\). El álgebra matricial representa
esta fibra finita; las rutas, los cocientes y la memoria que una
proyección de fase omita siguen perteneciendo al estado de partida.

La factorización de coordenadas induce el unitario
\begin{equation}
 \mathcal H_{d+1}\longrightarrow\mathcal H_d\otimes\CC^9,
 \qquad \delta_{(x,r)}\longmapsto\delta_x\otimes\delta_r.
 \label{iv:eq:factorizacion-hilbert}
\end{equation}
Definimos la traza normalizada
\(\tau_d(a)=9^{-d}\operatorname{Tr}(a)\). La cuenta finita y la
normalización se distinguen: \(\operatorname{Tr}(I)=9^d\), mientras
que \(\tau_d(I)=1\). Asimismo, el grupo \((\ZZ/9\ZZ)^3\) y el
espacio vectorial \(\FF_3^6\) tienen ambos \(729\) elementos, pero
conservan leyes algebraicas diferentes. La realización de una fibra de
fase no sustituye la codificación ternaria de las emisiones APP.

\subsection{Traslación y fase: representación finita de Weyl}

En la realización compleja posterior a las construcciones de
\(\pi\) y de la estructura \(\mathrm i^2=-1\), fijamos
\(\omega=\exp(2\pi\mathrm i/9)\). Sobre \(\mathcal H_1\), sean
\begin{equation}
 (Xf)(r)=f(r-1),\qquad (Zf)(r)=\omega^r f(r).
 \label{iv:eq:weyl-operadores}
\end{equation}
La primera operación desplaza la coordenada de fase; la segunda la
evalúa mediante un carácter. Ambas son unitarias y satisfacen
\begin{equation}
 X^9=Z^9=I,\qquad ZX=\omega XZ.
 \label{iv:eq:weyl-relacion}
\end{equation}
En efecto, \((ZXf)(r)=\omega^r f(r-1)\), mientras que
\((XZf)(r)=\omega^{r-1}f(r-1)\). El carácter aparece aquí como
realización de la fase ya construida y no selecciona los registros
numéricos del TPK.

% RESULTADO_RECUPERADO: 14_numero_heisenberg_y_d108.tex:741--788.
% Prueba algebraica desarrollada con la convención X^a Z^b ya fijada aquí.
\subsubsection{Área orientada y extensión central de las traslaciones}
\label{iv:weyl-area-extension}

El soporte de traslaciones es \(X_9=(\ZZ/9\ZZ)^2\).
La curvatura mixta de APP construida en la
sección~\ref{iv:extension-curvatura-mixta} asigna \(+1\) a la
plaqueta orientada de la pareja suma--producto y \(-1\) a la
pareja de orden contrario. Su extensión bilineal a dos desplazamientos
\(u=(a,b)\) y \(v=(c,d)\) es la forma alternante
\begin{equation}
 \sigma_{\mathrm{ar}}(u,v)=ad-bc\quad\text{en }\ZZ/9\ZZ.
 \label{iv:eq:weyl-area-orientada}
\end{equation}
En efecto, la alternancia anula la evaluación sobre dos desplazamientos
del mismo eje. Los valores sobre la base orientada son
\(\sigma_{\mathrm{ar}}((1,0),(0,1))=1\) y
\(\sigma_{\mathrm{ar}}((0,1),(1,0))=-1\); distribuir ambos
argumentos da \eqref{iv:eq:weyl-area-orientada}. La misma fórmula es
la suma de plaquetas de cualquier cadena celular orientada que realice
ese paralelogramo: las aristas interiores se cancelan. Modificar un
representante entero por un múltiplo de nueve conserva su publicación
en \(\ZZ/9\ZZ\).

El orden \(X^aZ^b\) de los operadores ya definidos fija una sección
de la extensión central. Denotemos su factor de composición por
\begin{equation}
 c_{\mathrm W}(u,v)=bc\quad\text{en }\ZZ/9\ZZ.
 \label{iv:eq:weyl-cociclo-polarizado}
\end{equation}
Este cociclo conserva el orden de las dos operaciones, mientras
\(\sigma_{\mathrm{ar}}\) conserva la orientación del paralelogramo.

\begin{proposition}[Extensión central y signo del conmutador]
\label{iv:prop:heisenberg-orientado}
El conjunto \(\mathfrak H_9=(\ZZ/9\ZZ)^3\), con producto
\begin{equation}
 (a,b,t)(c,d,s)=(a+c,b+d,t+s+bc),
 \label{iv:eq:heisenberg-ley}
\end{equation}
es un grupo de orden \(729\). La aplicación
\[
 \varrho(a,b,t)=\omega^tX^aZ^b
\]
es una representación unitaria fiel sobre \(\mathcal H_1\).
Con el convenio \([g,h]=ghg^{-1}h^{-1}\), se tiene
\begin{equation}
 [\varrho(a,b,0),\varrho(c,d,0)]
 =\omega^{bc-ad}I
 =\omega^{-\sigma_{\mathrm{ar}}((a,b),(c,d))}I.
 \label{iv:eq:weyl-conmutador-orientado}
\end{equation}
El centro y el grupo de traslaciones forman la sucesión exacta
\[
 0\longrightarrow\ZZ/9\ZZ
 \xrightarrow{\ t\mapsto(0,0,t)\ }\mathfrak H_9
 \xrightarrow{\ (a,b,t)\mapsto(a,b)\ }X_9
 \longrightarrow0.
\]
\end{proposition}
\begin{proof}
Para \(w=(e,f)\), la identidad de cociclo se escribe
\[
 c_{\mathrm W}(u,v)+c_{\mathrm W}(u+v,w)
 =bc+(b+d)e
 =de+b(c+e)
 =c_{\mathrm W}(v,w)+c_{\mathrm W}(u,v+w).
\]
Prueba la asociatividad de \eqref{iv:eq:heisenberg-ley}; la identidad
es \((0,0,0)\) y la inversa de \((a,b,t)\) es
\((-a,-b,-t+ab)\). Los tres residuos son libres, de donde el orden
\(9^3=729\). De \(ZX=\omega XZ\) resulta
\[
 (X^aZ^b)(X^cZ^d)=\omega^{bc}X^{a+c}Z^{b+d},
\]
que demuestra la propiedad de representación. Si
\(\omega^tX^aZ^b=I\), su acción sobre \(\delta_r\) obliga primero
a \(a=0\). La evaluación en \(r=0\) da \(t=0\), y la
evaluación en \(r=1\) da \(b=0\), pues \(\omega\) tiene orden
nueve. Por ello la representación es fiel. La comparación de los
productos en los dos órdenes da el factor \(bc-ad\).
Un elemento central debe satisfacer \(bc-ad=0\) para todo \((c,d)\).
Tomando sucesivamente \((c,d)=(1,0),(0,1)\), se obtiene
\(b=a=0\). Esto identifica el centro y prueba la exactitud.
\end{proof}

La fase de Wilson para la orientación
\eqref{iv:eq:weyl-area-orientada} es
\(\omega^{\sigma_{\mathrm{ar}}(u,v)}\); el conmutador de la
sección ordenada \(X^aZ^b\) publica su inversa. En particular,
\(c_{\mathrm W}(u,v)-c_{\mathrm W}(v,u)
=-\sigma_{\mathrm{ar}}(u,v)\).
Esta identidad relaciona dos convenciones determinadas, sin identificar
el cociclo polarizado con la forma alternante. La extensión pertenece
al grupo de traslaciones del soporte APP. Su representación actúa sobre
la fibra de fase; el registro de la trayectoria y los acarreos conservan
su residencia en el estado TPK que produjo esa fibra. El orden \(729\)
del grupo y el cardinal \(729\) del ambiente ternario conservan los
tipos distintos establecidos en \eqref{iv:eq:hilbert-finito}.
El centro \(\ZZ/9\ZZ\) de esta extensión registra la fase
del conmutador. La memoria \(C_{12}\) de la sucesión
\(0\to C_{12}\to C_{108}\to C_9\to0\), construida en la
proposición~\ref{prop:k-extension}, pertenece al calendario
dodecafásico. Cada extensión conserva su ley de composición y
su proyección propias.

\begin{proposition}[Álgebra generada por fase y traslación]
\label{iv:prop:weyl-completo}
Los \(81\) operadores \(X^aZ^b\), con \(0\leq a,b<9\), forman
una base de \(\mathcal A_1=M_9(\CC)\). En particular, la
representación de \eqref{iv:eq:weyl-operadores} es irreducible.
\end{proposition}
\begin{proof}
Si \(a\neq0\pmod9\), la matriz \(X^aZ^b\) carece de entradas
diagonales. Si \(a=0\), su traza es
\(\sum_{r=0}^8\omega^{br}\), igual a \(9\) para \(b=0\) y a
\(0\) en los demás casos. Por \eqref{iv:eq:weyl-relacion}, el
producto \((X^aZ^b)^*X^{a'}Z^{b'}\) es un factor de módulo uno
por \(X^{a'-a}Z^{b'-b}\). Por tanto,
\[
 \tau_1\bigl((X^aZ^b)^*X^{a'}Z^{b'}\bigr)
 =\delta_{a,a'}\delta_{b,b'}.
\]
Son \(81\) matrices linealmente independientes en un espacio de
dimensión \(81\). Generan toda el álgebra matricial, cuyo conmutante
está formado por los escalares; de ello resulta la irreducibilidad.
\end{proof}

\begin{theorem}[Unicidad con carácter central primitivo fijado]
\label{iv:thm:weyl-unicidad}
Toda representación unitaria irreducible del grupo presentado por
\(X^9=Z^9=c^9=I\), \(c\) central y \(ZX=cXZ\), en la que
\(c\) actúa como \(\omega I\), es unitariamente equivalente a
\eqref{iv:eq:weyl-operadores}.
\end{theorem}
\begin{proof}
Sea \(v\neq0\) un autovector de \(Z\), de autovalor \(\lambda\).
Como \(Z^9=I\), \(\lambda\) es una potencia de \(\omega\).
Los vectores \(X^kv\), \(0\leq k<9\), tienen autovalores
\(\omega^k\lambda\), distintos porque el carácter es primitivo.
Son ortogonales, y su espacio generado es invariante por \(X\),
\(Z\) y el centro. La irreducibilidad obliga a que sea todo el
espacio. Reordenando la base para comenzar en el autovalor \(1\),
\(Z\) es diagonal con entradas \(1,\omega,\ldots,\omega^8\),
y \(X\) es la traslación cíclica. La normalización de la base
produce la equivalencia unitaria afirmada.
\end{proof}

La primitividad del carácter y la irreducibilidad forman parte del
enunciado. A profundidad \(d\), los operadores que actúan en factores
distintos conmutan, y los productos tensoriales de las bases anteriores
generan \(\mathcal A_d\). La estructura de fase queda así
representada a cada nivel sin identificar el retorno de fase con el
retorno del estado genealógico completo.

\subsection{Dos elevaciones y dos normalizaciones diferentes}

La factorización \eqref{iv:eq:factorizacion-hilbert} permite conservar
toda la fibra nueva o fijar una de sus etiquetas. Para
\(p_j=|\delta_j\rangle\langle\delta_j|\), definimos
\begin{equation}
 \iota_d(a)=a\otimes I_9,
 \qquad \jmath_{d,j}(a)=a\otimes p_j.
 \label{iv:eq:elevaciones}
\end{equation}
\begin{proposition}[Conservación de unidad y traza]
\label{iv:prop:unidad-traza}
La inclusión \(\iota_d\) es unital y conserva la traza normalizada.
La inclusión \(\jmath_{d,j}\) es una inclusión en una esquina y
satisface
\[
 \tau_{d+1}(\iota_d(a))=\tau_d(a),\quad \iota_d(I)=I,
 \qquad
 \tau_{d+1}(\jmath_{d,j}(a))=\frac{\tau_d(a)}9,
 \quad \jmath_{d,j}(I)=I\otimes p_j.
\]
\end{proposition}
\begin{proof}
La traza de un producto tensorial es el producto de trazas.
Como \(\operatorname{Tr}(I_9)=9\) y
\(\operatorname{Tr}(p_j)=1\), dividir por \(9^{d+1}\)
produce las dos fórmulas. Las imágenes de la identidad se leen
directamente en \eqref{iv:eq:elevaciones}.
\end{proof}

La inclusión unital conserva el conjunto de las \(9\) fases y su
normalización global. La esquina conserva una continuación marcada y
su proyector de soporte. Ambas operaciones son legítimas, pero
transmiten datos distintos a la siguiente profundidad.

\subsection{Cierre inductivo y realización tracial}

\begin{theorem}[Álgebra uniformemente hiperfinita nonádica]
\label{iv:thm:uhf}
El cierre normado de la unión mediante las inclusiones \(\iota_d\)
es
\[
 \mathcal A_{9^\infty}
 =\varinjlim(\mathcal A_d,\iota_d)
 \cong\bigotimes_{n=1}^{\infty}M_9(\CC).
\]
Es un álgebra UHF de tipo \(9^\infty=3^\infty\), con una única
traza normalizada \(\tau\), cuya restricción a \(\mathcal A_d\)
es \(\tau_d\).
\end{theorem}
\begin{proof}
Las inclusiones son inyectivas y unitales. La clausura de la unión
creciente de álgebras matriciales es precisamente la construcción UHF.
La compatibilidad de la proposición \ref{iv:prop:unidad-traza}
define \(\tau\) sobre la unión; positividad y norma uno permiten
extenderla al cierre. Toda traza normalizada del límite se restringe
a la única traza normalizada de cada álgebra matricial. Como la unión
es densa, esas restricciones determinan la traza de manera única.
\end{proof}

Sea \(\mathcal H_\tau\) la completación de
\(\mathcal A_{9^\infty}\) para
\(\langle a,b\rangle_\tau=\tau(a^*b)\), y sea \(\pi_\tau\)
su representación por multiplicación izquierda. El vector definido
por la identidad se denota \(\Omega_\tau\).

\begin{theorem}[Factor tracial del refinamiento completo]
\label{iv:thm:factor-tracial}
La clausura débil
\(\mathcal R_9=\pi_\tau(\mathcal A_{9^\infty})''\)
es un factor hiperfinito de tipo \(II_1\). Las inclusiones
nonádicas conservan en él la unidad y la traza normalizada.
\end{theorem}
\begin{proof}
La traza producto se extiende a una traza normal fiel y finita
en la representación tracial. Para ver que el centro es trivial,
sea \(E_d\) la expectativa tracial sobre \(\pi_\tau(\mathcal A_d)\).
En un tensor de longitud \(n>d\), toma la traza de los últimos
\(n-d\) factores. Estas aplicaciones se prolongan por continuidad
a las proyecciones ortogonales correspondientes en \(L^2\).
Si \(z\) es central, \(E_d(z)\) conmuta con
\(\pi_\tau(\mathcal A_d)\), por la bimodularidad de \(E_d\).
Es por ello escalar e igual a \(\tau(z)I\). La densidad de la
unión matricial implica \(E_d(z)\to z\) en \(L^2\), de donde
\(z=\tau(z)I\). El álgebra es, pues, un factor finito.
Contiene unitariamente matrices de tamaños \(9^d\) sin cota,
por lo que no es un factor matricial finito. Es de tipo \(II_1\).
Es hiperfinita porque esa misma unión de matrices es débilmente
densa. La conservación de unidad y traza procede de cada inclusión
finita y se mantiene al tomar las clausuras.
\end{proof}

\subsection{Rangos enteros y dimensión tracial continua}

En \(\mathcal A_d\), una proyección \(P\) tiene dimensión
normalizada \(\tau_d(P)=\operatorname{rank}(P)/9^d\).
Los rangos siguen siendo enteros a toda profundidad finita. Sus
normalizaciones compatibles permiten el siguiente paso al límite.

\begin{proposition}[Realización de cada dimensión tracial]
\label{iv:prop:dimension-tracial}
Para cada \(t\in[0,1]\) existe una proyección
\(P_t\in\mathcal R_9\) con \(\tau(P_t)=t\).
\end{proposition}
\begin{proof}
Sea \(m_d=\lfloor9^dt\rfloor\). Se tiene
\(0\leq m_{d+1}-9m_d\leq8\). Elegida una proyección diagonal
de rango \(m_d\), su inclusión tiene rango \(9m_d\); se
añaden \(m_{d+1}-9m_d\) posiciones diagonales en su complemento.
Se obtiene así una sucesión creciente de proyecciones dentro de
\(\mathcal R_9\). Su límite fuerte es una proyección \(P_t\),
y la normalidad de la traza da
\[
 \tau(P_t)=\lim_{d\to\infty}\frac{\lfloor9^dt\rfloor}{9^d}=t.
\]
\end{proof}

Aquí \(t\) designa la coordenada de la dimensión que se representa
en el cierre, no una entrada del generador APP. La afirmación
describe la amplitud de la realización tracial resultante. Para
compararla con un observable determinado sigue siendo necesario
identificar la proyección y la información genealógica que conserva.

Por contraste, las isometrías
\(V_{d,j}\xi=\xi\otimes\delta_j\) realizan
\(\jmath_{d,j}(a)=V_{d,j}aV_{d,j}^*\). En el límite de una
ruta marcada, los espacios finitos son densos. Las esquinas
aproximan todos los operadores de rango finito; su cierre normado
es \(\mathcal K(\mathcal H_\infty)\), y su bicommutante es
\(B(\mathcal H_\infty)\), de tipo \(I_\infty\). Esta
diferencia conserva el contenido operatorio de las dos elecciones
de \eqref{iv:eq:elevaciones}: la selección de una continuación y
la prolongación unital de toda la fibra tienen límites diferentes.

La realización hilbertiana, sus operadores de fase y el cierre
tracial quedan construidos en esta sección, con sus elecciones de
inclusión explícitas. Constituyen el
soporte operatorio que se utilizará después para el intercambio de
coordenadas y para el Hamiltoniano compacto.


% SOURCE sections/trit_projection.tex
\section{Orientación trítica, prolongación compatible y doble proyección}
\label{sec:trit-projection}

El carácter generado de las coordenadas reales se expresa en la relación
entre profundidad aritmética y evaluación. El dato primitivo es una historia
admisible de operaciones con sus residuos, cocientes y orientaciones. Una
lectura arquimediana asocia a sus prefijos intervalos cada vez más estrechos;
una lectura incidencial conserva relaciones de frontera del mismo estado.
La doble proyección del corpus \cite{hmtI,HMT-IX} reúne esas dos operaciones
sin convertir una evaluación escalar en la totalidad de la historia. El
artículo utiliza ese antecedente para estudiar las realizaciones posteriores
del registro dodecafásico.

\subsection{Cilindros posicionales y determinación de la coordenada real}

Para un bloque ternario \(b=(b_1,\ldots,b_6)\), con representantes
\(b_j\in\{0,1,2\}\) de \(\mathbb F_3\), la evaluación entera
\(\nu(b)=\sum_{j=1}^6b_j3^{6-j}\) pertenece a
\(\{0,\ldots,728\}\). Si \(b_{n+1}\) es el bloque admisible siguiente,
la prolongación posicional satisface
\[
 N_{n+1}=729N_n+\nu(b_{n+1}),\qquad
 a_n=\frac{N_n}{729^n},\qquad b_n^+=\frac{N_n+1}{729^n}.
\]
Aquí \(N_0\) es la parte entera que determina el lector regional. La letra
\(b_n^+\) designa un extremo real y se distingue del bloque ternario \(b\).
El cilindro de representación es \([a_n,b_n^+)\); su clausura se utiliza
en el argumento de existencia del límite. El acarreo conserva el paso de
un prefijo al anterior mediante división euclídea.

\begin{proposition}[Límite de los centros cilíndricos]
Para una prolongación compatible, las clausuras de los cilindros están
encajadas y determinan una única coordenada \(x\). Si
\(c_n=(a_n+b_n^+)/2\), entonces
\[
 x=\lim_{n\to\infty}a_n=\lim_{n\to\infty}b_n^+
 =\lim_{n\to\infty}c_n,\qquad
 |x-c_n|\leq\frac1{2\,729^n}.
\]
\end{proposition}
\begin{proof}
La desigualdad \(0\leq\nu(b_{n+1})<729\) sitúa el cilindro hijo en
la clausura del padre. Sus extremos son racionales, su anchura es
\(729^{-n}\) y permanecen dentro del primer intervalo cerrado acotado.
El encajamiento y la anchura decreciente producen una única clase de
Cauchy de extremos. En la realización arquimediana, esa clase tiene el
representante \(x\); la distancia al punto medio es a lo sumo la mitad
de la anchura. La convención de frontera asigna una escritura cuando dos
expansiones posicionales representan el mismo límite.
\end{proof}

La media aritmética interviene, por tanto, en una evaluación precisa: el
centro de cada intervalo que acota la coordenada. El límite conserva la
exactitud del valor aunque la anchura desaparezca. La representación
enriquecida conserva adicionalmente la sucesión de prefijos y su incidencia.
La proyección numérica y la proyección excepcional son dos aplicaciones
con codominios diferentes; una conserva el valor límite y la otra transporta
las relaciones que habilitan la construcción reticular. El paso a
\(\RR\) se sitúa así en la realización del lector, después de la
generación aritmética de los datos que evalúa.

Para hacer explícita la compatibilidad, sean \(p_{n+1,n}\) las restricciones
del estado enriquecido y \(x_n=p_{n+1,n}(x_{n+1})\). Si \(I_n(x_n)\)
es el cilindro numérico y \(\mathcal B_n(x_n)\) la incidencia de frontera,
los dos transportes conservan simultáneamente
\[
 \overline{I_{n+1}(x_{n+1})}\subseteq\overline{I_n(x_n)},\qquad
 r_{n+1,n}\mathcal B_{n+1}(x_{n+1})=\mathcal B_n(x_n).
\]
La segunda igualdad utiliza los mapas incidenciales del corpus; la
proposición~\ref{prop:lf-incidence-naturality} explicita su realización
mediante marcos y códigos transportados. Sus cinco
construcciones consustanciales ---espectral, solenoidal, gauge, cohomológica
y determinantal--- pertenecen a ese mismo sistema compatible. El estudio
geométrico posterior conserva esta unidad de origen al especificar cuál
de sus observaciones utiliza en cada teorema.

\subsection{Trayectorias orientadas y temporalidad de la actualización}

Sea \(x\xrightarrow{\rm adm}y\) la relación de admisibilidad del TPK en
el espacio enriquecido \(\widehat X\). Una trayectoria parametrizada por
un conjunto discretamente ordenado \(\mathcal T\) es
\[
 \gamma:\mathcal T\longrightarrow\widehat X,
 \qquad\gamma(t)\xrightarrow{\rm adm}\gamma(t^+).
\]
La relación de compatibilidad determina las transiciones; la parametrización
permite ordenarlas. En el espacio \(E=\mathcal T\times\widehat X\),
con \(p(t,x)=t\), la trayectoria determina la sección
\(\widetilde\gamma(t)=(t,\gamma(t))\), que satisface
\(p\circ\widetilde\gamma=\id_{\mathcal T}\).
Esta precisión conserva la diferencia entre un estado, su trayectoria y
la coordenada utilizada para describir la evolución.

La firma trítica \(\tau\in\{+1,0,-1\}\) tipa los generadores mediante
\[
 J_\tau^2=-\tau I.
\]
En el lector temporal del corpus, \(+1\) representa retorno y memoria;
\(0\), actualización de umbral; y \(-1\), apertura de prolongaciones
conjugadas. Su interpretación como pasado, presente y futuro se refiere
a las funciones del antecedente conservado, la transición efectiva y la
continuación admisible. Los tres tiempos verbales adquieren así un
contenido operatorio sobre historias.

Si \(\mathcal P_n(x)\) es el conjunto de continuaciones admisibles de
longitud \(n\), la supresión del último paso define
\(r_n:\mathcal P_{n+1}(x)\to\mathcal P_n(x)\). Una continuación
ilimitada es una familia \((\gamma_n)\) con
\(r_n(\gamma_{n+1})=\gamma_n\). La restricción reconstruye antecedentes;
la actualización \(U_t\) actúa en la transición; la prolongación mantiene
la frontera del prefijo. La existencia de estas familias se utiliza en
los dominios supervivientes fijados por el teorema de prolongación del
corpus.

Sobre una firma finita \(\sigma=(\tau_1,\ldots,\tau_n)\), la inversión
orientada es
\[
 C_{\rm rev}\sigma=(-\tau_n,\ldots,-\tau_1),
 \qquad C_{\rm rev}^2\sigma=\sigma.
\]
La trayectoria conjugada conserva también las transformaciones de hoja,
posición y frontera de su transporte. Por ello la bidireccionalidad se
expresa en la historia completa, mientras esta fórmula publica su firma.
El retorno de fase de \(\Gamma_9\) admite un incremento de memoria:
estados con fase idéntica pueden ocupar profundidades diferentes.

\subsection{Álgebras cuadráticas y actividad del umbral parabólico}

En un régimen fijo, la conjugación de \(aI+bJ_\tau\) cambia el signo
de \(b\), y su norma algebraica es
\[
 (aI+bJ_\tau)(aI-bJ_\tau)=(a^2+\tau b^2)I.
\]
Los tipos elíptico, parabólico e hiperbólico corresponden a las tres
relaciones cuadráticas. Sus grupos locales se representan por
\[
 e^{tJ_{+1}}=\cos t\,I+\sin t\,J_{+1},\quad
 e^{tJ_0}=I+tJ_0,\quad
 e^{tJ_{-1}}=\cosh t\,I+\sinh t\,J_{-1}.
\]
La prueba separa términos pares e impares en la serie exponencial; en el
tipo nilpotente sólo permanecen los grados cero y uno. Cada familia
satisface la ley de composición en su parámetro y la inversión
\(e^{-tJ_\tau}=(e^{tJ_\tau})^{-1}\).

El valor trítico cero permite, en particular, un generador distinto de
cero con cuadrado nulo. Para
\(N=\left(\begin{smallmatrix}0&1\\0&0\end{smallmatrix}\right)\),
\[
 e^{tN}(q,p)=(q+tp,p).
\]
El umbral tiene una acción lineal efectiva sobre su fibra. Su parámetro,
su desplazamiento y la memoria de la transición permanecen determinados
aunque la firma de régimen sea cero. Éste es el contenido preciso que
permite desarrollar narrativamente el presente como actualización activa.

La forma elíptica positiva tiene niveles circulares en coordenadas
ortonormales y elípticos tras una transformación lineal; la forma hiperbólica
posee dos direcciones isotrópicas y niveles hiperbólicos; el tipo parabólico
posee una forma degenerada y transporte unipotente. Una geometría esférica
requiere además el espacio homogéneo o la métrica que la realiza. La
clasificación anterior fija los generadores que pueden intervenir en esa
realización y mantiene separados el tipo algebraico y la curvatura de un
espacio posterior.

El calendario de modos concreta la conexión con la autoescala. La regla
\(+1:m\mapsto\Sigma\), \(-1:m\mapsto\Pi\), \(0:m\mapsto m\)
aplicada al bloque cíclico \((+1,-1,0)\) produce
\((\Sigma,\Pi,\Pi)\). Contar sus tres transiciones, en el orden
\((\Sigma,\Pi)\), determina
\begin{equation}
 F_{\rm av}=\begin{pmatrix}0&1\\1&1\end{pmatrix}.
 \label{eq:geo-fav}
\end{equation}
Su rayo propio positivo, normalizado como \((1,x)^\mathsf T\), satisface
\(x=\lambda\), \(1+x=\lambda x\), y por ello
\(x^2-x-1=0\). La raíz positiva es la coordenada de autoescala publicada
como \(\varphi\). La matriz nace del calendario de operaciones; la
expresión radical y las identidades exponenciales siguientes son sus
realizaciones posteriores.

\input{sections/euler_trit.tex}


% SOURCE sections/euler_trit.tex
% Articulación posterior a la generación de pi, phi y e.
% Contenido acumulativo; no reemplaza el desarrollo previo del TRIT.
\subsection{Realizaciones exponenciales de las coordenadas generadas}
\label{geo-euler-trit-articulacion}

Las coordenadas regionales ya determinadas permiten reconocer las
realizaciones concretas de la familia cuadrática del TRIT. El orden es
sustantivo: la construcción de las regiones precede a la evaluación de una
exponencial en los parámetros \(\pi\) o \(2\log\varphi\). Estas expresiones
identifican la función de las coordenadas en operadores construidos desde
APP y TPK; no intervienen en la selección de los prefijos que las producen.

\subsubsection{Ciclo ternario, transporte nilpotente e incidencia de hojas}

La multiplicación \(a(r)=4r\) en \(\mathbb Z/9\mathbb Z\) tiene orden
\(3\). Sobre las unidades determina las órbitas \(\{1,4,7\}\) y
\(\{2,8,5\}\). En el módulo de aumento de cualquiera de ellas, formado por
las combinaciones enteras cuyos coeficientes suman cero, la base
\((e_r-e_{4r},e_{4r}-e_{7r})\) da
\[
 C=\begin{pmatrix}0&-1\\1&-1\end{pmatrix},
 \qquad C^2+C+I=0.
\]
La forma restringida tiene matriz de Gram
\(\bigl(\begin{smallmatrix}2&-1\\-1&2\end{smallmatrix}\bigr)\).
En las evaluaciones residuales módulo \(9\), la acción distingue las hojas:
\(S(ax,ay)=aS(x,y)\), mientras \(P(ax,ay)=a^{-1}P(x,y)\).
La realización cíclica conserva, por tanto, una orientación contragrediente
entre suma y producto. Los cocientes del levantamiento entero se conservan
en el registro aritmético correspondiente.

El transporte parabólico elemental se representa mediante
\[
 N=\begin{pmatrix}0&1\\0&0\end{pmatrix}.
\]
La incidencia areal--volumétrica ya construida en \eqref{eq:geo-fav} proporciona
\[
 F_{\rm av}=\begin{pmatrix}0&1\\1&1\end{pmatrix},
 \qquad A=F_{\rm av}^2=\begin{pmatrix}1&1\\1&2\end{pmatrix}.
\]
Los tres operadores actúan en espacios reales declarados: el módulo de
aumento tensorizado con \(\mathbb R\), el plano del transporte nilpotente y
el plano de incidencia modal. Compartir dimensión matricial no identifica
estos espacios ni sus coordenadas.

\begin{proposition}[Generadores concretos de los tres tipos]
\label{geo-euler-trit-generadores}
Los operadores
\[
 J_{\mathrm e}=\frac{2C+I}{\sqrt3},\qquad
 J_{\mathrm p}=N,\qquad
 J_{\mathrm h}=\frac{2A-3I}{\sqrt5}
\]
satisfacen, respectivamente,
\[
 J_{\mathrm e}^{\,2}=-I,\qquad
 J_{\mathrm p}^{\,2}=0,\qquad
 J_{\mathrm h}^{\,2}=I.
\]
\end{proposition}
\begin{proof}
La relación de \(C\) implica \((2C+I)^2=-3I\).
La multiplicación directa da \(N^2=0\).
Finalmente, \(A\) tiene traza \(3\) y determinante \(1\), por lo que
\(A^2-3A+I=0\) y \((2A-3I)^2=5I\).
\end{proof}

\begin{proposition}[Tres evaluaciones de la identidad exponencial]
\label{geo-euler-trit-evaluaciones}
En las realizaciones anteriores,
\begin{equation}
 \boxed{
 \exp(\pi J_{\mathrm e})+I=0,\qquad
 \exp(J_{\mathrm p})=I+J_{\mathrm p},\qquad
 \exp(2\log\varphi\,J_{\mathrm h})=F_{\rm av}^{\,2}.}
 \label{geo-euler-trit-tres-evaluaciones}
\end{equation}
\end{proposition}
\begin{proof}
La especialización elíptica de la exponencial trítica da
\(\cos\pi\,I+\sin\pi\,J_{\mathrm e}=-I\).
La nilpotencia elimina exactamente los términos de grado al menos \(2\)
en la segunda expresión. Para la tercera, \(\varphi^2=\varphi+1\) implica
\[
 \cosh(2\log\varphi)=\frac32,\qquad
 \sinh(2\log\varphi)=\frac{\sqrt5}{2}.
\]
Por tanto, la exponencial vale
\(\frac32I+\frac{\sqrt5}{2}J_{\mathrm h}=A\).
\end{proof}

La primera igualdad representa la media vuelta elíptica; la vuelta completa
satisface \(\exp(2\pi J_{\mathrm e})=I\). La segunda mantiene un término
lineal no nulo: el régimen parabólico expresa transporte, no ausencia de
evolución. La tercera realiza un avance hiperbólico unimodular. Sus
autovalores son \(\varphi^2\) y \(\varphi^{-2}\); una dirección se expande y
la otra se contrae conservando el determinante. Las tres evaluaciones
emplean parámetros distintos de una familia exponencial común.

En esa familia, \(e\) interviene mediante la evaluación escalar
\(\exp(I)=eI\) y la ley \(\exp(sI)\exp(tI)=\exp((s+t)I)\).
Su papel no se restringe al tipo parabólico. La coordenada \(e\), obtenida
previamente de su región, es reconocida aquí como la evaluación unitaria
de la exponencial; esta identidad no sustituye su genealogía coinductiva.

\subsubsection{Resolución polinómica de los regímenes}

Las tres realizaciones admiten una presentación compacta sobre la suma
directa de sus espacios. Defínanse
\[
 \mathbb J=J_{\mathrm e}\oplus J_{\mathrm p}\oplus J_{\mathrm h},
 \qquad Q=\mathbb J^2.
\]
Las relaciones anteriores implican
\[
 \mathbb J^6=\mathbb J^2,\qquad Q^3=Q.
\]
El uso de \(Q\) mantiene separado este cuadrado operatorial del registro
dodecafásico \(K\).

\begin{proposition}[Proyectores polinómicos de régimen]
\label{geo-euler-trit-proyectores}
Las aplicaciones
\[
 p_{\mathrm e}=\frac{Q^2-Q}{2},\qquad
 p_{\mathrm p}=I-Q^2,\qquad
 p_{\mathrm h}=\frac{Q^2+Q}{2}
\]
son idempotentes mutuamente ortogonales, suman \(I\) y recuperan los tres
sumandos de la representación.
\end{proposition}
\begin{proof}
Los tres polinomios toman los valores \((1,0,0)\), \((0,1,0)\) y
\((0,0,1)\) en los puntos \(-1,0,+1\), respectivamente. Las diferencias
entre sus cuadrados y ellos mismos, así como sus productos cruzados, son
múltiplos de \(X^3-X\). Sustituir \(Q\), que anula ese polinomio, demuestra
las identidades. En los tres bloques, \(Q\) actúa como \(-I,0,I\).
\end{proof}

Se obtiene así
\begin{align}
 \exp(t\mathbb J)
 ={}&p_{\mathrm e}(\cos t\,I+\sin t\,\mathbb J)
   +p_{\mathrm p}(I+t\mathbb J)\notag\\
   &+p_{\mathrm h}(\cosh t\,I+\sinh t\,\mathbb J).
 \label{geo-euler-trit-exponencial-total}
\end{align}
Cada sumando conserva su relación cuadrática y la conjugación invierte
su evolución. La representación total reúne los tres regímenes, pero no
define por sí sola un operador de transición entre ellos. Tal transición
requiere los mapas de transporte del TPK con sus dominios, orientación y
memoria. La articulación exponencial permite reconocer sus tipos sin
reemplazar esa dinámica por una identificación de las fibras.

% PROCEDENCIA FOCAL (no impresa)
% Resultado recuperado: tres generadores y tres evaluaciones exponenciales.
% Propietario completo:
% BASE/manuscrito/sections/hmt/09b_algebras_cuadraticas.tex:289-516.
% Antecedente causal de C y acción contragrediente:
% BASE/manuscrito/sections/hmt/03d_esqueleto_ciclotomico_app.tex:13-47,145-246.
% Los tres tipos y la orientación están en:
% BASE/manuscrito/sections/hmt/02_trit_geometrias.tex:83-196.
% Resolución polinómica recuperada:
% output/INVESTIGACION_HOLONOMIA_NONADICA_CRISTAL_APERIODICO_20260903/
% ALGEBRA_HOLONOMICA_UNIVERSAL_Y_EULER_TRITICO.md, secciones 8-11 y 22.
% Los cuatro propietarios se leyeron completos. BASE designa:
% /Users/ruben/Documents/HMT2/
% EL_CIERRE_HOLOGRAFICO_DEL_INFINITO_SUCESOR_102_CAPITULOS_2026-08-28_EN_TRABAJO/
% No se incorpora la identidad P9(alpha)=delta del documento especializado:
% esa identificación no pertenece a este bloque.
% Tampoco se identifica una dirección nilpotente con la vacancia electrónica,
% ni se presume una representación global del transporte por la suma directa.
% COMPROBACIONES: Python estándar, enteros y Fraction.
% C^2+C+I=0; (2C+I)^2=-3I; N^2=0; (2Fav^2-3I)^2=5I.
% 81 pares de residuos verifican las dos acciones contragredientes.
% Los 3 proyectores se verificaron en Q=-1,0,+1; prueba general en texto.
% 6 labels únicos y entornos equilibrados. Sin compilación.


% SOURCE sections/incidence_transport.tex
\section{Transporte del bloque modular al cociente incidencial}
\label{sec:incidence-transport-new}
El bloque central de APP y la incidencia de caras y vértices proporcionan
dos realizaciones de una misma involución. Su relación se expresa mediante
un entrelazador rectangular, antes de efectuar cualquier interpretación
métrica. La congruencia diagonal adyacente
\[
 k^2\equiv(k+1)^2\pmod9
 \quad\Longleftrightarrow\quad 2k+1\equiv0\pmod9
\]
selecciona \(k=4\) en la ventana declarada. El bloque correspondiente es
\[
 B_c=\begin{pmatrix}7&2\\2&7\end{pmatrix}
 =7I_2+2R_c=9P_++5P_-,\qquad
 R_c=\begin{pmatrix}0&1\\1&0\end{pmatrix},\quad
 P_\pm=\frac{I_2\pm R_c}{2}.
\]
La diferencia entre los dos valores propios es cuatro. El transporte
incidencial determina a continuación cómo se conserva esa diferencia
cuando cambian las multiplicidades de los sectores.

Las caras del cubo se indexan por \((i,\epsilon)\), con
\(i\in\{1,2,3\}\) y \(\epsilon\in\{-1,1\}\); sus vértices se
indexan por \(\nu\in\{-1,1\}^3\). Definimos
\[
 M_{(i,\epsilon),\nu}=\mathbf1_{\{\nu_i=\epsilon\}}.
\]
Sean \(R_F\) la oposición de caras y \(R_V\) la inversión de vértices
\(\nu\mapsto-\nu\). En cada entrada se verifica
\[
 \mathbf1_{\{\nu_i=-\epsilon\}}
 =\mathbf1_{\{(-\nu)_i=\epsilon\}},\qquad R_FM=MR_V.
\]
La matriz \(M\) transporta, por tanto, la involución completa.

\begin{theorem}[Descenso incidencial del bloque central]
Sean \(B_F=7I_6+2R_F\) y \(B_V=7I_8+2R_V\). Entonces
\[
 B_FM=MB_V.
\]
Si \(u=(1,1,1,1,1,1)^{\mathsf T}\),
\(P_u=uu^{\mathsf T}/6\) y \(P_-=(I_6-R_F)/2\), el cociente
\[
 Q_\square=\RR^6/\ker(MM^{\mathsf T})
 \simeq\RR u\oplus E_-
\]
tiene dimensión \(1+3\), y el operador inducido es
\[
 \overline B_F=9P_u+5P_-.
\]
\end{theorem}
\begin{proof}
La primera identidad resulta de multiplicar \(R_FM=MR_V\) por dos
y sumar \(7M\). Cada cara contiene cuatro vértices; dos caras adyacentes
comparten dos y dos caras opuestas no comparten ninguno. De esta cuenta se
obtiene, entrada por entrada,
\[
 MM^{\mathsf T}=12P_u+4P_-.
\]
El primer proyector tiene rango uno y el segundo rango tres; sus imágenes
son ortogonales. Los valores propios son \(12,4,0\), con multiplicidades
\(1,3,2\). La oposición actúa como \(+1\) en \(\RR u\) y como
\(-1\) en \(E_-\), y conserva el núcleo de la matriz de Gram. Así,
\(B_F\) desciende al cociente y actúa por nueve y cinco en esos dos
sumandos. Su polinomio característico es
\((\lambda-9)(\lambda-5)^3\).
\end{proof}

La asignación temporal al eje \(\RR u\) fija la forma firmada
\(J_L=P_--P_u\). Sobre el cociente, \(I_Q=P_u+P_-\), de modo que
\[
 J_L=\frac{7I_Q-\overline B_F}{2}.
\]
Esta igualdad relaciona la forma de signatura \((3,1)\), con la
convención temporal indicada, y el operador modular transportado. La
matriz de Gram inicial conserva su positividad; la firma se introduce
en la realización posterior mediante el signo asignado a cada sector.
Las dilataciones nueve y cinco siguen siendo las del bloque: la identidad
afín determina la forma y no convierte al bloque en una isometría.

\subsection{Agregación aritmética, norma reticular y funcional de ruta}
La conservación del entero previo a la reducción se hace visible al sumar
las dos hojas. Las matrices de agregación modular son
\[
 M_\Pi=\begin{pmatrix}4&5&6\\5&4&6\\6&6&9\end{pmatrix},\qquad
 M_\Sigma=\begin{pmatrix}5&6&4\\6&4&5\\4&5&6\end{pmatrix}.
\]
Sus sumas son cincuenta y uno y cuarenta y cinco. Al reconstruir los
nueve elementos de cada bloque, el entero y su acarreo satisfacen
\[
 2025-810=(459-405)+9(174-45)=54+9\cdot129.
\]
El agregado residual y el cociente constituyen dos coordenadas de esta
igualdad. En la construcción reticular posterior, el valor cincuenta y
cuatro aparece como norma de la cámara marcada
\(v_\alpha=4\rho_1+\rho_2+\cdots+\rho_{12}\), para vectores
ortogonales \(\rho_j^2=2\). En la realización material aparece como
evaluación del funcional de ruta \(D_1\) del perfil \((80,54,6)\).
Cada aparición conserva su aplicación de origen: agregación con acarreo,
forma reticular y evaluación orientada de ruta. La coincidencia de sus
valores adquiere contenido estructural mediante esas aplicaciones, cuyos
dominios y pruebas se desarrollan en sus secciones respectivas.


% SOURCE sections/memory_geometry.tex
\section{Reconstrucción de la geometría positiva desde la memoria nonádica}
\label{sec:memory}

La extracción dodecafásica conserva más información que una dirección de
deformación. La dirección selecciona un movimiento; el registro completo
determina también la partición sobre la que ese movimiento actúa. La
recuperación del registro a partir de sus memorias permite formular esta
distinción como un teorema de reconstrucción. El resultado procede de la
composición de las compresiones nonádicas del corpus \cite{hmtX}; su aplicación
geométrica conserva la carga uniforme y los complementos de ambas etapas.

En esta sección se numeran las coordenadas de cero a once. Sea
\(\mathsf S:\RR^{12}\to\RR^{12}\) el desplazamiento
\((\mathsf Sx)_i=x_{i+1}\), con índices módulo doce. En la carta euclídea
posterior al estado enriquecido, definimos
\[
 T_j=\frac{8I+\mathsf S^j}{9},\qquad
 D_j=\frac{\sqrt8}{9}(\mathsf S^j-I),\qquad j\in\{3,4\}.
\]
El coeficiente ocho cuenta las posiciones ordinarias de la carta nonádica;
la novena posición contiene el transporte. La ortogonalidad del desplazamiento
implica por multiplicación directa
\[
 T_j^*T_j+D_j^*D_j=I.
\]
La primera memoria es el complemento de la primera compresión; la segunda
se toma sobre el estado que ya atravesó esa compresión. Por tanto,
\begin{equation}
 q=\mathbf1^*K,\qquad m_0=D_3K,\qquad m_1=D_4T_3K,
 \qquad MK=(m_0,m_1).
 \label{eq:mem-data}
\end{equation}
El orden de la composición forma parte del objeto. La aplicación
\(M:\RR^{12}\to\RR^{24}\) conserva dos observaciones correlacionadas,
producidas por una misma trayectoria lineal.

\begin{theorem}[Inversión exacta de la publicación de memoria]
\label{thm:memory-inversion}
La aplicación \(K\mapsto(q,MK)\) es un isomorfismo lineal entre
\(\RR^{12}\) y \(\RR\times\operatorname{im}M\). Su inversa se obtiene
mediante
\begin{align}
 T_3^{-1}&=\frac{512I-64\mathsf S^3+8\mathsf S^6-\mathsf S^9}{455},
 \label{eq:memory-T-inverse}\\
 b&=-\frac9{\sqrt8}m_0,\qquad
 c=-\frac9{\sqrt8}T_3^{-1}m_1,\qquad w_i=c_i-b_{i+1},\notag\\
 B_0&=0,\qquad B_i=\sum_{j=0}^{i-1}w_j,\qquad
 K_i=\frac{q+\sum_{j=0}^{11}B_j}{12}-B_i.
 \label{eq:memory-inverse}
\end{align}
\end{theorem}
\begin{proof}
Escribiendo \(X=\mathsf S^3\), se tiene \(X^4=I\) y
\((8I+X)(512I-64X+8X^2-X^3)=4095I\). Dividir por nueve
prueba \eqref{eq:memory-T-inverse}. Todos los operadores utilizados son
polinomios del mismo desplazamiento y conmutan. Así,
\[
 b=(I-\mathsf S^3)K,\qquad c=(I-\mathsf S^4)K,
 \qquad w_i=K_i-K_{i+1}.
\]
La suma telescópica da \(B_i=K_0-K_i\); sumar esta igualdad sobre
las doce posiciones proporciona \eqref{eq:memory-inverse}. Además,
\(MK=0\) equivale a invariancia bajo \(\mathsf S^3\) y
\(\mathsf S^4\). Como \(\mathsf S=\mathsf S^4(\mathsf S^3)^{-1}\),
el núcleo es exactamente \(\RR\mathbf1\). La carga recupera esa
componente y demuestra la inyectividad. La descomposición
\(K=(q/12)\mathbf1+K^\perp\) prueba la sobreyectividad sobre el
codominio declarado.
\end{proof}

La forma geométrica queda determinada por la información que una lectura
comprimida parecía haber desplazado. El complemento vectorial contiene las
diferencias entre posiciones; la carga contiene la componente común. Su
reunión reconstruye las doce coordenadas antes de normalizarlas como
longitudes. Esta afirmación relaciona operaciones efectivas: compresión,
archivo de complementos e inversión. La conservación de una norma expresa
una propiedad distinta y, por sí sola, más débil que esa recuperabilidad.

\subsection{Cono de memoria compatible y estabilidad de los menores}

Sea \(L\) la pseudoinversa de \(M\), con imagen
\(\mathbf1^\perp\). Equivalentemente,
\[
 L=\left(M^*M\big|_{\mathbf1^\perp}\right)^{-1}M^*,
 \qquad LM=P_{11},\qquad K=\frac q{12}\mathbf1+Lm.
\]
La notación de inversa restringida se interpreta como un operador que toma
valores en \(\mathbf1^\perp\). La positividad estricta del registro
equivale a pertenecer al cono abierto
\begin{equation}
 \mathcal M_+=\left\{(q,m)\in\RR\times\operatorname{im}M:
 \frac q{12}\mathbf1+Lm>0\right\}.
 \label{eq:memory-positive-cone}
\end{equation}
Las desigualdades son coordenadas. El cono tiene dimensión doce; su sección
de carga positiva fija tiene dimensión once. En particular, las veinticuatro
entradas de las dos memorias satisfacen relaciones de compatibilidad.

\begin{proposition}[Constante óptima de reconstrucción en norma euclídea]
\label{prop:memory-stability}
Se cumple \(\|L\|=9/4\). Para perturbaciones arbitrarias de las observaciones,
\begin{equation}
 \|\delta K\|^2\leq\frac{|\delta q|^2}{12}
 +\frac{81}{16}\bigl(\|\delta m_0\|^2+\|\delta m_1\|^2\bigr).
 \label{eq:memory-stability}
\end{equation}
\end{proposition}
\begin{proof}
La base de Fourier diagonaliza \(\mathsf S\). Si \(z^{12}=1\),
el autovalor de \(M^*M\) en el modo correspondiente es
\[
 \lambda(z)=\frac8{81}|z^3-1|^2
 +\frac8{81}|z^4-1|^2\left|\frac{8+z^3}{9}\right|^2.
\]
Evaluar los doce modos da
\[
 \operatorname{Spec}(M^*M)=
 \left\{0^{[1]},(16/81)^{[2]},(8/27)^{[2]},(32/81)^{[1]},
 (952/2187)^{[4]},(1256/2187)^{[2]}\right\}.
\]
Los exponentes entre corchetes indican multiplicidad. El menor autovalor
positivo es \(16/81\), de donde \(\|L\|=9/4\). La componente uniforme
\((\delta q/12)\mathbf1\) es ortogonal a \(L\delta m\); la identidad
pitagórica y la norma de \(L\) producen \eqref{eq:memory-stability}.
\end{proof}

Para pasar a la geometría intervalar recuperamos la numeración de uno a doce
y ponemos \(d_j=K_j/q\), \(t_0=0\),
\(t_j=\sum_{r=1}^j d_r\). El cono \eqref{eq:memory-positive-cone}
determina exactamente \(0=t_0<t_1<\cdots<t_{12}=1\), y por tanto
\[
 \Delta_{ij}=t_j-t_i=\sum_{i<r\leq j}d_r>0\qquad(i<j).
\]
Son los menores de la matriz de caminos que se construye en la sección
\ref{app:network-cluster}. Su dependencia respecto de la memoria es
explícita. Los puntos \((t_j,t_j^2)\) quedan igualmente determinados y
producen la realización amplituhedral de rango uno desarrollada después.

Si la raíz del miembro derecho de \eqref{eq:memory-stability} es menor que
\(\min_j K_j\), la reconstrucción perturbada permanece en el cono positivo.
Para el registro publicado, \(\min_jK_j=58\); a carga fija basta
\[
 \sqrt{\|\delta m_0\|^2+\|\delta m_1\|^2}<\frac{232}{9}.
\]
La cota conserva el orden de los vértices y el signo de todos los menores
ordenados. Describe estabilidad de esta carta matemática, con la
normalización indicada. Para observaciones fuera de \(\operatorname{im}M\),
la reconstrucción utiliza la proyección \(MLm\); el residuo
\((I-ML)m\) se conserva separadamente cuando interesa recuperar la
observación íntegra.

\subsection{Recuperación nodal y reducción a la frontera}

La partición obtenida tiene también una realización cuadrática:
\[
 \mathcal E_d(f)=\sum_{j=1}^{12}\frac{|f_j-f_{j-1}|^2}{d_j}
 =f^*L_df.
\]
El laplaciano nodal recupera cada longitud por
\(d_j=-1/(L_d)_{j-1,j}\). Si se conservan únicamente los extremos
globales, la minimización elimina los nodos interiores y produce, en
cambio, la matriz \(\left(\begin{smallmatrix}1&-1\\-1&1\end{smallmatrix}\right)\),
puesto que \(\sum_jd_j=1\). El lector nodal distingue la forma intervalar;
el lector de dos extremos sólo observa su longitud total.

La demostración de esta eliminación se formula íntegramente en la sección
energética: la prolongación afín minimizante \(Hf\) y el detalle \(z\),
nulo en los nodos antiguos, satisfacen
\[
 \mathcal E_{\widetilde d}(Hf+z)
 =\mathcal E_d(f)+\mathcal E_{\widetilde d}(z).
\]
El segundo término es no negativo y se anula exactamente para el detalle
nulo. La lectura efectiva conserva la energía mínima; el par \((f,z)\)
conserva además la configuración refinada. Así, positividad de menores,
positividad energética y recuperabilidad de memoria se relacionan mediante
mapas definidos sobre la misma partición, manteniendo sus respectivos
espacios de llegada.


% SOURCE appendices/network_cluster.tex
% Procedencia HMT: output/EXTREMA_MEDIA_RAZON_NARRACION_Y_REVISION_20260916/
% ARTICULO/ES/source/libro/01_acoplamiento.tex, lib01:cadena-registro y lib01:K.
% Dirección: ARTICULO/ES/source/sections/k_direccion_dimensional.tex.
% Estatuto de las redes, cartas y flujos de este apéndice: FORMALIZACION_NUEVA.
% La extracción de K y el proyector P_3 son RESULTADO_RECUPERADO.
\section{Medición planar de frontera, coordenadas de Plücker y mutaciones positivas}
\label{app:network-cluster}

La publicación dodecafásica del estado enriquecido admite una realización
positiva cuya red, coordenadas y cambios de carta pueden escribirse
completamente. La composición de partida es
\[
 \mathrm{APP}\longrightarrow\mathrm{TRIT}\longrightarrow\mathrm{TPK}
 \longrightarrow x_{\rm term}\longrightarrow\mathcal R_{12}
 \longrightarrow(b^{90},b^{120},Q_{\rm TPK})
 \longrightarrow U_{\rm sgn}\longrightarrow K.
\]
Esta lectura actúa sobre el estado enriquecido que conserva las dos hojas,
la orientación, el acarreo, la memoria y la frontera. El registro ya generado
es
\[
 K=(234,543,140,729,659,824,621,58,914,794,146,601),
 \qquad Q=\sum_{r=1}^{12}K_r=6263.
\]
Su extracción se recupera exactamente de los dos canales y la carga total:
\[
 b^{90}=(-495,-116,-684,108,601,-90,-173,-88,313,560,-397,461),
\]
\[
 b^{120}=(-425,-281,-481,671,-255,30,475,-543,680,251,6,-128).
\]
En las órbitas de posiciones \((r,r+3,r+6,r+9)\), para \(r=1,2,3\),
sean \(B_r=\sum_{j=0}^3b^{120}_{r+3j}\) y
\(d_{r,j}=b^{90}_{r+3j}\). Se obtienen
\[
 s_1=(Q+2B_1+B_2)/3,\quad s_2=s_1-B_1,\quad s_3=s_2-B_2,
\]
\[
 U^{(r)}=(s_r,d_{r,1}-d_{r,3},d_{r,0}+d_{r,2},d_{r,0}-d_{r,2}),
 \qquad K|_r=\tfrac14H_4U^{(r)},
\]
\[
 H_4=\begin{pmatrix}1&1&1&1\\1&1&-1&-1\\1&-1&1&-1\\1&-1&-1&1\end{pmatrix},
 \qquad H_4^2=4I_4.
\]
La red construida a continuación realiza los valores positivos producidos
por esa extracción. Su forma combinatoria es una formalización posterior
del registro, con origen y orientación declarados.

\subsection{Red dirigida acíclica y determinantes de caminos no intersectantes}

Para un registro positivo de longitud \(L\), ponemos
\[
 d_r=K_r/Q>0,\qquad t_0=0,\qquad
 t_j=\sum_{r=1}^jd_r\quad(1\le j\le L).
\]
Así \(t_L=1\). Las \(L\) separaciones producen \(L+1\) puntos marcados;
en el caso dodecafásico son trece puntos. La matriz que los realiza es
\begin{equation}
 C(d)=\begin{pmatrix}1&1&\cdots&1\\t_0&t_1&\cdots&t_L\end{pmatrix}.
 \label{eq:nc-path-matrix}
\end{equation}
El Grassmanniano \(\operatorname{Gr}(2,L+1)\) identifica matrices de rango
dos por cambios invertibles de las dos filas. Su parte positiva está formada
por las clases que admiten todos sus menores ordenados de orden dos
estrictamente positivos.

Construimos una red dirigida con vértices \(a_j,b_j\), para
\(0\le j\le L\), y sumideros \(q_j\). Hay aristas horizontales
\(a_{j-1}\to a_j\) y \(b_{j-1}\to b_j\), de peso uno; aristas verticales
\(a_j\to b_j\), de peso \(d_j\), para \(j\ge1\); y aristas
\(b_j\to q_j\), de peso uno. Las fuentes ordenadas son \(b_0,a_0\).
Las dos filas se dibujan horizontalmente, las aristas verticales descienden
y los sumideros salen por debajo de la fila inferior. Este dibujo es planar
y la orientación es acíclica. La medición de frontera \(M_{ij}\) es la
suma de los productos de pesos de todos los caminos desde la fuente \(i\)
al sumidero \(q_j\).
La red de caminos tiene dos fuentes y \(L+1\) sumideros, es decir,
\(L+3\) conexiones exteriores. La matriz que realiza tiene \(L+1\)
columnas. La semilla plábica introducida después tiene precisamente
\(L+1\) vértices de frontera y constituye otra presentación de sus
coordenadas; estas dos cuentas de frontera se mantienen explícitas.

\begin{theorem}[Medición de frontera y menores del registro]
\label{thm:nc-network}
La medición de frontera de esta red es \(M=C(d)\). Para
\(0\le i<j\le L\),
\begin{equation}
 \Delta_{ij}(C(d))=t_j-t_i=\sum_{r=i+1}^{j}d_r>0.
 \label{eq:nc-minor-gap}
\end{equation}
El mismo número es la suma de pesos de las parejas de caminos
vértice-disjuntos \(b_0\to q_i\), \(a_0\to q_j\).
\end{theorem}
\begin{proof}
La fuente inferior tiene un único camino a cada sumidero, de peso uno.
Un camino desde \(a_0\) a \(q_j\) desciende exactamente una vez, por una
arista de índice \(1\le r\le j\); su peso es \(d_r\). Esto prueba la
fórmula de la matriz. El camino inferior hasta \(q_i\) ocupa todos los
vértices \(b_0,\ldots,b_i\). Por tanto, un camino de la otra fuente a
\(q_j\) es disjunto de él precisamente cuando desciende en
\(i<r\le j\). Su suma de pesos es el miembro derecho de
\eqref{eq:nc-minor-gap}. La pareja con destinos intercambiados siempre
comparte un vértice. El desarrollo del determinante cancela todas las
parejas que se cruzan y conserva exactamente las disjuntas. Aquí se ha
verificado directamente, para la red concreta, el mecanismo determinantal
de caminos no intersectantes.
\end{proof}

La identificación de las mediciones de redes planas con las celdas del
Grassmanniano no negativo proporciona su reconocimiento geométrico posterior;
véase A. Postnikov,
\href{https://arxiv.org/abs/math/0609764}{\emph{Total positivity,
Grassmannians, and networks}}. La prueba anterior da, además, una fórmula
finita para cada menor sin utilizar un valor metrológico.

La familia \(d_r>0\), \(\sum_rd_r=1\), tiene dimensión \(L-1\).
Su aplicación al Grassmanniano es inyectiva: si las filas de \(C(d)\) y
\(C(d')\) generan el mismo plano, las segundas filas difieren por una
transformación afín; las condiciones \(t_0=t'_0=0\) y
\(t_L=t'_L=1\) fijan esa transformación. En particular, para \(L=12\)
se obtiene una familia de dimensión once dentro de la celda positiva de
dimensión veintidós. El punto está en la celda superior; la familia conserva
su dimensión propia. Los trece puntos y las doce separaciones tienen así
papeles matemáticos distintos.

Permitir \(d_r\ge0\) conserva la no negatividad. Cuando \(d_r=0\), las
columnas consecutivas \(r-1\) y \(r\) coinciden y forman elementos
paralelos del matroide. El rango sigue siendo dos porque \(t_0=0\) y
\(t_L=1\). Las bases son exactamente los pares \(i<j\) que atraviesan
alguna separación positiva. Esta degeneración produce estratos de incidencia
consecutiva. Una columna nula, que daría un lazo, corresponde a una operación
distinta de la coincidencia de dos columnas no nulas.

\subsection{Semilla de triangulación y relación de intercambio}

Sea \(N=L+1\) y considérense los vértices \(0,\ldots,L\) de un polígono
convexo en su orden cíclico. A cada lado o diagonal \((i,j)\) se asigna
la coordenada \(A_{ij}=\Delta_{ij}(C(d))\), con índices ordenados.
La triangulación en abanico desde el vértice cero tiene diagonales
\((0,j)\), \(2\le j\le L-1\). Sus lados son variables congeladas y
sus diagonales variables mutables. El quiver se obtiene orientando
cíclicamente los tres lados de cada triángulo y cancelando los pares de
flechas opuestas. Esta es una semilla de tipo \(A_{N-3}\); para el registro
dodecafásico, su tipo es \(A_{10}\).

La semilla tiene una realización plábica explícita. Se coloca un vértice
negro dentro de cada triángulo; en los vértices del polígono se coloca un
vértice blanco si incide en alguna diagonal, y negro en caso contrario.
Se unen los centros a las tres esquinas de sus triángulos, se retiran los
lados originales y se contraen las aristas entre vértices del mismo color.
Se añade una hoja de frontera a cada vértice marcado. Esta construcción
recoge el algoritmo de triangulación de Kodama y Williams,
\href{https://people.math.harvard.edu/~williams/papers/KW-ArxivVersion.pdf}
{\emph{KP solitons and total positivity for the Grassmannian}, algoritmo 12.1}.
Las etiquetas de coordenadas de esta semilla se evalúan en los menores ya
producidos por la red de dos filas. Para una cara mutable, la coordenada
de tipo \(\mathcal X\) es el cociente de los productos de las
coordenadas \(\mathcal A\) incidentes que prescribe su columna del
quiver; en un cuadrilátero se escribirá explícitamente como una razón
doble. Cambiar el grafo cambia la carta de esos menores, manteniendo su
procedencia.

\begin{theorem}[Intercambio positivo de la realización del registro]
\label{thm:nc-ptolemy}
Para cuatro vértices \(a<b<c<d\),
\begin{equation}
 A_{ac}A_{bd}=A_{ab}A_{cd}+A_{ad}A_{bc}.
 \label{eq:nc-ptolemy}
\end{equation}
En consecuencia, el cambio de diagonal \((a,c)\leftrightarrow(b,d)\)
se evalúa mediante
\[
 A_{bd}=\frac{A_{ab}A_{cd}+A_{ad}A_{bc}}{A_{ac}}>0.
\]
Todas las triangulaciones proporcionan cartas positivas compatibles sobre
la misma realización \(C(d)\).
\end{theorem}
\begin{proof}
Al sustituir \(A_{ij}=t_j-t_i\), la igualdad se reduce a
\[
 (t_c-t_a)(t_d-t_b)
 =(t_b-t_a)(t_d-t_c)+(t_d-t_a)(t_c-t_b).
\]
La expansión de ambos miembros prueba la identidad. Las diferencias son
positivas; por tanto, cada flip conserva una evaluación positiva y el
cociente está definido. La conexión de las triangulaciones por flips
permite pasar entre ellas mediante estos intercambios. La interpretación
cluster de la relación de Plücker es el reconocimiento posterior desarrollado
por J. Scott en
\href{https://arxiv.org/abs/math/0311148}{\emph{Grassmannians and Cluster
Algebras}}. La prueba algebraica de la evaluación concreta ya está contenida
en la identidad mostrada.
\end{proof}

\subsection{Refinamiento de las separaciones y naturalidad de los intercambios}

Sea \(B\ge2\). En profundidad \(n\), cada separación \(d_r\) se divide
en \(B^n\) separaciones consecutivas iguales:
\[
 d^{(n)}_{(r-1)B^n+c}=d_r/B^n,
 \qquad 1\le c\le B^n.
\]
Esto define \(LB^n+1\) puntos y una matriz \(C_n\). Si \(\iota_n(j)=Bj\)
identifica un punto con su descendiente de igual posición, entonces
\begin{equation}
 C_{n+1}|_{\iota_n(\{0,\ldots,LB^n\})}=C_n,
 \qquad
 \Delta_{\iota_n(i),\iota_n(j)}(C_{n+1})=\Delta_{ij}(C_n).
 \label{eq:nc-refinement}
\end{equation}
La red refinada reemplaza cada arista vertical de peso \(d\) por una
secuencia de \(B\) posiciones de descenso, con pesos \(d/B\). La suma
de caminos entre extremos heredados sigue siendo \(d\). La composición de
refinamientos conserva esta igualdad porque las sumas finitas se agrupan
asociativamente.

\begin{proposition}[Naturalidad sobre extremos heredados]
\label{prop:nc-inherited}
Todo intercambio \eqref{eq:nc-ptolemy} entre cuatro extremos heredados
conmuta exactamente con el refinamiento. Toda triangulación antigua se
extiende a una del polígono refinado preservando sus triángulos interiores:
los polígonos añadidos entre dos extremos consecutivos se triangulan por
fuera de esos triángulos. Los flips entre diagonales antiguas conservan
entonces sus relaciones de intercambio.
\end{proposition}
\begin{proof}
La primera conclusión resulta al sustituir cada menor por su igualdad en
\eqref{eq:nc-refinement}. Para la segunda, se conservan las cuerdas que
unen extremos antiguos consecutivos y se triangula cada bolsillo exterior
formado por las nuevas marcas. Estas regiones tienen interiores disjuntos.
Los dos triángulos incidentes en una diagonal antigua permanecen iguales,
y, por tanto, también su cuadrilátero y su identidad de Ptolomeo.
\end{proof}

Esta naturalidad especifica los extremos sobre los que actúa. Un lado
antiguo, al recibir vértices nuevos en su arco de frontera, pasa a ser una
diagonal del polígono mayor. Por ello una variable antiguamente congelada
puede ser mutable en el problema ampliado. La compatibilidad anterior es
una compatibilidad exacta del sistema de evaluaciones y de los flips
heredados. Una inclusión de patrones de semillas con igual régimen de
coeficientes se obtiene manteniendo congeladas esas cuerdas y las nuevas
diagonales auxiliares. Esto declara la operación que conserva la semilla;
el incremento de profundidad y una mutación siguen siendo operaciones
distintas.

\subsection{Acción diagonal proyectiva e invariancia de las razones cruzadas}

Si \(\lambda_j>0\), la acción por columnas
\(C\mapsto C\operatorname{diag}(\lambda_0,\ldots,\lambda_L)\) da
\[
 A_{ij}\longmapsto\lambda_i\lambda_j A_{ij}.
\]
Preserva la positividad y todos los estratos definidos por anulación de
menores. La razón doble asociada a un cuadrilátero es
\begin{equation}
 X_{abcd}=\frac{A_{ab}A_{cd}}{A_{ad}A_{bc}}>0.
 \label{eq:nc-crossratio}
\end{equation}
Los cuatro factores \(\lambda_a\lambda_b\lambda_c\lambda_d\) se
cancelan exactamente. En consecuencia, toda acción diagonal positiva por
columnas deja fijo \(X_{abcd}\), incluso si depende de un parámetro y
tiene determinante uno. La modificación del peso de las columnas conserva
el punto proyectivo de cada columna; sus razones dobles expresan esa
conservación. Para que el registro cambie la forma proyectiva debe actuar
sobre las separaciones ordenadas.

\subsection{Deformación positiva de las separaciones dirigida por el registro}

El corpus fija la representación de las doce posiciones y su proyector
ortogonal \(P_3\) sobre el sector tridimensional. Con
\(P_{11}=I_{12}-J_{12}/12\), la dirección recuperada es
\[
 u=P_3P_{11}K=P_3K\in\mathbf1^\perp.
\]
En la carta declarada, poniendo \(r=\sqrt5\), su evaluación exacta es
\begin{align*}
u=(&-495/4-233r/10,\ 403/4+169r/10,\ -403/4-169r/10,\\
   &495/4+233r/10,\ 601/4+1031r/10,\ 203/4+211r/5,\\
   &-203/4-211r/5,\ -601/4-1031r/10,\ 313/4+569r/10,\\
   &162+219r/20,\ -162-219r/20,\ -313/4-569r/10).
\end{align*}
Se conservan \(\sum_ru_r=0\) y
\(\|u\|^2=(6638585+2275584\sqrt5)/20>0\).

Para \(s\in\mathbb R\), definimos la deformación de los pesos mediante
\begin{equation}
 d_r(s)=\frac{K_re^{su_r}}{\sum_{\ell=1}^{12}K_\ell e^{su_\ell}},
 \qquad t_j(s)=\sum_{r=1}^jd_r(s),\qquad C(s)=C(d(s)).
 \label{eq:nc-gap-flow}
\end{equation}
La fórmula utiliza el registro y su dirección ya generados. Es una nueva
realización analítica del lector; el parámetro \(s\) mide su deformación.
Cuando se usa la salida interna \(\rho_A=\pi_{\rm HMT}/\varphi_{\rm HMT}^2\),
la escritura \(s=\tau\log\rho_A\) expresa exactamente la misma familia.
Esa reparametrización actúa después de la generación de las constantes.
Una identificación de \(s\) o \(\tau\) con un tiempo físico requiere la
ley temporal del lector correspondiente.

\begin{theorem}[Cambio de forma, positividad y compatibilidad con la profundidad]
\label{thm:nc-shape-flow}
La familia \eqref{eq:nc-gap-flow} pertenece a
\(\operatorname{Gr}_{>0}(2,13)\) para todo \(s\) real y conserva
\(t_0=0\), \(t_{12}=1\). Para \(i<j\), sean
\[
 \mu_{ij}(s)=
 \frac{\sum_{r=i+1}^jK_ru_re^{su_r}}{\sum_{r=i+1}^jK_re^{su_r}},
 \qquad
 \mu(s)=\frac{\sum_rK_ru_re^{su_r}}{\sum_rK_re^{su_r}}.
\]
Entonces
\begin{equation}
 \frac{d}{ds}\log A_{ij}(s)=\mu_{ij}(s)-\mu(s),
 \quad
 \frac{d}{ds}\log X_{abcd}(s)
 =\mu_{ab}(s)+\mu_{cd}(s)-\mu_{ad}(s)-\mu_{bc}(s).
 \label{eq:nc-crossratio-derivative}
\end{equation}
La deformación conmuta con cada refinamiento que asigna a los hijos de la
separación \(r\) el peso \(d_r/B\) y la misma velocidad \(u_r\).
\end{theorem}
\begin{proof}
Todos los exponenciales y todos los \(K_r\) son positivos. La normalización
hace que la suma de separaciones sea uno; los menores son sumas de
separaciones positivas por \eqref{eq:nc-minor-gap}. Derivar el logaritmo de
\[
 A_{ij}(s)=\frac{\sum_{r=i+1}^jK_re^{su_r}}
                      {\sum_rK_re^{su_r}}
\]
da la primera identidad. La combinación de cuatro logaritmos en
\eqref{eq:nc-crossratio} cancela las cuatro medias globales con sus signos
y da la segunda. Finalmente, reemplazar \(K_r\) por \(B\) copias de
\(K_r/B\) con idéntica \(u_r\) conserva el numerador y el denominador
cuando se agrupan los hijos. La igualdad vale para todos los parámetros,
para cualquier profundidad y para toda composición de refinamientos.
\end{proof}

Hay un testigo exacto de que esta familia cambia la forma. Para los primeros
cuatro extremos,
\begin{equation}
 X_{0123}(0)=\frac{1560}{23711},\qquad
 \left.\frac{d}{ds}\log X_{0123}(s)\right|_{s=0}
 =-\frac{309899}{917}-\frac{268596}{4585}\sqrt5<0.
 \label{eq:nc-nontrivial-witness}
\end{equation}
La cuenta utiliza únicamente las primeras tres separaciones
\((234,543,140)\) y sus velocidades ya recuperadas. La desigualdad
estricta prueba un movimiento proyectivo efectivo. Las identidades de
Ptolomeo siguen siendo exactas durante ese movimiento, pues se aplican a
las columnas \((1,t_j(s))\) en cada parámetro. El registro controla así
una deformación positiva de la forma y, simultáneamente, un sistema de
cartas que permanece compatible con ella.

El alcance demostrado reúne una red planar de rango dos, su medición
explícita, una semilla plábica de triangulación, todas sus mutaciones por
flips, la naturalidad de los extremos heredados y una deformación de
razones dobles dirigida por \(K\). Cada operación tiene un dominio y un
resultado verificables. La comparación de esta realización con la
dinámica temporal enriquecida y con la rama material utiliza sus respectivos
mapas de lectura; ninguna identidad de coordenadas sustituye esos mapas.


% SOURCE appendices/polygon_amplituhedron.tex
\section{Geometrías amplituhedrales de rango uno y compatibilidad de residuos}
\label{sec:k-poligono-amplituhedral}

El registro dodecafásico permite reunir explícitamente datos positivos,
dominio grassmanniano, imagen cinemática, cobertura, grado sobre las celdas
y forma canónica. El primer resultado concierne al amplituedro
bidimensional de parámetros \((n,k,m)=(13,1,2)\); después se prolonga la
construcción a los polítopos cíclicos de rango uno y dimensión superior.
Se trata de una realización geométrica posterior del registro producido
por APP--TRIT--TPK. La operación adicional declarada es la elevación por
potencias de una partición racional ordenada. Esta elección fija una
realización particular; las correspondencias físicas posteriores conservan
sus lectores y condiciones propios.

La procedencia del registro conserva la composición
\[
 x_{\mathrm{term}}\longmapsto\mathcal R_{12}(x_{\mathrm{term}})
 \longmapsto(b^{90},b^{120},Q_{\mathrm{TPK}})
 \longmapsto U_{\mathrm{sgn}}\longmapsto K.
\]
APP conserva ambas hojas y sus residuos y cocientes; TRIT orienta el
régimen de cada transición; TPK transporta y actualiza el estado con
acarreo, frontera, ruta y memoria. La publicación \(K\) se lee en el mismo
sistema de estados enriquecidos que genera conjuntamente el continuo.
Su extracción firmada y su inversión
de Hadamard producen
\begin{equation}
 K=(234,543,140,729,659,824,621,58,914,794,146,601),
 \qquad S_K=\sum_{j=1}^{12}K_j=6263.
 \label{eq:polygon-K}
\end{equation}
La prueba geométrica utiliza la positividad de esas doce coordenadas, su
orden y su total. Sus coeficientes proceden de aquella extracción y no de
la selección retrospectiva de un valor físico.

\subsection{Partición racional y elevación positiva}

Definimos trece puntos racionales de partición:
\begin{equation}
 t_0=0,\qquad t_i=\frac{K_1+\cdots+K_i}{S_K}\quad(1\leq i\leq12),
 \qquad Z_i=(1,t_i,t_i^2)^\mathsf T,\qquad Z=(Z_0\ \cdots\ Z_{12}).
 \label{eq:polygon-Z}
\end{equation}
Sus numeradores, con denominador común \(6263\), son
\[
 (0,234,777,917,1646,2305,3129,3750,3808,4722,5516,5662,6263).
\]
Así, \(0=t_0<t_1<\cdots<t_{12}=1\). La partición etiquetada y el total
conservan el registro mediante la inversa exacta
\begin{equation}
 K_i=S_K(t_i-t_{i-1}).
 \label{eq:polygon-recover-K}
\end{equation}
La partición conserva las proporciones; retener \(S_K\) conserva también
la escala entera.

\begin{proposition}[Datos cinemáticos estrictamente positivos]
\label{prop:polygon-vandermonde}
La matriz \(Z\) tiene rango tres y todos sus menores ordenados de tamaño
tres son positivos. Para \(0\leq a<b<c\leq12\),
\begin{equation}
 \langle abc\rangle:=\det(Z_a,Z_b,Z_c)
 =(t_b-t_a)(t_c-t_a)(t_c-t_b)>0.
 \label{eq:polygon-minors}
\end{equation}
\end{proposition}
\begin{proof}
Se resta la primera columna de las dos últimas y se desarrolla por la
primera fila. Cada diferencia \(t_j^2-t_a^2\) factoriza como
\((t_j-t_a)(t_j+t_a)\), lo que da el producto indicado. Las sumas parciales
estrictamente crecientes prueban su positividad. Cualquiera de esos
menores acredita el rango tres.
\end{proof}

\subsection{Imagen, fibras y celdas de grado uno}

El Grassmanniano \(\operatorname{Gr}_{\geq0}(1,13)\) es el espacio
proyectivo de vectores \(c=(c_0,\ldots,c_{12})\) no negativos y no nulos,
módulo escala positiva. La carta \(\sum c_j=1\) lo identifica con el
símplex cerrado de dimensión doce. Definimos
\begin{equation}
 \Phi_Z:\operatorname{Gr}_{\geq0}(1,13)\longrightarrow\mathbb P^2,
 \qquad[c]\longmapsto\left[\sum_{j=0}^{12}c_jZ_j\right].
 \label{eq:polygon-map}
\end{equation}
En la carta \(Y_0=1\), sean \(v_j=(t_j,t_j^2)\) y
\(P_K=\operatorname{conv}(v_0,\ldots,v_{12})\).

\begin{theorem}[Amplituedro poligonal del registro]
\label{thm:polygon-amplituhedron}
La imagen cerrada de \(\Phi_Z\) es exactamente el polígono convexo
\(P_K=\mathcal A_{13,1,2}(Z)\), de dimensión dos y trece vértices. La imagen
del Grassmanniano estrictamente positivo es el interior de \(P_K\).
Las once celdas
\begin{equation}
 \mathcal S_i=\{[c]:c_0,c_i,c_{i+1}>0,\
 c_j=0\text{ si }j\notin\{0,i,i+1\}\},\qquad 1\leq i\leq11,
 \label{eq:polygon-cells}
\end{equation}
se aplican biyectivamente, con grado uno y orientación positiva, sobre los
interiores de \(T_i=[v_0,v_i,v_{i+1}]\). Sus clausuras cubren \(P_K\) y sus
interiores son disjuntos dos a dos.
\end{theorem}
\begin{proof}
En la normalización \(\sum c_j=1\), la aplicación es
\(c\mapsto\sum c_jv_j\), cuya imagen es la envolvente convexa. Los menores
de \eqref{eq:polygon-minors} muestran que cada triple ordenado tiene
orientación positiva. Cada lado consecutivo deja todos los otros vértices
estrictamente a su izquierda; también lo hace el lado final de \(v_{12}\)
a \(v_0\). Todos los puntos son, por tanto, vértices extremos de un polígono
convexo.

Toda combinación con \(c_j>0\) es interior, pues en cualquier recta de
soporte alguno de los vértices está estrictamente en el semiplano
interior. Recíprocamente, sean \(y\) interior y
\(b=13^{-1}\sum v_j\). Para \(\varepsilon>0\) suficientemente pequeño,
\(y'=(y-\varepsilon b)/(1-\varepsilon)\) permanece en el polígono.
Escribiendo \(y'=\sum a_jv_j\), \(a_j\geq0\), \(\sum a_j=1\), se obtiene
\(y=\sum((1-\varepsilon)a_j+\varepsilon/13)v_j\), con todos los
coeficientes positivos.

Las diagonales desde \(v_0\) dividen el polígono en los once triángulos.
Sobre cada soporte triple, las coordenadas baricéntricas son únicas,
porque \(\langle0,i,i+1\rangle>0\); su inversa explícita aparece en
\eqref{eq:polygon-barycentric}. Cada soporte define una celda positroidal
de rango uno con coordenadas \([1:a:b]\), \(a,b>0\), en esas posiciones.
Una red dirigida planar con una fuente y tres trayectorias de pesos
\(1,a,b\) hacia los tres destinos realiza esas coordenadas por medición
de frontera. El determinante positivo prueba orientación y grado uno.
\end{proof}

La fibra genérica de la aplicación completa tiene dimensión diez:
\(\sum c_j=1\), \(\sum c_jt_j=x\) y \(\sum c_jt_j^2=y\) son tres
ecuaciones independientes en trece coordenadas. Para \(y\) interior la
fibra contiene un punto estrictamente positivo, por lo que su intersección
con el símplex tiene dimensión \(13-3=10\). El grado uno concierne a las
celdas triangulares; la forma bidimensional se obtiene mediante esas
celdas.

\subsection{Forma canónica y residuos orientados}

Para \(Y=(1,x,y)^\mathsf T\), definimos
\begin{equation}
 \ell_{ab}(Y)=\det(Z_a,Z_b,Y)
 =(t_b-t_a)\{y-(t_a+t_b)x+t_at_b\}.
 \label{eq:polygon-lines}
\end{equation}
La expresión vale también cuando \(b<a\). Para \(a<b<c\), sean
\(D_{abc}=\langle abc\rangle\) y
\begin{equation}
 \lambda_a=\frac{\ell_{bc}}{D_{abc}},\qquad
 \lambda_b=\frac{\ell_{ca}}{D_{abc}},\qquad
 \lambda_c=\frac{\ell_{ab}}{D_{abc}},\qquad
 \lambda_a+\lambda_b+\lambda_c=1.
 \label{eq:polygon-barycentric}
\end{equation}
La orientación afín es \(dx\wedge dy\). Se fija la convención de frontera
\(\operatorname{Res}^{\partial}(\eta\wedge du/u)=\eta|_{u=0}\):
la diferencial normal ocupa la última posición. En dimensión dos, el
residuo de Poincaré que coloca \(du/u\) primero tiene el signo opuesto.
Las cancelaciones siguientes usan de forma uniforme la convención declarada.

\begin{theorem}[Forma racional y cancelación de las diagonales]
\label{thm:polygon-form}
La forma canónica del triángulo orientado \([v_a,v_b,v_c]\) es
\begin{equation}
 \Omega_{abc}=
 \frac{D_{abc}^2\,dx\wedge dy}
 {\ell_{ab}(Y)\ell_{bc}(Y)\ell_{ca}(Y)}.
 \label{eq:polygon-triangle-form}
\end{equation}
La forma canónica del polígono es
\begin{equation}
 \boxed{\displaystyle
 \Omega_{P_K}(x,y)=
 \sum_{i=1}^{11}
 \frac{t_it_{i+1}(t_{i+1}-t_i)\,dx\wedge dy}
 {(y-t_ix)\{y-(t_i+t_{i+1})x+t_it_{i+1}\}(t_{i+1}x-y)}.}
 \label{eq:polygon-full-form}
\end{equation}
Los polos de las diagonales interiores se cancelan. Los únicos divisores
polares son las trece rectas de los lados exteriores. Parametrizado cada
lado desde su vértice inicial al final por \(s\in[0,1]\), su residuo de
frontera es \(ds/[s(1-s)]\).
\end{theorem}
\begin{proof}
La forma del símplex baricéntrico es
\(d\lambda_b\wedge d\lambda_c/(\lambda_a\lambda_b\lambda_c)\).
El jacobiano satisface
\(dx\wedge dy=D_{abc}\,d\lambda_b\wedge d\lambda_c\).
Sustituir \eqref{eq:polygon-barycentric} da
\eqref{eq:polygon-triangle-form}. Para \(a=0\), cancelar los tres factores
escalares de las rectas da cada sumando de \eqref{eq:polygon-full-form}.

En el lado \(ab\), pongamos \(\lambda_c=u\),
\(\lambda_b=(1-u)s\), \(\lambda_a=(1-u)(1-s)\). Entonces
\[
 \Omega_{abc}=\frac{ds}{s(1-s)}\wedge\frac{du}{u(1-u)}.
\]
El residuo es la forma del intervalo orientado. Dos triángulos contiguos
inducen orientaciones contrarias sobre su diagonal y producen residuos
opuestos. Al ser simples todos los polos, anular el residuo elimina el
polo de la diagonal como divisor. Cada lado exterior aparece una vez.
La fórmula homogénea es de grado cero como forma proyectiva y no tiene
otros divisores polares. Si dos formas racionales proyectivas tuvieran los
mismos residuos y ningún otro polo, su diferencia sería una dos-forma
holomorfa sobre \(\mathbb P^2\), que es necesariamente cero. Así quedan
probadas caracterización y unicidad.
\end{proof}

Se emplea la noción convencional de geometría positiva mediante polos
logarítmicos y residuos canónicos de frontera, desarrollada por
N.~Arkani-Hamed, Y.~Bai y T.~Lam, \emph{Positive Geometries and Canonical
Forms}, JHEP 11 (2017), 039,
\url{https://arxiv.org/abs/1703.04541}, secciones 2, 3, 5.2, 6.1 y 7.2.
La selección de los coeficientes del presente polígono precede a ese
reconocimiento y procede de \(K\).

\subsection{Refinamiento por acarreo sobre la geometría fija}

En un triángulo \(T=[A,B,C]\) de orientación positiva, escribimos
\begin{equation}
 Y(u,r)=(1-u)A+u\{(1-r)B+rC\},\qquad 0<u,r<1.
 \label{eq:polygon-blowdown}
\end{equation}
El vértice \(A\) corresponde a \(u=0\) y el lado opuesto a \(u=1\).
La forma queda
\begin{equation}
 \Omega_T=\frac{du}{u(1-u)}\wedge\frac{dr}{r(1-r)}.
 \label{eq:polygon-product-form}
\end{equation}
Sea \(B_0\geq2\) una base entera declarada, distinta del vértice \(B\).
La coordenada \(r\) admite la actualización por residuo y acarreo
\begin{equation}
 c=\lfloor B_0r\rfloor,\qquad r'=B_0r-c,\qquad
 r=\frac{c+r'}{B_0}.
 \label{eq:polygon-carry}
\end{equation}
Los extremos conservan sus marcas de frontera. El dígito \(c\) permanece
como memoria y permite invertir la actualización. Conviene precisar las
cartas de los dos transportes afines. El zoom de la subcelda actúa sobre
\(r\in[c/B_0,(c+1)/B_0]\) mediante \(r'=B_0r-c\).
En cambio, el transporte de la pareja levantada
\(L_c(x,y)=(B_0x-c,B_0y+c)\), con \(x+y=M\), cambia la masa total
a \(B_0M\) y publica la coordenada
\[
 \frac{B_0x-c}{B_0M}=\frac{x}{M}-\frac{c}{B_0M}.
\]
Son aplicaciones de dominios y normalizaciones diferentes. Ambas conservan
el acarreo en el estado enriquecido y admiten recuperación con sus datos
de escala. La demostración siguiente utiliza el zoom de subcelda y la
partición geométrica que éste identifica.

\begin{theorem}[Invariancia exacta bajo subdivisión con memoria]
\label{thm:polygon-carry}
Sean \(s_j=j/B_0\) y \(W_j=(1-s_j)B+s_jC\).
Los triángulos \(T_j=[A,W_j,W_{j+1}]\), \(0\leq j<B_0\), subdividen
\(T\) y satisfacen
\begin{equation}
 \sum_{j=0}^{B_0-1}\Omega_{T_j}=\Omega_T.
 \label{eq:polygon-carry-additivity}
\end{equation}
La igualdad permanece válida a toda profundidad finita y para cualquier
partición racional ordenada del lado. Subdividir los once triángulos
de \(P_K\) deja \(\Omega_{P_K}\) exactamente invariante.
\end{theorem}
\begin{proof}
Las regiones \(r\in[s_j,s_{j+1}]\) cubren \(T\) con interiores disjuntos.
En las coordenadas comunes,
\[
 \Omega_{T_j}=\frac{du}{u(1-u)}\wedge
 \frac{(s_{j+1}-s_j)\,dr}{(r-s_j)(s_{j+1}-r)}.
\]
La identidad racional
\[
 \frac{s_{j+1}-s_j}{(r-s_j)(s_{j+1}-r)}
 =\frac1{r-s_j}+\frac1{s_{j+1}-r}
\]
telescopa: cada término interior cancela el del subintervalo contiguo.
Quedan \(1/r+1/(1-r)=1/[r(1-r)]\). Esto prueba igualdad de formas
racionales; iterarla proporciona toda profundidad finita.
La prueba sólo utiliza \(s_0=0<s_1<\cdots<s_N=1\), de modo que cubre
particiones no uniformes. El cambio
\(r'=(r-s_j)/(s_{j+1}-s_j)\) devuelve localmente la forma unitaria.
La palabra de acarreos identifica la subcelda y su inversa recupera la
coordenada anterior. Los nuevos polos interiores se cancelan mediante
la misma identidad telescópica.
\end{proof}

La suma anterior reúne la partición geométrica completa. Cuando un lector
de supervivencia selecciona sólo algunas historias, su imagen es la unión
de las subceldas seleccionadas y su forma es la suma correspondiente;
coincide con la forma del polígono completo precisamente cuando esa unión
lo cubre con multiplicidad orientada uno. Por tanto, el criterio de
cobertura para una selección concreta de historias queda expresado por
una condición comprobable sobre su cadena de celdas.

Este refinamiento conserva el polígono fijo. Insertar un nuevo punto
\((t,t^2)\) de la parábola entre dos parámetros consecutivos \(a<t<b\)
produce otra geometría: la altura de la cuerda menos \(t^2\) es
\((t-a)(b-t)>0\), y el nuevo punto queda fuera del polígono previo.
La invariancia demostrada corresponde a la subdivisión rectilínea de
celdas, cuya frontera exterior permanece fija.

\subsection{Cambios de triangulación y reparametrización positiva}

Para cuatro vértices ordenados \(A,B,C,D\), se cumple
\begin{equation}
 \Omega_{ABC}+\Omega_{ACD}=\Omega_{ABD}+\Omega_{BCD}.
 \label{eq:polygon-flip}
\end{equation}
Los dos miembros tienen los mismos residuos exteriores y cancelan su
diagonal interior; su igualdad procede de la unicidad demostrada.
También se comprueba multiplicando por sus denominadores lineales.
Todo cambio entre triangulaciones de un polígono convexo se descompone
en estos cambios de diagonal, por lo que la forma es independiente
de la triangulación. Una identificación con coordenadas cluster debe
conservar adicionalmente la variedad y sus relaciones de cambio.

Si \(D=\operatorname{diag}(d_0,\ldots,d_{12})\), \(d_j>0\), entonces
\begin{equation}
 \mathcal A_{13,1,2}(ZD)=\mathcal A_{13,1,2}(Z).
 \label{eq:polygon-rescale}
\end{equation}
El automorfismo \([c]\mapsto[Dc]\) reparametriza la aplicación cinemática.
En la carta afín, \(c_j\) cambia a
\(c_jd_j/\sum_\ell c_\ell d_\ell\), que recorre nuevamente el símplex.
La imagen y su forma canónica permanecen iguales. El desplazamiento
de una historia parametrizada y el cambio del conjunto geométrico
completo son así operaciones distinguibles.

\subsection{Extensión de rango uno a dimensiones superiores}

Para \(2\leq m\leq12\), definimos
\begin{equation}
 Z_i^{(m)}=(1,t_i,t_i^2,\ldots,t_i^m)^\mathsf T,\qquad
 P_K^{(m)}=\operatorname{conv}
 \{(t_i,t_i^2,\ldots,t_i^m):0\leq i\leq12\}.
 \label{eq:polygon-higher-lift}
\end{equation}
Los menores máximos ordenados son los productos de Vandermonde
\(\prod_{a<b}(t_{i_b}-t_{i_a})>0\). La matriz tiene rango \(m+1\).
La imagen de \(\operatorname{Gr}_{\geq0}(1,13)\) es el polítopo cíclico
\(P_K^{(m)}=\mathcal A_{13,1,m}(Z^{(m)})\), de dimensión \(m\); su fibra
genérica tiene dimensión \(12-m\). Aquí se usa la extensión matemática
de la definición amplituhedral a cualquier \(m\); las aplicaciones
habituales a amplitudes imponen sus propios valores de \(m\).

La triangulación se determina por el orden fijo de etiquetas. Se elige
el menor vértice de un polítopo, se triangulan recursivamente todas sus
facetas que no lo contienen y se toman conos desde ese vértice. La
inducción comienza con los puntos. Una recta desde el vértice inicial
hasta un punto interior corta la frontera opuesta en una faceta; ello
prueba la cobertura. La unicidad del punto de intersección y la
compatibilidad de las restricciones inducidas por el mismo orden prueban
que los interiores son disjuntos y que las intersecciones son caras
comunes. En cada paso disminuye la dimensión, por lo que el procedimiento
termina. Esta es una triangulación por tracción determinada por las
etiquetas.

Para un símplex orientado
\(T=[p_0,\ldots,p_m]\) de esa triangulación, sean
\[
 A_T=(p_1-p_0\ \cdots\ p_m-p_0),\qquad
 (\lambda_1,\ldots,\lambda_m)^\mathsf T=A_T^{-1}(x-p_0),
 \qquad\lambda_0=1-\sum_{j=1}^m\lambda_j.
\]
La orientación de sus vértices se toma de modo que \(\det A_T>0\).
Su forma canónica es
\begin{equation}
 \Omega_T(x)=
 \frac{dx_1\wedge\cdots\wedge dx_m}
 {\det(A_T)\prod_{j=0}^m\lambda_j(x)},\qquad
 \Omega_{P_K^{(m)}}=\sum_{T}\Omega_T.
 \label{eq:polygon-higher-form}
\end{equation}
El cambio baricéntrico prueba la fórmula y que la restricción cinemática
sobre el soporte de \(T\) tiene grado uno. Los residuos son las formas
de las caras con los signos inducidos por la orientación. En una cara
interior, las dos orientaciones son contrarias y el residuo total se
anula. En una cara exterior queda la suma de las formas de la
triangulación de esa cara. La inducción sobre la dimensión demuestra
que dicha suma es su forma canónica. Repetir el residuo a lo largo de
una bandera de caras alcanza los vértices, donde la normalización es
\(\pm1\). Se obtiene así la correspondencia de residuos en todas las
codimensiones para esta familia de rango uno.

La misma demostración cubre cualquier subdivisión simplicial compatible
del polítopo fijo: sus caras interiores se cancelan y las exteriores
conservan los residuos. La inexistencia de formas holomorfas de grado
máximo sobre \(\mathbb P^m\) asegura la unicidad de la forma racional.
En particular, la estabilidad bajo refinamiento, la cobertura y los
residuos quedan cerrados en esta familia de dimensiones \(2,\ldots,12\).
Los grados grassmannianos \(k>1\) conservan sus aplicaciones y sus
obligaciones demostrativas específicas.


% SOURCE sections/higher_rank.tex
\section{Matrices de momentos y realizaciones amplituhedrales de rango superior}
\label{sec:higher-rank}
La partición del registro dodecafásico permite prolongar la construcción
de rango uno a puntos explícitos de rango superior. Los mismos trece
nodos ordenados proporcionan simultáneamente menores de Plücker positivos
y una forma de momentos definida positiva. Estas dos positividades se
obtienen de una sola lista de intervalos; la segunda prueba que la
composición conserva el rango requerido.

Sea \(d_j=K_j/Q\), donde \(Q=\sum_{j=1}^{12}K_j=6263\), y definamos
\[
 t_0=0,\qquad t_j=\sum_{r=1}^jd_r.
\]
La positividad de los doce bloques da
\(0=t_0<t_1<\cdots<t_{12}=1\). Para cada \(2\le k\le9\), ponemos
\[
 C^{(k)}_{ai}=t_i^a\quad(0\le a<k),\qquad
 Z^{(k)}_{bi}=t_i^b\quad(0\le b<k+4),\qquad 0\le i\le12.
\]
La convención \(t_i^0=1\) incluye el nodo cero. Así, \(C^{(k)}\)
es una matriz de tamaño \(k\times13\), mientras \(Z^{(k)}\) tiene
tamaño \((k+4)\times13\).

\begin{theorem}[Positividad y rango de la realización de momentos]
La matriz \(C^{(k)}\) representa un punto de
\(\Gr_{>0}(k,13)\), y todos los menores máximos ordenados de
\(Z^{(k)}\) son positivos. La matriz
\[
 Y^{(k)}=C^{(k)}(Z^{(k)})^{\mathsf T},\qquad
 Y^{(k)}_{ab}=\sum_{i=0}^{12}t_i^{a+b},
\]
tiene rango \(k\). En particular, para el dato externo positivo
\(Z^{(k)}\) ya construido a partir del registro,
\[
 [Y^{(k)}]\in\mathcal A_{13,k,4}(Z^{(k)}),\qquad 2\le k\le9,
\]
donde
\[
 \mathcal A_{13,k,4}(Z^{(k)})
 =\{[C(Z^{(k)})^{\mathsf T}]:[C]\in\Gr_{\ge0}(k,13)\}.
\]
\end{theorem}
\begin{proof}
Un menor máximo de cualquiera de las dos matrices iniciales tiene la forma
\[
 \det(t_{i_j}^{a})_{a=0}^{r-1}
 =\prod_{p<q}(t_{i_q}-t_{i_p})>0,
\]
para índices crecientes. La desigualdad prueba las dos afirmaciones sobre
Plücker. El bloque inicial \(k\times k\) de \(Y^{(k)}\) es la matriz
de Hankel \(H^{(k)}\). Para \(x\ne0\),
\[
 x^{\mathsf T}H^{(k)}x
 =\sum_{i=0}^{12}\left(\sum_{a=0}^{k-1}x_at_i^a\right)^2>0.
\]
En efecto, un polinomio no nulo de grado inferior a \(k\le9\) no se
anula en trece puntos distintos. Luego \(H^{(k)}\) es definida positiva,
su determinante es estrictamente positivo y \(Y^{(k)}\) tiene rango
\(k\). La pertenencia amplituhedral es entonces la composición de las
dos matrices positivas anteriores, con el rango ya comprobado.
\end{proof}

En rango dos, la prueba adopta la forma particularmente explícita
\[
 \det H^{(2)}=13\sum_i t_i^2-\left(\sum_i t_i\right)^2
 =\sum_{i<j}(t_j-t_i)^2>0.
\]
La separación de los nodos, producida por la positividad del registro,
queda así convertida en coercividad de una forma cuadrática. A cada
registro fijo corresponde un punto determinado para cada rango. La
cobertura completa, las triangulaciones y los residuos demostrados en la
sección anterior mantienen su dominio de rango uno; el presente teorema
añade la familia explícita de puntos de rangos dos a nueve y sus matrices
de momentos. La distinción permite reutilizar cada demostración con su
alcance exacto.


% SOURCE sections/refinamiento_dimensional_momentos.tex
% Sección aditiva. No modifica la edición de 211 páginas.
% Procedencia y comprobaciones: documentos internos de este mismo directorio.
\section{Refinamiento nonádico, banderas positivas y compatibilidad de los momentos}
\label{sec:dim-cofinal}

La longitud de una publicación finita determina cuántos menores pueden
formarse con sus columnas. El refinamiento conserva las columnas heredadas
y añade otras con procedencia conocida. Por ello, la dimensión de una
realización matricial puede crecer sin sustituir el registro que fija sus
separaciones. La construcción siguiente reúne este crecimiento con dos
nociones precisas de compatibilidad: restricción de evaluaciones a las
marcas anteriores y promedio de observables sobre las celdas descendientes.
La primera preserva las banderas polinomiales; la segunda conserva una
medida y separa el incremento de información en detalles ortogonales.

El dato de partida es la salida dodecafásica ya producida por
APP--TRIT--TPK. Las dos hojas de APP conservan sus divisiones completas;
TRIT orienta el régimen y el operador
\(U_t=\operatorname{Upd}_t\circ\operatorname{Tra}_t\circ
\operatorname{Sel}_t\) transporta el estado enriquecido. La prolongación
\(w_6\to w_{12}\to w_{18}\to w_{24}\to w_{30}\to R_{36}\to G_9\)
conserva retorno de fase y avance de memoria. La estructura discreta conjunta
del continuo conserva correlativamente las publicaciones de ese estado,
con sus hojas, orientación, memoria y fronteras. De ese mismo estado proceden
las publicaciones regionales y el registro firmado cuya inversión da
\[
 K=(234,543,140,729,659,824,621,58,914,794,146,601),
 \qquad Q=\sum_{j=1}^{12}K_j=6263.
\]
Se utiliza esta salida con sus marcas y orientación; las matrices de
momentos son realizaciones posteriores. El índice de refinamiento que se
introduce abajo cuenta subdivisiones del lector geométrico. Su identificación
con una duración física o con un número de microtransiciones del TPK
requiere conservar el reloj de esa otra realización.

\subsection{Particiones marcadas y crecimiento cofinal}

Escribamos \(d_j=K_j/Q\), \(a_0=0\) y
\(a_j=\sum_{h=1}^j d_h\). Para \(r\geq0\), definimos
\[
 L_r=12\,9^r,\qquad N_r=L_r+1,\qquad
 t_{r,(j-1)9^r+c}=a_{j-1}+\frac{c}{9^r}d_j,
 \quad 0\leq c\leq9^r.
\]
Los extremos comunes se cuentan una sola vez. Las \(L_r\) celdas se
indexan por \((j,c)\), con \(1\leq j\leq12\) y
\(0\leq c<9^r\), y tienen longitud
\(\ell_{r,j,c}=d_j/9^r\). Se conserva también la escritura de \(c\)
en \(r\) cifras nonádicas, incluidos los ceros iniciales. Un hijo tiene
índice \(9c+h\), con \(0\leq h\leq8\); la división euclídea recupera
su padre y la última cifra. Los extremos compartidos retienen su marca
de frontera; la convención de intervalos semiabiertos sólo organiza
las celdas y no elimina ninguna marca.

\begin{proposition}[Cofinalidad y recuperación del registro]
\label{prop:dim-grid}
Las marcas satisfacen
\[
 0=t_{r,0}<t_{r,1}<\cdots<t_{r,L_r}=1,
 \qquad t_{r+1,9i}=t_{r,i}.
\]
El máximo de las longitudes es
\(h_r=(\max_jK_j)/(Q9^r)\), que tiende a cero. Para todo par de
enteros \(k\geq1\), \(m\geq1\), existe \(r_0\) tal que
\(k+m\leq N_r\) para \(r\geq r_0\). El registro se recupera en
cualquier nivel por
\[
 K_j=Q\sum_{c=0}^{9^r-1}\ell_{r,j,c}.
\]
\end{proposition}
\begin{proof}
Cada diferencia consecutiva es \(d_j/9^r>0\). Al reemplazar \(c\)
por \(9c\), la fórmula del nivel siguiente da exactamente la marca
anterior. Las nueve longitudes de los hijos suman la longitud del padre;
la iteración prueba la recuperación. Finalmente \(N_r=12\,9^r+1\)
crece sin cota. Una elección explícita es cualquier \(r_0\geq0\) con
\(9^{r_0}\geq(k+m-1)/12\). Todo el argumento conserva \(Q\), el
índice de bloque y las cifras de descendencia.
\end{proof}

La normalización determina las razones \(K_j/Q\). La recuperación de
los enteros utiliza, además, la carga \(Q\) conservada por el registro.
Así se distinguen la geometría normalizada y la información de escala
que acompaña a su publicación.

\subsection{Banderas de evaluaciones y realizaciones en toda dimensión finita}

Para \(1\leq q\leq N_r\), sea
\[
 V_r^{(q)}=(t_{r,i}^{\,a})_{0\leq a<q,\,0\leq i<N_r},
 \qquad F_r^q=\operatorname{fila}V_r^{(q)}.
\]
La potencia de exponente cero vale uno también en el nodo cero.

\begin{theorem}[Bandera positiva y composición de rango arbitrario]
\label{thm:dim-all-ranks}
Los espacios \(F_r^q\) forman una bandera
\(F_r^1\subset F_r^2\subset\cdots\subset F_r^{N_r}\), con
\(\dim F_r^q=q\), y cada matriz \(V_r^{(q)}\) tiene todos sus
menores máximos ordenados estrictamente positivos. Para
\(k\geq1\), \(m\geq1\), \(k+m\leq N_r\), las matrices
\[
 C_r=V_r^{(k)},\quad Z_r=V_r^{(k+m)},\quad
 Y_r=C_rZ_r^{\mathsf T}
\]
satisfacen
\[
 (Y_r)_{ab}=\sum_{i=0}^{L_r}t_{r,i}^{a+b},\qquad
 \operatorname{rank}Y_r=k,\qquad
 [Y_r]\in\mathcal A_{N_r,k,m}(Z_r).
\]
La construcción alcanza todo par finito \((k,m)\) mediante un nivel
suficientemente profundo de la misma partición marcada.
\end{theorem}
\begin{proof}
Para \(i_1<\cdots<i_q\), el menor correspondiente es
\[
 \det(t_{r,i_b}^{a-1})_{a,b=1}^{q}
   =\prod_{a<b}(t_{r,i_b}-t_{r,i_a})>0.
\]
Esto prueba rango, positividad e inclusiones de la bandera. El bloque
inicial \(k\times k\) de \(Y_r\) es el Gram de las evaluaciones
de \(1,t,\ldots,t^{k-1}\). Para \(x\neq0\),
\[
 x^{\mathsf T}(Y_r)_{[0,k-1],[0,k-1]}x
 =\sum_i\left(\sum_{a=0}^{k-1}x_at_{r,i}^{a}\right)^2>0,
\]
porque un polinomio no nulo de grado inferior a \(k\) no se anula
en los \(N_r\geq k+m\) nodos distintos. El bloque es invertible,
y \(Y_r\), que tiene \(k\) filas, tiene rango \(k\).
La pertenencia indicada es la aplicación de la matriz externa positiva
\(Z_r\) al punto positivo representado por \(C_r\).
\end{proof}

Se emplea la convención matemática
\[
 \mathcal A_{N,k,m}(Z)=
 \{[CZ^{\mathsf T}]:[C]\in\operatorname{Gr}_{\geq0}(k,N)\}.
\]
La definición para \(m\) general puede consultarse en S.~N.~Karp,
L.~K.~Williams y Y.~X.~Zhang,
\href{https://people.math.harvard.edu/~williams/papers/amplituhedron-plus.pdf}
{\emph{Decompositions of amplituhedra}}, definición~1.1.
Aquí el dato externo ya procede de la partición publicada por HMT.
La conclusión construye una familia determinada de puntos y banderas;
las dimensiones del Grassmanniano ambiente y de un atlas de su imagen
son objetos adicionales a los rangos probados.

Sea \(E_r:\mathbb R^{N_{r+1}}\to\mathbb R^{N_r}\) la restricción
a las posiciones \(9i\). Para todo \(q\leq N_r\),
\[
 E_r\bigl((V_{r+1}^{(q)})^{\mathsf T}x\bigr)
       =(V_r^{(q)})^{\mathsf T}x.
\]
Por tanto, \(E_r:F_{r+1}^q\to F_r^q\) es un isomorfismo sobre
estos espacios de evaluaciones: sus dos realizaciones identifican el
mismo polinomio de grado inferior a \(q\). La composición de las
restricciones coincide con la restricción directa. Los menores sobre
índices heredados conservan exactamente su valor.

\begin{proposition}[Recuperación desde una bandera calibrada]
Si se conservan el orden de las columnas, las marcas extremas y la
calibración \(t_{r,0}=0\), \(t_{r,L_r}=1\), el plano \(F_r^2\)
determina los nodos y las separaciones. Junto con \(Q\) y las
marcas de descendencia determina \(K\).
\end{proposition}
\begin{proof}
El plano contiene el vector constante y el vector de nodos. Cualquier
otra coordenada que, junto con el vector constante, genere el mismo
plano es \(a+bt\), con \(b\neq0\). Las dos calibraciones fijan
\(a=0\), \(b=1\). Las diferencias consecutivas recuperan las
longitudes, y la proposición~\ref{prop:dim-grid} recupera el registro.
\end{proof}

\subsection{Cuerpos convexos de rango uno y crecimiento del dominio}

Para \(m\geq1\) y \(m+1\leq N_r\), definimos
\[
 P_r^{(m)}=\operatorname{conv}
 \{(t_{r,i},t_{r,i}^2,\ldots,t_{r,i}^m):0\leq i\leq L_r\}.
\]
La normalización de una fila positiva de \(C\) como coeficientes
baricéntricos prueba
\(\mathcal A_{N_r,1,m}(Z_r)=P_r^{(m)}\). El determinante de
Vandermonde proporciona \(m+1\) puntos afínmente independientes,
de modo que el cuerpo tiene dimensión \(m\).

\begin{proposition}[Cobertura simplicial a dimensión finita y límite convexo]
\label{prop:dim-polytopes}
Cada \(P_r^{(m)}\) admite una triangulación por sus vértices, con
grado uno sobre el interior de cada celda simplicial. La suma de sus
formas canónicas tiene residuos interiores cancelados. Además,
\[
 P_r^{(m)}\subseteq P_{r+1}^{(m)},\qquad
 \overline{\bigcup_{r\geq r_0}P_r^{(m)}}=
 \operatorname{conv}\{(t,t^2,\ldots,t^m):0\leq t\leq1\}.
\]
\end{proposition}
\begin{proof}
Se ordenan los vértices por sus marcas. En cada cara se elige su menor
vértice; se triangulan recursivamente las facetas que no lo contienen
y se forman los conos correspondientes. La inducción sobre la dimensión
prueba cobertura e interiores disjuntos, y la misma regla en caras
comunes garantiza compatibilidad. En cada simplex orientado
\([v_0,\ldots,v_m]\), con matriz
\(A=(v_1-v_0,\ldots,v_m-v_0)\) de determinante positivo y
coordenadas baricéntricas \(\lambda_0,\ldots,\lambda_m\), la
aplicación baricéntrica es afín e invertible. Su forma es
\[
 \Omega_{[v_0,\ldots,v_m]}=
 \frac{dx_1\wedge\cdots\wedge dx_m}
 {\det A\prod_{j=0}^{m}\lambda_j}.
\]
La restricción residual a cada cara es su forma orientada. Dos simplices
adyacentes inducen orientaciones opuestas sobre su cara común, de modo
que los residuos interiores se cancelan. Esta es la misma operación de
triangulación y residuo aplicada en la realización de rango uno anterior,
con el número de vértices ahora aumentado por el refinamiento.

La inclusión de nodos heredados prueba la inclusión de cuerpos. La malla
tiende a cero, por lo que las curvas muestreadas son densas en la curva
compacta de momentos. Toda combinación convexa finita de puntos de la
curva se aproxima por combinaciones con nodos de un nivel común.
Recíprocamente, todos los cuerpos están contenidos en la envolvente
convexa compacta de la curva. Ambas inclusiones prueban la igualdad.
\end{proof}

La cancelación de residuos opera sobre una triangulación o subdivisión
de un cuerpo fijo. El paso de \(P_r^{(m)}\) a \(P_{r+1}^{(m)}\)
añade vértices y puede ampliar el cuerpo. Sus formas pertenecen a esos
dominios respectivos. El límite convexo precedente especifica la
convergencia de los cuerpos, sin imponer una identidad entre sus formas.

\subsection{Medida del calendario y convergencia de los momentos nodales}

Sea \(F_K:[0,1]\to[0,1]\) la función creciente afín por tramos
que lleva \((j-1)/12\) a \(a_{j-1}\) y \(j/12\) a \(a_j\).
En el intervalo correspondiente su pendiente es \(12d_j>0\).
La medida del calendario transportado es
\[
 \mu_K=(F_K)_*(ds),\qquad
 \frac{d\mu_K}{dt}=\frac1{12d_j}\quad
 (a_{j-1}<t<a_j).
\]
Cada intervalo inicial tiene masa \(1/12\), y cada celda de nivel
\(r\) tiene masa \(1/L_r\). Esta elección asigna el mismo peso
a las posiciones del calendario antes de realizar sus longitudes.
La medida de longitud \(dt\) es otro lector; ambas coinciden cuando
todas las separaciones iniciales son iguales.

Sus momentos son racionales y están determinados por \(K\):
\begin{equation}
 M_p=\int_0^1t^p\,d\mu_K(t)
   =\frac1{12(p+1)}\sum_{j=1}^{12}
       \frac{a_j^{p+1}-a_{j-1}^{p+1}}{d_j},\qquad p\geq0.
 \label{eq:dim-moments}
\end{equation}
\begin{proposition}[Recuperación por momentos y cuantiles]
\label{prop:dim-moment-recovery}
La sucesión completa \((M_p)_{p\geq0}\) determina \(\mu_K\).
Si \(H_K(t)=\mu_K([0,t])\) es su función de distribución, entonces
\[
 a_j=H_K^{-1}(j/12),\qquad
 K_j=Q\bigl(H_K^{-1}(j/12)-H_K^{-1}((j-1)/12)\bigr).
\]
En consecuencia, los momentos completos, la carga \(Q\) y el orden
del calendario recuperan el registro inicial.
\end{proposition}
\begin{proof}
Dos probabilidades sobre \([0,1]\) con los mismos momentos tienen las
mismas integrales para todos los polinomios. La aproximación uniforme
de funciones continuas por polinomios y la masa finita extienden esa
igualdad a todas las funciones continuas; éstas determinan la medida
de Borel. La densidad de \(\mu_K\) es positiva y constante en cada
intervalo inicial. Su distribución es, por tanto, continua y
estrictamente creciente, y cumple \(H_K(a_j)=j/12\). La inversión
de esta igualdad y la diferencia de extremos prueban las fórmulas.
\end{proof}

La probabilidad nodal correspondiente al teorema
\ref{thm:dim-all-ranks} es
\(\nu_r=N_r^{-1}\sum_{i=0}^{L_r}\delta_{t_{r,i}}\).

\begin{theorem}[Límite de la publicación nodal]
\label{thm:dim-nodal-limit}
Se cumple \(\nu_r\rightharpoonup\mu_K\). Si \(f\) es continua,
\[
 \left|\int f\,d\nu_r-\int f\,d\mu_K\right|
 \leq \omega_f(h_r)+\frac{\operatorname{osc}(f)}{N_r},
\]
donde \(\omega_f\) es su módulo de continuidad. Para \(p\geq1\),
el error del momento de orden \(p\) está acotado por
\(p h_r+1/N_r\). Para cada par fijo \((k,m)\),
\[
 \frac1{N_r}Y_r\longrightarrow (M_{a+b})_{
     0\leq a<k,\,0\leq b<k+m}=:Y_\infty^{(k,m)},
 \qquad\operatorname{rank}Y_\infty^{(k,m)}=k.
\]
\end{theorem}
\begin{proof}
La medida que asigna masa \(1/L_r\) a cada extremo izquierdo de
celda aproxima su integral con error a lo sumo \(\omega_f(h_r)\).
En efecto, cada celda tiene esa masa y diámetro no mayor que \(h_r\).
La medida \(\nu_r\) es \(L_r/N_r\) veces dicha medida más
\(\delta_1/N_r\), lo que añade a lo sumo
\(\operatorname{osc}(f)/N_r\). Para \(f(t)=t^p\), la constante
de Lipschitz es \(p\). La convergencia de cada entrada se sigue de
estas cotas. El bloque inicial de la matriz límite es definido positivo:
la integral del cuadrado de un polinomio no nulo es positiva, porque
\(\mu_K\) tiene densidad estrictamente positiva en todo intervalo
abierto de \([0,1]\). Así el límite conserva rango \(k\).
\end{proof}

Los momentos nodales varían exactamente por la incorporación de nodos.
Si \(\mathcal J_r\) es el conjunto de las nuevas posiciones y
\(z_q(t)=(1,t,\ldots,t^{q-1})^{\mathsf T}\), entonces
\[
 V_{r+1}^{(q)}(V_{r+1}^{(q)})^{\mathsf T}
 =V_r^{(q)}(V_r^{(q)})^{\mathsf T}
  +\sum_{i\in\mathcal J_r}z_q(t_{r+1,i})z_q(t_{r+1,i})^{\mathsf T}.
\]
Cada sumando es positivo semidefinido. Para las matrices rectangulares
\(Y_r\) vale la misma actualización con
\(z_kz_{k+m}^{\mathsf T}\). La igualdad expresa el cambio de la
medida de conteo; la normalización y el teorema anterior describen su
límite. El objeto límite de rango \(k\) tiene coordenadas en un
Grassmanniano finito, mientras los datos externos \(Z_r\) pertenecen
a las distintas publicaciones de cada nivel.

\subsection{Promedios celulares, operadores compatibles y detalle ortogonal}

Para una celda \(I_{r,i}=[b_{r,i},b_{r,i}+\ell_{r,i}]\), ponemos
\[
 a_{r,i}^{(p)}=
 \frac{(b_{r,i}+\ell_{r,i})^{p+1}-b_{r,i}^{p+1}}
 {(p+1)\ell_{r,i}},\qquad
 A_r^{(q)}=(a_{r,i}^{(p)})_{0\leq i<L_r,\,0\leq p<q}.
\]
Son las medias de los monomios respecto de la probabilidad condicional
de \(\mu_K\) en cada celda. Esta probabilidad es uniforme en la
celda, aunque \(\mu_K\) no sea la medida de longitud global.

Sea \(\mathcal H_r=\mathbb R^{L_r}\) con producto
\(\langle x,y\rangle_r=L_r^{-1}\sum_i x_i y_i\). La aplicación
\(J_r:\mathcal H_r\to\mathcal H_{r+1}\) repite el valor de un
padre en sus nueve hijos. Su adjunta \(R_r=J_r^*\) promedia esos
nueve valores. Se cumple \(R_rJ_r=I\).

\begin{theorem}[Compatibilidad exacta y balance de los momentos]
\label{thm:dim-cell-balance}
Para todo \(p\geq0\),
\[
 R_r a_{r+1}^{(p)}=a_r^{(p)},\qquad
 \frac1{L_r}\sum_i a_{r,i}^{(p)}=M_p.
\]
Definamos
\[
 D_r^{(q)}=A_{r+1}^{(q)}-J_rA_r^{(q)},\quad
 G_r^{(q)}=\frac1{L_r}(A_r^{(q)})^{\mathsf T}A_r^{(q)}.
\]
Entonces
\[
 R_rD_r^{(q)}=0,\qquad
 G_{r+1}^{(q)}=G_r^{(q)}+
       \frac1{L_{r+1}}(D_r^{(q)})^{\mathsf T}D_r^{(q)},
\]
y
\[
 G_r^{(q)}\longrightarrow H_K^{(q)}=(M_{a+b})_{0\leq a,b<q}.
\]
Los detalles y el nivel inicial reconstruyen todos los niveles. Para
\(q\leq L_r\), la matriz \(A_r^{(q)}\) tiene rango \(q\) y
todos los menores ordenados de orden \(q\) son positivos.
\end{theorem}
\begin{proof}
Las masas de los nueve hijos son iguales y sus intervalos particionan
el intervalo padre. La aditividad de la integral da la primera identidad;
la suma sobre todas las celdas da la segunda. La descomposición
\(A_{r+1}=J_rA_r+D_r\) es ortogonal en cada columna porque
\(R_rD_r=0\). Al expandir su Gram se anulan los términos cruzados,
y \(J_r\) es una isometría. Esto prueba el balance matricial.

El observable constante por celdas que toma el valor \(a_{r,i}^{(p)}\)
aproxima uniformemente \(t^p\), con error no mayor que \(p h_r\)
para \(p\geq1\). Por tanto converge en \(L^2(\mu_K)\), y los
productos internos convergen a \(M_{a+b}\). La reconstrucción se
obtiene iterando \(A_{r+1}=J_rA_r+D_r\).

Para probar los menores, elijamos \(q\) celdas ordenadas. Por
multilinealidad y Fubini, el determinante de sus promedios es la media
sobre el producto de esas celdas de
\(\prod_{a<b}(x_b-x_a)\). Este integrando es positivo en el
interior del producto porque las celdas son disjuntas y están ordenadas.
Sus fronteras tienen medida nula. La media es estrictamente positiva.
\end{proof}

La identidad se prolonga al par de grados \(k\) y \(k+m\):
\[
 \overline Y_r=\frac1{L_r}(A_r^{(k)})^{\mathsf T}A_r^{(k+m)},
 \qquad
 \overline Y_{r+1}-\overline Y_r
 =\frac1{L_{r+1}}(D_r^{(k)})^{\mathsf T}D_r^{(k+m)}.
\]
Para \(k+m\leq L_r\), las matrices de promedios definen una
segunda realización amplituhedral positiva, de \(L_r\) columnas,
y \(\overline Y_r\) tiene rango \(k\). Son promedios de celdas,
mientras \(Y_r\) usa evaluaciones en \(N_r=L_r+1\) nodos. Ambas
realizaciones convergen a \(Y_\infty^{(k,m)}\); su coincidencia
en el límite procede de las demostraciones anteriores.

La versión operatoria es igualmente explícita. Para una función continua
\(f\), sea \(\Psi_r(f)\) el operador diagonal de sus medias
condicionales en \(\mathcal H_r\). Es lineal, positivo y unital, y
\[
 R_r\Psi_{r+1}(f)J_r=\Psi_r(f).
\]
Esta igualdad se prueba aplicando ambos miembros a cada vector base.
Para \(X(t)=t\), los dos primeros operadores determinan, celda a
celda,
\[
 \Psi_r(X)_i=b_{r,i}+\frac{\ell_{r,i}}2,\qquad
 \Psi_r(X^2)_i-\Psi_r(X)_i^2=\frac{\ell_{r,i}^2}{12}.
\]
Si \(v_{r,i}\) denota esa varianza, la fórmula
\(\ell_{r,i}=\sqrt{12v_{r,i}}\) recupera la longitud; la
primera identidad recupera el extremo inicial. Con orientación, marcas y \(Q\),
se recupera nuevamente \(K\). El promedio es una compresión positiva:
la varianza escrita cuantifica por qué \(\Psi_r(X^2)\) y
\(\Psi_r(X)^2\) son operaciones distintas.

\subsection{Transporte de la medida y alcance del crecimiento dimensional}

Dos listas positivas de separaciones \(d\) y \(d'\), con las
mismas marcas, determinan un homeomorfismo creciente afín por intervalos
\(F=F_{K'}\circ F_K^{-1}\). Lleva cada celda a su celda homóloga
y satisface \(F_*\mu_K=\mu_{K'}\). Por ello
\[
 \mathcal U_F:L^2(\mu_K)\longrightarrow L^2(\mu_{K'}),
 \qquad\mathcal U_F f=f\circ F^{-1}
\]
es una isometría sobreyectiva. Si \(P_r\) y \(P'_r\) son las
proyecciones a funciones constantes por celdas, entonces
\(P'_r\mathcal U_F=\mathcal U_FP_r\): las medias condicionales
se transportan porque las masas de las celdas correspondientes coinciden.
Esto aplica en particular al flujo positivo de separaciones dirigido por
\(K\), conservando la dirección de cada bloque en sus hijos. Los
momentos de la coordenada espacial pueden cambiar; la isometría transporta
el observable junto con la medida y sus proyecciones.

El crecimiento demostrado es cofinal en toda dimensión externa finita.
Conserva un sistema de banderas positivas, una familia explícita de
puntos amplituhedrales, cuerpos convexos completos en rango uno y una
medida cuyos operadores de promedio son exactamente compatibles. La
memoria de los detalles permite invertir la compresión entre niveles.
Las pruebas de cobertura y residuos de rango uno conservan su dominio;
para rangos mayores, el resultado presente fija matrices, positividad,
recuperación, compatibilidad y límite con tipos explícitos. Esta
separación sitúa el alcance de cada construcción en el mismo desarrollo
sin confundir crecimiento de dimensión, parametrización de un punto y
clasificación de todas las celdas de una imagen.


% SOURCE sections/jacobi_spectral.tex
\section{Momentos genealógicos, operadores de Jacobi y residuos espectrales}
\label{sec:jacobi-generated}

La medida inducida por el calendario proporciona un enlace entre la
geometría de las celdas y una realización espectral positiva. Este enlace
utiliza toda la partición marcada: cada uno de sus doce intervalos conserva
la misma masa cronológica, mientras su longitud procede de la coordenada
correspondiente del registro. La densidad compensa exactamente esa
longitud. Por consiguiente, la distribución espacial de la medida retiene
la anisotropía del registro aunque la distribución entre ventanas sea
uniforme.

Sean \(K_j>0\), \(Q=\sum_{j=1}^{12}K_j\), \(d_j=K_j/Q\) y
\(t_j=\sum_{i\leq j}d_i\), con \(t_0=0\). La medida de probabilidad
utilizada en esta sección es
\[
 \dd\mu_K(t)=\sum_{j=1}^{12}
 \frac{\mathbf1_{[t_{j-1},t_j]}(t)}{12d_j}\,\dd t.
\]
Los extremos compartidos tienen medida nula. Sus momentos son racionales
para todo registro entero:
\begin{equation}
 m_r=\int_0^1t^r\,\dd\mu_K(t)
 =\frac1{12(r+1)}\sum_{j=1}^{12}
   \frac{t_j^{r+1}-t_{j-1}^{r+1}}{d_j},\qquad r\geq0.
 \label{eq:jacobi-moments}
\end{equation}
La igualdad \(m_0=1\) expresa conservación de la masa total. Para cada
\(n\geq1\), la matriz de Hankel \(H_n=(m_{a+b})_{0\leq a,b<n}\)
es definida positiva: su forma cuadrática es la integral del cuadrado
de un polinomio no nulo sobre un intervalo donde la densidad es positiva.
Esta observación permite pasar de los momentos a un operador sin elegir
un espectro objetivo.

\begin{theorem}[Recurrencia tridiagonal determinada por la partición]
\label{thm:jacobi-recurrence}
Existe una única sucesión de polinomios mónicos \(p_n\), de grado \(n\),
ortogonales para \(\mu_K\). Si
\[
 h_n=\int p_n^2\,\dd\mu_K>0,\qquad
 a_n=\frac{\int t p_n^2\,\dd\mu_K}{h_n},\qquad
 \beta_n=\frac{h_n}{h_{n-1}}>0\quad(n\geq1),
\]
entonces \(p_0=1\), \(p_1=t-a_0\) y
\begin{equation}
 p_{n+1}(t)=(t-a_n)p_n(t)-\beta_np_{n-1}(t),\qquad n\geq1.
 \label{eq:jacobi-recurrence}
\end{equation}
Todos los coeficientes \(a_n,\beta_n\) son racionales para el registro
entero declarado. En la base ortonormal \(\psi_n=p_n/\sqrt{h_n}\),
la compresión de la multiplicación por \(t\) a los polinomios de grado
menor que \(n\) tiene matriz simétrica
\[
 J_n=\begin{pmatrix}
 a_0&\sqrt{\beta_1}&&0\\
 \sqrt{\beta_1}&a_1&\ddots&\\
 &\ddots&\ddots&\sqrt{\beta_{n-1}}\\
 0&&\sqrt{\beta_{n-1}}&a_{n-1}
 \end{pmatrix}.
\]
\end{theorem}
\begin{proof}
La positividad de \(H_n\) garantiza que el sistema lineal que impone
ortogonalidad a un polinomio mónico tiene solución única. Como sus
coeficientes son momentos racionales, la solución es racional. Se
descompone \(tp_n\) en la base mónica hasta grado \(n+1\). Para
\(j\leq n-2\), la simetría del producto interno da
\(\langle tp_n,p_j\rangle=\langle p_n,tp_j\rangle=0\).
Quedan únicamente los términos de grados \(n+1,n,n-1\). Sus
coeficientes son, respectivamente, \(1,a_n,h_n/h_{n-1}\), lo que
demuestra la recurrencia. La normalización por \(\sqrt{h_n}\)
convierte los dos coeficientes adyacentes en \(\sqrt{\beta_n}\).
\end{proof}

Para el registro publicado en la sección~\ref{sec:register}, la primera
posición diagonal y el cuadrado del primer acoplamiento son
\[
 a_0=m_1=\frac{71195}{150312},\qquad
 \beta_1=m_2-m_1^2=\frac{2206226167}{22593697344}.
\]
La primera igualdad calcula el centro de masa de la partición cronológica;
la segunda calcula su dispersión y determina el primer enlace de la
matriz tridiagonal. Ambas cifras se obtienen integrando las longitudes
generadas, y su significado estructural antecede a cualquier evaluación
decimal.

\begin{theorem}[Espectro simple, cuadratura y residuos positivos]
\label{thm:jacobi-residues}
Los autovalores \(\lambda_1<\cdots<\lambda_n\) de \(J_n\) son
simples y pertenecen a \((0,1)\). Existen pesos únicos
\(w_i>0\), de suma uno, tales que
\begin{equation}
 \int f(t)\,\dd\mu_K(t)=\sum_{i=1}^{n}w_i f(\lambda_i)
 \qquad(\deg f\leq2n-1).
 \label{eq:jacobi-quadrature}
\end{equation}
El resolvente observado en la constante unitaria \(e_0\) satisface
\begin{equation}
 g_n(z)=\langle e_0,(zI-J_n)^{-1}e_0\rangle
       =\sum_{i=1}^{n}\frac{w_i}{z-\lambda_i}
       =\cfrac1{z-a_0-\cfrac{\beta_1}{z-a_1-
          \cfrac{\beta_2}{\ddots-\cfrac{\beta_{n-1}}{z-a_{n-1}}}}}.
 \label{eq:jacobi-resolvent}
\end{equation}
En particular, sus polos y residuos se obtienen del registro por las
operaciones de momento, ortogonalización y compresión.
\end{theorem}
\begin{proof}
Para cualquier polinomio no nulo \(p\) de grado menor que \(n\),
\[
 0<\int_0^1 t|p(t)|^2\,\dd\mu_K(t)
   <\int_0^1 |p(t)|^2\,\dd\mu_K(t).
\]

La caracterización variacional coloca el espectro de \(J_n\) en
\((0,1)\). Cada subdiagonal es positiva. Las ecuaciones de un autovector
determinan recursivamente todas sus coordenadas a partir de la primera;
si ésta se anulase, se anularía el vector entero. Cada autoespacio tiene
así dimensión uno. La simetría de la matriz da diagonalización ortogonal
y simplicidad del espectro. Para autovectores unitarios \(v_i\), se toma
\(w_i=|\langle e_0,v_i\rangle|^2>0\). La completitud de la base
propia da \(\sum_iw_i=1\) y la primera expresión del resolvente.

La recurrencia tridiagonal muestra por expansión del determinante que
\(\det(tI-J_n)=p_n(t)\). Para \(\deg f\leq2n-1\), la división
euclídea escribe \(f=qp_n+r\), con \(\deg q,\deg r<n\).
La ortogonalidad anula la integral de \(qp_n\), y su evaluación en los
\(\lambda_i\) también se anula. Para \(r\) de grado menor que
\(n\), las sucesivas multiplicaciones de la constante por \(t\)
permanecen dentro del espacio de polinomios hasta el grado requerido;
por ello \(\int r\,\dd\mu_K=\langle e_0,r(J_n)e_0\rangle\).
La resolución espectral prueba la fórmula de cuadratura. Su unicidad
se deduce de la invertibilidad de la matriz de Vandermonde evaluada
en las raíces distintas. Finalmente, el complemento de Schur de la
primera fila y columna de \(zI-J_n\), iterado sobre las submatrices
terminales, produce la fracción continua.
\end{proof}

\begin{figure}[htbp]
\centering
\begin{tikzpicture}[>=Latex, node distance=8mm and 8mm,
 every node/.style={font=\small}, box/.style={draw,rounded corners=1pt,
 text width=37mm,minimum height=13mm,align=center}]
\node[box] (k) {Registro marcado\\\(K,Q\)};
\node[box,right=of k] (m) {Medida cronológica\\\(\mu_K\)};
\node[box,right=of m] (h) {Momentos positivos\\\(H_n\)};
\node[box,below=of h] (j) {Operador tridiagonal\\\(J_n\)};
\node[box,left=of j] (g) {Resolvente\\\(g_n(z)\)};
\node[box,left=of g] (w) {Polos y residuos\\\(\lambda_i,w_i>0\)};
\draw[->] (k)--(m);\draw[->] (m)--(h);\draw[->] (h)--(j);
\draw[->] (j)--(g);\draw[->] (g)--(w);
\end{tikzpicture}
\caption{Realización espectral de la partición HMT. La geometría de las
celdas determina la medida; la positividad de sus momentos produce el
operador y los residuos positivos de su resolvente.}
\label{fig:jacobi-generated}
\end{figure}

La fórmula del resolvente aporta una lectura espectral de la memoria
espacial. Los pesos positivos describen cuánto observa la coordenada
constante de cada modo; los polos son las frecuencias algebraicas del
operador de multiplicación comprimido. Los momentos de orden creciente
recuperan detalles sucesivos de la medida, y los coeficientes de la
fracción continua conservan esa información mediante una recurrencia
local de tres términos. La apariencia clásica de una matriz de Jacobi
se reconoce después de haber construido sus coeficientes desde las
celdas HMT. La cadena causal queda fijada por
\[
 K\longmapsto(d,t)\longmapsto\mu_K\longmapsto(m_r)_{r\geq0}
 \longmapsto\bigl((a_n)_{n\geq0},(\beta_n)_{n\geq1}\bigr)
 \longmapsto(J_n,g_n).
\]

\begin{proposition}[Compatibilidad de dimensión y de refinamiento]
\label{prop:jacobi-refinement}
Las matrices \(J_n\) son compresiones principales compatibles de la
misma multiplicación acotada por \(t\) en \(L^2(\mu_K)\). Sus
medidas espectrales \(\sum_iw_i\delta_{\lambda_i}\) convergen
débilmente a \(\mu_K\). Si cada subintervalo nonádico se representa
por su punto medio y por su masa \(1/(12\cdot9^r)\), la medida
atómica resultante \(\mu_{K,r}\) aproxima los momentos de
\(\mu_K\). Para un grado fijo \(n\), la construcción desde esos
momentos recupera los coeficientes \(a_j,\beta_j\) de la misma
matriz cuando el refinamiento tiende a profundidad infinita.
\end{proposition}
\begin{proof}
La ortogonalización de la sucesión completa utiliza las mismas fórmulas
al aumentar \(n\), de modo que cada bloque principal permanece fijo.
La cuadratura reproduce exactamente cualquier polinomio cuando
\(2n-1\) supera su grado. La aproximación uniforme de funciones
continuas en \([0,1]\) por polinomios, junto con la masa total uno,
prueba la convergencia débil. Las sumas de puntos medios convergen
en cada uno de los doce intervalos con su densidad constante, y por
tanto aproximan todos los momentos. En dimensión fija, los coeficientes de
ortogonalización son funciones racionales de un número finito de
momentos, con denominadores formados por determinantes positivos de
Gram. La convergencia de esos momentos y la positividad del límite
aseguran la convergencia de los coeficientes.
\end{proof}

\begin{falsifier}
Una longitud nula invalida la densidad declarada y puede destruir el
rango de los momentos. La sustitución de la medida fija por conteos
no normalizados en mallas distintas cambia el operador. Los residuos
de \(g_n\) son pesos espectrales del resolvente, mientras los residuos
de una forma canónica se definen sobre divisores geométricos: una
identificación entre ambas nociones exige especificar la forma y el
mapa correspondiente. Del mismo modo, \(J_n\) es adimensional; su
identificación con un Hamiltoniano físico requiere la representación
y la escala de acción y reloj que utilice ese sector.
\end{falsifier}

\begin{theorem}[Fidelidad de la realización espectral marcada]
\label{thm:jacobi-marked-faithfulness}
El operador infinito de Jacobi \(J\), junto con su vector cíclico
unitario \(e_0\) y la carga \(Q\), determina el registro \(K\).
Dos realizaciones de este tipo son unitariamente equivalentes mediante
una aplicación que conserva el vector marcado y la carga si y sólo
si sus registros coinciden. El espectro como conjunto es, en cambio,
\([0,1]\) para todo registro positivo del dominio.
\end{theorem}
\begin{proof}
Los polinomios son densos en \(L^2(\mu_K)\): primero se aproximan
funciones continuas y después se usa su densidad para la medida finita.
La base ortonormal construida define una aplicación unitaria que lleva
\(e_n\) a \(\psi_n\). Bajo ella, \(J\) es la multiplicación por
\(t\), un operador autoadjunto acotado, y \(e_0\) representa la
función constante uno. Por tanto,
\[
 \langle e_0,J^re_0\rangle=m_r\qquad(r\geq0).
\]
Estos momentos determinan la medida: dos medidas con los mismos
momentos coinciden en polinomios, después en funciones continuas por
aproximación uniforme y finalmente como medidas de Borel.
Su distribución \(F(t)=\mu_K([0,t])\) es continua y estrictamente
creciente. Cada intervalo inicial tiene masa \(1/12\), de modo que
\[
 t_j=F^{-1}(j/12),\qquad
 K_j=Q\bigl(F^{-1}(j/12)-F^{-1}((j-1)/12)\bigr).
\]
Una equivalencia unitaria que conserva \(e_0\) conserva los momentos
y por ello el registro. La recíproca se sigue de la construcción única.

Para el último enunciado, la densidad de \(\mu_K\) es positiva
en todo subintervalo. Cada \(\lambda\in[0,1]\) admite funciones
unitarias concentradas en intervalos cada vez menores alrededor de
\(\lambda\), sobre las que la norma de \((J-\lambda I)f\)
tiende a cero. Así todo el intervalo pertenece al espectro. Fuera de
él, la multiplicación por \(1/(t-\lambda)\) es acotada y define
el resolvente. Esto prueba la igualdad del espectro.
\end{proof}

El teorema distingue la extensión del espectro y la información que
éste conserva respecto de una observación marcada. El mismo intervalo
espectral admite las distintas distribuciones de la familia; la medida
observada desde la unidad recupera las separaciones originales. El
enlace con la representación grassmanniana es ahora reversible:
los menores marcados reconstruyen \(K\), el registro construye
\((J,e_0,Q)\), y sus momentos y cuantiles recuperan el mismo registro.
Cada lector que factoriza por \((K,\xi)\) posee, por consiguiente,
una expresión espectral marcada una vez conservadas las restantes
coordenadas \(\xi\). La realización espectral contiene estructura,
además de posiciones de polos o valores de autovalores.


% SOURCE appendices/figures.tex
\begin{figure}[htbp]
\centering
\begin{tikzpicture}[x=1.9cm,y=1.05cm,>=Stealth,font=\small]
\foreach \j/\lab in {0/0,1/1,2/2,3/3,5/12}{
  \node[circle,fill=black,inner sep=1.6pt,label=above:{$a_{\lab}$}] (a\j) at (\j,1.4) {};
  \node[circle,fill=black,inner sep=1.6pt,label=below right:{$b_{\lab}$}] (b\j) at (\j,0) {};
  \node (q\j) at (\j,-1) {$q_{\lab}$};
  \draw[->] (b\j) -- (q\j);
}
\foreach \j/\prev in {1/0,2/1,3/2}{
 \draw[->] (a\prev)--(a\j);\draw[->] (b\prev)--(b\j);
 \draw[->,blue!65!black,thick] (a\j)--node[right]{$d_{\j}$}(b\j);
}
\draw[->,blue!65!black,thick] (a5)--node[right]{$d_{12}$}(b5);
\draw[dashed,->] (a3)--(a5);\draw[dashed,->] (b3)--(b5);
\node at (4,.7) {$\cdots$};
\node[left] at (a0) {$s_2$};\node[left] at (b0) {$s_1$};
\end{tikzpicture}
\caption{Red planar de dos filas. Las aristas horizontales y de salida
tienen peso uno. Cada camino desde la fuente superior desciende una sola
vez. Para dos destinos \(i<j\), las parejas disjuntas descienden después
de \(i\) y antes de \(j\); su peso total es
\(\Delta_{ij}=d_{i+1}+\cdots+d_j\). Los trazos discontinuos conservan
las ocho posiciones intermedias omitidas únicamente en el dibujo.}
\end{figure}

\begin{figure}[htbp]
\centering
\begin{tikzpicture}[x=9cm,y=6cm,font=\small]
\foreach \j/\num in {0/0,1/234,2/777,3/917,4/1646,5/2305,6/3129,7/3750,8/3808,9/4722,10/5516,11/5662,12/6263}{
 \pgfmathsetmacro{\tx}{\num/6263}
 \coordinate (v\j) at (\tx,\tx*\tx);
}
\fill[blue!5] (v0)--(v1)--(v2)--(v3)--(v4)--(v5)--(v6)--(v7)--(v8)--(v9)--(v10)--(v11)--(v12)--cycle;
\foreach \j in {2,...,11}{\draw[blue!45,thin] (v0)--(v\j);}
\draw[blue!80!black,thick] (v0)--(v1)--(v2)--(v3)--(v4)--(v5)--(v6)--(v7)--(v8)--(v9)--(v10)--(v11)--(v12)--cycle;
\foreach \j in {0,...,12}{\fill (v\j) circle[radius=.035cm];}
\node[below left] at (v0) {$v_0$};
\node[below right] at (v3) {$v_3$};
\node[below right] at (v6) {$v_6$};
\node[below right] at (v9) {$v_9$};
\node[above] at (v12) {$v_{12}$};
\end{tikzpicture}
\caption{Polígono racional determinado por los trece parámetros
\(t_j\) de \eqref{eq:polygon-Z}, dibujado con sus coordenadas efectivas
\((t_j,t_j^2)\). La triangulación desde \(v_0\) tiene once celdas.
Cada diagonal interior recibe dos residuos opuestos. Se rotulan cinco
vértices para facilitar la lectura; los trece están representados.}
\end{figure}


% SOURCE appendices/convex_refinement.tex
\section{Flujo de pesos, geometría de información y dualidad de Legendre}
\label{sec:convex-refinement}

La deformación de los intervalos admite una descripción variacional
exacta. Sean \(d_j=K_j/Q>0\), \(\sum_jd_j=1\), y \(u_j\) las
coordenadas de la dirección \eqref{eq:u-main}. Para \(s\in\RR\),
se definen
\begin{equation}
 Z(s)=\sum_{j=1}^{12}d_j e^{su_j},\qquad
 \psi(s)=\log Z(s),\qquad
 d_j(s)=\frac{d_j e^{su_j}}{Z(s)}.
 \label{eq:convex-partition}
\end{equation}
Todos los argumentos de la exponencial son adimensionales. El parámetro
\(s\) es la coordenada de esta deformación geométrica. La asignación de
una duración física pertenece a un lector temporal declarado.

\begin{theorem}[Convexidad estricta y coordenada dual]
La función \(\psi\) es analítica y estrictamente convexa sobre \(\RR\).
Si \(u_- =\min_j u_j\) y \(u_+=\max_j u_j\), la aplicación
\(v=\psi'(s)\) es un difeomorfismo de \(\RR\) sobre \((u_-,u_+)\).
Para cada \(v\) en ese intervalo existe una única coordenada \(s(v)\), y
\[
 \psi^*(v)=s(v)v-\psi(s(v)),\qquad
 (\psi^*)'(v)=s(v),\qquad
 (\psi^*)''(v)=\frac1{\psi''(s(v))}>0.
\]
\end{theorem}
\begin{proof}
Las sumas finitas permiten diferenciar término a término. Se obtiene
\begin{align}
 \psi'(s)&=\sum_jd_j(s)u_j=:\overline u(s),\notag\\
 \psi''(s)&=\sum_jd_j(s)(u_j-\overline u(s))^2.
 \label{eq:variance}
\end{align}
Los pesos son estrictamente positivos y los \(u_j\) no son todos iguales;
la varianza es por ello estrictamente positiva. Cuando \(s\) tiende a
\(+\infty\), dividir numerador y denominador de \(\psi'\) por
\(e^{su_+}\) muestra que sólo sobreviven los índices con \(u_j=u_+\).
Así \(\psi'(s)\to u_+\). El argumento con \(u_-\) da el límite
opuesto. La monotonía y el teorema de la función inversa proporcionan la
biyectividad suave. La función \(sv-\psi(s)\) tiene su máximo único
en \(s(v)\); diferenciar el valor máximo demuestra las dos identidades
duales.
\end{proof}

La evolución de cada intervalo es
\begin{equation}
 \dot d_j(s)=d_j(s)\bigl(u_j-\overline u(s)\bigr).
 \label{eq:replicator}
\end{equation}
La suma de las derivadas es cero. Por tanto la longitud total permanece
unitaria mientras se redistribuyen las longitudes internas. Esta conservación
de escala global permite atribuir el cambio de las razones dobles a una
deformación de posiciones y no a una dilatación común.

En el simplex abierto de pesos, la métrica
\(g_d(\xi,\zeta)=\sum_j\xi_j\zeta_j/d_j\), sobre vectores tangentes
de suma cero, es positiva. La función lineal
\(F(d)=\sum_jd_ju_j\) tiene como gradiente para esta métrica
precisamente el campo de \eqref{eq:replicator}, pues
\[
 g_d\bigl(d_j(u_j-\overline u),\xi\bigr)
 =\sum_j(u_j-\overline u)\xi_j=\sum_ju_j\xi_j=dF(\xi).
\]
En consecuencia \(dF/ds=\psi''(s)>0\) a lo largo del flujo. El
carácter ascendente corresponde a esta función y a esta métrica concretas.

\begin{theorem}[Invariancia variacional por subdivisión heredada]
Para cada intervalo \(j\), sean \(r_{j,c}>0\),
\(\sum_c r_{j,c}=1\), los pesos relativos de una subdivisión fija.
Asigne a los hijos
\[
 d_{j,c}=r_{j,c}d_j,\qquad u_{j,c}=u_j.
\]
La función de partición, su logaritmo, todas sus derivadas y la transformada
de Legendre coinciden exactamente con los de la partición madre.
Además, el flujo conmuta con la suma de hijos:
\[
 \sum_c d_{j,c}(s)=d_j(s),\qquad
 d_{j,c}(s)=r_{j,c}d_j(s).
\]
\end{theorem}
\begin{proof}
Por agrupación de la suma finita,
\[
 Z_{\rm ref}(s)=\sum_{j,c}r_{j,c}d_j e^{su_j}
 =\sum_j d_j e^{su_j}=Z(s).
\]
Todas las conclusiones se deducen de esta identidad y de dividir cada
peso por el mismo \(Z(s)\). El mismo argumento vale después de cualquier
número finito de subdivisiones. Para una subdivisión numerable, la suma de
términos positivos converge al mismo valor mediante su agrupación por padres.
\end{proof}

La hipótesis de dirección heredada tiene contenido verificable. Si un
intervalo con dirección \(a\) se divide en dos mitades de direcciones
\(a+b\) y \(a-b\), su contribución cambia de
\(d e^{sa}\) a \(d e^{sa}\cosh(sb)\). Para \(b\ne0\) la igualdad
de funciones de partición falla aunque el peso total en \(s=0\) sea
el mismo. Así, la conservación geométrica y variacional exige conservar
también la información direccional pertinente. La identidad demostrada
se aplica a la subdivisión heredada; su transporte a una actualización
TPK que cambie esas direcciones se decide componiendo el lector de dicha
actualización.

Existe una formulación hamiltoniana regular de la coordenada dual. En el
plano cotangente con coordenadas \((q,p)\), tome
\(H(q,p)=\psi(p)\). Las ecuaciones son
\(\dot q=\psi'(p)\), \(\dot p=0\). Para velocidades
\(\dot q\in(u_-,u_+)\), la transformación de Legendre da
\[
 L(q,\dot q)=\psi^*(\dot q),\qquad
 \frac{\partial^2L}{\partial\dot q^2}>0.
\]
Este sistema realiza variacionalmente el potencial del flujo de intervalos.
Su ecuación \(\dot p=0\) muestra que el tiempo hamiltoniano es distinto
del parámetro \(s\) que recorre la familia \eqref{eq:convex-partition}.
La evaluación de \(H\) y \(L\) se conserva bajo la subdivisión del
teorema, porque \(\psi\) y \(\psi^*\) permanecen idénticas.

Por su parte, cualquier flujo \(\dot d=F(d)\) admite en una carta local
del simplex una elevación cotangente con
\(\mathscr H(d,p)=p\cdot F(d)\). Esta segunda construcción reproduce
el propio flujo \eqref{eq:replicator} en su base, pero su Hessiano respecto
de \(p\) es nulo. Las dos realizaciones responden a preguntas distintas:
una produce una dualidad convexa regular y la otra transporta una dinámica
de base dada. La distinción preserva el significado de Hamiltoniano y
Lagrangiano al comparar esta geometría con los lectores físicos del corpus.


% SOURCE sections/energy_refinement.tex
% Desarrollo acotado: lector energético intervalar del registro K.
% Estatuto: formalización añadida sobre K y sobre las identidades de Schur
% ya presentes en el corpus. No se atribuye novedad histórica a Schur.
\section{Formas de Dirichlet y equivalencia variacional del refinamiento}
\label{sec:energia-refinamiento-k}

La publicación positiva del registro dodecafásico define una partición
ordenada mediante
\[
 K=(234,543,140,729,659,824,621,58,914,794,146,601),
 \qquad d_j=K_j/6263,
 \qquad t_i=\sum_{j=1}^id_j.
\]
El registro se recibe de la composición
\(x_{\rm term}\mapsto\mathcal R_{12}\mapsto
(b^{90},b^{120},Q_{\rm TPK})\mapsto U_{\rm sgn}\mapsto K\),
posterior a APP--TRIT--TPK y a su estado enriquecido, conforme a los lectores
reconstruidos en \cite{hmtX}. Esta sección aplica
un lector cuadrático a la partición resultante. La comparación entre
profundidades se efectuará por restricciones a extremos heredados, de modo
que cada coeficiente de la energía conserve su origen en una separación
positiva ya producida.

El corpus contiene una energía ponderada de diferencias para el grafo de
las dos memorias y una reducción de Schur de las formas completas. Aquí se
reúnen esas operaciones sobre la red intervalar del registro. La
especialización a conductancias \(1/d_j\), el operador de interpolación
y su compatibilidad con la deformación de las separaciones constituyen
la formalización desarrollada en esta sección.

\subsection{Reconstrucción de las separaciones a partir de la matriz de rigidez}

Sea \(d=(d_1,\ldots,d_m)\), con \(d_j>0\), y sean
\(f=(f_0,\ldots,f_m)\in\mathbb C^{m+1}\) los valores de un observable
en los extremos de la partición. Definimos
\begin{equation}
 \mathcal E_d(f)=\sum_{j=1}^{m}\frac{|f_j-f_{j-1}|^2}{d_j},
 \qquad L_d=D^*\operatorname{diag}(d_1^{-1},\ldots,d_m^{-1})D,
 \label{eq:energia-intervalar}
\end{equation}
donde \(D:\mathbb C^{m+1}\to\mathbb C^m\) es el operador de
incidencia orientada, \((Df)_j=f_j-f_{j-1}\), y el asterisco denota la
adjunción hermítica. Así \(\mathcal E_d(f)=f^*L_df\).
La matriz de rigidez tiene entradas
\[
 (L_d)_{00}=d_1^{-1},\quad (L_d)_{mm}=d_m^{-1},\quad
 (L_d)_{ii}=d_i^{-1}+d_{i+1}^{-1}\quad(1\le i<m),
\]
\[
 (L_d)_{i-1,i}=(L_d)_{i,i-1}=-d_i^{-1},
\]
y las demás son nulas. Es hermítica positiva semidefinida, su núcleo
consiste en las funciones constantes y su forma es positiva definida
en el cociente por ese núcleo. La matriz completa recupera las
separaciones mediante
\begin{equation}
 d_j=-\frac1{(L_d)_{j-1,j}}.
 \label{eq:recuperacion-d-rigidez}
\end{equation}

\begin{proposition}[Equivalencia entre separaciones y lector nodal completo]
\label{prop:energia-bidireccional}
La aplicación \(d\mapsto L_d\) es una biyección entre
\((0,\infty)^m\) y el conjunto de matrices reales simétricas
tridiagonales cuyas entradas inmediatamente adyacentes a la diagonal son
negativas y cuya suma en cada fila es cero. Su inversa es
\eqref{eq:recuperacion-d-rigidez}. La normalización \(\sum_jd_j=1\)
equivale a \(\sum_{j=1}^m-1/(L_d)_{j-1,j}=1\).
\end{proposition}
\begin{proof}
Las fórmulas explícitas de \(L_d\) verifican todas las condiciones.
Recíprocamente, las entradas adyacentes negativas determinan separaciones
positivas únicas por \eqref{eq:recuperacion-d-rigidez}. La suma nula de
cada fila determina sus entradas diagonales como la suma de las
conductancias incidentes. La tridiagonalidad fija en cero las restantes
entradas. Se recupera exactamente la matriz de
\eqref{eq:energia-intervalar}; su semidefinitud positiva sigue entonces
de la factorización \(D^*\operatorname{diag}(1/d_j)D\).
\end{proof}

La derivada discreta ponderada
\(J_j=(f_j-f_{j-1})/d_j\) es el flujo orientado de la arista \(j\).
En los vértices interiores, \((L_df)_i=J_i-J_{i+1}\); en los extremos,
\((L_df)_0=-J_1\) y \((L_df)_m=J_m\). La ecuación interior
\(L_df=0\) expresa que el flujo tiene el mismo valor en cada arista.
Esta definición es la de un observable cuadrático del lector intervalar;
su realización como magnitud física requiere la carta dimensional pertinente.

\subsection{Interpolación armónica, complemento de Schur y reconstrucción del detalle}

Refinemos cada separación como
\(d_j=\sum_{c=1}^{r_j}\delta_{j,c}\), con \(\delta_{j,c}>0\).
Sea \(\widetilde d\) la concatenación de esas separaciones, \(R\) la
restricción a los extremos antiguos y \(H\) la interpolación definida por
\begin{equation}
 (Hf)_{j,c}=(1-\theta_{j,c})f_{j-1}+\theta_{j,c}f_j,
 \qquad
 \theta_{j,c}=\frac{\sum_{a=1}^c\delta_{j,a}}{d_j}.
 \label{eq:interpolacion-armonica-k}
\end{equation}
Los extremos compartidos se cuentan una sola vez. Se cumple \(RH=I\).

\begin{theorem}[Refinamiento isométrico de la energía]
\label{thm:energia-schur-refinamiento}
Para todo vector de valores antiguos \(f\),
\begin{equation}
 H^*L_{\widetilde d}H=L_d,
 \qquad L_{\widetilde d}H=R^*L_d,
 \qquad
 \min_{Rg=f}\mathcal E_{\widetilde d}(g)=\mathcal E_d(f).
 \label{eq:isometria-energia-k}
\end{equation}
El minimizador es único y vale \(Hf\). Todo \(g\) admite la
descomposición única
\begin{equation}
 g=H(Rg)+z,\qquad Rz=0,
 \qquad
 \mathcal E_{\widetilde d}(g)
 =\mathcal E_d(Rg)+\mathcal E_{\widetilde d}(z).
 \label{eq:detalle-energia-k}
\end{equation}
\end{theorem}
\begin{proof}
En una arista antigua escribamos \(h_c=g_{j,c}-g_{j,c-1}\), de modo
que \(\sum_ch_c=f_j-f_{j-1}=a\). Si
\(h_c=\delta_{j,c}a/d_j+r_c\), entonces \(\sum_cr_c=0\) y
\[
 \sum_c\frac{|h_c|^2}{\delta_{j,c}}
 =\frac{|a|^2}{d_j}+\sum_c\frac{|r_c|^2}{\delta_{j,c}}.
\]
La igualdad se obtiene expandiendo; el término cruzado es proporcional a
\(\sum_cr_c\) y se anula. El mínimo se alcanza exactamente cuando
todos los \(r_c\) son cero, condición que determina
\eqref{eq:interpolacion-armonica-k}. Al sumar las aristas se obtienen
el mínimo y la descomposición energética. El flujo de \(Hf\) es
constante dentro de cada arista antigua, por lo que sus componentes
interiores de \(L_{\widetilde d}Hf\) son cero y sus componentes heredadas
coinciden con \(L_df\). Esto prueba la segunda identidad matricial; al
multiplicar por \(H^*\), usando \(RH=I\), se obtiene la primera.
\end{proof}

El término \(\mathcal E_{\widetilde d}(z)\) es no negativo y se anula
exactamente para \(z=0\): un detalle de energía nula sería constante y,
al anularse en los extremos heredados, se anularía en toda la red.
La pareja \((Rg,z)\) conserva el vector refinado completo y permite
recuperarlo por \(g=H(Rg)+z\). La minimización selecciona su
componente armónica; el archivo del detalle \(z\) mantiene la
información que la sola forma reducida no distingue. La conservación
de la energía observada y la conservación del estado refinado quedan
así expresadas por mapas diferentes y explícitos.

Ordenando los vértices como extremos heredados \(E\) e interiores \(I\),
escribimos
\[
 L_{\widetilde d}=\begin{pmatrix}A&B^*\\B&C\end{pmatrix}.
\]
Para un refinamiento con vértices interiores, el bloque \(C\) es
positivo definido: una función nula en los extremos
que tenga energía cero es constante en cada subintervalo y, por tanto,
idénticamente nula. Se obtiene
\begin{equation}
 H=\begin{pmatrix}I\\-C^{-1}B\end{pmatrix},\qquad
 \operatorname{Schur}_{I}(L_{\widetilde d})
 =A-B^*C^{-1}B=L_d.
 \label{eq:schur-intervalar-k}
\end{equation}
Si el refinamiento conserva todos los vértices y no introduce ninguno,
la misma afirmación se reduce a \(H=I\) y \(L_{\widetilde d}=L_d\).

\begin{proposition}[Naturalidad de las prolongaciones y de la eliminación]
\label{prop:energia-naturalidad-sucesiva}
Para tres particiones anidadas, las interpolaciones armónicas y los
complementos de Schur satisfacen
\begin{equation}
 H_{n+2\leftarrow n}
 =H_{n+2\leftarrow n+1}H_{n+1\leftarrow n},
 \qquad
 \operatorname{Schur}_{n+1\to n}
 \bigl(\operatorname{Schur}_{n+2\to n+1}(L_{n+2})\bigr)=L_n.
 \label{eq:energia-naturalidad-sucesiva}
\end{equation}
\end{proposition}
\begin{proof}
La interpolación de valores afines sobre un intervalo sigue siendo la
misma función afín después de cualquier subdivisión de ese intervalo.
Esto prueba la primera identidad. Minimizar primero sobre los vértices
incorporados en el nivel último y después sobre los del nivel intermedio
es minimizar sobre todos ellos, con los mismos extremos heredados fijados.
La identidad variacional de \eqref{eq:isometria-energia-k} y la unicidad
del minimizador prueban la segunda identidad. La restricción del mínimo
se realiza, por tanto, mediante la misma composición que la prolongación
armónica de los datos.
\end{proof}

\subsection{Respuesta de Dirichlet a Neumann y defecto energético de una extensión}

Si sólo se conservan los extremos inicial y final, sea
\(D_0=\sum_{j=1}^md_j\). Para valores \((a,b)\), la solución
armónica es
\[
 f_i=a+(b-a)\frac{\sum_{j=1}^id_j}{D_0},\qquad
 J_j=(b-a)/D_0.
\]
Su operador de Dirichlet a Neumann, que transforma valores de frontera en
flujos salientes, es
\begin{equation}
 \Lambda_d=\frac1{D_0}
 \begin{pmatrix}1&-1\\-1&1\end{pmatrix},\qquad
 \min\mathcal E_d=\frac{|b-a|^2}{D_0}.
 \label{eq:dtn-intervalar-k}
\end{equation}
Cuando se prescribe el valor en todos los vértices heredados, su operador
de respuesta multipuntual es, por el mismo argumento, \(\Lambda=L_d\).

Una extensión aproximada \(g=(f,y)\), donde \(y\) reúne los valores
en los vértices interiores, tiene residual interior
\(r=Bf+Cy\). La identidad exacta
\begin{equation}
 \mathcal E_{\widetilde d}(g)-f^*\Lambda f
 =r^*C^{-1}r\ge0
 \label{eq:residual-dtn-k}
\end{equation}
se deduce de \(y+C^{-1}Bf=C^{-1}r\) y de la descomposición
\[
 g^*L_{\widetilde d}g
 =f^*(A-B^*C^{-1}B)f
 +(y+C^{-1}Bf)^*C(y+C^{-1}Bf).
\]
Así, el valor efectivo y el detalle mantienen una descomposición
hermítica completa. En particular, si una
extensión lineal es \(\widetilde H=\begin{pmatrix}I\\Z\end{pmatrix}\), su matriz de
residuales es \(\mathscr R=B+CZ\) y
\[
 \widetilde H^*L_{\widetilde d}\widetilde H-\Lambda
 =\mathscr R^*C^{-1}\mathscr R\succeq0.
\]
El residual nulo caracteriza la interpolación armónica exacta. Este
residual energético vive en el espacio de vértices interiores; la
cancelación de residuos de formas meromorfas pertenece a otro mapa de
frontera y conserva su definición propia.

\subsection{Conmutación del refinamiento y del flujo de separaciones}

Sea \(u=P_3P_{11}K\) la dirección algebraica ya recuperada de los
lectores del registro \cite{hmtX}. Para el parámetro real \(s\), definamos
\[
 d_j(s)=\frac{K_je^{su_j}}{\sum_\ell K_\ell e^{su_\ell}}.
\]
Fijadas fracciones positivas \(\vartheta_{j,c}\) de suma uno para cada
\(j\), el refinamiento heredado es
\[
 \delta_{j,c}(s)=\vartheta_{j,c}d_j(s),\qquad
 u_{j,c}=u_j.
\]
Producir primero los pesos refinados \(K_j\vartheta_{j,c}\) y
deformarlos da la misma familia: los hijos tienen idéntico exponencial y
sus coeficientes suman \(K_j\). Para todo \(s\),
\begin{equation}
 \operatorname{Schur}_{I}(L_{\widetilde d(s)})=L_{d(s)},
 \qquad H^*L_{\widetilde d(s)}H=L_{d(s)}.
 \label{eq:schur-flujo-k}
\end{equation}
En esta comparación padre--hijos, \(H\) es constante: sus coeficientes
son las sumas parciales de \(\vartheta_{j,c}\). En cambio, la
interpolación desde los dos extremos globales usa
\(t_i(s)=\sum_{j\le i}d_j(s)\) y cambia con el parámetro. Ambas
afirmaciones corresponden a restricciones diferentes de la misma red.

La normalización \(\sum_jd_j(s)=1\) hace que el operador global de dos
extremos permanezca igual a
\(\left(\begin{smallmatrix}1&-1\\-1&1\end{smallmatrix}\right)\).
La matriz de todos los nodos \(L_{d(s)}\) sí conserva la deformación y
la recupera mediante \eqref{eq:recuperacion-d-rigidez}. La pérdida de
información en el lector de dos extremos queda identificada exactamente:
retiene la longitud total y sus flujos, mientras el lector nodal retiene
todas las separaciones. El archivo de detalles proporciona el complemento
necesario para reconstruir el observable refinado.


% SOURCE sections/energy_limit.tex
\section{Reconstrucción analítica de la energía y memoria del refinamiento}
\label{sec:limite-energetico-k}

Las separaciones positivas construidas a partir del registro dodecafásico
determinan una partición ordenada y, por la sección anterior, una forma
de Dirichlet sobre sus valores nodales. Su refinamiento posee dos
operaciones complementarias: la interpolación armónica conserva la energía
de los datos heredados, mientras cada detalle registra la variación
incorporada entre esos datos. El paso a profundidad ilimitada consiste en
reunir esta familia compatible mediante una norma que controle la suma de
sus detalles. El resultado concierne a la realización analítica
unidimensional de la energía; conserva como antecedente la construcción
HMT del registro y de la estructura discreta conjunta del continuo.

\subsection{Particiones anidadas y espacios de interpolación}

Fijemos la partición inicial
\(P_0=\{0=t_0<t_1<\cdots<t_{12}=1\}\), cuyas longitudes son
\(d_j=t_j-t_{j-1}=K_j/\sum_iK_i>0\). Sea
\((P_n)_{n\geq0}\) una sucesión de particiones finitas anidadas de
\([0,1]\), obtenidas por subdivisión de los intervalos anteriores.
Escribiremos
\[
 |P_n|=\max\{b-a:[a,b]\text{ es un intervalo de }P_n\},
 \qquad \lim_{n\to\infty}|P_n|=0.
\]
La subdivisión reiterada de cada intervalo en un número fijo de hijos
iguales, al menos dos, satisface esta condición. Para otras reglas, la
contracción de la malla es una hipótesis explícita de este lector.
Las etiquetas de los hijos y sus datos de acarreo acompañan la
subdivisión y conservan la procedencia del refinamiento.

Sea \(V_n\) el espacio complejo de funciones continuas sobre \([0,1]\)
que son afines en cada intervalo de \(P_n\). La energía y la norma
energética se definen por
\begin{equation}
 \mathcal E_n(f)
 =\sum_{[a,b]\in P_n}\frac{|f(b)-f(a)|^2}{b-a}
 =\int_0^1|f'(x)|^2\,dx,
 \qquad
 \|f\|_{\mathcal E}^2=|f(0)|^2+\mathcal E_n(f).
 \label{eq:energia-afines-k}
\end{equation}
La igualdad integral resulta de que la derivada de una función afín es
constante en cada intervalo. La energía coincide así exactamente con
la forma nodal ya construida. La inclusión \(V_n\subset V_{n+1}\)
interpreta la misma función afín en una partición más fina; conserva
tanto \(f(0)\) como la energía y es isométrica para
\(\|\cdot\|_{\mathcal E}\).

El valor inicial tiene una función precisa: fija la componente constante
que la forma de Dirichlet identifica en su núcleo. La energía controla
las diferencias, y el par formado por el valor inicial y esas diferencias
determina el observable. En el espacio de Sobolev
\(H^1([0,1];\mathbb C)\), cuyos elementos tienen representante
absolutamente continuo con derivada débil en \(L^2\), esta norma es
equivalente a la norma usual. En efecto,
\(f(x)=f(0)+\int_0^xf'(t)\,dt\) controla la norma de \(f\) a
partir de \(|f(0)|\) y \(\|f'\|_{L^2}\); la integración de la
misma identidad controla a su vez \(|f(0)|\) mediante
\(\|f\|_{L^2}+\|f'\|_{L^2}\).

\subsection{Reconstrucción de una función a partir de sus valores compatibles}

\begin{theorem}[Reconstrucción energética y completación de las interpolaciones]
\label{thm:limite-energia-h1-k}
Sea \(f_n\in V_n\) una familia compatible por restricción a los
nodos heredados:
\[
 f_{n+1}(a)=f_n(a)
 \quad\text{para todo nodo }a\text{ de }P_n.
\]
Si \(\sup_n\mathcal E_n(f_n)<\infty\), existe una única función
\(f\in H^1([0,1];\mathbb C)\), considerada en su representante
continuo, cuyos valores en todos esos nodos son los prescritos. Además,
\(f_n\to f\) uniformemente y en \(H^1\), y
\begin{equation}
 \int_0^1|f'(x)|^2\,dx
 =\lim_{n\to\infty}\mathcal E_n(f_n).
 \label{eq:limite-energia-h1-k}
\end{equation}
Recíprocamente, la interpolación nodal de cualquier elemento de \(H^1\)
produce una familia de esta clase. Por tanto, la completación de
\(\bigcup_n V_n\) en la norma energética es \(H^1([0,1];\mathbb C)\).
\end{theorem}

\begin{proof}
Pongamos \(g_n=f_n'\in L^2([0,1];\mathbb C)\). Para
\(m\geq n\), la compatibilidad de los valores en los extremos de cada
intervalo \([a,b]\) de \(P_n\) proporciona
\[
 \frac1{b-a}\int_a^bg_m(x)\,dx
 =\frac{f_m(b)-f_m(a)}{b-a}
 =\frac{f_n(b)-f_n(a)}{b-a}
 =g_n\big|_{(a,b)}.
\]
Sea \(W_n\) el espacio de funciones de \(L^2\) constantes en cada
intervalo de \(P_n\), y sea \(Q_n:L^2\to W_n\) la proyección
ortogonal dada por los promedios sobre dichos intervalos. La identidad
anterior afirma \(Q_ng_m=g_n\). En consecuencia,
\(g_m-g_n\) es ortogonal a \(g_n\) y
\begin{equation}
 \|g_m-g_n\|_{L^2}^2
 =\mathcal E_m(f_m)-\mathcal E_n(f_n).
 \label{eq:pitagoras-gradientes-k}
\end{equation}
Las energías son crecientes y acotadas. La identidad de Pitágoras
prueba que \((g_n)\) es una sucesión de Cauchy en \(L^2\), con
límite \(g\). Definamos
\[
 f(x)=f_0(0)+\int_0^xg(t)\,dt.
\]
Esta función pertenece a \(H^1\), tiene derivada débil \(g\) y
satisface, por la desigualdad de Cauchy--Schwarz,
\begin{equation}
 \sup_{x\in[0,1]}|f_n(x)-f(x)|
 \leq\|g_n-g\|_{L^2}\longrightarrow0.
 \label{eq:convergencia-uniforme-energia-k}
\end{equation}
La convergencia uniforme conserva cada valor nodal prescrito; junto con
la convergencia de las derivadas en \(L^2\), da la convergencia en
\(H^1\). La convergencia de las normas de las derivadas prueba
\eqref{eq:limite-energia-h1-k}. La unión de los nodos es densa porque
\(|P_n|\to0\); dos representantes continuos con los mismos valores
en esa unión coinciden. Se obtiene la unicidad.

Para la recíproca, tomemos \(f\in H^1\) en su representante
continuo y sea \(f_n\) su interpolación en los nodos de \(P_n\).
Su derivada es \(g_n=Q_nf'\). Las proyecciones \(Q_n\) son
contractivas en \(L^2\). Para una función continua \(v\), el
error \(|Q_nv-v|\) está acotado por su módulo de continuidad
evaluado en \(|P_n|\), que tiende a cero. La densidad de las
funciones continuas en \(L^2\) y la contractividad implican
\(Q_ng\to g\) para todo \(g\in L^2\): dado un aproximante
continuo \(v\), se emplea
\[
 \|Q_ng-g\|_{L^2}
 \leq 2\|g-v\|_{L^2}+\|Q_nv-v\|_{L^2}.
\]
Aplicando esta convergencia a \(g=f'\), se obtienen la convergencia
de las derivadas y, mediante \eqref{eq:convergencia-uniforme-energia-k},
las de las funciones. La contractividad acota sus energías por
\(\|f'\|_{L^2}^2\), y la convergencia de las normas prueba de
nuevo la identidad del límite. Toda función de \(H^1\) es, por
tanto, límite energético de elementos de \(\bigcup_nV_n\).
La equivalencia de normas y la completitud de \(H^1\) identifican
la completación enunciada.
\end{proof}

El argumento describe concretamente qué se conserva al pasar de la red
a la función. Cada derivada finita registra el promedio de todas las
derivadas posteriores en su intervalo; al refinar, se incorporan
componentes ortogonales que mantienen ese promedio. La energía acotada
controla la suma total de esas incorporaciones. El valor continuo se
recupera entonces integrando la derivada límite y conservando el valor
inicial. La realización analítica reúne, por esta vía explícita, la
información compatible de los niveles finitos.

\subsection{Descomposición ortogonal de la memoria por escalas}

Sea \(H_n\) la interpolación armónica de los valores de \(P_n\)
sobre los nodos de \(P_{n+1}\), interpretada como función de
\(V_{n+1}\), y definamos el detalle
\[
 z_{n+1}=f_{n+1}-H_nf_n.
\]
Por \eqref{eq:detalle-energia-k}, el detalle se anula en los nodos
heredados. Su derivada tiene promedio cero en cada intervalo antiguo
y es ortogonal a \(W_n\). Para niveles distintos, la derivada de
un detalle anterior pertenece a ese espacio de funciones constantes;
las derivadas de los detalles son, por tanto, ortogonales entre sí.
La misma ortogonalidad vale frente a \(f_0'\).

\begin{proposition}[Identidad de energía y recuperación de los detalles]
\label{prop:memoria-energetica-escalas-k}
Para toda profundidad finita \(N\),
\begin{equation}
 \mathcal E_N(f_N)
 =\mathcal E_0(f_0)
 +\sum_{n=0}^{N-1}\mathcal E_{n+1}(z_{n+1}).
 \label{eq:energia-telescopica-detalles-k}
\end{equation}
Si la familia satisface las hipótesis del
Teorema~\ref{thm:limite-energia-h1-k}, la serie de la derecha
converge y su suma es \(\int_0^1|f'|^2\). Los datos \(f_0\)
y todos los detalles determinan la familia compatible completa y su
realización analítica.
\end{proposition}
\begin{proof}
La identidad finita resulta de sumar
\(\mathcal E_{n+1}(f_{n+1})
=\mathcal E_n(f_n)+\mathcal E_{n+1}(z_{n+1})\), que es la
descomposición ortogonal de la sección anterior. Los sumandos son
no negativos y las sumas parciales convergen por
\eqref{eq:limite-energia-h1-k}. La reconstrucción finita se efectúa
recursivamente mediante \(f_{n+1}=H_nf_n+z_{n+1}\); el teorema
precedente determina después la función límite. Recíprocamente, una
función de \(H^1\) determina sus interpolaciones y, por diferencia,
todos estos detalles. Cada flecha conserva así sus datos propios.
\end{proof}

La eliminación variacional de nodos interiores selecciona el detalle
nulo en el nivel eliminado. La conservación de los detalles permite
recuperar cualquier configuración compatible de energía finita. La
respuesta efectiva y la configuración completa poseen, de esta manera,
una relación precisa: la primera registra el mínimo sobre las variables
interiores; la segunda añade las componentes ortogonales y su energía.

\subsection{Transporte de la realización analítica bajo el flujo positivo}

La dirección \(u=P_3P_{11}K\), recuperada antes de la evaluación
energética \cite{hmtX}, determina la familia de separaciones de
\eqref{eq:schur-flujo-k}. Su parámetro \(s\in\mathbb R\)
describe esta deformación analítica. Escribamos
\[
 d_j(s)=\frac{K_je^{su_j}}{\sum_\ell K_\ell e^{su_\ell}},
 \qquad
 t_i(s)=\sum_{j=1}^id_j(s),
 \qquad
 a_j(s)=\frac{d_j(s)}{d_j(0)}>0.
\]
Sea \(F_s:[0,1]\to[0,1]\) la aplicación que envía cada
\(t_i(0)\) a \(t_i(s)\) y es afín en los intervalos iniciales.
Es un homeomorfismo creciente, de pendiente \(a_j(s)\) en el
intervalo \([t_{j-1}(0),t_j(0)]\). Su inversa es también afín
por intervalos y ambas aplicaciones son Lipschitz.

Supongamos que cada subdivisión hereda las fracciones de longitud y la
dirección de su intervalo padre, como en la sección anterior. Todos
los nodos refinados se transportan entonces mediante el mismo \(F_s\).
Denotemos sus particiones por \(P_n(s)=F_s(P_n(0))\) y sus
espacios de interpolación por \(V_n(s)\). La aplicación
\[
 U_s f=f\circ F_s^{-1}
\]
transporta \(V_n(0)\) a \(V_n(s)\), preserva los valores
nodales correspondientes y prolonga esta acción a \(H^1\).

\begin{theorem}[Equivalencia energética y conmutación con el refinamiento]
\label{thm:transporte-h1-flujo-k}
Para todo \(f\in H^1([0,1];\mathbb C)\),
\begin{equation}
 \int_0^1|(U_sf)'(y)|^2\,dy
 =\sum_{j=1}^{12}\frac1{a_j(s)}
 \int_{t_{j-1}(0)}^{t_j(0)}|f'(x)|^2\,dx.
 \label{eq:transporte-energia-h1-k}
\end{equation}
Para \(s\) en cualquier intervalo compacto, las normas energéticas
antes y después del transporte son uniformemente equivalentes y
\(|P_n(s)|\to0\) uniformemente en ese compacto. La interpolación
armónica y la restricción a nodos heredados conmutan con \(U_s\);
la reconstrucción y la minimización de los detalles se transportan
tanto a profundidad finita como en la realización \(H^1\).
\end{theorem}

\begin{proof}
En el intervalo inicial \(j\), la regla de la cadena y el cambio
de variable \(y=F_s(x)\) proporcionan
\((U_sf)'(F_s(x))=f'(x)/a_j(s)\) y \(dy=a_j(s)\,dx\).
La integración y la suma sobre los intervalos prueban
\eqref{eq:transporte-energia-h1-k}. La fórmula vale para los
representantes absolutamente continuos, y su lado derecho es finito.

Fijado un intervalo compacto de parámetros, la continuidad y positividad
de los doce factores dan constantes \(0<a_-\leq a_+<\infty\)
tales que \(a_-\leq a_j(s)\leq a_+\) para todo \(j\) y
todo \(s\) de ese compacto. Por tanto,
\[
 \frac1{a_+}\mathcal E(f)
 \leq\mathcal E(U_sf)
 \leq\frac1{a_-}\mathcal E(f),
 \qquad
 |P_n(s)|\leq a_+|P_n(0)|.
\]
Como \((U_sf)(0)=f(0)\), se obtiene además
\[
 \min\{1,a_+^{-1}\}\|f\|_{\mathcal E}^2
 \leq\|U_sf\|_{\mathcal E}^2
 \leq\max\{1,a_-^{-1}\}\|f\|_{\mathcal E}^2.
\]
Estas cotas prueban la equivalencia uniforme y la conservación de la
condición de malla decreciente.

En cada arista padre, los coeficientes de interpolación son las sumas
parciales de las fracciones heredadas, independientes de \(s\).
Si \(H_n(s)\) denota su interpolación, se tiene
\(U_sH_n(0)=H_n(s)U_s\) en los espacios correspondientes.
La restricción conmuta porque los nodos se transportan con sus valores.
La descomposición en función armónica y detalle se conserva, y la
identidad de Schur de \eqref{eq:schur-flujo-k} preserva su
caracterización variacional. Finalmente, la continuidad acotada de
\(U_s\) permite pasar al límite \(H^1\). En particular, una
familia reconstruida antes del transporte converge después a \(U_sf\),
con los mismos valores nodales transportados.
\end{proof}

La ley \eqref{eq:transporte-energia-h1-k} identifica los factores
métricos exactos de la deformación: el cambio de longitud de cada
intervalo modifica inversamente su contribución energética. La
equivalencia de normas conserva las configuraciones de energía finita,
y el archivo de detalles conserva su reconstrucción. Se obtiene así una
realización analítica completa de este lector del registro positivo.
Su compatibilidad con las formas canónicas se articula sobre las
subdivisiones de un dominio fijo: la suma de formas conserva el objeto
meromorfo, mientras la reducción variacional y los detalles conservan
la respuesta energética y el observable refinado. La identificación
con formas amplituhedrales de rango superior conserva las condiciones
específicas de cobertura, grado y correspondencia de residuos
establecidas para aquella realización.


% SOURCE sections/nonadic_transport.tex
\section{Transportes nonádicos y conservación de la memoria}
\label{sec:nonadic}

La reconstrucción de la sección~\ref{sec:memory} proporciona un enlace
efectivo entre la memoria del transporte y la geometría del registro
dodecafásico. El registro generado por APP--TRIT--TPK determina sus
lecturas de incidencia; las dos compresiones conservan, en componentes
separadas, la publicación terminal y sus complementos; la reunión de
las memorias con la carga restituye las doce coordenadas. La dirección
geométrica y el proyector dimensional se obtienen después de esa
restitución. El desarrollo siguiente establece el balance operatorio
que sostiene esta cadena y su comportamiento bajo refinamiento,
realización espectral y cambio de métrica
\cite{HMT-VIII,hmtX}.

\subsection{Calendario de nueve fases y descomposición ortogonal}

Sea \(\mathcal H\) un espacio de Hilbert complejo y sea
\(C\in\mathcal B(\mathcal H)\). La suma ortogonal de nueve copias,
\(\mathcal H^{\oplus9}\), conserva por separado las nueve posiciones
del calendario. Definimos
\begin{equation}
 \begin{split}
 S_C(v_0,\ldots,v_8)&=(Cv_8,v_0,\ldots,v_7),\\
 J:\mathcal H\longrightarrow\mathcal H^{\oplus9},
 \qquad Jv&=\frac13(v,\ldots,v).
 \end{split}
 \label{nt:calendar}
\end{equation}
La normalización \(1/3=9^{-1/2}\) da \(J^*J=I\). El operador
\(P=JJ^*\) es la proyección ortogonal sobre la sección uniforme;
su acción sustituye las nueve componentes por su media. Cada
componente atraviesa una vez la unión que contiene \(C\) durante una
vuelta, de modo que
\begin{equation}
 S_C^9=I_9\otimes C.
 \label{nt:full-return}
\end{equation}
El retorno de la posición del calendario realiza así una acción de
\(C\) sobre la fibra de memoria.

\begin{definition}
La compresión uniforme del transporte y su complemento son
\begin{equation}
 T_C=J^*S_CJ=\frac{8I+C}{9},
 \qquad
 \eta_C=(I-P)S_CJ.
 \label{nt:compression}
\end{equation}
La primera aplicación tiene codominio \(\mathcal H\); la segunda
tiene codominio \(J(\mathcal H)^\perp\subset\mathcal H^{\oplus9}\).
\end{definition}

Para \(z=(C-I)v\), el complemento se escribe
\begin{equation}
 \eta_Cv=\frac1{27}(8z,-z,-z,-z,-z,-z,-z,-z,-z).
 \label{nt:complement}
\end{equation}
La suma de sus componentes es cero. Los coeficientes de la
descomposición proceden, por tanto, del promedio de ocho
continuaciones idénticas y una continuación transportada.

\begin{proposition}[Balance de la compresión y de la unión de memoria]
Para todo \(C\in\mathcal B(\mathcal H)\),
\begin{align}
 \eta_C^*\eta_C
   &=\frac8{81}(I-C)^*(I-C),\label{nt:eta-gram}\\
 I-T_C^*T_C
   &=\eta_C^*\eta_C+\frac19(I-C^*C).\label{nt:general-balance}
\end{align}
Si \(C^*C=I\), la aplicación
\(v\mapsto(T_Cv,\eta_Cv)\) es isométrica.
\end{proposition}

\begin{proof}
La expresión~\eqref{nt:complement} da
\[
 \|\eta_Cv\|^2
 =\frac{8^2+8}{27^2}\|(C-I)v\|^2
 =\frac8{81}\|(C-I)v\|^2.
\]
La polarización prueba~\eqref{nt:eta-gram}. Por otra parte,
\[
 T_C^*T_C
 =\frac{64I+8C+8C^*+C^*C}{81}.
\]
Sumando \eqref{nt:eta-gram} se obtiene
\[
 T_C^*T_C+\eta_C^*\eta_C=\frac{8I+C^*C}{9}.
\]
Restar este resultado de \(I\) da~\eqref{nt:general-balance}.
Cuando \(C^*C=I\), la suma de los dos productos de Gram es la
identidad, lo que conserva todos los productos interiores.
\end{proof}

El término \(\eta_C^*\eta_C\) mide la dispersión entre las nueve
componentes. El término \((I-C^*C)/9\) registra, por separado,
la variación de la norma en la unión de memoria. La realización
unitaria sitúa toda la variación visible en el primer término y
conserva la norma en la suma ortogonal completa.

\begin{proposition}[Conservación de un observable]
Sean \(C\) unitario y \(A=A^*\in\mathcal B(\mathcal H)\) tales que
\(C^*AC=A\). Entonces
\begin{equation}
 A=T_C^*AT_C+
   \eta_C^*(I_9\otimes A)\eta_C
  =T_C^*AT_C+\frac8{81}(I-C)^*A(I-C).
 \label{nt:observable-balance}
\end{equation}
Para \(A\succeq0\), ambos términos de la suma son positivos.
\end{proposition}

\begin{proof}
La fórmula del complemento, aplicada al producto sesquilineal
definido por \(A\), prueba la segunda igualdad. La expansión de
los dos términos da
\[
 \frac1{81}\bigl(
 64A+8AC+8C^*A+C^*AC
 +8A-8AC-8C^*A+8C^*AC\bigr)
 =\frac{8A+C^*AC}{9}=A.
\]
La positividad se conserva bajo una aplicación de la forma
\(B\mapsto L^*BL\).
\end{proof}

\subsection{Iteración de la compresión y conservación del sector fijo}

Fijamos en esta subsección un operador unitario \(C\) y escribimos
\(T=T_C\), \(\eta=\eta_C\). El transporte \(S_C^N\) mantiene todas
las fases durante \(N\) pasos. La potencia \(T^N\), en cambio,
aplica sucesivamente la compresión uniforme; en cada etapa
\(\eta T^jv\) conserva el complemento de la publicación anterior.
La distinción queda expresada por los siguientes operadores.

\begin{proposition}[Registro isométrico a horizonte finito]
Para todo entero \(N\geq1\), la aplicación
\begin{equation}
 V_Nv=(T^Nv,\eta v,\eta Tv,\ldots,\eta T^{N-1}v)
 \label{nt:finite-record}
\end{equation}
de \(\mathcal H\) a
\(\mathcal H\oplus(\mathcal H^{\oplus9})^{\oplus N}\)
es isométrica. Si \(A\) satisface las hipótesis de
\eqref{nt:observable-balance}, entonces
\begin{equation}
 A=T^{*N}AT^N+
 \sum_{j=0}^{N-1}T^{*j}\eta^*(I_9\otimes A)\eta T^j.
 \label{nt:telescopy}
\end{equation}
\end{proposition}

\begin{proof}
Multiplicar~\eqref{nt:observable-balance} por \(T^{*j}\) y
\(T^j\) da
\[
 T^{*j}AT^j-T^{*(j+1)}AT^{j+1}
 =T^{*j}\eta^*(I_9\otimes A)\eta T^j.
\]
La suma telescópica prueba~\eqref{nt:telescopy}. Para \(A=I\),
el miembro derecho es el producto de Gram de \(V_N\).
El paso de \(N\) a \(N+1\) sustituye únicamente \(T^Nv\) por
\((T^{N+1}v,\eta T^Nv)\), conservando todas las componentes
de memoria previamente producidas.
\end{proof}

\begin{theorem}[Límite de la compresión y memoria ilimitada]
Sea \(P_{\mathrm{fix}}\) la proyección ortogonal sobre
\(\ker(I-C)\). Entonces
\begin{align}
 \operatorname*{s\!-\!lim}_{N\to\infty}T^N
   &=P_{\mathrm{fix}},\label{nt:strong-limit}\\
 I&=P_{\mathrm{fix}}+
 \sum_{j=0}^{\infty}T^{*j}\eta^*\eta T^j.\label{nt:infinite-record}
\end{align}
La serie converge fuertemente. Para todo observable acotado
autoadjunto que satisfaga \(C^*AC=A\), se conserva además
\begin{equation}
 A=P_{\mathrm{fix}}AP_{\mathrm{fix}}+
 \sum_{j=0}^{\infty}
 T^{*j}\eta^*(I_9\otimes A)\eta T^j.
 \label{nt:infinite-observable}
\end{equation}
\end{theorem}

\begin{proof}
El teorema espectral del operador unitario realiza \(C\) como
la función coordenada \(z\) sobre la circunferencia unidad.
La compresión corresponde a \(t(z)=(8+z)/9\), y
\begin{equation}
 1-|t(z)|^2=\frac8{81}|1-z|^2,
 \qquad |z|=1.
 \label{nt:circular-defect}
\end{equation}
Así, \(t(z)^N\) converge a cero en cada \(z\ne1\) y conserva
el valor uno en \(z=1\). La cota \(|t(z)^N|\leq1\) permite
aplicar convergencia dominada a cada medida espectral vectorial.
Se obtiene~\eqref{nt:strong-limit} y, por el mismo argumento,
\(T^{*N}\to P_{\mathrm{fix}}\) fuertemente. La acotación de
\(A\), junto con \(\|T^N\|\leq1\), da
\(T^{*N}AT^N\to P_{\mathrm{fix}}AP_{\mathrm{fix}}\).
El límite de~\eqref{nt:telescopy} prueba las dos identidades.
\end{proof}

El registro
\[
 V_\infty v=(P_{\mathrm{fix}}v,\eta v,\eta Tv,\eta T^2v,\ldots)
\]
es, por consiguiente, una isometría. El sector fijo conserva la
parte que persiste en la publicación terminal; las demás
componentes permanecen en la memoria de los complementos.

Para el desplazamiento bilateral
\(C e_n=e_{n+1}\) en \(\ell^2(\mathbb Z)\), una sucesión fija
es constante. La condición de cuadrado-sumabilidad fuerza su
anulación, por lo que \(P_{\mathrm{fix}}=0\). Toda la norma
queda entonces representada en los complementos. El espectro de
este desplazamiento es la circunferencia completa, y el máximo
espectral de \(|t(z)|^N\) es uno; de ahí
\(\|T^N\|=1\) para cada \(N\). La convergencia fuerte conserva
este comportamiento no uniforme.

Para una permutación cíclica de \(d\) posiciones, el sector fijo
es la recta de los vectores uniformes. Si la permutación tiene
varias órbitas, cada órbita aporta su propia dirección uniforme.
Esta dimensión determina la parte terminal que debe conservarse.

\begin{proposition}[Compresión de un transporte de orden tres]
Sea \(C\) unitario y \(C^3=I\). Entonces
\begin{equation}
 P_{\mathrm{fix}}=\frac{I+C+C^2}{3},
 \qquad
 T_C^*T_C=P_{\mathrm{fix}}+
 \frac{19}{27}(I-P_{\mathrm{fix}}).
 \label{nt:order-three}
\end{equation}
\end{proposition}

\begin{proof}
El promedio de las tres potencias es el proyector ortogonal
sobre los vectores fijos. Como \(C^*=C^2\),
\[
 T_C^*T_C=\frac{65I+8(C+C^2)}{81}
 =\frac{57I+24P_{\mathrm{fix}}}{81},
\]
que coincide con~\eqref{nt:order-three}.
\end{proof}

En la realización de doce posiciones de
\(\mathbb R^{12}\), \((Sx)_i=x_{i+1}\) con índices módulo
doce, \(S^3\) tiene orden cuatro y \(S^4\) tiene orden tres.
Las cuatro órbitas de \(S^4\) dan
\(\operatorname{rank}P_{\mathrm{fix}}=4\).
La realización excepcional \(g_c\) del artículo~X satisface
\(g_c^3=I\) y \(\ker(I-g_c)=0\) en su propio portador;
en ella~\eqref{nt:order-three} se reduce a
\(T_{g_c}^*T_{g_c}=19I/27\)
\cite{hmtX}. El coeficiente y el sector fijo se transportan
conjuntamente con el dominio.

La contracción radial \(M=qC\), \(0<q<1\), posee otro registro
exacto:
\begin{equation}
 L_Nv=\bigl(\sqrt{1-q^2}\,v,
 \sqrt{1-q^2}\,qCv,\ldots,
 \sqrt{1-q^2}\,q^{N-1}C^{N-1}v,q^NC^Nv\bigr).
 \label{nt:radial-record}
\end{equation}
Su producto de Gram es la identidad porque
\((1-q^2)\sum_{j=0}^{N-1}q^{2j}+q^{2N}=1\).
El terminal tiene norma \(q^N\|v\|\), y el registro infinito
conserva el producto interior. El factor \(q=5/9\) del modo
diferencial central del artículo~VIII se aplica a una vuelta
completa; su reparto uniforme entre nueve pasos utiliza
\(q^{1/9}S_C\), cuya novena potencia es
\(q(I_9\otimes C)\) \cite{HMT-VIII}.

\subsection{Reducción del registro y conservación de la traza}

La dirección fija del complemento en las nueve componentes
permite una realización mínima en dos sumandos. Definamos
\[
 D_C=\frac{\sqrt8}{9}(I-C).
\]
Las identidades anteriores dan
\(T_C^*T_C+D_C^*D_C=I\) para \(C\) unitario.
El cambio de signo respecto de la expresión de
\(\eta_C\) corresponde a una orientación de la componente
complementaria y conserva su producto de Gram.

\begin{proposition}[Canal positivo de la compresión con memoria]
Para \(C\) unitario, la aplicación
\begin{equation}
 \mathcal Q_C(X)=T_CXT_C^*+D_CXD_C^*
              =\frac{8X+CXC^*}{9}
 \label{nt:channel}
\end{equation}
es completamente positiva en \(\mathcal B(\mathcal H)\) y
conserva la identidad. Su restricción a los operadores de
clase traza conserva la traza.
\end{proposition}

\begin{proof}
Al expandir los dos términos de la primera expresión,
los productos \(CX\) y \(XC^*\) se cancelan. Los coeficientes
restantes son \(72/81=8/9\) para \(X\) y \(9/81=1/9\)
para \(CXC^*\), lo que prueba~\eqref{nt:channel}.
Para cualquier entero \(m\geq1\), la ampliación a matrices
de orden \(m\) tiene la forma
\[
 X\longmapsto
 (I_m\otimes T_C)X(I_m\otimes T_C)^*
 +(I_m\otimes D_C)X(I_m\otimes D_C)^*.
\]
Cada sumando conserva positividad; por ello la aplicación es
completamente positiva. La segunda expresión de
\eqref{nt:channel} da \(\mathcal Q_C(I)=I\).
Si \(X\) es de clase traza, la invariancia de la traza bajo
conjugación unitaria implica
\(\operatorname{Tr}\mathcal Q_C(X)=\operatorname{Tr}X\).
\end{proof}

La aplicación \eqref{nt:channel} se obtiene igualmente
preparando la isometría
\(v\mapsto T_Cv\otimes e_0+D_Cv\otimes e_1\) y efectuando
la traza parcial sobre \(\mathbb C^2\).
La conservación del vector completo pertenece a la isometría;
la conservación de la traza del estado reducido pertenece a
\(\mathcal Q_C\). Ambas afirmaciones proceden del mismo
balance ortogonal.

En la prolongación nonádica de álgebras finitas aparece una
conservación complementaria. Para
\(\mathcal A_d=M_{9^d}(\mathbb C)\), con
\(\tau_d(a)=9^{-d}\operatorname{Tr}a\), las aplicaciones
\[
 \iota_d(a)=a\otimes I_9,\qquad
 \jmath_{d,j}(a)=a\otimes p_j,
 \quad p_j=|e_j\rangle\langle e_j|
\]
satisfacen
\begin{align}
 \tau_{d+1}(\iota_d(a))&=\tau_d(a),
 &\iota_d(I)&=I,\label{nt:unital-trace}\\
 \tau_{d+1}(\jmath_{d,j}(a))&=\frac{\tau_d(a)}9,
 &\jmath_{d,j}(I)&=I\otimes p_j.\label{nt:corner-trace}
\end{align}
La demostración usa
\(\operatorname{Tr}(a\otimes b)=
\operatorname{Tr}(a)\operatorname{Tr}(b)\),
\(\operatorname{Tr}I_9=9\) y \(\operatorname{Tr}p_j=1\).
La primera inclusión conserva toda la fibra y la unidad;
la segunda conserva una continuación marcada y su soporte.

Esta distinción también se manifiesta en el árbol de
historias enriquecidas. Sean \(\mathcal H_n\) los prefijos
supervivientes de longitud \(n\), y sea \(d(h)\geq1\) el número de
emisiones salientes conservadas de \(h\). Cada emisión con multiplicidad
se representa por hijos marcados distintos; la suma siguiente recorre
ese conjunto de \(d(h)\) hijos.
La ley de pesos
\[
 \mu_{n+1}(h\epsilon)=\mu_n(h)\,p_h(\epsilon),
 \qquad p_h(\epsilon)=d(h)^{-1}
\]
conserva
\(\sum_{\epsilon}\mu_{n+1}(h\epsilon)=\mu_n(h)\).
Por tanto, en
\(\mathscr L_n=L^2(\mathcal H_n,\mu_n)\), el pullback
\[
 R_nf(h\epsilon)=f(h)
\]
es isométrico y conserva la función unidad. Su adjunto es
\[
 E_ng(h)=\sum_{\epsilon}p_h(\epsilon)g(h\epsilon),
 \qquad E_nR_n=I,\quad E_n1=1.
\]
En efecto, agrupar la suma de \(\|R_nf\|^2\) por padre
produce \(\sum_h\mu_n(h)|f(h)|^2\); el mismo agrupamiento
en el producto interior determina \(E_n\).
La memoria y las multiplicidades quedan presentes en el
dominio de las historias. Estos mapas constituyen la
realización funcional del refinamiento descrito en el
artículo~X \cite{hmtX}.

\subsection{Carta espectral y reciprocidad del radio}

El transporte anterior admite realizaciones espectrales
con dominios propios. La carta de Möbius de VIII actúa,
en coordenadas finitas, por
\begin{equation}
 \tau(z)=\frac{z-1}{z},
 \qquad z\in\mathbb C\setminus\{0\}.
 \label{nt:mobius}
\end{equation}
La inversa es \(z=(1-\tau)^{-1}\), para \(\tau\ne1\).
La extensión a la esfera de Riemann envía \(0\) a
\(\infty\) e \(\infty\) a \(1\).

\begin{proposition}[Espejo espectral y radio recíproco]
Sea \(z^\sharp=1-\overline z\). Para
\(z\notin\{0,1\}\),
\begin{equation}
 \tau(z^\sharp)=\frac1{\overline{\tau(z)}}.
 \label{nt:mirror}
\end{equation}
Si \(z=\sigma+it\), entonces
\begin{equation}
 |\tau(z)|^2
 =\frac{(\sigma-1)^2+t^2}{\sigma^2+t^2};
 \qquad
 |\tau(z)|=1\ \Longleftrightarrow\ \sigma=\frac12.
 \label{nt:critical-circle}
\end{equation}
\end{proposition}

\begin{proof}
La sustitución directa da
\[
 \tau(1-\overline z)
 =\frac{-\overline z}{1-\overline z}
 =\frac{\overline z}{\overline z-1}
 =\frac1{\overline{(z-1)/z}}.
\]
El cociente de módulos cuadrados produce
\eqref{nt:critical-circle}; la igualdad entre numerador y
denominador equivale a \(-2\sigma+1=0\).
\end{proof}

Escribiendo \(\tau(z)=r e^{i\theta}\), con \(r>0\), el
espejo preserva \(\theta\) y transforma \(r\) en \(r^{-1}\).
Su conjunto fijo en esta carta es la circunferencia unidad.
La conjugación \(z\mapsto\overline z\), por su parte,
produce \(\tau\mapsto\overline\tau\) y cambia la orientación
angular. Las dos involuciones conservan, respectivamente,
la coordenada angular y la coordenada radial.

Sobre \(z=1/2+i\gamma\), con \(\gamma\in\mathbb R\),
\begin{equation}
 \tau_\gamma=\frac{\gamma+i/2}{\gamma-i/2},
 \qquad
 \gamma=\frac{i}{2}\frac{\tau_\gamma+1}{\tau_\gamma-1}
       =\frac12\cot\frac{\theta}{2},
 \quad \tau_\gamma=e^{i\theta}.
 \label{nt:spectral-phase}
\end{equation}
La compactificación conserva la altura mediante la inversa;
el punto \(\tau=1\) corresponde al infinito.

Sea \(\mathcal D=C_c^\infty(\mathbb R;\mathbb C)\), con su topología
usual de funciones de prueba. Para la realización positiva de la forma
de Weil se utilizan expresamente su positividad semidefinida, su
invariancia por traslación y su continuidad en esa topología, en el
dominio demostrado por el lector espectral correspondiente. Si
\(\mathscr W(f,g)=\langle Ef,Eg\rangle\), su núcleo es
\(\mathcal N=\{f\in\mathcal D:\mathscr W(f,f)=0\}\), y se define
\[
 \mathcal H_{\mathscr W}
 =\overline{\mathcal D/\mathcal N}.
\]
Las traslaciones inducen un grupo unitario fuertemente
continuo \(U_t[f]=[f(\,\cdot-t)]\): la invariancia preserva la norma,
y la continuidad de traslaciones en \(\mathcal D\) y de la forma
prueba la continuidad fuerte sobre un subespacio denso y después en
toda la completación. Con la convención
\(U_t=e^{itH}\), su generador autoadjunto actúa sobre
las clases de pruebas por \(H[f]=[if']\). El Cayley de
esta realización es el operador acotado
\begin{equation}
 C_{\mathscr W}
 =(H+iI/2)(H-iI/2)^{-1}
 =I+i(H-iI/2)^{-1}.
 \label{nt:cayley-hilbert}
\end{equation}
El cálculo funcional lo hace unitario, pues el cociente
en~\eqref{nt:spectral-phase} tiene módulo uno para
\(\gamma\) real. La identidad de compresión se aplica
a \(C_{\mathscr W}\) en \(\mathcal H_{\mathscr W}\).
El espacio, el dominio del generador y la positividad
que permite su construcción son las premisas de esta
realización espectral \cite{HMT-VIII}.

La memoria bilateral actúa, en cambio, en
\(\ell^2(\mathbb Z)\), y las permutaciones \(S^3,S^4\)
actúan en \(\mathbb R^{12}\). Cada realización determina
su espectro, sus sectores fijos y la acción de una vuelta.
Los respectivos mapas de transporte especifican las
composiciones entre estos espacios.

\subsection{Cayley sobre pruebas y conservación de la forma aritmética}

La misma publicación aritmética admite un lector de Cayley
definido directamente sobre funciones de prueba. Fijemos
\(a>b>1/2\) y
\[
 \mathscr E_b=
 \left\{f\in C^\infty(\mathbb R):
   \sup_x e^{b|x|}|f^{(j)}(x)|<\infty
   \text{ para todo }j\geq0\right\}.
\]
Las funciones suaves de soporte compacto pertenecen a
\(\mathscr E_b\). Definimos
\begin{align}
 (R_a^+f)(x)&=\int_0^\infty e^{-at}f(x+t)\,dt,
 &(R_a^-f)(x)&=\int_0^\infty e^{-at}f(x-t)\,dt,\label{nt:resolvents}\\
 C_af&=f-2aR_a^+f,
 &C_a^{-1}f&=f-2aR_a^-f.\label{nt:test-cayley}
\end{align}
Cada derivada de \(R_a^\pm f\) satisface la cota del
correspondiente seminorma de \(f\), multiplicada por
\((a-b)^{-1}\). Esto sigue de
\(e^{b|x|}|f(x\pm t)|\leq M_0e^{bt}\).
Ambas aplicaciones preservan \(\mathscr E_b\).

Con la transformada
\(\widehat f(\xi)=\int_{\mathbb R}f(x)e^{-i\xi x}\,dx\),
la aplicación de Fubini da
\begin{equation}
 \widehat{C_af}(\xi)
 =c_a(\xi)\widehat f(\xi),\qquad
 c_a(\xi)=\frac{\xi-ia}{\xi+ia}.
 \label{nt:cayley-multiplier}
\end{equation}
El multiplicador de la segunda fórmula
de~\eqref{nt:test-cayley} es \(c_a^{-1}\), lo que prueba
la inversa. Para \(\xi\in\mathbb R\),
\(|c_a(\xi)|=1\); Plancherel proporciona su extensión
unitaria a \(L^2(\mathbb R)\).

Para precisar la conservación aritmética, sean
\[
 \mathcal C_{f,g}(u)
 =\int_{\mathbb R}f(x)\overline{g(x-u)}\,dx,
 \qquad
 P_f^\pm=\int_{\mathbb R}f(x)e^{\pm x/2}\,dx.
\]
Las longitudes \(\ell_{p,k}=k\log p\) y los pesos
\(w_{p,k}=(\log p)p^{-k/2}\) son las publicaciones
primo--potencia del lector aritmético. La forma completa
del artículo~VIII se escribe
\begin{equation}
 \begin{split}
 \mathscr W(f,g)={}&c_\Gamma\mathcal C_{f,g}(0)\\
 &+\int_0^\infty\frac{e^{-u/2}}{1-e^{-2u}}
 \bigl[2\mathcal C_{f,g}(0)
       -\mathcal C_{f,g}(u)-\mathcal C_{f,g}(-u)\bigr]\,du\\
 &-\sum_{p,k}w_{p,k}
   \bigl[\mathcal C_{f,g}(\ell_{p,k})
        +\mathcal C_{f,g}(-\ell_{p,k})\bigr]
 +P_f^+\overline{P_g^-}+P_f^-\overline{P_g^+},
 \end{split}
 \label{nt:weil-complete}
\end{equation}
donde \(c_\Gamma=\psi(1/4)-\log\pi\), y
\(\psi=\Gamma'/\Gamma\) es la derivada logarítmica de la
función gamma. Estas funciones y constantes intervienen
en la realización analítica posterior del lector
aritmético del corpus.

La cota
\[
 |\mathcal C_{f,g}(u)|
 \leq M_fM_g\bigl(|u|+b^{-1}\bigr)e^{-b|u|}
\]
controla la suma prima por una serie convergente del tipo
\(\sum_{n\geq2}(\log n)(1+\log n)n^{-b-1/2}\).
La diferencia simétrica de correlaciones es \(O(u^2)\)
en el origen, mientras que el peso gamma es \(O(u^{-1})\).
En infinito ese peso decrece exponencialmente.
El margen \(b>1/2\) asegura asimismo la convergencia de
los dos lectores polares. Así queda fijado el dominio
de todos los términos de~\eqref{nt:weil-complete}.

\begin{proposition}[Invariancia aritmética y balance nonádico]
Para \(f,g\in\mathscr E_b\),
\begin{equation}
 \mathscr W(C_af,C_ag)=\mathscr W(f,g).
 \label{nt:weil-invariance}
\end{equation}
Con \(\mathcal T_a=(8I+C_a)/9\), se cumple
\begin{equation}
 \mathscr W(f,g)
 =\mathscr W(\mathcal T_af,\mathcal T_ag)
 +\frac8{81}\mathscr W((I-C_a)f,(I-C_a)g).
 \label{nt:weil-balance}
\end{equation}
\end{proposition}

\begin{proof}
El multiplicador \(c_a\) es unitario en la recta real y
conmuta con las traslaciones. En consecuencia,
\(\mathcal C_{C_af,C_ag}(u)=\mathcal C_{f,g}(u)\)
para todo \(u\). Los lectores polares se transforman
por Fubini según
\[
 P_{C_af}^+
 =-\frac{a-1/2}{a+1/2}P_f^+,\qquad
 P_{C_af}^-
 =-\frac{a+1/2}{a-1/2}P_f^-.
\]
Sus factores son reales y recíprocos, de modo que
cada producto polar cruzado permanece inalterado.
La reunión de los términos de
\eqref{nt:weil-complete} prueba~\eqref{nt:weil-invariance}.
La expansión sesquilineal del segundo miembro de
\eqref{nt:weil-balance} cancela los términos mixtos y
produce
\(\bigl(8\mathscr W(f,g)+
\mathscr W(C_af,C_ag)\bigr)/9\), igual a
\(\mathscr W(f,g)\).
\end{proof}

La iteración finita conserva la forma completa:
\[
 \mathscr W(f,g)
 =\mathscr W(\mathcal T_a^Nf,\mathcal T_a^Ng)
 +\frac8{81}\sum_{j=0}^{N-1}
 \mathscr W((I-C_a)\mathcal T_a^jf,
            (I-C_a)\mathcal T_a^jg).
\]
Esta identidad sesquilineal conserva simultáneamente
gamma, primos y términos polares. Su evaluación en una
realización positiva permite interpretarla como reparto
de energía; su validez algebraica precede a esa
interpretación. El margen \(a>b>1/2\) pertenece al
dominio exponencial de pruebas de esta construcción.
El Cayley~\eqref{nt:cayley-hilbert} con parámetro \(1/2\)
pertenece al dominio hilbertiano de su generador.

\subsection{Covariancia métrica del flujo dodecafásico}

La reconstrucción de la sección~\ref{sec:memory} determina
\(u_K\in\mathbb R^{12}\), con
\(\mathbf1^{\mathsf T}u_K=0\) y \(u_K\ne0\).
Escribimos \(A_K=\operatorname{diag}(u_K)\), con la dirección y la
normalización declaradas en \eqref{eq:u-main}, y
\[
 g_s=e^{sA_K}\in\operatorname{GL}(12,\mathbb R).
\]
La traza nula de \(A_K\) implica \(\det g_s=1\).
Cada \(g_s\) transporta los pesos y las coordenadas
geométricas según la realización correspondiente.
Para conservar también el balance de memoria entre
cartas hay que transportar el producto interior.

\begin{proposition}[Transporte métrico de compresión y complemento]
Sean \(V\) un espacio hermítico finito, \(C_0\) unitario
y \(g:V\to V\) invertible. Escribimos
\(g^{-*}=(g^{-1})^*\) y definimos
\begin{equation}
 C_g=gC_0g^{-1},\qquad
 G_g=g^{-*}g^{-1},\qquad
 \langle v,w\rangle_g=\langle G_gv,w\rangle.
 \label{nt:metric-chart}
\end{equation}
Entonces \(C_g^*G_gC_g=G_g\), y
\begin{align}
 T_{C_g}g&=gT_{C_0},\label{nt:metric-naturality}\\
 \eta_{C_g}g&=g^{\oplus9}\eta_{C_0},\label{nt:eta-naturality}\\
 G_g&=T_{C_g}^*G_gT_{C_g}
 +\frac8{81}(I-C_g)^*G_g(I-C_g).\label{nt:metric-balance}
\end{align}
\end{proposition}

\begin{proof}
La positividad de \(G_g\) se obtiene de
\(\langle G_gv,v\rangle=\|g^{-1}v\|^2\).
Sustituyendo~\eqref{nt:metric-chart},
\[
 C_g^*G_gC_g
 =g^{-*}C_0^*g^*g^{-*}g^{-1}gC_0g^{-1}
 =g^{-*}C_0^*C_0g^{-1}=G_g.
\]
La identidad \(C_gg=gC_0\) da
\eqref{nt:metric-naturality} al sumar \(8g\) y dividir
por nueve. También da
\((C_g-I)g=g(C_0-I)\); su sustitución en
\eqref{nt:complement} prueba~\eqref{nt:eta-naturality}.
Finalmente, la expansión de
\eqref{nt:observable-balance}, con \(A=G_g\) y
\(C_g^*G_gC_g=G_g\), prueba
\eqref{nt:metric-balance}.
\end{proof}

Aplicando la proposición a \(g=g_s\) y a los operadores
\(C_0=S^3,S^4\) de la realización dodecafásica, las
compresiones, las memorias y sus productos interiores
se transportan por cuadrados conmutativos explícitos.
La identificación \(g_s:(V,\langle\cdot,\cdot\rangle)
\to(V,\langle\cdot,\cdot\rangle_{g_s})\) es isométrica.
Por tanto, la fórmula transporta también los sectores
fijos y los registros a toda profundidad:
\[
 P_{\mathrm{fix},s}
 =g_sP_{\mathrm{fix},0}g_s^{-1}.
\]
El proyector del miembro derecho es ortogonal respecto
de \(G_{g_s}\).

La condición de volumen y la condición métrica tienen
efectos diferenciados. Para la matriz racional
\[
 g=\begin{pmatrix}2&0\\0&1/2\end{pmatrix},
 \qquad \det g=1,
\]
la sustitución directa en la compresión da
\[
 \frac{8I+g}{9}
 =\begin{pmatrix}10/9&0\\0&17/18\end{pmatrix},
 \qquad
 T_g^*T_g+\frac8{81}(I-g)^*(I-g)
 =\frac{8I+g^*g}{9}.
\]
La primera coordenada se expande en la métrica
euclídea. El balance~\eqref{nt:metric-balance}
conserva en cambio la unidad métrica cuando se
transporta un operador unitario \(C_0\) junto con
su producto interior. El determinante controla
el volumen; \(G_g\) controla las normas y los
productos cruzados.

El flujo satisface \(g_{-s}=g_s^{-1}\).
La reciprocidad del radio de
\eqref{nt:mirror} pertenece a la carta espectral,
y la inversión de \(g_s\) pertenece a la
realización dimensional del registro. El enlace
constructivo utilizado aquí conserva las
aplicaciones explícitas de memoria a \(K\) y de
\(K\) a su dirección geométrica, seguido por
el transporte métrico de cada realización.
Esta composición mantiene, en un mismo desarrollo,
la reconstrucción exacta, el refinamiento positivo,
el registro de fronteras y la conservación
de los productos interiores.


% SOURCE sections/leech_foundations.tex
\section{Incidencia ternaria, pegado de discriminantes y construcción de Leech}
\label{sec:lf-foundations}

La construcción reticular utiliza la información de incidencia que
acompaña al registro dodecafásico. Su dominio conserva las palabras
ternarias, el marco en que se representan, la orientación previa al
acarreo y la bandera que distingue una de las doce posiciones. El
paso a un retículo transmite esos datos mediante operaciones
sucesivas: codificación lineal, selección incidencial, pegado de
discriminantes y formación de un vecino. La clasificación reticular
interviene después de probar la positividad, paridad, autodualidad
y ausencia de raíces de la construcción obtenida.

Estas operaciones continúan la misma generación
APP--TRIT--TPK y el mismo refinamiento compatible del continuo
conjunto que anteceden a la lectura numérica \cite{HMT-IV}.
La matriz finita, el código y la métrica que se definen a
continuación son realizaciones con dominios precisos. La prueba
reticular recibe el estado marcado; el escalar final de una
constante no desempeña el papel de código de pegado ni de longitud.

\subsection{Carta elíptica y transporte de la incidencia}
\label{subsec:lf-incidence}

Las seis unidades \(\{1,2,4,5,7,8\}\) módulo nueve, conservadas
por el transporte modular, admiten una carta de seis posiciones.
La identificamos con
\(\mathbb P^1(\mathbb F_5)=\{\infty,0,1,2,3,4\}\), en ese orden.
Esta identificación fija coordenadas para realizar la relación
cuadrática del régimen elíptico del TRIT.

Sea \(\chi:\mathbb F_5\to\{0,1,-1\}\) el carácter cuadrático:
\(\chi(0)=0\), \(\chi(1)=\chi(4)=1\) y
\(\chi(2)=\chi(3)=-1\). Definimos la matriz entera
\(C_{\mathrm P}\) con diagonal cero,
\((C_{\mathrm P})_{\infty,a}=(C_{\mathrm P})_{a,\infty}=1\),
y entradas finitas \((C_{\mathrm P})_{a,b}=\chi(b-a)\).
Su cuadrado es \(5I_6\). En efecto, cada entrada diagonal del
cuadrado suma cinco unidades. Las entradas que involucran
\(\infty\) fuera de la diagonal se anulan porque la suma de
\(\chi\) es cero. Para dos índices finitos distintos, la
contribución uno de la posición \(\infty\) se cancela con
\(\sum_c\chi(c-a)\chi(c-b)=-1\). Esta última identidad se verifica
directamente sobre los cinco residuos, o mediante la sustitución
afín que lleva \((a,b)\) a \((0,1)\).

La reducción módulo tres produce la matriz simétrica
\begin{equation}
 A_W=
 \begin{pmatrix}
 0&1&1&1&1&1\\
 1&0&1&2&2&1\\
 1&1&0&1&2&2\\
 1&2&1&0&1&2\\
 1&2&2&1&0&1\\
 1&1&2&2&1&0
 \end{pmatrix},
 \qquad A_W^2=2I_6=-I_6
 \quad\text{en }\mathbb F_3.
 \label{eq:lf-aw}
\end{equation}
La igualdad conserva el tipo elíptico en este codominio finito.
El cuerpo finito y la posterior realización arquimediana de la
unidad imaginaria tienen dominios diferentes.

Para una matriz \(A\in M_6(\mathbb F_3)\) y una palabra fila
\(w\in\mathbb F_3^6\), escribimos
\[
 \gamma_A(w)=(w,wA)\in\mathbb F_3^{12}.
\]
Si \(f\in\operatorname{GL}_6(\mathbb F_3)\), entonces
\begin{equation}
 \gamma_{fAf^{-1}}(wf^{-1})
   =\gamma_A(w)(f^{-1}\oplus f^{-1}).
 \label{eq:lf-frame-covariance}
\end{equation}
La primera mitad es \(wf^{-1}\); la segunda es
\(wf^{-1}fAf^{-1}=wAf^{-1}\), lo que prueba la identidad.

Sea \(X_n^{\mathrm{mar}}\) el espacio de prefijos enriquecidos de
longitud \(n\), con marco inicial declarado, y sean
\(p_{n,m}:X_n^{\mathrm{mar}}\to X_m^{\mathrm{mar}}\) sus
truncaciones. En el paso \(j\), el prefijo contiene una palabra
visible \(w_j\) y un marco \(f_j\). La actualización del marco
tiene la forma \(f_{j+1}=g_j(x)f_j\), donde \(g_j(x)\) depende
únicamente de información ya contenida en el prefijo. Una
prolongación compatible conserva, por ello, todos los marcos
anteriores.

\begin{proposition}[Naturalidad del transporte incidencial]
\label{prop:lf-incidence-naturality}
La aplicación
\begin{equation}
 \mathcal Q_n(x)=
 \bigl((f_j,\gamma_{f_jA_Wf_j^{-1}}(w_j))\bigr)_{0\leq j<n}
 \label{eq:lf-Qn}
\end{equation}
satisface \(r_{n,m}\mathcal Q_n=\mathcal Q_mp_{n,m}\), donde \(r_{n,m}\) trunca
las primeras \(m\) componentes codificadas. Por tanto induce una
aplicación entre los límites inversos de los sistemas de
prefijos y de incidencias marcadas.
\end{proposition}
\begin{proof}
Dos prefijos compatibles tienen las mismas palabras y los mismos
operadores de transporte hasta el paso \(m-1\). Partiendo del
mismo marco inicial, la relación recursiva da por inducción los
mismos \(f_j\) hasta ese paso. Coinciden entonces las primeras
\(m\) componentes de \eqref{eq:lf-Qn}. La propiedad universal
del límite inverso produce la aplicación de prolongación.
\end{proof}

En cada marco, la primera proyección de \(\gamma_A\) recupera
exactamente la palabra visible. Este alcance es local: la
memoria profunda se conserva en el estado enriquecido y su
recuperación exige las coordenadas correspondientes. Al pasar
de palabras a soportes se retiene todavía menos información.
Asimismo, un cambio general de base en
\(\operatorname{GL}_6(\mathbb F_3)\) transporta también la métrica
de Hamming y la incidencia de la carta. Los cambios monomiales
las preservan directamente. Los soportes utilizados en esta
sección se calculan en la carta explícita \eqref{eq:lf-aw}.

La naturalidad anterior concreta la lectura incidencial que
acompaña a los cilindros numéricos. El encajamiento de cilindros
y la restricción de \eqref{eq:lf-Qn} actúan sobre un mismo
prefijo compatible, conservando respectivamente su evaluación
posicional y sus relaciones marcadas. El ambiente
\(\mathbb F_3^6\), de \(729\) palabras, proporciona el dominio
lineal de la codificación; conserva su distinción respecto de
las \(243\) palabras visibles y las \(468\) emisiones del censo
admisible de APP--TRIT--TPK.

\subsection{Código ternario autodual y enumeración de sus pesos}
\label{subsec:lf-code}

\begin{definition}[Código de la carta dodecafásica]
El código lineal asociado a \eqref{eq:lf-aw} es
\[
 C_W=\{(w,wA_W):w\in\mathbb F_3^6\}\subset\mathbb F_3^{12}.
\]
El peso \(\operatorname{wt}(c)\) cuenta las coordenadas no nulas
de \(c\), y \(\operatorname{supp}(c)\subset\{1,\ldots,12\}\)
es su conjunto de posiciones no nulas. Una hexada es el soporte
de una palabra de peso seis; su familia se denota \(\mathcal H\).
\end{definition}

\begin{theorem}[Autodualidad, distancia mínima y enumerador]
\label{thm:lf-code}
El código \(C_W\) es autodual de parámetros \([12,6,6]_3\) y
tiene enumerador
\begin{equation}
 W_{C_W}(z)=1+264z^6+440z^9+24z^{12}.
 \label{eq:lf-enumerator}
\end{equation}
\end{theorem}
\begin{proof}
La matriz generadora \([I_6\mid A_W]\) tiene rango seis.
Por la simetría de \(A_W\) y \eqref{eq:lf-aw},
\[
 [I_6\mid A_W][I_6\mid A_W]^{\mathsf T}
       =I_6+A_W^2=0.
\]
Así, el código es autoortogonal de dimensión mitad de la
longitud y, en consecuencia, autodual.

La enumeración restante es una operación entera finita sobre
la matriz explícita. Para \(r\in\{0,\ldots,6\}\), agrupamos
las palabras por \(\operatorname{wt}(w)=r\). La tabla siguiente
da el número de palabras con cada peso total:
\[
\begin{array}{c|rrrr|r}
r&0&6&9&12&\text{total}\\ \hline
0&1&0&0&0&1\\
1&0&12&0&0&12\\
2&0&60&0&0&60\\
3&0&120&40&0&160\\
4&0&60&180&0&240\\
5&0&12&180&0&192\\
6&0&0&40&24&64\\ \hline
\text{total}&1&264&440&24&729
\end{array}
\]
Cada fila contiene \(\binom6r2^r\) palabras. Para verificar
todas sus entradas, sea \(q\) un entero de \(0\) a \(728\), y
definamos sus seis cifras ternarias
\[
 d_i(q)=\left\lfloor\frac{q}{3^{6-i}}\right\rfloor\bmod3,
 \qquad
 e_j(q)=\left(\sum_{i=1}^{6}d_i(q)(A_W)_{ij}\right)\bmod3.
\]
La fila y la columna de la contribución de \(q\) son,
respectivamente,
\[
 r(q)=\sum_i\mathbf1_{\{d_i(q)\ne0\}},\qquad
 h(q)=r(q)+\sum_j\mathbf1_{\{e_j(q)\ne0\}}.
\]
La recurrencia de acumulación empieza con
\(T^{(0)}_{r,h}=0\) y aplica
\[
 T^{(q+1)}_{r,h}
 =T^{(q)}_{r,h}
   +\mathbf1_{\{r=r(q)\}}\mathbf1_{\{h=h(q)\}}
 \quad(0\leq q<729).
\]
Produce exactamente la tabla mostrada. La expansión ternaria
da una biyección entre esos \(729\) enteros y todas las palabras
de \(\mathbb F_3^6\), por lo que la operación agota el código.
La suma de sus filas prueba \eqref{eq:lf-enumerator}; el primer
peso no nulo es seis.
\end{proof}

\begin{remark}[Evaluación reproducible de la recurrencia finita]
El siguiente bucle implementa literalmente las sumas de la prueba.
Sus únicos datos son los enteros de la matriz
\eqref{eq:lf-aw}; calcula el enumerador antes de compararlo con
la igualdad obtenida.
\end{remark}
{\small
\begin{verbatim}
from collections import Counter
from itertools import product
A = ((0,1,1,1,1,1), (1,0,1,2,2,1),
     (1,1,0,1,2,2), (1,2,1,0,1,2),
     (1,2,2,1,0,1), (1,1,2,2,1,0))
rows = [Counter() for _ in range(7)]
for w in product(range(3), repeat=6):
    v = tuple(sum(w[i]*A[i][j] for i in range(6)) % 3
              for j in range(6))
    r = sum(x != 0 for x in w)
    h = r + sum(x != 0 for x in v)
    rows[r][h] += 1
total = sum(rows, Counter())
if dict(total) != {0:1, 6:264, 9:440, 12:24}:
    raise ArithmeticError("Enumerador incompatible")
\end{verbatim}
}

\begin{proposition}[Hexadas y diseño de Witt]
\label{prop:lf-witt}
La familia \(\mathcal H\) contiene \(132\) hexadas y cada
quintupla de posiciones pertenece a una única hexada.
\end{proposition}
\begin{proof}
Dos palabras de peso seis con el mismo soporte difieren por
un escalar. En efecto, entre sus seis cocientes de coordenadas
no nulas, uno de los signos \(1,-1\) aparece al menos tres
veces. Restando el múltiplo correspondiente, si las palabras
no fueran proporcionales, se obtendría una palabra no nula
de peso a lo sumo tres, en contradicción con
el teorema~\ref{thm:lf-code}. Cada soporte corresponde así a
la pareja \(\{c,-c\}\), y hay \(264/2=132\) soportes.

Si dos soportes distintos compartieran cinco posiciones,
en esa intersección uno de los dos cocientes aparecería al
menos tres veces. La combinación lineal que cancela dichas
coordenadas tendría soporte en una unión de tamaño siete
y peso a lo sumo cuatro. La distancia mínima lo impide.
Una quintupla pertenece, por tanto, a lo sumo a una hexada.
Como
\[
 132\binom65=792=\binom{12}{5},
\]
pertenece a exactamente una.
\end{proof}

El objeto construido es el código ternario extendido de Golay
con su pequeño diseño de Witt. El reconocimiento se efectúa
después de la codificación, la autodualidad y el cómputo de
distancia. Para la construcción reticular sólo se necesitan
las propiedades ya demostradas y la bandera marcada que se
obtiene a continuación.

\subsection{Bandera dodecafásica y selección del origen}
\label{subsec:lf-flag}

En el marco del registro generado se conservan las palabras
\[
 w_\pi=010211,\qquad
 w_e=201101,\qquad
 w_\varphi=121200,\qquad
 w_{\mathrm{APP}}=022020.
\]
Los subíndices \(\pi,e,\varphi\) identifican las correspondientes
regiones de lectura de las constantes; \(w_e\) conserva aquí
la región de Euler. Aplicar \(\gamma_{A_W}\) y tomar soportes da
\[
\begin{aligned}
 P_\pi&=\{2,4,5,6,7,8,9,10,11\},\\
 H_e&=\{1,3,4,6,9,10\},\\
 H_\varphi&=\{1,2,3,4,7,9\},\\
 H_{\mathrm{APP}}&=\{2,3,5,10,11,12\}.
\end{aligned}
\]
El primer soporte tiene peso nueve; los tres restantes son
hexadas. El registro
\[
 K=(234,543,140,729,659,824,621,58,914,794,146,601)
\]
selecciona
\[
 B_K=\{j:K_j\geq729\}=\{4,6,9,10\}=H_e\cap P_\pi.
\]
El umbral \(729\) pertenece al formato de bloques de la
lectura declarada. En este registro concreto, los umbrales
enteros de \(660\) a \(729\) producen el mismo soporte, de
modo que el soporte conserva la incidencia seleccionada,
mientras el registro completo conserva también los bloques
y el umbral de su lectura.

El registro firmado anterior al acarreo es
\[
\begin{split}
 Z_0={}&(7,297,353,-430,-715,-1199,\\
       &\hspace{9mm}-3,286,-894,-528,380,663).
\end{split}
\]
Su soporte negativo es la hexada
\[
 H_\alpha^-=\{4,5,6,7,9,10\}.
\]
Los signos se reciben del estado previo al acarreo y no se
extraen de la expansión del escalar final.

\begin{proposition}[Coincidencia de selección firmada e incidencial]
\label{prop:lf-flag}
La hexada \(H_\alpha^-\) es la única \(H\in\mathcal H\) que cumple
\[
\begin{gathered}
 B_K\subset H\subset P_\pi,\\
 (|H\cap H_e|,|H\cap H_\varphi|,
                   |H\cap H_{\mathrm{APP}}|)=(4,3,2).
\end{gathered}
\]
\end{proposition}
\begin{proof}
La codificación de las palabras de peso seis, o la recurrencia
finita anterior aplicada a sus soportes, da las cuatro
hexadas sobre \(B_K\). Sus pares complementarios son
\[
 \{1,3\},\qquad\{2,11\},\qquad\{5,7\},\qquad\{8,12\}.
\]
Sólo el segundo y el tercero están contenidos en \(P_\pi\).
La hexada del segundo tiene intersecciones de tamaño tres
con \(H_\varphi\) y con \(H_{\mathrm{APP}}\), y por tanto
incumple la última condición. La del tercero tiene
intersecciones \(4,3,2\), y coincide con \(H_\alpha^-\).
\end{proof}

\begin{lemma}[Origen de la bandera marcada]
\label{lem:lf-origin}
El origen seleccionado por la incidencia es
\[
 o_*=(H_e\cap H_\varphi)
                 \setminus(H_\alpha^-\cup H_{\mathrm{APP}})
     =\{1\}.
\]
La selección conmuta con toda biyección aplicada
simultáneamente a los soportes.
\end{lemma}
\begin{proof}
La intersección \(H_e\cap H_\varphi\) es
\(\{1,3,4,9\}\). La unión sustraída contiene \(3,4,9\)
y excluye \(1\). La equivariancia procede de que las
biyecciones conmutan con unión, intersección y diferencia.
\end{proof}

La bandera distingue una componente entre doce antes del
paso métrico. El vector radial se determina después,
dentro de una cámara positiva declarada. Así quedan
separadas y articuladas la selección combinatoria y la
operación reticular.

\subsection{Pegado ternario de las doce componentes}
\label{subsec:lf-glue}

Sea \(A_2=\mathbb Za_1\oplus\mathbb Za_2\), con forma
bilineal de matriz
\[
 B_2=\begin{pmatrix}2&-1\\-1&2\end{pmatrix},
 \qquad \rho=a_1+a_2,\qquad \rho^2=2.
\]
Para \(x=ua_1+va_2\),
\[
 x^2=2(u^2-uv+v^2).
\]
La identidad
\(4(u^2-uv+v^2)=(2u-v)^2+3v^2\) muestra
la positividad. El determinante es tres. Las tres
clases de \(A_2^*/A_2\) se representan mediante
\[
 \varpi_0=0,\qquad
 \varpi_1=\frac{2a_1+a_2}{3},\qquad
 \varpi_2=\frac{a_1+2a_2}{3}.
\]
El emparejamiento con \(a_1,a_2\) verifica que estos
vectores pertenecen a \(A_2^*\), y
\(2\varpi_1-\varpi_2=a_1\). Por ello
\(j\mapsto\varpi_j+A_2\) identifica aditivamente
\(\mathbb F_3\) con el discriminante.

Para \(c\in C_W\), definimos
\(\phi(c)=(\varpi_{c_1},\ldots,\varpi_{c_{12}})\)
y
\begin{equation}
 N=\bigcup_{c\in C_W}\bigl(A_2^{12}+\phi(c)\bigr).
 \label{eq:lf-glue}
\end{equation}

\begin{theorem}[Retículo de pegado par y unimodular]
\label{thm:lf-niemeier}
El conjunto \(N\) es un retículo positivo, par y
unimodular de rango \(24\). Sus vectores de norma dos
son exactamente las raíces de \(A_2^{12}\).
\end{theorem}
\begin{proof}
La linealidad del código y la aditividad del discriminante
hacen que la unión sea un subgrupo de
\((A_2^*)^{12}\) que contiene \(A_2^{12}\).
Es, por tanto, un retículo de rango \(24\).

El emparejamiento de las clases discriminantes es
\(2cd/3\) módulo \(\mathbb Z\) en cada componente.
La autoortogonalidad de \(C_W\) anula su suma y prueba
la integralidad de \(N\). La norma de una clase
representada por \(\phi(c)\) es
\(2\operatorname{wt}(c)/3\) módulo \(2\mathbb Z\).
Todos los pesos de \eqref{eq:lf-enumerator} son
múltiplos de tres, luego \(N\) es par.

Su índice sobre \(A_2^{12}\) es \(|C_W|=3^6\).
Con el determinante de la matriz de Gram,
\[
 \det N=\frac{\det(A_2^{12})}{[N:A_2^{12}]^2}
       =\frac{3^{12}}{(3^6)^2}=1.
\]
La forma ambiente es la suma ortogonal de doce formas
positivas; queda demostrada la positividad.
La integralidad y el determinante uno implican
\(N^*=N\).

En una clase no nula del discriminante local, un
vector tiene la forma \((ua_1+va_2)/3\), con residuos
\((u,v)\equiv(2,1)\) o \((1,2)\pmod3\).
El entero \(u^2-uv+v^2\) es positivo y divisible
por tres, luego es al menos tres. Los representantes
\(\varpi_1,\varpi_2\) alcanzan ese valor: la norma
mínima de cada clase no nula es \(2/3\).
Una palabra no nula del código tiene por lo menos
seis coordenadas no nulas. Todo vector de una clase
de pegado no trivial tiene, por tanto, norma
al menos \(6(2/3)=4\).

En \(A_2^{12}\), una norma dos sólo puede proceder
de una única componente, porque cada componente
tiene norma par positiva mínima dos. Las soluciones
de \(u^2-uv+v^2=1\) son
\((\pm1,0),(0,\pm1),(1,1),(-1,-1)\).
Son las seis raíces de \(A_2\), y completan el censo
de raíces de \(N\).
\end{proof}

Las propiedades anteriores identifican \(N\) como el
retículo de Niemeier de sistema de raíces \(A_2^{12}\).
La descripción de todas sus raíces es necesaria para
controlar cuáles sobreviven al paso al vecino.

\subsection{Cámara radial y vecino sin raíces}
\label{subsec:lf-neighbour}

Se fija la cámara radial
\[
 v(a)=\sum_{j=1}^{12}a_j\rho_j,\qquad
 a_j=1+3b_j,\qquad b_j\in\mathbb Z_{\geq0},
\]
donde \(\rho_j\) es la copia de \(\rho\) en la
componente \(j\). La paridad de un vecino de índice
tres requiere \(v(a)^2\equiv0\pmod{18}\).
Como
\[
 v(a)^2=2\sum_j(1+3b_j)^2
       =24+12\sum_jb_j+18\sum_jb_j^2,
\]
la congruencia equivale a \(\sum_jb_j\equiv1\pmod3\).
La norma mínima dentro de la cámara se alcanza con
un único \(b_j=1\) y los demás cero, y vale \(54\).
El origen \(o_*=1\) del lema~\ref{lem:lf-origin}
selecciona
\begin{equation}
 v_\alpha=4\rho_1+\rho_2+\cdots+\rho_{12},
 \qquad v_\alpha^2=54.
 \label{eq:lf-valpha}
\end{equation}
El calificativo de mínimo se refiere a la cámara
declarada. En la clase residual completa,
\((a_1-2a_2,\rho,\ldots,\rho)\) representa la
misma clase módulo tres y tiene norma
\(14+11\cdot2=36\). Esta distinción preserva
la condición que determina el coeficiente cuatro.

Definimos
\begin{equation}
\begin{split}
 N_0&=\{x\in N:\langle x,v_\alpha\rangle\equiv0\pmod3\},\\
 \Lambda&=N_0+\mathbb Z\,v_\alpha/3.
\end{split}
\label{eq:lf-lambda}
\end{equation}

\begin{theorem}[Vecino par autodual sin raíces]
\label{thm:lf-leech}
La construcción \eqref{eq:lf-lambda} es un
retículo positivo, par y autodual de rango \(24\)
sin vectores de norma dos. En consecuencia, es
isométrica al retículo de Leech.
\end{theorem}
\begin{proof}
El carácter \(x\mapsto\langle x,v_\alpha\rangle\bmod3\)
es no nulo. Una raíz simple de la componente \(j\)
tiene emparejamiento \(a_j\equiv1\pmod3\) con
\(a_j\rho_j\). Por tanto \([N:N_0]=3\), de donde
\(\det N_0=9\).

Para \(x\in N_0\), el emparejamiento
\(\langle x,v_\alpha/3\rangle\) es entero, y así
\(v_\alpha/3\in N_0^*\). Además,
\(\langle v_\alpha,v_\alpha\rangle=54\) implica
\(v_\alpha\in N_0\). Si \(v_\alpha/3\in N\),
la integralidad de \(N\) obligaría a que todos los
emparejamientos de \(N\) con \(v_\alpha\) fuesen
divisibles por tres, contradiciendo el carácter
no nulo. La clase de \(v_\alpha/3\) sobre \(N_0\)
tiene entonces orden exactamente tres. Resulta
\[
 [\Lambda:N_0]=3,\qquad
 \det\Lambda=\frac9{3^2}=1.
\]
Los emparejamientos de \(N_0\) con el nuevo
generador son enteros y la norma de este generador
es seis. La extensión es integral. Para
\(x\in N_0\) y \(m\in\mathbb Z\),
\[
 \left(x+m\frac{v_\alpha}{3}\right)^2
 =x^2+2m\frac{\langle x,v_\alpha\rangle}{3}
       +6m^2\in2\mathbb Z.
\]
Por ello \(\Lambda\) es par; su integralidad y
determinante uno prueban la autodualidad.
La positividad y el rango se heredan del espacio
euclídeo ambiente.

Para excluir raíces, primero consideramos la
clase \(N_0\). Todas las raíces de \(N\) pertenecen
a \(A_2^{12}\) por el teorema~\ref{thm:lf-niemeier}.
Los emparejamientos de sus seis raíces locales
con \(a_j\rho_j\) son \(\pm a_j,\pm2a_j\),
congruentes con \(\pm1,\pm2\pmod3\).
Ninguna raíz pertenece a \(N_0\).

Las otras dos clases de \(\Lambda/N_0\) pueden
representarse con \(\pm v_\alpha/3\).
En cada componente, sus vectores pertenecen a
una de las clases
\[
 A_2+\varpi_j\pm\rho/3,\qquad j=0,1,2.
\]
La componente distinguida tiene la misma
descripción porque \(4\rho/3\equiv\rho/3\pmod{A_2}\).
Para los numeradores \((u,v)\) en
\((ua_1+va_2)/3\), los seis residuos son
\[
 (1,1),(0,2),(2,0),(2,2),(1,0),(0,1)\pmod3.
\]
Todos son no nulos. La positividad de
\(u^2-uv+v^2\) obliga a que sea al menos uno,
y los representantes
\((1,1),(0,-1),(-1,0),(-1,-1),(1,0),(0,1)\)
alcanzan ese mínimo en las seis clases.
Cada componente tiene norma al menos \(2/9\).
Un vector de cualquiera de las dos clases
nuevas tiene norma al menos
\[
 12\cdot\frac29=\frac83>2.
\]
Se han excluido raíces en las tres clases.
El teorema de clasificación de los retículos
positivos pares unimodulares de rango veinticuatro
identifica el único caso sin raíces con el
retículo de Leech \cite{ConwaySloane}.
\end{proof}

La identificación final utiliza una clasificación
externa sobre un objeto cuyas propiedades ya han
sido demostradas. La selección de la componente
distinguida conserva la bandera procedente del
registro firmado; el subíndice \(\alpha\) expresa
esa procedencia. La norma y los coeficientes del
vector radial se obtienen de la métrica y de
la cámara mediante la congruencia indicada.

El retículo \(\Lambda\), su forma bilineal, sus
doce planos ambientes y el vector \(v_\alpha/3\)
quedan así determinados con todas las propiedades
utilizadas por la deformación posterior. El
código contiene asimismo la palabra
\[
 c_*=(1,1,1,1,1,1,2,1,1,1,1,1),
\]
pues \(q_*=(1,1,1,1,1,1)\) satisface
\(q_*A_W=(2,1,1,1,1,1)\). Éste es el dato
de soporte completo que permite construir y
probar la estabilidad de la isometría de orden
tres en la sección~\ref{sec:leech-dualidad}.


% SOURCE sections/leech_dualidad.tex
\section{Polarización dodecafásica, dualidad reticular y reciprocidad theta}
\label{sec:leech-dualidad}

La incidencia excepcional y la polarización del registro dodecafásico
proceden de dos realizaciones del mismo estado enriquecido. La primera
conserva las relaciones de frontera, el código de pegado y el origen
marcado que determinan el retículo; la segunda hace actuar una
representación sobre las doce coordenadas enteras ya obtenidas por los
lectores del Topological Phase Kernel. Su composición permite estudiar
una deformación métrica cuyo parámetro modifica las longitudes relativas
de doce planos y conserva simultáneamente el volumen total y una
simetría ternaria explícita.

El punto de partida causal sigue siendo la construcción
APP--TRIT--TPK del estado enriquecido y del continuo conjunto.
El registro \(K\) se recibe después del registro dodecafásico y de sus
operaciones de incidencia y orientación. La representación empleada
actúa sobre esa salida y conserva sus doce valores como datos ya
determinados. La realización reticular se recibe con su forma bilineal,
sus clases de pegado y su vector marcado. Estas estructuras, desarrolladas
respectivamente en \cite{hmtX,HMT-IV}, se mantienen al componerlas.
La familia que se construye a continuación es una realización geométrica
de esos datos con una carta común de doce posiciones.

\subsection{Polarización del registro y marcación de los doce planos}
\label{subsec:leech-polarizacion}

El registro dodecafásico y su suma son
\begin{equation}
 \begin{split}
 K={}&(234,543,140,729,659,824,\\
     &\hspace{13mm}621,58,914,794,146,601),\\
 \sum_{j=1}^{12}K_j={}&6263.
 \end{split}
 \label{eq:leech-K}
\end{equation}
Sea \(A_5\) el grupo de permutaciones pares de cinco elementos y sea
\(C_5=\langle(1\,2\,3\,4\,5)\rangle\). Sus doce clases izquierdas
determinan una representación por permutaciones
\(\varrho:A_5\to O(\mathbb R^{12})\). La carta marcada se fija mediante
los representantes siguientes, escritos como listas de imágenes:
\[
\begin{array}{c|c@{\qquad}c|c}
j&r_j&j&r_j\\ \hline
1&(4,3,5,2,1)&7&(5,4,3,2,1)\\
2&(4,2,1,3,5)&8&(1,3,4,2,5)\\
3&(1,2,4,5,3)&9&(4,2,3,5,1)\\
4&(2,5,3,4,1)&10&(3,2,5,4,1)\\
5&(4,3,1,5,2)&11&(4,5,2,3,1)\\
6&(5,1,2,3,4)&12&(5,3,2,4,1).
\end{array}
\]
Concretamente, \(\varrho(a)e_j=e_\ell\) si
\(ar_jC_5=r_\ell C_5\). El carácter tridimensional \(\chi_3\) toma
los valores
\[
\begin{array}{c|ccccc}
\text{clase}&1&2^2&3&5_a&5_b\\ \hline
\chi_3&3&-1&0&(1+\sqrt5)/2&(1-\sqrt5)/2.
\end{array}
\]
La clase \(5_a\) contiene el generador de \(C_5\) indicado y \(5_b\)
contiene su cuadrado. Esta elección distingue los dos caracteres
tridimensionales conjugados. Si \(I_{12}\) es la identidad y
\(\mathbf1\) es el vector de doce unos, se definen
\begin{equation}
 P_3=\frac3{60}\sum_{a\in A_5}\chi_3(a^{-1})\varrho(a),
 \qquad
 P_{11}=I_{12}-\frac1{12}\mathbf1\mathbf1^{\mathsf T}.
 \label{eq:leech-projectors}
\end{equation}
El campo \(\mathbb Q(\sqrt5)\) aparece aquí en la realización de la
representación que actúa sobre el registro generado.

\begin{lemma}[Polarización de suma nula]
\label{lem:leech-u}
Los operadores de \eqref{eq:leech-projectors} satisfacen
\[
 P_3=P_3^{\mathsf T}=P_3^2,\qquad
 \operatorname{rank}P_3=3,\qquad P_3\mathbf1=0,\qquad
 P_3P_{11}=P_{11}P_3=P_3.
\]
Por consiguiente,
\begin{equation}
 u=P_3P_{11}K=P_3K,\qquad
 \sum_{j=1}^{12}u_j=0,\qquad
 \|u\|^2=\frac{6638585+2275584\sqrt5}{20}>0.
 \label{eq:leech-u-norm}
\end{equation}
\end{lemma}
\begin{proof}
La suma del carácter es el proyector isótipo de la representación.
La realidad del carácter y la sustitución \(a\mapsto a^{-1}\)
dan su simetría. La multiplicidad del carácter tridimensional en la
representación de \(A_5/C_5\) es la dimensión de sus invariantes por
\(C_5\), que vale
\[
 \frac15\left(
 3+2\frac{1+\sqrt5}{2}+2\frac{1-\sqrt5}{2}
 \right)=1.
\]
El proyector tiene por ello rango tres. Su ortogonalidad al carácter
trivial implica \(P_3\mathbf1=0\) y las identidades con \(P_{11}\).
La aplicación de la suma finita al registro
\eqref{eq:leech-K} produce las doce coordenadas
\begin{equation}
\begin{aligned}
u={}&(-495/4-233\sqrt5/10,\quad403/4+169\sqrt5/10,\\
&-403/4-169\sqrt5/10,\quad495/4+233\sqrt5/10,\\
&601/4+1031\sqrt5/10,\quad203/4+211\sqrt5/5,\\
&-203/4-211\sqrt5/5,\quad-601/4-1031\sqrt5/10,\\
&313/4+569\sqrt5/10,\quad162+219\sqrt5/20,\\
&-162-219\sqrt5/20,\quad-313/4-569\sqrt5/10).
\end{aligned}
\label{eq:leech-u-coordinates}
\end{equation}
Sumarlas y elevarlas al cuadrado en \(\mathbb Q(\sqrt5)\) da
\eqref{eq:leech-u-norm}. Equivalentemente,
\[
 \|u\|^2=K^{\mathsf T}P_3K
 =\frac1{20}\sum_{a\in A_5}
          \chi_3(a^{-1})K^{\mathsf T}\varrho(a)K.
\]
Los coeficientes positivos de su evaluación prueban que \(u\ne0\).
\end{proof}

La suma nula tiene una función geométrica concreta: las expansiones de
unos planos podrán compensar las contracciones de los restantes en el
determinante total. La no nulidad de \(u\), en cambio, garantiza que esa
compensación contiene una deformación efectiva.

El espacio euclídeo ambiente de la realización reticular es
\begin{equation}
 E=\bigoplus_{j=1}^{12}E_j,\qquad
 B=\operatorname{diag}(B_2,\ldots,B_2),\qquad
 B_2=\begin{pmatrix}2&-1\\-1&2\end{pmatrix}.
 \label{eq:leech-ambient}
\end{equation}
Cada \(E_j\) es el plano real generado por el retículo de raíces
\(A_2\), expresado en su base de raíces simples. La descomposición
ortogonal pertenece al espacio ambiente. El retículo de Leech se
obtiene mediante pegado y paso a un vecino, operaciones que relacionan
los doce factores.

Sea \(C_W\subset\mathbb F_3^{12}\) el código ternario de dimensión seis
de la construcción Paley--Witt. Identificando el discriminante de cada
factor \(A_2\) con \(\mathbb F_3\), el pegado es
\[
 N=N(C_W)=
 \bigcup_{c\in C_W}\bigl(A_2^{12}+\phi(c)\bigr),
\]
donde \(\phi(c)\) elige representantes de las clases discriminantes.
La autoortogonalidad y autodualidad del código transmiten paridad e
integralidad a \(N\); su determinante es
\[
 \det N=\frac{3^{12}}{(3^6)^2}=1.
\]
La construcción excepcional fija el origen marcado \(j=1\) y el
vector
\begin{equation}
 v_\alpha=4\rho_1+\rho_2+\cdots+\rho_{12},
 \qquad \rho_j=(1,1)_j,\qquad
 \|v_\alpha\|_B^2=54.
 \label{eq:leech-marked-vector}
\end{equation}
Con él se forman
\begin{equation}
\begin{split}
 N_{\alpha,0}
  &=\{x\in N:\langle x,v_\alpha\rangle_B\equiv0\pmod3\},\\
 \Lambda&=N_{\alpha,0}+\mathbb Z\,v_\alpha/3.
\end{split}
\label{eq:leech-neighbour}
\end{equation}
La realización excepcional establece que \(\Lambda\) es un retículo
positivo, par, unimodular, sin raíces y de rango \(24\), de modo que es
isométrico al retículo de Leech \cite{HMT-IV}. En particular,
\(\Lambda^*=\Lambda\), con dual respecto de \(B\).
Estas propiedades y la pertenencia \(v_\alpha/3\in\Lambda\) son las
hipótesis reticulares precisas empleadas en los teoremas siguientes.

\subsection{Isometría ternaria y estabilidad del vecino}
\label{subsec:leech-c3}

Definimos las matrices bidimensionales
\[
 C=\begin{pmatrix}0&-1\\1&-1\end{pmatrix},
 \qquad
 R=\begin{pmatrix}0&1\\1&0\end{pmatrix},
 \qquad
 \varpi=\frac13\begin{pmatrix}2\\1\end{pmatrix}.
\]
La multiplicación directa verifica
\[
 C^3=I_2,\quad C^2+C+I_2=0,\quad
 C^{\mathsf T}B_2C=B_2,\quad
 RCR=C^{-1},\quad R^{\mathsf T}B_2R=B_2.
\]
Además, \(C\varpi-\varpi=(-1,0)^{\mathsf T}\), por lo que \(C\)
actúa trivialmente en \(A_2^*/A_2\). La reflexión \(R\) cambia el signo
de la clase discriminante. El código \(C_W\) contiene la palabra de
soporte completo
\[
 c_*=(1,1,1,1,1,1,2,1,1,1,1,1).
\]
En la presentación \(C_W=\{(q,qA_W):q\in\mathbb F_3^6\}\), dicha
palabra procede de \(q_*=(1,1,1,1,1,1)\) y
\(q_*A_W=(2,1,1,1,1,1)\).

Para cada plano se fija el marco
\begin{equation}
 Q_j=\begin{cases}
 R,&(c_*)_j=1,\\
 I_2,&(c_*)_j=2,
 \end{cases}
 \qquad
 g_c=\bigoplus_{j=1}^{12}Q_jCQ_j^{-1}.
 \label{eq:leech-c3}
\end{equation}
Los símbolos \(Q_j\) designan marcos de planos; el símbolo \(P_3\)
queda reservado para el proyector sobre el sector tridimensional.

\begin{proposition}[Simetría de orden tres del retículo marcado]
\label{prop:leech-c3}
La transformación \(g_c\) es una isometría de orden tres de \(E\),
tiene como único vector fijo el cero y satisface
\[
 g_cN=N,\qquad g_cN_{\alpha,0}=N_{\alpha,0},
 \qquad g_c\Lambda=\Lambda.
\]
\end{proposition}
\begin{proof}
Cada bloque es \(C\) o \(C^{-1}\), cuyo polinomio mínimo es
\(X^2+X+1\). Esto prueba el orden y la afirmación sobre vectores fijos,
así como la conservación de \(B\). Cada bloque actúa trivialmente en
el discriminante y conserva, por tanto, todas las clases de pegado de
\(N\).

Escribamos \(v=v_\alpha=\sum_j a_j\rho_j\), con \(a_1=4\) y
\(a_j=1\) para \(j>1\). Todos los \(a_j\) son congruentes con uno
módulo tres. Las identidades
\[
 (C-I_2)\rho=(-2,-1)^{\mathsf T},
 \qquad (C^{-1}-I_2)\rho=(-1,-2)^{\mathsf T}
\]
muestran que
\[
 y=\frac{g_cv-v}{3},\qquad z=\frac{g_c^{-1}v-v}{3}
\]
tienen clases discriminantes \(c_*\) y \(-c_*\), respectivamente.
Ambas pertenecen a \(C_W\), luego \(y,z\in N\). En cada factor,
\[
 \left\langle
 \frac{(C^{\pm1}-I_2)a_j\rho}{3},a_j\rho
 \right\rangle_{B_2}=-a_j^2.
\]
Por suma, \(\langle y,v\rangle_B=\langle z,v\rangle_B=-27\);
por ello \(y,z\in N_{\alpha,0}\).
Si \(x\in N_{\alpha,0}\), entonces
\[
 \langle g_cx,v\rangle_B
  =\langle x,g_c^{-1}v\rangle_B
  =\langle x,v\rangle_B+3\langle x,z\rangle_B
  \equiv0\pmod3.
\]
La misma operación con \(g_c^{-1}\) demuestra la igualdad
\(g_cN_{\alpha,0}=N_{\alpha,0}\). Finalmente,
\(g_c(v/3)=v/3+y\) conserva la clase que genera
\(\Lambda=N_{\alpha,0}+\mathbb Z(v/3)\).
\end{proof}

La simetría ternaria actúa, por consiguiente, sobre el retículo
pegado completo. Su conservación utiliza la palabra del código y
la congruencia del vector marcado, además de la rotación de cada
plano. Esta precisión permite transportar la simetría al flujo sin
perder la información aritmética que hizo posible su construcción.

\subsection{Flujo anisótropo y conservación del covolumen}
\label{subsec:leech-flow}

\begin{definition}[Deformación determinada por la polarización]
Para \(s\in\mathbb R\), parámetro adimensional, se define
\begin{equation}
 H_u=\bigoplus_{j=1}^{12}u_jI_2,\qquad
 D_s=\exp(sH_u)=\bigoplus_{j=1}^{12}e^{su_j}I_2,\qquad
 \Lambda_s=D_s\Lambda.
 \label{eq:leech-flow}
\end{equation}
La coordenada \(u_j\) se asigna al plano \(E_j\) de la misma carta
marcada.
\end{definition}

El factor positivo \(e^{su_j}\) conserva la orientación de cada
plano. Las doce etiquetas transmiten la correspondencia entre la
polarización y la realización reticular. Esta correspondencia
forma parte de la estructura de partida; una permutación de las
etiquetas se trata mediante su transporte simultáneo en ambos objetos.

\begin{theorem}[Covolumen y simetría del flujo]
\label{thm:leech-flow}
Para todos \(s,t\in\mathbb R\) se cumple
\[
\begin{gathered}
 D_{s+t}=D_sD_t,\qquad D_0=I,\qquad D_s^{-1}=D_{-s},\\
 D_s^\dagger=D_s,\qquad D_sg_c=g_cD_s,\qquad \det D_s=1,
\end{gathered}
\]
donde \(A^\dagger=B^{-1}A^{\mathsf T}B\).
Cada \(\Lambda_s\) es una red discreta de rango \(24\) y covolumen
uno, preservada por la isometría \(g_c\) de orden tres.
\end{theorem}
\begin{proof}
Las identidades de grupo se verifican en cada bloque exponencial.
Los bloques escalares son autoadjuntos para \(B_2\), y conmutan
con los bloques de \(g_c\). La suma nula de
\eqref{eq:leech-u-norm} da
\[
 \det D_s=\prod_{j=1}^{12}e^{2su_j}
 =\exp\left(2s\sum_{j=1}^{12}u_j\right)=1.
\]
La imagen de una red por un operador invertible es una red del
mismo rango; su covolumen se multiplica por el valor absoluto del
determinante. El covolumen de \(\Lambda\) es uno, de modo que el
de \(\Lambda_s\) también lo es. La conmutación y
la proposición~\ref{prop:leech-c3} dan
\[
 g_c\Lambda_s=g_cD_s\Lambda=D_sg_c\Lambda=D_s\Lambda.
\]
El orden y el espacio de puntos fijos de \(g_c\) permanecen
inalterados.
\end{proof}

\begin{remark}[Parámetro y simetría conservada]
La dirección \(u\) queda determinada por el registro; \(s\) recorre
una familia de realizaciones. Reemplazar \(u\) por \(u/\|u\|\)
reparametriza la misma familia. Una interpretación física que
seleccione un valor particular de \(s\) añade su regla de realización.
El teorema conserva la simetría explícita \(g_c\); el transporte de
otros automorfismos requiere sus correspondientes relaciones con
\(H_u\).
\end{remark}

\subsection{Dual euclídeo y deformación efectiva}
\label{subsec:leech-dual}

\begin{definition}[Dual euclídeo]
Para una red \(L\subset E\), su dual respecto de \(B\) es
\[
 L^*=\{y\in E:\langle y,x\rangle_B\in\mathbb Z
                         \text{ para todo }x\in L\}.
\]
\end{definition}

\begin{theorem}[Reciprocidad del dual]
\label{thm:leech-dual}
La familia \eqref{eq:leech-flow} satisface
\begin{equation}
 \Lambda_s^*=\Lambda_{-s}\qquad(s\in\mathbb R).
 \label{eq:leech-dual}
\end{equation}
Su autodualidad en un parámetro \(s\) equivale a
\(D_{2s}\Lambda=\Lambda\).
\end{theorem}
\begin{proof}
Para un operador invertible \(A\),
\[
 (AL)^*=A^{-\dagger}L^*.
\]
En efecto, \(\langle y,Ax\rangle_B=\langle A^\dagger y,x\rangle_B\)
es entero para todo \(x\in L\) precisamente cuando
\(A^\dagger y\in L^*\). Tomando \(A=D_s\), la autoadjunción,
la identidad \(D_s^{-1}=D_{-s}\) y \(\Lambda^*=\Lambda\) prueban
\eqref{eq:leech-dual}. La igualdad
\(D_s\Lambda=D_{-s}\Lambda\) equivale, multiplicando por \(D_s\),
a \(D_{2s}\Lambda=\Lambda\).
\end{proof}

El covolumen controla el volumen de un dominio fundamental; la
autodualidad controla todos los productos enteros con la red.
La primera propiedad es continua a lo largo de esta familia. La
segunda impone la condición aritmética exacta del teorema anterior.
Se empleará por ello la expresión \emph{covolumen uno} para el
flujo euclídeo y se especificará la integralidad cuando intervenga
la unimodularidad aritmética.

\begin{proposition}[Testigo de pérdida local de integralidad]
\label{prop:leech-nonintegral}
Existe \(\delta>0\) tal que, si \(0<|s|<\delta\), la red
\(\Lambda_s\) es no integral y no es isométrica al retículo de Leech.
\end{proposition}
\begin{proof}
El vector \(v_\alpha/3\) pertenece a \(\Lambda\). Su norma
transportada es
\begin{equation}
 f(s)=\left\|D_s\frac{v_\alpha}{3}\right\|_B^2
 =\frac29\left(16e^{2su_1}+
                       \sum_{j=2}^{12}e^{2su_j}\right).
 \label{eq:leech-witness}
\end{equation}
Por \(\sum_j u_j=0\) y por la primera coordenada de
\eqref{eq:leech-u-coordinates},
\begin{equation}
\begin{split}
 f(0)&=6,\\
 f'(0)&=\frac49\left(16u_1+\sum_{j=2}^{12}u_j\right)
       =\frac{20}{3}u_1
       =-825-\frac{466}{3}\sqrt5\ne0.
\end{split}
\label{eq:leech-witness-derivative}
\end{equation}
La continuidad de \(f'\) permite elegir un intervalo
\((-\delta,\delta)\) en el que \(f\) sea estrictamente monótona.
Reduciendo \(\delta\), se obtiene además
\(11/2<f(s)<13/2\). Para \(0<|s|<\delta\), resulta \(f(s)\ne6\);
en ese intervalo, seis es el único entero. Por tanto, \(\Lambda_s\)
contiene un vector de norma no entera. Una red integral tiene
norma entera en todos sus vectores. Puesto que Leech es integral
y la norma se conserva por isometrías, queda demostrada también
la segunda afirmación.
\end{proof}

La deformación preserva volumen y simetría ternaria al tiempo que
modifica efectivamente la estructura métrica aritmética. Los pares
opuestos de coordenadas de \(u\) definen una permutación del espacio
ambiente que intercambia dilataciones y contracciones. Su promoción
a un automorfismo del retículo pegado exigiría conservar el código
y el vecino. La reciprocidad del dual ya demostrada utiliza
únicamente la autoadjunción del flujo y la autodualidad inicial.

\subsection{Convergencia y reciprocidad de las funciones theta}
\label{subsec:leech-theta}

Para \(s\in\mathbb R\) y \(t>0\), definimos
\begin{equation}
 \Theta_s(t)=\sum_{\lambda\in\Lambda}
                \exp\bigl(-\pi t\|D_s\lambda\|_B^2\bigr).
 \label{eq:leech-theta-real}
\end{equation}
Es la suma gaussiana de todas las longitudes de la red deformada,
con sus multiplicidades. Si \(M=\max_j|u_j|\), se tiene
\begin{equation}
 e^{-2|s|M}\|x\|_B^2
 \leq\|D_sx\|_B^2
 \leq e^{2|s|M}\|x\|_B^2.
 \label{eq:leech-quadratic-bounds}
\end{equation}
Sobre un compacto de \(\mathbb R\times(0,\infty)\), la cota
inferior proporciona una gaussiana dominante sobre la red fija.
La serie converge absoluta y uniformemente. Las derivadas de
orden finito añaden factores polinomiales en las coordenadas de
\(\lambda\); dichos factores quedan dominados por otra gaussiana.
Las derivaciones término a término son, por consiguiente,
legítimas en esos compactos.

\begin{theorem}[Reciprocidad gaussiana]
\label{thm:leech-theta}
Con el volumen euclídeo asociado a \(B\),
\begin{equation}
 \Theta_s(t)=t^{-12}\Theta_{-s}(t^{-1}),
 \qquad s\in\mathbb R,\quad t>0.
 \label{eq:leech-theta-reciprocity}
\end{equation}
En particular, \(\Theta_s(1)=\Theta_{-s}(1)\).
\end{theorem}
\begin{proof}
Fijamos la transformada de Fourier
\[
 \widehat f(y)=\int_E
      f(x)e^{-2\pi i\langle x,y\rangle_B}\,dx.
\]
Para una red \(L\subset E\) y
\(f_t(x)=e^{-\pi t\|x\|_B^2}\), periodizamos
\[
 F_t(x)=\sum_{\ell\in L}f_t(x+\ell).
\]
La convergencia normal de la serie y sus derivadas da una función
suave sobre \(E/L\). Si \(\xi\in L^*\), su coeficiente de Fourier es
\[
\begin{split}
 c_\xi
 &=\frac1{\operatorname{covol}(L)}
   \int_{E/L}F_t(x)e^{-2\pi i\langle x,\xi\rangle_B}\,dx\\
 &=\frac{\widehat f_t(\xi)}{\operatorname{covol}(L)}.
\end{split}
\]
La segunda igualdad descompone \(E\) en traslaciones de un dominio
fundamental; el carácter tiene valor uno sobre \(L\). El producto
de las integrales gaussianas unidimensionales, en una base
ortonormal para \(B\), da
\[
 \widehat f_t(\xi)=
 t^{-12}\exp\bigl(-\pi\|\xi\|_B^2/t\bigr).
\]
La serie de Fourier es absolutamente convergente. Evaluándola en
\(x=0\), obtenemos la suma de Poisson
\[
 \sum_{\ell\in L}e^{-\pi t\|\ell\|_B^2}
 =\frac{t^{-12}}{\operatorname{covol}(L)}
      \sum_{\xi\in L^*}e^{-\pi\|\xi\|_B^2/t}.
\]
Ahora \(L=\Lambda_s\) tiene covolumen uno y
\(L^*=\Lambda_{-s}\) por el teorema~\ref{thm:leech-dual}.
La sustitución prueba \eqref{eq:leech-theta-reciprocity}.
El caso \(t=1\) es inmediato.
\end{proof}

La periodización anterior es la demostración gaussiana de Poisson;
puede cotejarse con \cite[teorema 2, pp.~10--11]{Elkies}.
La convención con signo positivo en Fourier conduce a la misma
transformada de la gaussiana, por su paridad. La igualdad de theta
afecta a las series completas y conserva el dominio analítico de
convergencia.

\begin{proposition}[Extensión holomorfa]
\label{prop:leech-theta-complex}
Para \(\tau\) en el semiplano superior
\(\mathbb H=\{\tau\in\mathbb C:\operatorname{Im}\tau>0\}\), la serie
\begin{equation}
 \theta_s(\tau)=\sum_{\lambda\in\Lambda}
             e^{\pi i\tau\|D_s\lambda\|_B^2}
 \label{eq:leech-theta-complex}
\end{equation}
converge normalmente y satisface
\begin{equation}
 \theta_s(-1/\tau)=(-i\tau)^{12}\theta_{-s}(\tau).
 \label{eq:leech-theta-modular}
\end{equation}
\end{proposition}
\begin{proof}
El valor absoluto de cada sumando es
\(e^{-\pi\operatorname{Im}\tau\|D_s\lambda\|_B^2}\).
Sobre compactos del semiplano superior, la parte imaginaria tiene
una cota inferior positiva, de modo que
\eqref{eq:leech-quadratic-bounds} prueba convergencia normal y
holomorfía. Para \(\tau=it\), con \(t>0\),
\eqref{eq:leech-theta-reciprocity} aplicada en \(t^{-1}\) da
\(\theta_s(i/t)=t^{12}\theta_{-s}(it)\).
Ambos miembros de \eqref{eq:leech-theta-modular} son holomorfos
en \(\mathbb H\); el teorema de identidad prolonga la igualdad
desde el eje imaginario positivo.
\end{proof}

La inversión modular relaciona los parámetros \(s\) y \(-s\).
La periodicidad adicional bajo \(\tau\mapsto\tau+1\) utiliza la
paridad de las normas. La proposición~\ref{prop:leech-nonintegral}
exhibe parámetros arbitrariamente próximos a cero para los que
la integralidad euclídea se pierde. Por ello la familia conserva
la reciprocidad analítica demostrada, mientras una realización
mediante formas modulares escalares o álgebras de vértices
reticulares incorpora además las hipótesis aritméticas
correspondientes. En la realización excepcional inicial, esas
hipótesis sostienen la continuación hacia las álgebras de vértices
y Monstrous Moonshine \cite{HMT-IV}.

\subsection{Forma partida y autodualidad en dimensión cuarenta y ocho}
\label{subsec:leech-split}

La pérdida de integralidad de una proyección euclídea conduce a
considerar ambas proyecciones de manera simultánea. En el espacio
\(E\oplus E\) se introduce la forma bilineal
\begin{equation}
 \mathcal B((x,p),(x',p'))=
       \langle x,p'\rangle_B+\langle p,x'\rangle_B.
 \label{eq:leech-split-form}
\end{equation}
Sus productos emparejan las dos coordenadas. Definimos
\begin{equation}
\begin{split}
 \mathcal D_s&=D_s\oplus D_{-s},\\
 \Gamma_s&=\mathcal D_s(\Lambda\oplus\Lambda)\\
 &=\{(D_s\lambda,D_{-s}\mu):\lambda,\mu\in\Lambda\}.
\end{split}
\label{eq:leech-split-lattice}
\end{equation}

\begin{theorem}[Autodualidad integral para la forma partida]
\label{thm:leech-split}
La forma \(\mathcal B\) tiene signatura \((24,24)\).
Para todo \(s\in\mathbb R\), \(\mathcal D_s\) es una isometría de
\(\mathcal B\), y \(\Gamma_s\) es una red integral, par y autodual
respecto de esta forma.
\end{theorem}
\begin{proof}
Introduciendo
\[
 p_L=\frac{x+p}{\sqrt2},\qquad
 p_R=\frac{x-p}{\sqrt2},
\]
se obtiene
\[
 \mathcal B((x,p),(x,p))=\|p_L\|_B^2-\|p_R\|_B^2.
\]
Como \(B\) es positiva en dimensión \(24\), la signatura es
\((24,24)\). La autoadjunción y reciprocidad de los bloques dan
\[
 \langle D_sx,D_{-s}p'\rangle_B=\langle x,p'\rangle_B,
\]
y análogamente para el segundo producto cruzado. Por ello
\(\mathcal B(\mathcal D_s\xi,\mathcal D_s\eta)=\mathcal B(\xi,\eta)\).

En \(\Gamma_0=\Lambda\oplus\Lambda\), todos los productos cruzados
son enteros y
\[
 \mathcal B((\lambda,\mu),(\lambda,\mu))
   =2\langle\lambda,\mu\rangle_B\in2\mathbb Z.
\]
Esto prueba integralidad y paridad. Si \((x,p)\) pertenece al dual
de \(\Gamma_0\) para \(\mathcal B\), el emparejamiento con
\((\lambda,0)\) y \((0,\mu)\) exige respectivamente
\(p,x\in\Lambda^*=\Lambda\). El dual es, por tanto, \(\Gamma_0\).
La isometría \(\mathcal D_s\) transporta estas tres propiedades
a \(\Gamma_s\).
\end{proof}

La métrica especifica el contenido de la conclusión. La familia
\(\Gamma_s\) conserva la integralidad para los productos cruzados
de \(\mathcal B\), incluso cuando \(\Lambda_s\) pierde integralidad
para \(B\). Como retículos abstractos con forma partida, los
\(\Gamma_s\) son isométricos. Sus dos proyecciones euclídeas,
consideradas con las métricas positivas respectivas, varían con
el parámetro.

\begin{proposition}[Intercambio de las proyecciones recíprocas]
\label{prop:leech-split-involution}
La involución \(\mathcal J(x,p)=(p,x)\) satisface
\[
 \mathcal J\Gamma_s=\Gamma_{-s},\qquad
 \mathcal J\mathcal D_s\mathcal J=\mathcal D_{-s}.
\]
En las coordenadas \((p_L,p_R)\), fija \(p_L\) y cambia el signo
de \(p_R\).
\end{proposition}
\begin{proof}
El intercambio convierte
\((D_s\lambda,D_{-s}\mu)\) en
\((D_{-s}\mu,D_s\lambda)\), que recorre \(\Gamma_{-s}\).
La conjugación de los dos bloques da la segunda identidad.
Finalmente, \(x+p\) se conserva y \(x-p\) cambia de signo.
\end{proof}

La involución expresa una dualidad de la construcción reticular
completa. Su eventual representación física conserva, además de
la forma partida, los operadores y los datos de identificación
propios de esa realización.

\subsection{Energía reticular y equivalencia unitaria}
\label{subsec:leech-energy}

La energía de radio recíproco
\(m^2/R^2+w^2R^2\) intercambia sus dos términos bajo
\((m,w,R)\mapsto(w,m,R^{-1})\), con una conjugación unitaria
que preserva los dominios operatorios \cite{HMT-IV}.
La realización sobre los doce planos marcados emplea
\begin{equation}
 \mathcal E_s(\lambda,\mu)=
       \|D_s\lambda\|_B^2+\|D_{-s}\mu\|_B^2,
 \qquad \lambda,\mu\in\Lambda.
 \label{eq:leech-energy}
\end{equation}
La energía es positiva y cumple
\(\mathcal E_s(\lambda,\mu)=\mathcal E_{-s}(\mu,\lambda)\).

\begin{theorem}[Operador positivo y conjugación de dominios]
\label{thm:leech-energy}
En el espacio de Hilbert \(\ell^2(\Lambda\oplus\Lambda)\), sea
\[
 (H_s\psi)(\lambda,\mu)=
       \mathcal E_s(\lambda,\mu)\psi(\lambda,\mu)
\]
con dominio
\begin{equation}
 \mathcal D(H_s)=
 \left\{\psi\in\ell^2(\Lambda\oplus\Lambda):
 \sum_{\lambda,\mu}\mathcal E_s(\lambda,\mu)^2
             |\psi(\lambda,\mu)|^2<\infty\right\}.
 \label{eq:leech-energy-domain}
\end{equation}
Entonces \(H_s\) es positivo, autoadjunto y de resolvente compacto.
La aplicación \(U\psi(\lambda,\mu)=\psi(\mu,\lambda)\) es unitaria y
\begin{equation}
 U\mathcal D(H_s)=\mathcal D(H_{-s}),\qquad
 UH_sU^{-1}=H_{-s}.
 \label{eq:leech-unitary}
\end{equation}
\end{theorem}
\begin{proof}
La multiplicación por una función real en un conjunto discreto,
con el dominio máximo \eqref{eq:leech-energy-domain}, es
autoadjunta. Puede comprobarse directamente mediante la base
ortonormal de los vectores de soporte unitario: la pertenencia
al dominio del adjunto exige que la sucesión obtenida al
multiplicar por \(\mathcal E_s\) sea cuadrado sumable, exactamente
la misma condición. Además,
\[
 \langle\psi,H_s\psi\rangle
   =\sum_{\lambda,\mu}\mathcal E_s(\lambda,\mu)
                         |\psi(\lambda,\mu)|^2\geq0.
\]
Por \eqref{eq:leech-quadratic-bounds},
\[
 \mathcal E_s(\lambda,\mu)\geq
 e^{-2|s|M}\bigl(\|\lambda\|_B^2+\|\mu\|_B^2\bigr).
\]
Una red tiene finitos puntos en cualquier conjunto acotado; por
ello cada subnivel de \(\mathcal E_s\) es finito. Los coeficientes
diagonales \((1+\mathcal E_s)^{-1}\) tienden a cero fuera de
conjuntos finitos, de modo que \((H_s+I)^{-1}\) es límite en norma
de operadores de rango finito y es compacto.

El intercambio de los dos índices conserva la norma de
\(\ell^2(\Lambda\oplus\Lambda)\), por lo que \(U\) es unitario
y \(U^{-1}=U\). La igualdad de energías al invertir \(s\) transforma
la condición de dominio de \(H_s\) en la de \(H_{-s}\), y la
evaluación sobre cada par \((\lambda,\mu)\) prueba
\eqref{eq:leech-unitary}.
\end{proof}

Quedan así articuladas tres conservaciones de distinto dominio.
En los planos euclídeos, el flujo mantiene covolumen y simetría
ternaria y relaciona cada red con el dual del parámetro opuesto.
En el espacio duplicado, los productos cruzados conservan una
red integral, par y autodual. En la realización operatoria, el
intercambio recíproco conserva la energía mediante una equivalencia
unitaria de operadores con sus dominios. Las fórmulas theta
expresan la misma reciprocidad en la suma convergente de las
longitudes de toda la red.


% SOURCE appendices/mass_bridge.tex
% Procedencia primaria (lectura de 18-09-2026):
% XROOT = output/EXTREMA_MEDIA_RAZON_NARRACION_Y_REVISION_20260916/ARTICULO/ES/source
% XROOT/libro/01_acoplamiento.tex:51-147,727-749.
% XROOT/sections/k_direccion_dimensional.tex:54-155.
% XROOT/nuclear/12.tex:9-88.
% VIROOT = output/AMPLIACION_FORMAL_LEAN_SERIE_20260916/delivery/USB_ES_EN_STAGE10_20260917/ES/PAQUETES/06_PARTICULAS_Y_MASAS_ES/documentacion_original/payload/source_es
% VIROOT/sections/05_registro_operador_masa.tex:11-119,193-294.
% VIROOT/sections/02_particula_persistencia.tex:210-235.
% XROOT/nuclear/07.tex:5-25,74-154: Hamiltoniano de transporte eliptico.
% XROOT/nuclear/11.tex:878-979: elevacion cotangente y accion discreta.
% Estatutos: cadena K->alpha->accion->masa = RESULTADO_RECUPERADO;
% carta 1+1+10 de eventos y criterio de entrelazamiento = FORMALIZACION_NUEVA;
% congruencia, espectro generalizado y control racional = pruebas completas aqui.

\section{Congruencias positivas, lectores de acción y dinámica hamiltoniana}
\label{sec:k-mass-bridge}

La relación entre el registro dodecafásico y la Ley General de Masas tiene
una composición efectiva. APP produce las dos hojas con sus residuos y
cocientes; TRIT orienta la transición y determina su régimen; TPK compone
selección, transporte y actualización conservando ruta, acarreo, frontera,
supervivencia y memoria. Del estado enriquecido y de su prolongación
compatible proceden conjuntamente las publicaciones regionales y el
registro de doce posiciones. La clausura numérica y la lectura material
comparten ese origen y conservan sus operaciones específicas.

Con origen de fase, orientación y base mil declarados, el registro es
\[
 K=(234,543,140,729,659,824,621,58,914,794,146,601).
\]
El lector posicional utiliza
\begin{equation}
 \kappa_{\rm ag}(K)=\frac{\sum_{j=1}^{12}K_j1000^{12-j}}{1000^{12}-1},
 \qquad
 \alpha=\pi_{\rm HMT}+e_{\rm HMT}-\varphi_{\rm HMT}-4-
              \kappa_{\rm ag}(K).
 \label{eq:mass-k-alpha}
\end{equation}
Los valores de la derecha son publicaciones internas del mismo proceso
coinductivo. La fórmula expresa el lector de clausura infinita compatible,
con la frontera de acarreo que lo define.

La acción recibe esas publicaciones mediante los lectores
\begin{align}
 H_5={}&\frac12\sqrt{1000\alpha/\varphi}-\alpha+
       \frac9{16}\alpha^2-\frac59\alpha^3+
       \frac7{48}\alpha^4-\frac1{54}\alpha^5,\notag\\
 D_A={}&\exp(-100\pi\alpha/9),\qquad
 \mathcal C_\pi=169\alpha^6+\frac{D_A\alpha^7}{1-D_A\alpha},\notag\\
 \eta_{\rm ret}={}&H_5-\frac{90}{\pi}\mathcal C_\pi,
 \qquad R_{\rm act}=\frac{H_5}{\eta_{\rm ret}}>0.
 \label{eq:mass-k-action}
\end{align}
Aquí y en las fórmulas siguientes, \(\pi,\varphi,e\) abrevian las
publicaciones HMT ya construidas. La norma determinantal local
\(\Sigma_H=\varphi\alpha^{16}\) fija la década \(-34\), y la
sección \(\hbar_{\rm ret}=\eta_{\rm ret}10^{-34}\mathcal U_S\)
pertenece a la recta de acción declarada. La razón \(R_{\rm act}\)
conserva la relación entre las dos secciones de acción al cambiar su base
dimensional común.

La ruta electrónica central proporciona, mediante sus lectores de
desplazamiento y autocorrelación, el triple \((80,54,6)\). Con
\[
 \Delta_4=\frac{\pi-e\log\pi}{270}\left(1+\frac\pi{729}\right),
 \qquad
 b_e^{\rm alg}=\frac{\sqrt3}{4}
       \bigl[\exp(\varphi/\pi^2)-22\alpha^3\bigr](1+15\Delta_4),
\]
su publicación material contiene
\begin{equation}
 \kappa_e=80-54\alpha+6\Delta_4,
 \qquad \varepsilon_e^\beta=b_e^{\rm alg}R_{\rm act}^{\kappa_e}.
 \label{eq:mass-k-electron}
\end{equation}
La composición \eqref{eq:mass-k-alpha}--\eqref{eq:mass-k-electron}
identifica la intervención concreta de \(K\): su lectura de clausura
participa en \(\alpha\), y esa salida actúa después en el lector de
acción y en el exponente de retorno. La elección central del electrón
pertenece a la involución de la carta APP, cuyo único punto fijo es
\((5,5)\). La centralidad de esa celda y su evaluación escalar son
dos resultados enlazados por la ruta y su registro.

\subsection{Coordenadas reversibles del registro de eventos materiales}

El registro material asigna a una ruta \(\gamma\) doce eventos
orientados \(\mathbf e(\gamma)=(e_0,\ldots,e_{11})\), uno por cada
ventana nonádica del calendario de ciento ocho posiciones. El mismo
orden temporal permite realizar estos eventos en el portador de doce
coordenadas utilizado por \(K\). Esta identificación requiere conservar
el origen y la orientación del calendario.

Sean \(\mathbf1\) el vector uniforme,
\(P_{11}=I-\mathbf1\mathbf1^*/12\), y \(P_3\) el proyector
tridimensional de la carta incidencial declarada. La construcción
geométrica del registro determina
\[
 u_K=P_3P_{11}K,\qquad
 \|u_K\|^2=\frac{6638585+2275584\sqrt5}{20}>0,
 \qquad u_K\perp\mathbf1,
\]
y el proyector transversal
\[
 P_{10}(K)=P_{11}-\frac{u_Ku_K^*}{\|u_K\|^2}.
\]

\begin{proposition}[Descomposición material relativa a la dirección de \(K\)]
Para todo vector de eventos \(\mathbf e\), la aplicación
\begin{equation}
 \mathbf e\longmapsto
 \left(s_C=\mathbf1^*\mathbf e,\quad
       a_K=\frac{u_K^*\mathbf e}{\|u_K\|^2},\quad
       v_K=P_{10}(K)\mathbf e\right)
 \label{eq:mass-k-event-chart}
\end{equation}
es invertible sobre su imagen y satisface
\begin{equation}
 \mathbf e=\frac{s_C}{12}\mathbf1+a_Ku_K+v_K,
 \qquad
 \|\mathbf e\|^2=\frac{|s_C|^2}{12}
          +|a_K|^2\|u_K\|^2+\|v_K\|^2.
 \label{eq:mass-k-event-inverse}
\end{equation}
\end{proposition}
\begin{proof}
Los proyectores \(\mathbf1\mathbf1^*/12\),
\(u_Ku_K^*/\|u_K\|^2\) y \(P_{10}(K)\) son ortogonales dos a
dos y suman la identidad. Aplicar esa igualdad a \(\mathbf e\)
demuestra la inversión. La ortogonalidad de los tres sumandos da la
identidad de normas.
\end{proof}

Esta carta conserva el vector de eventos completo. Por ello conserva
también los funcionales que dependen de sus signos, sus sumas y sus
correlaciones, una vez realizada la reconstrucción
\eqref{eq:mass-k-event-inverse}. El registro de predecesores de los nueves,
la cuenta de acción y las coordenadas de torre mantienen su residencia
en la trayectoria enriquecida: la carta anterior modifica únicamente
las coordenadas de los doce eventos. El carácter se evalúa sobre la misma
firma recuperada. La construcción proporciona así una representación
geométrica de la información utilizada por la masa, con una dirección
distinguida por \(K\) y diez componentes transversales.

La cantidad de información que retiene cada publicación puede probarse
exactamente. En el dominio algebraico de registros, \(K'=K+\mathbf1\)
satisface
\[
 u_{K'}=u_K,\qquad P_{10}(K')=P_{10}(K),\qquad
 \kappa_{\rm ag}(K')-\kappa_{\rm ag}(K)=\frac1{999}.
\]
La última identidad resulta de sumar la progresión geométrica de los
doce pesos posicionales. Si las publicaciones regionales se mantienen
fijas, la lectura de \(\alpha\) cambia en \(-1/999\).
Este control algebraico determina el alcance exacto de la dirección:
su proyector conserva una polarización, mientras la clausura escalar
utiliza además la información uniforme y posicional del registro.
La admisibilidad de \(K'\) como otra historia terminal requiere su
propia construcción APP--TRIT--TPK.

\subsection{Congruencia positiva y espectro del operador material}

Para una multisección \(M\), sea \(\mathscr H_M\) el espacio finito
con base ortonormal etiquetada por sus rutas. El lector de retorno
\(r\in\{D,\Omega\}\) produce un operador autoadjunto
\(B_M^r e_\gamma=\kappa_\gamma^r e_\gamma\). El carácter refinado,
evaluado sobre diferencias respecto del electrón, da el operador basal
\[
 M_M^{(0)}e_\gamma=
 m_e^{(0)}\chi_{\rm ref}(\sigma_\gamma-\sigma_e,
                  \eta_\gamma-\eta_e)e_\gamma.
\]
Los coeficientes de torre y su convergencia pertenecen al lector de cada
ruta. Su publicación beta es la congruencia
\begin{equation}
 D_M^r=R_{\rm act}^{B_M^r/2},\qquad
 M_M^{\beta,r}=D_M^rM_M^{(0)}D_M^r.
 \label{eq:mass-k-congruence}
\end{equation}

\begin{proposition}[Invariantes exactos de la congruencia material]
El operador \(D_M^r\) es positivo e invertible. La congruencia
\eqref{eq:mass-k-congruence} conserva rango e inercia; en particular,
conserva la positividad. Si \(M_M^{(0)}\) y \(B_M^r\) son diagonales
en la base de rutas, entonces
\[
 m_\gamma^\beta=m_\gamma^{(0)}R_{\rm act}^{\kappa_\gamma^r}.
\]
Para dos rutas en la misma carta dimensional se obtiene
\[
 \frac{m_Y^\beta}{m_X^\beta}=
 \chi_{\rm ref}(\sigma_Y-\sigma_X,\eta_Y-\eta_X)
 R_{\rm act}^{\kappa_Y^r-\kappa_X^r}.
\]
\end{proposition}
\begin{proof}
El cálculo funcional da
\(D_M^r=\exp((\log R_{\rm act})B_M^r/2)\), con espectro
estrictamente positivo. Su inversa es la exponencial del opuesto.
Para todo \(v\), la forma cuadrática transformada vale
\((D_M^rv)^*M_M^{(0)}(D_M^rv)\). El cambio invertible de variable
preserva las dimensiones máximas de subespacios positivos, negativos
y nulos; éstas determinan la inercia y el rango. En una base común
diagonal, los dos factores aportan conjuntamente
\(R_{\rm act}^{\kappa_\gamma^r}\). Dividir dos valores cancela
la sección dimensional común y utiliza la ley multiplicativa del
carácter.
\end{proof}

La distinción entre transporte de una forma y modificación de un
operador respecto de una métrica fija es comprobable sin aproximaciones.
Para
\[
 M_0=\begin{pmatrix}1&0\\0&4\end{pmatrix},\qquad
 D=\begin{pmatrix}2&0\\0&1/2\end{pmatrix},\qquad
 M'=D^*M_0D=\begin{pmatrix}4&0\\0&1\end{pmatrix},
\]
el determinante de \(D\) es uno y los valores asociados a las dos
rutas se intercambian. Incluso el espectro no ordenado cambia si
\(M_0=I\): se transforma de \(\{1,1\}\) en \(\{4,1/4\}\).
Ambos controles preservan positividad, rango e inercia.

\begin{proposition}[Invarianza espectral con métrica transportada]
Sean \(G_0>0\) una métrica y \(D\) invertible. Si
\[
 M'=D^*M_0D,\qquad G'=D^*G_0D,
\]
los pares \((M',G')\) y \((M_0,G_0)\) tienen el mismo espectro
generalizado, con las mismas multiplicidades.
\end{proposition}
\begin{proof}
Se tiene
\[
 \det(M'-\lambda G')=|\det D|^2\det(M_0-\lambda G_0),
 \qquad (G')^{-1}M'=D^{-1}G_0^{-1}M_0D.
\]
La primera igualdad preserva las raíces y sus multiplicidades; la
segunda exhibe una semejanza de los operadores asociados.
\end{proof}

El flujo diagonal positivo inducido por \(K\) posee, por tanto, dos
realizaciones matemáticas precisas sobre una forma material. Cuando
transporta simultáneamente la forma y su métrica describe el mismo
objeto en otra carta. Cuando actúa sobre la forma conservando la métrica
de rutas, produce una familia operatoria que puede cambiar las masas
relativas. Su identificación física se decide por el lector de rutas y
por la acción declarada.

\subsection{Compatibilidad del flujo dodecafásico con el retorno material}

Sea
\[
 A_K=(\log\rho_A)\operatorname{diag}(\widehat u_{K,1},\ldots,
       \widehat u_{K,12}),\quad
 \widehat u_K=u_K/\|u_K\|,\quad \rho_A=\pi/\varphi^2,
\]
de modo que \(g_K(t)=e^{tA_K}\). La familia de retorno material
en la multisección \(M\) tiene generador
\[
 C_M^r=\tfrac12(\log R_{\rm act})B_M^r,\qquad
 S_M^r(t)=e^{tC_M^r}.
\]

\begin{proposition}[Transporte de la congruencia por un lector entrelazante]
Una aplicación lineal declarada
\(J_M:\mathbb C^{12}\to\mathscr H_M\) que cumple
\begin{equation}
 C_M^rJ_M=J_MA_K
 \label{eq:mass-k-intertwiner}
\end{equation}
satisface
\[
 S_M^r(t)J_M=J_Mg_K(t),
\]
y, para \(M_M(t)=S_M^r(t)M_M^{(0)}S_M^r(t)\),
\begin{equation}
 J_M^*M_M(t)J_M
 =g_K(t)^*\bigl(J_M^*M_M^{(0)}J_M\bigr)g_K(t).
 \label{eq:mass-k-pullback}
\end{equation}
\end{proposition}
\begin{proof}
La identidad \eqref{eq:mass-k-intertwiner} se itera a cada potencia
y se suma en la serie exponencial convergente. Tomar adjuntos y
sustituir en la forma material demuestra
\eqref{eq:mass-k-pullback}.
\end{proof}

El criterio es finito y admite una prueba de esfuerzo directamente
operatoria. En las bases diagonales declaradas, si \(a_j\) y
\(c_\gamma\) son los autovalores de \(A_K\) y \(C_M^r\),
respectivamente, se debe cumplir
\[
 (c_\gamma-a_j)(J_M)_{\gamma j}=0
 \qquad\text{para cada }(\gamma,j).
\]
Cada coeficiente no nulo del lector exige la igualdad correspondiente
de autovalores. La composición recuperada de
\eqref{eq:mass-k-alpha}--\eqref{eq:mass-k-electron} y la carta reversible
\eqref{eq:mass-k-event-chart} son resultados efectivos. La identificación
adicional de sus dos flujos requiere un lector \(J_M\) con la propiedad
\eqref{eq:mass-k-intertwiner}; esta proposición determina exactamente
qué debe comprobar una realización que la sostenga.

\subsection{Transporte elíptico y transformación regular de Legendre}

La realización variacional del transporte tiene un antecedente explícito
en la elipse de acción. El par angular ya publicado determina
\(\eta=\operatorname{artanh}(C^*/A)\), en la cámara
\(A\ne0\), \(|C^*/A|<1\). Las coordenadas \((Q,P)\), ambas de
dimensión raíz de acción, realizan la forma
\[
 \mathcal I_\eta(Q,P)=\frac12\bigl(e^{-\eta}Q^2+e^\eta P^2\bigr),
 \qquad \lambda=P\,dQ,\quad \Omega=dQ\wedge dP.
\]
El transporte entre formas, con
\(M_\eta=\operatorname{diag}(e^{\eta/2},e^{-\eta/2})\) y
\(J_2=\left(\begin{smallmatrix}0&-1\\1&0\end{smallmatrix}\right)\),
es \(M_{\eta'}e^{-\vartheta J_2}M_\eta^{-1}\). Una
interpolación \(\eta(t)\), \(\vartheta(t)\), con
\(\dot\vartheta=\omega(t)>0\), tiene el Hamiltoniano recuperado
del artículo X:
\begin{equation}
 H_{\rm tr}(Q,P,t)=\omega(t)\mathcal I_{\eta(t)}(Q,P)
                 +\frac{\dot\eta(t)}2QP.
 \label{eq:mass-k-hamiltonian}
\end{equation}
El segundo término procede de \(\dot M_\eta M_\eta^{-1}\).
Las ecuaciones
\[
 \dot Q=\omega e^\eta P+\frac{\dot\eta}2Q,
 \qquad
 \dot P=-\omega e^{-\eta}Q-\frac{\dot\eta}2P
\]
conservan \(\mathcal I_{\eta(t)}(Q(t),P(t))\): la variación
explícita de \(\eta\) y la contribución del término de conexión
se cancelan. Este resultado y la identidad
\(P\dot Q-H_{\rm tr}=p\dot q-\omega(q^2+p^2)/2\), para
\(Q=e^{\eta/2}q\), \(P=e^{-\eta/2}p\), están demostrados
en la sección de conservación variacional de la obra de extrema y
media razón holográfica.

\begin{proposition}[Lagrangiano regular del transporte elíptico]
Supongamos \(\eta\in C^2\), \(\omega\in C^1\) y \(\omega>0\).
La transformación de Legendre respecto de \(P\) es invertible y
produce
\begin{equation}
 L_{\rm tr}(Q,\dot Q,t)=
 \frac{\bigl(\dot Q-\dot\eta Q/2\bigr)^2}{2\omega e^\eta}
 -\frac{\omega e^{-\eta}}2Q^2.
 \label{eq:mass-k-lagrangian}
\end{equation}
\end{proposition}
\begin{proof}
Escribamos \(a=\dot\eta/2\), \(b=\omega e^\eta>0\) y
\(c=\omega e^{-\eta}\). Entonces
\(H_{\rm tr}=bP^2/2+aQP+cQ^2/2\), y
\[
 \dot Q=\partial_PH_{\rm tr}=bP+aQ,
 \qquad P=\frac{\dot Q-aQ}{b}.
\]
Al sustituir esta expresión en \(P\dot Q-H_{\rm tr}\), los
términos que contienen \(P\) se reúnen en
\(P(\dot Q-aQ)-bP^2/2=(\dot Q-aQ)^2/(2b)\), lo cual
demuestra \eqref{eq:mass-k-lagrangian}. Además,
\[
 \partial_{\dot Q}L_{\rm tr}=P,
 \qquad
 \partial_P^2H_{\rm tr}=\omega e^\eta>0,
 \qquad
 \partial_{\dot Q}^2L_{\rm tr}=\frac1{\omega e^\eta}>0.
\]
La transformación es regular y sus ecuaciones de Euler--Lagrange
equivalen a las ecuaciones hamiltonianas anteriores.
\end{proof}

Esta eliminación algebraica reúne el Hamiltoniano ya construido con
su forma lagrangiana de segundo orden. La regularidad procede del
término cuadrático de la realización elíptica. Para un levantamiento
cotangente puramente lineal \(H_K(q,p)=p^{\mathsf T}A_Kq\), la
ecuación \(\dot q=A_Kq\) es independiente de \(p\) y el Hessiano
\(\partial_p^2H_K\) es nulo. Su formulación variacional utiliza la
acción de primer orden \(\int p^{\mathsf T}(\dot q-A_Kq)\,dt\).
El corpus contiene también su antecedente discreto
\(L_n=p_{n+1}^{\mathsf T}(q_{n+1}-U_nq_n)\), cuyas ecuaciones
producen el transporte \(q_{n+1}=U_nq_n\) y el transporte dual
\(p_{n+1}=U_n^{-\mathsf T}p_n\). Las dos formulaciones poseen
dominios y operadores determinados. La interpolación temporal y la
sección de acción se conservan como parte de la realización física
declarada.


% SOURCE sections/lectores_formas_compatibles.tex
% Incorporación aditiva. No modifica las fuentes ni el PDF de 211 páginas.
% Procedencia, dependencias y alcance en preservation_map.json.
\section{Representación grassmanniana de los lectores y formas con valores en la recta de acción}
\label{sec:compat-lectores-formas}

Las realizaciones positivas conservan información que puede volver a ser
utilizada por los lectores físicos. Esta conservación tiene una formulación
precisa: una representación marcada del registro permite reconstruir sus
coordenadas y transportar, por composición, cada función que las utiliza.
La geometría representa un resultado de la construcción; su elección de
coordenadas conserva la procedencia de ese resultado.

El dominio de partida está formado por los estados producidos por APP,
TRIT y TPK en las secciones precedentes. APP conserva sus hojas aditiva y
multiplicativa mediante residuo y cociente; TRIT selecciona el régimen
local con su orientación; la composición
\(U_t=\operatorname{Upd}_t\circ\operatorname{Tra}_t\circ
\operatorname{Sel}_t\) actualiza posición, registro y memoria.
La prolongación
\[
 w_6\longrightarrow w_{12}\longrightarrow w_{18}\longrightarrow
 w_{24}\longrightarrow w_{30}\longrightarrow R_{36}\longrightarrow G_9
\]
conserva el prefijo y distingue retorno de fase de retorno del estado.
Las realizaciones espectral, solenoidal, gauge, cohomológica y determinantal
permanecen como operaciones correlativas de la estructura discreta
conjunta del continuo. La extracción del registro y las publicaciones
regionales se realizan en ese mismo dominio. En lo que sigue, todas ellas
son salidas HMT previamente construidas.

\subsection{La familia marcada y su inversa}

Sea \(L\geq2\), y sea \(K=(K_1,\ldots,K_L)\) un registro positivo
con carga \(Q=\sum_jK_j\). En la aplicación dodecafásica,
\(L=12\) y \(Q=6263\). Definimos
\[
 d_j=K_j/Q,\qquad t_0=0,\qquad
 t_j=\sum_{a=1}^{j}d_a,\qquad
 B(K)=\begin{pmatrix}1&\cdots&1\\t_0&\cdots&t_L\end{pmatrix}.
\]
La familia \(\mathcal G_L^{\rm aff}\) es la imagen de estos planos
de filas en \(\operatorname{Gr}_{>0}(2,L+1)\), conservando las
columnas ordenadas \(0,\ldots,L\), las marcas extremas y la carga
\(Q\). La normalización afín fija \(t_0=0,t_L=1\).
Escribimos \(\Delta_{ij}\) para el menor de las columnas \(i,j\)
de cualquier representante del plano.

\begin{proposition}[Inversa de la realización marcada]
\label{prop:compat-inversa-marcada}
La aplicación \(K\mapsto([B(K)],Q)\) es inyectiva. Su inversa
sobre la imagen está dada por
\begin{equation}
 K_j=Q\frac{\Delta_{j-1,j}}{\Delta_{0L}},
 \qquad
 t_j=\frac{\Delta_{0j}}{\Delta_{0L}},
 \quad 1\leq j\leq L.
 \label{eq:compat-inversa-plucker}
\end{equation}
Los cocientes son invariantes bajo cualquier cambio invertible de base
de las dos filas.
\end{proposition}
\begin{proof}
En el representante declarado,
\(\Delta_{ij}=t_j-t_i\), y \(\Delta_{0L}=1\).
Las fórmulas recuperan las diferencias consecutivas y sus sumas.
Si se sustituye \(B\) por \(AB\), con \(A\) invertible,
todos los menores se multiplican por \(\det A\); sus cocientes
permanecen iguales. La carga recupera la escala anterior a la
normalización. Por tanto, dos registros con iguales datos marcados
tienen las mismas coordenadas.
\end{proof}

La condición de pertenencia a \(\mathcal G_L^{\rm aff}\) importa.
Los cocientes consecutivos de un plano positivo arbitrario pueden
incumplir la identidad \(\sum_j\Delta_{j-1,j}=\Delta_{0L}\).
En la familia construida esa identidad es telescópica. La inversa actúa
en la imagen marcada que acaba de definirse.

\subsection{Álgebra de lectores y cambios de carta}

Sea \(\mathcal D\) la imagen de una publicación del estado enriquecido
de la forma \(x\mapsto(K(x),\xi(x))\). El símbolo \(\xi\)
reúne exactamente las restantes coordenadas utilizadas por el lector:
publicaciones regionales, origen y orientación, ruta, firma, condiciones
de frontera y secciones dimensionales, según el objeto considerado.
No se supone que estas coordenadas sean independientes. Definimos
\[
 \mathcal F:\mathcal D\longrightarrow\widetilde{\mathcal G},
 \qquad (K,\xi)\longmapsto([B(K)],Q,\xi),
 \qquad \widetilde{\mathcal G}:=\mathcal F(\mathcal D).
\]
La proposición anterior proporciona la inversa \(\mathcal I\) de
\(\mathcal F\) en su imagen.

\begin{theorem}[Representación fiel del álgebra de lectores retenidos]
\label{thm:compat-algebra-lectores}
Para cada lector \(\mathcal R:\mathcal D\to V\), con codominio
declarado \(V\), existe un único lector geométrico
\(\mathcal R^{\rm geom}:\widetilde{\mathcal G}\to V\) tal que
\[
 \mathcal R^{\rm geom}=\mathcal R\circ\mathcal I,
 \qquad \mathcal R^{\rm geom}\circ\mathcal F=\mathcal R.
\]
Para lectores escalares, esta correspondencia preserva sumas, productos,
unidad y conjugación. Es un isomorfismo del álgebra de funciones sobre
\(\mathcal D\) con el álgebra correspondiente sobre
\(\widetilde{\mathcal G}\).
\end{theorem}
\begin{proof}
La composición está definida porque \(\mathcal I\mathcal F\) es
la identidad. Si dos lectores geométricos tienen la misma composición
con \(\mathcal F\), coinciden en su imagen, que es todo
\(\widetilde{\mathcal G}\). Las operaciones algebraicas se verifican
punto a punto; por ejemplo,
\((\mathcal R_1\mathcal R_2)\circ\mathcal I=
(\mathcal R_1\circ\mathcal I)(\mathcal R_2\circ\mathcal I)\).
La composición inversa es la restricción por \(\mathcal F\).
\end{proof}

La fidelidad se refiere a la publicación \((K,\xi)\). Para un
observable \(\mathcal O\) del estado enriquecido completo, el criterio
de descenso es que \(\mathcal O(x)=\mathcal O(x')\) siempre que
\((K(x),\xi(x))=(K(x'),\xi(x'))\). La necesidad se obtiene por
composición; la suficiencia permite definir el valor en una fibra
mediante cualquiera de sus representantes. Este criterio conserva
explícitamente qué información debe acompañar al registro.

Sea ahora \(\mathcal T\) una triangulación del polígono etiquetado
que determina una carta de Plücker positiva. Se incluyen sus coordenadas
de frontera y una normalización homogénea, por ejemplo
\(\Delta_{0L}=1\). Si se intercambia la diagonal \(ac\) por
\(bd\), con \(a<b<c<d\), el cambio de carta es
\begin{equation}
 \Delta_{bd}=
 \frac{\Delta_{ab}\Delta_{cd}+\Delta_{ad}\Delta_{bc}}{\Delta_{ac}}.
 \label{eq:compat-cambio-carta}
\end{equation}
El denominador es positivo y la relación de Plücker prueba que ambas
cartas evalúan el mismo punto.

\begin{corollary}[Compatibilidad de los lectores con las mutaciones]
\label{cor:compat-mutaciones}
Sean \(r_{\mathcal T}\) y \(r_{\mathcal T'}\) las expresiones
de \(\mathcal R^{\rm geom}\) en dos cartas que conservan las marcas,
y \(\mu_{\mathcal T\mathcal T'}\) su cambio positivo. Entonces
\[
 r_{\mathcal T'}\circ\mu_{\mathcal T\mathcal T'}
 =r_{\mathcal T}
\]
en el dominio común. Una sucesión cerrada de cambios de carta que
conserva el punto marcado conserva todos estos lectores.
\end{corollary}
\begin{proof}
Las dos reconstrucciones mediante
\eqref{eq:compat-inversa-plucker} recuperan el mismo \(K\).
Las coordenadas \(\xi\) se transportan con sus marcas. Aplicar
\(\mathcal R\) prueba la igualdad; iterarla prueba la afirmación
sobre una sucesión cerrada.
\end{proof}

Los lectores pueden expresarse en cartas distintas y seguir siendo
funciones del mismo objeto. Una permutación de las etiquetas actúa
adicionalmente sobre la marca posicional; una transformación que cambia
el punto actúa sobre el argumento del lector. Ambas operaciones tienen
su ley de transporte específica.

La acción diagonal por columnas proporciona un control explícito.
Para \(\lambda_i>0\),
\(\Delta_{ij}\mapsto\lambda_i\lambda_j\Delta_{ij}\).
La razón doble \(\Delta_{ab}\Delta_{cd}/
(\Delta_{ad}\Delta_{bc})\) se conserva, mientras
\(\Delta_{j-1,j}/\Delta_{0L}\) recibe el factor
\(\lambda_{j-1}\lambda_j/(\lambda_0\lambda_L)\).
Con \(t=(0,1/4,1/2,1)\) y \(\lambda=(1,2,1,1)\), las
separaciones reconstruidas serían \((1/2,1/2,1/2)\), cuya suma
es \(3/2\). Las razones dobles solas conservan un cociente diferente
del requerido por la inversa marcada.

\subsection{Composición explícita con acción y respuesta constitutiva}

En la carta dodecafásica, escribimos \(b=1000\) y
\begin{equation}
 \kappa^{\rm geom}=
 \frac{Q}{b^{12}-1}\sum_{j=1}^{12}
 b^{12-j}\frac{\Delta_{j-1,j}}{\Delta_{0,12}},
 \qquad
 \alpha^{\rm geom}=\pi_{\rm HMT}+e_{\rm HMT}
 -\varphi_{\rm HMT}-4-\kappa^{\rm geom}.
 \label{eq:compat-alpha-geometrica}
\end{equation}
Son las expresiones geométricas del lector de clausura
\eqref{eq:mass-k-alpha}. Las tres publicaciones regionales y \(K\)
comparten generación coinductiva; la expresión utiliza su orden
constructivo interno. El lector de acción
\eqref{eq:mass-k-action}, evaluado en \(\alpha^{\rm geom}\),
determina \(\eta_{\rm ret}^{\rm geom}\), y por tanto
\[
 \hbar_{\rm ret}^{\rm geom}
 =\eta_{\rm ret}^{\rm geom}10^{-34}\mathcal U_S
 =\hbar_{\rm ret}.
\]
La igualdad tiene lugar en la recta de acción \(\mathcal L_S\)
con su sección \(\mathcal U_S\); conserva la década procedente
del lector determinantal \(\varphi_{\rm HMT}\alpha^{16}\).
La composición no modifica la sección ni la aplica una segunda vez.

En la carta angular levantada que conserva el origen y la rama real
de la publicación, el par angular posterior determina
\[
 x=\frac{\pi_{\rm HMT}}{180}(1000\alpha),\qquad
 y=\frac{\pi_{\rm HMT}}{90}(\eta_{\rm ret}+\alpha),\qquad
 q_\pm=e^{-x\mp y}.
\]
La rama forma parte de las marcas retenidas: la reducción de un ángulo
módulo \(360\) por sí sola conserva otra información. Las igualdades
que siguen emplean el levantamiento declarado, en el cual la primera
coordenada es \(1000\alpha\).
En la cámara \(x>|y|\), con proyectores de hoja fijados
\(P_\pm\), sea \(T=q_+P_++q_-P_-\). La respuesta del corpus
compara las incidencias \(90\) y \(120\) mediante
\[
 \mathcal V=(I-T^{90})(I-T^{120})^{-1}
 =r_+P_++r_-P_-,\qquad
 r_\pm=f(q_\pm^{30}),\quad
 f(s)=\frac{1+s+s^2}{1+s+s^2+s^3}.
\]
La regla satisface
\(\mathcal V\Sigma_4(T^{30})=\Sigma_3(T^{30})\), donde
\(\Sigma_j(S)=\sum_{a=0}^{j-1}S^a\). Su unicidad se deduce
de la invertibilidad de \(\Sigma_4(T^{30})\).
Las cuatro lecturas normalizadas son
\[
 \widehat\varepsilon=r_+^2,\quad
 \widehat\mu=r_-^2,\quad
 \widehat Z=r_-/r_+,\quad
 \widehat c=(r_+r_-)^{-1}.
\]
En las rectas dimensionales de acción, carga y velocidad se realizan como
\[
 Z=\widehat Z\,u_Su_Q^{-2},\quad c=\widehat c\,u_C,
 \qquad
 \mu=\widehat\mu\,u_Su_C^{-1}u_Q^{-2},\quad
 \varepsilon=\widehat\varepsilon\,u_S^{-1}u_C^{-1}u_Q^2.
\]
Las identidades \(\mu\varepsilon=c^{-2}\) y
\(\mu/\varepsilon=Z^2\) se obtienen por producto y cociente.
El corolario de cambios de carta conserva esta respuesta completa,
con las marcas de hoja y las secciones dimensionales incluidas en
\(\xi\). El mismo argumento se aplica a un lector de masa con su
ruta, firma y operador de retorno retenidos, en el dominio de la
sección~\ref{sec:k-mass-bridge}.

\begin{theorem}[Recuperación del registro desde la respuesta etiquetada]
\label{thm:compat-recuperacion-constitutiva}
Fijadas las publicaciones regionales, la orientación de las hojas, la
base \(b=1000\), el origen, la longitud doce y el levantamiento angular
declarado, la aplicación
\(K\mapsto(\widehat\varepsilon,\widehat\mu)\) es inyectiva
sobre los registros enteros de su imagen física que satisfacen
\(0\leq K_j<b\) y la cámara contractiva indicada. Su inversa exacta
es la composición de las siguientes operaciones:
\begin{align*}
 r_+&=\sqrt{\widehat\varepsilon},&
 r_-&=\sqrt{\widehat\mu},\\
 r_\pm s_\pm^3+(r_\pm-1)(1+s_\pm+s_\pm^2)&=0,
 &0<s_\pm<1,\\
 x&=-\frac1{60}\log(s_+s_-),&
 \alpha&=\frac{180x}{1000\pi_{\rm HMT}},\\
 \kappa&=\pi_{\rm HMT}+e_{\rm HMT}-\varphi_{\rm HMT}-4-\alpha,
 &N&=(b^{12}-1)\kappa,\\
 K_j&=\left\lfloor N/b^{12-j}\right\rfloor\bmod b.
\end{align*}
\end{theorem}
\begin{proof}
Las raíces positivas recuperan los coeficientes etiquetados.
La derivada
\[
 f'(s)=-\frac{s^2(3+2s+s^2)}{(1+s+s^2+s^3)^2}<0
 \qquad(0<s<1)
\]
y sus límites \(1\) y \(3/4\) prueban que la ecuación cúbica
posee exactamente una raíz en ese intervalo. Puesto que
\(s_\pm=e^{-30x\mp30y}\), su producto es \(e^{-60x}\).
Se recuperan \(x\), \(\alpha\) y \(\kappa\).
En la imagen considerada, \(N=\sum_jK_jb^{12-j}\) es un entero
entre cero y \(b^{12}-1\). La división euclídea sucesiva recupera
de manera única sus doce cifras, incluidos los ceros iniciales.
\end{proof}

La precisión necesaria para una recuperación numérica también puede
expresarse. Con publicaciones regionales exactas, un error
\(|\widetilde\alpha-\alpha|<1/[2(b^{12}-1)]\) implica
\(|\widetilde N-N|<1/2\); el redondeo al entero más próximo
recupera \(N\). Si las publicaciones regionales tienen errores, se
aplica la misma condición a la suma de errores de
\(\pi+e-\varphi-\alpha\). La inyectividad exacta y la resolución
metrológica son así propiedades diferenciadas.

El dominio entero es esencial. En una prolongación analítica de los
registros positivos, la perturbación de tres posiciones consecutivas
proporcional a \((1,-(b+1),b)\) conserva tanto la suma como la
lectura posicional, pues
\(1-(b+1)+b=0\) y \(b^2-b(b+1)+b=0\).
Para perturbaciones suficientemente pequeñas conserva la positividad
y cambia el punto grassmanniano. La inyectividad precedente utiliza
la condición de doce cifras enteras; la familia analítica continua
posee otras fibras.

\subsection{Refinamiento y lectores heredados}

Si cada separación \(d_j\) se subdivide en separaciones positivas
\(d_{j,a}\) con \(\sum_ad_{j,a}=d_j\), los extremos originales
conservan su posición. Sea \(\rho\) la operación de sumar los hijos
de cada intervalo y \(H\) la inclusión de sus extremos en el nuevo
conjunto ordenado. Entonces
\[
 \mathcal I_{\rm heredada}(B(\widetilde d),Q)
 =\bigl(Q\sum_ad_{j,a}\bigr)_j=K,
 \qquad
 \mathcal R_{\rm heredada}=\mathcal R(K,\xi).
\]
La igualdad prueba la conmutatividad entre reconstrucción heredada y
subdivisión. En una sucesión de subdivisiones, las sumas se asocian;
el resultado persiste en cada composición finita. Esto conserva el
lector de longitud doce y sus marcas. Un lector definido sobre los
nuevos bloques tendrá su propio dominio y su ley de comparación.

\subsection{Formas canónicas con valores en una recta dimensional}

Sea \(P\) uno de los politopos de rango uno construidos en el
manuscrito y \(P=\bigcup_\sigma P_\sigma\) una subdivisión
orientada del mismo soporte, con multiplicidad uno y caras compatibles.
Escribimos \(\Omega_\sigma\) y \(\Omega_P\) para sus formas
canónicas. Sus identidades son
\(\Omega_P=\sum_\sigma\Omega_\sigma\) y
\(\operatorname{Res}_F\Omega_P=\Omega_F\), con la convención
de orientación fijada en el desarrollo amplituhedral.
La recta de acción \(\mathcal L_S\) se considera constante sobre
este politopo; su sección procede del estado HMT fijado.

\begin{theorem}[Transporte de la forma por una sección de acción]
\label{thm:compat-forma-accion}
Para \(a\in\mathcal L_S\), constante respecto de las coordenadas
del politopo,
\[
 \boldsymbol\Omega_P:=a\otimes\Omega_P
 =\sum_\sigma a\otimes\Omega_\sigma,
 \qquad
 \operatorname{Res}_F\boldsymbol\Omega_P=a\otimes\Omega_F.
\]
Las identidades se conservan bajo subdivisiones sucesivas y residuos
iterados sobre banderas de caras. Se aplican, en particular, a
\(a=\hbar_{\rm ret}\).
\end{theorem}
\begin{proof}
El producto tensorial con una recta es lineal. El residuo actúa sobre
el factor diferencial y conmuta con el coeficiente constante. La
compatibilidad de subdivisiones y de residuos iterados se transporta
por las mismas operaciones lineales. En un cambio de base dimensional
\(u'_S=\lambda\,u_S\), con \(\lambda>0\), el coeficiente numérico
cambia por \(\lambda^{-1}\);
el tensor y sus residuos permanecen iguales.
\end{proof}

\begin{theorem}[Condición de pegado de los coeficientes de frontera]
\label{thm:compat-pegado}
Sean \(a_\sigma\) secciones de la misma recta, regulares en un
entorno de las caras y meromorfas en una realización ambiental común.
Supóngase que sus productos con las formas canónicas tienen, a lo
largo de las caras consideradas, polos a lo sumo simples. En el
interior relativo de una cara común \(F\) de codimensión uno de dos celdas
\(\sigma,\tau\), con orientaciones opuestas,
\begin{equation}
 \operatorname{Res}_F\left(a_\sigma\Omega_\sigma+a_\tau\Omega_\tau\right)
 =\left(a_\sigma|_F-a_\tau|_F\right)\Omega_F.
 \label{eq:compat-defecto-frontera}
\end{equation}
Por tanto, el polo de esta pareja sobre el divisor de \(F\) desaparece
exactamente cuando las dos trazas coinciden. Si los coeficientes son
constantes y el grafo de adyacencia de las celdas es conexo, la cancelación
de cada pareja de celdas a través de su cara común equivale a la existencia
de un único coeficiente común \(a\).
\end{theorem}
\begin{proof}
Con coordenada local transversal \(z\), una forma logarítmica se
escribe \(dz/z\wedge\omega+\eta\), donde \(\eta\) es regular
respecto de ese divisor. La regularidad de \(a_\sigma\) da
\(\operatorname{Res}_F(a_\sigma\Omega_\sigma)
=a_\sigma|_F\operatorname{Res}_F\Omega_\sigma\).
Las dos orientaciones producen \(\Omega_F\) y \(-\Omega_F\).
La forma \(\Omega_F\) es no nula en el interior relativo de la
cara; la igualdad de residuos equivale a la igualdad de las trazas.
Para un polo simple, la anulación del residuo elimina el polo.
En el caso constante, la igualdad se propaga a lo largo de cada
arista del grafo conexo de adyacencia.
\end{proof}

Para la suma global deben contarse todos los términos con un polo en el
mismo divisor ambiental. Si \(H\) es el hiperplano que contiene \(F\),
y \(I_H\) reúne esos términos, la identidad general es
\[
 \operatorname{Res}_{H}\left(\sum_\nu a_\nu\Omega_\nu\right)
 =\sum_{\nu\in I_H}a_\nu|_H\operatorname{Res}_H\Omega_\nu,
\]
siempre que los coeficientes sean regulares sobre \(H\). Una faceta
coplanar, aunque no sea adyacente a la cara real \(F\), participa en esta
suma. La fórmula de dos términos representa el residuo global cuando
los restantes términos son regulares sobre ese divisor. La cancelación
por parejas de todas las facetas interiores elimina sus contribuciones;
las facetas exteriores conservan los residuos de frontera correspondientes.

En una cara de codimensión superior, se aplica esta fórmula sucesivamente
con las orientaciones inducidas. Los signos son los de la incidencia
celular y satisfacen la cancelación de borde de borde. La familia de
lectores locales define un pegado compatible por parejas cuando sus
trazas coinciden después del transporte a la misma recta de coeficientes.
El residuo global se obtiene sumando las contribuciones de cada divisor.
Este enunciado da un criterio efectivo para distribuir una lectura
física sobre cartas o celdas sin introducir discontinuidades espurias.

\subsection{Cancelación de polos y normalización global}

La cancelación interior determina una condición local. La identificación
con una forma canónica global requiere asimismo el divisor de polos y
los residuos exteriores normalizados. En espacio proyectivo, dos
formas racionales de grado máximo con los mismos polos logarítmicos
y los mismos residuos difieren por una forma holomorfa de grado máximo;
\(H^0(\mathbb P^m,\Omega^m)=0\) demuestra su igualdad.
Ésta es la condición global utilizada en la prueba del politopo
\cite{positive}.

Un ejemplo intervalar muestra por qué los coeficientes variables
requieren dicha condición. Para \(0<c<1\), sean
\[
 \Omega_{0c}=\frac{c\,dx}{x(c-x)},\qquad
 \Omega_{c1}=\frac{(1-c)\,dx}{(x-c)(1-x)},\qquad
 \Omega_{01}=\frac{dx}{x(1-x)}.
\]
Con \(a_1(x)=1+x(c-x)\) y \(a_2(x)=1\), sus trazas coinciden
en todos los extremos pertinentes, y sin embargo
\[
 a_1\Omega_{0c}+a_2\Omega_{c1}=\Omega_{01}+c\,dx.
\]
El término adicional es regular en la carta afín y tiene un polo en
el infinito proyectivo. Los residuos de los extremos finitos y la
cancelación en \(c\) se conservan, mientras la forma resultante
tiene otro comportamiento global. Multiplicar el ejemplo por una
sección de \(\mathcal L_S\) produce el mismo control dimensional.

La consecuencia de conjunto es precisa: la representación grassmanniana
marcada transporta los lectores retenidos; las mutaciones conservan sus
valores cuando sólo cambian la carta; el refinamiento conserva sus
evaluaciones heredadas; y las formas con valores en la recta de acción
se ensamblan mediante igualdad de trazas y control de polos. Cada
igualdad expresa la operación que conserva el objeto correspondiente.


% SOURCE appendices/multimorphology_figure.tex
\begin{figure}[htbp]
\centering
\begin{tikzpicture}[
 box/.style={draw=black!70,rounded corners=1pt,align=center,
             font=\small,inner sep=5pt,text width=40mm},
 arrow/.style={-{Stealth[length=2mm]},thick}]
\node[box] (memory) at (0,0) {Carga y memorias\\\((q,m_0,m_1)\)};
\node[box] (register) at (0,-1.65) {Registro reconstruido\\\(K=(q/12)\mathbf1+Lm\)};
\node[box] (weights) at (-3.8,-3.55) {Partición ordenada\\\(d_j=K_j/q\)};
\node[box] (direction) at (3.8,-3.55) {Dirección de suma nula\\\(u=P_3P_{11}K\)};
\node[box] (geometry) at (-4.6,-6) {Medición de frontera\\y forma canónica\\\(C(d),\ \Omega_{P_K}\)};
\node[box] (energy) at (0,-6) {Energía nodal\\y detalles\\\(L_d,\ \mathcal E(z)\)};
\node[box] (lattice) at (4.6,-6) {Deformación reticular\\y red dual\\\(\Lambda_s^*=\Lambda_{-s}\)};
\draw[arrow] (memory)--(register);
\draw[arrow] (register)--(weights);
\draw[arrow] (register)--(direction);
\draw[arrow] (weights)--(geometry);
\draw[arrow] (weights)--(energy);
\draw[arrow] (direction)--(lattice);
\draw[arrow] (direction) to[bend left=12] node[right,font=\scriptsize,align=left,fill=white,inner sep=1pt]{flujo\\heredado} (energy);
\end{tikzpicture}
\caption{Realizaciones posteriores del registro recuperable. La rama
intervalar determina los menores de caminos, la forma canónica del
polígono y el laplaciano nodal. La dirección algebraica determina el
flujo de intervalos y, junto con la realización reticular marcada, la
familia dual de veinticuatro dimensiones. Cada flecha conserva el dominio
y las hipótesis de su construcción; los codominios inferiores son
distintos.}
\label{fig:multimorphology}
\end{figure}


% SOURCE sections/ym_bridge.tex
\section{Positividades, espacios físicos y transporte del espectro}
La geometría positiva y la teoría gauge reciben estructuras distintas de
un mismo proceso genealógico. La primera realiza los pesos y sus
incidencias mediante menores ordenados, celdas y formas canónicas.
La segunda conserva la incidencia gauge, la holonomía, las
representaciones, las rutas y la memoria dentro de una medida compatible.
Su forma de positividad relevante para la reconstrucción física es la
positividad por reflexión. El nexo conceptual adquiere contenido
matemático cuando se conservan los mapas que producen cada forma.

En la parte geométrica se ha demostrado que la subdivisión de una celda
conserva su forma canónica mediante cancelación de polos interiores.
Para las energías nodales, la reducción de Schur conserva una forma
efectiva y registra el detalle. En la parte gauge, la compatibilidad del
refinamiento debe conservar además el vacío, el espacio de invariantes,
el generador y el testigo espectral. Ésta es la cadena desarrollada a
continuación: escala aritmético--geométrica, medida común, semigrupo,
proyección física, refinamiento y reconstrucción continua.

La norma de un vector, la positividad de una matriz de momentos y la
existencia de una brecha de un Hamiltoniano poseen dominios diferentes.
Una congruencia con factor invertible conserva la positividad y la
inercia de una forma; un entrelazamiento isométrico del generador, con su
dominio, transporta su realización sobre la imagen reducida de la
inclusión. La cota uniforme y el testigo persistente fijan después la
brecha límite. Por ello el desarrollo mantiene visibles las
inclusiones y las identidades de semigrupo que justifican el paso al
límite. La escala de la brecha procede de la incidencia areal--volumétrica
y de su unidad dimensional declarada.

Los cuerpos siguientes reúnen la construcción proyectiva y su
desarrollo posterior dependiente del acoplamiento \(g\). Las primeras
comparaciones con una parametrización Gibbs se completan después mediante
la acción, la medida y el transporte explícitos. El orden expositivo
conserva las demostraciones del proceso y sitúa la formulación completa
tras sus dependencias. Las cuatro reflexiones corresponden a las cuatro
direcciones euclídeas de la reconstrucción. El desarrollo aquí reunido
tiene por objeto esa teoría gauge y su brecha, con el grupo compacto,
conexo y simple declarado en los teoremas.

El sector supersimétrico planar en el que se estudian amplitudes mediante
amplituedros conserva una pregunta física específica. Para enlazarlo con
la presente realización se requieren los operadores de supercarga, sus
dominios y la representación correspondiente. La dimensión cuatro de la
carta euclídea y la diferencia cuatro del bloque modular permanecen
resultados ya definidos de sus respectivos operadores.


% SOURCE sections/ym_scale.tex
\section{Transición areal--volumétrica y escala espectral}
\label{sec:ym-areal-volumetric-scale}

La constante que fija el semigrupo gauge se obtiene antes de introducir el
Hamiltoniano. Su procedencia combina dos lectores independientes del mismo
estado genealógico: la frecuencia de retorno que publica
\(\pi_{\rm HMT}\) y la distribución estacionaria de las dos hojas
areal--volumétrica.

\subsection{Dinámica de las dos hojas}

Sea
\begin{equation}
 F_{\rm av}=\begin{pmatrix}0&1\\1&1\end{pmatrix},
 \qquad e_{\rm ar}=\binom10,
 \qquad e_{\rm vol}=\binom01.
 \label{eq:ym-av-transition-matrix}
\end{equation}
La recurrencia de Fibonacci da, por inducción,
\begin{equation}
 F_{\rm av}^{n}
 =\begin{pmatrix}F_{n-1}&F_n\\F_n&F_{n+1}\end{pmatrix}.
 \label{eq:ym-av-fibonacci-powers}
\end{equation}
El vector de Perron es proporcional a \((1,\varphi)\). Como
\(F_{\rm av}\) es simétrica, la medida de Parry asigna a las dos hojas
pesos proporcionales a \((1,\varphi^2)\), y por tanto
\begin{equation}
 \frac{\mu_{\rm ar}}{\mu_{\rm vol}}=\varphi^{-2},
 \qquad
 \mu_{\rm vol}=\frac{\varphi^2}{1+\varphi^2}.
 \label{eq:ym-av-parry-weights}
\end{equation}
Aquí \(\mu_{\rm vol}\) es un carácter de frecuencia adimensional; no
denota un volumen de dimensión \(L^4\).

El lector regional de \(\pi\) se introduce una sola vez mediante
\begin{equation}
 D_\pi
 =\pi_{\rm HMT}|e_{\rm ar}\rangle\langle e_{\rm ar}|
  +|e_{\rm vol}\rangle\langle e_{\rm vol}|.
 \label{eq:ym-pi-regional-reader}
\end{equation}
La razón publicada a profundidad \(n\) es
\begin{equation}
 \Theta_n
 :=\frac{\langle e_{\rm ar},D_\pi F_{\rm av}^{n}e_{\rm ar}\rangle}
 {\langle e_{\rm vol},F_{\rm av}^{n}e_{\rm vol}\rangle}
 =\pi_{\rm HMT}\frac{F_{n-1}}{F_{n+1}}
 \longrightarrow
 r_{\rm av}:=\frac{\pi_{\rm HMT}}{\varphi_{\rm HMT}^{2}}.
 \label{eq:ym-target-free-av-limit}
\end{equation}
La matriz, la ruta y el cociente se fijan antes de consultar una brecha de
masa o un parámetro de teoría gauge. En la profundidad nativa \(n=108\),
\begin{equation}
 F_{107}=10284720757613717413913,
 \qquad
 F_{109}=26925748508234281076009,
 \label{eq:ym-fibonacci-107-109}
\end{equation}
y
\begin{equation}
 \left|\Theta_{108}
 -\frac{\pi_{\rm HMT}}{\varphi_{\rm HMT}^{2}}\right|
 <2\cdot10^{-45}.
 \label{eq:ym-av-depth-108-error}
\end{equation}

\subsection{Transducción dimensional}

Sea \(\ell_0\) la unidad longitudinal del lector dimensional. La escala
inversa de longitud y el parámetro de una losa de espesor \(a_m\) son
\begin{equation}
 \kappa_{\rm av}
 :=\frac{\mu_{\rm vol}}{\ell_0}
 \left(\frac{\pi_{\rm HMT}}{\varphi_{\rm HMT}^{2}}-1\right)>0,
 \qquad
 q_m:=e^{-3a_m\kappa_{\rm av}}.
 \label{eq:ym-dimensional-gap-scale}
\end{equation}
Por tanto \([\kappa_{\rm av}]=L^{-1}\) y la brecha energética es
\begin{equation}
 \Delta_E=3\hbar c\,\kappa_{\rm av};
 \label{eq:ym-energy-gap-units}
\end{equation}
en unidades \(\hbar=c=1\), se escribe \(\Delta_E=3\kappa_{\rm av}\).
La memoria electrónica \(\Delta_4\) pertenece a otro dominio y a otro
lector; no interviene en \eqref{eq:ym-dimensional-gap-scale}.

En los semigrupos de las secciones siguientes, \(t,s,a_m\) son
longitudes euclídeas y los generadores \(D\) y \(H^{\rm ov}\) tienen
dimensión \(L^{-1}\). El Hamiltoniano energético es
\(H_E=\hbar_{\rm int}c_{\rm int}D\). Si
\(\tau=t/c_{\rm int}\), entonces
\[
 e^{-tD}=e^{-\tau H_E/\hbar_{\rm int}},\qquad
 \Delta=3\kappa_{\rm av},\quad
 \Delta_E=\hbar_{\rm int}c_{\rm int}\Delta,\quad
 m_{\rm gap}=\frac{\hbar_{\rm int}\Delta}{c_{\rm int}}.
\]
La misma brecha queda así expresada, respectivamente, como longitud
inversa, energía y masa en la carta dimensional declarada.

Como \(a_m=9a_{m+1}\), la definición da inmediatamente
\begin{equation}
 q_{m+1}^{9}=q_m.
 \label{eq:ym-nonadic-gap-root}
\end{equation}
Esta identidad es la raíz nonádica que transportará el mismo semigrupo y la
misma escala espectral a través de todos los refinamientos.


% SOURCE sections/ym_construction.tex
\section[Proceso gauge euclídeo Yang--Mills en dimensión 4]{Construcción
proyectiva de un proceso gauge euclídeo Yang--Mills, reconstrucción de
Osterwalder--Schrader y brecha colectiva en dimensión \(4\)}


\noindent La especialización de la dinámica nonádica enriquecida,
construida en las secciones precedentes, es
\[
 (\widehat X_\infty,\Gamma_9)
 \xrightarrow{\ \mathcal R_{\rm YM}\ }
 (\mu_m,T_m,H_m,\Omega_m)
 \longrightarrow(\mu_\infty,\mathrm{OS},E(4))
 \longrightarrow\operatorname{gap}(H_{\rm phys})>0.
\]
El lector clover que publica \(F^2\) reconoce la curvatura sobre este mismo
proceso proyectivo; no construye retrospectivamente \(\mu_m\). La teoría
gauge, la reconstrucción OS, el vacío simple y la brecha ya quedan fijados
por esta cadena. La identidad
\eqref{eq:amp-ym-identidad-accion-requerida} compara después una
parametrización opcional por \(g_R\); es un control no gobernante y no una
condición para identificar la teoría Yang--Mills construida.

\subsection*{Objeto gauge y dominio de refinamiento}

Sea \(G\) un grupo de Lie compacto, conectado y simple. Para cada
\(m\geq0\), se considera la red euclídea
\begin{equation}\label{eq:realizaciones-ym-red}
 \Lambda_m=a_m\mathbb Z^4,
 \qquad a_{m+1}=\frac{a_m}{9}.
\end{equation}
La álgebra cilíndrica gauge se genera por variables de enlace en \(G\),
operadores de entrelazamiento de representación y las coordenadas
enriquecidas de color,
orientación, holonomía, acarreo de fusión, ruta y memoria. Sea \(\mu_m\) la
medida invariante común y sea
\[
 \mathcal H_m=L^2(\mathcal A_m,\mu_m)^{\rm gauge}.
\]
El vector constante \(\Omega_m\) define
\begin{equation}\label{eq:realizaciones-ym-proyectores}
 P_{\Omega_m}=|\Omega_m\rangle\langle\Omega_m|,
 \qquad Q_m^{\rm phys}=I-P_{\Omega_m}.
\end{equation}
El segundo proyector es el complemento físico del vacío; no se identifica
con un residual empleado únicamente para organizar la escala de
renormalización.

El circuito de actualización es local:
\begin{equation}\label{eq:realizaciones-ym-circuito}
 K_m^{\rm loc}
 =\prod_{r=1}^{9\cdot81}
   \prod_{B\in\mathcal B_{m,r}}K_{m,B}.
\end{equation}
La partición \(\{\mathcal B_{m,r}\}_r\) enumera cada coordenada una sola vez; los factores
del circuito se identificarán abajo con el semigrupo de sus expectativas
condicionales, y no con nueve dinámicas independientes.
Los bloques de un mismo color poseen soportes disjuntos y el diámetro físico
de cada soporte está uniformemente acotado al variar \(m\). Esta localidad
forma parte del objeto, de modo que el refinamiento no convierte una
actualización elemental en una operación instantáneamente global.

Sea \(T_m\) la transferencia inducida por
\eqref{eq:realizaciones-ym-circuito}. Las inclusiones isométricas
\(J_m:\mathcal H_m\to\mathcal H_{m+1}\) conservan ancestro,
representación, orientación, acarreo y memoria. Para las cuatro reflexiones
euclídeas \(\Theta_{\nu,m}\), \(\nu=1,\ldots,4\), se cumple
\begin{equation}\label{eq:realizaciones-ym-entrelazamiento}
 (T_{m+1})^9J_m=J_mT_m,
 \qquad
 \Theta_{\nu,m+1}J_m=J_m\Theta_{\nu,m}
 \quad(\nu=1,\ldots,4).
\end{equation}
Las cuatro direcciones pertenecen, por tanto, al mismo refinamiento, la misma
medida y la misma dinámica.

\subsection*{Semigrupo común, corona relativa y cuatro formas de reflexión}

En la carta triangular del estado completo, sean \(E_{m,c}\) las esperanzas
condicionales ortogonales y conmutantes asociadas a coordenadas
independientes. El generador positivo de solapamiento se construye primero en
un nivel semilla \(m_0\):
\begin{equation}\label{eq:realizaciones-ym-generador}
 \widehat H_{m_0}^{\rm ov}
 =3\kappa_{\rm av}
  \sum_{c\in\mathcal I_{m_0}^{\rm seed}}(I-E_{m_0,c}),
\end{equation}
y se prolonga mediante
\begin{equation}\label{eq:realizaciones-ym-generador-refinado}
 \begin{aligned}
 \widehat H_{m+1}^{\rm ov}
 &=\widehat\Gamma_m^*
   \bigl(\widehat H_m^{\rm ov}\otimes I+
         I\otimes\widehat H_{m+1}^{\rm det}\bigr)
   \widehat\Gamma_m,\\
 \widehat H_{m+1}^{\rm det}
 &=3\kappa_{\rm av}
   \sum_{\iota\in\mathcal I_{m+1/m}^{\rm det}}
   (I-E_{m+1,\iota}^{\rm det}),
 \end{aligned}
\end{equation}
donde
\begin{equation}\label{eq:realizaciones-ym-kappa}
 \kappa_{\rm av}
 =\frac{\mu_{\rm vol}}{\ell_0}
  \left(\frac{\pi_{\rm HMT}}{\varphi_{\rm HMT}^{\,2}}-1\right)>0.
\end{equation}
La frecuencia volumétrica \(\mu_{\rm vol}\) es adimensional; la dimensión
procede de \(\ell_0\).

El semigrupo, la transferencia de un paso físico y su compatibilidad de
escala se definen con tipos visibles:
\begin{equation}\label{eq:realizaciones-ym-semigrupo}
 \widehat K_{m,t}^{\rm ov}=e^{-t\widehat H_m^{\rm ov}},
 \qquad
 \widehat T_m=\widehat K_{m,a_m}^{\rm ov},
 \qquad a_{m+1}=a_m/9,
\end{equation}
Como la partición enumera cada coordenada una sola vez y las expectativas
condicionales de sus bloques conmutan, para cada bloque y para el circuito
completo se obtiene
\begin{equation}\label{eq:realizaciones-ym-circuito-semigrupo}
 K_{m,B}=e^{-3\kappa_{\rm av}a_m(I-E_{m,B})},
 \qquad
 K_m^{\rm loc}=e^{-a_m\widehat H_m^{\rm ov}}=\widehat T_m.
\end{equation}
\begin{equation}\label{eq:realizaciones-ym-semigrupo-entrelazado}
 \widehat H_{m+1}^{\rm ov}\widehat J_m
 =\widehat J_m\widehat H_m^{\rm ov},
 \qquad
 (\widehat T_{m+1})^9\widehat J_m
 =\widehat J_m\widehat T_m.
\end{equation}
Así, la novena potencia se compara mediante \(\widehat J_m\); no se igualan
operadores que actúan en Hilbert distintos.

La medida común de ocho caras procede de ese mismo kernel reversible:
\begin{equation}\label{eq:realizaciones-ym-medida-ocho}
 \widehat\Omega_m^{(8)}(dy_1\cdots dy_8)
 =\int\widehat\mu_m^{\rm ov}(dz)
   \prod_{r=1}^{8}
   \widehat K_{m,a_m/2}^{\rm ov}(z,dy_r).
\end{equation}
Para cada reflexión \(\nu=1,\ldots,4\) y todo observable cilíndrico \(F\)
en el semiespacio positivo,
\begin{equation}\label{eq:realizaciones-ym-os}
 \int\overline{F(\Theta_{\nu,m}y)}F(y)\,
 d\widehat\Omega_m^{(8)}
 =\|\widehat{\mathcal B}_{m,\nu}F\|^2\geq0,
\end{equation}
y la marginal de caras opuestas da la identidad de enlace
\begin{equation}\label{eq:realizaciones-ym-os-transferencia}
 \widehat T_{m,\nu}^{\rm OS}
 =\widehat K_{m,a_m}^{\rm ov}
 =e^{-a_m\widehat H_m^{\rm ov}}
 =\widehat T_m,
 \qquad \nu=1,\ldots,4.
\end{equation}
Las cuatro reflexiones publican, por tanto, el mismo par
transferencia--medida.

\subsection*{Persistencia celular del terminal y brecha exacta}

Sean
\[
 K_9=C_*(Q_{9,3};\mathbb Z),\qquad
 K_{10}=C_*(Q_{10,3};\mathbb Z).
\]
La corona relativa tiene dimensiones
\[
 C_*(Q_{10,3},Q_{9,3})=(271,756,702,217),
 \qquad 271+702=756+217=973.
\]
La inclusión \(j:K_9\to K_{10}\), el colapso celular
\(r:K_{10}\to K_9\) y el homotopio integral \(H_{\rm cell}\) satisfacen
\begin{equation}\label{eq:realizaciones-ym-sdr-corona}
 rj=I,\qquad
 \partial H_{\rm cell}+H_{\rm cell}\partial=I-jr,\qquad
 rH_{\rm cell}=0,\quad H_{\rm cell}j=0,\quad H_{\rm cell}^2=0.
\end{equation}
Para todo mapa de cadenas
\(F:K_{10}\otimes V_{q=6}\to\mathcal T\), con
\(P=(jr)\otimes I_{q=6}\),
\begin{equation}\label{eq:realizaciones-ym-homotopia-terminal}
 F-FP
 =d_{\mathcal T}\bigl(F(H_{\rm cell}\otimes I)\bigr)
  +F(H_{\rm cell}\otimes I)(\partial\otimes I).
\end{equation}
La corona completa es así nulhomotópica y el terminal \(FP\) permanece
literalmente. Los \(271\) vértices no sustituyen las aristas, caras y cubos
de la corona; el factor celular tridimensional tampoco se identifica con el
factor gauge físico \(q=6\).

Para cada coordenada persistente \(c\), en la carta local
\(q=(q_1,\ldots,q_6)\in\mathbb F_3^6\), se fija el testigo centrado
\begin{equation}\label{eq:realizaciones-ym-testigo-definicion}
 \bigcap_c\operatorname{Ran}E_{m,c}=\mathbb C\Omega_m,
 \qquad
 W_c(q)=\mathbf 1_{\{q_1+\cdots+q_6=0\}}-\frac13,
\end{equation}
donde la suma del indicador se toma en \(\mathbb F_3\). La descomposición de
Hoeffding de las expectativas conmutantes prueba
\begin{equation}\label{eq:realizaciones-ym-gap-finito}
 \ker\widehat H_m^{\rm ov}=\mathbb C\Omega_m,
 \qquad
 \widehat H_m^{\rm ov}\succeq
 3\kappa_{\rm av}(I-P_{\Omega_m}).
\end{equation}
La igualdad no se deduce sólo de que el núcleo sea trivial: cada sector no
constante de la descomposición recibe al menos un proyector
\(I-E_{m,c}\) y, por \eqref{eq:realizaciones-ym-generador}, su coste mínimo
es \(3\kappa_{\rm av}\). El testigo material \(W_c\) satisface
\begin{equation}\label{eq:realizaciones-ym-testigo-gap}
 \widehat H_m^{\rm ov}W_c=3\kappa_{\rm av}W_c,
 \qquad \|W_c\|^2=\frac29,
\end{equation}
de modo que el complemento es no vacío y la brecha es exactamente
\(3\kappa_{\rm av}\).

La escala se publica por
\begin{equation}\label{eq:realizaciones-ym-qm}
 q_m=e^{-3\kappa_{\rm av}a_m},
 \qquad q_{m+1}^{\,9}=q_m,
 \qquad -a_m^{-1}\log q_m=3\kappa_{\rm av}.
\end{equation}
Aunque \(q_m\to1\) al refinar, el cociente espectral por unidad física
permanece constante. La contracción celular
\eqref{eq:realizaciones-ym-homotopia-terminal} conserva el terminal; la
brecha procede del generador positivo de solapamiento y de su descomposición de
Hoeffding, no de la aciclicidad por sí sola.

Las inclusiones de \eqref{eq:realizaciones-ym-semigrupo-entrelazado}
forman el límite inductivo
\begin{equation}\label{eq:realizaciones-ym-colimite}
 \mathcal H_\infty=\varinjlim(\mathcal H_m,\widehat J_m).
\end{equation}
Sobre la unión inductiva algebraica se define la forma
\begin{equation}\label{eq:realizaciones-ym-forma-limite}
 \mathfrak h_\infty[\widehat J_{m,\infty}\psi]
 :=\langle\psi,\widehat H_m^{\rm ov}\psi\rangle_{\mathcal H_m}.
\end{equation}
El entrelazamiento de los generadores hace que esta definición sea
independiente del representante. La forma es positiva y cerrable; su
clausura determina el generador autoadjunto \(\widehat H_\infty^{\rm ov}\).
La compatibilidad de las formas transporta
\begin{equation}\label{eq:realizaciones-ym-gap-forma-limite}
 \overline{\mathfrak h}_\infty[\psi]
 \geq3\kappa_{\rm av}
 \|(I-P_{\Omega_\infty})\psi\|^2,
\end{equation}
y el vector persistente
\(\widehat J_{m,\infty}W_c\) alcanza la igualdad. En consecuencia, la
brecha exacta del generador límite es \(3\kappa_{\rm av}\).

\subsection*{Hipótesis exactas y encadenamiento del teorema}

La publicación gauge utiliza:
\begin{enumerate}
 \item el grupo compacto, conectado y simple y la red cuatro-dimensional
       \eqref{eq:realizaciones-ym-red};
 \item el circuito local \eqref{eq:realizaciones-ym-circuito}, con diámetro
       físico uniforme, y una medida invariante común;
 \item las inclusiones isométricas y las cuatro reflexiones entrelazadas por
       \eqref{eq:realizaciones-ym-entrelazamiento} y
       \eqref{eq:realizaciones-ym-semigrupo-entrelazado};
 \item las cuatro factorizaciones
       \eqref{eq:realizaciones-ym-os} y el enlace exacto
       \eqref{eq:realizaciones-ym-os-transferencia} sobre el mismo par
       transferencia--medida;
 \item el generador de expectativas conmutantes
       \eqref{eq:realizaciones-ym-generador}--
       \eqref{eq:realizaciones-ym-generador-refinado} y su descomposición de
       Hoeffding;
 \item el SDR integral de la corona y la conservación del terminal
       \eqref{eq:realizaciones-ym-sdr-corona}--
       \eqref{eq:realizaciones-ym-homotopia-terminal};
 \item el terminal \(P_{\Omega_m}\), el complemento físico
       \(Q_m^{\rm phys}\) y la ley de escala
       \eqref{eq:realizaciones-ym-qm};
 \item el límite inductivo y la clausura compatible de formas
       \eqref{eq:realizaciones-ym-colimite}--
       \eqref{eq:realizaciones-ym-gap-forma-limite};
 \item las cartas internas de acción y velocidad
       \(\hbar_{\rm int}\), \(c_{\rm int}\), usadas sólo para publicar la
       dimensión de masa.
\end{enumerate}

\begin{theorem}[Teoría gauge local con vacío simple y brecha colectiva]
\label{thm:realizaciones-ym}
Bajo las premisas anteriores, el límite de la familia gauge es local y
positivo por reflexión en las cuatro direcciones euclídeas. Su generador
posee vacío simple y satisface, en el sector colectivo gauge-invariante,
\begin{equation}\label{eq:realizaciones-ym-gap-limite}
 \Delta_{\rm YM}^{\rm HMT}=3\kappa_{\rm av}>0,
 \qquad
 m_{\rm gap}^{\rm HMT}
 =\frac{\hbar_{\rm int}}{c_{\rm int}}
  \Delta_{\rm YM}^{\rm HMT}>0.
\end{equation}
La brecha pertenece al sector colectivo; no asigna una masa elemental al
gluón.
\end{theorem}

\begin{proof}
La localidad uniforme del circuito conserva la red de álgebras locales. Las
igualdades \eqref{eq:realizaciones-ym-semigrupo-entrelazado} identifican la
novena potencia después de aplicar la inclusión correcta. La construcción
\eqref{eq:realizaciones-ym-medida-ocho} usa ese mismo semigrupo en las ocho
caras y \eqref{eq:realizaciones-ym-os-transferencia} identifica sus cuatro
marginales OS con \(\widehat T_m\). Por tanto, las cuatro formas positivas
\eqref{eq:realizaciones-ym-os} son compatibles en el límite y conservan una
sola medida y una sola dinámica.

Las expectativas de \eqref{eq:realizaciones-ym-generador} son ortogonales y
conmutantes. Su descomposición de Hoeffding separa el modo constante de cada
sector no constante: el primero es \(\mathbb C\Omega_m\), mientras todo
sector del complemento recibe al menos un factor \(I-E_{m,c}\) con peso
\(3\kappa_{\rm av}\). Esto prueba
\eqref{eq:realizaciones-ym-gap-finito}. El testigo
\eqref{eq:realizaciones-ym-testigo-gap} alcanza la cota, de modo que no es
sólo un piso: es la primera energía no nula.

La homotopía \eqref{eq:realizaciones-ym-homotopia-terminal} contrae el
complemento celular y deja intacto \(FP\); por ello la corona no introduce un
segundo vacío ni borra el acarreo. La compatibilidad isométrica transporta el
vacío y su complemento al nivel siguiente sin mezclarlos con el residual de
escala. La ley
\eqref{eq:realizaciones-ym-qm} muestra que el cociente espectral por unidad
física es exactamente \(3\kappa_{\rm av}\) en todo nivel. Las formas
\eqref{eq:realizaciones-ym-forma-limite}--
\eqref{eq:realizaciones-ym-gap-forma-limite} transportan la cota y el testigo
al generador autoadjunto del límite, por lo que la brecha sigue siendo exacta.
La reconstrucción por positividad de
Osterwalder--Schrader identifica el Hamiltoniano con el generador de esta
misma transferencia, conserva el vacío simple y produce
\eqref{eq:realizaciones-ym-gap-limite}. La carta dimensional final transforma
la brecha en la masa colectiva indicada.
\end{proof}

\input{sections/ym_refuerzo.tex}

\subsection*{Lectura de Wilson}

El observable de Wilson se obtiene después del generador. Para una carta
arquimediana de enlace,
\begin{equation}\label{eq:realizaciones-ym-wilson}
 A_e=\sum_a\operatorname{ev}_{\mathbb R}(x_{e,a})T_a,
 \qquad
 U_e=e^{a_mA_e},
 \qquad
 W_\gamma=\operatorname{Tr}\!\prod_{e\in\gamma}U_e.
\end{equation}
La lectura conserva la representación, el operador de entrelazamiento y el
orden de la holonomía. Sirve como observable terminal de la teoría publicada; no define
retrospectivamente la medida, el semigrupo ni la brecha.

\subsection*{Necesidad de las hipótesis y obstrucciones exactas}

\begin{ablation}[Necesidad de la dinámica y la medida comunes]
\label{abl:realizaciones-ym}
El teorema \ref{thm:realizaciones-ym} no se conserva si se sustituyen las
nueve subfases por actualizaciones independientes, se emplean medidas o
núcleos diferentes en las cuatro reflexiones, se pierde la localidad física
uniforme, \(\widehat J_m\) deja de entrelazar transferencia y reflexión, se
borran acarreo o memoria, se elimina \(P_{\Omega_m}\), se identifica
\(Q_m^{\rm phys}\) con un residual de escala, se exige
\(q_m\leq q_*<1\) en unidades de retícula, se reemplaza la corona completa
por sus \(271\) vértices, se borra \(FP\), o se usa
\eqref{eq:realizaciones-ym-wilson} para seleccionar el generador.
\end{ablation}

\begin{proof}
Actualizaciones, núcleos o medidas distintos destruyen la factorización
simultánea \eqref{eq:realizaciones-ym-os}. Sin localidad o entrelazamiento,
la positividad de un corte no define una red local compatible. Borrar acarreo,
memoria o \(FP\) elimina información que
\eqref{eq:realizaciones-ym-homotopia-terminal} conserva. Usar sólo los
vértices cambia el complejo relativo y puede fabricar modos cero. Confundir
los dos complementos permite que el sector físico desaparezca por una
operación puramente escalar. La cota fija en unidades de retícula contradice
\eqref{eq:realizaciones-ym-qm}. Finalmente, Wilson sólo es una lectura del
estado construido y no demuestra por sí mismo ni
\eqref{eq:realizaciones-ym-os} ni
\eqref{eq:realizaciones-ym-gap-finito}.
\end{proof}

Los residuos independientes son: el fallo de
\(\widehat T_{m,\nu}^{\rm OS}=e^{-a_m\widehat H_m^{\rm ov}}\), el fallo del
entrelazamiento \eqref{eq:realizaciones-ym-semigrupo-entrelazado}, una forma
OS negativa, un vector de núcleo ortogonal a \(\Omega_m\), o un sector de
Hoeffding no constante con energía inferior a \(3\kappa_{\rm av}\).

\subsection*{Alcance de la construcción gauge}

La prueba combina la conexión conjunta, la conservación del terminal, la
doble lectura, el circuito local, la medida común, las cuatro factorizaciones
de reflexión, el refinamiento entrelazado y el generador positivo. En el
dominio gauge regular declarado, la
localidad, la teoría Yang--Mills cuatridimensional, el vacío simple y la
brecha colectiva son resultados exactos. La reconstrucción de
Osterwalder--Schrader y el observable de Wilson son
publicaciones posteriores del mismo estado; ninguna de ellas actúa como
selector retrospectivo de la medida, del semigrupo, del vacío ni de la
brecha. La comparación opcional con una familia parametrizada por \(g_R\) se estudia
mediante \eqref{eq:amp-ym-identidad-accion-requerida}; esa comprobación no
gobierna ni reabre la medida, el semigrupo, el vacío o la brecha ya
demostrados.


% SOURCE sections/ym_refuerzo.tex
\subsection*{Compresión gauge y testigo físico de la brecha}

Sea
\[
 \widehat{\mathscr H}_m
 :=L^2(\widehat X_m,\widehat\mu_m^{\rm ov})
\]
el espacio de la realización material completa. Además de la acción gauge
usual sobre enlaces y marcos, cada collar \(c\) porta la acción finita de
incidencia
\begin{equation}\label{eq:amp-ym-accion-gauge-incidencia}
 q_c\longmapsto q_c+Bx_c,
 \qquad x_c\in\mathbb F_3^4,
\end{equation}

donde las seis filas de \(B\) están indexadas por
\((01,02,03,12,13,23)\). El producto de ambas acciones preserva
\(\widehat\mu_m^{\rm ov}\). Con \(\mathcal C_m\) el conjunto finito de
collares presentes en el nivel, se escribe
\[
 \mathcal G_m^{\rm full}
 :=G^{V(\Lambda_m)}
   \times\prod_{c\in\mathcal C_m}\mathbb F_3^4.
\]
Su promedio de Haar define la proyección
ortogonal
\begin{equation}\label{eq:amp-ym-proyeccion-fisica-haar}
 \mathbb P_m^{\rm phys}F
 :=\int_{\mathcal G_m^{\rm full}}F(g^{-1}\!\cdot\,)\,dg,
 \qquad
 \widehat{\mathscr H}_m^{\rm phys}
 :=\operatorname{Ran}\mathbb P_m^{\rm phys}.
\end{equation}
Ésta es una compresión sobre la subálgebra gauge del estado completo; la
marginal que olvida las coordenadas de incidencia es otro mapa y no se
confunde con \(\mathbb P_m^{\rm phys}\).

\begin{lemma}[Reducción del Hamiltoniano a la subálgebra física]
\label{lem:amp-ym-compresion-fisica}
La proyección \(\mathbb P_m^{\rm phys}\) reduce el generador y su
semigrupo. Las inclusiones de refinamiento reducen simultáneamente las
subálgebras físicas:
\begin{equation}\label{eq:amp-ym-compresion-entrelazada}
 \begin{gathered}
 \relax
 [\mathbb P_m^{\rm phys},\widehat H_m^{\rm ov}]=0,
 \qquad
 [\mathbb P_m^{\rm phys},e^{-t\widehat H_m^{\rm ov}}]=0,\\
 \mathbb P_{m+1}^{\rm phys}\widehat J_m
 =\widehat J_m\mathbb P_m^{\rm phys},\qquad
 H_m^{\rm phys}
 :=\widehat H_m^{\rm ov}\big|_{\widehat{\mathscr H}_m^{\rm phys}},\\
 H_{m+1}^{\rm phys}\widehat J_m
 =\widehat J_mH_m^{\rm phys}.
 \end{gathered}
\end{equation}
\end{lemma}

\begin{proof}
El generador físico sobre enlaces es gauge equivariante por la
biinvariancia de Haar. En la carta producto, \(E_{q_c}\) es la esperanza
uniforme en \(\mathbb F_3^6\); por invariancia de la medida uniforme bajo
las traslaciones de \eqref{eq:amp-ym-accion-gauge-incidencia}, conmuta con
esa acción. Las expectativas de las subcolas de marcación actúan en
factores disjuntos. Por tanto cada sumando de
\[
 \widehat H_m^{\rm ov}
 =H_m^{\rm ov}\otimes I
 +3\kappa_{\rm av}\sum_c
 \bigl[(I-E_{q_c})+(I-E_{u_c^{\rm P}})\bigr]
\]
conmuta con el promedio \eqref{eq:amp-ym-proyeccion-fisica-haar}. El
cálculo funcional da la conmutación con el semigrupo. Finalmente,
\(\widehat J_m\) conserva los factores ancestrales e inserta la constante
en los factores nuevos; promediar antes o después de esa inclusión produce
la misma función. Esto prueba las cuatro identidades de
\eqref{eq:amp-ym-compresion-entrelazada}.
\end{proof}

\begin{lemma}[Testigo gauge no nulo y entrelazado]
\label{lem:amp-ym-testigo-fisico-q6}
Para todo collar \(c\), la función
\begin{equation}\label{eq:amp-ym-testigo-fisico-def}
 W_c(q)
 :=\mathbf 1_{\{q_1+\cdots+q_6=0\}}-\frac13
\end{equation}
pertenece a \(\widehat{\mathscr H}_m^{\rm phys}\), no es nula y verifica
\begin{equation}\label{eq:amp-ym-testigo-fisico-momentos}
 \mathbb E W_c=0,\qquad
 \|W_c\|^2=\frac29,\qquad
 \mathbb E(W_c^3)=\frac2{27}.
\end{equation}
Además,
\begin{equation}\label{eq:amp-ym-testigo-fisico-gap}
 H_m^{\rm phys}W_c=3\kappa_{\rm av}W_c,\qquad
 H_{m+1}^{\rm phys}\widehat J_mW_c
 =3\kappa_{\rm av}\widehat J_mW_c.
\end{equation}
En consecuencia, el valor \(3\kappa_{\rm av}\) pertenece al espectro de la
restricción física y no sólo al de una dilatación material.
\end{lemma}

\begin{proof}
La acción gauge sobre enlaces actúa en el primer factor y deja fijo
\(W_c\). Para la acción de incidencia, cada \(x_i\) aparece en tres de las
seis coordenadas, luego
\[
 \sum_{i<j}(Bx)_{ij}=3\sum_i x_i=0
 \quad\text{en }\mathbb F_3.
\]
La condición de \eqref{eq:amp-ym-testigo-fisico-def} es, por tanto,
invariante; también lo es bajo las permutaciones de las seis incidencias.
Así \(\mathbb P_m^{\rm phys}W_c=W_c\).

Bajo la medida uniforme, la suma de los seis trits es uniforme en
\(\mathbb F_3\). La función toma el valor \(2/3\) con probabilidad \(1/3\)
y \(-1/3\) con probabilidad \(2/3\), lo que da
\eqref{eq:amp-ym-testigo-fisico-momentos}. El primer sumando de
\(\widehat H_m^{\rm ov}\) anula \(W_c\); para \(c'\ne c\),
\((I-E_{q_{c'}})W_c=0\), mientras
\(E_{q_c}W_c=0\). Todas las expectativas de marcación anulan la
diferencia correspondiente porque \(W_c\) es constante en esas subcolas.
Queda exactamente
\(\widehat H_m^{\rm ov}W_c=3\kappa_{\rm av}W_c\). La primera identidad de
\eqref{eq:amp-ym-testigo-fisico-gap} sigue de la reducción; la segunda,
del entrelazamiento de \eqref{eq:amp-ym-compresion-entrelazada}.
\end{proof}

\subsection*{Semigrupo gauge, vacío y cota espectral exacta}

Sobre la subálgebra física se escribe
\[
 D_m^{\rm YM}:=H_m^{\rm phys},\qquad
 K_{m,t}:=\exp(-tD_m^{\rm YM}),\qquad
 P_{\Omega_m}:=|\Omega_m\rangle\langle\Omega_m|.
\]
La descomposición de Hoeffding del mismo generador y el cálculo funcional
dan, sin cambiar de medida ni de proceso,
\begin{equation}\label{eq:amp-ym-owner-semigrupo-vacio}
 K_{m,t}P_{\Omega_m}=P_{\Omega_m},
 \qquad
 \bigl\|K_{m,t}(I-P_{\Omega_m})\bigr\|
 \leq e^{-3\kappa_{\rm av}t}.
\end{equation}
El testigo de \eqref{eq:amp-ym-testigo-fisico-def} es no nulo y alcanza la
cota:
\begin{equation}\label{eq:amp-ym-owner-testigo-wc}
 D_m^{\rm YM}W_c=3\kappa_{\rm av}W_c,
 \qquad \|W_c\|^2=\frac29.
\end{equation}
Por tanto el vacío es simple y la brecha no sólo está acotada inferiormente,
sino que es exactamente
\begin{equation}\label{eq:amp-ym-owner-brecha-exacta}
 \Delta_{\rm YM}^{\rm HMT}=3\kappa_{\rm av}>0.
\end{equation}
El semigrupo, el proyector de vacío, la contracción del complemento y
\(W_c\) pertenecen a la misma familia proyectiva y constituyen la cadena
propietaria del cierre.

\subsection*{Reconstrucción euclídea, campos de Schwinger y lector de curvatura}

La medida proyectiva común del teorema determina inclusiones isométricas
\(\widehat J_m\) entre los espacios de nivel. Para cada una de las cuatro
reflexiones coordenadas \(\theta_\mu\), \(\mu=1,\ldots,4\), existe una
factorización
\begin{equation}\label{eq:amp-ym-cuatro-os}
 \int\overline{F(\theta_\mu y)}F(y)\,
 d\widehat\Omega_m(y)
 =\|\mathcal B_{m,\mu}F\|_2^2\geq0.
\end{equation}
Las cuatro formas proceden de la misma medida y del mismo semigrupo. Los
conectores OS
\[
 \mathcal J_{m,\mu}^{\rm OS}
 =\mathcal V_{m+1,\mu}^{-1}\widehat J_m\mathcal V_{m,\mu}
\]
entrelazan sus Hamiltonianos. En el colímite,
\begin{equation}\label{eq:amp-ym-hamiltoniano-comun}
 \mathcal V_{\infty,\mu}H_{{\rm OS},\infty}^{(\mu)}
 \mathcal V_{\infty,\mu}^{-1}
 =\widehat H_\infty,
 \qquad \mu=1,\ldots,4.
\end{equation}

La forma límite se diagonaliza por la descomposición de Hoeffding:
\begin{equation}\label{eq:amp-ym-forma-limite}
 \mathcal E_\infty(u)=3\kappa_{\rm av}
 \sum_{S\Subset\mathcal I_\infty}|S|\,\|u_S\|^2.
\end{equation}
Las sumas parciales proporcionan una sucesión de recuperación y la
semicontinuidad produce el límite inferior. La convergencia de Mosco de las
formas cerradas conserva el núcleo unidimensional y el umbral espectral. El
vector material \(W_c\) satisface
\begin{equation}\label{eq:amp-ym-testigo-gap}
 \widehat H_mW_c=3\kappa_{\rm av}W_c,
 \qquad \|W_c\|^2=\frac29,
 \qquad \mathbb E(W_c^3)=\frac2{27},
\end{equation}
de modo que la cota inferior se alcanza y la brecha es exactamente
\(3\kappa_{\rm av}\).

\begin{theorem}[Terminación euclídea y publicación de curvatura]
\label{thm:amp-ym-terminacion}
La familia proyectiva del teorema \ref{thm:realizaciones-ym} determina una
acción fuertemente continua de \(E(4)\), una red local gauge no escalar,
funcionales de Schwinger compatibles y una publicación de campo en
\(H^{-s}_{\rm loc}(\mathbb R^4)\) para \(s>2\). El producto superficial
ordenado de los enlaces define lectores clover y \(F^2\); en el dominio
suave, el lector de acción converge a
\begin{equation}\label{eq:amp-ym-f2}
 \mathcal S_{\rm YM}(A)
 =\frac1{4g_R^2}\int_{\mathbb R^4}\|F_A(x)\|^2\,d^4x.
\end{equation}
Estas publicaciones conservan el vacío simple y la brecha exacta del mismo
Hamiltoniano común.
\end{theorem}

\begin{proof}
Los Gram cofinales de los lectores suavizados y la conmutatividad de los
cuadrados de rotación y refinamiento preservan los ideales nulos; su
colímite construye la representación fuerte de \(E(4)\) y la red local.
Las estimaciones martingala de los incrementos de enlace dan tightness en
\(H^{-s}_{\rm loc}\), \(s>2\), y las características mixtas identifican los
funcionales de Schwinger con el mismo proceso cilíndrico. El producto
orientado alrededor de cada plaqueta construye el clover. Su expansión en
el régimen suave y la convergencia fuerte del lector marginal dan
\eqref{eq:amp-ym-f2}. Ninguna de estas operaciones modifica la medida, el
semigrupo ni la forma límite usados para obtener
\eqref{eq:amp-ym-testigo-gap}; por ello el vacío y la brecha se conservan.
\end{proof}

\subsection*{Identificación exacta del límite gauge}

La identificación de la teoría no descansa en que cuatro reflexiones
generen por sí solas el grupo euclídeo.  Las reflexiones proporcionan la
positividad OS; la acción continua procede, separadamente, del grupoide de
pantallas y de la compatibilidad de refinamiento.  Para
\(g=(g_v)_{v\in\Lambda_m}\in G^{\Lambda_m}\), la acción sobre una variable
de enlace es
\begin{equation}\label{eq:amp-ym-accion-gauge-enlace}
 (g\cdot U)_e=g_{s(e)}U_e g_{t(e)}^{-1}.
\end{equation}
La invariancia de Haar, la desintegración triangular del kernel y la
cancelación de las aristas interiores prueban
\begin{equation}\label{eq:amp-ym-invariancia-gauge-medida}
 (g\cdot)_*\widehat\mu_m^{\rm ov}=\widehat\mu_m^{\rm ov},
 \qquad
 K_{m,t}^{\rm ov}(g\cdot U,g\cdot dV)
 =K_{m,t}^{\rm ov}(U,dV).
\end{equation}
Por diferenciación sobre el núcleo cilíndrico, todo generador vertical
\(X\) satisface la identidad de Ward
\begin{equation}\label{eq:amp-ym-ward}
 \int \nabla_XF\,d\widehat\mu_m^{\rm ov}=0.
\end{equation}

La curvatura se obtiene del mismo sistema de enlaces.  Si una plaqueta
gruesa \(p\) se subdivide en sus 81 subplaquetas \(p'\), entonces
\begin{equation}\label{eq:amp-ym-producto-superficial}
 \operatorname{Hol}^{(m)}_{p}
 =\prod_{p'\subset p}^{\longrightarrow}
 T_{p'}\operatorname{Hol}^{(m+1)}_{p'}T_{p'}^{-1}.
\end{equation}
Las contribuciones de toda arista interior aparecen con orientaciones
opuestas y se cancelan exactamente.  La expresión es gauge-covariante y
conmuta con las inclusiones \(\widehat J_m\).  En una carta suave se define
\begin{equation}\label{eq:amp-ym-clover-definido}
 F_{{\rm cl},m}(x)
 =\frac1{4a_m^2}\sum_{p\ni x}
 \operatorname{Ad}_{T_p}\log\operatorname{Hol}_{p},
 \qquad
 F_{{\rm cl},m}=F_A+O(a_m^2).
\end{equation}
De aquí, para \(A\in C^4_{\rm loc}\) y
\(a_m^2/g_m^2\to0\), se obtiene por suma de Riemann
\begin{equation}\label{eq:amp-ym-f2-convergencia-calculada}
 \frac{a_m^4}{4g_m^2}
 \sum_{x\in\Lambda_m}\|F_{{\rm cl},m}(x)\|^2
 \longrightarrow
 \frac1{4g_R^2}\int_{\mathbb R^4}\|F_A(x)\|^2\,d^4x.
\end{equation}

\begin{theorem}[Proceso gauge euclídeo y lector de curvatura de
Yang--Mills]
\label{thm:amp-ym-identificacion-gauge}
El límite construido en el teorema
\ref{thm:realizaciones-ym} es un proceso gauge euclídeo en dimensión
\(4\) para el grupo compacto, conectado y simple fijado. Posee
covariancia gauge local, una red local no escalar, positividad por
reflexión, una representación fuertemente continua de \(E(4)\),
funcionales de Schwinger compatibles, curvatura gauge-covariante con límite
clásico \eqref{eq:amp-ym-f2-convergencia-calculada}, vacío simple y brecha
espectral
\(3\kappa_{\rm av}>0\).
\end{theorem}

\begin{proof}
Las identidades \eqref{eq:amp-ym-invariancia-gauge-medida}--
\eqref{eq:amp-ym-ward} construyen la simetría gauge y sus identidades de
Ward en todos los niveles; su naturalidad las transporta al límite.  La
localidad uniforme de los kernels produce la red isótona y local.  Las
cuatro formas \eqref{eq:amp-ym-cuatro-os}, junto con la acción del grupoide
de pantallas sobre el núcleo Weyl, producen respectivamente la positividad
OS y la acción fuerte de \(E(4)\).  La tightness en
\(H^{-s}_{\rm loc}\), \(s>2\), identifica una única ley límite y sus
funcionales de Schwinger.  Finalmente,
\eqref{eq:amp-ym-producto-superficial}--
\eqref{eq:amp-ym-f2-convergencia-calculada} identifican el lector de campo
y su límite clásico, mientras
\eqref{eq:amp-ym-forma-limite}--
\eqref{eq:amp-ym-testigo-gap} prueban vacío simple y brecha exacta para la
misma ley.

La medida queda determinada constructivamente por sus marginales
proyectivas, su kernel reversible y sus funcionales cilíndricos. El lector
clover no repondera esa medida: publica una función del mismo proceso y su
límite clásico.
\end{proof}

\begin{proposition}[Control no gobernante de comparación Gibbs--acción]
\label{prop:amp-ym-puente-accion-requerido}
La ley proyectiva, su vacío simple y su brecha exacta ya están construidos
por \eqref{eq:amp-ym-owner-semigrupo-vacio}--
\eqref{eq:amp-ym-owner-brecha-exacta}. Como comparación exterior opcional de
tipo Gibbs--acción, una parametrización por \(g_R\) puede estudiar una familia
\[
 g_R\longmapsto
 \bigl(\widehat\mu_{m,g_R}^{\rm ov},H_{m,g_R}^{\rm phys}\bigr)
\]
y una relación entre la medida o su forma de Dirichlet y el funcional de
acción, por ejemplo
\begin{equation}\label{eq:amp-ym-identidad-accion-requerida}
 \frac{d\widehat\mu_{m,g_R}^{\rm ov}}{d\nu_m}(U)
 =Z_{m,g_R}^{-1}e^{-S_{m,g_R}^{F^2}(U)},
 \qquad
 \mathfrak h_{m,g_R}[F]
 =\langle F,H_{m,g_R}^{\rm phys}F\rangle,
\end{equation}
con compatibilidad de refinamiento y paso al límite. Aquí \(\nu_m\) puede
ser la medida producto de Haar de la red finita. Esta comprobación se rotula
\texttt{CONTROL\_NO\_GOBERNANTE}: no es premisa de la teoría gauge HMT ni
de su brecha.
\end{proposition}

\begin{proof}
En \eqref{eq:amp-ym-f2-convergencia-calculada}, variar \(g_R\) reescala el
lector de curvatura, mientras las fórmulas que construyen
\(\widehat\mu_m^{\rm ov}\), \(\widehat H_m^{\rm ov}\) y la brecha
\(3\kappa_{\rm av}\) no contienen \(g_R\). Son, por ello, construcciones
independientes. La compresión
\eqref{eq:amp-ym-compresion-entrelazada} y el testigo
\eqref{eq:amp-ym-testigo-fisico-gap} prueban la persistencia física de la
brecha del proceso ya construido; el control sólo compara después una
dependencia dinámica adicional en \(g_R\).
\end{proof}

La secuencia demostrada en esta sección es
\[
 \text{medida común}\to\text{cuatro OS}\to\text{Hamiltoniano común}
 \to E(4),\ \text{Schwinger},\ H^{-s}_{\rm loc}
 \to\text{clover y }F^2.
\]
Esta secuencia demuestra una teoría Yang--Mills cuatridimensional con vacío
simple y brecha positiva exacta \(3\kappa_{\rm av}\) en el dominio HMT
registrado. La identidad
\eqref{eq:amp-ym-identidad-accion-requerida} queda fuera de la cadena
rectora como comparación posterior independiente.

La realización dependiente de \(g\) se desarrolla en
\eqref{eq:pdf2-ym-gibbs-measure} y
\eqref{eq:pdf2-ym-g-dynamics}: la medida, las expectativas
condicionales y el generador se transportan conjuntamente,
y \eqref{eq:pdf2-ym-g-gap-before-limit} conserva la cota y su testigo.


% SOURCE sections/ym_gauge_object.tex

\subsection{Incidencia proyectiva de calibre del constructor común y descenso al cociente}

La imagen gauge no es la marginal de enlaces. Su objeto completo por nivel es
\begin{equation}\label{eq:pdf2-g-gau-explicito}
 \mathfrak G_{{\rm gau},m}(d_{\rm gau})=
 \bigl(
 \begin{gathered}
  \mathfrak G_{\rm gau}^{\rm int};\\
  \Lambda_m,\widehat X_m,\widehat\mu_m^{\rm ov},
  \{\widehat\mu_{m,g}^{\rm YM},\nu_{m,g},K_{m,g}^{\rm phys}\}_{g>0},\\
  \mathcal G_m^{\rm full},\mathbb P_m^{\rm phys},
  \widehat J_m,\widehat\Gamma_m,
  \{E_{m,c}\}_c,\widehat H_m^{\rm ov},D_m^{\rm YM},\\
  \{\Theta_{\nu,m}\}_{\nu=1}^4,
  \{\mathfrak A_{+,\nu,m},\operatorname{ev}_{t}^{(\nu)},
      \mathcal V_{m,\nu}^{\rm OS}\}_{\nu=1}^4,\\
  \widehat\Omega_m^{(8)},P_{\Omega_m},W_c,
  \operatorname{Carr}_m,\operatorname{Mem}_m,
  \operatorname{Surv}_m,\tau_m,\mathcal P_m^{F^2}
 \end{gathered}
 \bigr).
\end{equation}
El grupo $\mathcal G_m^{\rm full}$ actúa simultáneamente en enlaces, marcos y
coordenadas de incidencia $q_c\in\mathbb F_3^6$. La proyección física es su promedio
de Haar. El generador reducido es
\begin{equation}\label{eq:pdf2-g-gau-generador-estructural}
 D_m^{\rm YM}
 =\left.3\kappa_{\rm av}\sum_c(I-E_{m,c})
  \right|_{\operatorname{Ran}\mathbb P_m^{\rm phys}},
\end{equation}
y $W_c$ pertenece a esa subálgebra. Las cuatro reflexiones actúan sobre una misma
medida, un mismo kernel y un mismo terminal.
El propietario exacto anterior construye la acción gauge de incidencia
\eqref{eq:amp-ym-accion-gauge-incidencia}, la proyección física de Haar
\eqref{eq:amp-ym-proyeccion-fisica-haar} y la compresión entrelazada
\eqref{eq:amp-ym-compresion-entrelazada}. De esa construcción proceden el
kernel completo equivariable $\widehat K_{m,t}^{\rm ov}$, su descenso al
cociente físico $K_{m,t}^{\rm phys}$, los cuatro operadores de reflexión y
la acción $U_m(g)$ de $E(4)$ en el núcleo cilíndrico suave. Este bloque
ensambla, sin introducirlos como datos, el generador, el vacío, el testigo
de curvatura y las formas OS:
\begin{equation}\label{eq:pdf2-g-gau-entrelazamiento-dato}
 D_{m+1}^{\rm YM}\widehat J_m=\widehat J_mD_m^{\rm YM},
 \quad
 \mathbb P_{m+1}^{\rm phys}\widehat J_m
  =\widehat J_m\mathbb P_m^{\rm phys},
 \quad
 \widehat J_mW_{m,c}=W_{m+1,c}.
\end{equation}
La primera y la tercera identidades son la publicación, en la notación
común, de \eqref{eq:amp-ym-compresion-entrelazada} y
\eqref{eq:amp-ym-testigo-fisico-gap}; el semigrupo, el vacío simple y la
brecha exacta ya se obtienen en
\eqref{eq:amp-ym-owner-semigrupo-vacio}--
\eqref{eq:amp-ym-owner-brecha-exacta}. Aquí
$\nu_{m,g}:=(\pi_m^{\rm phys})_*\widehat\mu_{m,g}^{\rm YM}$ y
$K_{m,g}^{\rm phys}(t)=e^{-tD_{m,g}^{\rm YM}}$ en el cociente físico.
La completitud gauge incluye además la desintegración por cada familia de
hiperplanos euclídeos. Si $A_j\in\mathfrak A_{+,\nu,m}$ y
$t_1\le\cdots\le t_n$, la marginal de la misma medida física satisface
\begin{equation}\label{eq:pdf2-g-gau-transfer-euclidean-dato}
\begin{aligned}
 &\int\prod_{j=1}^n
   (\tau_{t_je_\nu}A_j)(x)\,d\widehat\mu_{m,g}^{\rm YM}(x)\\
 &\quad=
 \int\prod_{j=1}^nA_j(y_j)\,
  \nu_{m,g}(dy_1)
  \prod_{j=1}^{n-1}K_{m,g}^{\rm phys}
       (t_{j+1}-t_j;y_j,dy_{j+1})\\
 &\quad=
 \left\langle\mathbf1,
  M_{A_1}e^{-(t_2-t_1)D_{m,g}^{\rm YM}}M_{A_2}\cdots
  e^{-(t_n-t_{n-1})D_{m,g}^{\rm YM}}M_{A_n}\mathbf1
 \right\rangle.
\end{aligned}
\end{equation}
La reconstrucción anterior prueba la identidad en las cuatro direcciones
mediante \eqref{eq:amp-ym-cuatro-os} y construye un único Hamiltoniano
entrelazado en \eqref{eq:amp-ym-hamiltoniano-comun}. En consecuencia, la
traslación OS se publica aquí como
\begin{equation}\label{eq:pdf2-g-gau-os-generator-dato}
 \mathcal V_{m,\nu}^{\rm OS}T_{m,\nu}^{\rm OS}(s)
 (\mathcal V_{m,\nu}^{\rm OS})^{-1}
 =e^{-sD_{m,g}^{\rm YM}}.
\end{equation}
Así, el generador de expectativas y el Hamiltoniano de traslaciones no se
identifican por nombre: son dos realizaciones ya construidas de la misma
desintegración, cuya forma límite y cuya publicación de curvatura constan en
\eqref{eq:amp-ym-forma-limite} y \eqref{eq:amp-ym-f2}.



% SOURCE sections/ym_complete.tex
\section{Imagen gauge y composición hacia Yang--Mills}

\subsection{Estado gauge enriquecido y generadores de transferencia}

El objeto ya generado que recibe este capítulo es la estructura gauge completa
$\mathfrak G_{\rm gau}$; no constituye un dato externo ni una hipótesis terminal.
En cada nivel $m$ conserva la red
$\Lambda_m=a_m\mathbb Z^4$, $a_{m+1}=a_m/9$, la medida $\mu_m$, los
entrelazadores $J_m$, cuatro reflexiones $\Theta_{\nu,m}$, holonomía, color,
representación, acarreo de fusión, ruta, memoria, terminal y lector de curvatura.
Sus dos dinámicas se tipan con símbolos distintos antes de cualquier límite.

El semigrupo de una sola fibra de memoria es
\begin{equation}\label{eq:pdf2-ym-fibra-semigrupo}
 K^{\rm fib}_{m,t}
 =P_{\Omega_m}+e^{-3\kappa_{\rm av}t}(I-P_{\Omega_m})
 =e^{-tD_m^{\rm fib}},
 \qquad
 D_m^{\rm fib}:=3\kappa_{\rm av}(I-P_{\Omega_m}).
\end{equation}
Su espectro está contenido en $\{0,3\kappa_{\rm av}\}$. El generador de la
carta producto completa es, en cambio,
\begin{equation}\label{eq:pdf2-ym-generador-producto}
 D_m^{\rm prod}:=3\kappa_{\rm av}
 \sum_{c\in\mathcal I_m}(I-E_{m,c}),
 \qquad
 K^{\rm prod}_{m,t}:=e^{-tD_m^{\rm prod}},
\end{equation}
donde las expectativas ortogonales $E_{m,c}$ conmutan. Si hay dos o más
coordenadas no constantes, $D_m^{\rm prod}$ posee niveles
$6\kappa_{\rm av},9\kappa_{\rm av},\ldots$; por ello
\eqref{eq:pdf2-ym-fibra-semigrupo} no puede declararse igual a
\eqref{eq:pdf2-ym-generador-producto}. El primero es la fibra simple; el segundo,
el producto de solapamiento.

\subsection{Teorema finito de Hoeffding, vacío y brecha alcanzada}

En la carta producto construida en
\eqref{eq:pdf2-ym-owner-product-factorization}, los $E_{m,c}$ son expectativas
condicionales ortogonales y conmutantes, y la prueba de
\eqref{eq:pdf2-ym-owner-vacuum} establece
\begin{equation}\label{eq:pdf2-ym-interseccion-vacio}
 \bigcap_{c\in\mathcal I_m}\operatorname{Ran}E_{m,c}
 =\mathbb C\Omega_m.
\end{equation}
La descomposición de Hoeffding es la suma ortogonal
\[
 \mathcal H_m=\bigoplus_{S\subseteq\mathcal I_m}\mathcal H_{m,S},
 \qquad u=\sum_Su_S,
\]
donde $E_{m,c}u_S=0$ si $c\in S$ y $E_{m,c}u_S=u_S$ si $c\notin S$.
Entonces
\begin{equation}\label{eq:pdf2-ym-hoeffding}
 D_m^{\rm prod}u_S=3\kappa_{\rm av}|S|u_S,
 \qquad
 \langle u,D_m^{\rm prod}u\rangle
 =3\kappa_{\rm av}\sum_S|S|\|u_S\|^2.
\end{equation}
La componente $S=\varnothing$ es exactamente el vacío de
\eqref{eq:pdf2-ym-interseccion-vacio}. Por tanto
\begin{equation}\label{eq:pdf2-ym-cota-finita}
 D_m^{\rm prod}\succeq
 3\kappa_{\rm av}(I-P_{\Omega_m}),
 \qquad
 \|K^{\rm prod}_{m,t}(I-P_{\Omega_m})\|
 \le e^{-3\kappa_{\rm av}t}.
\end{equation}

En una coordenada de incidencia
$q=(q_1,\ldots,q_6)\in\mathbb F_3^6$, el testigo
\begin{equation}\label{eq:pdf2-ym-wc}
 W_c(q):=\mathbf1_{\{q_1+\cdots+q_6=0\}}-\frac13
\end{equation}
satisface
\begin{equation}\label{eq:pdf2-ym-wc-momentos}
 \mathbb EW_c=0,\qquad \|W_c\|^2=\frac29,\qquad
 D_m^{\rm prod}W_c=3\kappa_{\rm av}W_c.
\end{equation}
Así, en cada nivel finito, el umbral de
\eqref{eq:pdf2-ym-cota-finita} se alcanza. Este cálculo fija el espectro del
generador producto finito; la medida cuatridimensional continua de Yang--Mills se
construye materialmente en
\eqref{eq:pdf2-ym-gibbs-measure}--\eqref{eq:pdf2-ym-g-dynamics}.

\subsection{Compresión gauge física}

Sea $G$ compacto, conectado y simple. El grupo que actúa en el nivel $m$ es
\[
 \mathcal G_m^{\rm full}
 =G^{V(\Lambda_m)}\times
  \prod_{c\in\mathcal C_m}\mathbb F_3^4,
\]
donde el segundo factor actúa sobre las seis incidencias mediante una matriz
que debe quedar visible. Si las filas se indexan por
$(01,02,03,12,13,23)$ y las columnas por $0,1,2,3$, se fija
\begin{equation}\label{eq:pdf2-ym-incidence-matrix}
 B^+_{(ij),k}:=\delta_{ik}+\delta_{jk}\in\mathbb F_3,
 \qquad
 (B^+x)_{ij}=x_i+x_j.
\end{equation}
Esta es la incidencia no orientada del collar; no se confunde con el
coborde firmado de una arista. La acción finita es
\begin{equation}\label{eq:pdf2-ym-incidence-action}
 q_c\longmapsto q_c+B^+x_c,
 \qquad x_c\in\mathbb F_3^4.
\end{equation}
La proyección de Haar
\begin{equation}\label{eq:pdf2-ym-proyeccion-fisica}
 \mathbb P_m^{\rm phys}F
 :=\int_{\mathcal G_m^{\rm full}}F(g^{-1}\!\cdot\,)\,dg,
 \qquad
 \mathcal H_m^{\rm phys}:=\operatorname{Ran}\mathbb P_m^{\rm phys}
\end{equation}
es distinta de la marginal que olvida incidencia. La invariancia de las
expectativas $E_{m,c}$ da
\begin{equation}\label{eq:pdf2-ym-reduccion}
 [\mathbb P_m^{\rm phys},D_m^{\rm prod}]=0,
 \qquad
 [\mathbb P_m^{\rm phys},K_{m,t}^{\rm prod}]=0.
\end{equation}
Además, cada variable de $x_c\in\mathbb F_3^4$ aparece tres veces en la suma de
las seis incidencias, de modo que
\begin{equation}\label{eq:pdf2-ym-incidence-conservation}
 \sum_{i<j}(B^+x_c)_{ij}
 =3\sum_{i=0}^3x_{c,i}=0
 \quad\text{en }\mathbb F_3.
\end{equation}
Por ello $W_c$ es gauge
invariante y pertenece a $\mathcal H_m^{\rm phys}$. La brecha finita no es un
artefacto de una dilatación no física.

\subsection{Positividad por reflexión en cada nivel}

Para un kernel reversible común, defínase la medida de ocho caras
\begin{equation}\label{eq:pdf2-ym-ocho-caras}
 \Omega_m^{(8)}(dy_1\cdots dy_8)
 :=\int\mu_m(dz)\prod_{r=1}^{8}K^{\rm prod}_{m,a_m/2}(z,dy_r).
\end{equation}
Si $\Theta_{\nu,m}$ intercambia las dos mitades respecto de la dirección
$\nu$, la propiedad de Markov y la reversibilidad factorizan, para todo observable
cilíndrico $F$ soportado en el semiespacio positivo,
\begin{equation}\label{eq:pdf2-ym-os-finita}
 \int\overline{F(\Theta_{\nu,m}y)}F(y)\,d\Omega_m^{(8)}
 =\|\mathcal B_{m,\nu}F\|_2^2\ge0,
 \qquad \nu=1,\ldots,4.
\end{equation}
La marginal de caras opuestas recupera el mismo
$K^{\rm prod}_{m,a_m}$. De este modo las cuatro formas OS finitas son
consecuencia de una sola medida y un solo semigrupo, no de cuatro procesos
independientes.

\subsection{Estado gauge completo y coordenadas físicas de incidencia}

El propietario vigente no toma como espacio físico la marginal que conserva sólo
los enlaces. Su dominio es $\mathfrak G_{\rm gau}$, con enlace, marco,
representación, incidencia, color, ruta, orientación,
acarreo, memoria, supervivencia y terminal. La proyección a enlaces es un mapa de
olvido posterior. Juzgar el testigo físico después de aplicar ese olvido sustituye el
objeto y no transfiere la conclusión.

\subsubsection{Medida proyectiva y publicación de enlaces}

Para una pantalla finita $\Gamma_m=(V_m,E_m)$ se parte del espacio levantado
\begin{equation}\label{eq:pdf2-ym-owner-lifted-space}
 \widetilde X_m:=G^{V_m}\times G^{E_m},\qquad
 d\widetilde\mu_m(h,z)
 =\prod_{v\in V_m}dh_v\prod_{e\in E_m}k_{\tau_{m,e}}(z_e)\,dz_e,
\end{equation}
y se publica cada enlace mediante
\begin{equation}\label{eq:pdf2-ym-owner-link-publication}
 U_e=h_{s(e)}z_eh_{t(e)}^{-1}.
\end{equation}
Una arista gruesa $e$ se refina en la cadena ordenada
$e_1\cdots e_9$. Póngase
\(\boldsymbol\tau=(\tau_1,\ldots,\tau_9)\), con
\(\tau_j=\tau_{m+1,e_j}>0\) y
\(\tau=\tau_{m,e}=\sum_{j=1}^9\tau_j\). Sea
\(d\lambda_z^{(9)}\) la desintegración de Haar sobre la fibra del producto,
caracterizada por
\[
 \int_{G^9}F(z_1,\ldots,z_9)\,dz_1\cdots dz_9
 =\int_G\int_{z_1\cdots z_9=z}
   F\,d\lambda_z^{(9)}\,dz.
\]
El puente correcto para tiempos no necesariamente iguales es
\begin{equation}\label{eq:pdf2-ym-owner-heat-bridge}
 dB_{z,\boldsymbol\tau}^{(9)}(z_1,\ldots,z_9)
 =\frac{\prod_{j=1}^9k_{\tau_j}(z_j)}{k_\tau(z)}
   \,d\lambda_z^{(9)},
 \qquad z_1\cdots z_9=z,
\end{equation}
que es una probabilidad porque
\(k_{\tau_1}*\cdots*k_{\tau_9}=k_\tau\). En el refinamiento nonádico uniforme se
recupera \(\tau_j=\tau/9\), pero la fórmula no presupone esa igualdad.

Para hacer independiente la innovación sin borrar el producto grueso, se fija, para
cada $(e,z)$, un transporte boreliano
\begin{equation}\label{eq:pdf2-ym-owner-bridge-triangularization}
 \chi_{m,e,z}:\bigl((0,1)^{\mathbb N},du^{\otimes\mathbb N}\bigr)
 \longrightarrow
 \bigl(\{(z_j):z_1\cdots z_9=z\},B_{z,\boldsymbol\tau}^{(9)}\bigr)
\end{equation}
que es isomorfismo módulo conjuntos nulos. Existe porque ambas son realizaciones
estándar no atómicas; una elección queda fijada una vez y se transporta bajo
relabelados de aristas. La variable gruesa $z$ y la innovación $u^{\rm det}_{m+1,e}$
son entonces factores independientes en la carta triangular. Los marcos nuevos se
integran con Haar. Ruta, acarreo, orientación y memoria se transportan por los
morfismos $\operatorname{Ext}_{\rm gau}$ del propio dato gauge.

La proyección gruesa multiplica los nueve incrementos y olvida exclusivamente las
innovaciones nuevas. Con esta definición se obtiene
\begin{equation}\label{eq:pdf2-ym-owner-projectivity}
 (\widehat\pi_{m+1,m})_*\widehat\mu_{m+1}^{\rm ov}
 =\widehat\mu_m^{\rm ov},\qquad
 \widehat J_mF=F\circ\widehat\pi_{m+1,m}.
\end{equation}
La misma arista aparece una sola vez en la medida; las seis caras incidentes son seis
lectores orientados de esa arista. Las rutas solapadas conservan el incremento
compartido y no se reemplazan por copias independientes.

\subsubsection{La fibra de incidencia es una coordenada física del estado completo}

La cola de incidencia incluida en $\mathfrak G_{\rm gau}$ se descompone sin perder
información:
\begin{equation}\label{eq:pdf2-ym-owner-seed-split}
 \Sigma:(\mathbb F_3^2)^{\mathbb N}\longrightarrow
 (\mathbb F_3^2)^{\mathbb N}\times\mathbb F_3^6,
 \qquad \Sigma_*\sigma=\sigma\otimes u_6.
\end{equation}
El primer factor continúa alimentando las cartas de innovación y el segundo conserva
las seis incidencias $q_c=(q_{01},q_{02},q_{03},q_{12},q_{13},q_{23})$ del collar
$c$. En vez de dejar implícito qué coordenadas gobierna el generador, se fija una
carta triangular recursiva
\begin{equation}\label{eq:pdf2-ym-owner-product-factorization}
 \widehat X_m
 =\widehat X_{m_0}^{\rm seed}
  \times\prod_{\ell=m_0+1}^{m}
       \mathcal U_{\ell/\ell-1}^{\rm det},
 \qquad
 \widehat\mu_m^{\rm ov}
 =\widehat\mu_{m_0}^{\rm seed}
  \otimes\bigotimes_{\ell=m_0+1}^{m}
       \upsilon_{\ell/\ell-1}^{\rm det}.
\end{equation}
El factor semilla enumera, una vez, los marcos y enlaces gruesos, las incidencias
$q_c$ y sus cartas de marcación. Cada
$\mathcal U_{\ell/\ell-1}^{\rm det}$ enumera, una vez, los marcos nuevos, los puentes
\eqref{eq:pdf2-ym-owner-bridge-triangularization}, las incidencias nuevas y las
subcolas de marcación de ese refinamiento. La factorización es interna a
$\mathfrak G_{\rm gau}$ y no añade otro dato de partida.

Se registra el índice exhaustivo, disjunto y numerable
\begin{equation}\label{eq:pdf2-ym-owner-exhaustive-index}
 \mathcal I_m
 :=\mathcal I_{m_0}^{\rm frame}
 \sqcup\mathcal I_{m_0}^{\rm link}
 \sqcup\mathcal I_{m_0}^{q}
 \sqcup\mathcal I_{m_0}^{\rm mark}
 \sqcup\bigsqcup_{\ell=m_0+1}^{m}
 \bigl(
   \mathcal I_{\ell/\ell-1}^{\rm frame}
   \sqcup\mathcal I_{\ell/\ell-1}^{\rm bridge}
   \sqcup\mathcal I_{\ell/\ell-1}^{q}
   \sqcup\mathcal I_{\ell/\ell-1}^{\rm mark}
 \bigr).
\end{equation}
Cada factor independiente de \eqref{eq:pdf2-ym-owner-product-factorization} aparece
exactamente una vez; una subcola infinita se enumera por sus cilindros escalares. No
queda un factor genealógico fuera de $\mathcal I_m$.
En particular, las abreviaturas de la recurrencia significan literalmente
\begin{equation}\label{eq:pdf2-ym-owner-index-abbreviations}
 \begin{aligned}
 \mathcal I_{m_0}^{\rm seed}
 &:={\mathcal I}_{m_0}^{\rm frame}\sqcup
    {\mathcal I}_{m_0}^{\rm link}\sqcup
    {\mathcal I}_{m_0}^{q}\sqcup
    {\mathcal I}_{m_0}^{\rm mark},\\
 \mathcal I_{\ell/\ell-1}^{\rm det}
 &:={\mathcal I}_{\ell/\ell-1}^{\rm frame}\sqcup
    {\mathcal I}_{\ell/\ell-1}^{\rm bridge}\sqcup
    {\mathcal I}_{\ell/\ell-1}^{q}\sqcup
    {\mathcal I}_{\ell/\ell-1}^{\rm mark}.
 \end{aligned}
\end{equation}

La marginal de enlaces olvida $q_c$; la subálgebra física, en cambio, se obtiene
promediando la acción gauge completa
\begin{equation}\label{eq:pdf2-ym-full-gauge-group}
 \mathcal G_m^{\rm full}
 :=G^{V(\Lambda_m)}\times
 \prod_{c\in\mathcal C_m}\mathbb F_3^4,
 \qquad q_c\longmapsto q_c+B^+x_c.
\end{equation}
Por tanto la proyección física es
\begin{equation}\label{eq:pdf2-ym-full-physical-projection}
 \mathbb P_m^{\rm phys}F
 :=\int_{\mathcal G_m^{\rm full}}F(g^{-1}\!\cdot x)\,dg,
 \qquad
 \widehat{\mathscr H}_m^{\rm phys}
 :=\operatorname{Ran}\mathbb P_m^{\rm phys}.
\end{equation}
No coincide con el mapa de olvido
$\widehat X_m\to X_m^{\rm link}$. Esta distinción es el cortafuegos de tipo
que impide expulsar $W_c$ del sector físico.

\subsection{Generador gauge y circuito local de transporte}

La carta triangular completa proporciona una unitaria
\begin{equation}\label{eq:pdf2-ym-owner-triangular-unitary}
 \widehat\Gamma_m:
 L^2(\widehat X_{m+1},\widehat\mu_{m+1}^{\rm ov})
 \longrightarrow
 L^2(\widehat X_m,\widehat\mu_m^{\rm ov})
 \widehat\otimes\mathcal K_{m+1}^{\rm det},
 \qquad
 \widehat\Gamma_m\widehat J_mF=F\otimes\mathbf1_{\rm det}.
\end{equation}
Para $\iota\in\mathcal I_m$, sea $E_{m,\iota}$ la esperanza que integra
exactamente el factor $\iota$ de
\eqref{eq:pdf2-ym-owner-product-factorization} y deja fijos todos los demás.
Son proyecciones Markov ortogonales, conservan la unidad y conmutan. Sobre el
núcleo denso $\mathscr C_m$ de cilindros que dependen de un número finito de
coordenadas se define la forma
\begin{equation}\label{eq:pdf2-ym-owner-full-form}
 \widehat{\mathcal E}_m[F]
 :=3\kappa_{\rm av}\sum_{\iota\in\mathcal I_m}
   \|(I-E_{m,\iota})F\|_2^2.
\end{equation}
La suma es finita para cada cilindro. Su clausura es una forma de Dirichlet y
determina el generador autoadjunto Markov
$\widehat H_m^{\rm ov}$. En un nivel semilla esto se escribe, en sentido de
formas,
\begin{equation}\label{eq:pdf2-ym-owner-base-generator}
 \widehat H_{m_0}^{\rm ov}
 =3\kappa_{\rm av}
  \sum_{\iota\in\mathcal I_{m_0}^{\rm seed}}(I-E_{m_0,\iota}),
\end{equation}
y cada refinamiento añade exactamente las expectativas de sus nuevas coordenadas:
\begin{equation}\label{eq:pdf2-ym-owner-refined-generator}
 \begin{aligned}
 \widehat H_{m+1}^{\rm ov}
 &=\widehat\Gamma_m^*
   (\widehat H_m^{\rm ov}\otimes I
    +I\otimes\widehat H_{m+1}^{\rm det})\widehat\Gamma_m,\\
 \widehat H_{m+1}^{\rm det}
 &=3\kappa_{\rm av}
   \sum_{\iota\in\mathcal I_{m+1/m}^{\rm det}}
   (I-E_{m+1,\iota}^{\rm det}).
 \end{aligned}
\end{equation}
La partición de colores no se deja nominal. Si una celda se indexa por
$n=(n_1,n_2,n_3,n_4)\in\mathbb Z^4$, se fija
\begin{equation}\label{eq:pdf2-ym-owner-729-coloring}
 \chi_m(n):=
 \bigl(n_1\bmod9,\ n_2\bmod9,\ n_3+3n_4\bmod9\bigr)
 \in(\mathbb Z/9\mathbb Z)^3,
 \qquad |(\mathbb Z/9\mathbb Z)^3|=9\cdot81.
\end{equation}
Si dos collares cerrados de radio celular uno se intersectan, su diferencia
$\delta$ satisface $|\delta_j|\le2$. La igualdad de colores fuerza
$\delta_1=\delta_2=0$ y
$\delta_3+3\delta_4\equiv0\pmod9$. Como el último entero pertenece a
$[-8,8]$, vale cero; las cotas $|\delta_3|,|\delta_4|\le2$ dan entonces
$\delta_3=\delta_4=0$. Por tanto, celdas distintas del mismo color tienen
collares disjuntos. Se define materialmente
\begin{equation}\label{eq:pdf2-ym-owner-729-blocks}
 \mathcal B_{m,r}:=\{\iota\in\mathcal I_m:
   \operatorname{cell}(\iota)=n, \chi_m(n)=r\},
 \qquad r\in(\mathbb Z/9\mathbb Z)^3.
\end{equation}
El transporte de una pantalla por $g\in E(4)$ transporta también esta
partición; no exige que una rotación preserve las coordenadas del rótulo.
Así, $\{\mathcal B_{m,r}\}_{r=1}^{9\cdot81}$ agrupa soportes espaciales
disjuntos;
el índice interno de cada bloque enumera sus factores de
\eqref{eq:pdf2-ym-owner-exhaustive-index}. Para un cilindro $F$, el producto sólo
contiene los bloques de los que depende $F$; la clausura define el producto fuerte
para toda la forma. El generador de un bloque conserva las multiplicidades de
Hoeffding; no se sustituye por un solo proyector que las borraría. Por
conmutatividad en cada carta,
\begin{equation}\label{eq:pdf2-ym-owner-local-circuit}
\begin{aligned}
 H_{m,B}&:=3\kappa_{\rm av}\sum_{\iota\in B}(I-E_{m,\iota}),\\
 K_{m,B}&:=e^{-a_mH_{m,B}}
 =\prod_{\iota\in B}e^{-3\kappa_{\rm av}a_m(I-E_{m,\iota})},\\
 K_m^{\rm loc}&:=\mathop{\prod_{r,B}}^{\rm strong}K_{m,B}
 =e^{-a_m\widehat H_m^{\rm ov}}.
\end{aligned}
\end{equation}
El diámetro físico de cada bloque permanece uniformemente acotado. De aquí procede
la localidad: el generador no introduce una actualización instantánea sobre toda la
red.
La partición es compatible con Bianchi, no sólo con disjunción geométrica. Para una
celda cúbica $c$ se fija el producto orientado, con todos los factores transportados
al mismo vértice,
\begin{equation}\label{eq:pdf2-ym-owner-bianchi-block-product}
 \mathscr B_{m,c}(x):=
 \prod_{p\subset\partial c}^{\longrightarrow}
 T_p\operatorname{Hol}_p(x)T_p^{-1}.
\end{equation}
Cada arista interna de la frontera aparece dos veces y con orientaciones contrarias;
por tanto $\mathscr B_{m,c}=e$ punto a punto. Un bloque
$\mathcal B_{m,r}$ contiene el collar completo de toda celda que actualiza: si toca
una cara de $c$, actualiza simultáneamente el par inverso que comparte esa arista.
La cancelación anterior sigue siendo literal después de cada factor del circuito:
\begin{equation}\label{eq:pdf2-ym-owner-bianchi-block-invariance}
 K_{m,B}\mathbf1_{\{\mathscr B_{m,c}=e\}}
 =\mathbf1_{\{\mathscr B_{m,c}=e\}},
 \qquad
 K_m^{\rm loc}\mathbf1_{\{\mathscr B_{m,c}=e\}}
 =\mathbf1_{\{\mathscr B_{m,c}=e\}}.
\end{equation}
Así el calendario de $9\cdot81$ capas preserva Bianchi en cada paso y en el producto
fuerte, no sólo al final de una barrida.

La acción gauge permuta o traslada las coordenadas integradas por cada
$E_{m,\iota}$ y
preserva sus medidas. Por ello
\begin{equation}\label{eq:pdf2-ym-owner-physical-reduction}
 U_m(g)\widehat H_m^{\rm ov}=\widehat H_m^{\rm ov}U_m(g)
 \quad(g\in\mathcal G_m^{\rm full}),
 \qquad
 [\mathbb P_m^{\rm phys},\widehat H_m^{\rm ov}]=0,
 \qquad
 \mathbb P_{m+1}^{\rm phys}\widehat J_m
 =\widehat J_m\mathbb P_m^{\rm phys}.
\end{equation}
El semigrupo completo y su kernel Markov se denotan por
\begin{equation}\label{eq:pdf2-ym-owner-full-kernel}
 \widehat K_{m,t}^{\rm ov}:=e^{-t\widehat H_m^{\rm ov}},
 \qquad
\widehat K_{m,t}^{\rm ov}(x,dy)
 \quad(x,y\in\widehat X_m).
\end{equation}
La primera identidad de \eqref{eq:pdf2-ym-owner-physical-reduction} implica
\begin{equation}\label{eq:pdf2-ym-owner-full-kernel-equivariance}
 \widehat K_{m,t}^{\rm ov}(gx,gA)
 =\widehat K_{m,t}^{\rm ov}(x,A)
 \quad(g\in\mathcal G_m^{\rm full}).
\end{equation}
Para tipar la reducción sin usar un kernel restringido sobre el espacio equivocado,
sea $q_m:\widehat X_m\to Y_m:=\widehat X_m/\mathcal G_m^{\rm full}$ el mapa de
órbitas y $\nu_m=(q_m)_*\widehat\mu_m^{\rm ov}$. El pullback
\begin{equation}\label{eq:pdf2-ym-owner-quotient-unitary}
 \mathcal W_m:L^2(Y_m,\nu_m)\longrightarrow
 \widehat{\mathscr H}_m^{\rm phys},
 \qquad \mathcal W_m\phi=\phi\circ q_m
\end{equation}
es unitario. La equivariancia del kernel completo permite definir
\begin{equation}\label{eq:pdf2-ym-owner-physical-kernel}
 K_{m,t}^{\rm phys}(q_mx,A)
 :=\widehat K_{m,t}^{\rm ov}(x,q_m^{-1}A),
 \qquad A\subseteq Y_m,
\end{equation}
independientemente del representante $x$. El objeto físico queda entonces tipado por
\begin{equation}\label{eq:pdf2-ym-owner-D}
 D_m^{\rm YM}
 :=\left.\widehat H_m^{\rm ov}\right|_{\widehat{\mathscr H}_m^{\rm phys}},
 \qquad
 e^{-tD_m^{\rm YM}}
 =\mathcal W_mK_{m,t}^{\rm phys}\mathcal W_m^{-1}.
\end{equation}
Éste es el generador producto completo; no es el generador de una sola fibra de
\eqref{eq:pdf2-ym-fibra-semigrupo}.

\subsection{Acción de Yang--Mills, medida dependiente de $g$ y transporte exacto}

Hasta aquí se ha escrito la carta cinemática de referencia. La denotaremos
$\widehat\mu_m^0:=\widehat\mu_m^{\rm ov}$ y
$\widehat H_m^0:=\widehat H_m^{\rm ov}$, con
$\widehat K_{m,t}^0:=e^{-t\widehat H_m^0}=\widehat K_{m,t}^{\rm ov}$.
Para convertir el lector estructural
$\mathcal P_m^{F^2}$ en una ley, y no en una comparación posterior, se descarga
primero su raíz positiva sobre la factorización triangular. La polarización de
$\mathcal P_m^{F^2}$ en un collar $b$ es el Gram positivo
\[
 Q_{m,b}(u,v):=\frac14\bigl(
 \mathcal P_{m,b}^{F^2}(u+v)-\mathcal P_{m,b}^{F^2}(u-v)
 \bigr).
\]
Se divide el espacio finito de lectores del collar por $\ker Q_{m,b}$, se completa
con $Q_{m,b}$ y se denota por $\mathcal E_{m,b}^{F}$ el Hilbert resultante. Si
$\zeta_{m,b}(x_b)$ es el vector de lectores centrados de la coordenada triangular
$x_b$, se fija materialmente
\begin{equation}\label{eq:pdf2-ym-action-root}
 \xi_{m,b}(x_b):=Q_{m,b}^{1/2}\zeta_{m,b}(x_b)
 \in\mathcal E_{m,b}^{F},
 \qquad
 p_{m,b}(x_b):=\|\xi_{m,b}(x_b)\|^2.
\end{equation}
El índice $b$ recorre los collares Gram de soporte mínimo. Un collar se asigna
al primer color de \eqref{eq:pdf2-ym-owner-729-coloring} que contiene su celda base;
así cada término aparece una vez y los collares de un mismo color son disjuntos.
Más precisamente, la diagonalización positiva agrupa los factores triangulares en
una partición disjunta
\begin{equation}\label{eq:pdf2-ym-action-factor-partition}
 \mathcal I_m^{F}
 =\bigsqcup_{r\in(\mathbb Z/9\mathbb Z)^3}
   \bigsqcup_{b\in\mathcal B_{m,r}^{F}}\mathcal I_{m,b},
 \qquad x_b:=(x_\iota)_{\iota\in\mathcal I_{m,b}},
 \qquad
 \mathcal B_{m,r}^{F}:=\{b:Q_{m,b}\ne0,\ \chi_m(b)=r\}.
\end{equation}
Ninguna coordenada triangular pertenece a dos $\mathcal I_{m,b}$; si un lector
geométrico toca dos collares, su modo de Gram se asigna al primero y su complemento
ortogonal al segundo. Ésta es la descomposición espectral finita de la forma positiva,
no una duplicación de la variable física.
La raíz de la suma ortogonal, calculada color por color, da la identidad, no una
aproximación,
\begin{equation}\label{eq:pdf2-ym-action-reader-sum}
 \mathcal P_m^{F^2}(x)
 =\sum_{r\in(\mathbb Z/9\mathbb Z)^3}
   \sum_{b\in\mathcal B_{m,r}^{F}}p_{m,b}(x_b),
 \qquad
 S_{m,g}(x):=\frac1{4g^2}\mathcal P_m^{F^2}(x),
 \quad g\in(0,\infty).
\end{equation}
Aquí $x_b:\widehat X_m\to X_{m,b}$ es la proyección de coordenadas del
collar, de modo que dominio, acción y codominio de cada sumando quedan fijados.
El lector es gauge invariante porque $Q_{m,b}$ usa el producto invariante de
$\mathfrak g$; es invariante por reflexión porque una reflexión sólo permuta los
seis pares orientados; y es natural por relabelado euclídeo.
La partición \eqref{eq:pdf2-ym-action-factor-partition} y la carta producto dan
explícitamente
\begin{equation}\label{eq:pdf2-ym-action-factor-measures}
 \widehat\mu_m^0=\lambda_m^{\rm inert}\otimes
   \bigotimes_{r,b}\lambda_{m,b}^0,
 \qquad
 d\lambda_{m,b}^g
 :=(Z_{m,b}^g)^{-1}e^{-p_{m,b}/(4g^2)}d\lambda_{m,b}^0,
 \qquad
 \widehat\mu_{m,g}^{\rm YM}=\lambda_m^{\rm inert}\otimes
   \bigotimes_{r,b}\lambda_{m,b}^g.
\end{equation}
El factor inerte contiene exactamente las coordenadas sobre las que el lector es
cero; no se elimina de la teoría.

La separación grueso--detalle de
\eqref{eq:pdf2-ym-owner-product-factorization} es ortogonal para la misma raíz:
la polarización directa da
$Q_{m+1}=Q_m\oplus Q_{m+1/m}^{\rm det}$ y
$\zeta_{m+1}=\zeta_m\oplus\zeta_{m+1/m}^{\rm det}$.
Si $x_{m+1}=\widehat\Gamma_m^{-1}(x_m,u)$, entonces
\begin{equation}\label{eq:pdf2-ym-action-cocycle}
 \xi_{m+1}(x_{m+1})
 =\xi_m(x_m)\oplus\xi_{m+1/m}^{\rm det}(u),
 \qquad
 S_{m+1,g}\circ\widehat\Gamma_m^{-1}
 =S_{m,g}+S_{m+1/m,g}^{\rm det}.
\end{equation}
Esta es la forma precisa en que carry y memoria impiden contar dos veces una cara
ancestral: la parte gruesa no se vuelve a sumar en el detalle. En particular,
$Z_{m+1,g}=Z_{m,g}Z_{m+1/m,g}^{\rm det}$ para los normalizadores y se obtiene la
familia de probabilidades genuinamente dependiente del acoplamiento
\begin{equation}\label{eq:pdf2-ym-gibbs-measure}
 d\widehat\mu_{m,g}^{\rm YM}(x)
 :=Z_{m,g}^{-1}e^{-S_{m,g}(x)}d\widehat\mu_m^0(x),
 \qquad
 (\widehat\pi_{m+1,m})_*\widehat\mu_{m+1,g}^{\rm YM}
 =\widehat\mu_{m,g}^{\rm YM}.
\end{equation}
Para escribir correctamente Schwinger--Dyson respecto de Haar, y no fingir que la
medida de referencia es plana, sea $\lambda_m$ el producto de Haar, Lebesgue y
conteo uniforme de la carta levantada, y sea
$R_m:=d\widehat\mu_m^0/d\lambda_m>0$. La acción total regulada es
\begin{equation}\label{eq:pdf2-ym-total-regulated-action}
 \mathscr S_{m,g}:=S_{m,g}-\log R_m,
 \qquad
 d\widehat\mu_{m,g}^{\rm YM}
 =\mathscr Z_{m,g}^{-1}e^{-\mathscr S_{m,g}}d\lambda_m.
\end{equation}
El término $-\log R_m$ contiene los kernels de calor y las cartas de puente de la
estructura gauge de partida; es independiente de $g$. El término de Yang--Mills que
cambia el acoplamiento y cuyo límite clásico se calcula abajo es exactamente
$S_{m,g}=\mathcal P_m^{F^2}/(4g^2)$.
La consistencia de la segunda igualdad y la estandaridad de los espacios finitos
permiten aplicar Kolmogórov y definen la única medida cilíndrica
$\widehat\mu_{\infty,g}^{\rm YM}$ cuyas marginales son
$\widehat\mu_{m,g}^{\rm YM}$.
No se ha usado aquí ningún autovalor: primero se han fijado $g$, la acción y
la medida.

Para definir la dinámica sobre esa medida sin perder producto, gauge ni las
identidades superficiales, se construye el transporte, no se lo postula. En cada
factor no atómico $X_{m,b}$ se fija una carta boreliana
$\rho_{m,b}:X_{m,b}/\mathcal G_{m,b}\to(0,1)$ y se desintegra primero sobre el
cociente gauge y después sobre la órbita de Haar. Sean $F_{m,b}^0$ y
$F_{m,b}^g$ las funciones de distribución continuas de $\rho_{m,b}$ para la
medida de referencia y la medida reponderada. El transporte triangular es
\begin{equation}\label{eq:pdf2-ym-gibbs-transport-factor}
 t_{m,b,g}:=(F_{m,b}^g)^{-1}\circ F_{m,b}^0,
 \qquad
 T_{m,b,g}(y,h):=(t_{m,b,g}(y),h),
 \qquad
 (T_{m,b,g})_*\lambda_{m,b}^0=\lambda_{m,b}^g.
\end{equation}
En un factor puramente atómico sobre el que $p_{m,b}=0$ se toma
$T_{m,b,g}=I$. En presencia de átomos y una coordenada continua, la carta es el
orden lexicográfico átomo--intervalo y la misma fórmula se aplica al espacio
estándar no atómico total. Los transportes de los $9\cdot81$ colores se componen
en el orden fijado por \eqref{eq:pdf2-ym-owner-729-coloring}; dentro de un color
conmutan porque los collares son disjuntos. Se obtiene
\begin{equation}\label{eq:pdf2-ym-gibbs-transport-global}
 T_{m,g}:=\prod_{r\in(\mathbb Z/9\mathbb Z)^3}
           \prod_{b\in\mathcal B_{m,r}^{F}}T_{m,b,g},
 \quad
 (T_{m,g})_*\widehat\mu_m^0=\widehat\mu_{m,g}^{\rm YM},
 \quad
 \widehat\pi_{m+1,m}T_{m+1,g}=T_{m,g}\widehat\pi_{m+1,m}.
\end{equation}
La última igualdad se comprueba factor a factor usando
\eqref{eq:pdf2-ym-action-cocycle}: el transporte del factor grueso es el ya fijado y
los transportes nuevos actúan sólo en la innovación. Las cartas de cociente se
eligen en una pantalla representante y se trasladan; por ello $T_{m,g}$ conmuta con
gauge y reflexión y, para $a\in E(4)$, satisface
$T_{a\alpha,g}=U_\alpha(a)T_{\alpha,g}U_\alpha(a)^{-1}$.

La composición
\begin{equation}\label{eq:pdf2-ym-full-probability-isomorphism}
 \mathcal T_{m,g}:L^\infty(\widehat X_m,\widehat\mu_{m,g}^{\rm YM})
 \longrightarrow L^\infty(\widehat X_m,\widehat\mu_m^0),
 \qquad \mathcal T_{m,g}F:=F\circ T_{m,g}
\end{equation}
es, por construcción, un isomorfismo de espacios de probabilidad y de
$*$-álgebras, y se prolonga unitariamente a $L^2$. En particular,
\begin{equation}\label{eq:pdf2-ym-full-algebra-identities}
 \mathcal T_{m,g}(FG)=\mathcal T_{m,g}F\,\mathcal T_{m,g}G,
 \quad
 \mathcal T_{m,g}(F^*)=(\mathcal T_{m,g}F)^*,
 \quad
 \int F\,d\widehat\mu_{m,g}^{\rm YM}
 =\int\mathcal T_{m,g}F\,d\widehat\mu_m^0.
\end{equation}
El dominio de cada $T_{m,b,g}$ es el propio espacio de holonomías que satisface
la ley de subdivisión y la restricción de Bianchi. Por tanto su imagen permanece
en ese espacio. Aplicando la multiplicatividad anterior generador a generador se
obtienen las identidades puntuales
\begin{equation}\label{eq:pdf2-ym-transport-surface-bianchi}
 \begin{aligned}
 \mathcal T_{m,g}(\operatorname{Hol}_p)
 &=\prod_{p'\subset p}^{\longrightarrow}
   \mathcal T_{m,g}(T_{p'}\operatorname{Hol}_{p'}T_{p'}^{-1}),\\
 \mathcal T_{m,g}(F_{{\rm cl},m})
 &=F_{{\rm cl},m}\circ T_{m,g},\\
 \prod_{p\subset\partial c}^{\longrightarrow}
 \mathcal T_{m,g}(T_p\operatorname{Hol}_pT_p^{-1})&=e,
 \qquad
 \operatorname{Alt}_{\lambda\mu\nu}
 D_\lambda(F_{{\rm cl},m}\circ T_{m,g})_{\mu\nu}=0.
 \end{aligned}
\end{equation}
La tercera línea es la cancelación arista por arista de la frontera orientada de
la frontera de $c$; la cuarta es su polarización infinitesimal. Así el transporte
preserva productos, estado, clover y Bianchi, no sólo el primer caos.

Finalmente se transportan expectativas, generador, kernel, proyección e inclusión:
\begin{equation}\label{eq:pdf2-ym-g-dynamics}
 \begin{aligned}
 E_{m,\iota}^{g}&:=\mathcal T_{m,g}^{-1}E_{m,\iota}\mathcal T_{m,g},&
 \widehat H_{m,g}^{\rm YM}&:=\mathcal T_{m,g}^{-1}\widehat H_m^0\mathcal T_{m,g},\\
 \widehat K_{m,t}^{g}&:=\mathcal T_{m,g}^{-1}\widehat K_{m,t}^{0}\mathcal T_{m,g},&
 \mathbb P_{m,g}^{\rm phys}&:=\mathcal T_{m,g}^{-1}
     \mathbb P_m^{\rm phys}\mathcal T_{m,g},\\
 \widehat J_{m,g}&:=\mathcal T_{m+1,g}^{-1}
     \widehat J_m\mathcal T_{m,g},&
 D_{m,g}^{\rm YM}&:=\left.\widehat H_{m,g}^{\rm YM}
     \right|_{\operatorname{Ran}\mathbb P_{m,g}^{\rm phys}}.
 \end{aligned}
\end{equation}
La primera línea no es sólo una conjugación abstracta. Usando
\eqref{eq:pdf2-ym-action-factor-measures}, su acción sobre un cilindro es la
actualización de baño térmico de la acción:
\begin{equation}\label{eq:pdf2-ym-gibbs-conditional-expectation}
 (E_{m,b}^{g}F)(x_{\ne b})
 =\frac{\displaystyle
   \int_{X_{m,b}}F(x_{\ne b},u)e^{-p_{m,b}(u)/(4g^2)}
       d\lambda_{m,b}^0(u)}{\displaystyle
   \int_{X_{m,b}}e^{-p_{m,b}(u)/(4g^2)}d\lambda_{m,b}^0(u)},
 \qquad
 \widehat H_{m,g}^{\rm YM}
 =3\kappa_{\rm av}\sum_b(I-E_{m,b}^{g}).
\end{equation}
Así tanto la medida como cada transición local contienen $g$ y la acción
$p_{m,b}/(4g^2)$ de forma explícita.
Son Markov porque $\mathcal T_{m,g}\mathbf1=\mathbf1$, locales porque cada
$T_{m,b,g}$ tiene el mismo collar que $p_{m,b}$, y dependen de $g$ por
\eqref{eq:pdf2-ym-gibbs-transport-factor}. La equivalencia es espectral y no
selecciona la medida a partir del espectro:
Para cerrar además las ecuaciones dinámicas, sea $X_{m,b}^a$ el campo
izquierdo-invariante en el factor de enlace $b$ asociado a
$e_a\in\mathfrak g$; en una carta continua de innovación se usa la derivada
direccional ordinaria, y las coordenadas finitas se tratan por diferencia exacta.
Para todo cilindro suave $F$ el dominio común es
$\mathscr C_m^1:=\bigcap_{b,a}\operatorname{Dom}X_{m,b}^a$, tomando soporte
compacto en las cartas abiertas de innovación; en los factores de grupo no hay
frontera. La invariancia de Haar y la regla del producto calculan
\begin{equation}\label{eq:pdf2-ym-schwinger-dyson-finite}
 \int X_{m,b}^aF\,d\widehat\mu_{m,g}^{\rm YM}
 =\int F\,X_{m,b}^a\mathscr S_{m,g}\,d\widehat\mu_{m,g}^{\rm YM},
 \qquad F\in\mathscr C_m^1.
\end{equation}
En una coordenada finita, para
$\Delta_\varepsilon F(q):=F(q+\varepsilon)-F(q)$, el cambio de variable
$q\mapsto q+\varepsilon$ da la versión exacta, sin derivada formal,
\begin{equation}\label{eq:pdf2-ym-schwinger-dyson-discrete}
 \int\Delta_\varepsilon F(q)\,d\widehat\mu_{m,g}^{\rm YM}
 =\int F(q)
  \left(e^{\mathscr S_{m,g}(q)-\mathscr S_{m,g}(q-\varepsilon)}-1\right)
  d\widehat\mu_{m,g}^{\rm YM}(q).
\end{equation}
Si $\delta_\varepsilon$ es la suma orientada de esos campos en las aristas que
inciden en un vértice, la invariancia gauge de la medida de referencia y de
$S_{m,g}$ da $\delta_\varepsilon\mathscr S_{m,g}=0$ y, por tanto, la identidad de Ward
\begin{equation}\label{eq:pdf2-ym-ward-finite}
 \int\delta_\varepsilon F\,d\widehat\mu_{m,g}^{\rm YM}=0.
\end{equation}
Para el factor gauge finito $\mathbb F_3^4$, la invariancia
$\mathscr S_{m,g}(q+\varepsilon)=\mathscr S_{m,g}(q)$ reduce
\eqref{eq:pdf2-ym-schwinger-dyson-discrete} a la misma identidad Ward con
$\delta_\varepsilon$ reemplazada por $\Delta_\varepsilon$.
Las dos expresiones son compatibles con $\widehat J_{m,g}$ porque
\eqref{eq:pdf2-ym-action-cocycle} separa las derivadas gruesas de las de detalle.
Todo cilindro del límite reside en algún nivel; así
\eqref{eq:pdf2-ym-schwinger-dyson-finite} y
\eqref{eq:pdf2-ym-ward-finite} pasan literalmente al estado cofinal. Ni la
ecuación Schwinger--Dyson ni Ward contienen la brecha como premisa.

\begin{equation}\label{eq:pdf2-ym-g-gap-before-limit}
 D_{m,g}^{\rm YM}\succeq3\kappa_{\rm av}
 (I-P_{\Omega_{m,g}}),
 \qquad
 W_{c,g}:=\mathcal T_{m,g}^{-1}W_c,
 \qquad
 D_{m,g}^{\rm YM}W_{c,g}=3\kappa_{\rm av}W_{c,g}.
\end{equation}
En lo que sigue se fija un $g\in(0,\infty)$ y se suprime ese subíndice; las
medidas, expectativas, inclusiones y generadores sin superíndice son los objetos
de \eqref{eq:pdf2-ym-g-dynamics}. Las holonomías geométricas permanecen como
funciones coordenadas canónicas sobre el espacio de configuraciones y, por ello, su
ley sí cambia con $g$; $\mathcal T_{m,g}$ las lleva a funciones compuestas de la
carta de referencia. Sólo el testigo espectral se transporta explícitamente en
\eqref{eq:pdf2-ym-g-gap-before-limit}. Esta convención no borra $g$: su residencia explícita es
\eqref{eq:pdf2-ym-gibbs-measure}--\eqref{eq:pdf2-ym-g-dynamics}.

\subsection{Vacío simple y testigo que alcanza la brecha}

La intersección de todas las expectativas del estado completo es la constante:
\begin{equation}\label{eq:pdf2-ym-owner-vacuum}
 \bigcap_{\iota\in\mathcal I_m}\operatorname{Ran}E_{m,\iota}
 =\mathbb C\Omega_m.
\end{equation}
Esta identidad no se introduce como hipótesis. En la carta producto
\eqref{eq:pdf2-ym-owner-product-factorization}, cada $E_{m,\iota}$ integra una
coordenada del índice exhaustivo
\eqref{eq:pdf2-ym-owner-exhaustive-index} y deja fijas las restantes. Si
$u$ pertenece a la intersección, sus expectativas condicionales respecto de todo
cilindro finito son constantes. La convergencia martingala de esos cilindros a la
$\sigma$-álgebra completa da que $u$ es constante casi en todas partes. La inclusión
inversa es inmediata porque toda expectativa conserva la unidad. Esto prueba
\eqref{eq:pdf2-ym-owner-vacuum} y, después de la proyección gauge, deja un vacío
físico unidimensional.
La descomposición de Hoeffding de
\eqref{eq:pdf2-ym-owner-base-generator}--
\eqref{eq:pdf2-ym-owner-refined-generator} da
\begin{equation}\label{eq:pdf2-ym-owner-hoeffding}
 \widehat{\mathcal E}_m(u)
 =3\kappa_{\rm av}\sum_{S\Subset\mathcal I_m}|S|\,\|u_S\|^2,
 \qquad
 D_m^{\rm YM}\succeq
 3\kappa_{\rm av}(I-P_{\Omega_m}).
\end{equation}
La función
\begin{equation}\label{eq:pdf2-ym-owner-wc}
 W_c(q):=\mathbf1_{\{q_1+\cdots+q_6=0\}}-\frac13
\end{equation}
es invariante bajo $G^{V(\Lambda_m)}$ porque no depende de enlaces ni marcos. Para
la acción de incidencia, cada $x_i$ aparece tres veces, luego
\begin{equation}\label{eq:pdf2-ym-owner-wc-invariance}
 \sum_{i<j}(B^+x)_{ij}=3\sum_i x_i=0
 \quad\text{en }\mathbb F_3.
\end{equation}
Así $\mathbb P_m^{\rm phys}W_c=W_c$. Bajo $u_6$, la suma de los seis trits es
uniforme; por cálculo directo,
\begin{equation}\label{eq:pdf2-ym-owner-wc-moments}
 \mathbb EW_c=0,\qquad
 \|W_c\|^2=\frac29,\qquad
 \mathbb E(W_c^3)=\frac2{27}.
\end{equation}
Sólo el proyector $I-E_{m,q_c}$ actúa no trivialmente sobre $W_c$. Por tanto
\begin{equation}\label{eq:pdf2-ym-owner-wc-eigenvalue}
 D_m^{\rm YM}W_c=3\kappa_{\rm av}W_c.
\end{equation}
La cota inferior de \eqref{eq:pdf2-ym-owner-hoeffding} no es sólo una cota: el
complemento físico contiene un vector no nulo que alcanza el umbral.

\subsection{Refinamiento nonádico y persistencia en el límite}

Las inclusiones físicas satisfacen
\begin{equation}\label{eq:pdf2-ym-owner-intertwining}
 D_{m+1}^{\rm YM}\widehat J_m
 =\widehat J_mD_m^{\rm YM},\qquad
 \bigl(e^{-a_{m+1}D_{m+1}^{\rm YM}}\bigr)^9\widehat J_m
 =\widehat J_me^{-a_mD_m^{\rm YM}},\qquad a_{m+1}=a_m/9.
\end{equation}
En el límite inductivo
\begin{equation}\label{eq:pdf2-ym-owner-inductive-limit}
 \widehat{\mathscr H}_\infty^{\rm phys}
 =\varinjlim(\widehat{\mathscr H}_m^{\rm phys},\widehat J_m)
\end{equation}
se define sobre la unión algebraica
\begin{equation}\label{eq:pdf2-ym-owner-limit-form}
 \mathcal E_\infty[\widehat J_{m,\infty}u]
 :=\langle u,D_m^{\rm YM}u\rangle.
\end{equation}
Su generador se denota
$D_{\infty,g}^{\rm YM}$; de acuerdo con la convención posterior a
\eqref{eq:pdf2-ym-g-gap-before-limit}, se escribe también
$D_\infty^{\rm YM}:=D_{\infty,g}^{\rm YM}$.
El primer entrelazamiento de \eqref{eq:pdf2-ym-owner-intertwining} hace que la
definición sea independiente del representante. Las truncaciones de Hoeffding son
un núcleo de forma y dan una sucesión de recuperación; Fatou en la suma no negativa
da la semicontinuidad inferior. La clausura posee, pues,
\begin{equation}\label{eq:pdf2-ym-owner-limit-gap}
 \overline{\mathcal E}_\infty[u]
 \ge3\kappa_{\rm av}\|(I-P_{\Omega_\infty})u\|^2.
\end{equation}
La inclusión conserva las coordenadas ancestrales, de modo que
$\widehat J_{m,\infty}W_c$ sigue siendo no nulo y satisface
\begin{equation}\label{eq:pdf2-ym-owner-limit-witness}
 D_\infty^{\rm YM}\widehat J_{m,\infty}W_c
 =3\kappa_{\rm av}\widehat J_{m,\infty}W_c.
\end{equation}
En consecuencia,
\begin{equation}\label{eq:pdf2-ym-owner-exact-gap}
 \inf\bigl(\operatorname{Spec}D_\infty^{\rm YM}\setminus\{0\}\bigr)
 =3\kappa_{\rm av}>0,
 \qquad \ker D_\infty^{\rm YM}=\mathbb C\Omega_\infty.
\end{equation}

\subsection{Cuatro reconstrucciones OS de la misma dinámica}

El kernel completo de \eqref{eq:pdf2-ym-owner-full-kernel} define una ley reversible
sobre $\widehat X_m$. Para la reconstrucción física se usa su descenso bien tipado
\eqref{eq:pdf2-ym-owner-physical-kernel}: sobre el cociente $Y_m$ se define
\begin{equation}\label{eq:pdf2-ym-owner-eight-face-law}
 \Omega_{m,{\rm phys}}^{(8)}(d\bar y_1\cdots d\bar y_8)
 =\int\nu_m(d\bar z)
  \prod_{r=1}^8K_{m,a_m/2}^{\rm phys}(\bar z,d\bar y_r).
\end{equation}
Para cada dirección $\nu=1,\ldots,4$, la reflexión intercambia las cuatro caras
positivas con las negativas. Condicionando respecto del estado central $z$ y usando
reversibilidad,
\begin{equation}\label{eq:pdf2-ym-owner-four-os}
 \int\overline{F(\Theta_{\nu,m}y)}F(y)\,
 d\Omega_{m,{\rm phys}}^{(8)}(y)
 =\|\mathcal B_{m,\nu}F\|_2^2\ge0.
\end{equation}
La marginal de la pareja opuesta normal a $\nu$ es
\begin{equation}\label{eq:pdf2-ym-owner-os-marginal}
 \nu_m(d\bar x)K_{m,a_m}^{\rm phys}(\bar x,d\bar y).
\end{equation}
Para identificar el Hamiltoniano de una dirección se toma el subespacio OS de
observables de una cara positiva simple,
$F(y)=F_\nu(\bar y_{\nu,+})$; los observables de varias caras conservan la
positividad anterior, pero no se confunden con este espacio de transferencia. En el
cociente de cara simple, el mapa
\begin{equation}\label{eq:pdf2-ym-owner-os-unitary}
 \mathcal V_{m,\nu}[F]
 :=e^{-a_mD_m^{\rm YM}/2}\mathcal W_mF
\end{equation}
es isométrico para la norma OS y tiene rango denso. Su extensión identifica el
cociente OS con $\widehat{\mathscr H}_m^{\rm phys}$. Póngase
$\overline D_m:=\mathcal W_m^{-1}D_m^{\rm YM}\mathcal W_m$ y
$\overline J_m:=\mathcal W_{m+1}^{-1}\widehat J_m\mathcal W_m$. El
entrelazamiento y $a_m=9a_{m+1}$ dan
\begin{equation}\label{eq:pdf2-ym-owner-os-range-inclusion}
 \overline J_me^{-a_m\overline D_m/2}
 =e^{-a_m\overline D_{m+1}/2}\overline J_m
 =\bigl(e^{-a_{m+1}\overline D_{m+1}/2}\bigr)^9\overline J_m.
\end{equation}
El último miembro pertenece al rango de
$e^{-a_{m+1}\overline D_{m+1}/2}$; por eso el inverso de $\mathcal V_{m+1,\nu}$
que aparece a continuación se aplica sólo en su rango y está bien definido. Los
conectores
\begin{equation}\label{eq:pdf2-ym-owner-os-connectors}
 \mathcal J_{m,\nu}^{\rm OS}
 :=\mathcal V_{m+1,\nu}^{-1}\widehat J_m\mathcal V_{m,\nu}
\end{equation}
son isometrías y entrelazan los cuatro Hamiltonianos reconstruidos:
\begin{equation}\label{eq:pdf2-ym-owner-common-os-hamiltonian}
 H_{{\rm OS},m+1}^{(\nu)}\mathcal J_{m,\nu}^{\rm OS}
 =\mathcal J_{m,\nu}^{\rm OS}H_{{\rm OS},m}^{(\nu)},
 \qquad
 \mathcal V_{\infty,\nu}H_{{\rm OS},\infty}^{(\nu)}
 \mathcal V_{\infty,\nu}^{-1}=D_\infty^{\rm YM}.
\end{equation}
Las cuatro reflexiones no fabrican cuatro medidas: son cuatro publicaciones del mismo
estado, del mismo kernel y del mismo generador.

\subsection{Localidad, acción euclídea y publicación de curvatura}

La categoría dirigida de pantallas incluye refinamientos, overlays y transportes
euclídeos. Si $a\in E(4)$ y $\alpha$ es una pantalla, el relabelado completo define
\begin{equation}\label{eq:pdf2-ym-owner-e4-finite-transport}
 U_\alpha(a):\widehat{\mathscr H}_\alpha^{\rm phys}
 \longrightarrow\widehat{\mathscr H}_{a\alpha}^{\rm phys},
 \qquad
 U_\beta(a)\widehat J_{\alpha\beta}
 =\widehat J_{a\alpha,a\beta}U_\alpha(a).
\end{equation}
El transporte preserva medida, generador, ideales nulos y soportes. La identidad de
naturalidad permite definir, sin elegir representante, una unitaria sobre el colímite
algebraico:
\begin{equation}\label{eq:pdf2-ym-owner-e4-colimit-action}
 U(a)\widehat J_{\alpha,\infty}F
 :=\widehat J_{a\alpha,\infty}U_\alpha(a)F.
\end{equation}
Los relabelados satisfacen $U_{a\alpha}(b)U_\alpha(a)=U_\alpha(ba)$ y preservan
norma; por ello \eqref{eq:pdf2-ym-owner-e4-colimit-action} se prolonga a una
representación unitaria de $E(4)$.

\subsubsection{Curvatura local del canal de incidencia}

Todo grupo compacto, conectado y simple posee un toro maximal no trivial. Se fija
$t_3\in G$ de orden exactamente tres. Para un collar $c$, sea
\begin{equation}\label{eq:pdf2-ym-owner-incidence-flux}
 \omega_c(q):=\sum_{i<j}q_{ij}\in\mathbb F_3.
\end{equation}
La identidad \eqref{eq:pdf2-ym-incidence-conservation} prueba que $\omega_c$ es
invariante bajo la acción finita. Si $h_c$ es el marco en el vértice base del collar,
se define la holonomía geométrica y, para el testigo espectral, su imagen transportada
en la ley a acoplamiento $g$:
\begin{equation}\label{eq:pdf2-ym-owner-incidence-holonomy}
 \operatorname{Hol}_{m,c}^{\rm inc,geom}:=h_ct_3^{\,\omega_c(q)}h_c^{-1}\in G,
 \qquad
 \operatorname{Hol}_{m,c,g}^{\rm inc,spec}
 :=\mathcal T_{m,g}^{-1}\operatorname{Hol}_{m,c}^{\rm inc,geom},
 \qquad
 \operatorname{Hol}_{m,c}^{\rm inc}:=\operatorname{Hol}_{m,c,g}^{\rm inc,spec}.
\end{equation}
Bajo $G^{V_m}$ se transforma por conjugación y bajo
$\prod_c\mathbb F_3^4$ permanece fija. Como $t_3$ tiene orden tres, el cálculo
funcional de esta variable de tres valores da la identidad literal
\begin{equation}\label{eq:pdf2-ym-owner-wc-curvature-functional}
 W_{c,g}:=\mathcal T_{m,g}^{-1}W_c^0
 =\mathbf1_{\{e\}}\!\left(\operatorname{Hol}_{m,c}^{\rm inc}\right)-\frac13.
\end{equation}
Se escribe $W_c:=W_{c,g}$ en el resto del capítulo.
Como $T_{m,b,g}$ actúa dentro del cociente del mismo collar, la holonomía espectral
es una función de su álgebra local de holonomías geométricas; no añade una variable
ni cambia el soporte. Las holonomías geométricas, usadas en $S_{m,g}$ y en el clover,
mantienen su distribución dependiente de $g$.
El proyector del miembro derecho es una función de clase local; por tanto $W_c$ no
es una coordenada auxiliar desacoplada, sino un observable de la holonomía de
incidencia de la misma restricción gauge.

Para un abierto acotado $O\subset\mathbb R^4$, sea
$\mathfrak A_{\rm YM}(O)$ el álgebra generada por funciones de clase acotadas de
\eqref{eq:pdf2-ym-owner-incidence-holonomy}, holonomías de plaqueta y lectores
clover con soporte en $O$. Se define el sector de curvatura física por
\begin{equation}\label{eq:pdf2-ym-owner-curvature-cyclic-sector}
 \mathscr H_{\rm YM}^{\rm curv}
 :=\overline{\bigcup_{O\Subset\mathbb R^4}
   \mathfrak A_{\rm YM}(O)\Omega_\infty}
 \subseteq\widehat{\mathscr H}_\infty^{\rm phys}.
\end{equation}
Por \eqref{eq:pdf2-ym-owner-wc-curvature-functional}, cada
$\widehat J_{m,\infty}W_c$ pertenece materialmente a este sector.
Las expectativas de coordenada envían cilindros de curvatura a cilindros de
curvatura; por tanto el semigrupo deja invariante
$\mathscr H_{\rm YM}^{\rm curv}$. Al ser autoadjunto, el sector es reductor. La
cota de \eqref{eq:pdf2-ym-owner-limit-gap} se restringe a él y el vector anterior
alcanza la igualdad:
\begin{equation}\label{eq:pdf2-ym-owner-curvature-sector-gap}
 D_{\rm curv,g}^{\rm YM}
 :=\left.D_{\infty,g}^{\rm YM}\right|_{\mathscr H_{\rm YM}^{\rm curv}},
 \qquad D_{\rm curv}^{\rm YM}:=D_{\rm curv,g}^{\rm YM},
 \qquad
 \inf\bigl(\operatorname{Spec}D_{\rm curv,g}^{\rm YM}\setminus\{0\}\bigr)
 =3\kappa_{\rm av}.
\end{equation}

\subsubsection{Sistema cofinal de publicación y productos cruzados}

La inclusión $\widehat J_{\alpha\beta}$ conserva una coordenada ancestral; no debe
confundirse con el promedio de bloque que publica un campo continuo. Para escribir
este segundo mapa sin alterar la genealogía, sea $\mathfrak P$ el conjunto dirigido
de pantallas celulares acotadas, cerrado bajo overlay, refinamiento nonádico y
transporte euclídeo. Para $\alpha\in\mathfrak P$ póngase
\begin{equation}\label{eq:pdf2-ym-owner-cell-step-space}
 V_\alpha:=\operatorname{span}\left\{
 e_{\alpha,c}:=|c|^{-1/2}\mathbf1_c:
 c\in\mathcal C_\alpha^{\rm cell}\right\}
 \subset L^2(\mathbb R^4),
 \qquad P_\alpha:L^2(\mathbb R^4)\to V_\alpha.
\end{equation}
Los overlays hacen que $P_\alpha h\to h$ para todo $h\in L^2$ a lo largo de
$\mathfrak P$. En el lado gauge se normaliza el testigo por
\begin{equation}\label{eq:pdf2-ym-owner-normalized-cell-witness}
 Z_{\alpha,c}:=\sqrt{\frac92}\,W_{\alpha,c},
 \qquad
 \langle Z_{\alpha,c},Z_{\alpha,d}\rangle=\delta_{cd},
 \qquad
 D_\alpha^{\rm YM}Z_{\alpha,c}=3\kappa_{\rm av}Z_{\alpha,c}.
\end{equation}
La segunda igualdad usa exactamente los momentos de
\eqref{eq:pdf2-ym-owner-wc-moments} y la factorización de coordenadas distintas.
Sea $\mathcal K_\alpha^{\rm cell}$ el espacio generado por esos vectores. Si
$\beta\succeq\alpha$, el entrelazador de publicación, distinto de
$\widehat J_{\alpha\beta}$, se define generador a generador por
\begin{equation}\label{eq:pdf2-ym-owner-block-intertwiner}
 I_{\alpha\beta}^{\rm curv}Z_{\alpha,c}
 :=\sum_{\substack{d\in\mathcal C_\beta^{\rm cell}\\d\subset c}}
       \sqrt{\frac{|d|}{|c|}}\,Z_{\beta,d}.
\end{equation}
Como las subceldas son disjuntas y suman $|c|$, se obtiene, sin tomar límites,
\begin{equation}\label{eq:pdf2-ym-owner-block-intertwiner-identities}
 (I_{\alpha\beta}^{\rm curv})^*I_{\alpha\beta}^{\rm curv}=I,
 \quad
 I_{\beta\gamma}^{\rm curv}I_{\alpha\beta}^{\rm curv}
 =I_{\alpha\gamma}^{\rm curv},
 \quad
 D_\beta^{\rm YM}I_{\alpha\beta}^{\rm curv}
 =I_{\alpha\beta}^{\rm curv}D_\alpha^{\rm YM}.
\end{equation}
El relabelado celular conserva volúmenes y descendientes; para $a\in E(4)$ se tiene además
\begin{equation}\label{eq:pdf2-ym-owner-block-e4-naturality}
 U_\beta(a)I_{\alpha\beta}^{\rm curv}
 =I_{a\alpha,a\beta}^{\rm curv}U_\alpha(a).
\end{equation}

Para $h\in C_c^\infty(\mathbb R^4)$ defínase ahora el lector de bloque
\begin{equation}\label{eq:pdf2-ym-owner-smeared-witness}
 \Phi_\alpha(h):=
 \sum_{c\in\mathcal C_\alpha^{\rm cell}}
 \langle h,e_{\alpha,c}\rangle_{L^2}Z_{\alpha,c}.
\end{equation}
En una malla uniforme esta fórmula es
$\sqrt{9/2}\,a_\alpha^2\sum_c(P_\alpha h)(x_c)W_{\alpha,c}$; el promedio celular,
no una evaluación puntual, fija el producto cruzado. En efecto, si $\gamma$ refina
dos pantallas $\alpha$ y $\beta$, las definiciones anteriores dan la identidad
calculada
\begin{equation}\label{eq:pdf2-ym-owner-smeared-cross-gram}
 \left\langle
  I_{\alpha\gamma}^{\rm curv}\Phi_\alpha(h),
  I_{\beta\gamma}^{\rm curv}\Phi_\beta(k)
 \right\rangle
 =\langle P_\alpha h,P_\beta k\rangle_{L^2}.
\end{equation}
Por polarización, con $k=h$,
\begin{equation}\label{eq:pdf2-ym-owner-smeared-cauchy-derived}
 \left\|
  I_{\alpha\gamma}^{\rm curv}\Phi_\alpha(h)
  -I_{\beta\gamma}^{\rm curv}\Phi_\beta(h)
 \right\|_2^2
 =\|P_\alpha h-P_\beta h\|_{L^2}^2\longrightarrow0.
\end{equation}
Así la propiedad de Cauchy es consecuencia de la convergencia fuerte de
proyecciones, no una identidad adicional postulada. El colímite de
\eqref{eq:pdf2-ym-owner-block-intertwiner} contiene, por tanto,
\begin{equation}\label{eq:pdf2-ym-owner-smeared-gram}
 \Phi(h):=\lim_\alpha I_{\alpha,\infty}^{\rm curv}\Phi_\alpha(h),
 \qquad
 \langle\Phi(h),\Phi(k)\rangle=\langle h,k\rangle_{L^2},
 \qquad
 D_{\rm pub}^{\rm YM}\Phi(h)=3\kappa_{\rm av}\Phi(h).
\end{equation}

Esta publicación no añade grados de libertad. En cada paso, la descomposición
$V_\beta=V_\alpha\oplus(V_\beta\ominus V_\alpha)$ tiene una base de Haar
nonádica. El orden exhaustivo de incidencias nuevas de
\eqref{eq:pdf2-ym-owner-exhaustive-index} empareja esa base, etiqueta por etiqueta,
con coordenadas $Z_{\beta,d}$ del factor de innovación, conservando el soporte de la
función de Haar. Las unitarias triangulares $\widehat\Gamma_\alpha$ pegan esos
emparejamientos y producen una isometría basada
\begin{equation}\label{eq:pdf2-ym-owner-publication-into-physical}
 \mathcal U_{\rm curv}:
 \varinjlim(\mathcal K_\alpha^{\rm cell},I_{\alpha\beta}^{\rm curv})
 \longrightarrow\mathscr H_{\rm YM}^{\rm curv},
 \quad
 D_{\rm curv}^{\rm YM}\mathcal U_{\rm curv}
 =\mathcal U_{\rm curv}D_{\rm pub}^{\rm YM},
 \quad U(a)\mathcal U_{\rm curv}=\mathcal U_{\rm curv}U_{L^2}(a).
\end{equation}
En adelante $\Phi(h)$ denota también su imagen por $\mathcal U_{\rm curv}$.
La compatibilidad de \eqref{eq:pdf2-ym-gibbs-transport-global} define
$\mathcal T_{\infty,g}$ en el límite. La publicación física al acoplamiento fijado
no es una segunda isometría escogida sólo sobre el primer caos, sino la restricción
\begin{equation}\label{eq:pdf2-ym-full-publication-restriction}
 \mathcal U_{{\rm curv},g}
 :=\mathcal T_{\infty,g}^{-1}\mathcal U_{{\rm curv},0},
 \qquad
 \widetilde{\mathcal U}_g
 :=\mathcal T_{\infty,g}^{-1}:
 L^\infty(\widehat X_\infty,\widehat\mu_\infty^0)
 \longrightarrow
 L^\infty(\widehat X_\infty,\widehat\mu_{\infty,g}^{\rm YM}).
\end{equation}
Sobre la $*$-álgebra polinómica generada por holonomías, plaquetas, clover y
$\Phi(h)$ se verifican, generador a generador y después por linealidad,
\begin{equation}\label{eq:pdf2-ym-full-publication-products}
 \widetilde{\mathcal U}_g(AB)
 =\widetilde{\mathcal U}_g(A)\widetilde{\mathcal U}_g(B),
 \quad
 \widetilde{\mathcal U}_g(A^*)=\widetilde{\mathcal U}_g(A)^*,
 \quad
 \omega_g(\widetilde{\mathcal U}_gA)=\omega_0(A).
\end{equation}
Las igualdades de superficie, clover y Bianchi son preservadas por
\eqref{eq:pdf2-ym-transport-surface-bianchi}; por densidad monótona, la extensión
vale en toda el álgebra local acotada. Así
\eqref{eq:pdf2-ym-owner-smeared-cross-gram} es la restricción de un isomorfismo
probabilístico y algebraico completo, no un sustituto de él.

\subsubsection{Núcleo denso, continuidad fuerte y red local}

Sea $M_{\Phi(h)}$ el operador de multiplicación afiliado al álgebra local y
\begin{equation}\label{eq:pdf2-ym-owner-curvature-weyl}
 \mathsf W(h,t):=\exp\!\bigl(itM_{\Phi(h)}\bigr),
 \qquad h\in C_c^\infty(\mathbb R^4),\quad t\in\mathbb R.
\end{equation}
Se completa $\mathfrak A_{\rm YM}(O)$ con estos unitarios para
$\operatorname{supp}h\subset O$ y con las funciones de clase acotadas de las
holonomías de plaqueta. Denótese por $\mathfrak A_{\rm YM}(O)^{\rm sm}$ la
$*$-álgebra así generada con funciones de test suaves. El núcleo
\begin{equation}\label{eq:pdf2-ym-owner-smooth-core}
 \mathscr D_{\rm sm}:=\operatorname{span}\left\{
 A_1\cdots A_n\Omega_\infty:
 A_j\in\mathfrak A_{\rm YM}(O_j)^{\rm sm},\quad
 O_j\Subset\mathbb R^4\right\}
\end{equation}
es denso por la definición cíclica de
\eqref{eq:pdf2-ym-owner-curvature-cyclic-sector}. Las ecuaciones
\eqref{eq:pdf2-ym-owner-block-e4-naturality} y
\eqref{eq:pdf2-ym-owner-publication-into-physical} dan, para $a\in E(4)$,
\begin{equation}\label{eq:pdf2-ym-owner-field-e4-covariance}
 U(a)\Phi(h)=\Phi(a\!\cdot h),
 \qquad
 \|(U(a)-I)\Phi(h)\|=\|a\!\cdot h-h\|_{L^2}.
\end{equation}
Además, $|e^{ix}-e^{iy}|\le|x-y|$ y
\eqref{eq:pdf2-ym-owner-smeared-gram} implican
\begin{equation}\label{eq:pdf2-ym-owner-weyl-continuity}
 \|\bigl(\mathsf W(a\!\cdot h,t)-\mathsf W(h,t)\bigr)\Omega_\infty\|
 \le |t|\,\|a\!\cdot h-h\|_{L^2}.
\end{equation}
La continuidad se comprueba también para los generadores de holonomía del core;
\eqref{eq:pdf2-ym-owner-weyl-continuity} por sí sola no los incluye. Para una ruta poligonal orientada $\gamma$,
una plaqueta $p$ y $f\in C^1(G)^{\rm class}$ se ponen
\begin{equation}\label{eq:pdf2-ym-owner-holonomy-core-generators}
 A_{\gamma,f}:=f(\operatorname{Hol}_\gamma),
 \qquad A_{p,f}:=f(\operatorname{Hol}_p).
\end{equation}
En un overlay de $\gamma$ y $a\gamma$, con $a\in E(4)$, la ley de puente
\eqref{eq:pdf2-ym-owner-heat-bridge} cancela los incrementos comunes. Para cada
incremento restante $z$ se usa
$|f(xz)-f(x)|\le\|\nabla f\|_\infty d_G(z,e)$ y la cota de segundo momento del
kernel de calor
\(
 \int_Gd_G(z,e)^2k_\tau(z)\,dz\le C_G\tau
\).
Para la ley reponderada, $0<e^{-p/(4g^2)}\le1$ y la continuidad de $p$ en el
factor compacto dan
\begin{equation}\label{eq:pdf2-ym-owner-gibbs-heat-moment}
 \frac{\int_Gd_G(z,e)^2e^{-p(z)/(4g^2)}k_\tau(z)\,dz}
      {\int_Ge^{-p(z)/(4g^2)}k_\tau(z)\,dz}
 \le C_{G,g}\tau,
\end{equation}
uniformemente en las pantallas de un compacto y para $0<\tau\le\tau_0$.
Sumando los tiempos de los incrementos de la diferencia simétrica se obtiene la
estimación calculada
\begin{equation}\label{eq:pdf2-ym-owner-holonomy-e4-estimate}
 \begin{aligned}
 \|U(a)A_{\gamma,f}\Omega_\infty-A_{\gamma,f}\Omega_\infty\|_2^2
 &\le C_{G,g}\|\nabla f\|_\infty^2
     \tau(\gamma\triangle a\gamma),\\
 \|U(a)A_{p,f}\Omega_\infty-A_{p,f}\Omega_\infty\|_2^2
 &\le C_{G,g}\|\nabla f\|_\infty^2
     \tau(\partial p\triangle\partial(ap)).
 \end{aligned}
\end{equation}
Aquí la diferencia simétrica no se mide como dos curvas disjuntas: los conectores
de superficie la llenan por la cinta orientada mínima
$\Sigma(\gamma,a\gamma)$ y
\begin{equation}\label{eq:pdf2-ym-owner-bridge-time-ribbon}
 \tau(\gamma\triangle a\gamma)
 :=\sum_{p\subset\Sigma(\gamma,a\gamma)}\tau_p,
 \qquad
 \tau(\gamma\triangle a\gamma)
 \le C_\gamma d_{E(4)}(a,e),
\end{equation}
con la misma definición para las fronteras de plaqueta. La desigualdad sale de
sumar los tiempos aditivos de \eqref{eq:pdf2-ym-owner-heat-bridge} sobre una cinta
de anchura $d_{E(4)}(a,e)$ y longitud acotada por la de $\gamma$.
La medida de tiempo de puente del miembro derecho tiende así a cero cuando $a\to e$;
esto es exactamente la continuidad de los conectores de ruta de la imagen gauge.
Cada lector clover acotado $\chi(F_{\rm cl})$, $\chi\in C_b^1$, es función de una
suma finita de plaquetas transportadas; la misma cota, multiplicada por
$\|\chi'\|_\infty$ y por el número de términos, prueba también su continuidad.

Para funciones de clase meramente acotadas se hace visible la regularización que
entra en el core. Si $d b$ es Haar izquierdo en $E(4)$ y
$\eta\in C_c^\infty(E(4))$, se define el integral de Bochner
\begin{equation}\label{eq:pdf2-ym-owner-orbit-smearing}
 A[\eta]:=\int_{E(4)}\eta(b)U(b)AU(b)^*\,db,
 \qquad A\in\{A_{\gamma,f},A_{p,f},\chi(F_{{\rm cl}})\},
 \quad \chi\in C_b^1.
\end{equation}
Con $(L_a\eta)(b):=\eta(a^{-1}b)$, el cambio de variable $b\mapsto ab$ da
\begin{equation}\label{eq:pdf2-ym-owner-orbit-smearing-continuity}
 U(a)A[\eta]U(a)^*=A[L_a\eta],
 \qquad
 \|(A[L_a\eta]-A[\eta])\Omega_\infty\|
 \le\|A\|\,\|L_a\eta-\eta\|_{L^1(E(4))}\longrightarrow0.
\end{equation}
Se entiende desde ahora que los generadores de holonomía, plaqueta y clover de
$\mathfrak A_{\rm YM}(O)^{\rm sm}$ son los de
\eqref{eq:pdf2-ym-owner-holonomy-core-generators} con $f\in C^1$, o sus
regularizaciones \eqref{eq:pdf2-ym-owner-orbit-smearing}, y que el soporte barrido
permanece compacto dentro de $O$. Las aproximaciones de identidad en $E(4)$,
junto con \eqref{eq:pdf2-ym-owner-holonomy-e4-estimate}, recuperan los generadores
no regularizados en $L^2$; por ello no se reduce el sector cíclico.

Un telescopio de productos que usa
\eqref{eq:pdf2-ym-owner-weyl-continuity} y
\eqref{eq:pdf2-ym-owner-orbit-smearing-continuity} prueba la continuidad de todos
los generadores de $\mathscr D_{\rm sm}$; por
densidad y unitariedad,
\begin{equation}\label{eq:pdf2-ym-owner-strong-e4}
 \lim_{a\to e}\|(U(a)-I)\Psi\|=0
 \qquad(\Psi\in\mathscr H_{\rm YM}^{\rm curv}).
\end{equation}
Los soportes disjuntos usan factores disjuntos y conmutan. Queda así la red
\begin{equation}\label{eq:pdf2-ym-owner-local-net}
\begin{aligned}
 O_1\subseteq O_2&\Rightarrow
 \mathfrak A_{\rm YM}(O_1)\subseteq\mathfrak A_{\rm YM}(O_2),\\
 [\mathfrak A_{\rm YM}(O_1),\mathfrak A_{\rm YM}(O_2)]&=0
 \quad (O_1\cap O_2=\varnothing),\\
 U(a)\mathfrak A_{\rm YM}(O)U(a)^*&=\mathfrak A_{\rm YM}(aO).
\end{aligned}
\end{equation}
La identidad de Gram para $h\ne0$ y
\eqref{eq:pdf2-ym-owner-wc-curvature-functional} prueban que la red no es escalar.

\subsubsection{Axiomas de Osterwalder--Schrader y reconstrucción continua}

Para una dirección $\nu$, sea $\mathfrak A_{+,\nu}$ la clausura de la unión de
las álgebras anteriores con soporte compacto en $\{x_\nu>0\}$. El estado
estático, la ley de pilas y el operador de transferencia no se definen por
separado. En una pantalla $\alpha$, sea
$\widehat\mu_{\alpha,g}^{\rm YM}$ la marginal de
\eqref{eq:pdf2-ym-gibbs-measure}, sea
$\nu_{\alpha,g}:=(\pi_\alpha^{\rm phys})_*
\widehat\mu_{\alpha,g}^{\rm YM}$ su ley en el cociente y sea
$K_{\alpha,g}(t)=e^{-tD_{\alpha,g}^{\rm YM}}$ su kernel.
Para tiempos ordenados $t_1\le\cdots\le t_n$, la marginal euclídea de la
misma medida es
\begin{equation}\label{eq:pdf2-ym-owner-euclidean-time-marginal}
 \mathbb P_{\alpha,g}^{E,(\nu;t_1,\ldots,t_n)}
 :=(\operatorname{ev}_{t_1}^{(\nu)},\ldots,
    \operatorname{ev}_{t_n}^{(\nu)})_*
    \widehat\mu_{\alpha,g}^{\rm YM}.
\end{equation}
La identidad de desintegración que forma parte de
\eqref{eq:pdf2-g-gau-transfer-euclidean-dato} la calcula, no la reemplaza por
una cadena auxiliar:
\begin{equation}\label{eq:pdf2-ym-owner-euclidean-markov-identification}
 d\mathbb P_{\alpha,g}^{E,(\nu;t_1,\ldots,t_n)}
 =\nu_{\alpha,g}(dy_1)
   \prod_{j=1}^{n-1}K_{\alpha,g}
      (t_{j+1}-t_j;y_j,dy_{j+1}).
\end{equation}
Si $M_A$ denota multiplicación por un lector acotado $A$, sus correlaciones
son por tanto
\begin{equation}\label{eq:pdf2-ym-owner-schwinger-transfer}
 \begin{aligned}
 S_{\alpha,n}^{(\nu,g)}(A_1,t_1;\ldots;A_n,t_n)
 &:=\int\prod_{j=1}^n(\tau_{t_je_\nu}A_j)(x)\,
       d\widehat\mu_{\alpha,g}^{\rm YM}(x)\\
 &=\langle\mathbf1,
 M_{A_1}K_{\alpha,g}(t_2-t_1)M_{A_2}\cdots
 K_{\alpha,g}(t_n-t_{n-1})M_{A_n}\mathbf1
 \rangle_{L^2(\nu_{\alpha,g})}.
 \end{aligned}
\end{equation}
La ley de la pila estacionaria en esos mismos tiempos es esa marginal
euclídea, escrita mediante su desintegración:
\begin{equation}\label{eq:pdf2-ym-owner-markov-stack-law}
 d\mathbb P_{\alpha,g}^{(\nu;t_1,\ldots,t_n)}
 :=d\mathbb P_{\alpha,g}^{E,(\nu;t_1,\ldots,t_n)}
 =\nu_{\alpha,g}(dy_1)
   \prod_{j=1}^{n-1}K_{\alpha,g}(t_{j+1}-t_j;y_j,dy_{j+1}).
\end{equation}
Integrar sucesivamente $y_n,y_{n-1},\ldots,y_1$ da la identidad exacta
\begin{equation}\label{eq:pdf2-ym-owner-stack-transfer-identity}
 S_{\alpha,n}^{(\nu,g)}(A_1,t_1;\ldots;A_n,t_n)
 =\int\prod_{j=1}^nA_j(y_j)\,
   d\mathbb P_{\alpha,g}^{(\nu;t_1,\ldots,t_n)}(y).
\end{equation}
Si todos los tiempos coinciden, los kernels son la identidad y
\begin{equation}\label{eq:pdf2-ym-owner-static-is-stack}
 S_{\alpha,n}^{(\nu,g)}(A_1,t;\ldots;A_n,t)
 =\int A_1\cdots A_n\,d\nu_{\alpha,g}.
\end{equation}
Así el estado Schwinger estático es la marginal de tiempo igual de la misma
medida euclídea, y su desintegración es la ley Markov; no son dos funcionales
que luego se identifican por nombre. La
compatibilidad de $\widehat J_{\alpha,g}$ con medida y kernel hace compatible
\eqref{eq:pdf2-ym-owner-stack-transfer-identity} con todo overlay. Su límite define
\begin{equation}\label{eq:pdf2-ym-owner-schwinger-functionals}
 S_n^{\rm YM}(A_1,t_1;\ldots;A_n,t_n)
 :=\lim_{\alpha}S_{\alpha,n}^{(\nu,g)}
 (A_{1,\alpha},t_1;\ldots;A_{n,\alpha},t_n).
\end{equation}

Las inserciones de $\Phi$ se obtienen derivando en $t=0$ los unitarios
\eqref{eq:pdf2-ym-owner-curvature-weyl}. Puesto que los $Z_{\alpha,c}$ transportados
son centrados, independientes y uniformemente acotados, el lema de Hoeffding da,
nivel a nivel,
\begin{equation}\label{eq:pdf2-ym-owner-subgaussian-bound}
 \left\langle\Omega_\alpha,
 e^{t\Phi_\alpha(h)}\Omega_\alpha\right\rangle
 \le \exp\!\left(\frac9{16}t^2\|P_\alpha h\|_{L^2}^2\right).
\end{equation}
La cota pasa al límite por
\eqref{eq:pdf2-ym-owner-smeared-cauchy-derived}; todas las derivadas son continuas en
la topología de Schwartz y definen distribuciones temperadas.

La forma de ocho caras se prolonga primero a toda pila finita. Si $F$ depende de
las $N$ capas positivas de una pantalla $\alpha$, la reflexión de
\eqref{eq:pdf2-ym-owner-markov-stack-law} y la condición en la capa central dan
\begin{equation}\label{eq:pdf2-ym-owner-finite-stack-os}
 \begin{aligned}
 &\int \overline{F(\Theta_\nu y_{-N},\ldots,
                         \Theta_\nu y_{-1})}
          F(y_1,\ldots,y_N)\,d\Omega_{\alpha,N}^{(\nu,g)}(y)\\
 &\qquad=
 \int\nu_{\alpha,g}(dz)\left|
  \int K_{\alpha,g}(a_\alpha/2;z,dy_1)
       \prod_{j=1}^{N-1}K_{\alpha,g}(a_\alpha;y_j,dy_{j+1})
       F(y_1,\ldots,y_N)
 \right|^2\ge0.
 \end{aligned}
\end{equation}
Todo cilindro de semiespacio pertenece a una pila finita; por ello
\eqref{eq:pdf2-ym-owner-finite-stack-os} materializa la forma OS antes de la
clausura. En el cociente OS, si $F$ tiene inserciones en tiempos positivos, sea
$(\tau_sF)(A_1,t_1;\ldots;A_n,t_n):=
F(A_1,t_1+s;\ldots;A_n,t_n+s)$. Se define el operador de traslado por
\begin{equation}\label{eq:pdf2-ym-owner-os-transfer-operator}
 T_{\rm OS}^{(\nu)}(s)[F]:=[\tau_sF],
 \qquad s\ge0.
\end{equation}
Para cilindros $F,G$ de semiespacio, la igualdad euclídea--Markov
\eqref{eq:pdf2-ym-owner-euclidean-markov-identification} y la propiedad de
semigrupo dan generador a generador
\begin{equation}\label{eq:pdf2-ym-owner-os-markov-intertwining}
 \begin{aligned}
 \langle[F],T_{\rm OS}^{(\nu)}(s)[G]\rangle_{\rm OS}
 &=S_g^{\rm YM}((\Theta_\nu F)^*\tau_sG)\\
 &=\left\langle
    \mathcal V_{\infty,\nu}[F],
    e^{-sD_{\rm curv,g}^{\rm YM}}
    \mathcal V_{\infty,\nu}[G]
   \right\rangle .
 \end{aligned}
\end{equation}
La identidad muestra a la vez que las clases nulas permanecen nulas y que,
por densidad de los cilindros,
\begin{equation}\label{eq:pdf2-ym-owner-os-generator-identification}
 \mathcal V_{\infty,\nu}T_{\rm OS}^{(\nu)}(s)
 \mathcal V_{\infty,\nu}^{-1}
 =e^{-sD_{\rm curv,g}^{\rm YM}}.
\end{equation}
Ésta es la flecha explícita medida euclídea $\to$ desintegración $\to$
transferencia OS $\to$ Hamiltoniano.

\begin{proposition}[Paquete Osterwalder--Schrader continuo]
\label{prop:pdf2-ym-owner-continuum-os-package}
Los funcionales \eqref{eq:pdf2-ym-owner-schwinger-functionals} satisfacen:
\begin{enumerate}
 \item normalización, hermiticidad y simetría bosónica;
 \item covariancia euclídea y regularidad temperada;
 \item positividad por reflexión en cada una de las cuatro direcciones,
 \begin{equation}\label{eq:pdf2-ym-owner-continuum-reflection-positive}
  \sum_{i,j}\overline{a_i}a_j\,
  S^{\rm YM}\!\left((\Theta_\nu F_i)^*F_j\right)\ge0,
  \qquad F_i\in\mathfrak A_{+,\nu};
 \end{equation}
 \item localidad y cluster exponencial: si $A,B$ son locales centrados y una
 traslación de longitud $r$ separa sus soportes, entonces
 \begin{equation}\label{eq:pdf2-ym-owner-continuum-cluster}
  |S^{\rm YM}(A\,\tau_rB)|
  \le e^{-3\kappa_{\rm av}r}\,
       \|A\Omega_\infty\|\,\|B\Omega_\infty\|.
 \end{equation}
\end{enumerate}
La reconstrucción OS de cualquiera de las cuatro semirredes produce el mismo
espacio cíclico, el mismo vacío y el Hamiltoniano
$D_{\rm curv,g}^{\rm YM}$. La continuación OS produce campos de Wightman locales y una
representación de Poincaré cuyo espectro satisface
\begin{equation}\label{eq:pdf2-ym-owner-wightman-spectrum}
 \operatorname{Spec}(P^0,\mathbf P)\subseteq\overline V_+,
 \qquad P^0=D_{\rm curv,g}^{\rm YM}\ge0.
\end{equation}
\end{proposition}

\begin{proof}
Normalización y hermiticidad proceden del estado; simetría y localidad, de la
conmutatividad euclídea a soportes disjuntos. La covariancia y continuidad son
\eqref{eq:pdf2-ym-owner-field-e4-covariance}--
\eqref{eq:pdf2-ym-owner-strong-e4}; la regularidad es
\eqref{eq:pdf2-ym-owner-subgaussian-bound}. Para
$F_i$ cilíndricos, todos viven en un overlay y una pila finitos;
\eqref{eq:pdf2-ym-owner-finite-stack-os} y la polarización prueban
\eqref{eq:pdf2-ym-owner-continuum-reflection-positive}. La aproximación en norma OS
y Cauchy--Schwarz la extienden a $\mathfrak A_{+,\nu}$.

Tras una rotación euclídea, la separación espacial se convierte en tiempo positivo.
La identidad pila--transferencia
\eqref{eq:pdf2-ym-owner-stack-transfer-identity} y el operador
\eqref{eq:pdf2-ym-owner-os-transfer-operator} escriben la correlación centrada como
\[
 \langle A^*\Omega_\infty,
 e^{-rD_{\rm curv,g}^{\rm YM}}B\Omega_\infty\rangle.
\]
La cota espectral \eqref{eq:pdf2-ym-owner-curvature-sector-gap} da
\eqref{eq:pdf2-ym-owner-continuum-cluster}. Finalmente, los conectores
\eqref{eq:pdf2-ym-owner-os-connectors} forman un sistema dirigido; su unión es densa
por \eqref{eq:pdf2-ym-owner-smooth-core}. La identidad
\eqref{eq:pdf2-ym-owner-common-os-hamiltonian} identifica las cuatro clausuras y sus
generadores con $D_{\rm curv,g}^{\rm YM}$. Las hipótesis ya verificadas de regularidad,
covariancia, reflexión y cluster son las entradas de la continuación OS; su teorema de
reconstrucción da localidad de Wightman y
\eqref{eq:pdf2-ym-owner-wightman-spectrum}.
\end{proof}

\subsubsection{Reconstrucción cuatridimensional de la teoría de Yang--Mills}

La identificación no se apoya sólo en evaluar un clover sobre una conexión dada.
Sea $\mathcal R_G^{\rm loc}$ la $*$-álgebra de lectores formada por todas las
funciones de clase de las holonomías de incidencia y plaqueta, sus transportes a un
vértice base y las polarizaciones del lector $\mathcal P^{F^2}$. Sea
\begin{equation}\label{eq:pdf2-ym-owner-curvature-fibre}
 \mathcal E_G:=\Lambda^2(\mathbb R^4)^*\otimes\mathfrak g
\end{equation}
con el producto invariante. En cada celda, la polarización de
$\mathcal P^{F^2}$ da un Gram positivo $G_{\alpha,c}^{F}$ sobre
$\mathcal E_G$. Se divide por su kernel y se aplica su raíz positiva; el lector
normalizado resultante se denota $Y_{\alpha,c}(v)$ y verifica
\begin{equation}\label{eq:pdf2-ym-owner-curvature-cell-gram}
 \langle Y_{\alpha,c}(v),Y_{\alpha,d}(w)\rangle
 =\delta_{cd}\langle v,w\rangle_{\mathcal E_G}.
\end{equation}
Si $d\subset c$, sea $T_{d\leftarrow c}$ el transporte de marco ya presente en la
holonomía superficial. El entrelazador de curvatura completa es
\begin{equation}\label{eq:pdf2-ym-owner-curvature-intertwiner}
 I_{\alpha\beta}^{F}Y_{\alpha,c}(v)
 :=\sum_{d\subset c}\sqrt{\frac{|d|}{|c|}}\,
 Y_{\beta,d}\!\left(\operatorname{Ad}_{T_{d\leftarrow c}}v\right).
\end{equation}
La invariancia del producto de $\mathfrak g$, la disjunción de subceldas y la
composición de transportes prueban
\begin{equation}\label{eq:pdf2-ym-owner-curvature-intertwiner-identities}
 (I_{\alpha\beta}^{F})^*I_{\alpha\beta}^{F}=I,
 \qquad
 I_{\beta\gamma}^{F}I_{\alpha\beta}^{F}=I_{\alpha\gamma}^{F}.
\end{equation}
Para $h\in C_c^\infty(\mathbb R^4;\mathcal E_G)$ se pone
\begin{equation}\label{eq:pdf2-ym-owner-f-alpha}
 \mathbf F_\alpha(h):=
 \sum_{c\in\mathcal C_\alpha^{\rm cell}}
 Y_{\alpha,c}\!\left(
   |c|^{-1/2}\int_c h(x)\,d^4x\right).
\end{equation}
En un overlay común, el mismo cálculo de
\eqref{eq:pdf2-ym-owner-smeared-cross-gram} da ahora
\begin{equation}\label{eq:pdf2-ym-owner-curvature-cross-gram}
 \left\langle I_{\alpha\gamma}^{F}\mathbf F_\alpha(h),
 I_{\beta\gamma}^{F}\mathbf F_\beta(k)\right\rangle
 =\langle P_\alpha h,P_\beta k\rangle_{L^2(\mathcal E_G)}.
\end{equation}
Por tanto $\mathbf F_\alpha(h)$ es Cauchy por convergencia fuerte de las proyecciones,
y no por una hipótesis nueva. El mismo emparejamiento de Haar e innovaciones de
\eqref{eq:pdf2-ym-owner-publication-into-physical}, ahora aplicado a la base
ortonormalizada de $G_{\alpha,c}^{F}$, realiza el colímite mediante una isometría
local y gauge-covariante en
$L^2(\widehat X_\infty,\widehat\mu_{\infty,g}^{\rm YM})$. Su límite define la
distribución aleatoria
gauge-covariante
\begin{equation}\label{eq:pdf2-ym-owner-distributional-curvature}
\begin{aligned}
 \mathbf F_{\mu\nu}&\in
 \mathcal S'(\mathbb R^4;\mathfrak g)\ \widehat\otimes\
 L^2(\widehat X_\infty,\widehat\mu_{\infty,g}^{\rm YM}),\\
 U(a)\mathbf F(h)U(a)^*&=\mathbf F(a\!\cdot h),\\
 u\mathbf F(h)u^{-1}&=\mathbf F(\operatorname{Ad}_u h),
 \qquad a\in E(4),\ u\in G.
\end{aligned}
\end{equation}
Su canal de clase espectral de orden tres es precisamente $\Phi$. La relación
superficial finita
\begin{equation}\label{eq:pdf2-ym-owner-surface-product}
 \operatorname{Hol}^{(\alpha)}_p
 =\prod_{p'\subset p}^{\longrightarrow}
  T_{p'}\operatorname{Hol}^{(\beta)}_{p'}T_{p'}^{-1}
\end{equation}
pasa a la clausura cofinal porque todos sus productos cruzados se calculan en un
overlay. En particular, el lector clover aleatorio, y no sólo su evaluación en una
conexión suave, satisface la identidad de pequeños lazos
\begin{equation}\label{eq:pdf2-ym-owner-holonomy-curvature-relation}
 \mathbf F(h)=L^2\!\mathord{-}\!\lim_{\alpha}
 a_\alpha^4\sum_x
 \left\langle(P_\alpha h)(x),F_{{\rm cl},\alpha}(x)\right\rangle,
 \qquad
 \delta_{\rm gauge}S_n^{\rm YM}=0,
 \qquad
 \operatorname{Alt}_{\lambda\mu\nu}D_\lambda\mathbf F_{\mu\nu}=0.
\end{equation}
La primera igualdad es \eqref{eq:pdf2-ym-owner-curvature-cross-gram} escrita en el
representante clover; la segunda es la identidad de Ward
\eqref{eq:pdf2-ym-ward-finite} pasada al límite, y la
tercera es la cancelación de la frontera de una frontera en el producto superficial.
Ninguna usa el término de masa.

Para evitar definir el blanco por la conclusión que se quiere obtener, sea
$\mathbf{YM}_4(G,g)$ la clase de quíntuplas locales
$(\mathfrak A,\mu_g,K_g,\operatorname{Hol},\mathbf F)$ que verifican las siguientes
ecuaciones comprobables: (i) la densidad de Gibbs es
\eqref{eq:pdf2-ym-gibbs-measure}; (ii) $K_g$ es Markov, reversible y satisface
\eqref{eq:pdf2-ym-schwinger-dyson-finite}--\eqref{eq:pdf2-ym-ward-finite};
(iii) holonomía, clover y curvatura satisfacen
\eqref{eq:pdf2-ym-transport-surface-bianchi} y
\eqref{eq:pdf2-ym-owner-holonomy-curvature-relation}; y (iv) sus funcionales son la
misma ley de pilas de \eqref{eq:pdf2-ym-owner-stack-transfer-identity} y cumplen el
paquete OS\@. El objeto construido es, componente a componente,
\begin{equation}\label{eq:pdf2-ym-owner-qym-object}
\begin{aligned}
 \mathfrak Q_{\rm YM}^{(4)}(G,g):={}&
 \bigl(\{\mathfrak A_{\rm YM}(O)\}_O,
 S_g^{\rm YM},U,\Theta,\Omega_{\infty,g},\\
 &\widehat\mu_{\infty,g}^{\rm YM},
 \{S_{m,g},\mathscr S_{m,g}\}_m,K_g,
 \operatorname{Hol},\mathbf F,\\
 &\mathscr H_{\rm OS},D_{\rm curv,g}^{\rm YM}\bigr)
 \in\mathbf{YM}_4(G,g).
\end{aligned}
\end{equation}
El mapa de publicación no omite ninguna de esas componentes:
\begin{equation}\label{eq:pdf2-ym-owner-typed-publication-map}
 \begin{aligned}
 \mathcal P_{{\rm YM},g}:\mathfrak G_{\rm gau}(G)&\longrightarrow
 \mathbf{YM}_4(G,g),\\
 (\widehat\mu^0,\mathbb P^{\rm phys},D^0,\Theta,W,
   \operatorname{Hol},\mathcal P^{F^2})
 &\longmapsto
 (\widehat\mu_g^{\rm YM},S_g,\mathscr S_g,K_g,\mathbb P_g^{\rm phys},
  S_g^{\rm YM},\mathfrak A_{\rm YM},D_{\rm curv,g}^{\rm YM},
  \mathbf F,\operatorname{Hol}).
 \end{aligned}
\end{equation}
Los $S_n^{\rm YM}$ son así exactamente las funciones de Schwinger de la medida
dependiente de $g$, la ley Markov y el operador de transferencia ya construidos.

Queda por comprobar que $S_{m,g}$ es la acción de Yang--Mills y no sólo una
función con ese nombre. Sea
$A\in C_c^3(\mathbb R^4;T^*\mathbb R^4\otimes\mathfrak g)$ una conexión suave en
una carta donde el logaritmo de los lazos pequeños es único. Para la plaqueta
$p_{\mu\nu}(x)$, la expansión ordenada de Magnus da
\begin{equation}\label{eq:pdf2-ym-owner-plaquette-magnus}
 \log\operatorname{Hol}_{p_{\mu\nu}(x)}(A)
 =a_m^2F_{A,\mu\nu}(x)
  +\frac{a_m^3}{2}(D_\mu+D_\nu)F_{A,\mu\nu}(x)
  +R_{m,\mu\nu}(x),
\end{equation}
con
$\|R_{m,\mu\nu}(x)\|\le C_G a_m^4
(\|D^2F_A\|_{\infty,p}+\|F_A\|_{\infty,p}^2)$.
El promedio de las cuatro orientaciones cancela el término impar y define
\begin{equation}\label{eq:pdf2-ym-owner-clover}
 F_{{\rm cl},m}(x)
 :=\frac1{4a_m^2}\sum_{p\ni x}
  \operatorname{Ad}_{T_p}\log\operatorname{Hol}_p,
 \qquad
 \|F_{{\rm cl},m}(x)-F_A(x)\|
 \le C_G a_m^2(\|D^2F_A\|_\infty+\|F_A\|_\infty^2).
\end{equation}
La definición de la raíz
\eqref{eq:pdf2-ym-action-root} es precisamente la polarización de esos seis
componentes. Sustituyendo \eqref{eq:pdf2-ym-owner-plaquette-magnus} en cada
generador de esa polarización se calcula
\begin{equation}\label{eq:pdf2-ym-owner-action-root-smooth-chart}
 \xi_{m,b}(A)
 =a_m^2\operatorname{vec}_{\mu<\nu}F_{A,\mu\nu}(x_b)
  +r_{m,b}(A),
 \qquad
 \|r_{m,b}(A)\|\le C_{G,A}a_m^4.
\end{equation}
Los términos cruzados impares se cancelan por las cuatro orientaciones del clover;
los pares están incluidos una sola vez por
\eqref{eq:pdf2-ym-action-factor-partition}. Por suma de Riemann y
Cauchy--Schwarz,
\begin{equation}\label{eq:pdf2-ym-owner-action-error}
 \left|S_{m,g}(A)-\frac1{4g^2}
  a_m^4\sum_x\|F_{{\rm cl},m}(x)\|^2\right|
 \le C_{G,A,g}a_m^2,
\end{equation}
y, usando la cota de \eqref{eq:pdf2-ym-owner-clover},
\begin{equation}\label{eq:pdf2-ym-owner-f2}
 S_{m,g}(A)\longrightarrow
 S_{{\rm YM},g}(A):=\frac1{4g^2}
 \int_{\mathbb R^4}\|F_A(x)\|^2\,d^4x.
\end{equation}
Las ecuaciones de Schwinger--Dyson y Ward ya se calcularon sobre la medida que usa
esta acción. La brecha, en cambio, se deduce después por la conjugación unitaria de
\eqref{eq:pdf2-ym-g-dynamics} y el testigo
\eqref{eq:pdf2-ym-g-gap-before-limit}; no fue usada para escoger $g$, $S_{m,g}$ ni
$\widehat\mu_{m,g}^{\rm YM}$. Éste es el enlace no circular entre acción, ley y gap.

\begin{theorem}[Cierre Yang--Mills para todo grupo de Lie compacto, conectado y simple]
\label{thm:pdf2-ym-cierre-global}
La estructura completa $\mathfrak G_{\rm gau}$ construye una familia proyectiva
local para todo grupo de Lie compacto, conectado y simple. El mapa
\eqref{eq:pdf2-ym-owner-typed-publication-map} publica una teoría cuántica de
Yang--Mills no trivial sobre $\mathbb R^4$: sus funcionales satisfacen el paquete OS
continuo y su continuación satisface los axiomas locales de Wightman por la
proposición~\ref{prop:pdf2-ym-owner-continuum-os-package}; sus
holonomías y su curvatura
distribucional satisfacen
\eqref{eq:pdf2-ym-owner-holonomy-curvature-relation}, y las cuatro reconstrucciones
coinducen el mismo Hamiltoniano $D_{\infty,g}^{\rm YM}$, con vacío simple. Su ley es
$\widehat\mu_{\infty,g}^{\rm YM}$, construida por la familia compatible de acciones
locales $\{S_{m,g},\mathscr S_{m,g}\}_m$, y
satisface las ecuaciones Schwinger--Dyson y Ward
\eqref{eq:pdf2-ym-schwinger-dyson-finite}--\eqref{eq:pdf2-ym-ward-finite}. El sector
cíclico de curvatura
\eqref{eq:pdf2-ym-owner-curvature-cyclic-sector} es reductor, contiene el testigo
\eqref{eq:pdf2-ym-owner-wc-curvature-functional} y posee
\begin{equation}\label{eq:pdf2-ym-final-gap}
 \boxed{\Delta_{\rm YM}^{\rm HMT}=3\kappa_{\rm av}>0},
 \qquad
 m_{\rm gap}^{\rm HMT}
 =\frac{\hbar_{\rm int}}{c_{\rm int}}\Delta_{\rm YM}^{\rm HMT}>0.
\end{equation}
La red local, la acción fuerte de $E(4)$, las distribuciones de Schwinger, la
holonomía y $\mathbf F$ proceden del mismo estado completo. El lector exacto
$\mathcal P^{F^2}$ define la densidad de la ley; el clover es su carta suave y
\eqref{eq:pdf2-ym-owner-f2} identifica su acción clásica. La brecha pertenece al sector colectivo
gauge-invariante de curvatura y no asigna una masa elemental al gluón.
\end{theorem}

\begin{proof}
La proyectividad cinemática es \eqref{eq:pdf2-ym-owner-projectivity}; la raíz
\eqref{eq:pdf2-ym-action-root}, el cociclo
\eqref{eq:pdf2-ym-action-cocycle} y la densidad
\eqref{eq:pdf2-ym-gibbs-measure} construyen primero la acción y la medida a
acoplamiento $g$. El transporte factor a factor
\eqref{eq:pdf2-ym-gibbs-transport-factor}--
\eqref{eq:pdf2-ym-full-algebra-identities} prueba proyectividad, preservación del
estado, productos y $*$; 
\eqref{eq:pdf2-ym-transport-surface-bianchi} descarga clover y Bianchi. El índice exhaustivo
\eqref{eq:pdf2-ym-owner-exhaustive-index}, la forma
\eqref{eq:pdf2-ym-owner-full-form} y el circuito
\eqref{eq:pdf2-ym-owner-local-circuit} construyen el generador común. La reducción
física y su kernel cociente son
\eqref{eq:pdf2-ym-owner-physical-reduction}--
\eqref{eq:pdf2-ym-owner-D}. La conjugación
\eqref{eq:pdf2-ym-g-dynamics} los transporta a la medida dependiente de $g$ sin
alterar su espectro. La descomposición de Hoeffding prueba el vacío y la cota, y
\eqref{eq:pdf2-ym-owner-wc-invariance}--
\eqref{eq:pdf2-ym-owner-wc-eigenvalue} exhiben un vector físico que alcanza el
umbral. El entrelazamiento nonádico transporta forma, vacío y testigo al límite, donde
\eqref{eq:pdf2-ym-owner-limit-gap}--
\eqref{eq:pdf2-ym-owner-exact-gap} fijan la brecha exacta. Las cuatro
factorizaciones \eqref{eq:pdf2-ym-owner-four-os}, la inclusión de rangos
\eqref{eq:pdf2-ym-owner-os-range-inclusion} y los conectores
\eqref{eq:pdf2-ym-owner-os-connectors} identifican sus reconstrucciones con ese
mismo generador. La holonomía
\eqref{eq:pdf2-ym-owner-incidence-holonomy} coloca materialmente $W_c$ en el sector
de curvatura, donde \eqref{eq:pdf2-ym-owner-curvature-sector-gap} prueba la brecha
exacta. El sistema \eqref{eq:pdf2-ym-owner-block-intertwiner} controla los productos
cruzados mediante \eqref{eq:pdf2-ym-owner-smeared-cross-gram};
\eqref{eq:pdf2-ym-full-publication-restriction}--
\eqref{eq:pdf2-ym-full-publication-products} elevan ese primer caos al isomorfismo
de toda el álgebra. La convergencia fuerte de $P_\alpha$ produce $\Phi$; las cotas
\eqref{eq:pdf2-ym-owner-holonomy-e4-estimate} y
\eqref{eq:pdf2-ym-owner-orbit-smearing-continuity} prueban además la continuidad
$E(4)$ de holonomías, plaquetas y clover. La desintegración euclídea
\eqref{eq:pdf2-ym-owner-euclidean-markov-identification}, la identidad exacta
\eqref{eq:pdf2-ym-owner-stack-transfer-identity} y el entrelazamiento
\eqref{eq:pdf2-ym-owner-os-markov-intertwining} enlazan la misma medida,
estado Schwinger, pila Markov, transferencia OS y generador del gap. La
proposición~\ref{prop:pdf2-ym-owner-continuum-os-package} prolonga las
cuatro formas finitas a las
semirredes completas, prueba regularidad y cluster y reconstruye el Hamiltoniano. Por
último, \eqref{eq:pdf2-ym-owner-distributional-curvature}--
\eqref{eq:pdf2-ym-owner-typed-publication-map} identifican esos mismos funcionales
como los Schwinger de la teoría gauge local; las integraciones por partes
\eqref{eq:pdf2-ym-schwinger-dyson-finite}--\eqref{eq:pdf2-ym-ward-finite} prueban
Schwinger--Dyson y Ward. Finalmente,
\eqref{eq:pdf2-ym-owner-plaquette-magnus}--\eqref{eq:pdf2-ym-owner-f2} identifican
la acción clásica antes de aplicar la equivalencia espectral. La cadena no usa el
gap para escoger la medida.
\end{proof}

\begin{falsifier}[Colapso de la dinámica producto]
Si se reemplaza $D_m^{\rm prod}$ por $D_m^{\rm fib}$, todos los sectores no
constantes reciben el mismo autovalor. Un vector de Hoeffding con $|S|=2$ debería
tener energía $6\kappa_{\rm av}$ por \eqref{eq:pdf2-ym-hoeffding}, pero el
semigrupo simple le asignaría $3\kappa_{\rm av}$. La identificación de ambos
generadores es falsa.
\end{falsifier}

\begin{falsifier}[Confusión entre inclusión genealógica y promedio de bloque]
Si se sustituye $I_{\alpha\beta}^{\rm curv}$ por
$\widehat J_{\alpha\beta}$ en
\eqref{eq:pdf2-ym-owner-smeared-cross-gram}, una coordenada gruesa persiste con
coeficiente uno mientras el coeficiente fino cambia por el cociente de escalas. Para
$a_{m+1}=a_m/9$, el producto cruzado ancestral lleva un factor $1/81$ y no converge
a la norma diagonal. La identidad
\eqref{eq:pdf2-ym-owner-smeared-cauchy-derived} falla; por eso los dos mapas se
mantienen separados.
\end{falsifier}

\begin{falsifier}[Ablación de la fibra de incidencia]
Si se sustituye $\widehat X_m$ por su marginal de enlaces, se borra $q_c$ y con ella
el observable físico $W_c$. El mapa de olvido no refleja la subálgebra gauge completa;
por tanto no puede utilizarse para negar el autovalor del objeto fuente. Objeto
sustituido; conclusión no transferible.
\end{falsifier}

\begin{falsifier}[Sustitución del lector exacto por su carta suave]
Si se usa directamente el término asintótico de
\eqref{eq:pdf2-ym-owner-f2} como densidad en una pantalla finita, se borra el resto
controlado de \eqref{eq:pdf2-ym-owner-action-error} y se pierde el cociclo exacto
\eqref{eq:pdf2-ym-action-cocycle}. La densidad finita gobernante es
$Z_{m,g}^{-1}e^{-\mathcal P_m^{F^2}/(4g^2)}$; el clover y el límite suave la
identifican con la acción clásica, pero no la sustituyen.
\end{falsifier}

\paragraph{Articulación de las demostraciones.}
La separación entre fibra simple y generador producto, la medida común, la compresión
gauge, las cuatro condiciones de Osterwalder--Schrader, la forma límite y el testigo físico son
resultados exactos previamente establecidos en el dominio gauge
regular HMT. La desintegración con tiempos no uniformes, el índice
exhaustivo de expectativas, el kernel cociente, la matriz $B^+$, la acción explícita
en el colímite, el sistema cofinal de bloque, el paquete OS continuo, la curvatura
distribucional y la holonomía de incidencia son
formalizaciones explícitas: determinan los dominios y demuestran los enlaces antes
enunciados, manteniendo la brecha como resultado. Su composición completa
establece el teorema precedente.
La raíz positiva del lector $F^2$, su densidad Gibbs, el transporte local compatible,
las ecuaciones Schwinger--Dyson/Ward, el isomorfismo algebraico completo, las cotas
de continuidad de holonomías y la identidad pila--transferencia son
formalizaciones explícitas de las cuatro correspondencias. La identificación
entre el observable clover y la acción clásica se obtiene como consecuencia
de la expansión de Magnus.


% SOURCE sections/hamiltonianos_curvatura.tex
\section{Generadores tríticos, acciones cuadráticas y transporte de energía}
\label{sec:hamiltonianos-curvatura}

La geometría dinámica conserva la distinción entre el estado que genera una
realización y las coordenadas con las que ésta expresa su evolución. APP
construye y conserva las hojas aditiva y multiplicativa; el TRIT orienta el
modo operativo; el TPK compone selección, transporte y actualización sobre
el estado enriquecido mediante
\[
 U_t=\operatorname{Upd}_t\circ\operatorname{Tra}_t\circ
 \operatorname{Sel}_t.
\]
La holonomía nonádica prolonga este estado a través de
\(w_6\to w_{12}\to w_{18}\to w_{24}\to w_{30}\to R_{36}\to G_9\).
El retorno de fase conserva el avance de memoria. Sus publicaciones
\(\pi_{\rm HMT},\varphi_{\rm HMT},e_{\rm HMT},\alpha_{\rm HMT}\)
pertenecen a la misma generación correlacionada, con los lectores y el
orden constructivo interno ya establecidos. La estructura discreta conjunta
del continuo sostiene después las representaciones geométricas y físicas.
En este punto de la construcción, un Hamiltoniano es el funcional que
realiza un transporte ya tipado sobre una fibra simpléctica; su dominio,
su orientación y su parámetro temporal forman parte de la afirmación.

El artículo sobre la extrema y media razón holográfica recupera tres
operadores concretos. La multiplicación por \(4\) módulo \(9\) actúa sobre
las órbitas ternarias de unidades; en su módulo de aumento tiene matriz
\[
 C=\begin{pmatrix}0&-1\\1&-1\end{pmatrix},\qquad C^2+C+I=0.
\]
El transporte parabólico y la incidencia areal--volumétrica están
representados, en sus respectivos planos, por
\[
 N=\begin{pmatrix}0&1\\0&0\end{pmatrix},\qquad
 F_{\rm av}=\begin{pmatrix}0&1\\1&1\end{pmatrix}.
\]
La incidencia procede del calendario de hojas; sus potencias cuentan
composiciones de esos modos. Los generadores reales publicados son
\begin{equation}
 J_+=\frac{2C+I}{\sqrt3},\qquad J_0=N,\qquad
 J_-=\frac{2F_{\rm av}^{2}-3I}{\sqrt5},\qquad
 J_\tau^2=-\tau I.
 \label{eq:hc-generadores}
\end{equation}
Aquí \(\tau=+1,0,-1\) discrimina el régimen elíptico, parabólico e
hiperbólico. Las fibras tienen procedencias y bases propias. Las fórmulas
\(e^{\pi J_+}=-I\), \(e^{J_0}=I+J_0\) y
\(e^{2\log\varphi J_-}=F_{\rm av}^{2}\) evalúan coordenadas HMT
previamente generadas en esos operadores~\cite{hmtX}.

\subsection{Forma alternante y Hamiltoniano del generador}

En una fibra real orientada de dimensión dos se fija
\[
 \omega=dQ\wedge dP,\qquad
 \Omega=\begin{pmatrix}0&1\\-1&0\end{pmatrix},\qquad z=(Q,P)^T.
\]
Se adopta \(\iota_{X_H}\omega=dH\); por tanto
\(\dot Q=\partial_PH\), \(\dot P=-\partial_QH\).
Esta convención fija también el signo de la acción orientada.

\begin{proposition}[Realización hamiltoniana de un generador de traza nula]
\label{prop:hc-cuadratica}
Para una matriz real \(J\) de tamaño dos y traza nula, la matriz
\(G_J=-\Omega J\) es simétrica y el Hamiltoniano
\begin{equation}
 H_J(z)=\frac12z^TG_Jz
 \label{eq:hc-hj}
\end{equation}
produce exactamente \(\dot z=Jz\). Si \(J^2=-\tau I\),
entonces \(\det G_J=\tau\). Para un isomorfismo simpléctico
\(B:V\to V'\), el transporte
\[
 J'=BJB^{-1},\qquad H'(z')=H_J(B^{-1}z')
\]
satisface \(G'=B^{-T}G_JB^{-1}\) y entrelaza ambos flujos.
\end{proposition}
\begin{proof}
Escribiendo \(J=\bigl(\begin{smallmatrix}a&b\\c&-a\end{smallmatrix}\bigr)\),
se obtiene \(G_J=\bigl(\begin{smallmatrix}-c&a\\a&b\end{smallmatrix}\bigr)\).
Así \(\Omega\nabla H_J=\Omega(-\Omega J)z=Jz\).
Cayley--Hamilton da \(J^2=-\det(J)I\), y
\(\det(-\Omega)=1\). Finalmente, \(B^T\Omega B=\Omega\)
permite sustituir \(z=B^{-1}z'\) tanto en la forma como en la ecuación.
\end{proof}

Aplicada a \eqref{eq:hc-generadores}, la proposición da
\begin{equation}
 H_{J_+}=-\frac{Q^2-QP+P^2}{\sqrt3},\qquad
 H_{J_0}=\frac{P^2}{2},\qquad
 H_{J_-}=\frac{-Q^2-QP+P^2}{\sqrt5}.
 \label{eq:hc-hamiltonianos-nativos}
\end{equation}
El signo elíptico registra la orientación de \(C\) en la base de aumento.
El flujo inverso \(-J_+\) posee la forma positiva opuesta. La elipse
de acción del corpus utiliza precisamente una orientación de fase
decreciente para su Hamiltoniano positivo. Conservar este signo hace
compatible el operador cíclico con la acción \(P\,dQ\).

La conjugación también proporciona un criterio exacto de comparación.
Si \(BJ_\tau=J_{\tau'}B\) y \(B\) es invertible, elevar al cuadrado
da \((\tau-\tau')B=0\), luego \(\tau=\tau'\). Un enlace puede
conservar \(\omega\) entre fibras de regímenes distintos y conservar
sus memorias respectivas; la conjugación de los generadores exige además
igualdad del régimen. Esta diferencia permite estudiar una transición
trítica sin sustituirla por un cambio de nombre del mismo flujo.

\subsection{Deformación de la elipse y transformación de Legendre}

La coordenada \(\eta=\operatorname{artanh}(C^*/A)\), en el dominio
\(|C^*/A|<1\), define el marco simpléctico
\(M_\eta=\operatorname{diag}(e^{\eta/2},e^{-\eta/2})\).
La deformación de las dos direcciones es recíproca y conserva su producto.
En la carta normalizada \((q,p)=M_\eta^{-1}(Q,P)\), consideremos
\[
 L_\tau=\begin{pmatrix}0&1\\-\tau&0\end{pmatrix},\qquad
 I_{\tau,\eta}=\frac12\bigl(\tau e^{-\eta}Q^2+e^\eta P^2\bigr).
\]
Estos representantes clasifican flujos cuadráticos orientados. Su empleo
en una fibra nativa se hace mediante un marco simpléctico declarado,
conforme a la proposición~\ref{prop:hc-cuadratica}. El representante
elíptico positivo es el usado por la elipse de acción; para él, \(I\)
tiene dimensión de acción.

Los cambios de marco pueden escribirse explícitamente. Si
\(b_+=\sqrt{\sqrt3/2}\) y \(b_-=\sqrt{\sqrt5/2}\), las matrices
\[
 B_+=\begin{pmatrix}b_+&b_+/\sqrt3\\0&2b_+/\sqrt3\end{pmatrix},
 \qquad
 B_-=\begin{pmatrix}b_-&-b_-/\sqrt5\\0&2b_-/\sqrt5\end{pmatrix}
\]
tienen determinante uno y satisfacen
\( (-J_+)B_+=B_+L_{+1}\) y
\(J_-B_-=B_-L_{-1}\). Para el régimen parabólico se toma
\(B_0=I\). La multiplicación de las dos columnas verifica las
igualdades, y el determinante uno prueba la conservación de la forma
alternante.

\begin{theorem}[Transporte variable y lagrangiano regular]
\label{thm:hc-legendre}
Sean \(\eta(t)\) de clase \(C^1\) y \(\omega(t)>0\). El Hamiltoniano
\begin{equation}
 H_{\tau,\rm tr}
 =\omega I_{\tau,\eta}+\frac{\dot\eta}{2}QP
 \label{eq:hc-htr}
\end{equation}
conserva \(I_{\tau,\eta(t)}\) a lo largo de sus soluciones. Su
transformación de Legendre respecto de \(P\) es invertible y da
\begin{equation}
 \mathcal L_{\tau,\rm tr}(Q,\dot Q,t)
 =\frac{(\dot Q-\dot\eta Q/2)^2}{2\omega e^\eta}
 -\frac{\omega\tau e^{-\eta}}2Q^2.
 \label{eq:hc-ltr}
\end{equation}
\end{theorem}
\begin{proof}
Las ecuaciones son
\[
 \dot Q=\omega e^\eta P+\frac{\dot\eta}2Q,
 \qquad
 \dot P=-\omega\tau e^{-\eta}Q-\frac{\dot\eta}2P.
\]
Derivar \(I_{\tau,\eta(t)}\) y sustituirlas cancela por separado
los términos \(\omega QP\) y los términos de variación de \(\eta\).
La derivada segunda \(\partial_P^2H=\omega e^\eta>0\)
permite resolver
\(P=(\dot Q-\dot\eta Q/2)/(\omega e^\eta)\).
Sustituir en \(P\dot Q-H\) prueba \eqref{eq:hc-ltr}.
Además,
\[
 P\dot Q-H_{\tau,\rm tr}
 =p\dot q-\frac\omega2(\tau q^2+p^2),
\]
lo que prueba que el término mixto es la conexión variacional del marco
móvil. La conservación procede del transporte completo y no de mantener
constantes los semiejes.
\end{proof}

El resultado elíptico y su acción están desarrollados en el corpus
variacional~\cite{hmtX}; la misma prueba exhibe sus dos realizaciones
cuadráticas conjugadas. La nulidad del determinante de la forma parabólica
convive con una transformación de Legendre regular: ésta depende de la
derivada segunda en el momento, no del determinante de la forma completa.
En cambio, para el levantamiento cotangente lineal
\(H_A(q,p)=p^TAq\), se tiene \(\partial_p^2H_A=0\).
Su acción de primer orden conserva \(\dot q=Aq\) mediante un
multiplicador; es un objeto variacional distinto de
\eqref{eq:hc-ltr}.

\subsection{Retorno de memoria y carga variacional conservada}

En el registro de memoria del calendario, los contadores \((g,s)\)
reciben el incremento \((4,72)\) en el retorno completo. La relación
\(72=18\cdot4\) fija
\[
 v=(1,18)^T,\qquad I_{\rm mem}=s-18g,
 \qquad B=\begin{pmatrix}1&0\\18&1\end{pmatrix}.
\]
El estabilizador nativo es
\begin{equation}
 N_v=BNB^{-1}=
 \begin{pmatrix}-18&1\\-324&18\end{pmatrix},\qquad
 N_vm=vI_{\rm mem}(m),\qquad N_v^2=0.
 \label{eq:hc-nv}
\end{equation}
En el plano realificado \(m=(g,s)^T\), dotado de \(dg\wedge ds\),
el Hamiltoniano \(H_{\rm mem}=I_{\rm mem}^2/2\) produce
\(dm/du=N_vm\). El parámetro \(u\) pertenece a esta realización
adimensional del registro. La carta
\(q=g,\ p=s-18g\) es simpléctica y transforma la uno-forma de acción
según
\begin{equation}
 s\,dg=p\,dq+d(9q^2).
 \label{eq:hc-memoria-borde}
\end{equation}
Por tanto, la eliminación de \(s\) en la acción de primer orden da
\begin{equation}
 \mathcal L_{\rm mem}
 =\frac12\left(\frac{dg}{du}\right)^2
 +18g\frac{dg}{du}
 =\frac12\left(\frac{dg}{du}\right)^2
 +\frac{d}{du}(9g^2).
 \label{eq:hc-lmem}
\end{equation}
La ecuación variacional es \(d^2g/du^2=0\), y el momento
conservado en la carta simpléctica es \(p=I_{\rm mem}\).
El retorno afín completo procede del Hamiltoniano lineal
\(H_{\rm ret}=4I_{\rm mem}\): su aplicación al tiempo \(u=1\)
es \((g,s)\mapsto(g+4,s+72)\). Ambos Hamiltonianos conmutan
porque son funciones de la misma carga, pero realizan operaciones
distintas: estabilización parabólica y traslación del registro.
Estas fórmulas recuperan el mecanismo variacional de la memoria, incluido
su término de frontera~\cite{hmtX}.

\subsection{Curvatura métrica y energía geodésica}

La realización métrica de los regímenes utiliza una carta radial y una
escala inversa de longitud \(\rho>0\):
\begin{equation}
 g_{\tau,\rho}=dr^2+S_{\tau,\rho}(r)^2d\theta^2,
 \qquad
 S_{\tau,\rho}(r)=
 \begin{cases}
 \sin(\rho r)/\rho,&\tau=+1,\\
 r,&\tau=0,\\
 \sinh(\rho r)/\rho,&\tau=-1.
 \end{cases}
 \label{eq:hc-metrique}
\end{equation}
La curvatura gaussiana es \(K_G=-S''/S=\tau\rho^2\).
La carta se considera en su interior regular, \(S>0\); para la
realización esférica, \(0<r<\pi/\rho\). Los polos se tratan con
cartas regulares diferentes. El parámetro \(\rho\) y la realización
dimensional de la métrica mantienen su declaración propia.

Para una masa realizada \(m>0\), el lagrangiano geodésico y su
Hamiltoniano son
\begin{equation}
 \mathcal L_{\rm geo}=\frac m2(\dot r^2+S^2\dot\theta^2),\qquad
 H_{\rm geo}=\frac1{2m}\left(p_r^2+\frac{p_\theta^2}{S^2}\right).
 \label{eq:hc-geodesique}
\end{equation}
En efecto, \(p_r=m\dot r\), \(p_\theta=mS^2\dot\theta\);
el hessiano de velocidades es \(m\operatorname{diag}(1,S^2)\)
y es positivo e invertible. La ecuación radial es
\(\ddot r-SS'\dot\theta^2=0\), mientras
\(d(mS^2\dot\theta)/dt=0\) conserva el momento angular.
La positividad de esta energía vale para las tres curvaturas. La forma
indefinida del generador hiperbólico de \eqref{eq:hc-hamiltonianos-nativos}
y la energía geodésica positiva de una métrica de curvatura negativa
son funciones sobre espacios distintos. Esta comparación fija exactamente
qué significado de curvatura interviene en cada afirmación.

\subsection{Respuesta constitutiva y realización modal de Maxwell}

Las secciones de acción y las coordenadas angulares preceden a la
respuesta constitutiva \cite{hmtIntegral}. En la orientación real fijada \(x>y>0\),
con los proyectores de las dos hojas,
\[
 T=q_+P_++q_-P_-,\quad q_\pm=e^{-x\mp y},\quad
 \mathcal V=(I-T^{90})(I-T^{120})^{-1},
 \qquad 0<q_+<q_-<1,
\]
el cálculo funcional produce
\[
 r_\pm=\frac{1-q_\pm^{90}}{1-q_\pm^{120}},\qquad
 \widehat\varepsilon=r_+^2,\quad \widehat\mu=r_-^2,
 \quad \widehat c=(r_+r_-)^{-1},\quad \widehat Z=r_-/r_+.
\]
Las cartas dimensionales correspondientes conservan exactamente
\(\varepsilon\mu=c^{-2}\) y \(\mu/\varepsilon=Z^2\).
La permitividad, la permeabilidad y la velocidad son así publicaciones
del mismo operador, con los sectores de incidencia fijados por el TPK;
la dinámica siguiente las emplea como resultados ya construidos.

Para exhibir el transporte hacia un Hamiltoniano de campo, se fija una
realización espacial orientada, condiciones de frontera que permitan
integración por partes, y un modo transversal real \(f\) normalizado por
\[
 \nabla\cdot f=0,\qquad \int|f|^2=1,\qquad
 \nabla\times\nabla\times f=k^2f,\qquad k>0.
\]
En calibre transversal y sin fuentes, \(A=af\) y el momento
\(\Pi=pf\) definen una inmersión del plano modal en el espacio de
campos. Para constantes constitutivas homogéneas, el Hamiltoniano y la
forma simpléctica de Maxwell se restringen a
\begin{equation}
 H_{\rm EM}(a,p)=\frac{p^2}{2\varepsilon}
 +\frac{k^2a^2}{2\mu},\qquad \omega_{\rm EM}=da\wedge dp.
 \label{eq:hc-maxwell-mode}
\end{equation}
La integración por partes da
\(\int|\nabla\times f|^2=k^2\), que prueba esta restricción.
Con \(\omega_k=ck\), el cambio
\begin{equation}
 Q=\sqrt{\varepsilon\omega_k}\,a,
 \qquad P=\frac{p}{\sqrt{\varepsilon\omega_k}}
 \label{eq:hc-maxwell-symplectic}
\end{equation}
preserva \(da\wedge dp=dQ\wedge dP\) y da
\(H_{\rm EM}=\omega_k(Q^2+P^2)/2\). Su inversa y sus ecuaciones
\(\dot Q=\omega_kP,\ \dot P=-\omega_kQ\)
constituyen el entrelazamiento modal con el régimen elíptico positivo.
La restricción es invariante porque el operador de Maxwell conserva el
modo propio transversal. La transformación de Legendre produce
\(\mathcal L_{\rm EM}=\varepsilon\dot a^2/2-k^2a^2/(2\mu)\).

Dos modos transversales de igual frecuencia, combinados en cuadratura,
realizan las dos orientaciones circulares del corpus. El campo satisface
\(\mathbf B=c^{-1}\widehat{\mathbf k}\times\mathbf E\) y
\(\omega=ck\). La relación con el oscilador se demuestra aquí mediante
la inmersión de campos, la normalización modal y el transporte de la forma
simpléctica, manteniendo explícitas las condiciones de frontera.
Un cambio de \(k\), de la geometría de la cavidad o de la preparación
modifica el modo realizado; mantiene intacta la genealogía de
\(\varepsilon,\mu,c,Z\).

En la carta lorentziana del mismo sector, el acoplamiento a una sección
material se escribe con \(F=dA\) y una carga \(q_e\) ya tipada.
Para esta realización variacional se toma \(\nabla\) igual a la
conexión de Levi--Civita de \(g\). Sobre
\(T^*\mathcal M\), el Hamiltoniano covariante
\[
 H_A=\frac1{2m}g^{\mu\nu}(p_\mu-q_eA_\mu)(p_\nu-q_eA_\nu)
\]
es la transformación de Legendre de
\(\mathcal L_A=\frac m2g_{\mu\nu}\dot x^\mu\dot x^\nu
+q_eA_\mu\dot x^\mu\). El hessiano \(mg\) es invertible.
La ecuación de Euler--Lagrange da
\(m\nabla_u u^\mu=q_eF^\mu{}_{\nu}u^\nu\).
La antisimetría de \(F\) conserva \(g(u,u)\). Se trata de la
realización covariante con parámetro afín, mientras
\eqref{eq:hc-maxwell-mode} utiliza el tiempo de una descomposición
espacial transversal: ambos parámetros y ambos dominios quedan definidos.

\subsection{Canales graduados, representación de Clifford y operador de masa}

El atlas material produce primero canales completos de ruta, hoja, sabor,
color y representación. En una profundidad finita, su realización
hermítica es \(\mathcal H_K=\bigoplus_{s\in S_K}V_s\), con base
ortonormal de canales y graduación de paridad transportada. El álgebra
exterior del sector impar genera las CAR: la inserción de un modo
\(a_j^*\) y su retirada adjunta \(a_j\) satisfacen
\(\{a_i,a_j^*\}=\delta_{ij}I\) y \(\{a_i,a_j\}=0\).
La prueba cuenta los signos de intercambio de modos completos; las masas
se evalúan después~\cite{hmtVI}.

En el sector exterior de un solo modo, la base \((|0\rangle,|1\rangle)\)
identifica el operador de retirada con
\(a=N\). El mapa \(e_1\mapsto|0\rangle,\ e_2\mapsto|1\rangle\)
es una isometría del plano complejo normalizado y entrelaza ambos
nilpotentes. Esta identificación emplea el producto interno y la
graduación declarados; conserva el distinto significado de sus bases.
Con el \(i\) interno ya realizado, se obtienen
\begin{equation}
 \sigma_1=N+N^*,\qquad
 \sigma_2=-i(N-N^*),\qquad
 \sigma_3=NN^*-N^*N.
 \label{eq:hc-pauli}
\end{equation}
La multiplicación directa da
\(\sigma_j^*=\sigma_j\) y
\(\sigma_i\sigma_j=\delta_{ij}I+i\epsilon_{ijk}\sigma_k\).
En particular, la pareja nilpotente y su adjunto producen una
representación de Clifford; el adjunto es el de la fibra hermítica y
se transporta junto con ella.

El operador de masa sobre una familia finita de canales \(\mathcal F\)
está construido por el registro, el carácter y sus refinamientos:
\begin{equation}
 M_\beta=S M_0 S,\qquad
 S=R_{\rm act}^{B_M/2},\quad
 R_{\rm act}=\hbar_{\rm pre}/\hbar_{\rm ret}>0,
 \qquad M_0>0.
 \label{eq:hc-mass}
\end{equation}
Aquí \(B_M\) es el operador diagonal del exponente de ruta y \(M_0\)
contiene el carácter refinado relativo al origen electrónico.
La positividad sigue de \(\langle v,M_\beta v\rangle
=\langle Sv,M_0Sv\rangle>0\) para \(v\ne0\).
El producto de hojas conserva el sector nulo por su realización propia;
la exponencial del carácter material permanece positiva.

En esta realización se escribe \(\hbar:=\hbar_{\rm ret}\),
la sección de acción ya construida, y \(c:=c_{\rm vac}\) en la
carta dimensional declarada. Sobre \(\mathbb C^4\otimes\mathcal F\), definamos
\[
 \alpha_j^{D}=\sigma_1\otimes\sigma_j,\qquad
 \beta^D=\sigma_3\otimes I_2.
\]
El superíndice distingue estas matrices de la constante de estructura fina.
La carta espacial \(1+3\) asigna tres componentes de momento
\(\mathbf p\). La aplicación de multiplicación
\begin{equation}
 H_D(\mathbf p)
 =c\sum_{j=1}^3\alpha_j^Dp_j\otimes I_{\mathcal F}
 +c^2\beta^D\otimes M_\beta
 \label{eq:hc-dirac}
\end{equation}
es la realización libre de Dirac de ese operador material.

\begin{theorem}[Factorización espinorial de la dispersión material]
\label{thm:hc-dirac-square}
La matriz \eqref{eq:hc-dirac} es autoadjunta y satisface
\begin{equation}
 H_D(\mathbf p)^2
 =I_4\otimes\bigl(c^2|\mathbf p|^2I_{\mathcal F}
 +c^4M_\beta^2\bigr).
 \label{eq:hc-dirac-square}
\end{equation}
En un autoespacio \(M_\beta v=m v\), sus energías son
\(\pm E_m(\mathbf p)\), donde
\(E_m=\sqrt{c^2|\mathbf p|^2+m^2c^4}\).
\end{theorem}
\begin{proof}
Las relaciones de \eqref{eq:hc-pauli} dan
\(\{\alpha_i^D,\alpha_j^D\}=2\delta_{ij}I_4\),
\(\{\alpha_i^D,\beta^D\}=0\) y \((\beta^D)^2=I_4\).
El factor de sabor conmuta con las matrices espinoriales.
Los términos cruzados se anulan y los cuadrados dan
\eqref{eq:hc-dirac-square}. La traza nula y las relaciones de Clifford
producen ambas ramas, con multiplicidad espinorial dos cuando \(E_m>0\).
\end{proof}

En representación de momentos, el dominio natural es
\(\{\psi\in L^2(\mathbb R^3;\mathbb C^4\otimes\mathcal F):
H_D\psi\in L^2\}\). Al ser una matriz hermítica mensurable actuando
por multiplicación, el operador es autoadjunto en ese dominio. La
representación espacial procede de la transformación unitaria de Fourier
posterior a la construcción de los canales. El teorema fija la realización
libre y su dependencia exacta de la masa HMT; la interacción conserva sus
operadores de conexión y su dominio propio.

En la rama positiva \(m>0\), la transformación de Legendre ordinaria
del símbolo \(E_m(\mathbf p)\) es
\[
 \mathbf v=\frac{c^2\mathbf p}{E_m},\qquad
 \mathcal L_m=-mc^2\sqrt{1-|\mathbf v|^2/c^2},
 \qquad |\mathbf v|<c.
\]
Resolver \(\mathbf p=m\mathbf v/\sqrt{1-|\mathbf v|^2/c^2}\)
y sustituir en \(\mathbf p\cdot\mathbf v-E_m\) demuestra la
igualdad. La acción espinorial, en cambio, es de primer orden en el campo:
\(\mathcal L_D=\frac{i\hbar}{2}(\psi^*\dot\psi-\dot\psi^*\psi)
-\psi^*H_D\psi\). Su variación da
\(i\hbar\dot\psi=H_D\psi\); su hessiano en velocidades del campo
es singular. Ambas afirmaciones son compatibles porque actúan sobre
espacios variacionales diferentes.

Para una conexión abeliana estática en una carta espacial plana, con
\(\Pi_j=-i\hbar\partial_j-q_eA_j\), masa constante y potencial
escalar nulo, se conserva sobre funciones lisas de soporte compacto
\[
 [\Pi_i,\Pi_j]=iq_e\hbar\epsilon_{ijk}B_k,
\]
y la misma factorización produce
\begin{equation}
 H_D(A)^2=c^2\boldsymbol\Pi^2+c^4M_\beta^2
 -q_e\hbar c^2\boldsymbol\Sigma\cdot\mathbf B,
 \qquad \Sigma_j=I_2\otimes\sigma_j.
 \label{eq:hc-dirac-curvature}
\end{equation}
Los factores identidad de sabor quedan implícitos. El término de curvatura
es el producto de los dos antisimétricos: el conmutador de derivadas
covariantes y la multiplicación de Pauli. Su coeficiente queda fijado
al elegir la representación de carga y la carta de acción; no se
selecciona mediante un valor espectral objetivo.

Si \(V\) es la transformación unitaria de sabor ya construida, entonces
\[
 (I_4\otimes V)H_D(M_\beta)(I_4\otimes V^*)
 =H_D(VM_\beta V^*).
\]
La mezcla unitaria conserva el espectro completo. La congruencia positiva
de \eqref{eq:hc-mass} transporta otra estructura y puede modificar los
autovalores en una métrica fija. Por ejemplo,
\(M_0=\operatorname{diag}(1,2)\), \(S=\operatorname{diag}(2,1)\)
da \(SM_0S=\operatorname{diag}(4,2)\). Si se transporta también el
producto interno a \(G'=S^*S\), el problema generalizado
\(SM_0Sv=\lambda G'v\) recupera los autovalores de \(M_0\)
mediante \(w=Sv\). De este modo, cambio de base de sabor, deformación
material y transporte de métrica tienen consecuencias espectrales precisas.

\subsection{Curvatura del potencial escalar en la realización electrodébil}

El sector electrodébil utiliza el determinante del mismo operador angular,
\(D_A=\det T=e^{-2A_{\rm rad}}\), junto con las firmas materiales ya
construidas. En la realización angular con refinamientos comunes de
\(W\) y \(Z\), sus firmas \((94,-1)\) y \((95,-1)\) dan
\(m_W^\angle/m_Z^\angle=e^{-A_{\rm rad}}\).
En la convención racionalizada de árbol del corpus,
\[
 g^2=\frac{4\pi_{\rm HMT}\alpha_{\rm HMT}}{1-D_A},\qquad
 \lambda_H=\frac{\pi_{\rm HMT}\alpha_{\rm HMT}}{2(1-D_A)}
 \exp\!\left(6A_{\rm rad}+\frac{10}3C^*_{\rm rad}\right)>0.
\]
La segunda fórmula resulta de la firma escalar \((97,9)\) y de
\(m_H^2/m_W^2=8\lambda_H/g^2\). La normalización del potencial es
\(V(\mathsf H)=-\mu_H^2\mathsf H^*\mathsf H
+\lambda_H(\mathsf H^*\mathsf H)^2\). Estas composiciones conservan
su escala y su convención de árbol; los refinamientos que se cancelan
en una razón quedan declarados en la propia realización.

En una carta radial real \(\mathsf H^*\mathsf H=h^2/2\), el potencial
es \(V(h)=-\mu_H^2h^2/2+\lambda_Hh^4/4\). Para
\(\mu_H^2>0\), sus puntos críticos son \(0\) y
\(\pm v\), con \(v^2=\mu_H^2/\lambda_H\), y
\[
 V''(0)=-\mu_H^2,\qquad V''(\pm v)=2\mu_H^2.
\]
En unidades naturales de esta carta, el modo homogéneo normalizado posee
\(H_h=p_h^2/2+V(h)\). Su linealización en un punto crítico \(h_0\)
tiene generador
\[
 A_{h_0}=\begin{pmatrix}0&1\\-V''(h_0)&0\end{pmatrix},
 \qquad A_{h_0}^2=-V''(h_0)I.
\]
Por tanto, el origen es hiperbólico y los mínimos son elípticos; la
frontera \(V''=0\) tiene el tipo parabólico. Esta afirmación se deduce
de la derivada segunda del potencial y de una normalización modal
explícita. La energía material positiva del mínimo y la dirección
inestable de otro fondo pertenecen así a la misma función con distintos
puntos de expansión. El modo linealizado es un transporte particular de
los tres tipos cuadráticos; sus frecuencias y unidades proceden del
potencial, no de identificar su hessiano con la curvatura gaussiana de
\eqref{eq:hc-metrique}.

\subsection{Modo de curvatura de Yang--Mills y oscilación elíptica}

La realización de Yang--Mills construida en las secciones precedentes
posee un generador positivo \(D^{\rm YM}\) de dimensión inversa de
longitud. La escala
\[
 \kappa_{\rm av}=\frac{\mu_{\rm vol}}{\ell_0}
 \left(\frac{\pi_{\rm HMT}}{\varphi_{\rm HMT}^2}-1\right)>0,
 \qquad \delta=3\kappa_{\rm av}
\]
desciende de la incidencia de hojas y del lector regional, con la unidad
longitudinal \(\ell_0\) declarada. El vacío es simple y el testigo de
curvatura en la carta producto es
\[
 W(q)=\mathbf1_{\{q_1+\cdots+q_6=0\}}-\frac13,
 \quad \mathbb EW=0,\quad \|W\|^2=\frac29,
 \quad D^{\rm YM}W=\delta W.
\]
El transporte unitario hacia la medida de interacción lleva \(W\) a
\(W_g\) y conserva esas identidades. Las inclusiones de refinamiento
\(J_m\) entrelazan los generadores y conservan el testigo, de modo que
éste permanece en el límite físico. Aquí se utiliza el resultado completo
de las pruebas de vacío, brecha uniforme y reconstrucción, no un
autovalor aislado de una matriz finita~\cite{HMT-YM}.

Sea \(w=W_g/\sqrt{2/9}\) el vector unitario en la fibra física y
\(\hbar=\hbar_{\rm ret}>0\). El Hamiltoniano energético es
\(H_E=\hbar cD^{\rm YM}\). Sobre el espacio de Hilbert realificado
se adopta la forma \(\omega_{\mathcal H}(u,v)=2\hbar
\operatorname{Im}\langle u,v\rangle\), con producto interno antilineal
en el primer argumento. Definamos
\begin{equation}
 \mathfrak F(Q,P)=\frac{Q+iP}{\sqrt{2\hbar}}\,w.
 \label{eq:hc-ym-embedding}
\end{equation}

\begin{theorem}[Inmersión simpléctica del modo físico de curvatura]
\label{thm:hc-ym-oscillator}
La aplicación \(\mathfrak F:\mathbb R^2\to\mathcal H^{\rm phys}\)
es una inmersión real con imagen invariante por la evolución física y
satisface
\begin{equation}
 \mathfrak F^*\omega_{\mathcal H}=dQ\wedge dP,
 \qquad
 \langle\mathfrak F,H_E\mathfrak F\rangle
 =\frac{c\delta}{2}(Q^2+P^2).
 \label{eq:hc-ym-energy-pullback}
\end{equation}
La ecuación de Schrödinger restringida entrelaza exactamente el flujo
elíptico \(\dot Q=c\delta P,\ \dot P=-c\delta Q\).
El mapa es compatible con las inclusiones de refinamiento.
\end{theorem}
\begin{proof}
Para dos variaciones \((a,b),(a',b')\), el producto interno de sus
imágenes tiene parte imaginaria \((ab'-ba')/(2\hbar)\).
Esto prueba la primera identidad. La ecuación propia y \(\|w\|=1\)
dan la segunda. Puesto que
\(\dot\psi=-icD^{\rm YM}\psi\), el coeficiente \(Q+iP\)
satisface \(\dot Q+i\dot P=-ic\delta(Q+iP)\).
Finalmente, \(J_mw_m=w_{m+1}\) implica
\(J_m\mathfrak F_m=\mathfrak F_{m+1}\). Estas identidades pasan al
límite por las isometrías ya construidas.
\end{proof}

El plano completo describe amplitudes de un modo; su círculo
\(Q^2+P^2=2\hbar\) representa los vectores normalizados de ese
autoespacio complejo. Sobre tal círculo la fase global evoluciona y el
rayo cuántico permanece fijo. Superposiciones de autoespacios distintos
conservan, además, sus fases relativas. La inmersión mantiene esta
diferencia entre espacio de amplitudes y espacio proyectivo.

La energía de brecha y la masa de brecha son
\(\Delta_E=\hbar c\delta\) y \(m_{\rm gap}=\hbar\delta/c\).
La frecuencia modal \(c\delta\) tiene dimensión de tiempo inverso.
Las longitudes euclídeas de \(e^{-sD^{\rm YM}}\) y el tiempo físico
de \(e^{-itH_E/\hbar}\) se relacionan mediante la carta y la
reconstrucción, conservando sus distintos grupos y semigrupos.
La composición \(\mathfrak F\circ M_\eta^{-1}\) transporta esta
misma energía a la elipse de acción; si el marco varía, el término de
conexión de \eqref{eq:hc-htr} completa su evolución.

El enlace obtenido es una realización exacta sobre el modo invariante de
curvatura y sus prolongaciones. La positividad del Hamiltoniano cuántico
genera una rotación en su realificación simpléctica; la denominación
hiperbólica de una métrica espacial pertenece al tensor métrico, y el
transporte unipotente del registro pertenece a su memoria. En dimensión
finita, un grupo \(I+tN\) unitario para todo \(t\in\mathbb R\), con
\(N^2=0\), fuerza \(N=0\): derivar la unitariedad da \(N^*=-N\),
y entonces \(N^*N=-N^2=0\). El flujo parabólico clásico conserva su
forma simpléctica y posee su cuantización en espacios apropiados; esta
identidad sólo delimita la identificación con una matriz unitaria finita.

Así queda articulado un conjunto de transportes concretos. La memoria
produce una carga y una acción parabólica; el marco de la elipse produce
la conexión variacional y su lagrangiano; las respuestas constitutivas
determinan los coeficientes de Maxwell; la graduación de canales construye
Clifford y permite factorizar el operador material; la reconstrucción de
Yang--Mills aporta una inmersión modal simpléctica que conserva exactamente
su brecha y su refinamiento. Cada identificación dispone de una aplicación,
una forma conservada y una ecuación entrelazada.


% SOURCE sections/exact_reproduction.tex
\section{Comprobación exacta de las construcciones finitas}
\label{sec:exact-reproduction}

La reproducción distingue la demostración general y la evaluación de
una especialización. Las proposiciones precedentes utilizan identidades
algebraicas, positividad e hipótesis de convergencia explícitas. Sus
controles computacionales se realizan con enteros y fracciones racionales;
el desarrollo reticular conserva, además, su representación exacta en
la extensión algebraica declarada. Los datos del registro son salidas
publicadas por la construcción anterior. Las ejecuciones de esta
sección comprueban su transporte geométrico y operatorio, y conservan
por separado la procedencia de su generación.

El primer control recupera las separaciones a partir de los menores,
comprueba la relación de Plücker en las triangulaciones y compara
cambios de carta con deformaciones efectivas. Un cambio invertible
de las filas multiplica todos los menores por el mismo determinante;
los cocientes marcados permanecen iguales. Un reescalado independiente
de las columnas conserva ciertas razones dobles y altera, en general,
las separaciones reconstruidas. Evaluar ambas operaciones sobre los
mismos datos proporciona un control negativo de la equivalencia
que se afirma.

El segundo control subdivide las longitudes por nueve, conserva las
marcas heredadas y verifica las identidades de promedios y de Gram.
Los determinantes se calculan por eliminación exacta y se contrastan
con sus productos de Vandermonde. La positividad en cualquier
dimensión finita procede de la fórmula demostrada; los casos finitos
ejecutados comprueban la implementación de esa fórmula, sus índices
y sus normalizaciones. Las medidas de conteo y las medidas
cronológicas se mantienen separadas para que una igualdad de
momentos no dependa de cambiar silenciosamente la normalización.

El tercer control reproduce los ocho primeros momentos mediante
el operador de Jacobi de orden cuatro. La matriz simétrica se
representa de forma racional en la base de polinomios mónicos;
su métrica es \(\operatorname{diag}(h_0,\ldots,h_3)\). Se
comprueba la autoadjunción en esa métrica, la recurrencia de tres
términos y el acuerdo entre el resolvente matricial y su fracción
continua en \(z=2\). El resultado mantiene explícitos el vector
constante observado y la carga \(Q\), que intervienen en el
teorema de fidelidad espectral marcada.

El cuarto grupo verifica los cambios simplécticos y sus
Hamiltonianos. La derivada respecto del momento recupera la
velocidad; la transformación de Legendre se comprueba como
identidad polinómica o de Laurent. La inercia de una forma
cuadrática y la positividad de un Hamiltoniano cuántico
pertenecen a sus espacios respectivos. La representación
de Clifford se comprueba mediante sus anticonmutadores,
y la composición con el operador de masa se prueba por
el cuadrado del operador de Dirac. Estas igualdades
identifican cada término y sus productos cruzados.

Finalmente se conservan las comprobaciones anteriores de
triangulación, forma canónica, reducción de Schur, congruencia
de masas y construcción reticular. En Yang--Mills, las
proposiciones sobre medida común, refinamiento y reconstrucción
se exponen en el cuerpo con sus premisas y falsadores. Un
control finito de matrices no sustituye el paso al límite
ni la reconstrucción de campos: sus dominios se registran
separadamente en la entrega reproducible.

El paquete de fuentes enlaza cada programa con el enunciado
que evalúa y conserva su salida íntegra. Las condiciones se
comprueban mediante bifurcaciones explícitas, de modo que
permanecen activas al ejecutar el intérprete en modo
optimizado. Se reproducen asimismo en modo aislado, sin
cargar paquetes de usuario. La correspondencia entre
construcción, prueba y ejecución puede inspeccionarse
paso a paso sin convertir la salida de un programa en
una afirmación de alcance superior a la que comprueba.


% SOURCE foundation/libro/censo_firmas_emision.tex
\section{Censo de firmas del emisor residual}
\label{app:firmas-emision}

Este censo especifica las fibras finitas utilizadas en la sección~\ref{sec:extension}.
Cada condición inicial se identifica por
\[
 \nu=4(9i+j)+d,\qquad 0\leq i,j\leq8,\quad
 d=0,1,2,3\ \longleftrightarrow\ N,E,S,O.
\]
La división euclidiana recupera \(d=[\nu]_4\), \(j=[\lfloor\nu/4\rfloor]_9\)
e \(i=\lfloor\nu/36\rfloor\). Por tanto, \(0\leq\nu<324\) enumera
cada condición exactamente una vez, y las parejas de tales identificadores
enumeran las \(104\,976\) condiciones de \(\Omega_{\rm seed}\).

Sea \(f_+(\nu)\), respectivamente \(f_\times(\nu)\), la firma de seis
componentes obtenida por la recurrencia de emisión. En cada paso activo se lee
la casilla antes de avanzar; los vectores de dirección son
\[
 v_N=(-1,0),\quad v_E=(0,1),\quad v_S=(1,0),\quad v_O=(0,-1).
\]
Las coordenadas se reducen módulo nueve y el vector se invierte tras los
pasos \(27\) y \(54\). Los pasos de fase cero no desplazan ninguno de los
cursores. Estas reglas, junto con \eqref{eq:emisor-seis}, determinan
\(f_+\) y \(f_\times\) sin un catálogo de valores objetivo.

La tabla da todas sus imágenes. La multiplicidad \(m_\sigma\) es el número
de identificadores en la fibra \(f_\sigma^{-1}(a)\); \(\nu_{\min}\) es su
menor elemento. Cada palabra conserva los ceros iniciales. La fibra completa
queda determinada por la condición
\[
 f_\sigma^{-1}(a)=\{\nu\in\{0,\ldots,323\}:f_\sigma(\nu)=a\}.
\]
\begin{center}\small
\begin{tabular}{@{}lrr@{\qquad}lrr@{}}
\toprule
\(f_+\) & \(m_+\) & \(\nu_{\min}\) & \(f_\times\) & \(m_\times\) & \(\nu_{\min}\)\\
\midrule
\texttt{122564} & 18 & 17 & \texttt{000000} & 108 & 8 \\
\texttt{122898} & 18 & 24 & \texttt{177393} & 8 & 1 \\
\texttt{212645} & 18 & 5 & \texttt{177555} & 8 & 54 \\
\texttt{212988} & 18 & 0 & \texttt{177771} & 8 & 11 \\
\texttt{221456} & 18 & 29 & \texttt{339555} & 8 & 6 \\
\texttt{221889} & 18 & 12 & \texttt{339771} & 8 & 15 \\
\texttt{456221} & 18 & 28 & \texttt{339933} & 8 & 59 \\
\texttt{465898} & 18 & 21 & \texttt{393177} & 8 & 0 \\
\texttt{546988} & 18 & 9 & \texttt{393393} & 8 & 45 \\
\texttt{564122} & 18 & 16 & \texttt{393555} & 8 & 40 \\
\texttt{645212} & 18 & 4 & \texttt{555177} & 8 & 52 \\
\texttt{654889} & 18 & 33 & \texttt{555339} & 8 & 4 \\
\texttt{889221} & 18 & 13 & \texttt{555393} & 8 & 41 \\
\texttt{889654} & 18 & 32 & \texttt{555555} & 24 & 84 \\
\texttt{898122} & 18 & 25 & \texttt{555717} & 8 & 14 \\
\texttt{898465} & 18 & 20 & \texttt{555771} & 8 & 18 \\
\texttt{988212} & 18 & 1 & \texttt{555933} & 8 & 24 \\
\texttt{988546} & 18 & 8 & \texttt{717555} & 8 & 12 \\
 & &  & \texttt{717717} & 8 & 21 \\
 & &  & \texttt{717933} & 8 & 19 \\
 & &  & \texttt{771177} & 8 & 9 \\
 & &  & \texttt{771339} & 8 & 13 \\
 & &  & \texttt{771555} & 8 & 16 \\
 & &  & \texttt{933339} & 8 & 57 \\
 & &  & \texttt{933555} & 8 & 26 \\
 & &  & \texttt{933717} & 8 & 17 \\
\bottomrule
\end{tabular}
\end{center}

La evaluación de la recurrencia para los \(324\) identificadores produce
las dos particiones de la tabla: \(18\cdot18=324\) y
\(24\cdot8+24+108=324\). El producto de sus fibras da
\(432\) fibras de tamaño \(144\), \(18\) de tamaño \(432\) y \(18\)
de tamaño \(1944\), por la inyectividad de \eqref{eq:acoplamiento}.
Finalmente, aplicar \eqref{eq:psi-catalogo} a ese producto de imágenes
da \(243\) palabras distintas. El censo es finito y pertenece al emisor
residual; el refinamiento coinductivo posterior conserva su propia ley.



% SOURCE sections/conclusions.tex
\clearpage
\section{Conclusiones}
\label{sec:conclusions-final}

La tesis geométrica queda establecida mediante reconstrucción y transporte:
una publicación positiva de la estructura discreta del continuo conserva,
con sus marcas, el registro del que procede y los lectores que éste
determina. La carga total y los cocientes de menores recuperan cada
coordenada; las dos memorias nonádicas proporcionan otra reconstrucción
operatoria. El valor y la forma se articulan así por inversas explícitas,
después de la generación APP--TRIT--TPK y de sus prolongaciones compatibles.
La doble proyección mantiene unidas la realización aritmética y la
incidencia excepcional dentro de esa genealogía.

El refinamiento supera la limitación de una matriz con trece columnas.
Para cualquier rango y dimensión externa finitos, un nivel nonádico
suficiente produce los menores positivos y la imagen amplituhedral
correspondiente. Los nodos heredados conservan sus banderas; los
promedios celulares conservan la medida y registran el detalle
ortogonal. En rango uno, la cobertura simplicial y la recursión de
residuos se extienden a toda dimensión finita. En rangos mayores,
el resultado determina familias positivas compatibles con rango y
límite probados; su amplitud dimensional queda separada de la
clasificación global de todas las celdas amplituhedrales.

La conservación admite tres expresiones complementarias. La forma
canónica se suma bajo subdivisión de un soporte fijo; la energía
efectiva se conserva mediante reducción de Schur; y la medida
cronológica determina momentos cuyo incremento se descompone en
detalles recuperables. Estos momentos construyen operadores de
Jacobi cuyas compresiones finitas tienen espectro simple y residuos positivos. Se obtiene así
un recorrido efectivo desde las separaciones aritméticas hasta una
representación espectral, con fracción continua determinada por la
misma medida. Las mutaciones que cambian únicamente la carta
conservan los lectores del punto marcado.

La realización física gana precisión mediante sus operadores. Los
Hamiltonianos de los tres regímenes tríticos poseen formas
simplécticas y transformaciones de Legendre explícitas. La carta
elíptica variable incorpora su término de conexión; la memoria
parabólica conserva su invariante; la apertura hiperbólica tiene
su evolución propia. Los lectores de acción y respuesta
constitutiva se transportan por la inversa grassmanniana. Con
las publicaciones regionales y las marcas enteras declaradas,
la respuesta constitutiva etiquetada recupera a su vez el registro.
La condición de igualdad de trazas gobierna el pegado de formas
con valores en la recta de acción.

La deformación excepcional conserva el volumen y la simetría
ternaria, y convierte la inversión del parámetro en dualidad
reticular. Su reciprocidad theta y su realización de signatura
partida enlazan la variación geométrica con la dualidad de la
red. En Yang--Mills, la escala areal--volumétrica, la medida común,
la coercividad uniforme y el testigo persistente sitúan la brecha
junto al Hamiltoniano que la realiza. El desarrollo conserva
el espacio físico y los entrelazadores necesarios para transportar
esa afirmación a través del refinamiento.

El resultado conjunto es una arquitectura de formas con memoria:
la dimensión puede aumentar, la representación puede cambiar y
la información pertinente permanece recuperable mediante mapas
determinados. Éste es el contenido diferencial del artículo.
La geometría positiva, la dualidad reticular y la realización
espectral quedan enlazadas por construcciones que especifican
qué generan, qué conservan y cómo se reconstruyen desde un mismo
origen discreto.
\clearpage


% SOURCE sections/bibliography_additions.tex

\phantomsection
\addcontentsline{toc}{section}{Referencias}
\bibitem{feynman1948} R. P. Feynman, ``Space-Time Approach to Non-Relativistic Quantum Mechanics'', \emph{Reviews of Modern Physics} \textbf{20}, 367--387 (1948). \href{https://doi.org/10.1103/RevModPhys.20.367}{doi:10.1103/RevModPhys.20.367}.
\bibitem{feynman1949} R. P. Feynman, ``The Theory of Positrons'', \emph{Physical Review} \textbf{76}, 749--759 (1949). \href{https://doi.org/10.1103/PhysRev.76.749}{doi:10.1103/PhysRev.76.749}.
\bibitem{feynman1965} R. P. Feynman, ``The Development of the Space-Time View of Quantum Electrodynamics'', conferencia Nobel, 11 de diciembre de 1965. \href{https://www.nobelprize.org/prizes/physics/1965/feynman/lecture/}{Texto de la conferencia}.
\bibitem{mustovicary} B. Musto y J. Vicary, ``Quantum Latin Squares and Unitary Error Bases'', \emph{Quantum Information and Computation} \textbf{16}, 1318--1332 (2016). \href{https://doi.org/10.26421/QIC16.15-16-4}{doi:10.26421/QIC16.15-16-4}.
\bibitem{delascuevas2020} G. De las Cuevas, T. Drescher y T. Netzer, ``Quantum Magic Squares: Dilations and Their Limitations'', \emph{Journal of Mathematical Physics} \textbf{61}, 111704 (2020). \href{https://doi.org/10.1063/5.0022344}{doi:10.1063/5.0022344}.
\bibitem{delascuevas2023} G. De las Cuevas, T. Netzer e I. Valentiner-Branth, ``Magic Squares: Latin, Semiclassical, and Quantum'', \emph{Journal of Mathematical Physics} \textbf{64}, 022201 (2023). \href{https://doi.org/10.1063/5.0127393}{doi:10.1063/5.0127393}.
\bibitem{dirac1933} P. A. M. Dirac, ``The Lagrangian in Quantum Mechanics'', \emph{Physikalische Zeitschrift der Sowjetunion} \textbf{3}, 64--72 (1933). \href{https://www.informationphilosopher.com/solutions/scientists/dirac/Lagrangian_1933.pdf}{Reproducción del artículo original}.
\bibitem{pauli1946} W. Pauli, ``Exclusion Principle and Quantum Mechanics'', conferencia Nobel, 13 de diciembre de 1946, en \emph{Nobel Lectures, Physics 1942--1962}, Elsevier, 1964, pp. 27--43, especialmente p. 42. \href{https://www.nobelprize.org/uploads/2018/06/pauli-lecture.pdf}{Texto de la conferencia}.
\bibitem{codata2018} E. Tiesinga, P. J. Mohr, D. B. Newell y B. N. Taylor, ``CODATA Recommended Values of the Fundamental Physical Constants: 2018'', \emph{Reviews of Modern Physics} \textbf{93}, 025010 (2021). \href{https://doi.org/10.1103/RevModPhys.93.025010}{doi:10.1103/RevModPhys.93.025010}. \href{https://physics.nist.gov/cuu/Constants/ArchiveASCII/allascii_2018.txt}{Tablas archivadas de 2018}.
\bibitem{codata2022} P. J. Mohr, D. B. Newell, B. N. Taylor y E. Tiesinga, ``CODATA Recommended Values of the Fundamental Physical Constants: 2022'', \emph{Reviews of Modern Physics} \textbf{97}, 025002 (2025). \href{https://doi.org/10.1103/RevModPhys.97.025002}{doi:10.1103/RevModPhys.97.025002}. \href{https://physics.nist.gov/cuu/Constants/Table/allascii.txt}{Tabla de constantes}.
\bibitem{flm} I. B. Frenkel, J. Lepowsky y A. Meurman, \emph{Vertex Operator Algebras and the Monster}, Pure and Applied Mathematics, vol. 134, Academic Press, 1988.
\bibitem{borcherds} R. E. Borcherds, ``Monstrous Moonshine and Monstrous Lie Superalgebras'', \emph{Inventiones Mathematicae} \textbf{109}, 405--444 (1992). \href{https://doi.org/10.1007/BF01232032}{doi:10.1007/BF01232032}.
\bibitem{conwaysloane} J. H. Conway y N. J. A. Sloane, \emph{Sphere Packings, Lattices and Groups}, 3.\textsuperscript{a} ed., Grundlehren der mathematischen Wissenschaften 290, Springer, 1999. \href{https://doi.org/10.1007/978-1-4757-6568-7}{doi:10.1007/978-1-4757-6568-7}.
\bibitem{bipm2018} Conférence générale des poids et mesures, \emph{Resolution 1 of the 26th CGPM: On the revision of the International System of Units}, 2018. \href{https://www.bipm.org/en/committees/cg/cgpm/26-2018/resolution-1}{Resolución oficial}. \href{https://doi.org/10.59161/CGPM2018RES1E}{doi:10.59161/CGPM2018RES1E}.
\bibitem{dirac1931} P. A. M. Dirac, ``Quantised singularities in the electromagnetic field'', \emph{Proceedings of the Royal Society of London. Series A} \textbf{133}, 60--72 (1931). \href{https://doi.org/10.1098/rspa.1931.0130}{doi:10.1098/rspa.1931.0130}.
\bibitem{lep2006} ALEPH, DELPHI, L3, OPAL y SLD Collaborations; LEP Electroweak Working Group; SLD Electroweak Group y SLD Heavy Flavour Group, ``Precision electroweak measurements on the Z resonance'', \emph{Physics Reports} \textbf{427}, 257--454 (2006). \href{https://doi.org/10.1016/j.physrep.2005.12.006}{doi:10.1016/j.physrep.2005.12.006}.
\bibitem{chenlamshimakura2016} H.-Y. Chen, C. H. Lam y H. Shimakura, ``\(\mathbb Z_3\)-orbifold construction of the Moonshine vertex operator algebra and some maximal \(3\)-local subgroups of the Monster'', prepublicación, arXiv:1606.05961 (2016), versión 2 (2017). \href{https://arxiv.org/abs/1606.05961}{arXiv:1606.05961}.
\bibitem{abelamyamada2017} T. Abe, C. H. Lam y H. Yamada, ``A remark on \(\mathbb Z_p\)-orbifold constructions of the Moonshine vertex operator algebra'', prepublicación, arXiv:1705.09022 (2017), versión 4. \href{https://arxiv.org/abs/1705.09022v4}{arXiv:1705.09022}.
\bibitem{kmconwaynorton1979} J. H. Conway y S. P. Norton, ``Monstrous Moonshine'', \emph{Bulletin of the London Mathematical Society} \textbf{11}, 308--339 (1979). \href{https://doi.org/10.1112/blms/11.3.308}{doi:10.1112/blms/11.3.308}.
\bibitem{bennett2003} C. H. Bennett, ``Notes on Landauer's principle, reversible computation, and Maxwell's Demon'', \emph{Studies in History and Philosophy of Modern Physics} \textbf{34}, 501--510 (2003). \href{https://doi.org/10.1016/S1355-2198(03)00039-X}{doi:10.1016/S1355-2198(03)00039-X}.
